\documentclass[10pt, letterpaper]{article}

% Inhaltsverzeichnis für Pakettypen (nur für Übersicht im Header, wird nicht im Dokument angezeigt)
% 1. Seitenlayout und Ränder
% 2. Sprache und Zeichensatz
% 3. Mathematik und Theorem-Umgebungen
% 4. Eigene Makros
% 5. Diagramme und Grafiken
% 6. Tabellen und Aufzählungen
% 7. Inhaltsverzeichnis
% 8. Abschnittsüberschriften
% 9. Abstrakt-Umgebung
% 10. Todos/Notizen
% 11. Rahmen/Box-Umgebungen
% 12. Python-Integration
% 13. Literaturverwaltung
% 14. Hyperlinks
% 15. Absatzeinstellungen
% 16. Umgebungen
% 17  Graphik
% 18  Extra
% 00. Titel und Autor

\usepackage{slashed}

% --- 1. Seitenlayout und Ränder ---
\usepackage[margin=3cm]{geometry}

% --- 2. Sprache und Zeichensatz ---
\usepackage[english]{babel}
\usepackage[T1]{fontenc}
\usepackage[utf8]{inputenc}

% --- 3. Mathematik und Theorem-Umgebungen ---
\usepackage{amsmath, amssymb, amsthm}
\usepackage{mathrsfs}
\DeclareMathOperator{\WF}{WF}

% --- 4. Eigene Makros ---
\usepackage{xcolor}
\newcommand{\SKP}{\langle\cdot,\cdot\rangle}
\newcommand{\R}{\mathbb{R}}
\newcommand{\N}{\mathbb{N}}
\newcommand{\Q}{\mathbb{Q}}
\newcommand{\Z}{\mathbb{Z}}
\newcommand{\C}{\mathbb{C}}
\newcommand{\entwurf}[1]{\textcolor{red}{#1}}
\newcommand{\GL}{\mathrm{GL}}
\newcommand{\SL}{\mathrm{SL}}
\newcommand{\SO}{\mathrm{SO}}
\newcommand{\SU}{\mathrm{SU}}
\newcommand{\Sp}{\mathrm{Sp}}
\newcommand{\Ogroup}{\mathrm{O}}
\newcommand{\U}{\mathrm{U}}

% --- 5. Diagramme und Grafiken ---
\usepackage{graphicx}
\graphicspath{{../../Library/Images/}}
\usepackage[export]{adjustbox}
\usepackage{tikz}
\usetikzlibrary{decorations.pathreplacing, arrows.meta, positioning}
\usepackage{tikz-cd}

% --- 6. Tabellen und Aufzählungen ---
\usepackage{enumitem}
\setlist[itemize]{left=0.5cm}

\newenvironment{romanenum}[1][]
  {%
    \ifx&#1&
    \else
      \textbf{#1}\quad
    \fi
    \begin{enumerate}[label=\roman*)]
  }
  {%
    \end{enumerate}%
  }

% --- 7. Inhaltsverzeichnis ---
\usepackage{tocloft}
\renewcommand{\cftsecfont}{\footnotesize}
\renewcommand{\cftsubsecfont}{\footnotesize}
\renewcommand{\cftsubsubsecfont}{\footnotesize}
\renewcommand{\cftsecpagefont}{\footnotesize}
\renewcommand{\cftsubsecpagefont}{\footnotesize}
\renewcommand{\cftsubsubsecpagefont}{\footnotesize}
\usepackage{etoc}

% --- 8. Abschnittsüberschriften ---
\usepackage{titlesec}
\titleformat{\section}{\normalfont\large\bfseries}{\thesection}{1em}{}
\titleformat{\subsection}{\normalfont\normalsize\bfseries}{\thesubsection}{0.5em}{}
\titleformat{\subsubsection}{\normalfont\normalsize\bfseries}{\thesubsubsection}{0.5em}{}
\setcounter{secnumdepth}{4}

% --- 9. Abstrakt-Umgebung ---
\usepackage{changepage}
\renewenvironment{abstract}
  {
    \begin{adjustwidth}{1.5cm}{1.5cm}
    \small
    \textsc{Abstract. –}%
  }
  {
    \end{adjustwidth}
  }

% --- 10. Todos/Notizen ---
\usepackage{todonotes}
% Hellblau definieren
\definecolor{hellblau}{rgb}{0.3,0.6,1}
\newenvironment{note}{\par\begingroup\color{hellblau}\itshape}{\par\endgroup}

% --- 11. Rahmen/Box-Umgebungen ---
\usepackage{mdframed}
\usepackage{tcolorbox}
\colorlet{shadecolor}{gray!25}

\newenvironment{customTheorem}
  {\vspace{10pt}%
   \begin{mdframed}[
     backgroundcolor=gray!20,
     linewidth=0pt,
     innertopmargin=10pt,
     innerbottommargin=10pt,
     skipabove=\dimexpr\topsep+\ht\strutbox\relax,
     skipbelow=\topsep,
   ]}
  {\end{mdframed}
   \vspace{10pt}%
  }

% --- 12. Python-Integration ---
% (Deaktiviert in dieser Version, aktiviere bei Bedarf)
% \usepackage{pythontex}
% \usepackage[makestderr]{pythontex}

% --- 13. Literaturverwaltung ---
\usepackage{csquotes}
\usepackage[backend=biber, style=alphabetic, citestyle=alphabetic]{biblatex}
\addbibresource{bibliography.bib}

% --- 14. Hyperlinks ---
\usepackage{hyperref}
\hypersetup{
  colorlinks   = true,
  urlcolor     = blue,
  linkcolor    = blue,
  citecolor    = blue,
  frenchlinks  = true
}

% --- 15. Absatzeinstellungen ---
\usepackage[parfill]{parskip}
\sloppy

% --- 16. Umgebungen ---
\usepackage{thmtools}

\newcommand{\CustomHeading}[3]{%
  \par\medskip\noindent%
  \textbf{#1 #2} \textnormal{(#3)}.\enskip%
}

\newenvironment{DEF}[2]{\begin{unitbox}\CustomHeading{Definition}{#1}{#2}}{\end{unitbox}}
\newenvironment{PROP}[2]{\begin{unitbox}\CustomHeading{Proposition}{#1}{#2}}{\end{unitbox}}
\newenvironment{THEO}[2]{\begin{unitbox}\CustomHeading{Theorem}{#1}{#2}}{\end{unitbox}}
\newenvironment{LEM}[2]{\begin{unitbox}\CustomHeading{Lemma}{#1}{#2}}{\end{unitbox}}
\newenvironment{KORO}[2]{\begin{unitbox}\CustomHeading{Corollar}{#1}{#2}}{\end{unitbox}}
\newenvironment{REM}[2]{\begin{unitbox}\CustomHeading{Remark}{#1}{#2}}{\end{unitbox}}
\newenvironment{EXA}[2]{\begin{unitbox}\CustomHeading{Example}{#1}{#2}}{\end{unitbox}}
\newenvironment{STUD}[2]{\begin{unitbox}\CustomHeading{Study}{#1}{#2}}{\end{unitbox}}
\newenvironment{CONC}[2]{\begin{unitbox}\CustomHeading{Concept}{#1}{#2}}{\end{unitbox}}
\newenvironment{OTH}[2]{\begin{unitbox}\CustomHeading{Other}{#1}{#2}}{\end{unitbox}}
\newenvironment{EXE}[2]{\begin{unitbox}\CustomHeading{Exercise}{#1}{#2}}{\end{unitbox}}
\newenvironment{MOT}[2]{\begin{unitbox}\CustomHeading{Motivation}{#1}{#2}}{\end{unitbox}}
\newenvironment{PROOF}[2]{\begin{unitbox}\CustomHeading{Proof}{#1}{#2}}{\end{unitbox}}

% --- Unit Umgebung für Source-Inhalte ---
\usepackage{mdframed}
\newmdenv[
  linewidth=1pt,
  topline=false,
  bottomline=false,
  rightline=false,
  leftmargin=0cm,
  rightmargin=0cm,
  skipabove=10pt,
  skipbelow=10pt,
  innertopmargin=0.5\baselineskip,
  innerbottommargin=0.5\baselineskip,
  backgroundcolor=gray!10,
  linecolor=gray
]{unitbox}

\newenvironment{unit}[1]
  {\begin{unitbox}\textbf{Unit #1}\par\smallskip}
  {\end{unitbox}}

% --- 17. ... ---

% --- 18. Extras ---
\usepackage{stmaryrd}
\usepackage{bbold}  % falls du \mathbb{1} nutzen willst

% --- 00. Titel und Autor ---
\title{Mein Titel}
\author{Tim Jaschik}
\date{\today}

\begin{document}

\maketitle
\rule{\textwidth}{0.5pt}
\begin{abstract}
Kurze Beschreibung …
\end{abstract}
\rule{\textwidth}{0.5pt}
\vspace{0.5cm}

\tableofcontents

\pagebreak

\section*{Part II}
Lie Algebras

\section*{Elementary Structure Theory of Lie Algebras}
Lie algebras form the infinitesimal counterpart of Lie groups. We have already seen in the preceding chapters how matrix groups give rise to Lie algebras of matrices, and we shall also see in Chapter 9 below how to associate a Lie algebra to any Lie group. This correspondence is the guiding motivation behind the theory of finite-dimensional Lie algebras to which we now turn in some depth.

In this part, we study Lie algebras as independent objects and thus provide tools to solve the linear algebra problems which we will encounter when translating Lie group questions into Lie algebra questions in Part IV. We start in the present chapter by working out the standard analysis of an algebraic structure for Lie algebras: What are the substructures? Under which condition does a substructure lead to a quotient structure? What are the simple structures? Does one have composition series? This leads to concepts like Lie subalgebras and ideals, nilpotent, solvable, and semisimple Lie algebras. Key results in this context are Engel's Theorems on nilpotent Lie algebras, Lie's Theorem for solvable Lie algebras, and Cartan's criteria for solvability and semisimplicity. The latter are first instances in which one recognizes the usefulness of the Cartan-Killing form, which is a specific structural element for Lie algebras.

\subsection*{Basic Concepts}
In this section, we provide the basic definitions and concepts concerning Lie algebras. In particular, we discuss ideals, quotients, homomorphisms, and the elementary connections between these concepts.

\subsubsection*{Definitions and Examples}
We start by recalling the definition of a Lie algebra from Definition 4.1.1.\\
Definition 5.1.1. Let $\mathfrak{g}$ be a vector space. A Lie bracket on $\mathfrak{g}$ is a bilinear map $[\cdot, \cdot]: \mathfrak{g} \times \mathfrak{g} \rightarrow \mathfrak{g}$ satisfying\\
(L1) $[x, y]=-[y, x]$ for $x, y \in \mathfrak{g}$,\\
(L2) $[x,[y, z]]+[y,[z, x]]+[z,[x, y]]=0$ for $x, y, z \in \mathfrak{g}$ (Jacobi identity).\\
For any Lie bracket on $\mathfrak{g}$, the pair ( $\mathfrak{g},[\cdot, \cdot]$ ) is called a Lie algebra.\\
Example 5.1.2. (cf. Definition 3.1.1) A vector space $\mathcal{A}$ together with a bilinear map $: \mathcal{A} \times \mathcal{A} \rightarrow \mathcal{A}$ is called an (associative) algebra, if

$$
a \cdot(b \cdot c)=(a \cdot b) \cdot c \quad \text { for } \quad a, b, c \in \mathcal{A}
$$

Then the commutator

$$
[a, b]:=a \cdot b-b \cdot a
$$

defines a Lie bracket on $\mathcal{A}$ (Exercise). We write $\mathcal{A}_{L}:=(\mathcal{A},[\cdot, \cdot])$ for this Lie algebra.

Example 5.1.3. (a) Let $V$ be a vector space and $\operatorname{End}(V)$ be the set of linear endomorphisms of $V$. Then $\operatorname{End}(V)$ is an associative algebra, and we write $\mathfrak{g} \mathfrak{l}(V):=\operatorname{End}(V)_{L}$ for the corresponding Lie algebra.\\
(b) The space $M_{n}(\mathbb{K})$ of ( $n \times n$ )-matrices with entries in $\mathbb{K}$ is an associative algebra with respect to matrix multiplication. We write $\mathfrak{g l}{ }_{n}(\mathbb{K}):=M_{n}(\mathbb{K})_{L}$ for the corresponding Lie algebra.

Definition 5.1.4. (a) Let $\mathfrak{g}$ and $\mathfrak{h}$ be Lie algebras. A linear map $\alpha: \mathfrak{g} \rightarrow \mathfrak{h}$ is called a homomorphism if

$$
\alpha([x, y])=[\alpha(x), \alpha(y)] \quad \text { for } \quad x, y \in \mathfrak{g} .
$$

An isomorphism of Lie algebras is a homomorphism $\alpha$ for which there exists a homomorphism $\beta: \mathfrak{h} \rightarrow \mathfrak{g}$ with $\alpha \circ \beta=\operatorname{id}_{\mathfrak{h}}$ and $\beta \circ \alpha=\operatorname{id}_{\mathfrak{g}}$. It is easy to see that this condition is equivalent to $\alpha$ being bijective (Exercise).

A representation of a Lie algebra $\mathfrak{g}$ on the vector space $V$ is a homomorphism $\alpha: \mathfrak{g} \rightarrow \mathfrak{g} \mathfrak{l}(V)$. We also write ( $\alpha, V$ ) for a representation $\alpha$ of $\mathfrak{g}$ on $V$.\\
(b) Let $\mathfrak{g}$ be a Lie algebra and $U, V$ be subsets of $\mathfrak{g}$. We write

$$
[U, V]:=\operatorname{span}\{[u, v]: u \in U, v \in V\}
$$

for the smallest subspace containing all brackets $[u, v]$ with $u \in U$ and $v \in V$.\\
(c) A linear subspace $\mathfrak{h}$ of $\mathfrak{g}$ is called a Lie subalgebra if $[\mathfrak{h}, \mathfrak{h}] \subseteq \mathfrak{h}$. We then write $\mathfrak{h}<\mathfrak{g}$. If we even have $[\mathfrak{g}, \mathfrak{h}] \subseteq \mathfrak{h}$, we call $\mathfrak{h}$ an ideal of $\mathfrak{g}$ and write $\mathfrak{h} \unlhd \mathfrak{g}$.\\
(d) The Lie algebra $\mathfrak{g}$ is called abelian if $[\mathfrak{g}, \mathfrak{g}]=\{0\}$, which means that all brackets vanish.

Remark 5.1.5. From the definitions it is immediately clear that the image of a homomorphism $\alpha: \mathfrak{g}_{1} \rightarrow \mathfrak{g}_{2}$ of Lie algebras is a subalgebra of $\mathfrak{g}_{2}$. Moreover, $\alpha^{-1}(\mathfrak{h})$ is an ideal in $\mathfrak{g}_{1}$ if $\mathfrak{h} \unlhd \mathfrak{g}_{2}$, and $\alpha^{-1}(\mathfrak{h})$ is a subalgebra if $\mathfrak{h}<\mathfrak{g}_{2}$. In particular, the kernel ker $\alpha$ of a Lie algebra homomorphism $\alpha$ is always an ideal.

Examples 5.1.6. (i) Let $\mathfrak{g}$ be Lie algebra. Then the center

$$
\mathfrak{z}(\mathfrak{g}):=\{x \in \mathfrak{g} \mid(\forall y \in \mathfrak{g})[x, y]=0\}
$$

of $\mathfrak{g}$ is an ideal in $\mathfrak{g}$.\\
(ii) For each Lie algebra $\mathfrak{g}$, the subspace $[\mathfrak{g}, \mathfrak{g}]$ is an ideal, called the commutator algebra of $\mathfrak{g}$.\\
(iii) Every one-dimensional subspace of a Lie algebra is a subalgebra since the Lie bracket is skew-symmetric.\\
(iv) The set

$$
\mathfrak{o}_{n}(\mathbb{K}):=\left\{x \in \mathfrak{g} \mathfrak{l}_{n}(\mathbb{K}) \mid x=-x^{\top}\right\}
$$

is a subalgebra of $\mathfrak{g l}(\mathbb{K})$, the orthogonal Lie algebra.\\
(v) The set

$$
\mathfrak{u}_{n}(\mathbb{C}):=\left\{x \in \mathfrak{g} \mathfrak{l}_{n}(\mathbb{C}) \mid x=-x^{*}\right\}
$$

is a real subalgebra of $\mathfrak{g l}(\mathbb{C})$, the unitary Lie algebra.\\
(vi) The set

$$
\mathfrak{s} \mathfrak{l}_{n}(\mathbb{K}):=\left\{x \in \mathfrak{g} \mathfrak{l}_{n}(\mathbb{K}) \mid \operatorname{tr}(x)=0\right\}
$$

is an ideal in $\mathfrak{g l}{ }_{n}(\mathbb{K})$, where $\operatorname{tr}(x)$ denotes the trace of $X$. It is called the special linear Lie algebra.\\
(vii) Let $J_{n}:=\left(\begin{array}{cc}\mathbf{0} & \mathbf{1}_{n} \\ -\mathbf{1}_{n} & \mathbf{0}\end{array}\right) \in \mathfrak{g l}_{2 n}(\mathbb{K})$ and note that $J_{n}^{\top}=-J_{n}$. Then the set

$$
\mathfrak{s} \mathfrak{p}_{n}(\mathbb{K}):=\left\{x \in \mathfrak{g} \mathfrak{l}_{2 n}(\mathbb{K}): x^{\top} J_{n}+J_{n} x=0\right\}
$$

is a Lie subalgebra of $\mathfrak{g} \mathfrak{l}_{2 n}(\mathbb{K})$, called the symplectic Lie algebra. Writing elements of $\mathfrak{g l}_{2 n}(\mathbb{K})$ as ( $2 \times 2$ )-block matrices, one easily verifies that

$$
\left(\begin{array}{ll}
A & B \\
C & D
\end{array}\right) \in \mathfrak{s p}_{n}(\mathbb{K}) \quad \Longleftrightarrow \quad B=B^{\top}, C=C^{\top}, A^{\top}=-D .
$$

(viii) The sets

$$
\mathfrak{n}=\left\{x=\left(x_{i j}\right) \in \mathfrak{g} \mathfrak{l}_{n}(\mathbb{K}) \mid(\forall i \geq j) x_{i j}=0\right\}
$$

and

$$
\mathfrak{b}=\left\{x=\left(x_{i j}\right) \in \mathfrak{g} \mathfrak{l}_{n}(\mathbb{K}) \mid(\forall i>j) x_{i j}=0\right\}
$$

are subalgebras of $\mathfrak{g l}{ }_{n}(\mathbb{K})$.\\
(ix) Let $V$ be a subspace of a Lie algebra $\mathfrak{g}$. The normalizer

$$
\mathfrak{n}_{\mathfrak{g}}(V)=\{x \in \mathfrak{g} \mid[x, V] \subseteq V\}
$$

of $V$ in $\mathfrak{g}$ is a subalgebra of $\mathfrak{g}$.

Example 5.1.7. Let $V$ be a vector space. A tuple $\mathcal{F}=\left(V_{0}, \ldots, V_{n}\right)$ of subspaces with

$$
\{0\}=V_{0} \subseteq V_{1} \subseteq \cdots \subseteq V_{n}=V
$$

is called a flag in $V$. Then

$$
\mathfrak{g}(\mathcal{F}):=\left\{x \in \mathfrak{g} \mathfrak{l}(V):(\forall j) x V_{j} \subseteq V_{j}\right\}
$$

is a Lie subalgebra of $\mathfrak{g l}(V)=\operatorname{End}(V)_{L}$.\\
To visualize this Lie algebra, we shall describe linear maps by suitable block matrices. If $V$ is a vector space which is a direct sum $V=W_{1} \oplus \cdots \oplus W_{n}$ of subspaces $W_{j}, j=1, \ldots, n$, then we write an endomorphism $A \in \operatorname{End}(V)$ as an ( $n \times n$ )-block matrix

$$
A=\left(A_{j k}\right)_{j, k=1, \ldots, n}=\left(\begin{array}{ccc}
A_{11} & \cdots & A_{1 n} \\
A_{21} & \cdots & A_{2 n} \\
\vdots & \ddots & \vdots \\
A_{n 1} & \cdots & A_{n n}
\end{array}\right) \text {, }
$$

where $A_{j k} \in \operatorname{Hom}\left(W_{k}, W_{j}\right)$ is uniquely determined by the requirement that the image of $v=\left(v_{1}, \ldots, v_{n}\right) \in V$ is

$$
A v=\left(\sum_{k=1}^{n} A_{j k} v_{k}\right)_{j=1, \ldots, n}
$$

Applying this kind of visualization to the Lie algebra $\mathfrak{g}(\mathcal{F})$, we choose in $V_{j}$ a subspace $W_{j}$ with $V_{j} \cong V_{j-1} \oplus W_{j}$. For each $j$, we then have $V_{j} \cong W_{1} \oplus \cdots \oplus W_{j}$, and in particular $V \cong W_{1} \oplus \cdots \oplus W_{n}$. Now the elements of $\mathfrak{g}(\mathcal{F})$ are those endomorphisms of $V$ corresponding to upper triangular matrices

$$
A=\left(\begin{array}{cccc}
A_{11} & A_{12} & \cdots & A_{1 n} \\
0 & A_{22} & \cdots & A_{2 n} \\
\vdots & \ddots & \ddots & \vdots \\
0 & \cdots & 0 & A_{n n}
\end{array}\right)
$$

Definition 5.1.8. Let $\mathfrak{g}$ be a Lie algebra. A linear map $\delta: \mathfrak{g} \rightarrow \mathfrak{g}$ is called a derivation if

$$
\delta([x, y])=[\delta(x), y]+[x, \delta(y)] \quad \text { for } \quad x, y \in \mathfrak{g} .
$$

The set of all derivations is denoted by $\operatorname{der}(\mathfrak{g})$.\\
Definition 5.1.9. Let $\mathfrak{g}$ be a Lie algebra and $X \in \mathfrak{g}$. Then the Jacobi identity implies that the linear map

$$
\operatorname{ad}(x): \mathfrak{g} \rightarrow \mathfrak{g}, \quad y \mapsto[x, y]
$$

is a derivation. Derivations of this form are called inner derivations. The map ad: $\mathfrak{g} \rightarrow \mathfrak{g} \mathfrak{l}(\mathfrak{g})$ is called the adjoint representation. That it is a representation follows directly from the Jacobi identity, which implies that

$$
\operatorname{ad}[x, y]=[\operatorname{ad} x, \operatorname{ad} y] .
$$

Proposition 5.1.10. For any Lie algebra $\mathfrak{g}$,\\
(i) $\operatorname{der}(\mathfrak{g})<\mathfrak{g} \mathfrak{l}(\mathfrak{g})$ and $\operatorname{ad}(\mathfrak{g}) \unlhd \operatorname{der}(\mathfrak{g})$ is an ideal. In particular,

$$
[D, \operatorname{ad} x]=\operatorname{ad}(D x) \quad \text { for } \quad D \in \operatorname{der}(\mathfrak{g}), x \in \mathfrak{g} .
$$

(ii) $\operatorname{ker}(\mathrm{ad})=\mathfrak{z}(\mathfrak{g})$.

Proof. (i) The first part is a special case of Exercise 5.1.4; and for the second, one verifies (5.1) by direct calculation.\\
(ii) is trivial. \(\square\)

\subsubsection*{Representations and Modules}
In this short subsection, we introduce some terminology concerning representations of Lie algebras and the corresponding concept of a Lie algebra module.

Definition 5.1.11. Let $\mathfrak{g}$ be a Lie algebra, and $V$ a vector space. Suppose that

$$
\mathfrak{g} \times V \rightarrow V, \quad(x, v) \mapsto x \cdot v
$$

is a bilinear map. If

$$
[x, y] \cdot v=x \cdot(y \cdot v)-y \cdot(x \cdot v) \quad \text { for } \quad x, y \in \mathfrak{g}, v \in V \text {, }
$$

then $V$ is called a $\mathfrak{g}$-module.\\
Definition 5.1.12. (a) Let $\mathfrak{g}$ be a Lie algebra and $V$ a $\mathfrak{g}$-module. A subspace $W \subseteq V$ is called a $\mathfrak{g}$-submodule if $\mathfrak{g} \cdot W \subseteq W$.\\
(b) A $\mathfrak{g}$-module $V$ is called simple if it is nonzero and there are no submodules except $\{0\}$ and $V$. It is called semisimple, if $V$ is the direct sum of simple submodules.\\
(c) If $V$ and $W$ are $\mathfrak{g}$-modules, then a linear map $\varphi: V \rightarrow W$ is called a homomorphism of $\mathfrak{g}$-modules if for all $x \in \mathfrak{g}$ and all $v \in V$,

$$
\varphi(x \cdot v)=x \cdot \varphi(v) .
$$

We write $\operatorname{Hom}_{\mathfrak{g}}(V, W)$ for the vector space of all $\mathfrak{g}$-module homomorphisms from $V$ to $W$ and note that the set $\operatorname{End}_{\mathfrak{g}}(V):=\operatorname{Hom}_{\mathfrak{g}}(V, V)$ of module endomorphisms of $V$ is an associative subalgebra of $\operatorname{End}(V)$.

If $\varphi \in \operatorname{Hom}_{\mathfrak{g}}(V, W)$ is bijective, then the inverse map $\psi: W \rightarrow V$ is also a homomorphism of $\mathfrak{g}$-modules (Exercise) satisfying $\varphi \circ \psi=\operatorname{id}_{W}$ and $\psi \circ \varphi=\operatorname{id}_{V}$. Therefore, $\varphi$ is called an isomorphism of $\mathfrak{g}$-modules. The set of isomorphisms $V \rightarrow V$ is the group $\operatorname{Aut}_{\mathfrak{g}}(V):=\operatorname{End}_{\mathfrak{g}}(V)^{\times}$of units in the algebra $\operatorname{End}_{\mathfrak{g}}(V)$.

Example 5.1.13. (a) Any Lie algebra $\mathfrak{g}$ carries a natural $\mathfrak{g}$-module structure defined by the adjoint representation $x \cdot y:=[x, y]$. The $\mathfrak{g}$-submodules of $\mathfrak{g}$ are precisely the ideals (cf. Definition 5.1.9).\\
(b) If $\mathfrak{g}=\mathbb{K}$ is the one-dimensional Lie algebra and $V$ a $\mathbb{K}$-vector space, then any endomorphism $D \in \operatorname{End}(V)$ determines a $\mathfrak{g}$-module structure on $V$ defined by $t \cdot v:=t D(v)$. Clearly, each $\mathfrak{g}$-module structure of $V$ is of this form for $D(v)=1 \cdot v$.

Remark 5.1.14. If $\pi: \mathfrak{g} \rightarrow \mathfrak{g l}(V)$ is a representation, then a $\mathfrak{g}$-module structure on $V$ is defined by $x \cdot v=\pi(x) v$. Conversely, for every $\mathfrak{g}$-module $V$, the map $\pi: \mathfrak{g} \rightarrow \mathfrak{g} \mathfrak{l}(V)$ defined by $\pi(x) v=x \cdot v$ is a representation. Thus representations of $\mathfrak{g}$ and $\mathfrak{g}$-modules are equivalent concepts.

Definition 5.1.15. A representation ( $\pi, V$ ) of a Lie algebra $\mathfrak{g}$ is called irreducible if $V$ is a simple $\mathfrak{g}$-module. It is called completely reducible if $V$ is a semisimple $\mathfrak{g}$-module.

\subsubsection*{Quotients and Semidirect Sums}
We have already seen that the kernel of a homomorphism of Lie algebras is an ideal. The following proposition implies in particular that each ideal is the kernel of a surjective homomorphism of Lie algebras.

Proposition 5.1.16. Let $\mathfrak{g}$ be a Lie algebra and $\mathfrak{n}$ be an ideal in $\mathfrak{g}$. Then the quotient space $\mathfrak{g} / \mathfrak{n}=\{x+\mathfrak{n}: x \in \mathfrak{g}\}$ is a Lie algebra with respect to the bracket

$$
[x+\mathfrak{n}, y+\mathfrak{n}]:=[x, y]+\mathfrak{n} .
$$

The quotient map $\pi: \mathfrak{g} \rightarrow \mathfrak{g} / \mathfrak{n}$ is a surjective homomorphism of Lie algebras with kernel $\mathfrak{n}$.

Proof. The decisive step in the proof is to show that the Lie bracket is well defined. But this immediately follows from the definition of an ideal. All other properties one can easily derive from the respective properties of the Lie bracket on $\mathfrak{g}$ (Exercise). \(\square\)

Theorem 5.1.17. Let $\mathfrak{g}$ and $\mathfrak{h}$ be Lie algebras.\\
(i) If $\alpha: \mathfrak{g} \rightarrow \mathfrak{h}$ is a homomorphism, then

$$
\alpha(\mathfrak{g}) \cong \mathfrak{g} / \operatorname{ker} \alpha .
$$

For any ideal $\mathfrak{i} \unlhd \mathfrak{g}$ with $\mathfrak{i} \subseteq \operatorname{ker} \alpha$, there is exactly one homomorphism $\beta: \mathfrak{g} / \mathfrak{i} \rightarrow \mathfrak{h}$ with $\beta \circ \pi=\alpha$, where $\pi: \mathfrak{g} \rightarrow \mathfrak{g} / \mathfrak{i}$ is the quotient map.\\
(ii) If $\mathfrak{i}, \mathfrak{j} \unlhd \mathfrak{g}$ are ideals with $\mathfrak{i} \subseteq \mathfrak{j}$, then $\mathfrak{j} / \mathfrak{i} \unlhd \mathfrak{g} / \mathfrak{i}$, and $(\mathfrak{g} / \mathfrak{i}) /(\mathfrak{j} / \mathfrak{i}) \cong \mathfrak{g} / \mathfrak{j}$.\\
(iii) If $\mathfrak{i}, \mathfrak{j} \unlhd \mathfrak{g}$ are two ideals, then $\mathfrak{i}+\mathfrak{j}$ and $\mathfrak{i} \cap \mathfrak{j}$ are ideals of $\mathfrak{g}$, and

$$
\mathfrak{i} /(\mathfrak{i} \cap \mathfrak{j}) \cong(\mathfrak{i}+\mathfrak{j}) / \mathfrak{j} .
$$

Proof. Exercise. \(\square\)

We have already seen that we obtain for each ideal $\mathfrak{n} \unlhd \mathfrak{g}$ a quotient algebra $\mathfrak{g} / \mathfrak{n}$, so that we may consider the two Lie algebras $\mathfrak{n}$ and $\mathfrak{g} / \mathfrak{n}$ as two pieces into which $\mathfrak{g}$ is decomposed. It is therefore a natural question how we may build a Lie algebra $\mathfrak{g}$ from two Lie algebras $\mathfrak{n}$ and $\mathfrak{h}$ in such a way that $\mathfrak{n} \unlhd \mathfrak{g}$ and $\mathfrak{g} / \mathfrak{n} \cong \mathfrak{h}$. The following definition describes one such construction (cf. Section 7.6 for more information).

Definition 5.1.18. Let $\mathfrak{n}$ and $\mathfrak{h}$ be Lie algebras and $\alpha: \mathfrak{h} \rightarrow \operatorname{der}(\mathfrak{n})$ be a homomorphism. Then the direct sum $\mathfrak{n} \oplus \mathfrak{h}$ of the vector spaces $\mathfrak{n}$ and $\mathfrak{h}$ is a Lie algebra with respect to the bracket

$$
\left[(x, y),\left(x^{\prime}, y^{\prime}\right)\right]:=\left(\alpha(y) x^{\prime}-\alpha\left(y^{\prime}\right) x+\left[x, x^{\prime}\right],\left[y, y^{\prime}\right]\right)
$$

for $x, x^{\prime} \in \mathfrak{n}, y, y^{\prime} \in \mathfrak{h}$. This Lie algebra is called the semidirect sum with respect to $\alpha$ of $\mathfrak{n}$ and $\mathfrak{h}$. It is denoted by $\mathfrak{n} \rtimes_{\alpha} \mathfrak{h}$. If $\alpha=0$, then $\mathfrak{n} \rtimes_{\alpha} \mathfrak{h}$ is called the direct sum of $\mathfrak{n}$ and $\mathfrak{h}$, and it is denoted by $\mathfrak{n} \oplus \mathfrak{h}$.

The subspace $\left\{(x, 0) \in \mathfrak{n} \rtimes_{\alpha} \mathfrak{h}\right\}$ is an ideal in $\mathfrak{n} \rtimes_{\alpha} \mathfrak{h}$ isomorphic to $\mathfrak{n}$ and $\left\{(0, y) \in \mathfrak{n} \rtimes_{\alpha} \mathfrak{h}\right\}$ is a subalgebra of $\mathfrak{n} \rtimes_{\alpha} \mathfrak{h}$ isomorphic to $\mathfrak{h}$ (cf. Exercise 5.1.2).

For a derivation $D \in \operatorname{der} \mathfrak{n}$, we simply write $\mathfrak{n} \rtimes_{D} \mathbb{K}$ for the semidirect sum defined by $\alpha(t):=t D$.

Example 5.1.19. Let $\mathfrak{h}_{3}(\mathbb{R})$ be the 3-dimensional vector space with the basis $p, q, z$ equipped with the skew-symmetric bracket determined by

$$
[p, q]=z, \quad[p, z]=[q, z]=0 .
$$

Then $\mathfrak{h}_{3}(\mathbb{R})$ is a Lie algebras called the three-dimensional Heisenberg algebra. It is isomorphic to the algebra $\mathfrak{n}$ in Example 5.1.6(viii) for $n=3$. The linear endomorphism of $\mathfrak{h}_{3}(\mathbb{R})$ defined by

$$
D z=0, \quad D p=q, \quad \text { and } \quad D q=-p
$$

is then a derivation of $\mathfrak{h}_{3}(\mathbb{R})$, so that we obtain a Lie algebra $\mathfrak{g}:= \mathfrak{h}_{3}(\mathbb{R}) \rtimes_{D} \mathbb{R}$, called the oscillator algebra. Writing $h:=(0,1)$ for the additional basis element in $\mathfrak{g}$, the nonzero brackets of basis elements are

$$
[p, q]=z, \quad[h, p]=q, \quad \text { and } \quad[h, q]=-p .
$$

Example 5.1.20. If $\mathcal{F}=\left(V_{0}, \ldots, V_{n}\right)$ is a flag in the vector space $V$ (Example 5.1.7), then we know already the associated Lie algebra

$$
\mathfrak{g}(\mathcal{F})=\left\{x \in \mathfrak{g} \mathfrak{l}(V):(\forall j) x V_{j} \subseteq V_{j}\right\} .
$$

It is easy to see that

$$
\mathfrak{g}_{n}(\mathcal{F}):=\left\{x \in \mathfrak{g} \mathfrak{l}(V):(\forall j>0) x V_{j} \subseteq V_{j-1}\right\}
$$

is an ideal of $\mathfrak{g}(\mathcal{F})$. Here the $n$ in $\mathfrak{g}_{n}(\mathcal{F})$ stands for "nilpotent".\\
To find a subalgebra complementary to this ideal, we choose subspaces $W_{0}, \ldots, W_{n-1}$ of $V$ with $V_{j+1} \cong V_{j} \oplus W_{j}$ for $j=0, \ldots, n-1$. Then

$$
\mathfrak{g}_{s}(\mathcal{F}):=\left\{X \in \mathfrak{g} \mathfrak{l}(V):(\forall j) X W_{j} \subseteq W_{j}\right\} \subseteq \mathfrak{g}(\mathcal{F})
$$

is a subalgebra with

$$
\mathfrak{g}(\mathcal{F}) \cong \mathfrak{g}_{n}(\mathcal{F}) \rtimes \mathfrak{g}_{s}(\mathcal{F}) \quad \text { and } \quad \mathfrak{g}_{s}(\mathcal{F}) \cong \bigoplus_{j=1}^{n} \mathfrak{g} \mathfrak{l}\left(W_{j}\right)
$$

The $s$ in $\mathfrak{g}_{s}(\mathcal{F})$ stands for "semisimple". Describing the elements of $\mathfrak{g}(\mathcal{F})$ as in Example 5.1.7 by block matrices, the semidirect decomposition of the Lie algebra $\mathfrak{g}(\mathcal{F})$ corresponds to the decomposition of an upper triangular matrix as a sum of a strictly upper triangular matrix and a diagonal matrix. For $n=3$, we have in particular:

$$
A=\left(\begin{array}{ccc}
A_{11} & A_{12} & A_{13} \\
0 & A_{22} & A_{23} \\
0 & 0 & A_{33}
\end{array}\right)=\underbrace{\left(\begin{array}{ccc}
A_{11} & 0 & 0 \\
0 & A_{22} & 0 \\
0 & 0 & A_{33}
\end{array}\right)}_{\in \mathfrak{g}_{s}(\mathcal{F})}+\underbrace{\left(\begin{array}{ccc}
0 & A_{12} & A_{13} \\
0 & 0 & A_{23} \\
0 & 0 & 0
\end{array}\right)}_{\in \mathfrak{g}_{n}(\mathcal{F})} .
$$

\subsubsection*{Complexification and Real Forms}
Up to now, the base field did not really play a role in our considerations. Nevertheless, it is important (just think about basic linear algebra) to note that, for instance, homomorphisms of complex Lie algebras are supposed to be complex linear. Therefore, we also consider the complexification of a real Lie algebra. For this, we briefly recall how to calculate with the complexification of a vector space.

Let $V$ be an $\mathbb{R}$-vector space. The complexification $V_{\mathbb{C}}$ of $V$ is the tensor product $\mathbb{C} \otimes_{\mathbb{R}} V$. This is a complex vector space with respect to the scalar multiplication $\lambda \cdot(z \otimes v):=\lambda z \otimes v$. Identifying $V$ with the subspace $1 \otimes V$ of $V_{\mathbb{C}}$, each element of $V_{\mathbb{C}}$ can be written in a unique fashion as $z=x+i y$ with $x, y \in V$. If $\left\{v_{1}, \ldots, v_{n}\right\}$ is a real basis for $V$, then $\left\{1 \otimes v_{1}, \ldots, 1 \otimes v_{n}\right\}$ is a complex basis for $V_{\mathbb{C}}$.

Proposition 5.1.21. Let $\mathfrak{g}$ be a real Lie algebra.\\
(i) $\mathfrak{g}_{\mathbb{C}}$ is a complex Lie algebra with respect to the complex bilinear Lie bracket, defined by

$$
\left[x+i y, x^{\prime}+i y^{\prime}\right]:=\left(\left[x, x^{\prime}\right]-\left[y, y^{\prime}\right]\right)+i\left(\left[x, y^{\prime}\right]+\left[y, x^{\prime}\right]\right)
$$

and satisfying

$$
\left[z \otimes v, z^{\prime} \otimes v^{\prime}\right]=z z^{\prime} \otimes\left[v, v^{\prime}\right]
$$

(ii) $\left[\mathfrak{g}_{\mathbb{C}}, \mathfrak{g}_{\mathbb{C}}\right] \cong[\mathfrak{g}, \mathfrak{g}]_{\mathbb{C}}$ as complex Lie algebras.

Proof. Elementary calculations. \(\square\)

Definition 5.1.22. Let $\mathfrak{g}$ be a complex Lie algebra. A real Lie algebra $\mathfrak{h}$ with $\mathfrak{h}_{\mathbb{C}} \cong \mathfrak{g}$ is called a real form of $\mathfrak{g}$.

We have seen that to every real Lie algebra we can assign a complexification in a unique way. However, nonisomorphic real algebras can have isomorphic complexifications, resp., complex Lie algebras can have nonisomorphic real forms.

Example 5.1.23. Let $\mathfrak{s o}_{3}(\mathbb{R})=\mathfrak{o}_{3}(\mathbb{R}) \cap \mathfrak{s} \mathfrak{l}_{3}(\mathbb{R})=\mathfrak{o}_{3}(\mathbb{R})$. Then the complexifications of $\mathfrak{s o}_{3}(\mathbb{R})$ and $\mathfrak{s} \mathfrak{l}_{2}(\mathbb{R})$ are both isomorphic to $\mathfrak{s} \mathfrak{l}_{2}(\mathbb{C})$. To see this, we consider the bases

$$
h=\left(\begin{array}{cc}
1 & 0 \\
0 & -1
\end{array}\right), \quad u=\left(\begin{array}{cc}
0 & 1 \\
-1 & 0
\end{array}\right), \quad t=\left(\begin{array}{ll}
0 & 1 \\
1 & 0
\end{array}\right)
$$

of $\mathfrak{s} \mathfrak{l}_{2}(\mathbb{R})$, and

$$
x=\left(\begin{array}{ccc}
0 & 1 & 0 \\
-1 & 0 & 0 \\
0 & 0 & 0
\end{array}\right), \quad y=\left(\begin{array}{ccc}
0 & 0 & 1 \\
0 & 0 & 0 \\
-1 & 0 & 0
\end{array}\right), \quad z=\left(\begin{array}{ccc}
0 & 0 & 0 \\
0 & 0 & -1 \\
0 & 1 & 0
\end{array}\right)
$$

of $\mathfrak{s o}_{3}(\mathbb{R})$. Then

$$
[h, u]=2 t, \quad[h, t]=2 u, \quad[u, t]=2 h
$$

and

$$
[x, y]=z, \quad[z, x]=y, \quad[y, z]=x .
$$

Let

$$
\mathfrak{h}=\mathbb{R} i h+\mathbb{R} u+\mathbb{R} i t .
$$

Then $\mathfrak{h}$ is a Lie algebra with $\mathfrak{h}_{\mathbb{C}}=\mathfrak{s l}_{2}(\mathbb{C})$ which is isomorphic to $\mathfrak{s o}_{3}(\mathbb{R})$ via

$$
i h \mapsto 2 x, u \mapsto 2 y, i t \mapsto 2 z .
$$

We note that $\mathfrak{h}$ coincides with $\mathfrak{s u}_{2}(\mathbb{C})$. Since, obviously, $\mathfrak{s} \mathfrak{l}_{2}(\mathbb{R})_{\mathbb{C}}=\mathfrak{s} \mathfrak{l}_{2}(\mathbb{C})$, it only remains to show that $\mathfrak{s} \mathfrak{l}_{2}(\mathbb{R})$ and $\mathfrak{s o}_{3}(\mathbb{R})$ are not isomorphic. For this, it suffices to check that $\mathbb{R} h+\mathbb{R}(u+t)$ is a two-dimensional subalgebra of $\mathfrak{s} \mathfrak{l}_{2}(\mathbb{R})$, while $\mathfrak{s o}_{3}(\mathbb{R})$ has no two-dimensional subalgebra. Namely, the latter is isomorphic to $\mathbb{R}^{3}$ with the vector product (Exercise 5.1.6), and since the vector product of two vectors is orthogonal to these vectors, a plane cannot be a subalgebra.

Example 5.1.24 (A complex Lie algebra with no real form). On the abelian Lie algebra $V:=\mathbb{C}^{2}$, we consider the linear operator $D$, defined by $D e_{1}=2 e_{1}$ and $D e_{2}=i e_{2}$ with respect to the canonical basis. Then we form the three-dimensional complex Lie algebra $\mathfrak{g}:=V \rtimes_{D} \mathbb{C}$ and note that $V=[\mathfrak{g}, \mathfrak{g}]$ is a 2-dimensional ideal of $\mathfrak{g}$.

Suppose that $\mathfrak{g}$ has a real form. Let $\sigma \in \operatorname{Aut}(\mathfrak{g})$ be the corresponding complex conjugation, which is an involutive automorphism of $\mathfrak{g}$. Then

$$
\sigma(V)=\sigma([\mathfrak{g}, \mathfrak{g}])=[\sigma(\mathfrak{g}), \sigma(\mathfrak{g})]=[\mathfrak{g}, \mathfrak{g}]=V,
$$

so that $\sigma$ induces an antilinear involution $\sigma_{V}$ on $V$. Let $\sigma(0,1)=\left(v_{0}, \lambda\right)$ and note that $\sigma(V)=V$ implies that $\lambda \neq 0$. Applying $\sigma$ again, we see that

$$
(0,1)=\sigma^{2}(0,1)=\sigma_{V}\left(v_{0}\right)+\bar{\lambda} \sigma(0,1)=\left(\sigma_{V}\left(v_{0}\right)+\bar{\lambda} v_{0}, \bar{\lambda} \cdot \lambda\right)
$$

We conclude that $|\lambda|=1$. Further, $\sigma \circ \operatorname{ad}(0,1) \circ \sigma=\operatorname{ad}(\sigma(0,1))=\operatorname{ad}\left(v_{0}, \lambda\right)$ implies by restricting to $V$ that

$$
\sigma_{V} \circ D \circ \sigma_{V}=\lambda D .
$$

If $v \in V$ is a $D$-eigenvector with $D v=\alpha v$, then

$$
D\left(\sigma_{V} v\right)=\sigma_{V}(\lambda D v)=\bar{\lambda} \bar{\alpha} \sigma_{V}(v)
$$

This means that $\sigma_{V}\left(V_{\alpha}(D)\right)=V_{\bar{\lambda} \bar{\alpha}}(D)$. In particular, $\sigma_{V}$ permutes the $D$ eigenspaces. Now $|\lambda|=1$ and $|i| \neq 2$ show that $\sigma_{V}$ preserves both eigenspaces. For $\alpha=2$, this leads to $\bar{\lambda}=1$, so that $\lambda=1$. For $\alpha=i$, we now arrive at the contradiction $-i=\bar{\lambda} \bar{\alpha}=\alpha=i$.

This example is minimal because each complex Lie algebra of dimension 2 has a real form (cf. Example 5.4.2).

\subsubsection*{Exercises for Section 5.1}
Exercise 5.1.1. Let $A$ be an associative algebra and $A_{L}$ be the associated Lie algebra (cf. Example 5.1.2).\\
(i) A derivation of $A$ is a linear map $\delta: A \rightarrow A$ such that

$$
\delta(a b)=\delta(a) b+a \delta(b) \quad \forall a, b \in A .
$$

Then $\operatorname{der}(A) \subseteq \operatorname{der}\left(A_{L}\right)$, i.e., every derivation of the associative algebra $A$ is a derivation of the Lie algebra $A_{L}$, too.\\
(ii) $[a, b c]=[a, b] c+b[a, c]$ for $a, b, c \in A$.\\
(iii) In general, $\operatorname{der}(A) \neq \operatorname{der}\left(A_{L}\right)$.\\
(iv) If $A$ is commutative, then $A \cdot \operatorname{der}(A) \subseteq \operatorname{der}(A)$.

Exercise 5.1.2. Let $U$ be an open subset of $\mathbb{R}^{2 n}$ and $\mathfrak{g}=C^{\infty}(U, \mathbb{R})$ be the set of smooth functions on $U$ and write $q_{1}, \ldots, q_{n}, p_{1}, \ldots, p_{n}$ for the coordinates with respect to a basis. Then $\mathfrak{g}$ is a Lie algebra with respect to the Poisson bracket

$$
\{f, g\}:=\sum_{i=1}^{n} \frac{\partial f}{\partial q_{i}} \frac{\partial g}{\partial p_{i}}-\frac{\partial f}{\partial p_{i}} \frac{\partial g}{\partial q_{i}} .
$$

Exercise 5.1.3. Let $U$ be an open subset of $\mathbb{R}^{n}, A=C^{\infty}(U, \mathbb{R})$, and $\mathfrak{g}= C^{\infty}\left(U, \mathbb{R}^{n}\right)$. For $f \in A$ and $X \in \mathfrak{g}$, we define

$$
\mathcal{L}_{X} f:=X f:=\sum_{i=1}^{n} X_{i} \frac{\partial f}{\partial x_{i}}
$$

(i) The maps $\mathcal{L}_{X}$ are derivations of the algebra $A$.\\
(ii) If $\mathcal{L}_{X}=0$, then $X=0$.\\
(iii) The commutator of two such operators has the form $\left[\mathcal{L}_{X}, \mathcal{L}_{Y}\right]=\mathcal{L}_{[X, Y]}$, where the bracket on $\mathfrak{g}$ is defined by (cf. Definition 8.1.1)

$$
[X, Y](p):=\mathrm{d} Y(p) X(p)-\mathrm{d} X(p) Y(p),
$$

resp.,

$$
[X, Y]_{i}=\sum_{j=1}^{n} X_{j} \frac{\partial Y_{i}}{\partial x_{j}}-Y_{j} \frac{\partial X_{i}}{\partial x_{j}}
$$

(iv) $(\mathfrak{g},[\cdot, \cdot])$ is a Lie algebra.\\
(v) To each $A \in \mathfrak{g l}_{n}(\mathbb{R})$, we associate the linear vector field $X_{A}(x):=A x$. Show that, for $A, B \in M_{n}(\mathbb{R})$, we have $X_{[A, B]}=-\left[X_{A}, X_{B}\right]$.

Exercise 5.1.4. Let $A$ be a vector space and $m: A \times A \rightarrow A$ a bilinear map. Then the space

$$
\operatorname{der}(A, m):=\{D \in \mathfrak{g} \mathfrak{l}(A):(\forall a, b \in A) D m(a, b)=m(D a, b)+m(a, D b)\}
$$

of derivations is a Lie subalgebra of $\mathfrak{g l}(A)$.\\
Exercise 5.1.5. Let $\mathfrak{g}$ be a Lie algebra, $\mathfrak{n} \unlhd \mathfrak{g}$ an ideal, and $\mathfrak{h}<\mathfrak{g}$ a Lie subalgebra with $\mathfrak{g}=\mathfrak{n}+\mathfrak{h}$ and $\mathfrak{n} \cap \mathfrak{h}=\{0\}$. Then

$$
\delta: \mathfrak{h} \rightarrow \operatorname{der} \mathfrak{n}, \quad \delta(x):=\left.\operatorname{ad} x\right|_{\mathfrak{n}}
$$

defines a homomorphism of Lie algebras and the map

$$
\Phi: \mathfrak{n} \rtimes_{\delta} \mathfrak{h} \rightarrow \mathfrak{g}, \quad(x, y) \mapsto x+y
$$

is an isomorphism of Lie algebras.

Exercise 5.1.6. On $\mathbb{R}^{3}$, we define the vector product by

$$
\left(\begin{array}{l}
v_{1} \\
v_{2} \\
v_{3}
\end{array}\right) \times\left(\begin{array}{l}
w_{1} \\
w_{2} \\
w_{3}
\end{array}\right):=\left(\begin{array}{l}
v_{2} w_{3}-v_{3} w_{2} \\
v_{3} w_{1}-v_{1} w_{3} \\
v_{1} w_{2}-v_{2} w_{1}
\end{array}\right) .
$$

Show that $\left(\mathbb{R}^{3}, \times\right)$ is a Lie algebra and that the map

$$
\Phi: \mathbb{R}^{3} \rightarrow \mathfrak{s o}_{3}(\mathbb{R}), \quad \Phi(v):=v_{1} x+v_{2} y+v_{3} z
$$

(in the notation of Example 5.1.23) is an isomorphism of Lie algebras.\\
Exercise 5.1.7. Show that $\left[\mathfrak{g} \mathfrak{l}_{n}(\mathbb{K}), \mathfrak{g} \mathfrak{l}_{n}(\mathbb{K})\right]=\mathfrak{s} \mathfrak{l}_{n}(\mathbb{K})$ and $\mathfrak{z}\left(\mathfrak{g} \mathfrak{l}_{n}(\mathbb{K})\right)=\mathbb{K} \mathbf{1}$.\\
Exercise 5.1.8. On the algebra $A:=C^{\infty}(\mathbb{R}, \mathbb{C})$, consider the operators

$$
P f(x):=i f^{\prime}(x), \quad Q f(x):=x f(x), \quad \text { and } \quad Z f(x)=i f(x) .
$$

Then the Lie subalgebra of $\mathfrak{g l}(A)$ generated by $P, Q$, and $Z$ is isomorphic to the Heisenberg algebra $\mathfrak{h}_{3}(\mathbb{R})$, i.e.,

$$
[P, Q]=Z \quad \text { and } \quad[P, Z]=[Q, Z]=0 .
$$

Adding also the operator

$$
H f(x):=\frac{i}{2}\left(-\frac{d^{2} f}{d x^{2}}(x)+x^{2} f(x)\right), \quad H=\frac{i}{2}\left(P^{2}+Q^{2}\right)
$$

we obtain a four-dimensional subalgebra, isomorphic to the oscillator algebra (Example 5.1.19).

Exercise 5.1.9. Show that for two ideals $\mathfrak{a}$ and $\mathfrak{b}$ of the Lie algebra $\mathfrak{g}$, the subspace $[\mathfrak{a}, \mathfrak{b}]$ also is an ideal.

Exercise 5.1.10. Let ( $\pi, V$ ) be a representation of the Lie algebra $\mathfrak{g}$ on $V$ and $W \subseteq V$ a $\mathfrak{g}$-invariant subspace, i.e., $\pi(\mathfrak{g}) W \subseteq W$. Then

$$
\bar{\pi}: \mathfrak{g} \rightarrow \mathfrak{g} \mathfrak{l}(V / W), \quad \bar{\pi}(x)(v+W):=\pi(x) v+W
$$

defines a representation of $\mathfrak{g}$ on the quotient space $V / W$.\\
Exercise 5.1.11. For the following Lie algebras, find a faithful, i.e., injective, finite-dimensional representation: $\mathfrak{s} \mathfrak{l}_{2}(\mathbb{K})$, the Heisenberg algebra, the oscillator algebra, and the abelian Lie algebra $\mathbb{R}^{n}$.

\subsection*{Nilpotent Lie Algebras}
In the following, we shall encounter several important classes of Lie algebras that play a central role in the structure theory of finite-dimensional Lie algebras. The first of these two classes, nilpotent Lie algebras, are those for which iterated brackets $\left[x_{1},\left[x_{2},\left[x_{3},\left[x_{4}, \ldots\right]\right]\right]\right]$ of sufficiently large order vanish. The most important result on nilpotent Lie algebras is Engel's Theorem which translates nilpotency of a Lie algebra into a pointwise condition. Typical examples of nilpotent Lie algebras are Lie algebras of strictly upper triangular (block) matrices.

Definition 5.2.1. Let $\mathfrak{g}$ be a Lie algebra. We define its descending (lower) central series inductively by

$$
C^{1}(\mathfrak{g}):=\mathfrak{g} \quad \text { and } \quad C^{n+1}(\mathfrak{g}):=\left[\mathfrak{g}, C^{n}(\mathfrak{g})\right] .
$$

In particular, $C^{2}(\mathfrak{g})=[\mathfrak{g}, \mathfrak{g}]$ is the commutator algebra. The Lie algebra $\mathfrak{g}$ is called nilpotent, if there is a $d \in \mathbb{N}_{0}$ with $C^{d+1}(\mathfrak{g})=\{0\}$. If $d$ is minimal with this property, then it is called the nilpotence degree of $\mathfrak{g}$. By induction, one immediately sees that each $C^{n}(\mathfrak{g})$ is an ideal of $\mathfrak{g}$, so that $C^{n+1}(\mathfrak{g}) \subseteq C^{n}(\mathfrak{g})$. Hence, for finite-dimensional Lie algebras, the nilpotency of $\mathfrak{g}$ is equivalent to the vanishing of the ideal $C^{\infty}(\mathfrak{g}):=\bigcap_{n \in \mathbb{N}} C^{n}(\mathfrak{g})$.

Example 5.2.2. (i) The Heisenberg algebra $\mathfrak{h}_{3}(\mathbb{R})$ is nilpotent.\\
(ii) Every abelian Lie algebra is nilpotent.\\
(iii) If $\mathcal{F}=\left(V_{0}, \ldots, V_{n}\right)$ is a flag in the vector space $V$ and we put $V_{i}:=\{0\}$ for $i<0$, then $\mathfrak{g}_{n}(\mathcal{F})$ is a nilpotent Lie algebra. In fact, an easy induction leads to

$$
C^{m}\left(\mathfrak{g}_{n}(\mathcal{F})\right) V_{n} \subseteq V_{n-m}
$$

and therefore to $C^{n}\left(\mathfrak{g}_{n}(\mathcal{F})\right)=\{0\}$.\\
Proposition 5.2.3. Let $\mathfrak{g}$ be a Lie algebra.\\
(i) If $\mathfrak{g}$ is nilpotent, then all subalgebras and all homomorphic images of $\mathfrak{g}$ are nilpotent.\\
(ii) If $\mathfrak{a}<\mathfrak{z}(\mathfrak{g})$ and $\mathfrak{g} / \mathfrak{a}$ is nilpotent, then $\mathfrak{g}$ is nilpotent.\\
(iii) If $\mathfrak{g} \neq\{0\}$ is nilpotent, then $\mathfrak{z}(\mathfrak{g}) \neq\{0\}$.\\
(iv) If $\mathfrak{g}$ is nilpotent, then there is an $n \in \mathbb{N}$ with $\operatorname{ad}(x)^{n}=0$ for all $x \in \mathfrak{g}$, i.e., the $\operatorname{ad}(x)$ are nilpotent as linear maps.\\
(v) If $\mathfrak{i} \unlhd \mathfrak{g}$, then all the spaces $C^{n}(\mathfrak{i})$ are ideals of $\mathfrak{g}$.

Proof. (i) If $\mathfrak{h}<\mathfrak{g}$, then $[\mathfrak{h}, \mathfrak{h}] \subseteq[\mathfrak{g}, \mathfrak{g}]$ and $C^{n}(\mathfrak{h}) \subseteq C^{n}(\mathfrak{g})$ follows by induction. Therefore, each subalgebra of a nilpotent Lie algebra is nilpotent.

For a homomorphism $\alpha: \mathfrak{g} \rightarrow \mathfrak{h}$, we obtain inductively

$$
C^{n}(\alpha(\mathfrak{g}))=\alpha\left(C^{n}(\mathfrak{g})\right) \quad \text { for each } n \in \mathbb{N} .
$$

Thus, if $C^{n}(\mathfrak{g})=\{0\}$, then $C^{n}(\operatorname{im} \alpha)=\{0\}$.\\
(ii) If $\mathfrak{g} / \mathfrak{a}$ is nilpotent, then there is an $n \in \mathbb{N}$ with $C^{n}(\mathfrak{g} / \mathfrak{a})=\{0\}$, so that (5.2), applied to the quotient homomorphism $q: \mathfrak{g} \rightarrow \mathfrak{g} / \mathfrak{a}$, leads to $C^{n}(\mathfrak{g}) \subseteq \mathfrak{a} \subseteq \mathfrak{z}(\mathfrak{g})$ and thus to $C^{n+1}(\mathfrak{g}) \subseteq[\mathfrak{g}, \mathfrak{z}(\mathfrak{g})]=\{0\}$.\\
(iii) If $\mathfrak{g} \neq\{0\}$ is nilpotent, for some $n \in \mathbb{N}$, we have $C^{n}(\mathfrak{g})=\{0\}$ and $C^{n-1}(\mathfrak{g}) \neq\{0\}$. Then $\left[\mathfrak{g}, C^{n-1}(\mathfrak{g})\right]=\{0\}$ implies that the nontrivial ideal $C^{n-1}(\mathfrak{g})$ is contained in the center.\\
(iv) If $C^{n+1}(\mathfrak{g})=\{0\}$, then $(\operatorname{ad} x)^{n} \mathfrak{g} \subseteq C^{n+1}(\mathfrak{g})=\{0\}$.\\
(v) In view of Exercise 5.1.9, this follows by induction.

In Proposition 5.2.3, we have seen that for every nilpotent Lie algebra, all the endomorphisms $\operatorname{ad}(x), x \in \mathfrak{g}$, are nilpotent. Now our aim is to show that a finite-dimensional Lie algebra, for which every ad $x$ is nilpotent, is nilpotent itself. We start with a simple lemma, the proof of which we leave to the reader as an exercise (cf. Exercises 4.1.3 and 5.1.10).

Lemma 5.2.4. (i) Let $V$ be a finite-dimensional vector space, $\mathfrak{g} \subseteq \mathfrak{g} \mathfrak{l}(V)$ a subalgebra, and $x \in \mathfrak{g}$. If $x \in \mathfrak{g} \mathfrak{l}(V)$ is nilpotent, then $\operatorname{ad}(x): \mathfrak{g} \rightarrow \mathfrak{g}$ is also nilpotent.\\
(ii) Let $\mathfrak{g}$ be a Lie algebra and $\mathfrak{h}<\mathfrak{g}$. Then

$$
\operatorname{ad}_{\mathfrak{g} / \mathfrak{h}}: \mathfrak{h} \rightarrow \mathfrak{g} \mathfrak{l}(\mathfrak{g} / \mathfrak{h}), \quad \operatorname{ad}_{\mathfrak{g} / \mathfrak{h}}(x)(y+\mathfrak{h}):=[x, y]+\mathfrak{h}
$$

defines a representation of $\mathfrak{h}$ on $\mathfrak{g} / \mathfrak{h}$.\\
Theorem 5.2.5 (Engel's Theorem on Linear Lie Algebras). Let $V \neq\{0\}$ be a finite-dimensional vector space and $\mathfrak{g} \subseteq \mathfrak{g} \mathfrak{l}(V)$ a Lie subalgebra. If all $x \in \mathfrak{g}$ are nilpotent, i.e., $x^{n}=0$ for some $n \in \mathbb{N}$, then there exists a nonzero $v_{o} \in V$ with $\mathfrak{g}\left(v_{o}\right)=\{0\}$.

Proof. We proceed by induction on $\operatorname{dim} \mathfrak{g}$. For $\operatorname{dim} \mathfrak{g}=0$, the assertion holds trivially for each nonzero $v_{o} \in V$.

Next we assume that $\operatorname{dim} \mathfrak{g}>0$ and pick a proper subalgebra $\mathfrak{h}<\mathfrak{g}$ of maximal dimension. According to Lemma 5.2.4, for each $x \in \mathfrak{h}$ the operators $\operatorname{ad}_{\mathfrak{g} / \mathfrak{h}}(x)$ are nilpotent. Now our induction hypothesis implies the existence of some $x_{o} \in \mathfrak{g} \backslash \mathfrak{h}$ with $\operatorname{ad}_{\mathfrak{g} / \mathfrak{h}}(\mathfrak{h})\left(x_{o}\right)=\{0\}$, i.e., $\left[\mathfrak{h}, x_{o}\right] \subseteq \mathfrak{h}$. This implies that $\mathbb{K} x_{o}+\mathfrak{h}$ is a subalgebra of $\mathfrak{g}$ and, again by maximality of $\mathfrak{h}$, it follows that $\mathbb{K} x_{o}+\mathfrak{h}=\mathfrak{g}$. The induction hypothesis also implies that the space $V_{o}:=\{v \in V: \mathfrak{h}(v)=\{0\}\}$ is nonzero. Moreover,

$$
y x(w)=x y(w)-[x, y](w) \in x \mathfrak{h}(w)+\mathfrak{h}(w)=\{0\}
$$

for $x \in \mathfrak{g}, y \in \mathfrak{h}, w \in V_{o}$, implies that $\mathfrak{g}\left(V_{o}\right) \subseteq V_{o}$. Since $\left.x_{o}\right|_{V_{o}}$ is also nilpotent, there exists a nonzero $v_{o} \in V_{o}$ with $x_{o}\left(v_{o}\right)=0$. Putting all this together, we arrive at $\mathfrak{g}\left(v_{o}\right)=\mathfrak{h}\left(v_{o}\right)+\mathbb{K} x_{o}\left(v_{o}\right)=\{0\}$.

Exercise 5.2.8 discusses an interesting Lie algebra of nilpotent endomorphisms of an infinite-dimensional space, showing in particular that Engel's Theorem does not generalize to infinite-dimensional spaces.

From Theorem 5.2.5, one also gets the induction step by which one shows that every subalgebra of $\mathfrak{g l}(V)$ which consists of nilpotent elements can be written as a Lie algebra of triangular matrices.

Definition 5.2.6. Let $V$ be an $n$-dimensional vector space. A complete flag in $V$ is a flag $\left(V_{0}, \ldots, V_{n}\right)$ with $\operatorname{dim} V_{k}=k$ for each $k$.

Corollary 5.2.7. Let $V$ be a finite-dimensional vector space and $\mathfrak{g}<\mathfrak{g} \mathfrak{l}(V)$ such that all elements of $\mathfrak{g}$ are nilpotent. Then there exists a complete flag $\mathcal{F}$ in $V$ with $\mathfrak{g} \subseteq \mathfrak{g}_{n}(\mathcal{F})$. In particular, there is a basis for $V$ with respect to which the elements of $\mathfrak{g}$ correspond to strictly upper triangular matrices. In particular, $\mathfrak{g}$ is nilpotent.

Proof. In view of Theorem 5.2.5, there exists a nonzero $v_{1} \in V$ with $\mathfrak{g}\left(v_{1}\right)= \{0\}$. We set $V_{1}:=\mathbb{K} v_{1}$. Then

$$
\alpha: \mathfrak{g} \rightarrow \mathfrak{g} \mathfrak{l}\left(V / V_{1}\right), \quad \alpha(x)\left(v+V_{1}\right):=x(v)+V_{1}
$$

is a representation of $\mathfrak{g}$ on $V / V_{1}$ (Exercise 5.1.10), and $\alpha(\mathfrak{g})$ consists of nilpotent endomorphisms. We now proceed by induction on $\operatorname{dim} V$, so that the induction hypothesis implies that $V / V_{1}$ possesses a complete flag $\mathcal{F}_{1}= \left(W_{1}, \ldots, W_{k}\right)$ with $\alpha(\mathfrak{g}) \subseteq \mathfrak{g}_{n}\left(\mathcal{F}_{1}\right)$. Then $\{0\}$, together with the preimage of the flag $\mathcal{F}_{1}$ in $V$, is a complete flag $\mathcal{F}$ in $V$ with $\mathfrak{g} \subseteq \mathfrak{g}_{n}(\mathcal{F})$. Since $\mathfrak{g}_{n}(\mathcal{F})$ is nilpotent (Example 5.2.2(iii)), the subalgebra $\mathfrak{g}$ is also nilpotent.

Now we are able to prove the previously announced criterion for the nilpotency of a Lie algebra.

Theorem 5.2.8 (Engel's Characterization Theorem for Nilpotent Lie Algebras). Let $\mathfrak{g}$ be a finite-dimensional Lie algebra. Then $\mathfrak{g}$ is nilpotent if and only if for each $x \in \mathfrak{g}$ the operator ad $x$ is nilpotent.

Proof. We have already seen in Proposition 5.2.3 that for each $x \in \mathfrak{g}$ the operator ad $x$ is nilpotent. It remains to show the converse.

If ad $x$ is nilpotent for each $x \in \mathfrak{g}$, then Corollary 5.2.7 implies that the Lie algebra $\operatorname{ad}(\mathfrak{g}) \cong \mathfrak{g} / \mathfrak{z}(\mathfrak{g})$ is nilpotent. Now Proposition 5.2.3(ii) shows that $\mathfrak{g}$ is also nilpotent.

Lemma 5.2.9. If $\mathfrak{a}$ and $\mathfrak{b}$ are nilpotent ideals of the Lie algebra $\mathfrak{g}$, then so is their sum $\mathfrak{a}+\mathfrak{b}$.

Proof. We claim that

$$
C^{2 m}(\mathfrak{a}+\mathfrak{b}) \subseteq C^{m}(\mathfrak{a})+C^{m}(\mathfrak{b}) \quad \text { for } \quad m \in \mathbb{N}
$$

This implies the assertion because $C^{m}(\mathfrak{a})=C^{m}(\mathfrak{b})=\{0\}$ holds if $m$ is sufficiently large.

The space $C^{2 m}(\mathfrak{a}+\mathfrak{b})$ is spanned by elements of the form

$$
y:=\left[x_{1},\left[x_{2},\left[x_{3}, \ldots\left[x_{2 m-1}, x_{2 m}\right] \ldots\right]\right]\right]
$$

with $x_{j} \in \mathfrak{a} \cup \mathfrak{b}$. If at least $m$ of the $x_{j}$ are contained in $\mathfrak{a}$, then $y \in C^{m}(\mathfrak{a})$. If this is not the case, then $m$ of the $x_{j}$ are contained in $\mathfrak{b}$, which leads to $y \in C^{m}(\mathfrak{b})$. This proves our claim and hence the lemma.

Definition 5.2.10. The main consequence of Lemma 5.2.9 is that it implies that in every finite-dimensional Lie algebra $\mathfrak{g}$, there is a largest nilpotent ideal. In fact, if $\mathfrak{n} \unlhd \mathfrak{g}$ is a nilpotent ideal of maximal dimension and $\mathfrak{m} \unlhd \mathfrak{g}$ any other nilpotent ideal, then the nilpotency of $\mathfrak{n}+\mathfrak{m}$ implies that $\mathfrak{m} \subseteq \mathfrak{n}$. Therefore, the ideal $\mathfrak{n}$ contains all other nilpotent ideals.

The maximal nilpotent ideal is called the nilradical of $\mathfrak{g}$, and is denoted by $\mathfrak{n i l}(\mathfrak{g})$.

Remark 5.2.11. One should be aware of the fact that some authors use the term "nilradical" with a different meaning, namely for the intersection of the kernels of all irreducible finite-dimensional representations. Since for an abelian Lie algebra $\mathfrak{g}$, the one-dimensional representations separate the points, we have $\mathfrak{n i l}(\mathfrak{g})=\mathfrak{g}$, but the intersection of the kernels of irreducible finite-dimensional representations is $\{0\}$. For more details on this ideal, see Proposition 5.7.4.

Remark 5.2.12. If $\mathfrak{g}$ is not finite dimensional, then the preceding lemma implies that for each finite collection $\mathfrak{n}_{1}, \ldots, \mathfrak{n}_{k} \unlhd \mathfrak{g}$ of nilpotent ideals, their sum $\mathfrak{n}_{1}+\cdots+\mathfrak{n}_{k}$ is a nilpotent ideal. Therefore, the sum $\mathfrak{n}$ of all nilpotent ideals of $\mathfrak{g}$ coincides with the union of all nilpotent ideals. However, this ideal need not be nilpotent because there may be nilpotent ideals of an arbitrary high nilpotence degree.

\subsubsection*{Exercises for Section 5.2}
Exercise 5.2.1. Let $V$ be a finite-dimensional complex vector space and $x \in \operatorname{End}(V)$ with eigenvalues $\lambda_{1}, \ldots, \lambda_{n}$. Then the eigenvalues of ad $x$ are all numbers of the form

$$
\lambda_{i}-\lambda_{j}, \quad i, j=1, \ldots, n
$$

Exercise 5.2.2. Let $\mathfrak{g}$ be a Lie algebra, $\mathfrak{h}$ a subalgebra and $x \in \mathfrak{n}_{\mathfrak{g}}(\mathfrak{h}) \backslash \mathfrak{h}$. Then $\mathfrak{h}+\mathbb{R} x \cong \mathfrak{h} \rtimes_{\alpha} \mathbb{R} x$ for $\alpha(t x)=\left.\operatorname{ad}(t x)\right|_{\mathfrak{h}}$.

Exercise 5.2.3. Let $\mathfrak{g}$ be the Heisenberg algebra. Determine a basis for $\mathfrak{g}$ with respect to which ad $\mathfrak{g}$ consists of upper triangular matrices.

Exercise 5.2.4. Let $\mathfrak{g}$ be a nilpotent Lie algebra and $\mathfrak{h}$ be a nonzero ideal in $\mathfrak{g}$. Show that the intersection of $\mathfrak{h}$ with the center of $\mathfrak{g}$ is not trivial.

Exercise 5.2.5. Show: If $\mathfrak{a}, \mathfrak{b}$ are nilpotent ideals of the Lie algebra $\mathfrak{g}$, then $\mathfrak{a}+\mathfrak{b}$ is also nilpotent.

Exercise 5.2.6. Give an example of a Lie algebra $\mathfrak{g}$ which contains a nilpotent ideal $\mathfrak{n}$ for which $\mathfrak{g} / \mathfrak{n}$ is nilpotent and $\mathfrak{g}$ is not nilpotent.

Exercise 5.2.7. For each Lie algebra $\mathfrak{g}$, we have

$$
\left[C^{n}(\mathfrak{g}), C^{m}(\mathfrak{g})\right] \subseteq C^{n+m}(\mathfrak{g}) \quad \text { for } \quad n, m \in \mathbb{N}
$$

Exercise 5.2.8. This exercise shows why Engel's Theorem does not generalize to infinite-dimensional spaces. We consider the vector space $V=\mathbb{K}^{(\mathbb{N})}$ with the basis $\left\{e_{i}: i \in \mathbb{N}\right\}$. In terms of the rank-one-operators $E_{i j} \in \operatorname{End}(V)$, defined by $E_{i j} e_{k}=\delta_{j k} e_{i}$, we consider the Lie algebra

$$
\mathfrak{g}:=\operatorname{span}\left\{E_{i j}: i>j\right\}
$$

(strictly lower triangular matrices). Show that:\\
(a) $C^{n}(\mathfrak{g})=\operatorname{span}\left\{E_{i j}: i \geq j+n\right\}, n \in \mathbb{N}$. In particular, we have $C^{\infty}(\mathfrak{g})= \bigcap_{n \in \mathbb{N}} C^{n}(\mathfrak{g})=\{0\}$, i.e., $\mathfrak{g}$ is residually nilpotent.\\
(b) $\mathfrak{g}$ consists of nilpotent endomorphisms of finite rank.\\
(c) $\mathfrak{z}(\mathfrak{g})=\{0\}$.\\
(d) $V^{\mathfrak{g}}=\{v \in V: \mathfrak{g} \cdot v=\{0\}\}=\{0\}$ (compare with Engel's Theorem).

\subsection*{The Jordan Decomposition}
In this section, we develop a tool that will be of crucial importance throughout the structure theory of Lie algebras: the Jordan decomposition of an endomorphism of a finite-dimensional vector space. Although the existence of the Jordan decomposition can be derived from the Jordan normal form, the proof of the Jordan decomposition is less involved because it does not specify the structure of the nilpotent component. Since we need various properties of the Jordan decomposition, we give a direct self-contained proof which does not require more than some elementary properties of polynomials.

Definition 5.3.1. Let $V$ be a vector space and $M \in \operatorname{End}(V)$.\\
(a) For $\lambda \in \mathbb{K}$, we define the eigenspace with respect to $\lambda$ as

$$
V_{\lambda}(M):=\operatorname{ker}(M-\lambda \mathbf{1})
$$

and the generalized eigenspace as

$$
V^{\lambda}(M):=\bigcup_{n \in \mathbb{N}} \operatorname{ker}(M-\lambda \mathbf{1})^{n}
$$

Note that the ascending sequence $\operatorname{ker}(M-\lambda \mathbf{1})^{n}$ is eventually constant if $V$ is finite-dimensional. We call $\lambda$ an eigenvalue if $V_{\lambda}(M) \neq\{0\}$.\\
(b) We call $M$ diagonalizable if $V=\bigoplus_{\lambda \in \mathbb{K}} V_{\lambda}(M)$, i.e., $V$ is a direct sum of the eigenspaces of $M$.\\
(c) We call $M$ nilpotent if there exists an $n \in \mathbb{N}$ with $M^{n}=0$. If $M$ is nilpotent, then $V=V^{0}(M)$.\\
(d) We call $M$ split if there is a nonzero polynomial $f \in \mathbb{K}[X]$ with $f(M)=0$ which decomposes as a product of linear factors. This is always the case for $\mathbb{K}=\mathbb{C}$.\\
(e) For $\mathbb{K}=\mathbb{R}$, we call $M$ semisimple if the endomorphism $M_{\mathbb{C}}$ of $V_{\mathbb{C}}$, defined by $M_{\mathbb{C}}(z \otimes v)=z \otimes M v$ is diagonalizable (cf. Exercise 5.3.5).

Lemma 5.3.2. For two commuting endomorphisms $M, N$ of the finite-dimensional vector space $V$, the following assertions hold:\\
(a) If $M$ and $N$ are diagonalizable, then $M+N$ is diagonalizable.\\
(b) If $M$ and $N$ are nilpotent, then $M+N$ is nilpotent.

Proof. (a) Exercise 2.1.1(a)-(c).\\
(b) Suppose that $M^{m}=N^{n}=0$. Then $[M, N]=0$ implies that

$$
(M+N)^{k}=\sum_{i+j=k}\binom{k}{i} M^{i} N^{j}
$$

If $k \geq n+m-1$, then either $i \geq m$ or $j \geq n$, so that all summands vanish. Hence $(M+N)^{k}=0$.

Theorem 5.3.3 (Jordan Decomposition Theorem). Let $V$ be a finitedimensional vector space and $M \in \operatorname{End}(V)$ a split endomorphism. Then there exists a diagonalizable endomorphism $M_{s}$ and a nilpotent endomorphism $M_{n}$ such that\\
(i) $M=M_{s}+M_{n}$.\\
(ii) $V^{\lambda}\left(M_{s}\right)=V_{\lambda}\left(M_{s}\right)=V^{\lambda}(M)$ for each $\lambda \in \mathbb{K}$.\\
(iii) There exist polynomials $P, Q \in \mathbb{K}[X]$ with $P(0)=Q(0)=0$ such that $M_{s}=P(M)$ and $M_{n}=Q(M)$.\\
(iv) If $L \in \operatorname{End}(V)$ commutes with $M$, then it also commutes with $M_{s}$ and $M_{n}$.\\
(v) (Uniqueness of the Jordan decomposition) If $S, N \in \operatorname{End}(V)$ commute, $S$ is diagonalizable, and $N$ nilpotent with $M=S+N$, then $S=M_{s}$ and $N=M_{n}$.

Proof. Let $f \in \mathbb{K}[X]$ be the minimal polynomial of $M$, i.e., a generator of the ideal $I_{M}:=\{f \in \mathbb{K}[X]: f(M)=0\}$ with leading coefficient 1 . By assumption, $I_{M}$ contains a nonzero polynomial which is a product of linear factors, so that Exercise 5.3.6 implies that $f$ also has this property. Hence\\
there exist pairwise different $\lambda_{1}, \ldots, \lambda_{m} \in \mathbb{K}$ and $k_{i} \in \mathbb{N}$ such that $f$ can be written as

$$
f=\left(X-\lambda_{1}\right)^{k_{1}}\left(X-\lambda_{2}\right)^{k_{2}} \cdots\left(X-\lambda_{m}\right)^{k_{m}} .
$$

Put $f_{i}:=f /\left(X-\lambda_{i}\right)^{k_{i}}$. Then the ideal

$$
I:=\left(f_{1}\right)+\cdots+\left(f_{m}\right) \subseteq \mathbb{K}[X]
$$

is generated by some element $g$ ( $\mathbb{K}[X]$ is a principal ideal domain, a simple consequence of Euclid's Algorithm) which is the greatest common divisor of the polynomials $f_{i}$. The fact that the $f_{1}, \ldots, f_{m}$ have no nontrivial common divisor (cf. Exercise 5.3.6) implies that $g$ is constant, so that $I=\mathbb{K}[X]$. Hence $1 \in I$, so that there exist polynomials $r_{1}, \ldots, r_{m} \in \mathbb{K}[X]$ with

$$
1=r_{1} f_{1}+\cdots+r_{m} f_{m}
$$

Put $E_{i}:=\left(r_{i} f_{i}\right)(M) \in \operatorname{End}(V)$ and note that $\sum_{i} E_{i}=\operatorname{id}_{V}$. If $i \neq j$, then $f$ divides $r_{i} f_{i} r_{j} f_{j}$, so that $f(M)=0$ leads to $E_{i} E_{j}=0$, and thus $E_{i}^{2}= E_{i}\left(\sum_{j=1}^{m} E_{j}\right)=E_{i}$. Therefore, the $E_{i}$ are pairwise commuting projections onto subspaces $V_{i}$ with $V=\bigoplus_{i=1}^{m} V_{i}$ (since $\sum_{i=1}^{m} E_{i}=\mathbf{1}$ ). Now $M_{s}:= \sum_{i=1}^{m} \lambda_{i} E_{i}$ is diagonalizable with $V_{i}=V_{\lambda_{i}}\left(M_{s}\right)$.

Since $M$ commutes with each $E_{i}$, it preserves the subspaces $V_{i}$, and therefore $f_{i}(M)$ preserves $V_{i}$. The relation

$$
\mathrm{id}_{V_{i}}=\left.E_{i}\right|_{V_{i}}=\left.r_{i}(M) f_{i}(M)\right|_{V_{i}}
$$

shows that the restriction of $f_{i}(M)$ to $V_{i}$ is invertible. Therefore, $f(M)=0$ leads to

$$
\left(M-\lambda_{i} \mathbf{1}\right)^{k_{i}}\left(V_{i}\right)=\left(M-\lambda_{i} \mathbf{1}\right)^{k_{i}} f_{i}(M)\left(V_{i}\right)=f(M)\left(V_{i}\right)=\{0\}
$$

i.e., $V_{i} \subseteq V^{\lambda_{i}}(M)$.\\
With $M_{n}:=M-M_{s}$ and $k_{0}:=\max \left\{k_{i}: i=1, \ldots, m\right\}$ we finally get $M_{n}^{k_{0}}=0$, proving (i).\\
(ii) We have to show that $V_{i}=V^{\lambda_{i}}(M)$. We know already that $V_{i} \subseteq V^{\lambda_{i}}(M)$. So let $v \in V^{\lambda_{i}}(M)$ and write $v$ as $v=\sum_{j=1}^{m} v_{j}$ with $v_{j} \in V_{j}$. Then the invariance of $V_{j}$ under $M$ implies that $v_{j} \in V^{\lambda_{i}}(M)$. If $v_{j} \neq 0$, then there exists a nonzero eigenvector $v_{j}^{\prime} \in V_{\lambda_{i}}(M) \cap V_{j}$ (put $v_{j}^{\prime}=\left(M-\lambda_{i} \mathbf{1}\right)^{k} v_{j}$, where $k$ is maximal with the property that this vector is nonzero). Then $\left(M-\lambda_{j} \mathbf{1}\right)^{n_{j}} v_{j}^{\prime}=\left(\lambda_{i}-\lambda_{j}\right)^{n_{j}} v_{j}^{\prime}=0$, hence $\lambda_{j}=\lambda_{i}$, i.e., $j=i$. This implies that $v=v_{i} \in V_{i}$, and therefore $V^{\lambda_{i}}(M)=V_{i}$.\\
(iii) By construction, $M_{s}=P_{1}(M)$ and $M_{n}=Q_{1}(M)$ for $P_{1}=\sum_{i} \lambda_{i} r_{i} f_{i}$ and $Q_{1}=X-P_{1}$. It remains to be seen that these polynomials can be chosen with trivial constant term. If one eigenvalue $\lambda_{j}$ vanishes, then $\{0\} \neq V_{0}:=$ ker $M \subseteq V_{j}$ and $\left.M_{s}\right|_{V_{0}}=0$ implies that $P_{1}$ has no constant term. Then $Q_{1}=X-P_{1}$ likewise has no constant term and (iii) holds with $P:=P_{1}$ and $Q:=Q_{1}$.

If all eigenvalues $\lambda_{j}$ are nonzero, then $f(0) \neq 0$ and (iii) holds with

$$
P:=P_{1}-\frac{P_{1}(0)}{f(0)} f \quad \text { and } \quad Q:=Q_{1}-\frac{Q_{1}(0)}{f(0)} f .
$$

(iv) is a direct consequence of (iii).\\
(v) Since $N$ and $S$ commute with $M=N+S$, (iii) shows that they both commute with $M_{s}$ and $M_{n}$. Then Lemma 5.3.2 shows that

$$
S-M_{s}=M_{n}-N
$$

is nilpotent as well as diagonalizable, which leads to $0=S-M_{s}= M_{n}-N$.

Definition 5.3.4. The decomposition $M=M_{s}+M_{n}$ is called the Jordan decomposition of $M, M_{s}$ is called the semisimple Jordan component and $M_{n}$ the nilpotent Jordan component of $M$.

Example 5.3.5. If $M$ is a Jordan block $\left(\begin{array}{cc}\lambda & 1 \\ 0 & \lambda\end{array}\right)$, then the Jordan decomposition is

$$
M=\underbrace{\left(\begin{array}{ll}
\lambda & 0 \\
0 & \lambda
\end{array}\right)}_{M_{s}}+\underbrace{\left(\begin{array}{ll}
0 & 1 \\
0 & 0
\end{array}\right)}_{M_{n}}
$$

The matrix $M=\left(\begin{array}{ll}1 & 2 \\ 0 & 3\end{array}\right)$ is diagonalizable, and therefore $M=M_{s}$. In this case,

$$
M=\left(\begin{array}{ll}
1 & 2 \\
0 & 3
\end{array}\right)=\left(\begin{array}{ll}
1 & 0 \\
0 & 3
\end{array}\right)+\left(\begin{array}{ll}
0 & 2 \\
0 & 0
\end{array}\right)
$$

is not the Jordan decomposition, even though the first summand is diagonalizable and the second summand is nilpotent. These summands do not commute.

The preceding theorem does not apply to all endomorphisms of real vector spaces. We now explain how this problem can be overcome, so that we also obtain a Jordan decomposition for endomorphisms of real vector spaces.

Definition 5.3.6 (Jordan decomposition in the real case). If $V$ is a finite-dimensional real vector space and $M \in \operatorname{End}(V)$, then $M_{\mathbb{C}} \in \operatorname{End}\left(V_{\mathbb{C}}\right)$, defined by $M_{\mathbb{C}}(z \otimes v):=z \otimes M v$ has a Jordan decomposition

$$
M_{\mathbb{C}}=M_{\mathbb{C}, s}+M_{\mathbb{C}, n}
$$

Let $\sigma: V_{\mathbb{C}} \rightarrow V_{\mathbb{C}}$ be the antilinear map defined by $\sigma(z \otimes v):=\bar{z} \otimes v$ and $\underline{\text { define for any complex linear } A \in \operatorname{End}\left(V_{\mathbb{C}}\right) \text { the complex linear endomorphism }} \bar{A}:=\sigma \circ A \circ \sigma \in \operatorname{End}\left(V_{\mathbb{C}}\right)$. Then $\overline{M_{\mathbb{C}}}=M_{\mathbb{C}}$ leads to

$$
M_{\mathbb{C}}=\overline{M_{\mathbb{C}}}=\overline{M_{\mathbb{C}, s}}+\overline{M_{\mathbb{C}, n}}
$$

where the summands on the right commute, the first is diagonalizable and the second is nilpotent (Exercise). Hence the uniqueness of the Jordan decomposition yields

$$
\overline{M_{\mathbb{C}, s}}=M_{\mathbb{C}, s} \quad \text { and } \quad \overline{M_{\mathbb{C}, n}}=M_{\mathbb{C}, n}
$$

In view of Exercise 5.3.1, this implies the existence of $M_{s} \in \operatorname{End}(V)$ and $M_{n} \in \operatorname{End}(V)$, with

$$
\left(M_{s}\right)_{\mathbb{C}}=M_{\mathbb{C}, s} \quad \text { and } \quad\left(M_{n}\right)_{\mathbb{C}}=M_{\mathbb{C}, n} .
$$

Then $M=M_{s}+M_{n}$, and this is called the Jordan decomposition of $M$. It is uniquely characterized by the properties that $\left[M_{s}, M_{n}\right]=0, M_{s}$ is semisimple and $M_{n}$ is nilpotent (Exercise).

Proposition 5.3.7 (Properties of the Jordan decomposition). Let $V$ be a finite-dimensional vector space and $M \in \operatorname{End}(V)$.\\
(i) If $M^{\prime} \in \operatorname{End}\left(V^{\prime}\right)$ and $f: V \rightarrow V^{\prime}$ satisfy $f \circ M=M^{\prime} \circ f$, then

$$
f \circ M_{s}=M_{s}^{\prime} \circ f \quad \text { and } \quad f \circ M_{n}=M_{n}^{\prime} \circ f
$$

(ii) If $W \subseteq V$ is an $M$-invariant subspace, then

$$
\left(\left.M\right|_{W}\right)_{s}=\left.M_{s}\right|_{W} \quad \text { and } \quad\left(\left.M\right|_{W}\right)_{n}=\left.M_{n}\right|_{W}
$$

In particular, $W$ is invariant under $M_{s}$ and $M_{n}$. If $\bar{M}$ denotes the induced endomorphism of $V / W$, then

$$
(\bar{M})_{s}=\overline{M_{s}} \quad \text { and } \quad(\bar{M})_{n}=\overline{M_{n}}
$$

(iii) If $U \subseteq W$ are subspaces of $V$ with $M W \subseteq U$, then $M_{s} W \subseteq U$ and $M_{n} W \subseteq U$.

Proof. (i) Let $W:=V \oplus V^{\prime}, L:=M \oplus M^{\prime}$, and consider the linear map $\varphi: W \rightarrow W$ defined by $\varphi\left(v, v^{\prime}\right)=(0, f(v))$. Then $\varphi \circ L=L \circ \varphi$ and thus $\varphi L_{s}=L_{s} \varphi$ and $\varphi L_{n}=L_{n} \varphi$. Further, $L_{s}=M_{s} \oplus M_{s}^{\prime}$ and $L_{n}=M_{n} \oplus M_{n}^{\prime}$ follows from the uniqueness of the Jordan decomposition (Theorem 5.3.3(v)) and the semisimplicity of $M_{s} \oplus M_{s}^{\prime}$. This shows that

$$
M_{s}^{\prime} \circ f=f \circ M_{s} \quad \text { and } \quad M_{n}^{\prime} \circ f=f \circ M_{n} .
$$

(ii) Apply (i) to the inclusion $j: W \rightarrow V$ and the quotient map $p: V \rightarrow V / W$.\\
(iii) For $\mathbb{K}=\mathbb{C}$, this follows from Theorem 5.3.3(iii) and the real case is obtained by complexification.

Proposition 5.3.8. Let $V$ be finite-dimensional vector space and $x \in \mathfrak{g} \mathfrak{l}(V)$. If $x$ is nilpotent (diagonalizable), then so is ad $x$.

This proposition can be obtained by combining Lemma 5.2.4(i) with Exercise 5.2.1. We give an alternative proof using the Jordan decomposition.

Proof. Put $L_{x}: \mathfrak{g} \mathfrak{l}(V) \rightarrow \mathfrak{g} \mathfrak{l}(V), y \mapsto x y$ and $R_{x}: \mathfrak{g} \mathfrak{l}(V) \rightarrow \mathfrak{g} \mathfrak{l}(V), y \mapsto y x$. Then $\operatorname{ad} x=L_{x}-R_{x}$ and $\left[L_{x}, R_{x}\right]=0$. In view of Lemma 5.3.2, it suffices to see that $L_{x}$ and $R_{x}$ are nilpotent, resp., diagonalizable, whenever $x$ has this property.

If $x^{n}=0$, then $L_{x}^{n}=L_{x^{n}}=0=R_{x}^{n}$. If $x$ is diagonalizable and $\lambda_{1}, \ldots, \lambda_{n}$ are the eigenvalues of $x$, we consider the decomposition $V=\bigoplus_{j=1}^{n} V_{\lambda_{j}}(x)$. We write any $y \in \operatorname{End}(V)$ as $y=\sum_{j, k=1}^{n} y_{j k}$ with $y_{j k} V_{\lambda_{k}}(x) \subseteq V_{\lambda_{j}}(x)$. Then $L_{x} y_{j k}=\lambda_{j} y_{j k}$ and $R_{x} y_{j k}=\lambda_{k} y_{j k}$ imply that $L_{x}$ and $R_{x}$ are diagonalizable on $\mathfrak{g l}(V)=\operatorname{End}(V)$.

Corollary 5.3.9. For each endomorphism $x \in \mathfrak{g} \mathfrak{l}(V)$ of the finite-dimensional vector space $V$, the Jordan decomposition of $\operatorname{ad} x$ is given by

$$
\operatorname{ad} x=\operatorname{ad}\left(x_{s}\right)+\operatorname{ad}\left(x_{n}\right) .
$$

Proof. Proposition 5.3.8 implies for $\mathbb{K}=\mathbb{C}$ that $\operatorname{ad}\left(x_{s}\right)$ is diagonalizable and for $\mathbb{K}=\mathbb{R}$ that $\left(\operatorname{ad} x_{s}\right)_{\mathbb{C}}=\operatorname{ad}\left(\left(x_{s}\right)_{\mathbb{C}}\right)$ is diagonalizable. Further, $\operatorname{ad}\left(x_{n}\right)$ is nilpotent, and $\left[\operatorname{ad}\left(x_{s}\right), \operatorname{ad}\left(x_{n}\right)\right]=\operatorname{ad}\left[x_{s}, x_{n}\right]=0$, so that the assertion follows from the uniqueness of the Jordan decomposition of ad $x$.

Proposition 5.3.10. If $A$ is a finite-dimensional algebra and $D \in \operatorname{der}(A)$, then the Jordan components $D_{s}$ and $D_{n}$ are also derivations of $A$.

Proof. First proof: Let $m: A \otimes A \rightarrow A$ denote the linear map defined by the algebra multiplication. Then $D \in \operatorname{der}(A)$ is equivalent to the relation

$$
D \circ m=m \circ\left(D \otimes \mathrm{id}_{A}+\mathrm{id}_{A} \otimes D\right)
$$

Next we observe that

$$
D \otimes \mathrm{id}_{A}+\mathrm{id}_{A} \otimes D=\left(D_{s} \otimes \mathrm{id}_{A}+\mathrm{id}_{A} \otimes D_{s}\right)+\left(D_{n} \otimes \mathrm{id}_{A}+\mathrm{id}_{A} \otimes D_{n}\right)
$$

is the Jordan decomposition (Exercise!), so that Proposition 5.3.7 implies that

$$
D_{s} \circ m=m \circ\left(D_{s} \otimes \mathrm{id}_{A}+\mathrm{id}_{A} \otimes D_{s}\right)
$$

which means that $D_{s} \in \operatorname{der}(A)$, and hence that $D_{n}=D-D_{s} \in \operatorname{der}(A)$ because $\operatorname{der}(A)$ is a linear space.

Second proof: Since $\operatorname{der}(A)$ is a vector space, it suffices to show that $D_{s} \in \operatorname{der}(A)$. Furthermore, $D \in \operatorname{der}(A)$ is equivalent to $D_{\mathbb{C}} \in \operatorname{der}\left(A_{\mathbb{C}}\right)$, so that we may assume that $\mathbb{K}=\mathbb{C}$.

For $a, b \in A$ and $\lambda, \mu \in \mathbb{K}$, we have for all $n \in \mathbb{N}$ the formula

$$
(D-(\lambda+\mu) \mathbf{1})^{n}(a b)=\sum_{k=0}^{n}\binom{n}{k}(D-\lambda \mathbf{1})^{k}(a) \cdot(D-\mu \mathbf{1})^{n-k}(b)
$$

(Exercise 5.3.7). It follows that for $a \in A_{\lambda}\left(D_{s}\right)=A^{\lambda}(D)$ and $b \in A_{\mu}\left(D_{s}\right)= A^{\mu}(D)$, we have $a b \in A^{\lambda+\mu}(D)=A_{\lambda+\mu}\left(D_{s}\right)$. Furthermore,

$$
D_{s}(a) b+a D_{s}(b)=\lambda a b+\mu a b=(\lambda+\mu) a b=D_{s}(a b)
$$

Since $A=\sum_{\lambda \in \mathbb{K}} A_{\lambda}\left(D_{s}\right)$, it follows that $D_{s} \in \operatorname{der}(A)$.\\
Lemma 5.3.11 (Fitting decomposition). Let $V$ be a finite-dimensional vector space and $T \in \operatorname{End}(V)$. If $V^{+}(T):=\bigcap_{n \in \mathbb{N}} T^{n}(V)$, then

$$
V=V^{0}(T) \oplus V^{+}(T)
$$

Proof. The sequence $T^{n}(V)$ is decreasing and $\operatorname{dim} V<\infty$ implies that there exists some $n$ with $T^{n+1}(V)=T^{n}(V)$. Then

$$
\left.T\right|_{T^{n}(V)}: T^{n}(V) \rightarrow T^{n}(V)
$$

is surjective, hence bijective, and we see that $V^{+}(T)=V^{n}(T)$. On the intersection $V^{0}(T) \cap V^{+}(T)$, the restriction of $T$ is nilpotent and bijective at the same time. This is only possible if the intersection is $\{0\}$. Therefore, it remains to see that $V=V^{0}(T)+V^{+}(T)$.

First, we assume that $T$ is split. We consider the generalized eigenspace decomposition

$$
V=V^{0}(T) \oplus \bigoplus_{\lambda \neq 0} V^{\lambda}(T)
$$

and note that each $V^{\lambda}(T)$ is $T$-invariant. For $\lambda \neq 0$, we have $V^{\lambda}(T) \cap \operatorname{ker} T= \{0\}$, so that $\left.T\right|_{V^{\lambda}(T)}$ is injective, hence surjective. This implies that

$$
\bigoplus_{\lambda \neq 0} V^{\lambda}(T) \subseteq V^{+}(T)
$$

and thus $V=V^{0}(T)+V^{+}(T)$.\\
If $T$ is not split, then $\mathbb{K}=\mathbb{R}$, and we consider the complexification $T_{\mathbb{C}} \in \operatorname{End}\left(V_{\mathbb{C}}\right)$. Then the preceding argument implies that

$$
V_{\mathbb{C}}=V^{0}\left(T_{\mathbb{C}}\right) \oplus V^{+}\left(T_{\mathbb{C}}\right)=V^{0}(T)_{\mathbb{C}} \oplus V^{+}(T)_{\mathbb{C}}
$$

where we use $T^{n}(V)_{\mathbb{C}}=T_{\mathbb{C}}^{n}\left(V_{\mathbb{C}}\right)$ for $n \in \mathbb{N}$ in the second equality. We conclude that $\left(V^{0}(T)+V^{+}(T)\right)_{\mathbb{C}}=V_{\mathbb{C}}$, and this entails that $V^{0}(T)+V^{+}(T)=V$.

The space $V^{+}(T)$ is called the Fitting one component. In this context, the generalized eigenspace $V^{0}(T)$ is called the Fitting null component.

\subsubsection*{Exercises for Section 5.3}
Exercise 5.3.1. Let $V$ be a real vector space and

$$
V_{\mathbb{C}}=\mathbb{C} \otimes_{\mathbb{R}} V=(1 \otimes V) \oplus(i \otimes V)
$$

its complexification. We identify $V$ with the real subspace $1 \otimes V$, so that

$$
V_{\mathbb{C}} \cong V \oplus i V
$$

Show that:\\
(i) $\sigma(z \otimes v):=\bar{z} \otimes v$ defines an antilinear involution of $V_{\mathbb{C}}$ whose fixed point space is $V$.\\
(ii) A complex subspace $E \subseteq V_{\mathbb{C}}$ is of the form $W_{\mathbb{C}}$ for some real subspace of $V$ if and only if $\sigma(E)=E$.\\
(iii) For each $M \in \operatorname{End}(V)$, the complexification $M_{\mathbb{C}} \in \operatorname{End}\left(V_{\mathbb{C}}\right)$, defined by $M_{\mathbb{C}}(z \otimes v):=z \otimes M v$ commutes with $\sigma$.\\
(iv) For $A \in \operatorname{End}\left(V_{\mathbb{C}}\right)$ the following are equivalent\\
(a) $A$ commutes with $\sigma$.\\
(b) $A$ preserves the real subspace $V$.\\
(c) $A=M_{\mathbb{C}}$ for some $M \in \operatorname{End}(V)$.

Exercise 5.3.2. Let $V$ be a complex vector space and $M \in \operatorname{End}(V)$. Show that $M$ is diagonalizable if and only if each $M$-invariant subspace $W \subseteq V$ possesses an $M$-invariant complement.

Exercise 5.3.3. Let $V$ be a real vector space, $A \in \operatorname{End}(V)$ and $z \in V_{\mathbb{C}}$ an eigenvector of $A_{\mathbb{C}}$ with respect to the eigenvalue $\lambda$. Show that if $z=x+i y$ with $x, y \in V$ and $\lambda=a+i b$, then

$$
A x=a x-b y \quad \text { and } \quad A y=a y+b x .
$$

In particular, the 2 -dimensional subspace $E:=\operatorname{span}\{x, y\} \subseteq V$ is invariant under $A$.

Exercise 5.3.4. Let $A \in M_{2}(\mathbb{R})$ with no real eigenvalue. Then there exists a basis $x, y \in \mathbb{R}^{2}$ and $a, b \in \mathbb{R}$ with

$$
A x=a x-b y \quad \text { and } \quad A y=a y+b x .
$$

Exercise 5.3.5. Let $V$ be a real vector space and $M \in \operatorname{End}(V)$. Show that $M$ is semisimple if and only if each $M$-invariant subspace $W \subseteq V$ possesses an $M$-invariant complement.\\
Exercise 5.3.6. Let $f \in \mathbb{K}[X]$ be a polynomial of the form

$$
f=\left(X-\lambda_{1}\right)^{k_{1}}\left(X-\lambda_{2}\right)^{k_{2}} \cdots\left(X-\lambda_{m}\right)^{k_{m}}
$$

and $g \in \mathbb{K}[X]$ a divisor of $f$ with leading coefficient 1 . Show that there exist $\ell_{i} \leq k_{i}$ with

$$
g=\left(X-\lambda_{1}\right)^{\ell_{1}}\left(X-\lambda_{2}\right)^{\ell_{2}} \cdots\left(X-\lambda_{m}\right)^{\ell_{m}}
$$

Exercise 5.3.7. Show that for each algebra $A$, a derivation $D \in \operatorname{der}(A)$ and $\lambda, \mu \in \mathbb{K}$, we have for $a, b \in A$ :

$$
(D-(\lambda+\mu) \mathbf{1})^{n}(a b)=\sum_{k=0}^{n}\binom{n}{k}(D-\lambda \mathbf{1})^{k}(a) \cdot(D-\mu \mathbf{1})^{n-k}(b) .
$$

Exercise 5.3.8. Let $V$ be a finite-dimensional vector space over $\mathbb{K}$ and $A \in \operatorname{End}(V)$. Then the multiplicity of the root 0 of its characteristic polynomial $\operatorname{det}(A-X \mathbf{1}) \in \mathbb{K}[X]$ coincides with $\operatorname{dim} V^{0}(A)$.

\subsection*{Solvable Lie Algebras}
In this section, we turn to the class of solvable Lie algebras. They are defined in a similar fashion as nilpotent ones, and indeed every nilpotent Lie algebra is solvable. The central results on solvable Lie algebras are Lie's Theorem on representations of solvable Lie algebras and Cartan's Solvability Criterion in terms of vanishing of

$$
\operatorname{tr}(\operatorname{ad}[x, y] \operatorname{ad} z) \quad \text { for } \quad x, y, z \in \mathfrak{g} .
$$

As we shall see later on, similar techniques apply to semisimple Lie algebras.\\
Definition 5.4.1. (a) Let $\mathfrak{g}$ be a Lie algebra. The derived series of $\mathfrak{g}$ is defined by

$$
D^{0}(\mathfrak{g}):=\mathfrak{g} \quad \text { and } \quad D^{n}(\mathfrak{g}):=\left[D^{n-1}(\mathfrak{g}), D^{n-1}(\mathfrak{g})\right]
$$

for $n \in \mathbb{N}$.\\
From $D^{1}(\mathfrak{g}) \subseteq \mathfrak{g}$, we inductively see that $D^{n}(\mathfrak{g}) \subseteq D^{n-1}(\mathfrak{g})$. Further, an easy induction shows that all $D^{n}(\mathfrak{g})$ are ideals of $\mathfrak{g}$ (Exercise 5.1.9). The derived series is a descending series of ideals.\\
(b) The Lie algebra $\mathfrak{g}$ is said to be solvable, if there exists an $n \in \mathbb{N}$ with $D^{n}(\mathfrak{g})=\{0\}$.

Example 5.4.2. (i) The oscillator algebra is solvable, but not nilpotent.\\
(ii) Every nilpotent Lie algebra is solvable because $D^{n}(\mathfrak{g}) \subseteq C^{n+1}(\mathfrak{g})$ follows easily by induction.\\
(iii) Consider $\mathbb{R}$ and $\mathbb{C}$ as abelian real Lie algebras and write $I \in \operatorname{End}_{\mathbb{R}}(\mathbb{C})$ for the multiplication with $i$. Then the Lie algebra $\mathbb{C} \rtimes_{I} \mathbb{R}$ is solvable, but not nilpotent. It is isomorphic to $\mathfrak{g} / \mathfrak{z}(\mathfrak{g})$, where $\mathfrak{g}$ is the oscillator algebra.\\
(iv) Let $\mathfrak{g}$ be a 2-dimensional nonabelian Lie algebra with basis $x, y$. Then $0 \neq[x, y]=a x+b y$ for $(a, b) \neq(0,0)$. Assume w.lo.g. that $b \neq 0$ and put $v:=b^{-1} x$ and $w:=a x+b y$. Then $[v, w]=w$. This implies in particular that all nonabelian two-dimensional Lie algebras are isomorphic because they have a basis ( $v, w$ ) with $[v, w]=w$. Then $D^{1}(\mathfrak{g})=\mathbb{K} w$ and $D^{2}(\mathfrak{g})=\{0\}$, so that $\mathfrak{g}$ is solvable. On the other hand, $C^{n}(\mathfrak{g})=\mathbb{K} w$ for each $n>1$, so that $\mathfrak{g}$ is not nilpotent.

A natural matrix realization of this Lie algebra is\\
$\mathfrak{a f f}_{1}(\mathbb{K})=\left(\begin{array}{cc}\mathbb{K} & \mathbb{K} \\ 0 & 0\end{array}\right) \quad$ with the basis $\quad v:=\left(\begin{array}{ll}1 & 0 \\ 0 & 0\end{array}\right), \quad w:=\left(\begin{array}{ll}0 & 1 \\ 0 & 0\end{array}\right)$.\\
Proposition 5.4.3. For a Lie algebra $\mathfrak{g}$, the following assertions hold:\\
(i) If $\mathfrak{g}$ is solvable, then all subalgebras and homomorphic images of $\mathfrak{g}$ are solvable.\\
(ii) Solvability is an extension property: If $\mathfrak{i}$ is a solvable ideal of $\mathfrak{g}$ and $\mathfrak{g} / \mathfrak{i}$ is solvable, then $\mathfrak{g}$ is solvable.\\
(iii) If $\mathfrak{i}$ and $\mathfrak{j}$ are solvable ideals of $\mathfrak{g}$, then the ideal $\mathfrak{i}+\mathfrak{j}$ is solvable.\\
(iv) If $\mathfrak{i} \unlhd \mathfrak{g}$ is an ideal, then the $D^{n}(\mathfrak{i})$ are ideals in $\mathfrak{g}$.

Proof. (i) If $\mathfrak{h} \subseteq \mathfrak{g}$ is a subalgebra, then $D^{n}(\mathfrak{h}) \subseteq D^{n}(\mathfrak{g})$ follows by induction, and if $\alpha: \mathfrak{g} \rightarrow \mathfrak{h}$ is a homomorphism of Lie algebras, then we obtain

$$
D^{n}(\alpha(\mathfrak{g}))=\alpha\left(D^{n}(\mathfrak{g})\right)
$$

by induction. This implies (i).\\
(ii) Let $\pi: \mathfrak{g} \rightarrow \mathfrak{g} / \mathfrak{i}$ be the quotient map. We have already seen in (i) that $\pi\left(D^{n}(\mathfrak{g})\right)=D^{n}(\pi(\mathfrak{g}))$ for each $n$. If $\mathfrak{g} / \mathrm{i}$ is solvable, then $\pi\left(D^{n}(\mathfrak{g})\right)$ vanishes for some $n \in \mathbb{N}$. Now $D^{n}(\mathfrak{g}) \subseteq \operatorname{ker} \pi=\mathfrak{i}$, so that $D^{n+k}(\mathfrak{g}) \subseteq D^{k}(\mathfrak{i})$ for each $k \in \mathbb{N}$. If $\mathfrak{i}$ is also solvable, we immediately derive that $\mathfrak{g}$ is solvable.\\
(iii) The ideal $\mathfrak{j}$ of $\mathfrak{i}+\mathfrak{j}$ is solvable and $(\mathfrak{i}+\mathfrak{j}) / \mathfrak{j} \cong \mathfrak{i} /(\mathfrak{i} \cap \mathfrak{j})$ (Theorem 5.1.17(iii)) is solvable by (i). Hence (ii) implies that $\mathfrak{i}+\mathfrak{j}$ is solvable.\\
(iv) We only have to observe that for each ideal $\mathfrak{i}$, its commutator algebra $[\mathrm{i}, \mathrm{i}]$ also is an ideal (Exercise 5.1.9). Then (iv) follows by induction.\\
Example 5.4.4. If $\mathcal{F}=\left(V_{0}, \ldots, V_{n}\right)$ is a complete flag in the $n$-dimensional vector space $V$, then $\mathfrak{g}(\mathcal{F})$ is a solvable Lie algebra. In fact,

$$
\mathfrak{g}(\mathcal{F}) \cong \mathfrak{g}_{n}(\mathcal{F}) \rtimes \mathfrak{g} \mathfrak{l}_{1}(\mathbb{K})^{n} \cong \mathfrak{g}_{n}(\mathcal{F}) \rtimes \mathbb{K}^{n}
$$

(Example 5.1.20).\\
Since $\mathbb{K}^{n} \cong \mathfrak{g}(\mathcal{F}) / \mathfrak{g}_{n}(\mathcal{F})$ is abelian and $\mathfrak{g}_{n}(\mathcal{F})$ nilpotent (Example 5.2.2(iii)), the solvability of $\mathfrak{g}(\mathcal{F})$ follows from Proposition 5.4.3(ii). Below we shall see that Lie's Theorem provides a converse for solvable subalgebras of $\mathfrak{g} \mathfrak{l}(V)$ and $\mathbb{K}=\mathbb{C}$; they are always contained in some $\mathfrak{g}(\mathcal{F})$ for a complete flag $\mathcal{F}$.

Definition 5.4.5. Proposition 5.4.3(iii) shows that in every finite-dimensional Lie algebra $\mathfrak{g}$, there is a largest solvable ideal. This ideal is called the radical of $\mathfrak{g}$, and is denoted by $\mathfrak{r a d}(\mathfrak{g})$.

Remark 5.4.6. The analog of Proposition 5.4.3(iii) for nilpotent Lie algebras is also valid (Exercise 5.2.5). Note that nilpotency is not an extension property, i.e., the analog of Proposition 5.4.3(ii) is false for nilpotent Lie algebras (cf. Example 5.4.2(iv)).

\subsubsection*{Lie's Theorem}
Now we turn to solvable Lie subalgebras $\mathfrak{g}$ of $\mathfrak{g} \mathfrak{l}(V)$. In this context, we do not want to make any assumption on the elements of $\mathfrak{g}$, as in Corollary 5.2.7.

Theorem 5.4.7. Let $V$ be a nonzero finite-dimensional complex vector space and $\mathfrak{g}$ be a solvable subalgebra of $\mathfrak{g} \mathfrak{l}(V)$. Then there exists a nonzero common eigenvector $v$ for $\mathfrak{g}$, i.e., $\mathfrak{g}(v) \subseteq \mathbb{C} v$.

Proof. We may w.lo.g. assume that $\mathfrak{g} \neq\{0\}$. We proceed by induction on the dimension of $\mathfrak{g}$. If $\mathfrak{g}=\mathbb{C} x$, then every eigenvector of $x$ (and such an eigenvector always exists) satisfies the requirement of the theorem. So let $\operatorname{dim}_{\mathbb{C}} \mathfrak{g}>1$ and $\mathfrak{h}$ be a complex hyperplane in $\mathfrak{g}$ which contains $[\mathfrak{g}, \mathfrak{g}]=D^{1}(\mathfrak{g})$. Here we use that $D^{1}(\mathfrak{g})$ is a proper subspace because $\mathfrak{g}$ is solvable. In view of $[\mathfrak{g}, \mathfrak{g}] \subseteq \mathfrak{h}$, the subspace $\mathfrak{h}$ is an ideal of $\mathfrak{g}$. Now the induction hypothesis provides a nonzero common eigenvector $v$ for $\mathfrak{h}$. If $x(v)=\lambda(x) v$ for $x \in \mathfrak{h}$, then $\lambda: \mathfrak{h} \rightarrow \mathbb{C}$ is a linear map and

$$
v \in V_{\lambda}(\mathfrak{h}):=\{w \in V \mid(\forall x \in \mathfrak{h}) x(w)=\lambda(x) w\} .
$$

Suppose that $V_{\lambda}(\mathfrak{h})$ is $\mathfrak{g}$-invariant and pick $y \in \mathfrak{g} \backslash \mathfrak{h}$. Then there exists a nonzero eigenvector $v_{o} \in V_{\lambda}(\mathfrak{h})$ for $y$. Then $v_{o}$ is a common eigenvector for $\mathfrak{g}=\mathfrak{h}+\mathbb{C} y$ and the proof is complete.

It remains to show that $V_{\lambda}(\mathfrak{h})$ is $\mathfrak{g}$-invariant. For this, we calculate as in the proof of Theorem 5.2.5:

$$
y x(w)=x y(w)-[x, y](w)=\lambda(y) x(w)-\lambda([x, y])(w)
$$

for $w \in V_{\lambda}(\mathfrak{h}), x \in \mathfrak{g}$ and $y \in \mathfrak{h}$. Hence it suffices to show that $[\mathfrak{g}, \mathfrak{h}] \subseteq \operatorname{ker} \lambda$. For fixed $w \in V_{\lambda}(\mathfrak{h}), x \in \mathfrak{g}$ and $k \in \mathbb{N}$, we consider the space

$$
W^{k}=\mathbb{C} w+\mathbb{C} x(w)+\cdots+\mathbb{C} x^{k}(w)
$$

Since

$$
y x^{k}(w)=x y\left(x^{k-1} w\right)-[x, y]\left(x^{k-1} w\right)
$$

and $y(w)=\lambda(y) w$ for $y \in \mathfrak{h}$, we see by induction on $k$ that $\mathfrak{h}\left(W^{k}\right) \subseteq W^{k}$ for each $k \in \mathbb{N}$.

Now we choose $k_{o} \in \mathbb{N}$ maximal with respect to the property that

$$
\left\{w, x(w), \ldots, x^{k_{o}}(w)\right\}
$$

is a basis for $W^{k_{o}}$. Then $W^{k_{o}+m}=W^{k_{o}}$ for all $m \in \mathbb{N}$, and

$$
\mathcal{F}=\left(\{0\}, W^{1}, \ldots, W^{k_{o}}\right)
$$

is a complete flag in $W^{k_{o}}$ which is invariant under $\mathfrak{h}$. Thus, every $y \in \mathfrak{h}$ corresponds to an upper triangular matrix ( $y_{i j}$ ) with respect to the above\\
basis for $W^{k_{o}}$. The diagonal entries $y_{i i}$ of this matrix are all equal to $\lambda(y)$ since $y(w)=\lambda(y) w$ and (5.5) imply by induction that

$$
y x^{k}(w) \in \lambda(y) x^{k}(w)+W^{k-1} .
$$

Since $x$ and $y$ leave the space $W^{k_{o}}$ invariant, we have

$$
\left.[x, y]\right|_{W^{k_{o}}}=\left[\left.x\right|_{W^{k_{o}}},\left.y\right|_{W^{k_{o}}}\right]
$$

In particular, $\left.[x, y]\right|_{W^{k_{o}}}$ is a commutator of two endomorphisms so that its trace vanishes. Finally, $[x, y] \in \mathfrak{h}$ leads to

$$
0=\operatorname{tr}\left(\left.[x, y]\right|_{W^{k_{o}}}\right)=\left(k_{o}+1\right) \lambda([x, y]),
$$

so that $\lambda([x, y])=0$.\\
Theorem 5.4.8 (Lie's Theorem). Let $V$ be a finite-dimensional complex vector space and $\mathfrak{g}$ be a solvable subalgebra of $\mathfrak{g} \mathfrak{l}(V)$. Then there exists a complete $\mathfrak{g}$-invariant flag in $V$.

Proof. We may assume that $V$ is nonzero. By Theorem 5.4.7, there is a nonzero common $\mathfrak{g}$-eigenvector $v_{1} \in V$. Put $V_{1}:=\mathbb{C} v_{1}$. Then

$$
\alpha: \mathfrak{g} \rightarrow \mathfrak{g} \mathfrak{l}\left(V / V_{1}\right), \quad \alpha(x)\left(v+V_{1}\right):=x(v)+V_{1}
$$

defines a representation of $\mathfrak{g}$ on the quotient space $V / V_{1}$ (Exercise 5.1.10) and $\alpha(\mathfrak{g})$ is solvable. Proceeding by induction on $\operatorname{dim} V$, we may assume that there exists an $\alpha(\mathfrak{g})$-invariant complete flag in $V / V_{1}$, and the preimage in $V$, together with $\{0\}$, is a complete $\mathfrak{g}$-invariant flag in $V$.

Remark 5.4.9. If we apply Lie's Theorem 5.4.8 to $V=\mathfrak{g}$ and $\operatorname{ad}(\mathfrak{g})$, where $\mathfrak{g}$ is solvable, we get a complete flag of ideals

$$
\{0\}=\mathfrak{g}_{0}<\mathfrak{g}_{1}<\cdots<\mathfrak{g}_{n}=\mathfrak{g}
$$

of $\mathfrak{g}$ with $\operatorname{dim} \mathfrak{g}_{k}=k$. Such a chain is called a Hölder series for $\mathfrak{g}$.\\
Definition 5.4.10. We call a representation ( $\pi, V$ ) of the Lie algebra $\mathfrak{g}$ nilpotent if there exists an $n \in \mathbb{N}$ with $\rho(\mathfrak{g})^{n}=\{0\}$.

Corollary 5.4.11. Let $\pi: \mathfrak{g} \rightarrow \mathfrak{g l}(V)$ be a representation of the solvable Lie algebra $\mathfrak{g}$. Then the restriction to $[\mathfrak{g}, \mathfrak{g}]$ is a nilpotent representation.

Proof. (a) First, we assume that $\mathbb{K}=\mathbb{C}$. Applying Lie's Theorem to the solvable subalgebra $\pi(\mathfrak{g})$ of $\mathfrak{g} \mathfrak{l}(V)$, we obtain a complete flag $\mathcal{F}$ with $\pi(\mathfrak{g}) \subseteq \mathfrak{g}(\mathcal{F})$. Then

$$
\pi([\mathfrak{g}, \mathfrak{g}]) \subseteq[\mathfrak{g}(\mathcal{F}), \mathfrak{g}(\mathcal{F})] \subseteq \mathfrak{g}_{n}(\mathcal{F})
$$

implies the assertion.\\
(b) If $\mathbb{K}=\mathbb{R}$, then

$$
\pi_{\mathbb{C}}: \mathfrak{g}_{\mathbb{C}} \rightarrow \mathfrak{g} \mathfrak{l}\left(V_{\mathbb{C}}\right) \cong \mathbb{C} \otimes \mathfrak{g} \mathfrak{l}(V), \quad \pi(z \otimes x)=z \otimes \pi(x)
$$

defines a representation of the complex Lie algebra $\mathfrak{g}_{\mathbb{C}}$ on the complex vector space $V_{\mathbb{C}}$. Using Proposition 5.1.21, we see by induction that $D^{k}(\mathfrak{g})_{\mathbb{C}} \cong D^{k}\left(\mathfrak{g}_{\mathbb{C}}\right)$ for each $k \in \mathbb{N}$, so that $\mathfrak{g}_{\mathbb{C}}$ is also solvable. Now (a) applies, and we obtain a complete flag $\mathcal{F}$ in $V_{\mathbb{C}}$ with $\pi_{\mathbb{C}}([\mathfrak{g}, \mathfrak{g}]) \subseteq \mathfrak{g}_{n}(\mathcal{F})$. In particular, each endomorphism $\pi(x), x \in[\mathfrak{g}, \mathfrak{g}]$, is nilpotent. Finally, Corollary 5.2.7 provides a complete flag $\mathcal{F}^{\prime}$ in $V$ with $\pi([\mathfrak{g}, \mathfrak{g}]) \subseteq \mathfrak{g}_{n}\left(\mathcal{F}^{\prime}\right)$, and the assertion follows.

Corollary 5.4.12. A Lie algebra $\mathfrak{g}$ is solvable if and only if its commutator algebra $[\mathfrak{g}, \mathfrak{g}]$ is nilpotent.

Proof. If $[\mathfrak{g}, \mathfrak{g}]$ is nilpotent, then $\mathfrak{g}$ is solvable because $\mathfrak{g} /[\mathfrak{g}, \mathfrak{g}]$ is abelian and solvability is an extension property (Proposition 5.4.3(ii)).

If, conversely, $\mathfrak{g}$ is solvable, then Corollary 5.4.11 implies that the adjoint representation of $[\mathfrak{g}, \mathfrak{g}]$ on $\mathfrak{g}$, and hence on $[\mathfrak{g}, \mathfrak{g}]$, is nilpotent. From that we derive in particular that $C^{\infty}([\mathfrak{g}, \mathfrak{g}])=\{0\}$, so that $[\mathfrak{g}, \mathfrak{g}]$ is nilpotent.

\subsubsection*{The Ideal $[\mathfrak{g}, \mathfrak{r a d}(\mathfrak{g})]$}
Lemma 5.4.13. Let $\mathfrak{g}$ be a Lie algebra and ( $\rho, V$ ) a representation of $\mathfrak{g}$. Let $\mathfrak{a} \subseteq \mathfrak{g}$ be a subspace for which there exists an $n \in \mathbb{N}$ with $\rho(\mathfrak{a})^{n}=\{0\}$ and $x \in \mathfrak{n}_{\mathfrak{g}}(\mathfrak{a})=\{y \in \mathfrak{g} \mid[y, \mathfrak{a}] \subseteq \mathfrak{a}\}$ such that $\rho(x)$ is nilpotent. Then there exists an $N \in \mathbb{N}$ with $\rho(\mathfrak{a}+\mathbb{K} x)^{N}=\{0\}$.

Proof. Replacing $\mathfrak{g}$ by $\rho(\mathfrak{g})$, we may w.lo.g. assume that $\mathfrak{g} \subseteq \mathfrak{g} \mathfrak{l}(V)$. Let $m \in \mathbb{N}$ with $x^{m}=0$. We claim that $(\mathbb{K} x+\mathfrak{a})^{n m}=\{0\}$.

Let $u=u_{1} \cdots u_{n m}$ be a product of elements of $\{x\} \cup \mathfrak{a}$. We have to show that all such products vanish. For $a \in \mathfrak{a}$, we have

$$
a x=x a+[a, x] \in x a+\mathfrak{a} .
$$

This leads to

$$
u_{1} \cdots u_{n m} \in \sum_{r=0}^{n m} x^{r} \mathfrak{a}^{t},
$$

where $t$ is the number of indices $j$ with $u_{j} \in \mathfrak{a}$. Hence this product vanishes for $t \geq n$. If $t<n$, then there exists a $j$ with $u_{j+1} \cdots u_{j+m}=x^{m}=0$ because in this case $n m-t>n(m-1)$ factors are not contained in $\mathfrak{a}$, so that we always find a consecutive product of $m$ such elements. We therefore have in all cases $u_{1} \cdots u_{n m}=0$.

Proposition 5.4.14. For any finite-dimensional representation ( $\rho, V$ ) of the Lie algebra $\mathfrak{g}$, the restriction to the ideal $[\mathfrak{g}, \mathfrak{r a d}(\mathfrak{g})]$ is nilpotent, i.e., there exists an $m \in \mathbb{N}$ with $\rho([\mathfrak{g}, \mathfrak{r a d}(\mathfrak{g})])^{m}=\{0\}$.

Proof. Let $\mathfrak{r}:=\mathfrak{r a d}(\mathfrak{g})$ and $\mathfrak{a}:=[\mathfrak{g}, \mathfrak{r}]$. According to Corollary 5.4.11, the representation is nilpotent on the ideal $[\mathfrak{r}, \mathfrak{r}]$. Now let $\mathfrak{t} \subseteq[\mathfrak{g}, \mathfrak{r}]$ be a subspace containing [ $\mathfrak{r}, \mathfrak{r}$ ], which is maximal with respect to the property that the representation on $V$ is nilpotent on $\mathfrak{t}$. Note that $\mathfrak{t}$ always is an ideal of $\mathfrak{r}$, hence in particular a subalgebra, because it contains the commutator algebra.

Assume that $\mathfrak{t} \neq[\mathfrak{g}, \mathfrak{r}]$. Then there exists an $x \in \mathfrak{g}$ and $y \in \mathfrak{r}$ with $[x, y] \notin \mathfrak{t}$. The subspace $\mathfrak{b}:=\mathfrak{r}+\mathbb{K} x$ is a subalgebra of $\mathfrak{g}, \mathfrak{r}$ is a solvable ideal of $\mathfrak{b}$, and $\mathfrak{b} / \mathfrak{r} \cong \mathbb{K}$ is abelian. Therefore, $\mathfrak{b}$ is solvable (Proposition 5.4.3).

Again, we use Corollary 5.4.11 to see that the representation is nilpotent on $[\mathfrak{b}, \mathfrak{b}]$ and hence that $\rho([x, y])$ is nilpotent. Since $\mathfrak{t} \subseteq \mathfrak{r}$ and $[x, y] \in[\mathfrak{g}, \mathfrak{r}] \subseteq \mathfrak{r}$, we have $[[x, y], \mathfrak{t}] \subseteq[\mathfrak{r}, \mathfrak{t}] \subseteq \mathfrak{t}$. Finally, the preceding Lemma 5.4.13 shows that the representation is nilpotent on the subspace $\mathbb{K}[x, y]+\mathfrak{t}$. This contradicts the maximality of $\mathfrak{t}$. We conclude that $\mathfrak{t}=[\mathfrak{g}, \mathfrak{r}]$, so that the representation is nilpotent on $[\mathfrak{g}, \mathfrak{r}]$.

Applying the preceding proposition to the adjoint representation, using Engel's Theorem 5.2.8 we get:

Corollary 5.4.15. The ideal $[\mathfrak{g}, \mathfrak{r a d}(\mathfrak{g})]$ is nilpotent. In particular, ad $x$ is nilpotent on $\mathfrak{g}$ for each $x \in[\mathfrak{g}, \mathfrak{r a d g}]$.

Remark 5.4.16. Since the ideal $[\mathfrak{g}, \mathfrak{r a d}(\mathfrak{g})]$ is nilpotent, it is contained in the nilradical $\mathfrak{n i l}(\mathfrak{g})$. That it may be strictly smaller follows from the case where $\mathfrak{g}$ is abelian. Then $\mathfrak{n i l}(\mathfrak{g})=\mathfrak{g}$ and $[\mathfrak{g}, \mathfrak{r a d}(\mathfrak{g})]=\{0\}$.

Corollary 5.4.17. If $\mathfrak{g}$ is a finite dimensional Lie algebra and $D \in \operatorname{der}(\mathfrak{g})$, then $D(\mathfrak{r a d}(\mathfrak{g})) \subseteq \mathfrak{n i l}(\mathfrak{g})$. In particular, $\mathfrak{n i l}(\mathfrak{g})$ is invariant under $\operatorname{der}(\mathfrak{g})$, i.e., a characteristic ideal.

Proof. For $D \in \operatorname{der}(\mathfrak{g})$, consider the Lie algebra $\widetilde{\mathfrak{g}}:=\mathfrak{g} \rtimes_{D} \mathbb{K}$ with the bracket

$$
\left[(x, t),\left(x^{\prime}, t^{\prime}\right)\right]:=\left(t D x^{\prime}-t^{\prime} D x+\left[x, x^{\prime}\right], 0\right)
$$

and identify $\mathfrak{g}$ with the subalgebra $\mathfrak{g} \times\{0\}$. Then $\mathfrak{r a d}(\mathfrak{g})$ is a solvable ideal of $\widetilde{\mathfrak{g}}$, hence contained in $\mathfrak{r a d}(\widetilde{\mathfrak{g}})$. We thus obtain

$$
D(\mathfrak{r a d}(\mathfrak{g}))=[(0,1), \mathfrak{r a d}(\mathfrak{g})] \subseteq \mathfrak{g} \cap[\widetilde{\mathfrak{g}}, \mathfrak{r a d}(\widetilde{\mathfrak{g}})] \subseteq \mathfrak{g} \cap \mathfrak{n i l}(\widetilde{\mathfrak{g}}) \subseteq \mathfrak{n i l}(\mathfrak{g})
$$

Corollary 5.4.18. If $\mathfrak{a} \unlhd \mathfrak{g}$ is an ideal, then $\mathfrak{n i l}(\mathfrak{a})=\mathfrak{a} \cap \mathfrak{n i l}(\mathfrak{g})$.\\
Proof. Clearly, $\mathfrak{a} \cap \mathfrak{n i l}(\mathfrak{g})$ is a nilpotent ideal of $\mathfrak{a}$, hence contained in $\mathfrak{n i l}(\mathfrak{a})$. Conversely, $\mathfrak{n i l}(\mathfrak{a})$ is an ideal of $\mathfrak{g}$ because it is invariant under all the derivations ad $\left.x\right|_{\mathfrak{a}}, x \in \mathfrak{g}$. This implies that $\mathfrak{n i l}(\mathfrak{a}) \subseteq \mathfrak{n i l}(\mathfrak{g}) \cap \mathfrak{a}$.

\subsubsection*{Cartan's Solvability Criterion}
This subsection is devoted to a characterization of solvable Lie algebras by properties of its elements. The result will be that $\mathfrak{g}$ is solvable if and only if $\operatorname{tr}(\operatorname{ad} x \operatorname{ad} y)=0$ for $x \in[\mathfrak{g}, \mathfrak{g}]$ and $y \in \mathfrak{g}$ (Cartan's criterion). Thus, we have to study the linear maps $\operatorname{ad}(x): \mathfrak{g} \rightarrow \mathfrak{g}$.

Lemma 5.4.19. Let $V$ be a finite-dimensional vector space and $E \subseteq F$ be subspaces of $\mathfrak{g l}(V)$. Further, let

$$
x \in M:=\{y \in \mathfrak{g} \mathfrak{l}(V) \mid[y, F] \subseteq E\} .
$$

If $\operatorname{tr}(x y)=0$ for all $y \in M$, then $x$ is nilpotent.\\
Proof. Complexifying all vector spaces involved, we may assume without loss of generality that $\mathbb{K}=\mathbb{C}$. In particular, we know that $x$ has a Jordan decomposition $x=x_{s}+x_{n}$. Now, let $\left\{v_{1}, \ldots, v_{n}\right\}$ be a basis for $V$ consisting of eigenvectors of $x_{s}$. We denote the corresponding eigenvalues by $\lambda_{j}$ for $j=1, \ldots, n$. Let $Q$ be the $\mathbb{Q}$-vector space in $\mathbb{C}$ which is spanned by the $\lambda_{j}$. We have to show that $Q=\{0\}$. To do this, we consider an $f \in Q^{*}:=\operatorname{Hom}_{\mathbb{Q}}(Q, \mathbb{Q})$, the dual space (over $\mathbb{Q}$ ) of $Q$. The matrix

$$
\left(\begin{array}{ccc}
f\left(\lambda_{1}\right) & & 0 \\
& \ddots & \\
0 & & f\left(\lambda_{n}\right)
\end{array}\right)
$$

represents an element $y \in \mathfrak{g} \mathfrak{l}(V)$ with respect to the basis $\left\{v_{1}, \ldots, v_{n}\right\}$. As in the proof of Proposition 5.3.8, if we choose a basis $\left\{x^{i j}\right\}$ for $\mathfrak{g l}(V)$ with $x^{i j}\left(v_{k}\right)=\delta_{j k} v_{i}$, we get

$$
\operatorname{ad}\left(x_{s}\right) x^{i j}=\left(\lambda_{i}-\lambda_{j}\right) x^{i j} \quad \text { and } \quad \operatorname{ad}(y) x^{i j}=f\left(\lambda_{i}-\lambda_{j}\right) x^{i j} .
$$

Now, choose a polynomial $P \in \mathbb{C}[t]$ with

$$
P(0)=0 \quad \text { and } \quad P\left(\lambda_{i}-\lambda_{j}\right)=f\left(\lambda_{i}-\lambda_{j}\right)
$$

for all pairs $(i, j)$ (Exercise 5.4.6). Then

$$
P\left(\operatorname{ad}\left(x_{s}\right)\right) x^{i j}=f\left(\lambda_{i}-\lambda_{j}\right) x^{i j}=\operatorname{ad}(y) x^{i j}
$$

i.e., $P\left(\operatorname{ad}\left(x_{s}\right)\right)=\operatorname{ad}(y)$. Since $\operatorname{ad}\left(x_{s}\right)$ is the semisimple part of $\operatorname{ad}(x)$ by Corollary 5.3.9, it follows from $x \in M$ and Proposition 5.3.7(iii) that $\operatorname{ad}\left(x_{s}\right) F \subseteq E$. But then $P(0)=0$ implies $\operatorname{ad}(y) F \subseteq E$, i.e., $y \in M$. Since $x y v_{i}=\lambda_{i} f\left(\lambda_{i}\right) v_{i}$ for each $i$, our assumption and $y \in M$ leads to

$$
\sum_{k=1}^{n} \lambda_{k} f\left(\lambda_{k}\right)=\operatorname{tr}(x y)=0
$$

and therefore,

$$
\sum_{k=1}^{n} f\left(\lambda_{k}\right)^{2}=f\left(\sum_{k=1}^{n} \lambda_{k} f\left(\lambda_{k}\right)\right)=0 .
$$

Hence $f\left(\lambda_{k}\right)=0$ for all $\lambda_{k}$, which yields $f=0$. But since $f \in Q^{*}$ was arbitrary, we see that $Q=\{0\}$.

Theorem 5.4.20 (Cartan's Solvability Criterion). Let $V$ be a finitedimensional vector space and $\mathfrak{g}<\mathfrak{g} \mathfrak{l}(V)$. Then the following are equivalent\\
(i) $\mathfrak{g}$ is solvable.\\
(ii) $\operatorname{tr}(x y)=0$ for all $x \in[\mathfrak{g}, \mathfrak{g}]$ and $y \in \mathfrak{g}$.

Proof. $(i i) \Rightarrow(i)$ : By Corollary 5.4.12, it suffices to show that $[\mathfrak{g}, \mathfrak{g}]$ is nilpotent. But to show that, by Corollary 5.2.7, we only have to prove that every element of $[\mathfrak{g}, \mathfrak{g}]$ is nilpotent. We want to apply Lemma 5.4.19 with $E=[\mathfrak{g}, \mathfrak{g}]$ and $F=\mathfrak{g}$, i.e., we set

$$
M:=\{y \in \mathfrak{g} \mathfrak{l}(V) \mid[y, \mathfrak{g}] \subseteq[\mathfrak{g}, \mathfrak{g}]\} .
$$

Since the trace is linear, it is enough to show that $\operatorname{tr}\left(\left[x, x^{\prime}\right] y\right)=0$ for $x, x^{\prime} \in \mathfrak{g}$ and $y \in M$. But this follows from $\left[x^{\prime}, y\right] \subseteq[\mathfrak{g}, \mathfrak{g}]$ and $(i i)$ :

$$
\operatorname{tr}\left(\left[x, x^{\prime}\right] y\right)=\operatorname{tr}\left(x\left[x^{\prime}, y\right]\right)=0
$$

(cf. Exercise 5.4.5).\\
$(i) \Rightarrow(i i)$ : Since $\operatorname{tr}\left(x_{\mathbb{C}}\right)=\operatorname{tr} x$ for $x \in \operatorname{End}(V)$, we may assume that $\mathbb{K}=\mathbb{C}$ (cf. also Exercise 5.4.2). Then by Lie's Theorem 5.4.8, there is a basis for $V$ with respect to which all $x \in \mathfrak{g}$ are upper triangular matrices. In particular, all elements of $[\mathfrak{g}, \mathfrak{g}]$ are given by strictly upper triangular matrices. But multiplying an upper triangular matrix with a strictly upper triangular matrix yields a strictly upper triangular matrix which has zero trace.

Corollary 5.4.21. Let $\mathfrak{g}$ be a Lie algebra. Then the following statements are equivalent\\
(i) $\mathfrak{g}$ is solvable.\\
(ii) $\operatorname{tr}(\operatorname{ad} x \operatorname{ad} y)=0$ for all $x \in[\mathfrak{g}, \mathfrak{g}]$ and all $y \in \mathfrak{g}$.

Proof. $(i i) \Rightarrow(i)$ : By the Cartan Criterion (Theorem 5.4.20), $\operatorname{ad}(\mathfrak{g})$ is solvable. Then it follows from $\operatorname{ad}(\mathfrak{g}) \cong \mathfrak{g} / \mathfrak{z}(\mathfrak{g})$ and Proposition 5.4.3(ii) that $\mathfrak{g}$ is solvable.\\
$(i) \Rightarrow(i i)$ : Proposition 5.4.3(i) shows that $\operatorname{ad}(\mathfrak{g})$ is solvable, so that (ii) is an immediate consequence of Theorem 5.4.20.

\subsubsection*{Exercises for Section 5.4}
Exercise 5.4.1. Let $\mathfrak{g}$ be a Lie algebra and $\alpha: \mathfrak{g} \rightarrow \mathfrak{g} \mathfrak{l}(V)$ be a representation of $\mathfrak{g}$ on $V$. Then $V \rtimes_{\alpha} \mathfrak{g}$ is a Lie algebra which contains $V$ as an abelian ideal.

Exercise 5.4.2. (i) For a real Lie algebra $\mathfrak{g}$, we have

$$
C^{n}\left(\mathfrak{g}_{\mathbb{C}}\right)=C^{n}(\mathfrak{g})_{\mathbb{C}} \quad \text { and } \quad D^{n}\left(\mathfrak{g}_{\mathbb{C}}\right)=D^{n}(\mathfrak{g})_{\mathbb{C}}
$$

(ii) A finite-dimensional Lie algebra $\mathfrak{g}$ is nilpotent (solvable) if and only if $\mathfrak{g}_{\mathbb{C}}$ is nilpotent (solvable).

Exercise 5.4.3. Show that for the Heisenberg algebra $\mathfrak{h}_{3}$, the derivation algebra is isomorphic to $\mathbb{K}^{2} \rtimes \mathfrak{g} \mathfrak{l}_{2}(\mathbb{K})$, where $\operatorname{ad}\left(\mathfrak{h}_{3}\right) \cong \mathbb{K}^{2}$ is an abelian ideal. Show that this Lie algebra is neither nilpotent nor solvable.

Exercise 5.4.4. Show that a representation ( $\pi, V$ ) of a Lie algebra $\mathfrak{g}$ on a vector space $V$ is nilpotent if and only if there is a flag $\mathcal{F}$ in $V$ with $\pi(V) \subseteq \mathfrak{g}_{n}(\mathcal{F})$.

Exercise 5.4.5. A symmetric bilinear form $\kappa: \mathfrak{g} \times \mathfrak{g} \rightarrow \mathbb{K}$ on a Lie algebra $\mathfrak{g}$ is called invariant if

$$
\kappa([x, y], z)=\kappa(x,[y, z]) \quad \text { for } \quad x, y, z \in \mathfrak{g}
$$

Show that:\\
(i) The form $\kappa(x, y):=\operatorname{tr}(x y)$ on $\mathfrak{g} \mathfrak{l}(V)$ is invariant for each finite-dimensional vector space $V$.\\
(ii) For each representation ( $\pi, V$ ) of the Lie algebra $\mathfrak{g}$, the form $\kappa_{\pi}(x, y):= \operatorname{tr}(\pi(x) \pi(y))$ is invariant.\\
(iii) For each Lie algebra $\mathfrak{g}$, the Cartan-Killing form $\kappa_{\mathfrak{g}}(x, y):=\operatorname{tr}(\operatorname{ad} x \operatorname{ad} y)$ is invariant.\\
(iv) For each invariant symmetric bilinear form $\kappa$ on $\mathfrak{g}$, its radical $\mathfrak{r a d}(\kappa)= \{x \in \mathfrak{g}: \kappa(x, \mathfrak{g})=\{0\}\}$ is an ideal.\\
(v) For any invariant symmetric bilinear form $\kappa$ on $\mathfrak{g}$, the trilinear map $\Gamma(\kappa)(x, y, z):=\kappa([x, y], z)$ is alternating, i.e.,

$$
\Gamma(\kappa)\left(x_{\sigma(1)}, x_{\sigma(2)}, x_{\sigma(3)}\right)=\operatorname{sgn}(\sigma) \Gamma(\kappa)\left(x_{1}, x_{2}, x_{3}\right)
$$

for $\sigma \in S_{3}$ and $x_{1}, x_{2}, x_{3} \in \mathfrak{g}$.\\
(vi) $\mathfrak{g}$ is solvable if and only if $\Gamma\left(\kappa_{\mathfrak{g}}\right)=0$.

Exercise 5.4.6 (Interpolation polynomials). Let $\mathbb{K}$ be a field, $x_{1}, \ldots, x_{n} \in \mathbb{K}$ pairwise different, and $\lambda_{1}, \ldots, \lambda_{n} \in \mathbb{K}$. Then there exists a polynomial $f \in \mathbb{K}[t]$ with $f\left(x_{i}\right)=\lambda_{i}$ for $i=1, \ldots, n$. Hint: Consider the polynomials $f_{i}(t):=\prod_{j \neq i} \frac{t-x_{j}}{x_{i}-x_{j}}$ of degree $n-1$.

Exercise 5.4.7. This exercise shows why Lie's Theorem does not generalize to infinite-dimensional spaces. We consider the vector space $V=\mathbb{K}^{(\mathbb{N})}$ with the basis $\left\{e_{i}: i \in \mathbb{N}\right\}$. In terms of the rank-one-operators $E_{i j} \in \operatorname{End}(V)$, defined by $E_{i j} e_{k}=\delta_{j k} e_{i}$, we consider the Lie algebra

$$
\mathfrak{g}:=\operatorname{span}\left\{E_{i j}: i \geq j\right\}
$$

(lower triangular matrices). Show that:\\
(a) $D^{n}(\mathfrak{g})=\operatorname{span}\left\{E_{i j}: i \geq j+2^{n-1}\right\}, n \in \mathbb{N}$. In particular, we have $D^{\infty}(\mathfrak{g}):= \bigcap_{n \in \mathbb{N}} D^{n}(\mathfrak{g})=\{0\}$, i.e., $\mathfrak{g}$ is residually solvable.\\
(b) $\mathfrak{g}=\bigcup_{n} \mathfrak{g}_{n}$ for an increasing sequence of finite-dimensional solvable subalgebras $\mathfrak{g}_{n}$ ( $\mathfrak{g}$ is locally solvable).\\
(c) $\mathfrak{g}$ has no common eigenvector in $V$ (compare with Lie's Theorem).

Exercise 5.4.8. Show that:\\
(a) A finite-dimensional Lie algebra $\mathfrak{g}$ is solvable if and only if there exists a sequence

$$
\{0\}=\mathfrak{g}_{0} \subseteq \mathfrak{g}_{1} \subseteq \cdots \subseteq \mathfrak{g}_{n}=\mathfrak{g}
$$

of subalgebras with $\left[\mathfrak{g}_{i}, \mathfrak{g}_{i}\right] \subseteq \mathfrak{g}_{i-1}$ for $i=1, \ldots, n$.\\
(b) If $\mathfrak{g}$ is solvable, then there exists a sequence as in (a), satisfying, in addition, $\operatorname{dim} \mathfrak{g}_{i}=i$. Conclude that

$$
\mathfrak{g}_{i+1} \cong \mathfrak{g}_{i} \rtimes_{D_{i}} \mathbb{K} \quad \text { for some } \quad D_{i} \in \operatorname{der}\left(\mathfrak{g}_{i}\right), \quad i=1, \ldots, n-1
$$

This means that

$$
\mathfrak{g} \cong\left(\cdots\left(\left(\mathbb{K} \rtimes_{D_{1}} \mathbb{K}\right) \rtimes_{D_{2}} \mathbb{K}\right) \cdots \rtimes_{D_{n-1}} \mathbb{K}\right)
$$

Exercise 5.4.9. (Is there a Cartan Criterion for nilpotent Lie algebras?)\\
(a) If $\mathfrak{g}$ is nilpotent, then $\kappa_{\mathfrak{g}}=0$.\\
(b) Consider the Lie algebra $\mathfrak{g}=\mathbb{C}^{2} \rtimes_{D} \mathbb{C}$, where $\mathbb{C}^{2}$ is considered as an abelian Lie algebra and $D=\operatorname{diag}(1, i)$. This Lie algebra is not nilpotent, but $\kappa_{\mathfrak{g}}=0$. If we consider $\mathfrak{g}$ as a 6 -dimensional real Lie algebra, then its Cartan-Killing form also vanishes. Conclude that it is NOT true that a Lie algebra is nilpotent if and only if its Cartan-Killing form vanishes.

\subsection*{Semisimple Lie Algebras}
In this section, we encounter a third class of Lie algebras. Semisimple Lie algebras are a counterpart to the solvable and nilpotent Lie algebras because their ideal structure is quite simple. They can be decomposed as a direct sum of simple ideals. On the other hand, they have a rich geometric structure which even makes a complete classification of finite-dimensional semisimple Lie algebras possible. One can even show that every finite-dimensional Lie algebra is a semidirect sum of its radical and a semisimple subalgebra (cf. Levi's Theorem 5.6.6).

Definition 5.5.1. Let $\mathfrak{g}$ be a finite-dimensional Lie algebra. Then $\mathfrak{g}$ is called semisimple if its radical is trivial, i.e., $\mathfrak{r a d}(\mathfrak{g})=\{0\}$. The Lie algebra $\mathfrak{g}$ is called simple if it is not abelian and it contains no ideals other than $\mathfrak{g}$ and $\{0\}$.

Lemma 5.5.2. Each simple Lie algebra is semisimple.\\
Proof. Let $\mathfrak{g}$ be a simple Lie algebra. Since the commutator algebra $[\mathfrak{g}, \mathfrak{g}]$ is a nonzero ideal of $\mathfrak{g}$, it coincides with $\mathfrak{g}$. Hence $\mathfrak{g}$ is not solvable, so that $\mathfrak{r a d}(\mathfrak{g})=\{0\}$.

We shall see in Proposition 5.5.11 that a Lie algebra is semisimple if and only if it is a direct sum of simple ideals.

\subsubsection*{Cartan's Semisimplicity Criterion}
In this subsection, we prove the characterization of semisimple Lie algebras in terms of the Cartan-Killing form which can be defined for any Lie algebra.

Definition 5.5.3. In connection with the Cartan criterion for solvable Lie algebras, we have seen that the bilinear form

$$
\kappa_{\mathfrak{g}}: \mathfrak{g} \times \mathfrak{g} \rightarrow \mathbb{K}, \quad \kappa_{\mathfrak{g}}(x, y):=\operatorname{tr}(\operatorname{ad} x \operatorname{ad} y)
$$

on a finite-dimensional Lie algebra is of interest. It is called the CartanKilling form of $\mathfrak{g}$. Its compatibility with the Lie algebra structure is expressed by its invariance

$$
\kappa_{\mathfrak{g}}([x, y], z)=\kappa_{\mathfrak{g}}(x,[y, z]) \quad \text { for } \quad x, y, z, \in \mathfrak{g}
$$

(Exercise 5.4.5). If $\mathfrak{g}$ is clear from the context, we sometimes write $\kappa$ instead of $\kappa_{\mathfrak{g}}$.

Example 5.5.4. (i) With respect to the basis ( $h, u, t$ ) for $\mathfrak{s l}_{2}(\mathbb{K})$ given in Example 5.1.23, the Cartan-Killing form has the matrix

$$
\kappa=\left(\begin{array}{ccc}
8 & 0 & 0 \\
0 & -8 & 0 \\
0 & 0 & 8
\end{array}\right)
$$

(ii) With respect to the basis $(x, y, z)$ for $\mathfrak{s o}_{3}(\mathbb{R})$, given in Example 5.1.23, the Cartan-Killing form has the matrix

$$
\kappa=\left(\begin{array}{ccc}
-2 & 0 & 0 \\
0 & -2 & 0 \\
0 & 0 & -2
\end{array}\right)
$$

(iii) With respect to the basis $(h, p, q, z)$ for the oscillator algebra given in Example 5.1.19, the Cartan-Killing form has the matrix

$$
\kappa=\left(\begin{array}{cccc}
-2 & 0 & 0 & 0 \\
0 & 0 & 0 & 0 \\
0 & 0 & 0 & 0 \\
0 & 0 & 0 & 0
\end{array}\right)
$$

In general, the Cartan-Killing form of a subalgebra $\mathfrak{h} \subseteq \mathfrak{g}$ cannot be calculated in terms of the Cartan-Killing form of $\mathfrak{g}$, but for ideals we have:

Lemma 5.5.5. For any ideal $\mathfrak{i} \unlhd \mathfrak{g}, \kappa_{\mathfrak{i}}=\left.\kappa_{\mathfrak{g}}\right|_{\mathfrak{i} \times \mathfrak{i}}$.\\
Proof. If the image of $A \in \operatorname{End}(\mathfrak{g})$ is contained in $\mathfrak{i}$, then we pick a basis for $\mathfrak{g}$ which starts with a basis for $\mathfrak{i}$. With respect to this basis, we can write $A$ as a block matrix

$$
A=\left(\begin{array}{cc}
\left.A\right|_{\mathrm{i}} & * \\
0 & 0
\end{array}\right)
$$

and this shows $\operatorname{tr}(A)=\operatorname{tr}\left(\left.A\right|_{\mathfrak{i}}\right)$. We apply this to $A=\operatorname{ad}(x) \operatorname{ad}(y)$ for $x, y \in \mathfrak{i}$ to obtain $\operatorname{tr}(\operatorname{ad}(x) \operatorname{ad}(y))=\operatorname{tr}\left(\left.\left.\operatorname{ad}(x)\right|_{\mathfrak{i}} \operatorname{ad}(y)\right|_{\mathfrak{i}}=\kappa_{\mathfrak{i}}(x, y)\right.$.

Remark 5.5.6. Let $\mathfrak{g}$ be a finite-dimensional $\mathbb{R}$-Lie algebra. Since a basis for $\mathfrak{g}$ is also a (complex) basis for $\mathfrak{g}_{\mathbb{C}}$, one immediately sees (cf. Exercise 5.5.3) that

$$
\kappa_{\mathfrak{g}}=\left.\kappa_{\mathfrak{g} C}\right|_{\mathfrak{g} \times \mathfrak{g}}
$$

Let $V$ be a vector space and $\beta: V \times V \rightarrow \mathbb{K}$ be a symmetric bilinear form. Then we denote the orthogonal set

$$
\{v \in V \mid(\forall w \in W) \beta(v, w)=0\}
$$

of a subspace $W$ with respect to $\beta$ by $W^{\perp, \beta}$. If $\beta$ is the Cartan-Killing form of a Lie algebra, we simply write $\perp$ instead of $\perp, \beta$. The set $\mathfrak{r a d}(\beta):=V^{\perp, \beta}$ is called the radical of $\beta$. The form is called degenerate if $\mathfrak{r a d}(\beta) \neq\{0\}$. Using this notation, we can reformulate the Cartan Criterion 5.4.20 as follows:

Remark 5.5.7. In terms of the Cartan-Killing form, Cartan's Solvability Criterion states that $\mathfrak{g}$ is solvable if and only if $[\mathfrak{g}, \mathfrak{g}] \subseteq \mathfrak{r a d}\left(\kappa_{\mathfrak{g}}\right)$ (cf. Exercise 5.5.7 for a description of the radical of $\kappa_{\mathfrak{g}}$ for a general Lie algebra).

Lemma 5.5.8. For any ideal $\mathfrak{j}$ of a Lie algebra $\mathfrak{g}$, the following assertions hold:\\
(i) Its orthogonal space $\mathfrak{j}^{\perp}$ with respect to $\kappa_{\mathfrak{g}}$ is also an ideal.\\
(ii) $\mathfrak{j} \cap \mathfrak{j}^{\perp}$ is a solvable ideal.\\
(iii) If $\mathfrak{j}$ is semisimple, then $\mathfrak{g}$ decomposes as a direct sum $\mathfrak{g}=\mathfrak{j} \oplus \mathfrak{j}^{\perp}$ of Lie algebras.

Proof. (i) For $x \in \mathfrak{j}^{\perp}, y \in \mathfrak{g}$ and $z \in \mathfrak{j}$, we find $\kappa_{\mathfrak{g}}([x, y], z)=\kappa_{\mathfrak{g}}(x,[y, z])=0$, so that $\mathfrak{j}^{\perp}$ is an ideal of $\mathfrak{g}$.\\
(ii) For $\mathfrak{i}:=\mathfrak{j} \cap \mathfrak{j}^{\perp}$, the Cartan-Killing form $\kappa_{\mathfrak{g}}$ vanishes on $\mathfrak{i} \times \mathfrak{i}$. Hence $\mathfrak{r a d}\left(\kappa_{\mathfrak{i}}\right)=\mathfrak{i}$ by Lemma 5.5.5. In particular, $\mathfrak{i}$ is solvable by Remark 5.5.7.\\
(iii) If $\mathfrak{j}$ is semisimple, then (ii) implies that $\mathfrak{j} \cap \mathfrak{j}^{\perp} \subseteq \mathfrak{r a d}(\mathfrak{j})=\{0\}$. Since $\mathfrak{j}^{\perp}$ is the kernel of the linear map $\mathfrak{g} \rightarrow \mathfrak{j}^{*}, x \mapsto \kappa_{\mathfrak{g}}(x, \cdot)$, we have $\operatorname{dim} \mathfrak{j}^{\perp} \geq \operatorname{dim} \mathfrak{g}-\operatorname{dim} \mathfrak{j}$, which implies $\mathfrak{j}+\mathfrak{j}^{\perp}=\mathfrak{g}$, so that $\mathfrak{g}$ is a direct sum of the vector subspaces $\mathfrak{j}$ and $\mathfrak{j}^{\perp}$. As both are ideals, $\left[\mathfrak{j}, \mathfrak{j}^{\perp}\right] \subseteq \mathfrak{j} \cap \mathfrak{j}^{\perp}=\{0\}$, and we obtain a direct sum of Lie algebras.

We can also characterize semisimplicity in terms of the Cartan-Killing form.

Theorem 5.5.9 (Cartan's Semisimplicity Criterion). A Lie algebra $\mathfrak{g}$ is semisimple if and only if $\kappa_{\mathfrak{g}}$ is nondegenerate, i.e., $\mathfrak{r a d}\left(\kappa_{\mathfrak{g}}\right)=\{0\}$.

Proof. With Lemma 5.5.8(ii), we see that $\mathfrak{g} \cap \mathfrak{g}^{\perp}=\mathfrak{r a d}\left(\kappa_{\mathfrak{g}}\right)$ is a solvable ideal, so that $\mathfrak{r a d}\left(\kappa_{\mathfrak{g}}\right) \subseteq \mathfrak{r a d}(\mathfrak{g})$. In particular, $\kappa_{\mathfrak{g}}$ is nondegenerate if $\mathfrak{g}$ is semisimple.

Suppose, conversely, that $\mathfrak{g}$ is not semisimple and put $\mathfrak{r}:=\mathfrak{r a d}(\mathfrak{g}) \neq\{0\}$. Let $n \in \mathbb{N}_{0}$ be maximal with $\mathfrak{h}:=D^{n}(\mathfrak{r}) \neq\{0\}$. Then $\mathfrak{h}$ is an abelian ideal of $\mathfrak{g}$. For $x \in \mathfrak{h}$ and $y \in \mathfrak{g}$, we then have $(\operatorname{ad} x \operatorname{ad} y) \mathfrak{g} \subseteq \mathfrak{h}$, and therefore $(\operatorname{ad} x \operatorname{ad} y)^{2}=0$. This implies that $\kappa_{\mathfrak{g}}(x, y)=\operatorname{tr}(\operatorname{ad} x \operatorname{ad} y)=0$. Since $y \in \mathfrak{g}$ was arbitrary, this means that $x \in \mathfrak{r a d}\left(\kappa_{\mathfrak{g}}\right)$, i.e., $\kappa_{\mathfrak{g}}$ is degenerate.

Remark 5.5.10. In view of $\mathfrak{r a d}\left(\kappa_{\mathfrak{g}}\right)_{\mathbb{C}}=\mathfrak{r a d}\left(\kappa_{\mathfrak{g} \mathbb{C}}\right)$ (cf. Exercise 5.5.2), Theorem 5.5.9 shows that a real Lie algebra $\mathfrak{g}$ is semisimple if and only if its complexification $\mathfrak{g}_{\mathbb{C}}$ is semisimple.

Proposition 5.5.11. Let $\mathfrak{g}$ be a semisimple Lie algebra. Then there are simple ideals $\mathfrak{g}_{1}, \ldots, \mathfrak{g}_{k}$ of $\mathfrak{g}$ with

$$
\mathfrak{g}=\mathfrak{g}_{1} \oplus \cdots \oplus \mathfrak{g}_{k}
$$

Every ideal $\mathfrak{i} \unlhd \mathfrak{g}$ is semisimple and a direct sum $\mathfrak{i}=\bigoplus_{j \in I} \mathfrak{g}_{j}$ for some subset $I \subseteq\{1, \ldots, k\}$. Conversely, each direct sum of simple Lie algebras is semisimple.

Proof. Let $\mathfrak{j} \unlhd \mathfrak{g}$. Since $\mathfrak{g}$ is semisimple, Lemma 5.5.8(ii) shows that $\mathfrak{j} \cap \mathfrak{j}^{\perp}= \{0\}$, so that $\kappa_{\mathfrak{j}}=\left.\kappa_{\mathfrak{g}}\right|_{\mathfrak{j} \times \mathfrak{j}}$ is nondegenerate. Hence $\mathfrak{j}$ is semisimple. According to Lemma 5.5.8(iii), $\mathfrak{j}^{\perp}$ also is an ideal of $\mathfrak{g}$ with $\mathfrak{g}=\mathfrak{j} \oplus \mathfrak{j}^{\perp}$ (direct sum of Lie algebras). As $\mathfrak{j}$ and $\mathfrak{j}^{\perp}$ are semisimple, an induction on $\operatorname{dim} \mathfrak{g}$ now implies that $\mathfrak{g}$ decomposes as a direct sum

$$
\mathfrak{g}=\mathfrak{g}_{1} \oplus \cdots \oplus \mathfrak{g}_{k}
$$

of simple ideals.

Finally, let $\mathfrak{i} \neq\{0\}$ be an ideal of $\mathfrak{g}$. Let $\pi_{j}: \mathfrak{g} \rightarrow \mathfrak{g}_{j}$ be the projections. Then we have $\pi_{j}(\mathfrak{i}) \neq\{0\}$ for at least one $j$. But since $\pi_{j}$ is surjective $\pi_{j}(\mathfrak{i})$ is an ideal of $\mathfrak{g}_{j}$, and therefore equal to $\mathfrak{g}_{j}$. Thus

$$
\mathfrak{g}_{j}=\left[\mathfrak{g}_{j}, \mathfrak{g}_{j}\right]=\left[\mathfrak{g}_{j}, \pi_{j}(\mathfrak{i})\right]=\left[\mathfrak{g}_{j}, \mathfrak{i}\right] \subseteq \mathfrak{i}
$$

because $\left[\mathfrak{g}_{j}, \pi_{l}(\mathfrak{i})\right]=\{0\}$ for $l \neq j$. The argument shows that every $\mathfrak{g}_{j}$ with $\pi_{j}(\mathfrak{i}) \neq\{0\}$ is contained in $\mathfrak{i}$. But then $\mathfrak{i}$ is the direct sum of these $\mathfrak{g}_{j}$.

The preceding argument shows in particular that no nonzero ideal of a direct sum of simple Lie algebras is solvable, hence that any such direct sum is semisimple.

Example 5.5.12. The Lie algebras $\mathfrak{s l}_{2}(\mathbb{K}), \mathfrak{s o}_{3}(\mathbb{K})$, and $\mathfrak{s u}_{2}(\mathbb{C})$ are simple. We have seen in Example 5.5.4 that the Cartan-Killing forms of $\mathfrak{s} \mathfrak{l}_{2}(\mathbb{K})$ and $\mathfrak{s o}_{3}(\mathbb{K})$ are nondegenerate, so that they are semisimple. Further, $\mathfrak{s u}_{2}(\mathbb{C})$ is a real form of $\mathfrak{s} \mathfrak{l}_{2}(\mathbb{C})$, hence semisimple by Remark 5.5.10. Now we apply Exercise 5.5.1.

Corollary 5.5.13. For a semisimple Lie algebra $\mathfrak{g}$, the following assertions hold:\\
(i) $\mathfrak{g}$ is perfect, i.e., $\mathfrak{g}=[\mathfrak{g}, \mathfrak{g}]$.\\
(ii) All homomorphic images of $\mathfrak{g}$ are semisimple.\\
(iii) Each ideal $\mathfrak{n} \unlhd \mathfrak{g}$ is semisimple and there exists another ideal $\mathfrak{c}$ with $\mathfrak{g}=\mathfrak{n} \oplus \mathfrak{c}$.

Proof. (i) Let $\mathfrak{g}=\mathfrak{g}_{1} \oplus \cdots \oplus \mathfrak{g}_{k}$ be the decomposition of $\mathfrak{g}$ from Proposition 5.5.11. Since $\mathfrak{g}_{j}$ is not abelian, we have $\left[\mathfrak{g}_{j}, \mathfrak{g}_{j}\right]=\mathfrak{g}_{j}$, and therefore

$$
[\mathfrak{g}, \mathfrak{g}]=\sum_{j=1}^{k}\left[\mathfrak{g}_{j}, \mathfrak{g}_{j}\right]=\sum_{j=1}^{k} \mathfrak{g}_{j}=\mathfrak{g} .
$$

(ii) This follows by combining Proposition 5.5.11 with the Isomorphism Theorem 5.1.17(i).\\
(iii) This follows immediately from Proposition 5.5.11.

In Example 5.1.9, we have seen that the adjoint representation gives derivations on the Lie algebra. In the case of semisimple Lie algebras, this representation, in fact, gives all derivations.

Theorem 5.5.14. For a semisimple Lie algebra $\mathfrak{g}$ all derivations are inner, i.e.,

$$
\operatorname{ad}(\mathfrak{g})=\operatorname{der}(\mathfrak{g})
$$

Proof. By Proposition 5.1.10(i), ad $\mathfrak{g} \unlhd \operatorname{der}(\mathfrak{g})$ is an ideal, and since $\mathfrak{z}(\mathfrak{g})= \{0\}$, the ideal $\operatorname{ad}(\mathfrak{g}) \cong \mathfrak{g}$ is semisimple. Therefore, der $\mathfrak{g}$ decomposes as a direct sum $\operatorname{ad} \mathfrak{g} \oplus \mathfrak{j}$ for the orthogonal complement $\mathfrak{j}$ of $\operatorname{ad}(\mathfrak{g})$ with respect to the

Cartan-Killing form of $\operatorname{der}(\mathfrak{g})$ (Lemma 5.5.8(iii)). For $\delta \in \mathfrak{j}$ and $x \in \mathfrak{g}$, we then have

$$
0=[\delta, \operatorname{ad} x](y)=\delta([x, y])-[x, \delta(y)]=[\delta(x), y]=\operatorname{ad}(\delta(x))
$$

This means that $\delta(x) \in \mathfrak{z}(\mathfrak{g})=\{0\}$, i.e., $\delta=0$. We conclude that $\mathfrak{j}=\{0\}$, so that $\operatorname{der}(\mathfrak{g})=\operatorname{ad} \mathfrak{g}$.

\subsubsection*{Weyl's Theorem on Complete Reducibility}
We have already seen how Engel's Theorem and Lie's Theorem provide important information on the structure of representations of nilpotent, resp., solvable, Lie algebras. For semisimple Lie algebras, Weyl's Theorem which asserts that each representation of a semisimple Lie algebra is completely reducible plays a similar role. The crucial tool needed for the proof of Weyl's Theorem is the Casimir element.

Definition 5.5.15. Let $\beta$ be a nondegenerate invariant symmetric bilinear form on the Lie algebra $\mathfrak{g}, x_{1}, \ldots, x_{n}$ a basis for $\mathfrak{g}$, and $x^{1}, \ldots, x^{n}$ the dual basis with respect to $\beta$, i.e., $\beta\left(x_{i}, x^{j}\right)=\delta_{i j}$ (Kronecker delta). For any Lie algebra homomorphism $\rho: \mathfrak{g} \rightarrow A_{L}, A$ an associative algebra, we define the Casimir element by

$$
\Omega(\beta, \rho):=\sum_{i=1}^{k} \rho\left(x_{i}\right) \rho\left(x^{i}\right) .
$$

The same argument that shows the independence of the trace of an operator (defined as the sum of its diagonal elements) of the choice of the basis, shows that $\Omega(\beta, \rho)$ does not depend on the choice of the basis $x_{1}, \ldots, x_{k}$ (cf. Exercise 5.5.9).

The Casimir element is a useful tool for the study of representations since it commutes with $\rho(\mathfrak{g})$ :

Lemma 5.5.16. For each nondegenerate invariant symmetric bilinear form $\beta$ on $\mathfrak{g}$ and each homomorphism $\rho: \mathfrak{g} \rightarrow A$, the Casimir element $\Omega(\beta, \rho) \in A$ commutes with $\rho(\mathfrak{g})$.

Proof. Let $z \in \mathfrak{g}$. Then we have

$$
\operatorname{ad} z\left(x_{j}\right)=\sum_{k=1}^{n} a_{k j} x_{k} \quad \text { and } \quad \operatorname{ad} z\left(x^{j}\right)=\sum_{k=1}^{n} a^{k j} x^{k}
$$

with two matrices ( $a_{i j}$ ) and ( $a^{i j}$ ) in $M_{k}(\mathbb{K})$. Then

$$
a_{k j}=\beta\left(\left[z, x_{j}\right], x^{k}\right)=-\beta\left(x_{j},\left[z, x^{k}\right]\right)=-a^{j k}
$$

and with this relation we obtain

$$
\begin{aligned}
& \Omega \rho(z)-\rho(z) \Omega=\sum_{j=1}^{n} \rho\left(x_{j}\right) \rho\left(x^{j}\right) \rho(z)-\rho(z) \rho\left(x_{j}\right) \rho\left(x^{j}\right) \\
& =\sum_{j=1}^{n} \rho\left(x_{j}\right)\left(\rho\left(x^{j}\right) \rho(z)-\rho(z) \rho\left(x^{j}\right)\right) \\
& \quad-\left(\rho(z) \rho\left(x_{j}\right)-\rho\left(x_{j}\right) \rho(z)\right) \rho\left(x^{j}\right) \\
& =\sum_{j=1}^{n} \rho\left(x_{j}\right) \rho\left(\left[x^{j}, z\right]\right)-\rho\left(\left[z, x_{j}\right]\right) \rho\left(x^{j}\right) \\
& =\sum_{j, k=1}^{n}\left(-a^{k j}\right) \rho\left(x_{j}\right) \rho\left(x^{k}\right)-a_{k j} \rho\left(x_{k}\right) \rho\left(x^{j}\right) \\
& =\sum_{j, k=1}^{n} a_{j k} \rho\left(x_{j}\right) \rho\left(x^{k}\right)-a_{k j} \rho\left(x_{k}\right) \rho\left(x^{j}\right) \\
& =\sum_{j, k=1}^{n} a_{j k} \rho\left(x_{j}\right) \rho\left(x^{k}\right)-a_{j k} \rho\left(x_{j}\right) \rho\left(x^{k}\right)=0 .
\end{aligned}
$$

Proposition 5.5.17. Let $\mathfrak{g}$ be a semisimple Lie algebra and ( $\rho, V$ ) a finitedimensional representation. Then $V$ is the direct sum of the $\mathfrak{g}$-modules

$$
V^{\mathfrak{g}}:=\bigcap_{x \in \mathfrak{g}} \operatorname{ker} \rho(x) \quad \text { and } \quad V_{\text {eff }}:=\sum_{x \in \mathfrak{g}} \rho(x)(V)
$$

Proof. Note that $\rho(\mathfrak{g})\left(V^{\mathfrak{g}}\right)=\{0\}$ and $\rho(\mathfrak{g})\left(V_{\text {eff }}\right) \subseteq V_{\text {eff }}$, so that $V^{\mathfrak{g}}$ and $V_{\text {eff }}$ are indeed $\mathfrak{g}$-invariant. We argue by induction on $\operatorname{dim} V$. The case $\operatorname{dim} V=\{0\}$ is trivial. Since the statement of the proposition is obvious for $\rho=0$, we may assume that $\rho \neq 0$.\\
Step 1: Let $\beta_{\rho}(x, y)=\operatorname{tr}(\rho(x) \rho(y))$ denote the trace form on $\mathfrak{g}$ and $\mathfrak{a}:= \mathfrak{r a d}\left(\beta_{\rho}\right)$ denote the radical of $\beta_{\rho}$, which is an ideal because $\beta_{\rho}$ is invariant (Exercise 5.4.5). Let $\mathfrak{b} \unlhd \mathfrak{g}$ be a complementary ideal, so that $\mathfrak{g}=\mathfrak{a} \oplus \mathfrak{b}$ is a direct sum of Lie algebras (Proposition 5.5.11). In view of Cartan's Solvability Criterion (Theorem 5.4.20), the Lie algebra $\rho(\mathfrak{a})$ is a solvable ideal of $\rho(\mathfrak{g})$ because the trace form vanishes on this Lie algebra. Since $\rho(\mathfrak{g})$ is semisimple, $\rho(\mathfrak{a}) \subseteq \mathfrak{r a d}(\rho(\mathfrak{g}))=\{0\}$, so that $\mathfrak{a} \subseteq \operatorname{ker} \rho$. Conversely, the ideal $\operatorname{ker} \rho$ is contained in $\mathfrak{r a d}\left(\beta_{\rho}\right)$, which leads to $\mathfrak{a}=\operatorname{ker} \rho$. It follows in particular that $\beta:=\left.\beta_{\rho}\right|_{\mathfrak{b} \times \mathfrak{b}}$ is nondegenerate on the semisimple Lie algebra $\mathfrak{b}$.\\
Step 2: Let

$$
\Omega:=\Omega\left(\beta,\left.\rho\right|_{\mathfrak{b}}\right):=\sum_{j} \rho\left(x_{j}\right) \rho\left(x^{j}\right) \in \operatorname{End}(V)
$$

be the associated Casimir element (Definition 5.5.15). Then Lemma 5.5.16 implies that

$$
\Omega \in \operatorname{End}_{\mathfrak{b}}(V):=\{A \in \operatorname{End}(V):(\forall x \in \mathfrak{b}) A \rho(x)=\rho(x) A\}
$$

Since $\mathfrak{a}=\operatorname{ker} \rho$ this implies

$$
\Omega \in \operatorname{End}_{\mathfrak{g}}(V):=\{A \in \operatorname{End}(V):(\forall x \in \mathfrak{g}) A \rho(x)=\rho(x) A\}
$$

Finally, we note that

$$
\operatorname{tr} \Omega=\sum_{j} \operatorname{tr}\left(\rho\left(x_{j}\right) \rho\left(x^{i}\right)\right)=\sum_{j} \beta\left(x_{j}, x^{j}\right)=\operatorname{dim} \mathfrak{b}
$$

Step 3: If $V$ is the direct sum of two nonzero $\mathfrak{g}$-invariant subspaces, then $V^{\mathfrak{g}}$ and $V_{\text {eff }}$ decompose accordingly, and we can use our induction hypothesis. Let $V=V^{0}(\Omega) \oplus V^{+}(\Omega)$ be the Fitting decomposition of $V$ with respect to $\Omega$. Since $\Omega$ commutes with $\mathfrak{g}$, both summands are $\mathfrak{g}$-invariant, so that we may assume that one of these summands is trivial.

Since we assume that $\mathfrak{b} \cong \rho(\mathfrak{g}) \neq\{0\}$, we have $\operatorname{tr} \Omega>0$, so that $\Omega$ is not nilpotent, and thus $V^{+}(\Omega)$ is nonzero. Hence $V^{0}(\Omega)=\{0\}$ and, consequently, $V=V^{+}(\Omega)$. Then $\Omega$ is invertible, so that $V=V^{+}(\Omega) \subseteq V_{\text {eff }}$ and $V^{\mathfrak{g}} \subseteq V^{0}(\Omega)=\{0\}$. This completes the proof.

For the following proposition recall Definitions 5.1.12 and 5.1.15.\\
Proposition 5.5.18. For a finite-dimensional representation ( $\rho, V$ ) of the Lie algebra $\mathfrak{g}$, the following are equivalent:\\
(i) Each $\mathfrak{g}$-invariant subspace of $V$ possesses a $\mathfrak{g}$-invariant complement.\\
(ii) ( $\rho, V$ ) is completely reducible.\\
(iii) $V$ is a sum of $\mathfrak{g}$-invariant subspaces on which the representation is irreducible.

Proof. (i) $\Rightarrow$ (ii): For dimensional reasons, each $V$ contains a nonzero $\mathfrak{g}$-invariant subspace $V_{1}$ of minimal dimension. Then there exists a $\mathfrak{g}$-invariant complement $W$, so that $V=V_{1} \oplus W$. We now argue by induction on $\operatorname{dim} V$ and apply the induction hypothesis to the representation of $\mathfrak{g}$ on $W$.\\
(ii) $\Rightarrow$ (iii) is trivial.\\
(iii) $\Rightarrow$ (ii): Let $V_{1}, \ldots, V_{n}$ be a maximal set of $\mathfrak{g}$-invariant subspaces on which the representation of $\mathfrak{g}$ is irreducible and whose sum $W:=\sum_{i=1}^{n} V_{i}$ is direct. We claim that $W=V$, which implies (ii). If $W$ is a proper subspace, then (iii) implies the existence of a minimal nonzero $\mathfrak{g}$-invariant subspace $U$ not contained in $W$. Then $W \cap U=\{0\}$ follows from the minimality of $U$, so that the sum $U+\sum_{i} V_{i}$ is direct, contradicting the maximality of the set $\left\{V_{1}, \ldots, V_{n}\right\}$. This proves $W=V$.\\
(ii) $\Rightarrow$ (i): Let $V=\bigoplus_{i=1}^{n} V_{i}$ be a direct sum of minimal nonzero $\mathfrak{g}$-invariant subspaces and $W \subseteq V$ a $\mathfrak{g}$-invariant subspace. Further, let $J \subseteq\{1, \ldots, n\}$ be maximal with $W \cap\left(\sum_{i \in J} V_{i}\right)=\{0\}$. Then $W^{\prime}:=\sum_{i \in J} V_{i}$ satisfies $W \cap W^{\prime}= \{0\}$ and it remains to see that $W+W^{\prime}=V$.

Pick $i \in I$. If $i \in J$, then $V_{i} \subseteq W^{\prime} \subseteq W+W^{\prime}$. If $i \notin J$, then the maximality of $J$ implies that $\left(W^{\prime}+V_{i}\right) \cap W \neq\{0\}$ and hence $\left(W+W^{\prime}\right) \cap V_{i} \neq\{0\}$. Hence the minimality of $V_{i}$ implies that $V_{i} \subseteq W+W^{\prime}$, and this proves $V=W+W^{\prime}$.

Lemma 5.5.19. If $\mathfrak{g}$ is a real Lie algebra and $V$ a finite-dimensional $\mathfrak{g}$-module, then the following are equivalent:\\
(i) $V$ is semisimple.\\
(ii) $V_{\mathbb{C}}$ is a semisimple complex $\mathfrak{g}$-module.

Proof. (i) $\Rightarrow$ (ii): If $V$ is semisimple, then $V$ is a direct sum of simple submodules $V_{i}$, then $V_{\mathbb{C}}$ is the direct sum of the submodules $\left(V_{i}\right)_{\mathbb{C}}$. Hence it suffices to show that the complexification $W_{\mathbb{C}}$ of a simple real $\mathfrak{g}$-module $W$ is semisimple. In fact, if $W_{\mathbb{C}}$ is not simple, then let $U \subseteq W_{\mathbb{C}}$ be a nonzero minimal complex submodule. This implies in particular that $U$ is simple. Let $\sigma: W_{\mathbb{C}} \rightarrow W_{\mathbb{C}}$ be the conjugation involution defined by $\sigma(z \otimes w)=\bar{z} \otimes w$. Then $\sigma$ commutes with the action of $\mathfrak{g}$ on $W_{\mathbb{C}}$, and therefore $\sigma(U)$ also is a simple complex submodule. Now $U+\sigma(U)$ is a complex $\sigma$-invariant submodule of $W_{\mathbb{C}}$, hence of the form $X_{\mathbb{C}}$ for $X:=W \cap(U+\sigma(U))$ (Exercise 5.3.1(ii)). Then $X$ is a nonzero $\mathfrak{g}$-submodule of $W$, so that the simplicity of $W$ yields $X=W$ and thus $U+\sigma(U)=W_{\mathbb{C}}$. Now Proposition 5.5.18 shows that $W_{\mathbb{C}}$ is semisimple because it is the sum of two simple submodules.\\
(ii) $\Rightarrow$ (i): Let $W \subseteq V$ be a submodule. We have to show that there exists a module complement $U$ (Proposition 5.5.18). Since $V_{\mathbb{C}}$ is semisimple, there exists a module complement $X$ of $W_{\mathbb{C}}$ in $V_{\mathbb{C}}$, i.e., a complex linear projection $p: V_{\mathbb{C}} \rightarrow W_{\mathbb{C}}$ commuting with $\mathfrak{g}$. Let $q: \mathbb{C} \rightarrow \mathbb{R}, q(z):=\operatorname{Re} z$, denote the canonical projection onto $\mathbb{R}$. Then $q \otimes \operatorname{id}_{V}: V_{\mathbb{C}} \rightarrow V$ is a real linear projection commuting with $\mathfrak{g}$. Hence

$$
P:=\left.\left(q \otimes \mathrm{id}_{V}\right) \circ p\right|_{V}: V \rightarrow W
$$

is a $\mathfrak{g}$-equivariant real linear map with $P(V)=W$ and $\left.P\right|_{W}=\operatorname{id}_{W}$. Hence $P$ is a projection onto $W$ and $\operatorname{ker} P$ is a complementary submodule.

Lemma 5.5.20. Let $V$ be a finite-dimensional vector space and $\mathfrak{g} \subseteq \mathfrak{g} \mathfrak{l}(V)$ a commutative subalgebra consisting of semisimple elements. Then $V$ is a semisimple $\mathfrak{g}$-module.

Proof. In view of Lemma 5.5.19, it suffices to show that $V_{\mathbb{C}}$ is a semisimple complex $\mathfrak{g}$-module, resp., $\mathfrak{g}_{\mathbb{C}}$. On $V_{\mathbb{C}}$, each $x \in \mathfrak{g}$ is diagonalizable, and since $\mathfrak{g}$ is abelian, $\mathfrak{g}$ is simultaneously diagonalizable (Exercise 2.1.1(d)), so that $V_{\mathbb{C}}$ is a direct sum of one-dimensional submodules, hence semisimple.

Theorem 5.5.21 (Weyl's Theorem on Complete Reducibility). Each finite-dimensional representation of a semisimple Lie algebra is completely reducible.

Proof. We argue by induction on the dimension of the representation ( $\rho, V$ ) of the semisimple Lie algebra $\mathfrak{g}$. In view of Proposition 5.5.18, it suffices to show that each $\mathfrak{g}$-invariant subspace $W \subseteq V$ possesses a $\mathfrak{g}$-invariant complement $U$.\\
Step 1: Let $W \subseteq V$ be a $\mathfrak{g}$-invariant subspace of codimension 1 . Then the representation ( $\bar{\rho}, V / W$ ), defined by $\bar{\rho}(x)(v+W):=\rho(x) v+W$, is onedimensional. Since $\mathfrak{g}=[\mathfrak{g}, \mathfrak{g}]$ is perfect and $\mathfrak{g} \mathfrak{l}_{1}(\mathbb{K}) \cong \mathbb{K}$ is abelian, $\bar{\rho}=0$, i.e., $\rho(\mathfrak{g}) V \subseteq W$. In view of Proposition 5.5.17, $V=V^{\mathfrak{g}}+V_{\text {eff }}$, and since $V_{\text {eff }}$ is contained in $W$, there exists some $v_{o} \in V^{\mathfrak{g}} \backslash W$. Then $\mathbb{K} v_{o}$ is a $\mathfrak{g}$-invariant complement of $W$.\\
Step 2: Now let $W \subseteq V$ be an arbitrary $\mathfrak{g}$-invariant subspace. We define a representation of $\mathfrak{g}$ on $\operatorname{Hom}(V, W)$ by

$$
\pi(x) \varphi:=\left.\rho(x)\right|_{W} \circ \varphi-\varphi \circ \rho(x)
$$

(Exercise). Then the subspace

$$
U:=\left\{\varphi \in \operatorname{Hom}(V, W):\left.\varphi\right|_{W} \in \mathbb{K} \operatorname{id}_{W}\right\}
$$

is $\mathfrak{g}$-invariant because we have for $\varphi \in U$ the relation $(\pi(x) \varphi)(W)=\{0\}$ : For $\left.\varphi\right|_{W}=\lambda \mathrm{id}_{W}$ and $w \in W$, we have

$$
(\pi(x) \varphi)(w)=\rho(x) \varphi(w)-\varphi(\rho(x) w)=\rho(x)(\lambda w)-\lambda \rho(x) w=0
$$

Therefore,

$$
U_{0}:=\{\varphi \in U: \varphi(W)=\{0\}\}
$$

is a $\mathfrak{g}$-invariant subspace of $U$ of codimension 1 . Step 1 now implies the existence of a $\mathfrak{g}$-invariant $\varphi_{0} \in U \backslash U_{0}$. The $\mathfrak{g}$-invariance of $\varphi_{0}$ means that $\varphi_{0} \in \operatorname{Hom}_{\mathfrak{g}}(V, W)$ and since $\left.\varphi_{0}\right|_{W} \in \mathbb{K}^{\times} \operatorname{id}_{W}$ is invertible, $\operatorname{ker} \varphi_{0}$ is a $\mathfrak{g}$-invariant subspace complementing $W$.

\subsubsection*{Exercises for Section 5.5}
Exercise 5.5.1. Show that the dimension of a simple Lie algebra is at least 3 . Conclude that every semisimple Lie algebra of dimension $\leq 5$ is simple.

Exercise 5.5.2. For a real Lie algebra $\mathfrak{g}$, we have:\\
(i) $\mathfrak{r a d}\left(\mathfrak{g}_{\mathbb{C}}\right)=\mathfrak{r a d}(\mathfrak{g})_{\mathbb{C}}$.\\
(ii) $\mathfrak{r a d}\left(\kappa_{\mathfrak{g}}\right)_{\mathbb{C}}=\mathfrak{r a d}\left(\kappa_{\mathfrak{g}_{\mathbb{C}}}\right)$.\\
(iii) $\mathfrak{g}$ is semisimple if and only if $\mathfrak{g}_{\mathbb{C}}$ is semisimple.

Exercise 5.5.3. Let $\mathfrak{g}$ be a real Lie algebra and $\mathfrak{g}_{\mathbb{C}}$ its complexification. Show that the Cartan-Killing forms of $\mathfrak{g}$ and $\mathfrak{g}_{\mathbb{C}}$ are related by

$$
\kappa_{\mathfrak{g}}(x, y)=\kappa_{\mathfrak{g} \mathbb{C}}(x, y) \quad \text { for } \quad x, y \in \mathfrak{g} .
$$

Exercise 5.5.4. Verify the computations of the Cartan-Killing forms of $\mathfrak{s} \mathfrak{l}_{2}(\mathbb{K}), \mathfrak{s o}_{3}(\mathbb{R})$, and of the oscillator algebra in Example 5.5.4.

Exercise 5.5.5. (i) Let $\alpha: \mathfrak{g} \rightarrow \mathfrak{g} \mathfrak{l}(V)$ be a representation of the Lie algebra $\mathfrak{g}$ on $V$ and $\mathfrak{n} \unlhd \mathfrak{g}$ be an ideal. Then the space

$$
V_{0}(\mathfrak{n}):=\{v \in V \mid(\forall x \in \mathfrak{n}) \alpha(x) v=0\}
$$

is $\mathfrak{g}$-invariant.\\
(ii) Let

$$
\mathfrak{a}_{0}=\{0\} \subseteq \mathfrak{a}_{1} \subseteq \cdots \subseteq \mathfrak{a}_{n}=\mathfrak{g}
$$

be a maximal chain of ideals of $\mathfrak{g}$ and $\mathfrak{n} \unlhd \mathfrak{g}$ a nilpotent ideal. Then $\left[\mathfrak{n}, \mathfrak{a}_{j}\right] \subseteq \mathfrak{a}_{j-1}$ for $j>0$.

Exercise 5.5.6. Let $\mathfrak{g}$ be a finite-dimensional Lie algebra. Every nilpotent ideal $\mathfrak{n}$ of $\mathfrak{g}$ is orthogonal to $\mathfrak{g}$ with respect to the Cartan-Killing form.

Exercise 5.5.7. Show that $[\mathfrak{g}, \mathfrak{g}]^{\perp}=\mathfrak{r a d}(\mathfrak{g})$ for every finite-dimensional Lie algebra. Here $\perp$ refers to the Cartan-Killing form.

Exercise 5.5.8. Each one-dimensional representation ( $\pi, V$ ) of a perfect Lie algebra is trivial.

Exercise 5.5.9. Let $V$ be a finite-dimensional vector space and $V^{*}$ its dual space. Show that:\\
(a) The map $\gamma: V \otimes V^{*} \rightarrow \operatorname{End}(V), \gamma(v, \alpha)(w):=\alpha(w) v$, is a linear isomorphism.\\
(b) If $v_{1}, \ldots, v_{n}$ is a basis for $V$ and $v_{1}^{*}, \ldots, v_{n}^{*}$ the dual basis for $V^{*}$, defined by $v_{j}^{*}\left(v_{i}\right)=\delta_{i j}$, then $\gamma\left(\sum_{i=1}^{n} v_{i} \otimes v_{i}^{*}\right)=\mathrm{id}_{V}$.\\
(c) If $\beta: V \times V \rightarrow \mathbb{K}$ is a nondegenerate symmetric bilinear form, then

$$
\widetilde{\gamma}: V \otimes V \rightarrow \operatorname{End}(V), \quad \widetilde{\gamma}(v \otimes w)(x):=\beta(x, w) v
$$

is a linear isomorphism. If $v_{1}, \ldots, v_{n}$ is a basis for $V$ and $v^{1}, \ldots, v^{n} \in V$ with $\beta\left(v^{i}, \cdot\right)=v_{i}^{*}, i=1, \ldots, n$, then $\widetilde{\gamma}\left(\sum_{i=1}^{n} v_{i} \otimes v^{i}\right)=\operatorname{id}_{V}$.

Exercise 5.5.10. (cf. Exercise 13.1.4)\\
(i) Let $\mathfrak{g}$ be a simple Lie algebra over $\mathbb{C}$. Show that each invariant bilinear form $\kappa$ on $\mathfrak{g}$ is a scalar multiple of the Cartan-Killing form $\kappa_{\mathfrak{g}}$.\\
(ii) Show that the result from (i) is false for simple algebras over $\mathbb{R}$.

\subsection*{The Theorems of Levi and Malcev}
In the preceding sections, we dealt in particular with solvable and semisimple Lie algebras separately. Now we shall address the question how a finitedimensional Lie algebra $\mathfrak{g}$ decomposes into its maximal solvable ideal $\mathfrak{r a d}(\mathfrak{g})$ and the semisimple quotient $\mathfrak{g} / \mathfrak{r a d}(\mathfrak{g})$. The theorems of Levi and Malcev are fundamental for the structure theory of finite-dimensional Lie algebras. Levi's

Theorem asserts the existence of a semisimple subalgebra $\mathfrak{s}$ of $\mathfrak{g}$ complementing the radical $\mathfrak{r a d}(\mathfrak{g})$, also called a Levi complement. As a consequence, $\mathfrak{g} \cong \mathfrak{r a d}(\mathfrak{g}) \rtimes \mathfrak{s}$ is a semidirect sum. Malcev's Theorem asserts that all Levi complements are conjugate under the group of inner automorphisms of $\mathfrak{g}$, which is a uniqueness result.

\subsubsection*{Levi's Theorem}
Lemma 5.6.1. The quotient Lie algebra $\mathfrak{g} / \mathfrak{r a d}(\mathfrak{g})$ is semisimple.\\
Proof. Let $q: \mathfrak{g} \rightarrow \mathfrak{g} / \mathfrak{r a d}(\mathfrak{g})$ be the quotient homomorphism and $\mathfrak{a} \unlhd \mathfrak{g} / \mathfrak{r a d}(\mathfrak{g})$ a solvable ideal. Then $\mathfrak{b}:=q^{-1}(\mathfrak{a}) \unlhd \mathfrak{g}$ is an ideal containing $\mathfrak{r a d}(\mathfrak{g})$, for which $\mathfrak{a} \cong \mathfrak{b} / \mathfrak{r a d}(\mathfrak{g})$ is solvable. Since solvability is an extension property, $\mathfrak{b}$ is solvable, hence $\mathfrak{b} \subseteq \mathfrak{r a d}(\mathfrak{g})$, and thus $\mathfrak{a}=\{0\}$. This proves that $\mathfrak{r a d}(\mathfrak{g} / \mathfrak{r a d}(\mathfrak{g}))=\{0\}$, i.e., $\mathfrak{g} / \mathfrak{r a d}(\mathfrak{g})$ is semisimple.

Lemma 5.6.2. If $\alpha: \mathfrak{g} \rightarrow \mathfrak{h}$ is a surjective homomorphism of Lie algebras, then $\alpha(\mathfrak{r a d g})=\mathfrak{r a d h}$.

Proof. Let $\mathfrak{r}:=\mathfrak{r a d g}$. First, we note that $\alpha(\mathfrak{r})$ is a solvable ideal of $\mathfrak{h}$, hence contained in $\mathfrak{r a d}(\mathfrak{h})$. Here we use that images of ideals under surjective homomorphisms are ideals: $[\mathfrak{h}, \alpha(\mathfrak{r})]=[\alpha(\mathfrak{g}), \alpha(\mathfrak{r})]=\alpha([\mathfrak{g}, \mathfrak{r}]) \subseteq \alpha(\mathfrak{r})$.

Let $\pi: \mathfrak{h} \rightarrow \mathfrak{h} / \alpha(\mathfrak{r})$ be the quotient homomorphism. We consider the homomorphism $\beta:=\pi \circ \alpha: \mathfrak{g} \rightarrow \mathfrak{h} / \alpha(\mathfrak{r})$. In view of $\mathfrak{r} \subseteq \operatorname{ker} \beta, \beta$ factors through a surjective homomorphism $\widetilde{\beta}: \mathfrak{g} / \mathfrak{r} \rightarrow \mathfrak{h} / \alpha(\mathfrak{r})$. Since $\mathfrak{g} / \mathfrak{r}$ is semisimple (Lemma 5.6.1), the homomorphic image $\mathfrak{h} / \alpha(\mathfrak{r})$ is also semisimple. Consequently $\pi(\mathfrak{r a d h}) \subseteq \mathfrak{r a d}(\mathfrak{h} / \alpha(\mathfrak{r}))=\{0\}$, i.e., $\mathfrak{r a d h} \subseteq \alpha(\mathfrak{r})$. We thus obtain $\mathfrak{r a d h}=\alpha(\mathfrak{r})$.

Definition 5.6.3. An ideal $\mathfrak{a} \unlhd \mathfrak{g}$ is called characteristic, if it is invariant under all derivations of $\mathfrak{g}$.

Lemma 5.6.4. For the radical of the Lie algebra $\mathfrak{g}$, the following assertions hold:\\
(i) $\mathfrak{r a d}(\mathfrak{g})$ is a characteristic ideal.\\
(ii) If $\mathfrak{a} \subseteq \mathfrak{g}$ is an ideal, then $\mathfrak{r a d}(\mathfrak{a})=\mathfrak{r a d}(\mathfrak{g}) \cap \mathfrak{a}$.

Proof. (i) First, we note that $[\mathfrak{g}, \mathfrak{g}]$ is a characteristic ideal of $\mathfrak{g}$ because for each derivation $D \in \operatorname{der} \mathfrak{g}$ and $x, y \in \mathfrak{g}$ we have $D([x, y])=[D x, y]+[x, D y] \in [\mathfrak{g}, \mathfrak{g}]$. Next we note that the Cartan-Killing form is invariant (cf. Exercise 5.4.5) under $\operatorname{der}(\mathfrak{g})$ :

$$
\begin{aligned}
\kappa_{\mathfrak{g}}(D x, y) & =\operatorname{tr}(\operatorname{ad}(D x) \operatorname{ad} y)=\operatorname{tr}([D, \operatorname{ad} x] \operatorname{ad} y) \\
& =-\operatorname{tr}(\operatorname{ad} x[D, \operatorname{ad} y])=-\operatorname{tr}(\operatorname{ad} x \operatorname{ad}(D y))=-\kappa_{\mathfrak{g}}(x, D y)
\end{aligned}
$$

Therefore, $\mathfrak{r a d}(\mathfrak{g})=[\mathfrak{g}, \mathfrak{g}]^{\perp, \kappa_{\mathfrak{g}}}$ (Exercise 5.5.7) is also invariant under $\operatorname{der}(\mathfrak{g})$.\\
(ii) Clearly, $\mathfrak{r a d}(\mathfrak{g}) \cap \mathfrak{a}$ is a solvable ideal of $\mathfrak{a}$, hence contained in $\mathfrak{r a d}(\mathfrak{a})$. Since $\mathfrak{r a d}(\mathfrak{a})$ is a characteristic ideal of $\mathfrak{a}$, it is invariant under the adjoint representation of $\mathfrak{g}$ on $\mathfrak{a}$, hence a solvable ideal of $\mathfrak{g}$. This proves that $\mathfrak{r a d}(\mathfrak{a}) \subseteq \mathfrak{r a d}(\mathfrak{g})$. \(\square\)

We will need the following technical lemma.\\
Lemma 5.6.5. Let ( $\rho, V$ ) be a representation of $\mathfrak{g}$ and $\mathfrak{n} \unlhd \mathfrak{g}$ an ideal. For $v \in V$, let $\mathfrak{z}_{\mathfrak{g}}(v):=\{x \in \mathfrak{g}: \rho(x) v=0\}$. If $v \in V$ satisfies

$$
\rho(\mathfrak{g}) v=\rho(\mathfrak{n}) v \quad \text { and } \quad \mathfrak{z}_{\mathfrak{g}}(v) \cap \mathfrak{n}=\{0\}
$$

then $\mathfrak{g} \cong \mathfrak{n} \rtimes \mathfrak{z}_{\mathfrak{g}}(v)$.\\
Proof. Our assumption implies that $\mathfrak{z}_{\mathfrak{g}}(v) \cap \mathfrak{n}=\mathfrak{z}_{\mathfrak{n}}(v)=\{0\}$. The linear map $\varphi: \mathfrak{g} \rightarrow V, x \mapsto \rho(x) v$ satisfies $\varphi(\mathfrak{g})=\varphi(\mathfrak{n})$, hence $\mathfrak{g}=\mathfrak{n}+\operatorname{ker} \varphi=\mathfrak{n}+\mathfrak{z}_{\mathfrak{g}}(v)$. Since $\mathfrak{z}_{\mathfrak{g}}(v)$ is a subalgebra, the assertion follows. \(\square\)

Theorem 5.6.6 (Levi's Theorem). If $\alpha: \mathfrak{g} \rightarrow \mathfrak{s}$ is a surjective homomorphism of Lie algebras and $\mathfrak{s}$ is semisimple, then there exists a homomorphism $\beta: \mathfrak{s} \rightarrow \mathfrak{g}$ with $\alpha \circ \beta=\operatorname{id}_{\mathfrak{s}}$.\\
\includegraphics[max width=\textwidth, center]{2025_10_18_3eaf04e77f3cb364fce5g-135}

Proof. We argue by induction on the dimension of $\mathfrak{n}:=\operatorname{ker} \alpha$. For $\mathfrak{n}=\{0\}$, there is nothing to show. So we assume that $\mathfrak{n} \neq\{0\}$.

Case 1: The ideal $\mathfrak{n} \unlhd \mathfrak{g}$ is not minimal, i.e., there exists a nonzero ideal $\mathfrak{n}_{1} \subseteq \mathfrak{n}$, different from $\mathfrak{n}$. Now $\alpha$ factors through a surjective homomorphism $\alpha_{1}: \mathfrak{g} / \mathfrak{n}_{1} \rightarrow \mathfrak{s}$ with

$$
\operatorname{dim}\left(\operatorname{ker} \alpha_{1}\right)=\operatorname{dim} \mathfrak{n}-\operatorname{dim} \mathfrak{n}_{1}<\operatorname{dim} \mathfrak{n} .
$$

Therefore, our induction hypothesis implies the existence of a homomorphism $\beta_{1}: \mathfrak{s} \rightarrow \mathfrak{g} / \mathfrak{n}_{1}$ with $\alpha_{1} \circ \beta_{1}=\operatorname{id}_{\mathfrak{s}}$. Let $q: \mathfrak{g} \rightarrow \mathfrak{g} / \mathfrak{n}_{1}$ be the quotient map and $\mathfrak{b}:=q^{-1}\left(\beta_{1}(\mathfrak{s})\right)$. Then $\mathfrak{b}$ is a subalgebra of $\mathfrak{g}$ and the homomorphism

$$
\alpha_{2}:=\left.q\right|_{\mathfrak{b}}: \mathfrak{b} \rightarrow \beta_{1}(\mathfrak{s}) \cong \mathfrak{s}, \quad x \mapsto x+\mathfrak{n}_{1}
$$

is surjective. In view of $\operatorname{dim}\left(\operatorname{ker} \alpha_{2}\right)=\operatorname{dim} \mathfrak{n}_{1}<\operatorname{dim} \mathfrak{n}$, the induction hypothesis implies the existence of a homomorphism $\beta_{2}: \beta_{1}(\mathfrak{s}) \rightarrow \mathfrak{b}$ with $\alpha_{2} \circ \beta_{2}=\operatorname{id}_{\beta_{1}(\mathfrak{s})}$. Now $\beta:=\beta_{2} \circ \beta_{1}: \mathfrak{s} \rightarrow \mathfrak{g}$ is a homomorphism satisfying

$$
\alpha \circ \beta=\alpha_{1} \circ \alpha_{2} \circ \widetilde{\beta}_{2} \circ \beta_{1}=\alpha_{1} \circ \beta_{1}=\operatorname{id}_{\mathfrak{s}} .
$$

Case 2: The ideal $\mathfrak{n}$ is minimal. Since $\mathfrak{s}$ is semisimple, the radical $\mathfrak{r}:= \mathfrak{r a d}(\mathfrak{g})$ of $\mathfrak{g}$ is contained in $\mathfrak{n}$ (Lemma 5.6.2). If $\mathfrak{r}=\{0\}$, then $\mathfrak{g}$ is semisimple,\\
and the assertion follows from Proposition 5.5.11 because $\mathfrak{g}$ contains an ideal complementing $\mathfrak{n}$. So let us assume that $\mathfrak{r} \neq\{0\}$. Then the minimality of $\mathfrak{n}$ shows that $\mathfrak{n}=\mathfrak{r}$ is abelian.

The representation $\rho: \mathfrak{g} \rightarrow \mathfrak{g} \mathfrak{l}(\mathfrak{n}), x \mapsto$ ad $\left.x\right|_{\mathfrak{n}}$ satisfies $\mathfrak{n} \subseteq \operatorname{ker} \rho(\mathfrak{n}$ is abelian), hence factors through a representation $\bar{\rho}$ of $\mathfrak{s}$ on $\mathfrak{n}$, determined by $\bar{\rho} \circ \alpha=\rho$. Since $\mathfrak{n}$ is a minimal ideal of $\mathfrak{g}$, we thus obtain on $\mathfrak{n}$ an irreducible representation of $\mathfrak{s}$. If $\bar{\rho}=0$, then $\mathfrak{n}$ is central in $\mathfrak{g}$, and the adjoint representation ad: $\mathfrak{g} \rightarrow \operatorname{der}(\mathfrak{g})$ factors through a representation of $\mathfrak{s}$ on $\mathfrak{g}$. According to Weyl's Theorem, there exists an ideal of $\mathfrak{g}$ complementing $\mathfrak{n}$ (Proposition 5.5.18), and the proof is complete. We may therefore assume that $\bar{\rho}$ is nonzero.

We are now at the point where we can use Lemma 5.6.5. On $V:=\operatorname{End}(\mathfrak{g})$, we consider the representation

$$
\pi(x) \varphi:=\operatorname{ad} x \circ \varphi-\varphi \circ \operatorname{ad} x=[\operatorname{ad} x, \varphi]
$$

We consider the following three subspaces of $V$ :

$$
\begin{aligned}
P:=\operatorname{ad} \mathfrak{n} & \subseteq Q:=\{\varphi \in V: \varphi(\mathfrak{g}) \subseteq \mathfrak{n}, \varphi(\mathfrak{n})=\{0\}\} \\
& \subseteq R:=\left\{\varphi \in V: \varphi(\mathfrak{g}) \subseteq \mathfrak{n},\left.\varphi\right|_{\mathfrak{n}} \in \mathbb{K} \operatorname{id}_{\mathfrak{n}}\right\}
\end{aligned}
$$

Since $Q \subseteq R$ is the kernel of the linear map $\chi: R \rightarrow \mathbb{K}$, defined by $\left.\varphi\right|_{\mathfrak{n}}= \chi(\varphi) \operatorname{id}_{\mathfrak{n}}$, we see that $\operatorname{dim}(R / Q)=1$.

We claim that $P, Q$ and $R$ are $\mathfrak{g}$-invariant. To this end, let $y \in \mathfrak{g}$. For $x \in \mathfrak{n}$, we then have $[\operatorname{ad} y, \operatorname{ad} x]=\operatorname{ad}[y, x] \in P$, so that $P$ is $\mathfrak{g}$-invariant. To see that $R$ and $Q$ are $\mathfrak{g}$-invariant, we show that $\pi(\mathfrak{g}) R \subseteq Q$. So let $x \in \mathfrak{g}$, $\varphi \in R$ and $\left.\varphi\right|_{\mathfrak{n}}=\lambda \mathrm{id}_{\mathfrak{n}}$. For $a \in \mathfrak{n}$, we then have

$$
(\pi(x) \varphi)(a)=[x, \varphi(a)]-\varphi([x, a])=[x, \lambda a]-\lambda[x, a]=0
$$

hence $\pi(x) \varphi \in Q$. For $y \in \mathfrak{n}$ we get

$$
[\operatorname{ad} y, \varphi]=\operatorname{ad} y \circ \varphi-\varphi \circ \operatorname{ad} y=-\lambda \operatorname{ad} y \in P
$$

This proves that $\pi(\mathfrak{n}) R \subseteq P$. The ideal $\mathfrak{n}$ acts trivially on the quotient space $R / P$, which therefore inherits a representation of $\mathfrak{s} \cong \mathfrak{g} / \mathfrak{n}$.

According to Weyl's Theorem 5.5.21, there exists an $\mathfrak{s}$-invariant subspace of $R / P$ complementing $Q / P$. Clearly, this complement is one-dimensional, hence generated by the image $\bar{v}$ of one element $v \in R \backslash Q$, of which we may assume that $\left.v\right|_{\mathfrak{n}}=\operatorname{id}_{\mathfrak{n}}$. As the one-dimensional representation of $\mathfrak{s}$ on $\mathbb{R} \bar{v}$ is trivial because $\mathfrak{s}$ is perfect (Exercise 5.5.8), we see that $\pi(\mathfrak{g}) v \subseteq P$. Next we verify the assumptions of Lemma 5.6.5.

For $x \in \mathfrak{n}$, we have already seen that $\pi(x) v=[\operatorname{ad} x, v]=-\operatorname{ad} x$. If $\pi(x) v=0$, then $\operatorname{ad} x=0$, i.e., $x \in \mathfrak{z}(\mathfrak{g})$. Since $\mathfrak{n}$ is a minimal ideal of $\mathfrak{g}$ which is not central, we derive that $x=0$. This leads to $\mathfrak{z}_{\mathfrak{n}}(v)=\{0\}$ and $\pi(\mathfrak{n}) v=\operatorname{ad} \mathfrak{n}=P=\pi(\mathfrak{g}) v$. Finally, we apply Lemma 5.6.5 to complete the proof.

Definition 5.6.7. If $\mathfrak{g}$ is a finite-dimensional Lie algebra and $\mathfrak{r a d}(\mathfrak{g}) \unlhd \mathfrak{g}$ its solvable radical, then we call a subalgebra $\mathfrak{s} \leq \mathfrak{g}$ complementing the radical $\mathfrak{r a d}(\mathfrak{g})$ a Levi complement. Note that $\mathfrak{g} \cong \mathfrak{r a d}(\mathfrak{g}) \rtimes \mathfrak{s}$ holds for any Levi complement.

Corollary 5.6.8. Each finite-dimensional Lie algebra $\mathfrak{g}$ contains a semisimple Levi complement.

Proof. Let $\mathfrak{s}:=\mathfrak{g} / \mathfrak{r a d}(\mathfrak{g})$ and $\alpha: \mathfrak{g} \rightarrow \mathfrak{s}$ the quotient map. According to Lemma 5.6.1, $\mathfrak{s}$ is semisimple. Hence Theorem 5.6.6 provides a homomorphism $\beta: \mathfrak{s} \rightarrow \mathfrak{g}$ with $\alpha \circ \beta=\operatorname{id}_{\mathfrak{s}}$. Then $\beta$ is injective, so that $\beta(\mathfrak{s}) \cap \mathfrak{r a d}(\mathfrak{g})= \{0\}$ as well as $\beta(\mathfrak{s})+\mathfrak{r a d}(\mathfrak{g})=\mathfrak{g}$. Thus $\beta(\mathfrak{s})$ is a semisimple Levi complement.

Corollary 5.6.9. If $\mathfrak{s}$ is a Levi complement in $\mathfrak{g}$, then

$$
[\mathfrak{g}, \mathfrak{g}]=[\mathfrak{g}, \mathfrak{r a d}(\mathfrak{g})] \rtimes \mathfrak{s} .
$$

If $\mathfrak{r a d}(\mathfrak{g})=\mathfrak{z}(\mathfrak{g})$, then $[\mathfrak{g}, \mathfrak{g}]$ is a Levi complement.\\
Proof. The second assertion immediately follows from the first and the fact that $\mathfrak{r a d}(\mathfrak{g})=\mathfrak{z}(\mathfrak{g})$ is equivalent to $[\mathfrak{g}, \mathfrak{r a d}(\mathfrak{g})]=\{0\}$.

For the first assertion, we note that $[\mathfrak{s}, \mathfrak{s}]=\mathfrak{s}$ leads to

$$
[\mathfrak{g}, \mathfrak{g}]=[\mathfrak{g}, \mathfrak{r a d}(\mathfrak{g})]+[\mathfrak{g}, \mathfrak{s}]=[\mathfrak{g}, \mathfrak{r a d}(\mathfrak{g})]+[\mathfrak{r a d}(\mathfrak{g}), \mathfrak{s}]+[\mathfrak{s}, \mathfrak{s}]=[\mathfrak{g}, \mathfrak{r a d}(\mathfrak{g})]+\mathfrak{s} .
$$

Corollary 5.6.10. If $q: \widehat{\mathfrak{g}} \rightarrow \mathfrak{g}$ is a surjective homomorphism of finitedimensional Lie algebras, $\mathfrak{s}$ is semisimple, and $\alpha: \mathfrak{s} \rightarrow \mathfrak{g}$ is a homomorphism, then there exists a homomorphism $\widehat{\alpha}: \mathfrak{s} \rightarrow \widehat{\mathfrak{g}}$ with $q \circ \widehat{\alpha}=\alpha$.

Proof. Apply Levi's Theorem to the surjective homomorphism $q: q^{-1}(\alpha(\mathfrak{s})) \rightarrow \alpha(\mathfrak{s})$ and note that the homomorphic image $\alpha(\mathfrak{s})$ of $\mathfrak{s}$ is semisimple.

Example 5.6.11. (a) Let $V$ be a finite-dimensional vector space and $\mathfrak{g}= \mathfrak{g} \mathfrak{l}(V)$. Then

$$
\mathfrak{r a d}(\mathfrak{g})=\mathfrak{z}(\mathfrak{g})=\mathbb{K} \mathbf{1}
$$

and $\mathfrak{s} \mathfrak{l}(V)$ is a Levi complement in $\mathfrak{g}$ (Exercise).\\
(b) Let $V$ be a finite-dimensional vector space and $\mathcal{F}=\left(V_{0}, \ldots, V_{n}\right)$ a flag in $V$. Then

$$
\mathfrak{r a d}(\mathfrak{g}(\mathcal{F}))=\left\{\varphi \in \mathfrak{g}(\mathcal{F}):(\forall i)\left(\exists \lambda_{i} \in \mathbb{K}\right)\left(\varphi-\lambda_{i} \mathbf{1}\right)\left(V_{i}\right) \subseteq V_{i-1}\right\}
$$

(Exercise). Note in particular that

$$
\mathfrak{r a d}(\mathfrak{g}(\mathcal{F})) \supseteq \mathfrak{g}_{n}(\mathcal{F})=\left\{\varphi \in \mathfrak{g}(\mathcal{F}):(\forall i) \varphi\left(V_{i}\right) \subseteq V_{i-1}\right\}
$$

Choosing subspaces $W_{1}, \ldots, W_{n} \subseteq V$ with $V_{i}=W_{1} \oplus \cdots \oplus W_{i}$, we find

$$
\begin{gathered}
\mathfrak{g}(\mathcal{F}) \cong \mathfrak{g}_{n}(\mathcal{F}) \rtimes \bigoplus_{i=1}^{n} \mathfrak{g} \mathfrak{l}\left(W_{i}\right), \\
\mathfrak{r a d}(\mathfrak{g}(\mathcal{F})) \cong \mathfrak{g}_{n}(\mathcal{F}) \rtimes \mathbb{K}^{n}, \quad \text { and } \quad \mathfrak{g}(\mathcal{F}) / \mathfrak{r a d}(\mathfrak{g}(\mathcal{F})) \cong \bigoplus_{i=1}^{n} \mathfrak{s} \mathfrak{l}\left(W_{i}\right) .
\end{gathered}
$$

Remark 5.6.12. If $\mathfrak{g}$ is a solvable Lie algebra, then $\mathfrak{g}$ is isomorphic to a nested semidirect sum

$$
\left(\cdots\left(\left(\mathfrak{g}_{1} \rtimes_{\alpha_{1}} \mathfrak{g}_{2}\right) \rtimes_{\alpha_{2}} \mathfrak{g}_{3}\right) \cdots \rtimes_{\alpha_{n-1}} \mathfrak{g}_{n}\right)
$$

in which every $\mathfrak{g}_{j}$ is isomorphic to $\mathbb{K}$.\\
Composing this with Levi's Theorem and using Proposition 5.5.11, we obtain a similar factorization for arbitrary finite-dimensional Lie algebras $\mathfrak{g}$, the only difference is that the $\mathfrak{g}_{j}$ are either one-dimensional or simple. In particular, we get a sequence

$$
\mathfrak{a}_{0}=\{0\} \subseteq \mathfrak{a}_{1} \subseteq \cdots \subseteq \mathfrak{a}_{k}=\mathfrak{g}
$$

of subalgebras of $\mathfrak{g}$ for which $\mathfrak{a}_{j-1}$ is an ideal in $\mathfrak{a}_{j}$ and the quotient $\mathfrak{a}_{j} / \mathfrak{a}_{j-1}$ is either one-dimensional or simple. Such a series is called Jordan-Hölder series of $\mathfrak{g}$.

\subsubsection*{Malcev's Theorem}
Theorem 5.6.13 (Malcev's Theorem). For two Levi complements $\mathfrak{s}$ and $\mathfrak{s}^{\prime}$ in $\mathfrak{g}$, there exists some $x \in[\mathfrak{g}, \mathfrak{r a d}(\mathfrak{g})]$ with $e^{\text {ad } x} \mathfrak{s}^{\prime}=\mathfrak{s}$.

Proof. Let $\mathfrak{r}:=\mathfrak{r a d}(\mathfrak{g})$. We first consider some special cases.\\
(a) If $[\mathfrak{g}, \mathfrak{r}]=\{0\}$, then $\mathfrak{g}=\mathfrak{r} \oplus \mathfrak{s}$ is a direct sum of Lie algebras and $\mathfrak{r}=\mathfrak{z}(\mathfrak{g})$ is abelian. Therefore, $\mathfrak{s}=[\mathfrak{s}, \mathfrak{s}]=[\mathfrak{g}, \mathfrak{g}]=\left[\mathfrak{s}^{\prime}, \mathfrak{s}^{\prime}\right]=\mathfrak{s}^{\prime}$, and there is nothing to show.\\
(b) If $[\mathfrak{g}, \mathfrak{r}] \neq\{0\}$ and $\mathfrak{r}$ is a minimal nonzero ideal of $\mathfrak{g}$, then $[\mathfrak{g}, \mathfrak{r}]=\mathfrak{r}$, $[\mathfrak{r}, \mathfrak{r}]=\{0\}$ (since $D^{1}(\mathfrak{r}) \neq \mathfrak{r}$ ), and $\mathfrak{z}(\mathfrak{g})=\{0\}$ (because $\mathfrak{r} \nsubseteq \mathfrak{z}(\mathfrak{g})$ ). We define a map $h: \mathfrak{s}^{\prime} \rightarrow \mathfrak{r}$ by $x+h(x) \in \mathfrak{s}$ for $x \in \mathfrak{s}^{\prime}$, i.e., $-h$ is the projection of $\mathfrak{s}^{\prime}$ to $\mathfrak{s}$ along $\mathfrak{r}$. Since $\mathfrak{s}$ is a subalgebra and $\mathfrak{r}$ is abelian, we have

$$
[x+h(x), y+h(y)]=[x, y]+[x, h(y)]+[h(x), y] \in \mathfrak{s}
$$

Therefore,

$$
h([x, y])=[x, h(y)]+[h(x), y]
$$

This implies that

$$
\pi(x)(r, t):=([x, r]+t h(x), 0)
$$

defines a representation of $\mathfrak{s}^{\prime}$ on $\mathfrak{r} \times \mathbb{K}$. The subspace $\mathfrak{r} \cong \mathfrak{r} \times\{0\}$ is $\mathfrak{s}^{\prime}$-invariant. According to Weyl's Theorem, there exists an $\mathfrak{s}^{\prime}$-invariant complement $\mathbb{K}(v, 1)$\\
of $\mathfrak{r}$ in $\mathfrak{r} \oplus \mathbb{K}$. As $\mathfrak{s}^{\prime}$ is semisimple, $\pi\left(\mathfrak{s}^{\prime}\right)(v, 1)=\{0\}$, and hence $h(x)+[x, v]=0$ for $x \in \mathfrak{s}^{\prime}$. Now we have

$$
e^{\operatorname{ad} v} x=x+[v, x]=x+h(x) \in \mathfrak{s}
$$

for $x \in \mathfrak{s}^{\prime}$ and thus $e^{\operatorname{ad} v}\left(\mathfrak{s}^{\prime}\right) \subseteq \mathfrak{s}$. Equality follows from $\operatorname{dim} \mathfrak{s}=\operatorname{dim} \mathfrak{g} / \mathfrak{r}= \operatorname{dim} \mathfrak{s}^{\prime}$. This proves the theorem if $[\mathfrak{g}, \mathfrak{r}]$ is a nonzero minimal ideal.\\
(c) Finally, we turn to the general case. We argue by induction on $n:= \operatorname{dim} \mathfrak{r}$. The case $n=0$ is trivial, so that we assume that $n>0$ and that the assertion holds for all Lie algebras $\mathfrak{h}$ with $\operatorname{dim} \mathfrak{r a d}(\mathfrak{h})<n$. In view of (a), we may assume that $[\mathfrak{g}, \mathfrak{r}] \neq\{0\}$. As the ideal $[\mathfrak{g}, \mathfrak{r}]$ is nilpotent (Corollary 5.4.15), its center $\mathfrak{c}:=\mathfrak{z}([\mathfrak{g}, \mathfrak{r}])$ is nonzero. Let $\mathfrak{m} \neq\{0\}$ be a minimal ideal of $\mathfrak{g}$ contained in $\mathfrak{c}$. If $\mathfrak{m}=\mathfrak{r}$, then we are in the situation of (b). We therefore assume $\mathfrak{m} \neq \mathfrak{r}$. Let $\pi: \mathfrak{g} \rightarrow \mathfrak{g}_{1}:=\mathfrak{g} / \mathfrak{m}$ be the quotient map. Then $\mathfrak{r}_{1}:=\pi(\mathfrak{r})$ is the radical of $\mathfrak{g}_{1}$ (Lemma 5.6.2), and $\pi(\mathfrak{s})$ and $\pi\left(\mathfrak{s}^{\prime}\right)$ are Levi complements in $\mathfrak{g} / \mathfrak{m}$ because both are semisimple (Corollary 5.5.13) and complementing $\pi(\mathfrak{r})$. Now our induction hypothesis provides an $x_{1} \in\left[\mathfrak{g}_{1}, \mathfrak{r}_{1}\right]$ with $e^{\text {ad } x_{1}} \pi\left(\mathfrak{s}^{\prime}\right)= \pi(\mathfrak{s})$. Using $\pi([\mathfrak{g}, \mathfrak{r}])=\left[\mathfrak{g}_{1}, \mathfrak{r}_{1}\right]$, we find an $x \in[\mathfrak{g}, \mathfrak{r}]$ with $\pi(x)=x_{1}$. Then $e^{\operatorname{ad} x_{1}} \pi\left(\mathfrak{s}^{\prime}\right)=\pi\left(e^{\operatorname{ad} x_{\mathfrak{s}^{\prime}}}\right) \subseteq \pi(\mathfrak{s})$, i.e.,

$$
e^{\operatorname{ad} x} \mathfrak{s}^{\prime} \subseteq \mathfrak{h}:=\mathfrak{s}+\mathfrak{m} .
$$

Now $e^{\operatorname{ad} x_{\mathfrak{s}^{\prime}}}$ and $\mathfrak{s}$ are two Levi complements in the Lie algebra $\mathfrak{h}$ with $\operatorname{dim} \mathfrak{r a d}(\mathfrak{h})=\operatorname{dim} \mathfrak{m}<n=\operatorname{dim} \mathfrak{r}$. Hence the induction hypothesis provides a $y \in \mathfrak{m}$ with $e^{\text {ady }} e^{\text {ad } x} \mathfrak{s}^{\prime} \subseteq \mathfrak{s}$. Since $\mathfrak{m}$ is central in $[\mathfrak{g}, \mathfrak{r}]$, we have $[x, y]=0$, and therefore $e^{\operatorname{ad} y} e^{\operatorname{ad} x} \mathfrak{s}^{\prime}=e^{\operatorname{ad}(x+y)} \mathfrak{s}^{\prime} \subseteq \mathfrak{s}$.

Malcev's Theorem has interesting consequences:\\
Corollary 5.6.14. Each semisimple subalgebra of $\mathfrak{g}$ is contained in a Levi complement. In particular, the Levi complements are precisely the maximal semisimple subalgebras of $\mathfrak{g}$.

Proof. Let $\mathfrak{r}:=\mathfrak{r a d}(\mathfrak{g})$ be the radical of $\mathfrak{g}, \mathfrak{h} \subseteq \mathfrak{g}$ a semisimple subalgebra, and $\mathfrak{a}:=\mathfrak{r}+\mathfrak{h}$. Then $\mathfrak{a}$ is a subalgebra of $\mathfrak{g}$ and $\mathfrak{r}$ is a solvable ideal of $\mathfrak{a}$. Since the solvable ideal $\mathfrak{r a d}(\mathfrak{a}) \cap \mathfrak{h}$ of the semisimple Lie algebra $\mathfrak{h}$ is trivial, we see that $\mathfrak{r}=\mathfrak{r a d}(\mathfrak{a})$. The ideal $\mathfrak{h} \cap \mathfrak{r}$ of $\mathfrak{h}$ is solvable and semisimple, hence trivial. This proves that $\mathfrak{h}$ is a Levi complement in $\mathfrak{a}$.

Let $\mathfrak{s}$ be a Levi complement in $\mathfrak{g}$. Then $\mathfrak{a}=\mathfrak{r}+(\mathfrak{a} \cap \mathfrak{s})$ is a semidirect sum and since $\mathfrak{a} \cap \mathfrak{s} \cong \mathfrak{a} / \mathfrak{r} \cong \mathfrak{h}$ is semisimple, $\mathfrak{a} \cap \mathfrak{s}$ is a Levi complement in $\mathfrak{a}$. According to Malcev's Theorem 5.6.13, there exists an $x \in[\mathfrak{a}, \mathfrak{r}]$ with $e^{\operatorname{ad} x}(\mathfrak{a} \cap \mathfrak{s})=\mathfrak{h}$, i.e., $\mathfrak{h}$ is a contained in the Levi complement $e^{\operatorname{ad} x}(\mathfrak{s})$ of $\mathfrak{g}$.

Corollary 5.6.15. If $\mathfrak{n} \unlhd \mathfrak{g}$ is an ideal of $\mathfrak{g}$ and $\mathfrak{g}=\mathfrak{r} \rtimes \mathfrak{s}$ a Levi decomposition, i.e., $\mathfrak{s}$ is a Levi complement, then $\mathfrak{n}=(\mathfrak{n} \cap \mathfrak{r}) \rtimes(\mathfrak{n} \cap \mathfrak{s})$ is a Levi decomposition of $\mathfrak{n}$.

Proof. We have already seen in Lemma 5.6.4 that $\mathfrak{n} \cap \mathfrak{r}=\mathfrak{r a d}(\mathfrak{n})$. If $\mathfrak{s}_{\mathfrak{n}}$ is a Levi complement in $\mathfrak{n}$, then Corollary 5.6.14 implies that the semisimple Lie algebra $\mathfrak{s}_{\mathfrak{n}}$ is contained in a Levi complement $\mathfrak{s}^{\prime}$ of $\mathfrak{g}$. For $x \in[\mathfrak{g}, \mathfrak{r}]$ with $e^{\text {ad } x} \mathfrak{s}^{\prime}=\mathfrak{s}$, we now see that $e^{\text {ad } x} \mathfrak{s}_{\mathfrak{n}} \subseteq \mathfrak{n} \cap \mathfrak{s}$ because

$$
e^{\operatorname{ad} x} \mathfrak{n} \subseteq \mathfrak{n}+[x, \mathfrak{n}] \subseteq \mathfrak{n}
$$

Since the ideal $\mathfrak{n} \cap \mathfrak{s}$ of $\mathfrak{s}$ is semisimple (Corollary 5.5.13) and $\mathfrak{s}_{\mathfrak{n}}$ is maximal semisimple in $\mathfrak{n}$, we obtain $e^{\text {ad } x} \mathfrak{s}_{\mathfrak{n}}=\mathfrak{n} \cap \mathfrak{s}$. This shows that $\mathfrak{n} \cap \mathfrak{s}$ is a Levi complement in $\mathfrak{n}$.

\subsubsection*{Exercises for Section 5.6}
Exercise 5.6.1. Let $\mathcal{F}=\left(V_{0}, \ldots, V_{n}\right)$ be a flag in the finite-dimensional vector space $V$. Determine a Levi decomposition of the Lie algebra $\mathfrak{g}(\mathcal{F})$ (Example 5.6.11).

Exercise 5.6.2. Show that for every Jordan-Hölder series

$$
\mathfrak{a}_{0}=\{0\} \subseteq \mathfrak{a}_{1} \subseteq \cdots \subseteq \mathfrak{a}_{k}=\mathfrak{g}
$$

of subalgebras of $\mathfrak{g}$ (i.e., $\mathfrak{a}_{j-1}$ is an ideal in $\mathfrak{a}_{j}$ and the quotient $\mathfrak{a}_{j} / \mathfrak{a}_{j-1}$ is either one-dimensional or simple), the set of quotients $\left\{\mathfrak{a}_{j} / \mathfrak{a}_{j-1}: j=1, \ldots, k\right\}$ does not depend on the Jordan-Hölder series (cf. Remark 5.6.12).

\subsection*{Reductive Lie Algebras}
We conclude this chapter with a brief discussion of reductive Lie algebras. This class of Lie algebras is only slightly larger than the class of semisimple Lie algebras, but they occur naturally. It contains the abelian Lie algebras. Reductive Lie algebras often occur as stabilizer subalgebras inside semisimple Lie algebras. Thus they appear frequently in proofs by induction on the dimension.

Definition 5.7.1. We call a finite-dimensional Lie algebra $\mathfrak{g}$ reductive if $\mathfrak{g}$ is a semisimple module with respect to the adjoint representation, i.e., for each ideal $\mathfrak{a} \unlhd \mathfrak{g}$, there exists an ideal $\mathfrak{b} \unlhd \mathfrak{g}$ with $\mathfrak{g}=\mathfrak{a} \oplus \mathfrak{b}$.

Lemma 5.7.2. For a reductive Lie algebra $\mathfrak{g}$, the following assertions hold:\\
(i) If $\mathfrak{n} \unlhd \mathfrak{g}$ is an ideal, then $\mathfrak{n}$ and $\mathfrak{g} / \mathfrak{n}$ are reductive.\\
(ii) $\mathfrak{g}=\mathfrak{z}(\mathfrak{g}) \oplus[\mathfrak{g}, \mathfrak{g}]$ and $[\mathfrak{g}, \mathfrak{g}]$ is semisimple.\\
(iii) $\mathfrak{g}$ is semisimple if and only if $\mathfrak{z}(\mathfrak{g})=\{0\}$.

Proof. (i) Since $\mathfrak{g}$ is reductive, there exists an ideal $\mathfrak{b} \unlhd \mathfrak{g}$ with $\mathfrak{g}=\mathfrak{n} \oplus \mathfrak{b}$. Then $[\mathfrak{b}, \mathfrak{n}]=\{0\}$, so that $\mathfrak{g}$ is a direct sum of Lie algebras. As submodules of the semisimple $\mathfrak{g}$-module $\mathfrak{g}$, the ideals $\mathfrak{n}$ and $\mathfrak{b}$ are semisimple $\mathfrak{g}$-modules, and since the complementary ideals do not act on each other, it follows that $\mathfrak{n}$ and $\mathfrak{b} \cong \mathfrak{g} / \mathfrak{n}$ are reductive Lie algebras.\\
(ii) Let $\mathfrak{a} \subseteq \mathfrak{g}$ be an ideal complement of $[\mathfrak{g}, \mathfrak{g}]$. Then $\mathfrak{g}=\mathfrak{a} \oplus[\mathfrak{g}, \mathfrak{g}]$, and $[\mathfrak{g}, \mathfrak{a}] \subseteq \mathfrak{a} \cap[\mathfrak{g}, \mathfrak{g}]=\{0\}$ implies that $\mathfrak{a}$ is central. Further, (i) implies that $[\mathfrak{g}, \mathfrak{g}]$ is reductive. To see that $\mathfrak{z}(\mathfrak{g})$ is not larger than $\mathfrak{a}$, we choose an ideal $\mathfrak{b}$ of $[\mathfrak{g}, \mathfrak{g}]$ complementing $\mathfrak{z}(\mathfrak{g}) \cap[\mathfrak{g}, \mathfrak{g}]$. Then $[\mathfrak{g}, \mathfrak{g}]=[\mathfrak{b}, \mathfrak{b}] \subseteq \mathfrak{b}$ yields $\mathfrak{z}(\mathfrak{g}) \cap[\mathfrak{g}, \mathfrak{g}]=\{0\}$, and hence $\mathfrak{z}(\mathfrak{g})=\mathfrak{a}$.

Since $[\mathfrak{g}, \mathfrak{g}]$ is reductive, it is a direct sum of simple modules $\mathfrak{g}_{1}, \ldots, \mathfrak{g}_{m}$ for the adjoint representation. The preceding argument implies that none of these ideals is abelian, hence they are simple Lie algebras and thus $[\mathfrak{g}, \mathfrak{g}]$ is semisimple.\\
(iii) If $\mathfrak{z}(\mathfrak{g})=\{0\}$, then (ii) implies that $\mathfrak{g}$ is semisimple. If, conversely, $\mathfrak{g}$ is semisimple, then $\mathfrak{z}(\mathfrak{g}) \subseteq \mathfrak{r a d}(\mathfrak{g})=\{0\}$.\\
Proposition 5.7.3. For a finite-dimensional Lie algebra $\mathfrak{g}$, the following are equivalent:\\
(i) $\mathfrak{g}$ is reductive.\\
(ii) $[\mathfrak{g}, \mathfrak{g}]$ is semisimple.\\
(iii) $\mathfrak{r a d}(\mathfrak{g})$ is central in $\mathfrak{g}$.

Proof. (i) $\Rightarrow$ (ii) follows from Lemma 5.7.2.\\
(ii) $\Rightarrow$ (iii): Let $\mathfrak{g}=\mathfrak{r} \rtimes \mathfrak{s}$ be a Levi decomposition of $\mathfrak{g}$ with $\mathfrak{r}=\mathfrak{r a d}(\mathfrak{g})$. Then $[\mathfrak{g}, \mathfrak{g}]=[\mathfrak{g}, \mathfrak{r}] \rtimes \mathfrak{s}$, so that the semisimplicity of $[\mathfrak{g}, \mathfrak{g}]$ implies that $[\mathfrak{g}, \mathfrak{r}]= \{0\}$.\\
(iii) $\Rightarrow$ (i): If $\mathfrak{r}$ is central in $\mathfrak{g}$, then any Levi decomposition $\mathfrak{g}=\mathfrak{r} \rtimes \mathfrak{s}$ is a direct sum $\mathfrak{g}=\mathfrak{r} \oplus \mathfrak{s}$, where $\mathfrak{r}$ is a central ideal. Since $\mathfrak{z}(\mathfrak{s})=\{0\}$, we immediately get $\mathfrak{z}(\mathfrak{g})=\mathfrak{r}$, so that $\mathfrak{g}=\mathfrak{z}(\mathfrak{g}) \oplus \mathfrak{s}$. Thus $\mathfrak{g}$ is a direct sum of simple submodules with respect to the adjoint representation, hence a semisimple module and this means that $\mathfrak{g}$ is reductive.\\
Proposition 5.7.4. The ideal $[\mathfrak{g}, \mathfrak{r a d}(\mathfrak{g})]$ coincides with the intersection of the kernels of finite dimensional irreducible representations.\\
Proof. If $\rho: \mathfrak{g} \rightarrow \mathfrak{g l}(V)$ is a finite dimensional irreducible representation and $\mathfrak{n}:=[\mathfrak{g}, \mathfrak{r a d}(\mathfrak{g})]$, then each subspace $V_{k}:=\rho(\mathfrak{n})^{k}(V)$ is $\mathfrak{g}$-invariant because

$$
\begin{aligned}
& \rho(x) \rho\left(x_{1}\right) \cdots \rho\left(x_{k}\right) v \\
& =\sum_{j=1}^{k} \rho\left(x_{1}\right) \cdots\left[\rho(x), \rho\left(x_{j}\right)\right] \cdots \rho\left(x_{k}\right) v+\rho\left(x_{1}\right) \cdots \rho\left(x_{k}\right) \rho(x) v
\end{aligned}
$$

for each $v \in V$. According to Proposition 5.4.14, there exists a minimal $m$ with $V_{m}=\{0\}$. Then $V_{m-1} \neq\{0\}$, so that the irreducibility of the representation implies that $V_{m-1}=V$. Hence $\rho(\mathfrak{n}) V=\rho(\mathfrak{n}) V_{m-1} \subseteq V_{m}=\{0\}$, i.e., $\mathfrak{n} \subseteq \operatorname{ker} \rho$.

Next we consider the quotient $\mathfrak{q}:=\mathfrak{g} /[\mathfrak{g}, \mathfrak{r a d}(\mathfrak{g})]$ and the quotient homomorphism $q: \mathfrak{g} \rightarrow \mathfrak{q}$. In view of Lemma 5.6.2, $q(\mathfrak{r a d}(\mathfrak{g}))=\mathfrak{r a d}(\mathfrak{q})$, so that

$$
[\mathfrak{q}, \mathfrak{r a d}(\mathfrak{q})]=[q(\mathfrak{g}), q(\mathfrak{r a d}(\mathfrak{g}))]=q([\mathfrak{g}, \mathfrak{r a d}(\mathfrak{g})])=\{0\}
$$

implies that $\mathfrak{r a d}(\mathfrak{q})$ is central. Now Proposition 5.7.3 shows that $\mathfrak{q}$ is reductive.\\
It remains to observe that for each non-zero $x \in \mathfrak{q}$, there exists an irreducible finite dimensional representation $\rho: \mathfrak{q} \rightarrow \mathfrak{g} \mathfrak{l}(V)$ with $\rho(x) \neq 0$. Then $\rho \circ q: \mathfrak{g} \rightarrow \mathfrak{g} \mathfrak{l}(V)$ is an irreducible representation of $\mathfrak{g}$, this implies the assertion. We know that $\mathfrak{q}=\mathfrak{z}(\mathfrak{q}) \oplus[\mathfrak{q}, \mathfrak{q}]$ and that $[\mathfrak{q}, \mathfrak{q}]$ is semisimple, hence a sum of simple ideals $\mathfrak{q}_{1}, \ldots, \mathfrak{q}_{k}$. If $x \notin \mathfrak{z}(\mathfrak{q})$, then $x$ projects to a non-zero element of some $\mathfrak{q}_{j}$, and then $\rho(y):=\left.\operatorname{ad} y\right|_{\mathfrak{q}_{j}}$ defines an irreducible representation with $\rho(x) \neq 0$. If $x \in \mathfrak{z}(\mathfrak{g})$, then it suffices to use a linear functional $\lambda: \mathfrak{q} \rightarrow \mathbb{K}$ with $[\mathfrak{q}, \mathfrak{q}] \subseteq \operatorname{ker} \lambda$ with $\lambda(x) \neq 0$ and to observe that $\rho(y):=\lambda(y) \mathbf{1}$ defines a one-dimensional representation with $\rho(x) \neq 0$.

\subsubsection*{Notes on Chapter 5}
The Jacobi identity was discovered in around 1830 by Carl Gustav Jacob Jacobi (1804-1851) as an identity for the Poisson bracket $\{\cdot, \cdot\}$ on smooth functions on $\mathbb{R}^{2 n}$ (Exercise 5.1.2).

The term Lie algebra was introduced in the 1920s by Hermann Weyl, following a suggestion of N. Jacobson. Lie himself was dealing mainly with Lie algebras of vector fields (Exercise 5.1.3), which he called (infinitesimal) transformation groups. The term Lie group was introduced by É. Cartan.

The Jordan decompositions and the Jordan normal form are due to Camille Jordan (1838-1922). In the 1870s, he wrote a textbook on Galois theory of polynomial equations, thus making Galois' ideas available to the mathematical world. This promoted group theoretical ideas considerably. In particular, it inspired Sophus Lie to work on a "Galois theory" for differential equations, using symmetries of differential equations to understand the structure of their solutions.

In the original proof of his theorem, Weyl used the famous "unitary trick". For $\mathbb{K}=\mathbb{C}$, one can derive Weyl's theorem on complete reducibility from the representation theory of compact groups (cf. Exercise 13.2.5). For $\mathfrak{g}=\mathfrak{s} \mathfrak{l}_{n}(\mathbb{R})$ this works roughly as follows. One shows that the complex representations of $\mathfrak{g}$ are in one-to-one correspondence with the complex representations of $\mathfrak{s} \mathfrak{l}_{n}(\mathbb{C})$, resp., its real form $\mathfrak{s u}_{n}(\mathbb{C})$, hence further with unitary representations of the compact simply connected Lie group $\mathrm{SU}_{n}(\mathbb{C})$. For unitary representations, complete irreducibility is a simple consequence of the existence of an invariant inner product. A purely algebraic proof was found later in the 1935 by H.B.G. Casimir and B.L. van der Waerden [CW35], after H.B.G. Casimir had dealt with the case $\mathfrak{s} \mathfrak{l}_{2}(\mathbb{C})$ using the operator named after him. Another algebraic proof was found in 1935 by R. Brauer [Br36]. A completely different approach based on Lie algebra cohomology has been developed by J. H. C. Whitehead (see the first Whitehead Lemma 7.5.26).

The original proof of Levi's Theorem for complex Lie algebras [Le05] was based on the classification of simple Lie algebras. The classification free proof for real Lie algebras given here goes back to J. H. C. Whitehead [Wh36]. The conjugacy of the Levi complements was shown by A. I. Malcev in [Ma42].

For a detailed account of the early history of Lie theory up to 1926, we refer to the book [Haw00] of Th. Hawkins.

\section*{Root Decomposition}
Since a simple Lie algebra $\mathfrak{g}$ has no other ideals than $\mathfrak{g}$ and $\{0\}$, we cannot analyze its structure by breaking it up into an ideal $\mathfrak{n}$ and the corresponding quotient algebra $\mathfrak{g} / \mathfrak{n}$. We therefore need refined tools to look inside simple Lie algebras. It turns out that Cartan subalgebras and the corresponding root decompositions provide such a tool.

Roots and root spaces have remarkable properties, some of which one turns into a system of axioms for abstract root systems. We derive a number of additional properties from these axioms. Moreover, we define certain objects associated with abstract roots systems like Weyl groups and Weyl chambers. Using these structural elements, one could proceed rather easily to a complete classification of complex simple Lie algebras, but we refrain from doing this since our emphasis is on structure rather than classification.

In the present chapter, we first develop the concept of a Cartan subalgebra and root decompositions for general Lie algebras. Then we turn to semisimple Lie algebras, and we finally discuss the geometry of the root system.

\subsection*{Cartan Subalgebras}
For a better understanding of the structure of a Lie algebra, one decomposes it into simultaneous eigenspaces of linear maps of the form ad $x$. Cartan subalgebras $\mathfrak{h} \leq \mathfrak{g}$ provide maximal subsets of $\mathfrak{g}$, for which there exist simultaneous generalized eigenspaces for all operators ad $x, x \in \mathfrak{h}$, whenever $\mathfrak{g}$ is a complex Lie algebra.

\subsubsection*{Weight and Root Decompositions}
Root decompositions are the simultaneous eigenspace decompositions of the type mentioned above. They are special cases of weight decompositions.

Definition 6.1.1. Let ( $\pi, V$ ) be a representation of the Lie algebra $\mathfrak{h}$. For a function $\lambda: \mathfrak{h} \rightarrow \mathbb{K}$, we define the corresponding weight space and the corresponding generalized weight space by

$$
V_{\lambda}(\mathfrak{h}):=\bigcap_{x \in \mathfrak{h}} V_{\lambda(x)}(\pi(x)) \quad \text { and } \quad V^{\lambda}(\mathfrak{h}):=\bigcap_{x \in \mathfrak{h}} V^{\lambda(x)}(\pi(x)) .
$$

Any function $\lambda: \mathfrak{h} \rightarrow \mathbb{K}$ for which $V^{\lambda}(\mathfrak{h}) \neq\{0\}$ is called a weight of the representation ( $\pi, V$ ). We write $\mathcal{P}_{\mathfrak{h}}(V)$ for the set of weights of ( $\pi, V$ ).

Remark 6.1.2. Suppose that we have a subalgebra $\mathfrak{h} \subseteq \mathfrak{g} \mathfrak{l}(V)$ for which

$$
V=\bigoplus_{\lambda} V^{\lambda}(\mathfrak{h})
$$

is a direct sum of generalized weight spaces which are $\mathfrak{h}$-invariant. We claim that $\mathfrak{h}$ is nilpotent. For each $x \in \mathfrak{h}$, let $x=x_{s}+x_{n}$ denote the Jordan decomposition, $\mathfrak{h}_{s}:=\left\{x_{s}: x \in \mathfrak{h}\right\}$ and $\mathfrak{h}_{n}:=\left\{x_{n}: x \in \mathfrak{h}\right\}$. For each weight $\lambda$ we have

$$
V^{\lambda}(\mathfrak{h})=\bigcap_{x \in \mathfrak{h}} V_{\lambda(x)}\left(x_{s}\right),
$$

so that the weight space decomposition yields a simultaneous diagonalization of the set $\mathfrak{h}_{s}$. Moreover, for each $x \in \mathfrak{h}$, all eigenspaces of $x_{s}$ are $\mathfrak{h}$-invariant because they are a sum of certain weight spaces, and this implies that $\left[\mathfrak{h}, x_{s}\right]=0$.

Let $\pi_{\lambda}: \mathfrak{h} \rightarrow \mathfrak{g} \mathfrak{l}\left(V^{\lambda}(\mathfrak{h})\right)$ denote the representation of $\mathfrak{h}$ on the weight space $V^{\lambda}(\mathfrak{h})$. For any $x \in \mathfrak{h}$, we then have $\pi_{\lambda}(x)_{s}=\lambda(x) \mathbf{1}$. With $n:=\operatorname{dim} V^{\lambda}(\mathfrak{h})$ we now see that

$$
\lambda(x)=\frac{1}{n} \operatorname{tr}\left(\pi_{\lambda}(x)\right)
$$

is a linear functional on $\mathfrak{h}$ vanishing on all brackets. In particular, $\operatorname{ker} \lambda$ is an ideal of $\mathfrak{h}$. For $x, y \in \mathfrak{h}$, we further derive that

$$
\left[\pi_{\lambda}(x)-\lambda(x) \mathbf{1}, \pi_{\lambda}(y)-\lambda(y) \mathbf{1}\right]=\left[\pi_{\lambda}(x), \pi_{\lambda}(y)\right]=\pi_{\lambda}([x, y]) \in \pi_{\lambda}(\operatorname{ker} \lambda)
$$

so that $\mathfrak{a}:=\left\{\pi_{\lambda}(x)-\lambda(x) \mathbf{1}: x \in \mathfrak{h}\right\}$ is a Lie subalgebra of $\mathfrak{g} \mathfrak{l}\left(V^{\lambda}(\mathfrak{h})\right)$ consisting of nilpotent elements, hence nilpotent by Corollary 5.2.7. Now

$$
\pi_{\lambda}(\mathfrak{h}) \subseteq \mathfrak{a}+\mathbb{K} \mathbf{1} \cong \mathfrak{a} \oplus \mathbb{K}
$$

is a subalgebra of a nilpotent Lie algebra, hence nilpotent.\\
Finally, we observe that we have an inclusion $\mathfrak{h} \hookrightarrow \bigoplus_{\lambda} \pi_{\lambda}(\mathfrak{h})$ of $\mathfrak{h}$ into a nilpotent Lie algebra, so that $\mathfrak{h}$ is nilpotent.

In the preceding remark, we have seen that we can only expect that a representation decomposes into invariant generalized weight spaces if $\mathfrak{h}$ is nilpotent.

Lemma 6.1.3. Let ( $\pi, V$ ) be a finite-dimensional representation of the nilpotent Lie algebra $\mathfrak{h}$ such that every $\pi(x), x \in \mathfrak{h}$, is split. Then each weight is linear and $V$ decomposes as

$$
V=\bigoplus_{\lambda \in \mathfrak{h}^{*}} V^{\lambda}(\mathfrak{h}) .
$$

Moreover, each generalized weight space $V^{\lambda}(\mathfrak{h})$ is $\mathfrak{h}$-invariant.

Proof. Since the assertion of the lemma only refers to the Lie algebra $\pi(\mathfrak{h})$, we may replace $\mathfrak{h}$ by $\pi(\mathfrak{h})$ and assume that $\mathfrak{h} \subseteq \mathfrak{g} \mathfrak{l}(V)$ is a nilpotent subalgebra consisting of split endomorphisms.

For each $x \in \mathfrak{h}$, we have $\operatorname{ad} x(\mathfrak{h}) \subseteq \mathfrak{h}$, so that we also get $(\operatorname{ad} x)_{s}(\mathfrak{h}) \subseteq \mathfrak{h}$. Since $\left.\operatorname{ad} x\right|_{\mathfrak{h}}$ is nilpotent, Corollary 5.3.9 shows that $0=\left.(\operatorname{ad} x)_{s}\right|_{\mathfrak{h}}=\left.\operatorname{ad} x_{s}\right|_{\mathfrak{h}}$. It follows that $\left[x_{s}, \mathfrak{h}\right]=\{0\}$. In view of Theorem 5.3.3(iv), we further have $\left[x_{s}, y_{s}\right]=0$ for $x, y \in \mathfrak{h}$, so that $\mathfrak{h}_{s}:=\left\{x_{s}: x \in \mathfrak{h}\right\}$ is a commutative set of diagonalizable endomorphisms, hence simultaneously diagonalizable (Exercise 2.1.1(d)). Let $\widetilde{\lambda}: \mathfrak{h}_{s} \rightarrow \mathbb{K}$ be a map and $V_{\widetilde{\lambda}}\left(\mathfrak{h}_{s}\right)=\bigcap_{x \in \mathfrak{h}} V_{\widetilde{\lambda}\left(x_{s}\right)}\left(x_{s}\right)$ be the corresponding simultaneous eigenspace, so that

$$
V=\bigoplus_{\tilde{\lambda}} V_{\tilde{\lambda}}\left(\mathfrak{h}_{s}\right)
$$

(cf. Exercise 2.1.1(e)).\\
In view of $\left[\mathfrak{h}, \mathfrak{h}_{s}\right]=\{0\}$ and Exercise 2.1.1(a), $V_{\widetilde{\lambda}}\left(\mathfrak{h}_{s}\right)$ is $\mathfrak{h}$-invariant. Let $\pi_{\widetilde{\lambda}}$ denote the representation of $\mathfrak{h}$ on this subspace. For any $x \in \mathfrak{h}$, we then have $\pi_{\widetilde{\lambda}}(x)_{s}=\widetilde{\lambda}\left(x_{s}\right) \mathbf{1}$. For $n:=\operatorname{dim} V_{\widetilde{\lambda}}\left(\mathfrak{h}_{s}\right)$, we now see that

$$
\lambda(x)=\widetilde{\lambda}\left(x_{s}\right)=\frac{1}{n} \operatorname{tr}\left(\pi_{\widetilde{\lambda}}(x)\right)
$$

defines a linear functional on $\mathfrak{h}$, satisfying

$$
V_{\widetilde{\lambda}}\left(\mathfrak{h}_{s}\right)=\bigcap_{x \in \mathfrak{h}} V_{\lambda(x)}\left(x_{s}\right)=\bigcap_{x \in \mathfrak{h}} V^{\lambda(x)}(x)=V^{\lambda}(\mathfrak{h}) .
$$

This completes the proof.\\
Definition 6.1.4. If $\mathfrak{h}$ is a nilpotent subalgebra of the Lie algebra $\mathfrak{g}$, then the weights of the representation $\pi=\left.\operatorname{ad}\right|_{\mathfrak{h}}$ which are different from zero are called roots of $\mathfrak{g}$ with respect to $\mathfrak{h}$. The set of all roots is denoted $\Delta(\mathfrak{g}, \mathfrak{h})$. The (generalized) weight spaces $\mathfrak{g}^{\lambda}(\mathfrak{h})$ are called root spaces. Sometimes we write $\mathfrak{g}^{\lambda}$ instead of $\mathfrak{g}^{\lambda}(\mathfrak{h})$. If $0 \neq \mu \in \mathfrak{h}^{*}$ is not a root, we put $\mathfrak{g}^{\mu}:=\{0\}$.\\
Proposition 6.1.5. Let $\mathfrak{g}$ be a finite-dimensional Lie algebra and $\mathfrak{h}$ a nilpotent subalgebra of $\mathfrak{g}$.\\
(i) $\left[\mathfrak{g}^{\lambda}, \mathfrak{g}^{\mu}\right] \subseteq \mathfrak{g}^{\lambda+\mu}$ for all $\lambda, \mu \in \mathfrak{h}^{*}$.\\
(ii) $\mathfrak{g}^{0}$ is a subalgebra of $\mathfrak{g}$.

Proof. (i) For $x \in \mathfrak{g}^{\lambda}, y \in \mathfrak{g}^{\mu}$ and $h \in \mathfrak{h}$, we have

$$
\begin{aligned}
& (\operatorname{ad}(h)-\lambda(h)-\mu(h))^{n}([x, y]) \\
& =\sum_{k=0}^{n}\binom{n}{k}\left[(\operatorname{ad}(h)-\lambda(h))^{k} x,(\operatorname{ad}(h)-\mu(h))^{n-k} y\right]
\end{aligned}
$$

(Exercise 5.3.7). If $n$ is sufficiently large, then for every summand either the left factor or the right factor in the bracket vanishes, so that the whole sum vanishes. This proves that $[x, y] \in \mathfrak{g}^{\lambda+\mu}$.\\
(ii) is a direct consequence of (i).

\subsubsection*{Cartan Subalgebras}
In the following, we want to decompose a Lie algebra into root spaces with respect to a nilpotent subalgebra $\mathfrak{h}$. Since we want such decompositions to be as fine as possible, $\mathfrak{g}^{0}(\mathfrak{h})$ should be as small as possible, hence equal to $\mathfrak{h}$. The following result prepares the definition of a Cartan subalgebra.

Proposition 6.1.6. For a subalgebra $\mathfrak{h}$ of a finite-dimensional Lie algebra $\mathfrak{g}$, the following are equivalent:\\
(C1) $\mathfrak{h}$ is nilpotent and self-normalizing, i.e., $\mathfrak{h}=\mathfrak{n}_{\mathfrak{g}}(\mathfrak{h})$.\\
(C2) $\mathfrak{h}=\mathfrak{g}^{0}(\mathfrak{h})$.\\
If these conditions are satisfied, then $\mathfrak{h}$ is a maximal nilpotent subalgebra of $\mathfrak{g}$.\\
Proof. (C1) $\Rightarrow$ (C2): Since $\mathfrak{h}$ is nilpotent, we have $\mathfrak{h} \subseteq \mathfrak{g}^{0}(\mathfrak{h})$. Assume that $\mathfrak{g}^{0}(\mathfrak{h})$ strictly contains $\mathfrak{h}$ and consider the representation

$$
\pi: \mathfrak{h} \rightarrow \mathfrak{g} \mathfrak{l}\left(\mathfrak{g}^{0}(\mathfrak{h}) / \mathfrak{h}\right), \quad \pi(h)(x+\mathfrak{h}):=[h, x]+\mathfrak{h}
$$

whose image consists of nilpotent endomorphisms. By Theorem 5.2.5, there exists some $x \in \mathfrak{g}^{0}(\mathfrak{h}) \backslash \mathfrak{h}$ with $\operatorname{ad}(h) x \in \mathfrak{h}$ for all $h \in \mathfrak{h}$. But this contradicts $\mathfrak{n}_{\mathfrak{g}}(\mathfrak{h})=\mathfrak{h}$.\\
(C2) $\Rightarrow$ (C1): From $\mathfrak{h} \subseteq \mathfrak{g}^{0}(\mathfrak{h})$ we derive with Engel's Theorem that $\mathfrak{h}$ is nilpotent. To see that $\mathfrak{h}$ is self-normalizing, let $x \in \mathfrak{n}_{\mathfrak{g}}(\mathfrak{h})$. Then $\operatorname{ad}(h) x \in \mathfrak{h}$ for all $h \in \mathfrak{h}$, and therefore $\operatorname{ad}(h)^{n} x=0$ for sufficiently large $n \in \mathbb{N}$. Hence $x \in \mathfrak{g}^{0}(\mathfrak{h})=\mathfrak{h}$.

If (C1) and (C2) are satisfied and $\mathfrak{n} \supseteq \mathfrak{h}$ is a nilpotent subalgebra, then $\mathfrak{n} \subseteq \mathfrak{g}^{0}(\mathfrak{h})=\mathfrak{h}$. Therefore, $\mathfrak{h}$ is maximally nilpotent.

Definition 6.1.7. Let $\mathfrak{g}$ be a finite-dimensional Lie algebra. A nilpotent subalgebra $\mathfrak{h}$ is called a Cartan subalgebra of $\mathfrak{g}$ if it is self-normalizing, i.e., $\mathfrak{n}_{\mathfrak{g}}(\mathfrak{h})=\mathfrak{h}$.

Remark 6.1.8. (a) If $\mathfrak{g}$ is nilpotent, then $\mathfrak{g}$ is the only Cartan subalgebra of $\mathfrak{g}$ because any Cartan subalgebra is maximally nilpotent (Proposition 6.1.6).\\
(b) If $\mathfrak{g}_{j}$ are Lie algebras with Cartan subalgebras $\mathfrak{h}_{j}$, then $\bigoplus_{j} \mathfrak{h}_{j}$ is a Cartan subalgebra of $\bigoplus_{j} \mathfrak{g}_{j}$.

Example 6.1.9. (i) Let $\mathfrak{g}=\mathbb{R} h+\mathbb{R} p+\mathbb{R} q+\mathbb{R} z$ be the oscillator algebra (cf. Example 5.1.19). Then $\mathfrak{h}=\mathbb{R} h+\mathbb{R} z$ is a Cartan subalgebra of $\mathfrak{g}$.\\
(ii) In $\mathfrak{g}=\mathfrak{g} \mathfrak{l}_{n}(\mathbb{K})$, the Lie subalgebra $\mathfrak{h}$ of the diagonal matrices is a Cartan subalgebra (cf. Exercise 6.1.1).\\
(iii) Every one-dimensional subspace of $\mathfrak{g}=\mathfrak{s o}_{3}(\mathbb{R})$ is a Cartan subalgebra.\\
(iv) Every Cartan subalgebra is maximally nilpotent (Proposition 6.1.6), but not every maximally nilpotent subalgebra is a Cartan subalgebra. The subalgebra

$$
\mathfrak{n}:=\mathbb{R}\left(\begin{array}{ll}
0 & 1 \\
0 & 0
\end{array}\right)<\mathfrak{g}=\mathfrak{s} \mathfrak{l}_{2}(\mathbb{R})
$$

is maximally nilpotent and not self-normalizing. Its normalizer is the only proper subalgebra

$$
\mathfrak{b}:=\mathbb{R}\left(\begin{array}{ll}
0 & 1 \\
0 & 0
\end{array}\right)+\mathbb{R}\left(\begin{array}{cc}
1 & 0 \\
0 & -1
\end{array}\right)
$$

containing $\mathfrak{n}$. It is solvable but not nilpotent.\\
Definition 6.1.10. Let $\mathfrak{g}$ be a finite-dimensional Lie algebra and $\mathfrak{h}<\mathfrak{g}$ be a Cartan subalgebra. We call $\mathfrak{h}$ a splitting Cartan subalgebra if ad $x$ is split for each $x \in \mathfrak{h}$. By Lemma 6.1.3 and Proposition 6.1.6, we then have the root space decomposition

$$
\mathfrak{g}=\mathfrak{h} \oplus \bigoplus_{\lambda \neq 0} \mathfrak{g}^{\lambda}
$$

with respect to the Cartan subalgebra $\mathfrak{h}$.\\
If, in addition, each ad $x, x \in \mathfrak{h}$, diagonalizable, we call $\mathfrak{h}$ a toral Cartan subalgebra. In this case, $\mathfrak{g}^{\lambda}=\mathfrak{g}_{\lambda}$ for each $\lambda \in \mathfrak{h}^{*}$, and in particular $\mathfrak{h}=\mathfrak{g}_{0}$ implies that $\mathfrak{h}$ is abelian.

Proposition 6.1.11. (i) A subalgebra $\mathfrak{h}$ of the real Lie algebra $\mathfrak{g}$ is a Cartan subalgebra if and only if $\mathfrak{h}_{\mathbb{C}}$ is a Cartan subalgebra of $\mathfrak{g}_{\mathbb{C}}$.\\
(ii) Let $\varphi: \mathfrak{g} \rightarrow \widetilde{\mathfrak{g}}$ be a surjective homomorphism and $\mathfrak{h}<\mathfrak{g}$ a Cartan subalgebra. Then $\varphi(\mathfrak{h})$ is a Cartan subalgebra of $\widetilde{\mathfrak{g}}$.\\
(iii) If $\mathfrak{h}$ is a Cartan subalgebra of $\mathfrak{g}$ contained in the subalgebra $\mathfrak{k}$, then $\mathfrak{h}$ also is a Cartan subalgebra of $\mathfrak{k}$.

Proof. (i) This follows immediately from $\mathfrak{n}_{\mathfrak{g}}(\mathfrak{h})_{\mathbb{C}}=\mathfrak{n}_{\mathfrak{g}_{\mathbb{C}}}\left(\mathfrak{h}_{\mathbb{C}}\right)$ and Exercise 5.4.2.\\
(ii) By (i), we may assume that $\mathbb{K}=\mathbb{C}$. Proposition 5.2.3 shows that $\widetilde{\mathfrak{h}}:=\varphi(\mathfrak{h})$ is nilpotent. By Proposition 6.1.6, it suffices to show $\widetilde{\mathfrak{g}}^{0}(\widetilde{\mathfrak{h}})=\widetilde{\mathfrak{h}}$. Clearly, $\varphi\left(\mathfrak{g}^{0}(\mathfrak{h})\right)=\varphi(\mathfrak{h})=\widetilde{\mathfrak{h}}$ (Proposition 6.1.6). Moreover, Lemma 6.1.3 yields decompositions

$$
\mathfrak{g}=\mathfrak{g}^{0}(\mathfrak{h}) \oplus \bigoplus_{0 \neq \lambda \in \mathfrak{h}^{*}} \mathfrak{g}^{\lambda}(\mathfrak{h}) \quad \text { and } \quad \widetilde{\mathfrak{g}}=\widetilde{\mathfrak{g}}^{0}(\widetilde{\mathfrak{h}}) \oplus \bigoplus_{\widetilde{\lambda} \in \widetilde{\mathfrak{h}}^{*}} \widetilde{\mathfrak{g}}^{\widetilde{\lambda}}(\widetilde{\mathfrak{h}}) \text {. }
$$

If $x \in \mathfrak{g}^{\lambda}(\mathfrak{h}), \lambda \neq 0$ and $h \in \mathfrak{h}$, there is an $m \in \mathbb{N}$ with $(\operatorname{ad} h-\lambda(h))^{m} x=0$, and since $\varphi$ is a homomorphism, $(\operatorname{ad} \varphi(h)-\lambda(h))^{m} \varphi(x)=0$. If $\varphi(x) \neq 0$ and $\varphi(h)=0$, then $\lambda(h)=0$, so that we can define $\widetilde{\lambda} \in \widetilde{\mathfrak{h}}^{*}$ via $\widetilde{\lambda}(\varphi(h)):=\lambda(h)$. We now have

$$
\varphi\left(\mathfrak{g}^{\lambda}(\mathfrak{h})\right) \subseteq \begin{cases}\{0\} & \text { if } \varphi\left(\mathfrak{g}^{\lambda}(\mathfrak{h})\right)=\{0\}, \\ \widetilde{\mathfrak{g}}^{\widetilde{\lambda}}(\widetilde{\mathfrak{h}}) & \text { if } \varphi\left(\mathfrak{g}^{\lambda}(\mathfrak{h})\right) \neq\{0\} .\end{cases}
$$

In particular, for each root $\lambda$, the image of $\mathfrak{g}^{\lambda}(\mathfrak{h})$ is contained in $\sum_{0 \neq \tilde{\lambda}} \widetilde{\mathfrak{g}}^{\widetilde{\lambda}}(\widetilde{\mathfrak{h}})$. Since $\varphi$ is surjective, it follows that $\widetilde{\mathfrak{h}}=\varphi(\mathfrak{h})=\widetilde{\mathfrak{g}}^{0}(\widetilde{\mathfrak{h}})$.\\
(iii) The subalgebra $\mathfrak{h}$ of $\mathfrak{k}$ is nilpotent and self-normalizing, hence a Cartan subalgebra.

Proposition 6.1.12. Let $\mathfrak{a} \subseteq \mathfrak{g}$ be an abelian subalgebra for which all operators $\operatorname{ad} x, x \in \mathfrak{a}$, are semisimple. Then the Cartan subalgebras of the centralizer $\mathfrak{z}_{\mathfrak{g}}(\mathfrak{a})$ are precisely the Cartan subalgebras of $\mathfrak{g}$ containing $\mathfrak{a}$. In particular, such Cartan subalgebras exist.

Proof. Let $\mathfrak{c}:=\mathfrak{z}_{\mathfrak{g}}(\mathfrak{a})$ and fix a Cartan subalgebra $\mathfrak{h}$ of $\mathfrak{c}$. Since $\mathfrak{a}$ is central in $\mathfrak{c}$, we find $\mathfrak{a} \subseteq \mathfrak{z}(\mathfrak{c}) \subseteq \mathfrak{n}_{\mathfrak{c}}(\mathfrak{h})=\mathfrak{h}$. For the normalizer $\mathfrak{n}:=\mathfrak{n}_{\mathfrak{g}}(\mathfrak{h})$, we have

$$
[\mathfrak{a}, \mathfrak{n}] \subseteq[\mathfrak{h}, \mathfrak{n}] \subseteq \mathfrak{h}
$$

Since each $\operatorname{ad}_{\mathfrak{g}} x, x \in \mathfrak{a}$, is semisimple, Lemma 5.5.20 shows that $\mathfrak{g}$ is a semisimple $\mathfrak{a}$-module, so that there exists an $\mathfrak{a}$-invariant subspace $\mathfrak{d} \subseteq \mathfrak{n}$ with $\mathfrak{n}=\mathfrak{h} \oplus \mathfrak{d}$. Then $[\mathfrak{a}, \mathfrak{d}] \subseteq \mathfrak{d} \cap \mathfrak{h}=\{0\}$ implies $\mathfrak{d} \subseteq \mathfrak{c}$. Therefore, $\mathfrak{n}$ is contained in $\mathfrak{c}$, and since $\mathfrak{h}$ is a Cartan subalgebra of $\mathfrak{c}$, we get $\mathfrak{n}=\mathfrak{h}$, showing that $\mathfrak{h}$ is a Cartan subalgebra of $\mathfrak{g}$.

If, conversely, $\mathfrak{h}$ is a Cartan subalgebra of $\mathfrak{g}$ containing $\mathfrak{a}$, then $\mathfrak{h} \subseteq \mathfrak{g}^{0}(\mathfrak{h}) \subseteq \mathfrak{g}^{0}(\mathfrak{a})=\mathfrak{g}_{0}(\mathfrak{a})=\mathfrak{c}$. Since $\mathfrak{h}$ is self-normalizing in $\mathfrak{g}$, it is also self-normalizing in $\mathfrak{c}$, hence a Cartan subalgebra of $\mathfrak{c}$.

Lemma 6.1.13. For a nilpotent Lie algebra $\mathfrak{h} \subseteq \mathfrak{g} \mathfrak{l}(V)$, the following assertions hold:\\
(i) The set $\mathfrak{h}_{s}:=\left\{x_{s}: x \in \mathfrak{h}\right\}$, consisting of all semisimple Jordan components of elements in $\mathfrak{h}$, is an abelian Lie algebra commuting with $\mathfrak{h}$.\\
(ii) The set $\mathfrak{h}_{n}:=\left\{x_{n}: x \in \mathfrak{h}\right\}$ of all nilpotent Jordan components of elements in $\mathfrak{h}$ is a nilpotent Lie subalgebra of $\mathfrak{g} \mathfrak{l}(V)$.\\
(iii) $\widetilde{\mathfrak{h}}:=\mathfrak{h}_{n}+\mathfrak{h}_{s}$ is a direct sum of Lie algebras.

Proof. (i) For each $x \in \mathfrak{h}$, we have $\operatorname{ad} x(\mathfrak{h}) \subseteq \mathfrak{h}$, so that we get $(\operatorname{ad} x)_{s}(\mathfrak{h}) \subseteq \mathfrak{h}$. Since $\operatorname{ad} x_{\mathfrak{h}}$ is nilpotent, and Proposition 5.3.7(ii)) and Corollary 5.3.9 show that

$$
0=\left(\left.\operatorname{ad} x\right|_{\mathfrak{h}}\right)_{s}=\left.(\operatorname{ad} x)_{s}\right|_{\mathfrak{h}}=\left.\operatorname{ad}\left(x_{s}\right)\right|_{\mathfrak{h}}
$$

it follows that $\left[x_{s}, \mathfrak{h}\right]=\{0\}$. In view of Theorem 5.3.3(iv), $\left[x_{s}, y_{s}\right]=0$ for $x, y \in \mathfrak{h}$, so that $\mathfrak{h}_{s}$ is a commutative subset of $\mathfrak{g} \mathfrak{l}(V)$.

To see that $\mathfrak{h}_{s}$ actually is a linear subspace, we may w.l.o.g. assume that $\mathbb{K}=\mathbb{C}$ (otherwise we consider the complexification $V_{\mathbb{C}}$ ). We recall from Lemma 6.1.3 the generalized weight space decomposition

$$
V=\bigoplus_{\lambda \in \mathcal{P}(V)} V^{\lambda}(\mathfrak{h}), \quad \text { where } \quad \mathcal{P}(V)=\left\{\lambda \in \mathfrak{h}^{*}: V^{\lambda}(\mathfrak{h}) \neq\{0\}\right\}
$$

With respect to this decomposition, $x_{s}$ corresponds to the diagonal operator acting on $V^{\lambda}(\mathfrak{h})$ by $\lambda(x)$. We may therefore identify $\mathfrak{h}_{s}$ with the image of $\mathfrak{h}$ under the linear map

$$
\mathfrak{h} \rightarrow \mathbb{K}^{\mathcal{P}(V)}, \quad h \mapsto(\lambda(h))_{\lambda},
$$

which is a linear subspace. In particular, we see that $(x+y)_{s}=x_{s}+y_{s}$ for $x, y \in \mathfrak{h}$.\\
(ii) As before, we may assume that $\mathbb{K}=\mathbb{C}$. In view of (i), $\widetilde{\mathfrak{h}}:=\mathfrak{h}+\mathfrak{h}_{s}$ is a Lie subalgebra of $\mathfrak{g l}(V)$ and all weights $\lambda: \mathfrak{h} \rightarrow \mathbb{C}$ extend to linear functionals $\widetilde{\lambda}: \widetilde{\mathfrak{h}} \rightarrow \mathbb{C}$, so that we arrive at the same generalized weight spaces $V^{\lambda}(\mathfrak{h})= V^{\widetilde{\lambda}}(\widetilde{\mathfrak{h}})$. In view of $x_{n}=x-x_{s}, \widetilde{\mathfrak{h}}$ also contains $\mathfrak{h}_{n}$. Moreover, each element of $\widetilde{\mathfrak{h}}$ can be written as a sum $x+y_{s}$ for $x, y \in \mathfrak{h}$. Since ad $x\left(y_{s}\right)=0$, the same holds for $\operatorname{ad} x_{n}\left(y_{s}\right)=(\operatorname{ad} x)_{n}\left(y_{s}\right)$, so that $x+y_{s}=x_{n}+\left(x_{s}+y_{s}\right)$ is the Jordan decomposition of $x+y_{s}$. We further know from (i) that $(x+y)_{s}=x_{s}+y_{s}$, so that $\widetilde{\mathfrak{h}}=\mathfrak{h}_{n}+\mathfrak{h}_{s}$.

An element of $\mathfrak{h}$ is of the form $x_{n}$ for some $x \in \mathfrak{h}$ if and only if its semisimple component vanishes, i.e., $x \in \bigcap \operatorname{ker} \widetilde{\lambda}$. Since the functionals $\widetilde{\lambda}: \widetilde{h} \rightarrow \mathbb{C}$ vanish on the commutator algebras (Corollary 5.4.11), the intersection of their kernels is an ideal in $\widetilde{\mathfrak{h}}$. Since $\mathfrak{h}_{s}$ is central in $\widetilde{\mathfrak{h}}$ and $\mathfrak{h}_{n} \cap \mathfrak{h}_{s}=\{0\}, \widetilde{\mathfrak{h}}=\mathfrak{h}_{n} \oplus \mathfrak{h}_{s}$ is a direct sum of Lie algebras.

\subsubsection*{Cartan Subalgebras and Regular Elements}
Next we want to ensure that Cartan subalgebras actually exist. To this end, for every $x \in \mathfrak{g}$, we consider the generalized eigenspace $\mathfrak{g}^{0}(\operatorname{ad} x)$. This space is always different from zero since it obviously contains $x$.

Definition 6.1.14. (a) The number

$$
\operatorname{rank}(\mathfrak{g}):=\min \left\{\operatorname{dim} \mathfrak{g}^{0}(\operatorname{ad} x) \mid x \in \mathfrak{g}\right\}
$$

is called the rank of $\mathfrak{g}$. An element $x \in \mathfrak{g}$ is called regular if $\operatorname{dim}_{\mathfrak{g}}{ }^{0}(\operatorname{ad} x)= \operatorname{rank}(\mathfrak{g})$. We write $\operatorname{reg}(\mathfrak{g})$ for the set of regular elements in $\mathfrak{g}$.\\
(b) Since $\operatorname{dim} \mathfrak{g}^{0}(\operatorname{ad} x)$ is the multiplicity of 0 as a root of the characteristic polynomial

$$
\operatorname{det}\left(\operatorname{ad} x-t \operatorname{id}_{\mathfrak{g}}\right)=\sum_{k=0}^{n} p_{k}(x) t^{k}
$$

of ad $x$, we have

$$
\operatorname{rank}(\mathfrak{g})=\min \left\{k \in \mathbb{N} \mid p_{k} \not \equiv 0\right\} .
$$

As the determinant is a polynomial function on $\operatorname{End}(\mathfrak{g})$, all the functions $p_{k}: \mathfrak{g} \rightarrow \mathbb{K}$ in (6.3) are polynomials (cf. Exercise 5.3.8).

Lemma 6.1.15. The set $\operatorname{reg}(\mathfrak{g})$ of regular elements has the following properties:\\
(i) $\mathfrak{g} \backslash \operatorname{reg}(\mathfrak{g})$ is the zero-set of a nonconstant polynomial.\\
(ii) $\operatorname{reg}(\mathfrak{g})$ is open and dense. For $\mathbb{K}=\mathbb{C}$, it is connected.\\
(iii) $\operatorname{reg}(\mathfrak{g})$ is invariant under the automorphism group $\operatorname{Aut}(\mathfrak{g})$ of $\mathfrak{g}$.\\
(iv) If $\mathbb{K}=\mathbb{R}$, then $\operatorname{reg}(\mathfrak{g})=\mathfrak{g} \cap \operatorname{reg}\left(\mathfrak{g}_{\mathbb{C}}\right)$ and $\operatorname{rank}_{\mathbb{R}}(\mathfrak{g})=\operatorname{rank}_{\mathbb{C}}\left(\mathfrak{g}_{\mathbb{C}}\right)$.

Proof. (i) Let $r:=\operatorname{rank}(\mathfrak{g})$. Then (i) follows from the fact that the functions $p_{k}: \mathfrak{g} \rightarrow \mathbb{K}$ occurring in (6.3) are polynomials and

$$
\operatorname{reg}(\mathfrak{g})=\left\{x \in \mathfrak{g}: p_{r}(x) \neq 0\right\}
$$

(ii) Clearly, $\operatorname{reg}(\mathfrak{g})=p_{r}^{-1}\left(\mathbb{K}^{\times}\right)$is open because $p_{r}$ is continuous. Further, the zero set of the nonconstant polynomial $p_{r}$ contains no open subset (Exercise 6.1.2), so that $\operatorname{reg}(\mathfrak{g})$ is also dense in $\mathfrak{g}$. If $\mathfrak{g}$ is a complex Lie algebra, then Exercise 6.1.3 further implies that $\operatorname{reg}(\mathfrak{g})$ is connected.\\
(iii) For $\gamma \in \operatorname{Aut}(\mathfrak{g})$ we have $\operatorname{ad} \gamma(x)=\gamma \circ \operatorname{ad} x \circ \gamma^{-1}$. Therefore, $\operatorname{ad} \gamma(x)$ and ad $x$ have the same characteristic polynomial. We conclude that all polynomials $p_{k}$ are invariant under $\operatorname{Aut}(\mathfrak{g})$. In particular, $\operatorname{reg}(\mathfrak{g})=p_{r}^{-1}\left(\mathbb{K}^{\times}\right)$is invariant.\\
(iv) Since $\operatorname{det}\left(M_{\mathbb{C}}\right)=\operatorname{det}(M)$ for each $M \in \operatorname{End}(\mathfrak{g})$, the polynomials $p_{k}^{\mathbb{C}}$ on the complexified Lie algebra $\mathfrak{g}_{\mathbb{C}}$ satisfy $\left.p_{k}^{\mathbb{C}}\right|_{\mathfrak{g}}=p_{k}$ for each $k$. Hence Exercise 6.1.4 shows that $p_{k}$ vanishes if and only $p_{k}^{\mathbb{C}}$ does. In particular, $\operatorname{rank}_{\mathbb{R}}(\mathfrak{g})=\operatorname{rank}_{\mathbb{C}}\left(\mathfrak{g}_{\mathbb{C}}\right)$, and (iv) follows.

For the following lemma, we recall from Example 4.2.5 that for each inner derivation ad $x, e^{\operatorname{ad} x}$ is an automorphism of $\mathfrak{g}$. The elements of the group

$$
\operatorname{Inn}(\mathfrak{g}):=\left\langle e^{\operatorname{ad} x}: x \in \mathfrak{g}\right\rangle
$$

generated by these automorphisms are called inner automorphisms.\\
Lemma 6.1.16. Let $\mathfrak{h} \subseteq \mathfrak{g}$ be a splitting Cartan subalgebra and $\Delta:=\Delta(\mathfrak{g}, \mathfrak{h})$. Then the following assertions hold:\\
(i) $\operatorname{reg}(\mathfrak{g}) \cap \mathfrak{h}=\mathfrak{h} \backslash \bigcup_{\alpha \in \Delta} \operatorname{ker} \alpha$.\\
(ii) $\operatorname{rank}(\mathfrak{g})=\operatorname{dim} \mathfrak{h}$.\\
(iii) $\operatorname{Inn}(\mathfrak{g})(\mathfrak{h} \cap \operatorname{reg}(\mathfrak{g}))$ is an open subset of $\mathfrak{g}$.

Proof. (i), (ii) For each $h \in \mathfrak{h}$, we have $\mathfrak{g}^{0}(\operatorname{ad} h)=\mathfrak{h}+\sum_{\alpha(h)=0} \mathfrak{g}^{\alpha}$, and this space is minimal if and only if no root vanishes on $h$. Since $\Delta$ is finite, such elements exist, and for any such element $\mathfrak{g}^{0}(\operatorname{ad} h)=\mathfrak{h}$.

It remains to show that $r:=\operatorname{rank}(\mathfrak{g})$ is not strictly smaller than $\operatorname{dim} \mathfrak{h}$. To this end, we consider the map

$$
\Phi: \mathfrak{g} \times \mathfrak{h} \rightarrow \mathfrak{g}, \quad(a, b) \mapsto e^{\operatorname{ad} a} b
$$

Then

$$
\mathrm{d} \Phi(0, b)(v, w)=[v, b]+w
$$

for $v \in \mathfrak{g}, w \in \mathfrak{h}$. We therefore have

$$
\operatorname{Im}(\mathrm{d} \Phi(0, b))=[b, \mathfrak{g}]+\mathfrak{h} \supseteq \mathfrak{h}+\sum_{\alpha(b) \neq 0} \mathfrak{g}^{\alpha}
$$

because $\alpha(b) \neq 0$ implies $\mathfrak{g}^{\alpha} \subseteq[b, \mathfrak{g}]$ since $\left.\operatorname{ad} b\right|_{\mathfrak{g}^{\alpha}}$ is invertible. If no root vanishes on $b$, then $\mathrm{d} \Phi(0, b)$ is surjective, and the Implicit Function Theorem implies that the image of $\Phi$ is a neighborhood of $b=\Phi(0, b)$. Since $\operatorname{reg}(\mathfrak{g})$ is a dense subset of $\mathfrak{g}$ (Lemma 6.1.15), the image of $\Phi$ contains a regular element $x$ which we write as $x=\Phi(a, h)=e^{\operatorname{ad} a} h$ for some $h \in \mathfrak{h}$ and $a \in \mathfrak{g}$. Then $h=e^{-\operatorname{ad} a} x$ is also regular (Lemma 6.1.15) and contained in $\mathfrak{h}$. Thus $\mathfrak{h}$ contains regular elements and the discussion above yields $r=\operatorname{dim} \mathfrak{h}$ and (i).\\
(iii) The argument above also shows that $\mathrm{d} \Phi(0, h)$ is surjective for each regular element $h \in \mathfrak{h}$. Since $\operatorname{reg}(\mathfrak{g}) \cap \mathfrak{h}$ is open, we see with the Implicit Function Theorem that there exist neighborhoods $U_{h}$ of $h$ in $\mathfrak{h} \cap \operatorname{reg}(\mathfrak{g})$ and $V$ of 0 in $\mathfrak{g}$ such that $\Phi\left(V \times U_{h}\right)=e^{\text {ad } V} U_{h}$ is an open subset of $\mathfrak{g}$. Since $e^{\text {ad } V} U_{h}$ consists of regular elements, $h$ is an interior point of

$$
\operatorname{Inn}(\mathfrak{g})(\operatorname{reg}(\mathfrak{g}) \cap \mathfrak{h})
$$

As $\operatorname{Inn}(\mathfrak{g})$ acts by homeomorphisms on $\mathfrak{g}$, this shows that each point in $\operatorname{Inn}(\mathfrak{g})(\operatorname{reg}(\mathfrak{g}) \cap \mathfrak{h})$ is inner, hence that $\operatorname{Inn}(\mathfrak{g})(\operatorname{reg}(\mathfrak{g}) \cap \mathfrak{h})$ is open in $\mathfrak{g}$.

Lemma 6.1.17. Let $\mathfrak{g}$ be a finite-dimensional Lie algebra and $\mathfrak{h}<\mathfrak{g}$. If $x \in \mathfrak{h}$ is regular in $\mathfrak{g}$, then $x$ is also regular in $\mathfrak{h}$.

Proof. For $h \in \mathfrak{h}$ and $V:=\mathfrak{g} / \mathfrak{h}$, consider the linear maps

$$
A(h): V \rightarrow V, \quad y+\mathfrak{h} \mapsto[h, y]+\mathfrak{h},
$$

and $B(h)=\left.\operatorname{ad}(h)\right|_{\mathfrak{h}}: \mathfrak{h} \rightarrow \mathfrak{h}$. We set

$$
\operatorname{det}(A(h)-t \mathrm{id})=\sum_{k=0}^{m} a_{k}(h) t^{k} \quad \text { and } \quad \operatorname{det}(B(h)-t \mathrm{id})=\sum_{j=0}^{n} b_{j}(h) t^{j}
$$

We further put

$$
d_{A}(h):=\operatorname{dim} V^{0}(A(h)), \quad d_{B}(h):=\operatorname{dim} \mathfrak{h}^{0}(B(h)),
$$

and

$$
r_{A}:=\min _{h \in \mathfrak{h}} d_{A}(h), \quad r_{B}:=\min _{h \in \mathfrak{h}} d_{B}(h) .
$$

Then $r_{A}=d_{A}(h)$ if and only if $a_{r_{A}}(h) \neq 0$, and similarly $r_{B}=d_{B}(h)$ if and only if $b_{r_{B}}(h) \neq 0$. We consider the set

$$
S:=\left\{h \in \mathfrak{h} \mid a_{r_{A}}(h) \cdot b_{r_{B}}(h) \neq 0\right\}=\left\{h \in \mathfrak{h} \mid r_{A}=d_{A}(h), r_{B}=d_{B}(h)\right\}
$$

which is nonempty because $b_{r_{B}}$ does not vanish on the open complement of the zero-set of $a_{r_{A}}$ (Exercise 6.1.2).

Every element of $S$ is clearly regular in $\mathfrak{h}$, so that it suffices to show that $x \in S$. If we identify $V=\mathfrak{g} / \mathfrak{h}$ with a vector space complement of $\mathfrak{h}$ in $\mathfrak{g}$, we can write ad $h$ in block form as

$$
\operatorname{ad} h=\left(\begin{array}{cc}
B(h) & * \\
0 & A(h)
\end{array}\right)
$$

(cf. Example 5.1.7). This leads for $h \in \mathfrak{h}$ to the factorization

$$
\operatorname{det}(\operatorname{ad} h-t \operatorname{id})=\operatorname{det}(A(h)-t \operatorname{id}) \operatorname{det}(B(h)-t \operatorname{id}),
$$

and hence to $r:=\operatorname{rank}(\mathfrak{g}) \leq r_{A}+r_{B}$ since $S \neq \emptyset$. On the other hand, the fact that $\mathfrak{h}$ contains the regular element $x$ leads to $r=r_{A}+r_{B}$, so that $a_{r_{A}}(h) b_{r_{B}}(h)=p_{r}(h)$. We conclude that $x \in S \subseteq \operatorname{reg}(\mathfrak{h})$.

The following theorem clarifies the connection between Cartan subalgebras and regular elements.

Theorem 6.1.18. Let $\mathfrak{g}$ be a finite-dimensional Lie algebra.\\
(i) For any regular element $x \in \mathfrak{g}, \mathfrak{g}^{0}(\operatorname{ad} x)$ is a Cartan subalgebra of $\mathfrak{g}$.\\
(ii) Every Cartan subalgebra $\mathfrak{h}$ contains regular elements and if $x \in \mathfrak{h}$ is regular, then $\mathfrak{h}=\mathfrak{g}^{0}(\operatorname{ad} x)$.\\
(iii) All Cartan subalgebras have the same dimension rank $(\mathfrak{g})$.\\
(iv) For any Cartan subalgebra $\mathfrak{h}$, the set $\operatorname{Inn}(\mathfrak{g})(\mathfrak{h} \cap \operatorname{reg}(\mathfrak{g}))$ is open in $\mathfrak{g}$.

Proof. (i) Clearly, $x \in \mathfrak{h}:=\mathfrak{g}^{0}(\operatorname{ad} x)$, so by Lemma 6.1.17, $x$ is also regular in $\mathfrak{h}$. Since $\mathfrak{h}=\mathfrak{h}^{0}(\operatorname{ad} x)$, we have

$$
\operatorname{rank}(\mathfrak{h})=\operatorname{dim}(\mathfrak{h})=\operatorname{rank}(\mathfrak{g}) .
$$

Hence $\mathfrak{h}^{0}(\operatorname{ad} y)=\mathfrak{h}$ for all $y \in \mathfrak{h}$, i.e., each $\left.\operatorname{ad}(y)\right|_{\mathfrak{h}}$ is nilpotent. Therefore, $\mathfrak{h}$ is nilpotent by Corollary 5.2.8. We thus arrive at

$$
\mathfrak{h} \subseteq \mathfrak{g}^{0}(\operatorname{ad} \mathfrak{h}) \subseteq \mathfrak{g}^{0}(\operatorname{ad} x)=\mathfrak{h},
$$

and Proposition 6.1.6 shows that $\mathfrak{h}$ is a Cartan subalgebra.\\
(ii) Let $\mathfrak{h} \subseteq \mathfrak{g}$ be a Cartan subalgebra. If $x \in \mathfrak{h} \cap \operatorname{reg}(\mathfrak{g})$, then (i) implies that $\mathfrak{g}^{0}(\operatorname{ad} x)$ is a Cartan subalgebra containing $\mathfrak{h}$, so that the fact that Cartan subalgebras are maximally nilpotent yields $\mathfrak{h}=\mathfrak{g}^{0}(\operatorname{ad} x)$.

If $\mathbb{K}=\mathbb{C}$, then $\mathfrak{h}$ is splitting and Lemma 6.1.16 implies that $\mathfrak{h}$ contains regular elements. For $\mathbb{K}=\mathbb{R}$, we note that $\mathfrak{h}_{\mathbb{C}} \subseteq \mathfrak{g}_{\mathbb{C}}$ is a Cartan subalgebra (Proposition 6.1.11), hence contains regular elements. This means that for $r:=\operatorname{rank}(\mathfrak{g})$, the polynomial $p_{r}: \mathfrak{g}_{\mathbb{C}} \rightarrow \mathbb{C}$ does not vanish on $\mathfrak{h}_{\mathbb{C}}$. Hence it does not vanish on $\mathfrak{h}$ (Exercise 6.1.4), and this means that $\mathfrak{h}$ contains regular elements of $\mathfrak{g}$.\\
(iii) follows immediately from (ii) and $\operatorname{dim} \mathfrak{g}^{0}(\operatorname{ad} x)=\operatorname{rank}(\mathfrak{g})$ for regular elements.\\
(iv) We know already that this is the case for $\mathbb{K}=\mathbb{C}$. In the real case, we also consider the map

$$
\Phi: \mathfrak{g} \times \mathfrak{h} \rightarrow \mathfrak{g}, \quad(a, b) \mapsto e^{\operatorname{ad} a} b
$$

with

$$
\mathrm{d} \Phi(0, x)(v, w)=[v, x]+w .
$$

If $x \in \mathfrak{h}$ is regular, then $\mathfrak{h}=\mathfrak{g}^{0}(\operatorname{ad} x)$, so that the induced map $\operatorname{ad}_{\mathfrak{g} / \mathfrak{h}}(y+\mathfrak{h}):= [x, y]+\mathfrak{h}$ is invertible because its kernel is trivial. This implies that $\mathfrak{g} \subseteq \mathfrak{h}+[x, \mathfrak{g}]$, which means that $\mathrm{d} \Phi(0, x)$ is surjective, so that the Implicit Function Theorem implies that $\Phi(\mathfrak{g} \times(\operatorname{reg}(\mathfrak{g}) \cap \mathfrak{h}))$ is a neighborhood of $x=\Phi(0, x)$. Now the same argument as in the complex case proves (iv) (Lemma 6.1.16).

Let $\mathfrak{g}$ be a finite-dimensional Lie algebra. Define an equivalence relation $R$ on the set $\operatorname{reg}(\mathfrak{g})$ of regular elements via

$$
x \sim y \quad: \Longleftrightarrow \quad \exists g \in \operatorname{Inn}(\mathfrak{g}) \text { with } g\left(\mathfrak{g}^{0}(\operatorname{ad} x)\right)=\mathfrak{g}^{0}(\operatorname{ad} y) .
$$

Lemma 6.1.19. The equivalence classes of $\sim$ are open subsets of $\operatorname{reg}(\mathfrak{g})$.\\
Proof. Fix $x \in \operatorname{reg}(\mathfrak{g})$ and set $\mathfrak{h}:=\mathfrak{g}^{0}(\operatorname{ad} x)$. In view of Theorem 6.1.18, the set $\operatorname{Inn}(\mathfrak{g})(\mathfrak{h} \cap \operatorname{reg}(\mathfrak{g}))$ is an open neighborhood of $x$. Each element in this set is of the form $g(y)$ for $y \in \operatorname{reg}(\mathfrak{g}) \cap \mathfrak{h}$, so that

$$
\mathfrak{g}^{0}(\operatorname{ad}(g(y)))=g\left(\mathfrak{g}^{0}(\operatorname{ad} y)\right)=g(\mathfrak{h})=g\left(\mathfrak{g}^{0}(\operatorname{ad} x)\right)
$$

implies that $g(y) \sim x$. This proves that all equivalence classes of $\sim$ are open.

Theorem 6.1.20. If $\mathfrak{g}$ is a finite-dimensional complex Lie algebra, then the group $\operatorname{Inn}(\mathfrak{g})$ acts transitively on the set of Cartan subalgebras of $\mathfrak{g}$.

Proof. According to Lemma 6.1.15, $\operatorname{reg}(\mathfrak{g})$ is connected. On the other hand, it is the disjoint union of the open equivalence classes of the relation $\sim$ (cf. Lemma 6.1.19). Hence only one such class exists. Since every Cartan subalgebra of $\mathfrak{g}$ is of the form $\mathfrak{g}^{0}(\operatorname{ad} x)$ by Theorem 6.1.18, the assertion follows.

\subsubsection*{Exercises for Section 6.1}
Exercise 6.1.1. The diagonal matrices form a Cartan subalgebra of $\mathfrak{g l}{ }_{n}(\mathbb{K})$.\\
Exercise 6.1.2. Let $V$ be a finite-dimensional vector space and $p: V \rightarrow \mathbb{K}$ a polynomial function vanishing on an open subset $U \subseteq V$. Show that $p=0$.

Exercise 6.1.3. Let $V$ be a complex finite-dimensional vector space and $p: V \rightarrow \mathbb{K}$ a nonzero polynomial function. Then $p^{-1}\left(\mathbb{C}^{\times}\right)$is connected.

Exercise 6.1.4. Let $V$ be a finite-dimensional real vector space and $V_{\mathbb{C}}$ its complexification. Show that a polynomial function $p: V_{\mathbb{C}} \rightarrow \mathbb{C}$ vanishes if and only if $\left.p\right|_{V}$ vanishes.

\subsection*{The Classification of Simple $\mathfrak{s l}_{2}(\mathbb{K})$-Modules}
As we shall see in Section 6.3 below, the Lie algebra $\mathfrak{s} \mathfrak{l}_{2}(\mathbb{K})$ is of particular importance because semisimple Lie algebras with splitting Cartan subalgebras contain many subalgebras isomorphic to $\mathfrak{s} \mathfrak{l}_{2}(\mathbb{K})$ and the collection of these subalgebras essentially determines the structure of the whole Lie algebra. Therefore, the representation theory of $\mathfrak{s} \mathfrak{l}_{2}(\mathbb{K})$ plays a key role in the structure theory of these Lie algebras.

In the following, we shall use the following basis for $\mathfrak{s l}_{2}(\mathbb{K})$ :

$$
h:=\left(\begin{array}{cc}
1 & 0 \\
0 & -1
\end{array}\right), \quad e:=\left(\begin{array}{ll}
0 & 1 \\
0 & 0
\end{array}\right), \quad f:=\left(\begin{array}{ll}
0 & 0 \\
1 & 0
\end{array}\right)
$$

satisfying

$$
[h, e]=2 e, \quad[h, f]=-2 f, \quad \text { and } \quad[e, f]=h .
$$

We start with a discussion of a concrete family of representations of $\mathfrak{s} \mathfrak{l}_{2}(\mathbb{K})$. It will turn out later that the study of this family already provides all irreducible finite-dimensional representations of $\mathfrak{s} \mathfrak{l}_{2}(\mathbb{K})$.

Example 6.2.1. Let $V:=\mathbb{K}\left[Z, Z^{-1}\right]$ be the algebra of Laurent polynomials in $Z$. For any $f \in V$, the operator $D:=f \frac{d}{d Z}$ is a derivation of $V$ (Product Rule) and any derivation of the algebra $V$ is of this kind (Exercise 6.2.2). The Lie bracket on $\operatorname{der}(V)$ satisfies

$$
\left[f \frac{d}{d Z}, g \frac{d}{d Z}\right]=\left(f g^{\prime}-f^{\prime} g\right) \frac{d}{d Z}
$$

For $\lambda \in \mathbb{K}$, we consider the operators

$$
E:=\frac{d}{d Z}, \quad F:=-Z^{2} \frac{d}{d Z}+\lambda Z \mathbf{1}, \quad H:=-2 Z \frac{d}{d Z}+\lambda \mathbf{1}
$$

With (6.6) we obtain

$$
\left[Z^{n} \frac{d}{d Z}, Z^{m} \frac{d}{d Z}\right]=(m-n) Z^{n+m-1} \frac{d}{d Z}
$$

and hence the commutator relations

$$
\begin{array}{llll}
{[H, E]=} & -2\left[Z \frac{d}{d Z}, \frac{d}{d Z}\right] & =2 \frac{d}{d Z} & =2 E \\
{[H, F]} & =-2\left[Z \frac{d}{d Z},-Z^{2} \frac{d}{d Z}+\lambda Z \mathbf{1}\right] & =2 Z^{2} \frac{d}{d Z}-2 \lambda Z \mathbf{1} & =-2 F \\
{[E, F]=} & {\left[\frac{d}{d Z},-Z^{2} \frac{d}{d Z}+\lambda Z \mathbf{1}\right]} & =-2 Z \frac{d}{d Z}+\lambda \mathbf{1} & =H
\end{array}
$$

Since these are precisely the commutator relations of $\mathfrak{g}:=\mathfrak{s l}_{2}(\mathbb{K})$, we obtain by

$$
e \mapsto E, \quad f \mapsto F, \quad h \mapsto H
$$

a representation $\rho_{\lambda}: \mathfrak{s} \mathfrak{l}_{2}(\mathbb{K}) \rightarrow \operatorname{End}(V)$, resp., an $\mathfrak{s} \mathfrak{l}_{2}(\mathbb{K})$-module structure on $V$.

To understand the structure of this module, we consider the action of the operators $H, E$, and $F$ on the canonical basis:

$$
H \cdot Z^{n}=(\lambda-2 n) Z^{n}, \quad E \cdot Z^{n}=n Z^{n-1}, \quad F \cdot Z^{n}=(\lambda-n) Z^{n+1}
$$

In particular, we see that $H$ is diagonalizable with one-dimensional eigenspaces. With this information, it is easy to determine all submodules. Any submodule is adapted to the eigenspace decomposition of $H$ (Exercise 2.1.1(b)). Hence each submodule is of the form

$$
V_{J}:=\operatorname{span}\left\{Z^{n}: n \in J\right\}
$$

for a subset $J \subseteq \mathbb{Z}$. From the formulas above, we see that $V_{J}$ is a submodule if and only if $J$ satisfies the following conditions:\\
(i) If $n \in J$ and $n \neq 0$, then $n-1 \in J$.\\
(ii) If $n \in J$ and $\lambda \neq n$, then $n+1 \in J$.

If $\lambda \notin \mathbb{Z}$, then $\mathbb{K}[Z]=V_{\mathbb{N}_{0}}$ is the only nontrivial submodule of $V$. If $\lambda \in \mathbb{Z}$, then there are two possibilities. For $\lambda<0$, the only proper subsets of $\mathbb{Z}$ satisfying (i) and (ii) are

$$
\begin{gathered}
\{\ldots, \lambda-1, \lambda\}, \quad \mathbb{N}_{0}, \quad \text { and } \quad\{\ldots, \lambda-1, \lambda\} \cup \mathbb{N}_{0} . \\
\circ \quad \circ \quad] \quad \circ \quad \circ \quad \cdots \quad \circ \quad[\circ \quad \circ \quad \cdots .
\end{gathered}
$$

For $\lambda \geq 0$, the subsets of $\mathbb{Z}$ defining submodules are

$$
\begin{aligned}
\mathbb{N}_{0}, & \{\ldots, \lambda-1, \lambda\}, \quad\{0,1, \ldots, \lambda-1, \lambda\} . \\
\circ & \circ[\quad \circ \quad \circ \quad \cdots \quad \circ \quad] \quad \circ \quad \circ \quad \cdots .
\end{aligned}
$$

In this case, we obtain in particular a finite-dimensional submodule

$$
L(\lambda):=\operatorname{span}\left\{1, Z, \ldots, Z^{\lambda}\right\}
$$

Since $L(\lambda)$ contains no nontrivial proper submodule, it is simple.\\
We have seen in the preceding example that for each $\lambda \in \mathbb{N}_{0}$, there exists a simple $\mathfrak{s} \mathfrak{l}_{2}(\mathbb{K})$-module of dimension $\lambda+1$. Our next goal is to show that all simple finite-dimensional modules are isomorphic to some $L(\lambda)$.

The following lemma specializes for $\mu=1$ to an assertion on the Lie algebra $\mathfrak{s l}_{2}(\mathbb{K})$.

Lemma 6.2.2. Let ( $e, h, f$ ) be a triple of elements of an associative algebra $A$, satisfying the commutator relations

$$
[h, e]=2 e, \quad[h, f]=-2 f, \quad \text { and } \quad[e, f]=h .
$$

Then the following assertions hold:\\
(i) $\left[h, e^{n}\right]=2 n e^{n}$ and $\left[h, f^{n}\right]=-2 n f^{n}$ for $n \in \mathbb{N}_{0}$.\\
(ii) For $n>0$,

$$
\left[f, e^{n}\right]=-n e^{n-1}(h+(n-1))=-n(h-(n-1)) e^{n-1}
$$

and

$$
\left[e, f^{n}\right]=n f^{n-1}(h-(n-1))=n(h+(n-1)) f^{n-1} .
$$

Proof. (i) Since $\operatorname{ad} h(a):=h a-a h$ is a derivation of $A$ and $[h, e]=2 e$ commutes with $e$, we obtain inductively $\left[h, e^{n}\right]=n[h, e] e^{n-1}=2 n e^{n}$. The second part of (i) is obtained similarly.\\
(ii) We calculate

$$
\begin{aligned}
& {\left[f, e^{n}\right]=\sum_{j=0}^{n-1} e^{j}[f, e] e^{n-j-1}=\sum_{j=0}^{n-1} e^{j}(-h) e^{n-j-1}} \\
& =-\sum_{j=0}^{n-1} e^{j}\left[h, e^{n-j-1}\right]-\sum_{j=0}^{n-1} e^{n-1} h \\
& =-\left(\sum_{j=0}^{n-1} 2(n-j-1) e^{n-1}\right)-n e^{n-1} h \\
& =-\left(\sum_{j=0}^{n-1} 2 j e^{n-1}\right)-n e^{n-1} h=-n(n-1) e^{n-1}-n e^{n-1} h
\end{aligned}
$$

In view of (i), this equals

$$
-n(n-1) e^{n-1}-n h e^{n-1}+n\left[h, e^{n-1}\right]=n(n-1) e^{n-1}-n h e^{n-1} .
$$

This is the first part of (ii). The second part is reduced to the first one by considering the triple ( $f,-h, e$ ), satisfying the same commutation relations as ( $e, h, f$ ).

Lemma 6.2.3. Let $V$ be a $\mathfrak{g}$-module and $\mathfrak{h} \subseteq \mathfrak{g}$ a subalgebra. Then, for any $\alpha, \beta \in \mathfrak{h}^{*}$,

$$
\mathfrak{g}_{\alpha}(\mathfrak{h}) \cdot V_{\beta}(\mathfrak{h}) \subseteq V_{\alpha+\beta}(\mathfrak{h}) .
$$

Proof. For $v_{\beta} \in V_{\beta}(\mathfrak{h}), x \in \mathfrak{h}$, and $y \in \mathfrak{g}_{\alpha}(\mathfrak{h})$, we have

$$
x \cdot\left(y \cdot v_{\beta}\right)=[x, y] \cdot v_{\beta}+y \cdot\left(x \cdot v_{\beta}\right)=\alpha(x) y \cdot v_{\beta}+\beta(x) y \cdot v_{\beta}=(\alpha+\beta)(x) y \cdot v_{\beta} .
$$

Proposition 6.2.4. Let $V$ be a finite-dimensional $\mathfrak{s l}_{2}(\mathbb{K})$-module and $v_{0} \in V$ an element with $e \cdot v_{0}=0$ and $h \cdot v_{0}=\lambda v_{0}$. Then\\
(i) $\lambda \in \mathbb{N}_{0}$.\\
(ii) $v_{0}$ generates a submodule isomorphic to $L(\lambda)$.

Proof. (i) Let $V_{\alpha}:=V_{\alpha}(h)$ be the $h$-eigenspace corresponding to the eigenvalue $\alpha$ on $V$, which is a weight space for the representation of the subalgebra $\mathfrak{h}=\mathbb{K} h$. From $v_{0} \in V_{\lambda}$ and $[h, f]=-2 f$, we obtain with Lemma 6.2.3 the relation $h \cdot\left(f^{n} \cdot v_{0}\right)=(\lambda-2 n)\left(f^{n} \cdot v_{0}\right)$.

We further obtain with Lemma 6.2.3:\\
$e \cdot\left(f^{n} \cdot v_{0}\right)=\left[e, f^{n}\right] \cdot v_{0}+f^{n} \cdot \underbrace{\left(e \cdot v_{0}\right)}_{=0}=n f^{n-1}(h-n+1) \cdot v_{0}=n(\lambda-n+1) f^{n-1} \cdot v_{0}$.\\
This shows that the submodule $W$ generated by $v_{0}$ is

$$
W=\operatorname{span}\left\{f^{n} \cdot v_{0}: n \in \mathbb{N}_{0}\right\}
$$

Since $V$ is finite-dimensional, $h$ has only finitely many eigenvalues on $V$. Hence there is a minimal $N \in \mathbb{N}_{0}$ with $f^{N+1} \cdot v_{0}=0$. From $e \cdot\left(f^{N+1} \cdot v_{0}\right)=0$, we derive that $\lambda=N \in \mathbb{N}_{0}$.\\
(ii) To see that $W \cong L(\lambda)$, we consider the following basis

$$
v_{k}:=\frac{f^{k} \cdot v_{0}}{\lambda(\lambda-1) \cdots(\lambda-k+1)}, \quad k=0, \ldots, \lambda,
$$

for $W$ (note that the denominator never vanishes because the product is empty for $\lambda=0$ ). For this basis, we have

$$
h \cdot v_{k}=(\lambda-2 k) v_{k}, \quad f \cdot v_{k}=(\lambda-k) v_{k+1}
$$

$e \cdot v_{0}=0$ and, for $k>0$,

$$
\begin{aligned}
e \cdot v_{k} & =\frac{k(\lambda-k+1)}{\lambda(\lambda-1) \cdots(\lambda-k+1)} f^{k-1} \cdot v_{0} \\
& =\frac{k}{\lambda(\lambda-1) \cdots(\lambda-k+2)} f^{k-1} \cdot v_{0}=k v_{k-1}
\end{aligned}
$$

With respect to this basis, $e, f$, and $h$ are represented by the same matrices as on $L(\lambda)$, and this shows that $W \cong L(\lambda)$.

Lemma 6.2.5. If $V$ is a finite-dimensional real vector space and $a, b \in \mathfrak{g} \mathfrak{l}(V)$ with $[a, b]=b$, then $b$ is nilpotent.

Proof. Apply Proposition 5.4.14 to the solvable subalgebra $\mathbb{K} a+\mathbb{K} b \subseteq \mathfrak{g} \mathfrak{l}(V)$. Then $[a, b]=b$ implies that $b$ is nilpotent.

Theorem 6.2.6 (Classification of Simple $\mathfrak{s} \mathfrak{l}_{2}(\mathbb{K})$-Modules). Each finitedimensional simple $\mathfrak{s} \mathfrak{l}_{2}(\mathbb{K})$-module is isomorphic to some $L(\lambda), \lambda \in \mathbb{N}_{0}$. For each $n \in \mathbb{N}$, there exists a simple $\mathfrak{s l}_{2}(\mathbb{K})$-module of dimension $n$ which is unique up to isomorphism.

Proof. Let ( $\rho, V$ ) be a simple $\mathfrak{s} \mathfrak{l}_{2}(\mathbb{K})$-module. We consider the solvable subalgebra $\mathfrak{b}:=\operatorname{span}\{e, h\}$. We apply Lemma 6.2.5 to $a:=\frac{1}{2} \rho(h)$ and $b:=\rho(e)$, to see that $\rho(e)$ is nilpotent. Let $d \in \mathbb{N}$ be minimal with $\rho(e)^{d}=0$. Then Lemma 6.2.2 yields

$$
0=\left[\rho(f), \rho(e)^{d}\right]=-d(\rho(h)-(d-1) \mathbf{1}) \rho(e)^{d-1}
$$

so that each nonzero $v_{0} \in \rho(e)^{d-1}(V)$ is an eigenvector of $\rho(h)$. In view of the simplicity of the module $V$, it is generated by $v_{0}$, and Proposition 6.2.4 shows that $V \cong L(\lambda)$. The remaining assertions are immediate from Example 6.2.1.

Remark 6.2.7. A particular interesting representation of $\mathfrak{s} \mathfrak{l}_{2}(\mathbb{K})$ is the $o s$ cillator representation. Here we consider the space

$$
\mathcal{P}=\mathbb{C}\left[x_{1}, \ldots, x_{n}\right]
$$

of complex-valued polynomials on $\mathbb{R}^{n}$. Let $\Delta=\sum_{j} \frac{\partial^{2}}{\partial x_{j}^{2}}$ be the Laplacian. We put $e:=\frac{1}{2} \Delta$ and $f=\frac{1}{2} m_{r^{2}}$ (multiplication operator with $r^{2}:=\sum_{j} x_{j}^{2}$ ), and $h:=E+\frac{n}{2} \mathbf{1}$, where $E=\sum_{j} x_{j} \frac{\partial}{\partial x_{j}}$ is the Euler operator, for which a homogeneous polynomial of degree $d$ is an eigenvector of degree $d$.

It is easily verified that $(h, e, f) \in \operatorname{End}(\mathcal{P})$ satisfies the commutation relations of $\mathfrak{s l}_{2}(\mathbb{R})$, so that $\mathcal{P}$ is an $\mathfrak{s l}_{2}(\mathbb{K})$-module (Exercise 6.2.3). This module plays an important role in quantum mechanics of systems on $\mathbb{R}^{n}$ with full rotational symmetry. An important example is the spherical harmonic oscillator on $\mathbb{R}^{3}$, corresponding to the hydrogen atom.

Proposition 6.2.8. For a finite-dimensional $\mathfrak{s l}_{2}(\mathbb{K})$-representation ( $\rho, V$ ), the following assertions hold:\\
(i) $\rho(h)$ is diagonalizable and the set $\mathcal{P}_{V}$ of all eigenvalues is contained in $\mathbb{Z}$.\\
(ii) $\mathcal{P}_{V}=-\mathcal{P}_{V}$.\\
(iii) (String property) If $\alpha, \alpha+2 k \in \mathcal{P}_{V}$ for some $k \in \mathbb{N}_{0}$, then $\alpha+2 j \in \mathcal{P}_{V}$ for $j=0,1, \ldots, k$.

Proof. In view of Weyl's Theorem 5.5.21, $V$ is a direct sum of simple submodules $V_{1}, \ldots, V_{m}$, and Theorem 6.2.6 implies that $V_{i} \cong \mathbf{L}\left(\lambda_{i}\right)$ for some $\lambda_{i} \in \mathbb{N}_{0}$.\\
(i) and (ii) now follow immediately from the corresponding property of the modules $L(\lambda)$ (Example 6.2.1).\\
(iii) In view of (ii), we may w.lo.g. assume that $\beta:=\alpha+2 k$ satisfies $|\beta| \geq|\alpha|$. Then we pick some simple submodule $V_{i} \cong L\left(\lambda_{i}\right)$ of $V$ such that $\beta$ is an eigenvalue $\left.\rho(h)\right|_{V_{i}}$. Then $\lambda_{i}-\beta \in 2 \mathbb{N}_{0}$ and all integers in

$$
[-|\beta|,|\beta|] \cap(\beta+2 \mathbb{Z})
$$

are eigenvalues of $\left.\rho(h)\right|_{V_{i}}$. This contains in particular the set of all integers of the form $\alpha+2 j, j=0,1, \ldots, k$, between $\alpha$ and $\beta$.

We consider the element

$$
\theta:=e^{\operatorname{ad} e} e^{-\operatorname{ad} f} e^{\operatorname{ad} e} \in \operatorname{Aut}\left(\mathfrak{s} \mathfrak{l}_{2}(\mathbb{K})\right)
$$

(Example 4.2.5) and

$$
\begin{aligned}
\sigma & :=\exp (e) \exp (-f) \exp (e)=\left(\begin{array}{ll}
1 & 1 \\
0 & 1
\end{array}\right)\left(\begin{array}{cc}
1 & 0 \\
-1 & 1
\end{array}\right)\left(\begin{array}{ll}
1 & 1 \\
0 & 1
\end{array}\right) \\
& =\left(\begin{array}{cc}
0 & 1 \\
-1 & 0
\end{array}\right) \in \mathrm{SL}_{2}(\mathbb{K})
\end{aligned}
$$

Then Lemma 3.4.1 implies for $z \in \mathfrak{s l}_{2}(\mathbb{K})$ the relation $\theta(z)=\sigma z \sigma^{-1}$, hence in particular

$$
\theta(h)=-h, \quad \theta(e)=-f, \quad \text { and } \quad \theta(f)=-e .
$$

Lemma 6.2.9. Let ( $\rho, V$ ) be a finite-dimensional representation of $\mathfrak{s l}_{2}(\mathbb{K})$ and $\sigma_{V}:=e^{\rho(e)} e^{-\rho(f)} e^{\rho(e)} \in \mathrm{GL}(V)$. Then

$$
\sigma_{V} \rho(z) \sigma_{V}^{-1}=\rho\left(\sigma z \sigma^{-1}\right)
$$

for $z \in \mathfrak{s l}_{2}(\mathbb{K})$,

$$
\sigma_{V}\left(V_{\alpha}(\rho(h))\right)=V_{-\alpha}(\rho(h))
$$

for the eigenspaces of $\rho(h)$, and in particular

$$
\operatorname{dim} V_{\alpha}(\rho(h))=\operatorname{dim} V_{-\alpha}(\rho(h))
$$

Proof. For $z \in \mathfrak{s l}_{2}(\mathbb{K})$, we obtain with Lemma 3.4.1 the relation

$$
\sigma_{V} \rho(z) \sigma_{V}^{-1}=\rho\left(\sigma z \sigma^{-1}\right)
$$

This implies that

$$
\sigma_{V}\left(V_{\alpha}(\rho(h))\right)=V_{-\alpha}(\rho(h))
$$

because we have for $v \in V_{\alpha}(\rho(h))$ :

$$
\rho(h)\left(\sigma_{V}(v)\right)=\sigma_{V}\left(\sigma_{V}^{-1} \rho(h) \sigma_{V}\right)(v)=\sigma_{V} \rho(-h)(v)=-\alpha \sigma_{V}(v) .
$$

This completes the proof.

\subsubsection*{Exercises for Section 6.2}
Exercise 6.2.1. We consider the 2 -dimensional nonabelian complex Lie algebra $\mathfrak{g}$ in which we choose a basis ( $h, e$ ) satisfying $[h, e]=e$. In the following, $V$ denotes a $\mathfrak{g}$-module and $\rho: \mathfrak{g} \rightarrow \mathfrak{g} \mathfrak{l}(V)$ the corresponding representation. Classify all finite-dimensional $\mathfrak{g}$-modules $V$ for which $\rho(h)$ is diagonalizable.

Exercise 6.2.2. Let $A=\mathbb{K}\left[Z, Z^{-1}\right]$ be the algebra of Laurent polynomials with coefficients in the field $\mathbb{K}$. Show that every derivation of $A$ is of the form $D:=f \frac{d}{d Z}$ for some $f \in A$.

Exercise 6.2.3. On the space $V=C^{\infty}\left(\mathbb{R}^{n}\right)$ we consider the operators

$$
e:=\frac{1}{2} \sum_{j} \frac{\partial^{2}}{\partial x_{j}^{2}}, \quad f=\frac{1}{2} m_{r^{2}}
$$

(multiplication operator with $r^{2}:=\sum_{j} x_{j}^{2}$ ), and $h:=\frac{n}{2} \mathbf{1}+\sum_{j} x_{j} \frac{\partial}{\partial x_{j}}$. Show that $h, e, f$ satisfy the $\mathfrak{s l}_{2}$-relations

$$
[h, e]=2 e, \quad[h, f]=-2 f, \quad \text { and } \quad[e, f]=h .
$$

\subsection*{Root Decompositions of Semisimple Lie Algebras}
The technique of root decompositions is particularly fruitful for semisimple Lie algebras $\mathfrak{g}$ because, for this class of Lie algebras, all elements $h$ of a Cartan subalgebra $\mathfrak{h}$ turn out to be ad-semisimple, i.e., ad $h$ is a semisimple endomorphism of $\mathfrak{g}$. For complex Lie algebras, we thus obtain a root space decomposition diagonalizing ad $\mathfrak{h}$.

Proposition 6.3.1. Let $\mathfrak{g}$ be a semisimple Lie algebra, $\mathfrak{h}$ a splitting Cartan subalgebra of $\mathfrak{g}$, and $m_{\lambda}:=\operatorname{dim} \mathfrak{g}^{\lambda}$.\\
(i) $\kappa_{\mathfrak{g}}\left(h, h^{\prime}\right)=\sum_{\lambda \in \Delta(\mathfrak{g}, \mathfrak{h})} m_{\lambda} \lambda(h) \lambda\left(h^{\prime}\right)$ for $h, h^{\prime} \in \mathfrak{h}$.\\
(ii) If $\lambda+\mu \neq 0$, then $\mathfrak{g}^{\lambda}$ and $\mathfrak{g}^{\mu}$ are orthogonal with respect to the CartanKilling form.

Proof. (i) Both sides of the equation define a symmetric bilinear form on $\mathfrak{h}$. By polarization, it suffices to verify the equality for $h=h^{\prime}$. But $\lambda(h)$ is the only eigenvalue of $\operatorname{ad} h$ on $\mathfrak{g}^{\lambda}$. Therefore, the Jordan canonical form, together with the root decomposition (6.1), shows that $\operatorname{tr}\left(\operatorname{ad}(h)^{2}\right)=\sum_{\lambda} m_{\lambda} \lambda(h)^{2}$.\\
(ii) By Proposition 6.1.5, we have $\operatorname{ad}(x) \operatorname{ad}(y) \mathfrak{g}^{\nu} \subseteq \mathfrak{g}^{\lambda+\mu+\nu}$ for $x \in \mathfrak{g}^{\lambda}$ and $y \in \mathfrak{g}^{\mu}$. Choose a basis for $\mathfrak{g}$ consisting of elements in the $\mathfrak{g}^{\nu}$. Then since $\nu+\lambda+\mu \neq \nu$, the trace of $\operatorname{ad}(x) \operatorname{ad}(y)$ equals zero.

Proposition 6.3.2. Let $\mathfrak{g}$ be a semisimple Lie algebra and $\mathfrak{h} \subseteq \mathfrak{g}$ be a splitting Cartan subalgebra.\\
(i) The Cartan-Killing form $\kappa$ of $\mathfrak{g}$ induces a nondegenerate pairing of $\mathfrak{g}^{\alpha}$ and $\mathfrak{g}^{-\alpha}$, i.e., for $x \in \mathfrak{g}^{\alpha}$ and $y \in \mathfrak{g}^{-\alpha}$,

$$
\kappa\left(x, \mathfrak{g}^{-\alpha}\right)=\{0\} \Rightarrow x=0 \quad \text { and } \quad \kappa\left(\mathfrak{g}^{\alpha}, y\right)=\{0\} \Rightarrow y=0
$$

In particular, $m_{\alpha}=m_{-\alpha}$ and $\left.\kappa\right|_{\mathfrak{h} \times \mathfrak{h}}$ is nondegenerate.\\
(ii) The Lie algebra $\mathfrak{h}$ is a toral Cartan subalgebra. In particular, $\mathfrak{h}$ is abelian and $\mathfrak{g}^{\alpha}=\mathfrak{g}_{\alpha}$ for each root $\alpha$.\\
(iii) $\Delta(\mathfrak{g}, \mathfrak{h})$ spans $\mathfrak{h}^{*}$.

Proof. (i) This claim immediately follows from Proposition 6.3.1(ii) since the Cartan-Killing form is nondegenerate.\\
(ii) Since the adjoint representation of $\mathfrak{g}$ is injective ( $\mathfrak{z}(\mathfrak{g})=\{0\}$ ) and the set $(\operatorname{ad} \mathfrak{h})_{s}$ is commutative because it is diagonalized by the root decomposition, it suffices to show that for each $h \in \mathfrak{h}$, the operator ad $h$ is semisimple.

In view of Proposition 5.3.10, the nilpotent Jordan component $(\operatorname{ad} h)_{n}$ is a derivation of $\mathfrak{g}$, and since all derivations of $\mathfrak{g}$ are inner (Theorem 5.5.14), there exists an $h_{n} \in \mathfrak{g}$ with $\operatorname{ad}\left(h_{n}\right)=(\operatorname{ad} h)_{n}$. Since the derivation $(\operatorname{ad} h)_{n}$ commutes with $(\operatorname{ad} \mathfrak{h})_{s}$, we have

$$
\operatorname{ad}\left((\operatorname{ad} \mathfrak{h})_{s} h_{n}\right)=\left[(\operatorname{ad} \mathfrak{h})_{s}, \operatorname{ad} h_{n}\right]=\{0\},
$$

so that the injectivity of ad entails that $h_{n} \in \mathfrak{g}^{0}(\mathfrak{h})=\mathfrak{h}$ (cf. Proposition 5.1.10 and Lemma 6.1.13). As $\operatorname{ad}\left(h_{n}\right)$ is nilpotent, all roots vanish on $h_{n}$, and the formula in Proposition 6.3.1(i) implies that $\kappa\left(h_{n}, \mathfrak{h}\right)=\{0\}$. Now $h_{n}=0$ follows from (i). This proves that ad $h$ is diagonalizable for each $h \in \mathfrak{h}$ and since ad is injective, $\mathfrak{h} \cong \operatorname{ad} \mathfrak{h}$ is abelian.\\
(iii) As a consequence of (ii) and the injectivity of the adjoint representation, $\Delta(\mathfrak{g}, \mathfrak{h}) \subseteq \mathfrak{h}^{*}$ separates the points of $\mathfrak{h}$, and this is equivalent to (iii).

Since the Cartan-Killing form is nondegenerate on the Cartan subalgebra, we can assign to every root $\alpha$ a uniquely determined element $h_{\alpha}^{\prime} \in \mathfrak{h}$ via the equation

$$
\kappa\left(h, h_{\alpha}^{\prime}\right)=\alpha(h)
$$

Further, we can introduce a bilinear form on $\mathfrak{h}^{*}$ via

$$
(\alpha, \beta):=\kappa\left(h_{\alpha}^{\prime}, h_{\beta}^{\prime}\right)=\alpha\left(h_{\beta}^{\prime}\right)=\beta\left(h_{\alpha}^{\prime}\right) .
$$

Lemma 6.3.3. For $\alpha \in \Delta(\mathfrak{g}, \mathfrak{h})$,

$$
[x, y]=\kappa(x, y) h_{\alpha}^{\prime} \quad \text { for } x \in \mathfrak{g}_{\alpha}, y \in \mathfrak{g}_{-\alpha}
$$

Proof. Recall that Proposition 6.3.2(ii) implies that $[h, x]=\alpha(h) x$ for $h \in \mathfrak{h}$ and $x \in \mathfrak{g}_{\alpha}$. Both sides of the equation are in $\mathfrak{h}$, hence (6.10) follows from

$$
\kappa(h,[x, y])=\kappa([h, x], y)=\alpha(h) \kappa(x, y)=\kappa\left(h, \kappa(x, y) h_{\alpha}^{\prime}\right)
$$

since the Cartan-Killing form is nondegenerate on $\mathfrak{h}$ by Proposition 6.3.2.

\subsection*{1 $\mathfrak{s l}_{2}$-Triples}
The following theorem is the starting point of a complete classification of simple Lie algebras. It emphasizes the special role which the algebra $\mathfrak{s l}_{2}(\mathbb{K})$ plays in the theory.

Theorem 6.3.4 ( $\mathfrak{s} \mathfrak{l}_{2}$-Theorem). Let $\mathfrak{g}$ be a semisimple Lie algebra, $\mathfrak{h} \subseteq \mathfrak{g}$ a splitting Cartan subalgebra, and $\alpha \in \Delta(\mathfrak{g}, \mathfrak{h})$.\\
(i) For every root $\alpha \in \Delta(\mathfrak{g}, \mathfrak{h})$, we have $(\alpha, \alpha) \neq 0$, and there are elements $e_{\alpha} \in \mathfrak{g}_{\alpha}, f_{\alpha} \in \mathfrak{g}_{-\alpha}$, and $h_{\alpha} \in\left[\mathfrak{g}_{\alpha}, \mathfrak{g}_{-\alpha}\right]$ such that

$$
\left[h_{\alpha}, e_{\alpha}\right]=2 e_{\alpha}, \quad\left[h_{\alpha}, f_{\alpha}\right]=-2 f_{\alpha}, \quad \text { and } \quad\left[e_{\alpha}, f_{\alpha}\right]=h_{\alpha}
$$

In particular, $\mathfrak{g}(\alpha):=\operatorname{span}\left\{h_{\alpha}, e_{\alpha}, f_{\alpha}\right\}$ is a subalgebra isomorphic to $\mathfrak{s} \mathfrak{l}_{2}(\mathbb{K})$.\\
(ii) $m_{\alpha}=\operatorname{dim} \mathfrak{g}_{\alpha}=\operatorname{dim}\left(\left[\mathfrak{g}_{\alpha}, \mathfrak{g}_{-\alpha}\right]\right)=1$ and $\mathbb{Z} \alpha \cap \Delta=\{ \pm \alpha\}$.\\
(iii) $\alpha\left(h_{\alpha}\right)=2$.

Proof. (i) From $\kappa\left(\mathfrak{g}_{\alpha}, \mathfrak{g}_{-\alpha}\right) \neq\{0\}$ and Lemma 6.3.3, we obtain elements $e_{ \pm \alpha} \in \mathfrak{g}_{ \pm \alpha}$ with $\left[e_{\alpha}, e_{-\alpha}\right]=h_{\alpha}^{\prime}$. To see that $(\alpha, \alpha)=\alpha\left(h_{\alpha}^{\prime}\right)$ is nonzero, let us assume the contrary and consider some $\beta \in \Delta$. Then the subspace

$$
V:=\bigoplus_{k \in \mathbb{Z}} \mathfrak{g}_{\beta+k \alpha}
$$

of $\mathfrak{g}$ is invariant under $\operatorname{ad}\left(e_{ \pm \alpha}\right)$, so that

$$
\begin{aligned}
0 & =\operatorname{tr}\left(\left[\left.\operatorname{ad} e_{\alpha}\right|_{V},\left.\operatorname{ad} e_{-\alpha}\right|_{V}\right]\right)=\operatorname{tr}\left(\left.\operatorname{ad} h_{\alpha}^{\prime}\right|_{V}\right)=\sum_{k \in \mathbb{Z}}(\beta+k \alpha)\left(h_{\alpha}^{\prime}\right) \cdot m_{\beta+k \alpha} \\
& =\beta\left(h_{\alpha}^{\prime}\right) \cdot \sum_{k \in \mathbb{Z}} m_{\beta+k \alpha}
\end{aligned}
$$

Since $\sum_{k \in \mathbb{Z}} m_{\beta+k \alpha} \geq m_{\beta}>0$, we get $\beta\left(h_{\alpha}^{\prime}\right)=0$ for all roots $\beta$. But since the roots span $\mathfrak{h}^{*}$ (Proposition 6.3.2), this contradicts $h_{\alpha}^{\prime} \neq 0$. We conclude that $(\alpha, \alpha)=\alpha\left(h_{\alpha}^{\prime}\right) \neq 0$.

Set $h_{\alpha}:=2 \frac{h_{\alpha}^{\prime}}{\alpha\left(h_{\alpha}^{\prime}\right)}$. This proves (iii). From Proposition 6.3.2, we get an element $f_{\alpha} \in \mathfrak{g}_{-\alpha}$ with $\kappa\left(e_{\alpha}, f_{\alpha}\right)=\frac{2}{\alpha\left(h_{\alpha}^{\prime}\right)}$, so that Lemma 6.3.3 implies that $\left[e_{\alpha}, f_{\alpha}\right]=h_{\alpha}$. Now (i) follows from (iii).\\
(ii) We consider the subspace

$$
V:=\mathbb{K} f_{\alpha}+\mathfrak{h}+\sum_{n=1}^{\infty} \mathfrak{g}_{n \alpha}
$$

of $\mathfrak{g}$. One verifies easily that this subspace is invariant under $\operatorname{ad}(\mathfrak{g}(\alpha))$ because it is invariant under ad $\mathfrak{h},\left[f_{\alpha}, \mathfrak{h}\right]=\mathbb{K} f_{\alpha}$, and $\left[e_{\alpha}, \mathfrak{g}_{\beta}\right] \subseteq \mathfrak{g}_{\beta+\alpha}$. According to Lemma 6.2.9, we therefore have

$$
\operatorname{dim} V_{m}\left(\operatorname{ad} h_{\alpha}\right)=\operatorname{dim} V_{-m}\left(\operatorname{ad} h_{\alpha}\right)
$$

for all $m \in \mathbb{Z}$. This leads to

$$
\operatorname{dim} \mathfrak{g}_{\alpha}=\operatorname{dim} V_{2}\left(\operatorname{ad} h_{\alpha}\right)=\operatorname{dim} V_{-2}\left(\operatorname{ad} h_{\alpha}\right)=1
$$

and $\operatorname{dim} \mathfrak{g}_{n \alpha}=0$ for $n>1$. Now (ii) follows from (i).\\
Let $\mathfrak{g}$ be a semisimple Lie algebra and $\mathfrak{h} \subseteq \mathfrak{g}$ a splitting Cartan subalgebra. Then $\mathfrak{h}$ is splitting, and we obtain a root decomposition of $\mathfrak{g}$ with respect to $\mathfrak{h}$. For brevity we put $\Delta:=\Delta(\mathfrak{g}, \mathfrak{h})$. For a root $\alpha \in \Delta$, we have already seen that $\left[\mathfrak{g}_{\alpha}, \mathfrak{g}_{-\alpha}\right]$ is a one-dimensional subspace of $\mathfrak{h}$ on which $\alpha$ does not vanish. Hence there is a unique element

$$
\check{\alpha}=h_{\alpha} \in\left[\mathfrak{g}_{\alpha}, \mathfrak{g}_{-\alpha}\right] \quad \text { with } \quad \alpha(\check{\alpha})=2,
$$

called the coroot corresponding to $\alpha$ (cf. Theorem 6.3.4).\\
Lemma 6.3.5 (Root String Lemma). Let $\alpha, \beta \in \Delta$.\\
(i) For $\beta \in \Delta \backslash\{ \pm \alpha\}$, the set $\{k \in \mathbb{Z}: \beta+k \alpha \in \Delta\}$ is an interval in $\mathbb{Z}$. If it is of the form $[-p, q] \cap \mathbb{Z}$ with $p, q \in \mathbb{Z}$, then $p-q=\beta(\check{\alpha})$. In particular, $\beta(\check{\alpha}) \in \mathbb{Z}$.\\
(ii) If $\beta(\check{\alpha})<0$, then $\beta+\alpha \in \Delta$, and if $\beta(\check{\alpha})>0$, then $\beta-\alpha \in \Delta$.\\
(iii) If $\left[\mathfrak{g}_{\alpha}, \mathfrak{g}_{\beta}\right]=\{0\}$, then $\beta(\check{\alpha}) \geq 0$.\\
(iv) If $\alpha+\beta \neq 0$, then $\left[\mathfrak{g}_{\alpha}, \mathfrak{g}_{\beta}\right]=\mathfrak{g}_{\alpha+\beta}$.

Proof. (i) We consider the subspace $V:=\sum_{k \in \mathbb{Z}} \mathfrak{g}_{\beta+k \alpha}$. Note that $\beta \neq\{ \pm \alpha\}$ implies that 0 is not contained in $\beta+\mathbb{Z} \alpha$ (Theorem 6.3.4(ii)). From $\left[\mathfrak{g}_{\gamma}, \mathfrak{g}_{\delta}\right] \subseteq \mathfrak{g}_{\gamma+\delta}$, we derive that $V$ is a $\mathfrak{g}(\alpha)$-submodule of $\mathfrak{g}$. The eigenvalues of $\check{\alpha}=h_{\alpha}$ on $V$ are given by

$$
\mathcal{P}_{V}:=\{(\beta+k \alpha)(\check{\alpha}): \beta+k \alpha \in \Delta\}=\beta(\check{\alpha})+2\{k: \beta+k \alpha \in \Delta\} .
$$

Hence the string property of $\mathfrak{s} \mathfrak{l}_{2}$-modules (Proposition 6.2.8) implies the string property of the root system.

Next we note that $\beta \in \Delta$ leads to $p \geq 0$. In view of Proposition 6.2.8, we have $\mathcal{P}_{V}=-\mathcal{P}_{V}$. Therefore,

$$
\beta(\check{\alpha})-2 p=(\beta-p \alpha)(\check{\alpha})=-(\beta+q \alpha)(\check{\alpha})=-\beta(\check{\alpha})-2 q .
$$

(ii) If $\beta(\check{\alpha})<0$, then (i) leads to $q>0$ and hence to $\beta+\alpha \in \Delta$. The second assertion follows similarly.\\
(iii) As all multiplicities of the eigenvalues of $\check{\alpha}$ on $V$ are 1 (Theorem 6.3.4(ii)), the $\mathfrak{s l}_{2}(\mathbb{K})$-module $V$ is simple and isomorphic to $L(\beta(\check{\alpha})+2 q)$ (apply Proposition 6.2.4 to a nonzero element of $\mathfrak{g}_{\beta+q \alpha}$ ). This immediately shows that

$$
\left[\mathfrak{g}_{\alpha}, \mathfrak{g}_{\beta+k \alpha}\right]=\mathfrak{g}_{\beta+(k+1) \alpha} \quad \text { for } \quad k=-p,-p+1, \ldots, q-1
$$

If $\left[\mathfrak{g}_{\alpha}, \mathfrak{g}_{\beta}\right]=\{0\}$, then $V:=\sum_{k \leq 0} \mathfrak{g}_{\beta+k \alpha}$ is invariant under $\mathfrak{g}(\alpha) \cong \mathfrak{s l}_{2}(\mathbb{K})$ and $\beta(\check{\alpha})$ is the maximal eigenvalue of $\operatorname{ad}(\check{\alpha})$ on $V$. Hence Proposition 6.2.8 shows that $\beta(\check{\alpha}) \geq 0$.\\
(iv) We may assume that $\beta+\alpha \in \Delta$ (otherwise $\mathfrak{g}_{\alpha+\beta}=\{0\}$ ), so that $q \geq 1$. Then (iv) follows from (6.11).

Lemma 6.3.6. Each root $\alpha \in \Delta$ satisfies $\mathbb{K} \alpha \cap \Delta=\{ \pm \alpha\}$.\\
Proof. Suppose that $\alpha$ and $c \alpha$ lie in $\Delta$ for some $c \in \mathbb{K}$. By Lemma 6.3.5, $2 c=c \alpha(\check{\alpha}) \in \mathbb{Z}$, so that $c \in \frac{1}{2} \mathbb{Z}$. For symmetry reasons, also $c^{-1} \in \frac{1}{2} \mathbb{Z}$, so that $c \in\left\{ \pm 2, \pm 1, \pm \frac{1}{2}\right\}$. From the $\mathfrak{s l}_{2}$-Theorem 6.3.4, we know already that $\mathbb{Z} \alpha \cap \Delta=\{ \pm \alpha\}$, which rules out the cases $c= \pm 2$. The cases $c= \pm \frac{1}{2}$ are likewise ruled out by applying the same argument to $c \alpha$ instead.

Lemma 6.3.7. The subspace $\mathfrak{h}_{\mathbb{R}}:=\operatorname{span}\{\check{\alpha}: \alpha \in \Delta\}$ has the following properties:\\
(i) $\alpha\left(\mathfrak{h}_{\mathbb{R}}\right)=\mathbb{R}$ for every root $\alpha \in \Delta$.\\
(ii) $\kappa$ is positive definite on $\mathfrak{h}_{\mathbb{R}}$.\\
(iii) $\mathfrak{h}_{\mathbb{R}}$ spans the space $\mathfrak{h}\left(\mathfrak{h}_{\mathbb{R}}=\mathfrak{h}\right.$ for $\left.\mathbb{K}=\mathbb{R}\right)$.\\
(iv) $\mathfrak{h}=\mathfrak{h}_{\mathbb{R}} \oplus i \mathfrak{h}_{\mathbb{R}}$ if $\mathbb{K}=\mathbb{C}$.

Proof. (i) This follows from $\beta(\check{\alpha}) \in \mathbb{Z}$ for $\alpha, \beta \in \Delta$, which is a consequence of the Root String Lemma 6.3.5.\\
(ii) This follows from Proposition 6.3.1 and (i).\\
(iii) If $\beta \in \mathfrak{h}^{*}$ vanishes on $\mathfrak{h}_{\mathbb{R}}$, then $0=\beta(\check{\alpha})=\kappa\left(h_{\beta}^{\prime}, \check{\alpha}\right)$ for each root $\alpha$, and since $\check{\alpha}$ is a nonzero multiple of $h_{\alpha}^{\prime}$, this implies that $\alpha\left(h_{\beta}^{\prime}\right)=0$. As $\Delta$ spans $\mathfrak{h}^{*}$, we get $\beta=0$, and this implies that $\operatorname{span}_{\mathbb{C}} \mathfrak{h}_{\mathbb{R}}=\mathfrak{h}$ for $\mathbb{K}=\mathbb{C}$.\\
(iv) $0 \leq \kappa(x, x)$ and $0 \leq \kappa(i x, i x)=-\kappa(x, x)$ for $x \in \mathfrak{h}_{\mathbb{R}} \cap i \mathfrak{h}_{\mathbb{R}}$. Hence $\kappa(x, x)=0$ and therefore $x=0$ by (ii).

\subsubsection*{Examples}
The following lemma is a useful tool to see that the examples discussed below are indeed semisimple Lie algebras.\\
Lemma 6.3.8. Suppose that $\mathfrak{g}=\mathfrak{h}+\sum_{\alpha \in \Delta} \mathfrak{g}_{\alpha}$ is a Lie algebra and $\mathfrak{h}$ a toral Cartan subalgebra such that\\
(i) $\mathfrak{g}(\alpha):=\mathfrak{g}_{\alpha}+\mathfrak{g}_{-\alpha}+\left[\mathfrak{g}_{\alpha}, \mathfrak{g}_{-\alpha}\right] \cong \mathfrak{s} \mathfrak{l}_{2}(\mathbb{K})$ for each root $\alpha$, and\\
(ii) $\mathfrak{z}(\mathfrak{g})=\{0\}$.

Then $\mathfrak{g}$ is semisimple.\\
Proof. Let $\mathfrak{r}:=\mathfrak{r a d}(\mathfrak{g})$ be the solvable radical of $\mathfrak{g}$. As an ideal, it is $\mathfrak{h}$ invariant, hence adapted to the root space decomposition $\mathfrak{r}=\mathfrak{r}_{0}+\sum_{\alpha} \mathfrak{r}_{\alpha}$ (Exercise 2.1.1). Since all semisimple subalgebras $\mathfrak{s}$ of $\mathfrak{g}$ intersect $\mathfrak{r}$ trivially (otherwise $\mathfrak{s} \cap \mathfrak{r}$ would be a nontrivial solvable ideal of $\mathfrak{s}$ ), $\mathfrak{r}_{\alpha} \subseteq \mathfrak{g}(\alpha) \cap \mathfrak{r}=\{0\}$. Hence $\mathfrak{r} \subseteq \mathfrak{h}$, and, in view of $\left[\mathfrak{r}, \mathfrak{g}_{\alpha}\right] \subseteq \mathfrak{r} \cap \mathfrak{g}_{\alpha}=\{0\}$, we get $\mathfrak{r} \subseteq \bigcap_{\alpha \in \Delta}$ ker $\alpha= \mathfrak{z}(\mathfrak{g})=\{0\}$.

Example 6.3.9 (The special linear Lie algebra). Let

$$
\mathfrak{g}:=\mathfrak{s} \mathfrak{l}_{n}(\mathbb{K}):=\left\{x \in \mathfrak{g} \mathfrak{l}_{n}(\mathbb{K}): \operatorname{tr} x=0\right\} .
$$

Then the subalgebra $\mathfrak{h}$ of diagonal matrices in $\mathfrak{g}$ is abelian, and for the matrix units $E_{i j}$ with a single nonzero entry 1 in position ( $i, j$ ), we have

$$
\left[\operatorname{diag}(h), E_{i j}\right]=\left(h_{i}-h_{j}\right) E_{i j} \quad \text { for } \quad h \in \mathbb{K}^{n}
$$

We conclude that $\mathfrak{h}$ is maximal abelian, $\mathfrak{g}^{0}(\mathfrak{h})=\mathfrak{g}_{0}(\mathfrak{h})=\mathfrak{h}$, so that $\mathfrak{h}$ is a Cartan subalgebra, and that the one-dimensional subspace $\mathbb{K} E_{i j}$ is a root space corresponding to the root $\varepsilon_{i}-\varepsilon_{j}$, where $\varepsilon_{i}(\operatorname{diag}(h)):=h_{i}$. The corresponding root system is

$$
A_{n-1}:=\left\{\varepsilon_{j}-\varepsilon_{k}: 1 \leq j \neq k \leq n\right\}
$$

Example 6.3.10 (The orthogonal Lie algebras). Let

$$
I_{n, n}:=\left(\begin{array}{cc}
\mathbf{1} & \mathbf{0} \\
\mathbf{0} & -\mathbf{1}
\end{array}\right) \in M_{2 n}(\mathbb{K}) \cong M_{2}\left(M_{n}(\mathbb{K})\right)
$$

and recall the Lie algebra

$$
\mathfrak{o}_{n, n}(\mathbb{K})=\left\{x \in \mathfrak{g} \mathfrak{l}_{2 n}(\mathbb{K}): x^{\top} I_{n, n}+I_{n, n} x=0\right\}
$$

In terms of block matrices, we then have

$$
\mathfrak{o}_{n, n}(\mathbb{K})=\left\{\left(\begin{array}{ll}
a & b \\
c & d
\end{array}\right) \in M_{2}\left(M_{n}(\mathbb{K})\right): a^{\top}=-a, d^{\top}=-d, b^{\top}=c\right\} .
$$

In this matrix presentation, it is quite inconvenient to describe a root decomposition of this Lie algebra. It is much simpler to use an equivalent description, based on the following observation. For

$$
g:=\frac{1}{\sqrt{2}}\left(\begin{array}{cc}
\mathbf{1} & \mathbf{1} \\
-\mathbf{1} & \mathbf{1}
\end{array}\right)
$$

we have

$$
g^{\top} I_{n, n} g=\left(\begin{array}{ll}
\mathbf{0} & \mathbf{1} \\
\mathbf{1} & \mathbf{0}
\end{array}\right)=: S .
$$

Hence $x \in \mathfrak{o}_{n, n}(\mathbb{K})$ is equivalent to $g^{-1} x g$ being contained in

$$
\mathfrak{g}:=\mathfrak{o}\left(\mathbb{K}^{2 n}, S\right):=\left\{x \in \mathfrak{g l}_{2 n}(\mathbb{K}): x^{\top} S+S x=0\right\}=g^{-1} \mathfrak{o}_{n, n}(\mathbb{K}) g .
$$

From

$$
\left(\begin{array}{ll}
a & b \\
c & d
\end{array}\right) \in \mathfrak{o}\left(\mathbb{K}^{2 n}, S\right) \quad \Leftrightarrow \quad d=-a^{\top}, b^{\top}=-b, c^{\top}=-c
$$

we immediately derive that $\mathfrak{g}$ has a root decomposition with respect to the maximal abelian subalgebra

$$
\mathfrak{h}=\operatorname{span}\left\{E_{j j}-E_{n+j, n+j}: j=1, \ldots, n\right\}=\left\{\operatorname{diag}(h,-h): h \in \mathbb{K}^{n}\right\} .
$$

The corresponding root system is

$$
D_{n}:=\left\{ \pm \varepsilon_{j} \pm \varepsilon_{k}: j, k=1, \ldots, n, j \neq k\right\}
$$

where $\varepsilon_{j}: \mathfrak{h} \rightarrow \mathbb{K}$ is the linear functional defined by $\varepsilon_{k}(\operatorname{diag}(h,-h)):=h_{k}$. Here the roots in the subsystem $A_{n-1}$ of $D_{n}$ correspond to the root spaces in the image of the embedding

$$
\mathfrak{s} \mathfrak{l}_{n}(\mathbb{K}) \rightarrow \mathfrak{o}\left(\mathbb{K}^{2 n}, S\right), \quad x \mapsto\left(\begin{array}{cc}
x & 0 \\
0 & -x^{\top}
\end{array}\right)
$$

Further, $\mathfrak{g}_{\varepsilon_{j}+\varepsilon_{k}}=\mathbb{K}\left(E_{j, n+k}-E_{k, n+j}\right)$ and $\mathfrak{g}_{-\varepsilon_{j}-\varepsilon_{k}}=\mathbb{K}\left(E_{n+k, j}-E_{n+j, k}\right)$.\\
For the symmetric matrix

$$
T:=\left(\begin{array}{lll}
\mathbf{0} & \mathbf{1} & 0 \\
\mathbf{1} & \mathbf{0} & 0 \\
0 & 0 & 1
\end{array}\right) \in M_{2 n+1}(\mathbb{K})
$$

we also obtain a Lie algebra

$$
\mathfrak{o}\left(\mathbb{K}^{2 n+1}, T\right):=\left\{x \in \mathfrak{g} \mathfrak{l}_{2 n+1}(\mathbb{K}): x^{\top} T+T x=0\right\}
$$

Then

$$
\begin{aligned}
& \left(\begin{array}{lll}
a & b & x \\
c & d & y \\
\widetilde{x} & \widetilde{y} & z
\end{array}\right) \in \mathfrak{o}\left(\mathbb{K}^{2 n+1}, T\right) \\
& \Leftrightarrow \quad\left(\begin{array}{ll}
a & b \\
c & d
\end{array}\right) \in \mathfrak{o}\left(\mathbb{K}^{2 n}, S\right), \widetilde{x}=-y^{\top}, \widetilde{y}=-x^{\top}, z=0
\end{aligned}
$$

implies that this Lie algebra has a root decomposition with respect to the maximal abelian subalgebra

$$
\mathfrak{h}=\operatorname{span}\left\{E_{j j}-E_{n+j, n+j}: j=1, \ldots, n\right\}=\left\{\operatorname{diag}(h,-h, 0): h \in \mathbb{K}^{n}\right\} .
$$

The corresponding root system is

$$
B_{n}:=\left\{ \pm \varepsilon_{j}, \pm \varepsilon_{j} \pm \varepsilon_{k}: j, k=1, \ldots, n, j \neq k\right\}
$$

where $\varepsilon_{j}: \mathfrak{h} \rightarrow \mathbb{K}$ is the linear functional defined by $\varepsilon_{k}(\operatorname{diag}(h,-h, 0)):= h_{k}$. Here the root spaces corresponding to roots in the subsystem $D_{n}$ of $B_{n}$ correspond to root spaces in the subalgebra $\mathfrak{o}\left(\mathbb{K}^{2 n}, S\right)$ (corresponding to $x=y=0$ ), and

$$
\mathfrak{g}_{\varepsilon_{j}}=\mathbb{K}\left(E_{j, 2 n+1}-E_{2 n+1, n+j}\right) \quad \text { and } \quad \mathfrak{g}_{-\varepsilon_{j}}=\mathbb{K}\left(E_{j+n, 2 n+1}-E_{2 n+1, j}\right)
$$

Example 6.3.11 (The symplectic Lie algebra). For the skew-symmetric matrix

$$
J:=\left(\begin{array}{cc}
\mathbf{0} & \mathbf{1} \\
-\mathbf{1} & \mathbf{0}
\end{array}\right) \in M_{2 n}(\mathbb{K})
$$

we obtain the symplectic Lie algebra

$$
\mathfrak{s} \mathfrak{p}_{n}(\mathbb{K})=\left\{x \in \mathfrak{g l}_{2 n}(\mathbb{K}): x^{\top} J+J x=0\right\} .
$$

Using

$$
\left(\begin{array}{ll}
a & b \\
c & d
\end{array}\right) \in \mathfrak{s p}_{n}(\mathbb{K}) \quad \Leftrightarrow \quad d=-a^{\top}, b^{\top}=b, c^{\top}=c
$$

we see that $\mathfrak{g}:=\mathfrak{s} \mathfrak{p}_{n}(\mathbb{K})$ has a root decomposition with respect to the maximal abelian subalgebra

$$
\mathfrak{h}=\operatorname{span}\left\{E_{j j}-E_{n+j, n+j}: j=1, \ldots, n\right\}=\left\{\operatorname{diag}(h,-h): h \in \mathbb{K}^{n}\right\} .
$$

The corresponding root system is

$$
C_{n}:=\left\{ \pm 2 \varepsilon_{j}, \pm \varepsilon_{j} \pm \varepsilon_{k}: j, k=1, \ldots, n, j \neq k\right\}
$$

where $\varepsilon_{j}: \mathfrak{h} \rightarrow \mathbb{K}$ is the linear functional defined by $\varepsilon_{k}(\operatorname{diag}(h,-h)):=h_{k}$. Again, the roots in the subsystem $A_{n-1}$ of $C_{n}$ correspond to the root spaces in the image of the embedding

$$
\mathfrak{s l}_{n}(\mathbb{K}) \rightarrow \mathfrak{s p}_{n}(\mathbb{K}), \quad x \mapsto\left(\begin{array}{cc}
x & 0 \\
0 & -x^{\top}
\end{array}\right) .
$$

Further,

$$
\mathfrak{g}_{\varepsilon_{j}+\varepsilon_{k}}=\mathbb{K}\left(E_{j, n+k}+E_{k, n+j}\right) \quad \text { and } \quad \mathfrak{g}_{-\varepsilon_{j}-\varepsilon_{k}}=\mathbb{K}\left(E_{n+k, j}+E_{n+j, k}\right)
$$

Example 6.3.12 (The odd symplectic Lie algebra). There is also an "odd-dimensional" version of the symplectic Lie algebra, which is not semisimple, due to the fact that no skew-symmetric form on an odd-dimensional space is nondegenerate.

For the skew-symmetric matrix

$$
J:=\left(\begin{array}{ccc}
\mathbf{0} & \mathbf{1} & 0 \\
-\mathbf{1} & \mathbf{0} & 0 \\
0 & 0 & 0
\end{array}\right) \in M_{2 n+1}(\mathbb{K})
$$

we define the odd symplectic Lie algebra

$$
\mathfrak{s p}_{n+\frac{1}{2}}(\mathbb{K}):=\left\{x \in \mathfrak{g l}_{2 n+1}(\mathbb{K}): x^{\top} J+J x=0\right\}
$$

Then

$$
\left(\begin{array}{ccc}
a & b & \widetilde{x} \\
c & d & \widetilde{y} \\
x & y & z
\end{array}\right) \in \mathfrak{s} \mathfrak{p}_{n+\frac{1}{2}}(\mathbb{K}) \quad \Leftrightarrow \quad\left(\begin{array}{cc}
a & b \\
c & d
\end{array}\right) \in \mathfrak{s} \mathfrak{p}_{n}(\mathbb{K}), \widetilde{x}=\widetilde{y}=0
$$

we see that $\mathfrak{g}:=\mathfrak{s p}_{n+\frac{1}{2}}(\mathbb{K})$ has a root decomposition with respect to the maximal abelian subalgebra

$$
\mathfrak{h}=\left\{\operatorname{diag}(h,-h, z): h \in \mathbb{K}^{n}, z \in \mathbb{K}\right\} .
$$

The corresponding root system is

$$
\left\{\widetilde{\varepsilon}_{j}^{ \pm}, \pm 2 \varepsilon_{j}, \pm \varepsilon_{j} \pm \varepsilon_{k}: j, k=1, \ldots, n, j \neq k\right\}
$$

where $\varepsilon_{j}: \mathfrak{h} \rightarrow \mathbb{K}$ is the linear functional defined by $\varepsilon_{k}(\operatorname{diag}(x,-x, z)):=x_{k}$ and $\widetilde{\varepsilon}_{k}^{ \pm}(\operatorname{diag}(x,-x, z)):=z \pm x_{k}$. Here the subsystem $C_{n}$ corresponds to the subalgebra $\mathfrak{s p}_{n}(\mathbb{K})$ (defined by $\widetilde{x}=\widetilde{y}=0$ ). Further,

$$
\mathfrak{g}_{\widetilde{\varepsilon}_{j}}=\mathbb{K} E_{2 n+1, n+j} \quad \text { and } \quad \mathfrak{g}_{-\widetilde{\varepsilon}_{j}}=\mathbb{K} E_{2 n+1, j}
$$

Note that

$$
\mathfrak{r}:=\left\{\left(\begin{array}{ccc}
0 & 0 & 0 \\
0 & 0 & 0 \\
x & y & z
\end{array}\right): x, y \in \mathbb{K}^{n}, z \in \mathbb{K}\right\}
$$

is a nilpotent ideal and that

$$
\mathfrak{s p}_{n+\frac{1}{2}}(\mathbb{K}) \cong \mathfrak{r} \rtimes \mathfrak{s p}_{n}(\mathbb{K})
$$

is a Levi decomposition. Restricting the roots to the Cartan subalgebra $\mathfrak{h}_{s}:=\mathfrak{h} \cap \mathfrak{s p}_{n}(\mathbb{K})$, we obtain the (nonreduced) root system

$$
B C_{n}:=\left\{ \pm \varepsilon_{j}, \pm 2 \varepsilon_{j}, \pm \varepsilon_{j} \pm \varepsilon_{k}: j, k=1, \ldots, n, j \neq k\right\}
$$

\subsection*{Abstract Root Systems and Their Weyl Groups}
In the previous section, we proved a number of results on the set of roots associated with a given splitting Cartan subalgebra of a semisimple Lie algebra. In this section, we distill some of these properties into the definition of an abstract root system and show how to derive further properties using this concept.

Definition 6.4.1. Let $E$ be a euclidian space, i.e., a finite-dimensional real vector space with an inner product ( $\cdot, \cdot$ ), i.e., a positive definite symmetric bilinear form. A reflection in $E$ is a linear map $\sigma \in \operatorname{End}(E)$ which induces the map - id on a line $\mathbb{R} \alpha$ and the identity on the hyperplane

$$
\alpha^{\perp}=\{\beta \in E \mid(\beta, \alpha)=0\}
$$

If $\alpha \in E \backslash\{0\}$ is given, then we write $P_{\alpha}:=\alpha^{\perp}$ and $\sigma_{\alpha}$ for $\sigma$. Then

$$
\sigma_{\alpha}(\beta)=\beta-2 \frac{(\beta, \alpha)}{(\alpha, \alpha)} \alpha
$$

for all $\beta \in E$.\\
Lemma 6.4.2. Let $E$ be a euclidian space and $\Phi \subseteq E$ be a finite subset which spans $E$, and which is invariant under all reflections $\sigma_{\alpha}$ with $\alpha \in \Phi$. Suppose that $\sigma \in \mathrm{GL}(E)$ fixes a hyperplane $P$ pointwise and maps some nonzero $\alpha \in \Phi$ to $-\alpha$. If $\sigma$ leaves $\Phi$ invariant, then $\sigma=\sigma_{\alpha}$.

Proof. Let $\tau:=\sigma \sigma_{\alpha} \in \mathrm{GL}(E)$. Then $\tau(\Phi)=\Phi$ (since $\Phi$ is finite) and $\tau(\alpha)= \alpha$. The linear automorphism $\widetilde{\tau} \in \mathrm{GL}(E / \mathbb{R} \alpha)$ induced by $\tau$ coincides with the linear automorphism $\widetilde{\sigma} \in \mathrm{GL}(E / \mathbb{R} \alpha)$ induced by $\sigma$. Since $E=P+\mathbb{R} \alpha$, our assumption yields $\widetilde{\sigma}=\mathrm{id}$. Together, we obtain that $\tau$ is split and 1 is its only eigenvalue. Therefore, $\tau_{s}=\mathrm{id}$, so that $\tau-\mathrm{id}$ is nilpotent, i.e., $\tau$ is unipotent. On the other hand, the $\tau$-invariance of the finite set $\Phi$ shows that there has to be a power $\tau^{k}$ which keeps $\Phi$ pointwise fixed. But $\Phi$ spans $E$, so that $\tau^{k}=\mathrm{id}$. As $\tau$ is unipotent, it follows that $\tau=\mathrm{id}$.

Definition 6.4.3. Let $E$ be a euclidian space and $\Delta \subseteq E \backslash\{0\}$ be a finite subset which spans $E$. Then $\Delta$ is called a reduced root system if it satisfies the following conditions\\
(R1) $\Delta \cap \mathbb{R} \alpha=\{ \pm \alpha\}$ for all $\alpha \in \Delta$,\\
(R2) $\sigma_{\alpha}(\Delta) \subseteq \Delta$ for all $\alpha \in \Delta$,\\
(R3) For $\alpha \in \Delta$ the coroot $\check{\alpha}:=\frac{2 \alpha}{(\alpha, \alpha)}$ satisfies $(\beta, \check{\alpha}) \in \mathbb{Z}$ for all $\beta \in \Delta$.\\
It is called root system if it only satisfies (R2) and (R3). If $\Delta$ is a root system, then we call the group $W=W(\Delta)$ generated by the reflections $\sigma_{\alpha}, \alpha \in \Delta$, the Weyl group of the root system.

Remark 6.4.4. Let $\mathfrak{g}$ be a semisimple Lie algebra and $\mathfrak{h} \subseteq \mathfrak{g}$ a splitting Cartan subalgebra. Recall the euclidian space $\mathfrak{h}_{\mathbb{R}}$ from Lemma 6.3.7 together with the finite subset $\Delta=\Delta(\mathfrak{g}, \mathfrak{h})$ of all roots with respect to $\mathfrak{h}$. Then Lemma 6.3.6 shows that $\Delta \cap \mathbb{R} \alpha=\{ \pm \alpha\}$ for all $\alpha \in \Delta$. Moreover, Lemma 6.3.5 yields $(\beta, \check{\alpha}) \in \mathbb{Z}$ for all $\alpha, \beta \in \Delta$. The same lemma also shows

$$
\beta-(\beta, \check{\alpha}) \alpha \in \Delta
$$

since, in the notation of Lemma 6.3.5, we have $r \leq-(\beta, \check{\alpha}) \leq s$. Thus $\Delta$ is a reduced root system. In particular, we see that to every pair ( $\mathfrak{g}, \mathfrak{h}$ ) consisting of a semisimple complex Lie algebra $\mathfrak{g}$ and a Cartan subalgebra $\mathfrak{h}$, a Weyl group $W(\mathfrak{g}, \mathfrak{h}):=W(\Delta)$ is assigned in a canonical way.

Remark 6.4.5. If $\Delta$ is a nonreduced root system and $\alpha, c \alpha \in \Delta$ for some $c>1$, then

$$
c \cdot(\alpha, \check{\alpha})=2 c \in \mathbb{Z}
$$

implies that $c \in \frac{1}{2} \mathbb{Z}$. Further, $(c \alpha)^{\sim}=\frac{1}{c} \check{\alpha}$ leads to $\frac{1}{c}(\alpha, \check{\alpha})=\frac{2}{c} \in \mathbb{Z}$, so that $c=2$. We therefore get

$$
\Delta \cap \mathbb{R} \alpha=\{ \pm \alpha, \pm 2 \alpha\} .
$$

Remark 6.4.6. The angle $\theta \in[0, \pi]$ between $\alpha$ and $\beta$ is defined by the identity

$$
\|\alpha\|\|\beta\| \cos \theta=(\alpha, \beta),
$$

where $\|\alpha\|=\sqrt{(\alpha, \alpha)}$ is the norm of the euclidian space $E$. We have

$$
2 \frac{\|\beta\|}{\|\alpha\|} \cos \theta=2 \frac{(\beta, \alpha)}{(\alpha, \alpha)}=(\beta, \check{\alpha}) \quad \text { and } \quad(\alpha, \check{\beta})(\beta, \check{\alpha})=4 \cos ^{2} \theta .
$$

Hence, if $(\beta, \check{\alpha}),(\alpha, \check{\beta}) \in \mathbb{Z}$, then we also have $4 \cos ^{2} \theta \in \mathbb{Z}$, and there are only the following possibilities, provided $\|\alpha\| \leq\|\beta\|$ :

\begin{center}
\begin{tabular}{|c||c|c|c|c|c|c|c|c|c|c|c|}
\hline
$(\alpha, \check{\beta})$ & 0 & 1 & -1 & 1 & -1 & 1 & -1 & 1 & -1 & 2 & -2 \\
\hline
$(\beta, \check{\alpha})$ & 0 & 1 & -1 & 2 & -2 & 3 & -3 & 4 & -4 & 2 & -2 \\
\hline
$\theta$ & $\frac{\pi}{2}$ & $\frac{\pi}{3}$ & $\frac{2 \pi}{3}$ & $\frac{\pi}{4}$ & $\frac{3 \pi}{4}$ & $\frac{\pi}{6}$ & $\frac{5 \pi}{6}$ & 0 & $\pi$ & 0 & $\pi$ \\
\hline
$\frac{\|\beta\|^{2}}{\|\alpha\|^{2}}$ & arb. & 1 & 1 & 2 & 2 & 3 & 3 & 4 & 4 & 1 & 1 \\
\hline
\end{tabular}
\end{center}

Proposition 6.4.7. Let $\Delta \subseteq E$ be a root system with Weyl group $W$. If $\tau \in \mathrm{GL}(E)$ leaves $\Delta$ invariant, then\\
(i) $\tau \sigma_{\alpha} \tau^{-1}=\sigma_{\tau \alpha}$ for all $\alpha \in \Delta$.\\
(ii) $(\beta, \check{\alpha})=(\tau(\beta), \tau(\alpha))$ for all $\alpha, \beta \in \Delta$.

Proof. (i) Note that $\tau \sigma_{\alpha} \tau^{-1}(\tau(\beta))=\tau \sigma_{\alpha}(\beta) \in \tau(\Delta) \subseteq \Delta$. Since $\Delta$ is finite and $\tau: \Delta \rightarrow \Delta$ is injective, we have $\tau(\Delta)=\Delta$. Hence we also have $\tau \sigma_{\alpha} \tau^{-1}(\Delta) \subseteq \Delta$. Further, $\tau \sigma_{\alpha} \tau^{-1}$ keeps the hyperplane $\tau\left(\alpha^{\perp}\right)$ pointwise fixed, and it maps $\tau \alpha$ to $-\tau \alpha$. Hence Lemma 6.4.2 shows that $\tau \sigma_{\alpha} \tau^{-1}=\sigma_{\tau \alpha}$.\\
(ii) In view of (i), this follows by comparison of the formulas

$$
\tau \sigma_{\alpha} \tau^{-1}(\tau(\beta))=\tau(\beta-(\beta, \check{\alpha}) \alpha)=\tau(\beta)-(\beta, \check{\alpha}) \tau(\alpha)
$$

and $\sigma_{\tau(\alpha)}(\tau(\beta))=\tau(\beta)-(\tau(\beta), \tau(\alpha)) \tau(\alpha)$. \(\square\)

Lemma 6.4.8. Let $\Delta$ be a root system, and suppose that $\alpha, \beta \in \Delta$ are not proportional. If $(\alpha, \beta)>0$, then $\alpha-\beta \in \Delta$.

Proof. Since $(\alpha, \beta)$ is positive if and only if $(\alpha, \check{\beta})$ is positive, Remark 6.4.6 shows that we have $(\alpha, \breve{\beta})=1$ or $(\beta, \check{\alpha})=1$. If $(\alpha, \breve{\beta})=1$, then $\alpha-\beta= \sigma_{\beta}(\alpha) \in \Delta$. Similarly, for $(\beta, \check{\alpha})=1$, we have $\beta-\alpha=\sigma_{\alpha}(\beta) \in \Delta$, hence $\alpha-\beta \in \Delta$. \(\square\)

\subsubsection*{Simple Roots}
Definition 6.4.9. Let $\Delta \subseteq E$ be a root system and $\Pi \subseteq \Delta$. Then $\Pi$ is called a basis for $\Delta$ if $\Pi$ is a basis for the vector space $E$, and if every root $\beta \in \Delta$ is of the form $\beta=\sum_{\alpha \in \Pi} k_{\alpha} \alpha$, where either all $k_{\alpha} \in \mathbb{N}_{0}$, or all $k_{\alpha} \in-\mathbb{N}_{0}$. In this case, the elements of $\Pi$ are called simple roots. The height of the root $\beta=\sum_{\alpha \in \Pi} k_{\alpha} \alpha$ is the number $\sum_{\alpha \in \Pi} k_{\alpha}$. Roots with positive height are called positive and roots with negative height are called negative. The set of positive roots is denoted $\Delta^{+}$, the set of negative roots by $\Delta^{-}$. We define a partial order $\prec$ on $E$ by

$$
\alpha \prec \beta \quad: \Longleftrightarrow \quad \beta-\alpha \in \sum_{\gamma \in \Delta^{+}} \mathbb{N}_{0} \gamma .
$$

Lemma 6.4.10. Let $\Delta$ be a root system and $\Pi$ be a basis for $\Delta$. Suppose that $\alpha, \beta \in \Pi$ with $\alpha \neq \beta$. Then $(\alpha, \beta) \leq 0$ and $\alpha-\beta$ is not a root.

Proof. Since $\Pi$ is a basis for $E, \alpha$ and $\beta$ cannot be proportional. If we have $(\alpha, \beta)>0$, then Lemma 6.4.8 shows that $\alpha-\beta \in \Delta$. But this contradicts the definition of a basis for $\Delta$.

Lemma 6.4.11. Let $M \subseteq E$ be contained in an open half space of $E$, i.e., there is a $\lambda \in E$ with $(\lambda, \alpha)>0$ for all $\alpha \in M$, and $(\alpha, \beta) \leq 0$ for all $\alpha, \beta \in M$ with $\alpha \neq \beta$. Then $M$ is linearly independent.

Proof. Suppose that $\sum_{\alpha \in M} r_{\alpha} \alpha=0$ with $r_{\alpha} \in \mathbb{R}$, and set

$$
M_{ \pm}:=\left\{\alpha \in M \mid \pm r_{\alpha}>0\right\}
$$

Then

$$
\nu:=\sum_{\alpha \in M_{+}}\left|r_{\alpha}\right| \alpha=\sum_{\beta \in M_{-}}\left|r_{\beta}\right| \beta,
$$

and therefore

$$
(\nu, \nu)=\sum_{\alpha \in M_{+}, \beta \in M_{-}}\left|r_{\alpha} r_{\beta}\right|(\alpha, \beta) \leq 0,
$$

which leads to $\nu=0$. But we then have $0=(\lambda, \nu)=\sum_{\alpha \in M_{ \pm}}\left|r_{\alpha}\right|(\lambda, \alpha)$, which implies $M_{ \pm}=\emptyset$ since we otherwise arrive at the contradiction $r_{\alpha}=0$ for some $\alpha \in M_{ \pm}$.

Definition 6.4.12. (a) Let $\Delta \subseteq E$ be a root system and $\lambda \in E$. Then $\lambda$ is called regular if $\lambda \notin \alpha^{\perp}$ holds for all $\alpha \in \Delta$. Otherwise, $\lambda$ is called singular. The connected components of the set of the regular elements are called Weyl chambers. The Weyl chamber which contains the regular element $\lambda \in E$ is denoted by $\mathcal{C}(\lambda)$.\\
(b) For any regular element $\lambda \in E$, the set

$$
\Delta^{+}(\lambda):=\{\alpha \in \Delta \mid(\lambda, \alpha)>0\}
$$

is called the corresponding positive system. An element $\alpha \in \Delta^{+}(\lambda)$ is called decomposable if there are $\beta_{1}, \beta_{2} \in \Delta^{+}(\lambda)$ with $\alpha=\beta_{1}+\beta_{2}$, otherwise, it is called indecomposable.

Theorem 6.4.13. For each regular element $\lambda \in E$, the set $\Pi:=\Pi(\lambda)$ of indecomposable elements in $\Delta^{+}(\lambda)$ is a basis for $\Delta$. Conversely, every basis is of this form.

Proof. Claim 1: $\Pi(\lambda)$ is a basis for $\Delta$.\\
First, we show that every element of $\Delta^{+}(\lambda)$ can be written as a linear combination of elements of $\Pi(\lambda)$ with coefficients in $\mathbb{N}_{0}$. For this, we suppose that $\alpha \in \Delta^{+}(\lambda)$ cannot be written in this form, and that it has the smallest $(\alpha, \lambda)$ of all elements with this property. In particular, $\alpha \notin \Pi(\lambda)$, and there exist $\beta_{1}, \beta_{2} \in \Delta^{+}(\lambda)$ with $\alpha=\beta_{1}+\beta_{2}$, hence, $(\alpha, \lambda)=\left(\beta_{1}, \lambda\right)+\left(\beta_{2}, \lambda\right)$. But since $\left(\beta_{i}, \lambda\right)>0$ for $i=1,2$, the minimality of ( $\alpha, \lambda$ ) shows that the $\beta_{i}$ can be written as a linear combination of elements of $\Pi(\lambda)$ with coefficients in $\mathbb{N}_{0}$. Then this also holds for $\alpha$, contradicting our assumption.

As a consequence, we see that every $\beta \in \Delta$ is of the form $\beta= \sum_{\alpha \in \Pi(\lambda)} k_{\alpha} \alpha$, where either all $k_{\alpha} \in \mathbb{N}_{0}$ or all $k_{\alpha} \in-\mathbb{N}_{0}$. Since $\Delta$ spans the space $E$, it remains to be shown that $\Pi(\lambda)$ is linearly independent.

Next we show that $(\alpha, \beta) \leq 0$ for all $\alpha, \beta \in \Pi(\lambda)$ with $\alpha \neq \beta$. In fact, if $(\alpha, \beta)>0$, then $\alpha \notin-\mathbb{R}^{+} \beta$. By Remark 6.4.5, $\alpha$ and $\beta$ can only be proportional if $\alpha=2 \beta$ or $\beta=2 \alpha$, which would contradict the indecomposability of these elements. Hence we can apply Lemma 6.4.8, and we get $\alpha-\beta \in \Delta= \Delta^{+}(\lambda) \cup-\Delta^{+}(\lambda)$. If $\alpha-\beta \in \Delta^{+}(\lambda)$, then $\alpha=\beta+(\alpha-\beta)$, which contradicts the assumption that $\alpha$ is indecomposable. Similarly, $\beta-\alpha \in \Delta^{+}(\lambda)$ gives a contradiction by $\beta=\alpha+(\beta-\alpha)$. Now Claim 1 follows by Lemma 6.4.11 applied to $M=\Pi(\lambda)$.\\
Claim 2: Every basis $\Pi$ for $\Delta$ is of the form $\Pi(\lambda)$ for some regular element $\lambda \in E$.

We arrange $\Pi$ in the form $\alpha_{1}, \ldots, \alpha_{n}$, and consider the dual basis $\widetilde{\alpha}_{1}, \ldots, \widetilde{\alpha}_{n}$ for $E$ which is given by $\left(\alpha_{i}, \widetilde{\alpha}_{j}\right)=\delta_{i j}$. Set $\lambda:=\sum_{j=1}^{n} \widetilde{\alpha}_{j}$. Then $\left(\lambda, \alpha_{j}\right)=1$ for all $j=1, \ldots, n$, so that $(\lambda, \alpha)>0$ for all $\alpha \in \Pi$. Since every $\beta \in \Delta$ can be written as a linear combination of the $\alpha \in \Pi$ with coefficients of the same sign, $\lambda$ is regular. Then the set $\Delta^{+}$of positive roots defined by $\Pi$ satisfies $\Delta^{+} \subseteq \Delta^{+}(\lambda)$, which leads to $\Delta^{+}=\Delta^{+}(\lambda)$ because of $\Delta^{+} \cup-\Delta^{+}=\Delta=\Delta^{+}(\lambda) \cup-\Delta^{+}(\lambda)$. From the definition of a basis for $\Delta$, we see that $\Pi$ consists of indecomposable elements of $\Delta^{+}=\Delta^{+}(\lambda)$, and therefore it is contained in $\Pi(\lambda)$. On the other hand, the cardinalities of $\Pi$ and $\Pi(\lambda)$ are both equal to $n=\operatorname{dim} E$ since both sets are bases for $E$. This proves that $\Pi=\Pi(\lambda)$. \(\square\)

\subsubsection*{Weyl Chambers}
Remark 6.4.14. Let $\Delta \subseteq E$ be a root system and $\lambda, \lambda^{\prime} \in E$ be regular elements.\\
(i) $\mathcal{C}(\lambda)=\mathcal{C}\left(\lambda^{\prime}\right) \Leftrightarrow \Delta^{+}(\lambda)=\Delta^{+}\left(\lambda^{\prime}\right) \Leftrightarrow \Pi(\lambda)=\Pi\left(\lambda^{\prime}\right)$.\\
(ii) By (i) and Theorem 6.4.13, there is a bijection between the set of Weyl chambers and the set of bases for $\Delta$.\\
(iii) If $\Pi=\Pi(\lambda)$, then we call $\mathcal{C}(\Pi):=\mathcal{C}(\lambda)$ the fundamental chamber associated with the basis $\Pi$. It is given by

$$
\begin{aligned}
\mathcal{C}(\Pi) & =\{\beta \in E \mid(\beta, \alpha)>0 \text { for all } \alpha \in \Pi\} \\
& =\left\{\beta \in E \mid(\beta, \alpha)>0 \text { for all } \alpha \in \Delta^{+}\right\} .
\end{aligned}
$$

Lemma 6.4.15. Let $\Delta \subseteq E$ be a root system and $\Pi$ be a basis for $\Delta$.\\
(i) For $\alpha \in \Delta^{+} \backslash \Pi$, there exists a $\beta \in \Pi$ with $\alpha-\beta \in \Delta^{+}$.\\
(ii) If $\Delta$ is reduced and $\alpha \in \Pi$, then $\sigma_{\alpha}$ permutes the set $\Delta^{+} \backslash\{\alpha\}$.\\
(iii) Let $\alpha_{1}, \ldots, \alpha_{r} \in \Pi$ and set $\sigma_{i}:=\sigma_{\alpha_{i}}$. If $\sigma_{1} \cdots \sigma_{r-1}\left(\alpha_{r}\right) \in \Delta^{-}$, then there is an $s \in\{1, \ldots, r-1\}$ such that

$$
\sigma_{1} \cdots \sigma_{r}=\sigma_{1} \cdots \sigma_{s-1} \sigma_{s+1} \cdots \sigma_{r-1}
$$

Proof. (i) Suppose for all $\beta \in \Pi$, we have $(\alpha, \beta) \leq 0$. Then the set $\Pi \cup\{\alpha\}$ satisfies the assumptions of Lemma 6.4.11, hence is linearly independent. Since $\Pi$ is a basis for $E$, this cannot be the case, i.e., there is a $\beta \in \Pi$ with $(\alpha, \beta)>0$.\\
Case 1: $\alpha$ and $\beta$ are not proportional. Then Lemma 6.4.8 shows that $\alpha-\beta$ is a root.\\
Case 2: $\alpha$ and $\beta$ are proportional. Then Remark 6.4.5 shows that $\alpha=2 \beta$ since $\beta$ is indecomposable. But then $\alpha-\beta=\beta$ is a root.\\
Since $\alpha \in \Delta^{+} \backslash \Pi$, it is a linear combination of elements in $\Pi$ with at least two positive (integral) coefficients. Subtracting $\beta$ leaves at least one positive coefficient, so $\alpha-\beta$, being a root, has to be positive.\\
(ii) Let $\beta \in \Delta^{+} \backslash\{\alpha\}$ and $\beta=\sum_{\gamma \in \Pi} k_{\gamma} \gamma$ with $k_{\gamma} \in \mathbb{N}_{0}$. Since $\Delta$ is reduced, $\beta$ is not proportional to $\alpha$. Hence there is a $\gamma \neq \alpha$ with $k_{\gamma}>0$. Since

$$
\sigma_{\alpha}(\beta)=\beta-(\beta, \check{\alpha}) \alpha=\left(k_{\alpha}-(\beta, \check{\alpha})\right) \alpha+\sum_{\gamma \in \Pi \backslash\{\alpha\}} k_{\gamma} \gamma
$$

$\sigma_{\alpha}(\beta)$ has a positive coefficient $k_{\gamma}$ in its representation as a linear combination of simple roots. Thus, all coefficients are nonnegative, and $\sigma_{\alpha}(\beta) \in \Delta^{+}$. By $\sigma_{\alpha}(\alpha)=-\alpha$, we also have $\sigma_{\alpha}(\beta) \neq \alpha$, i.e., $\sigma_{\alpha}(\beta) \in \Delta^{+} \backslash\{\alpha\}$. Since the latter set is finite, the claim follows.\\
(iii) Set

$$
\beta_{i}:= \begin{cases}\sigma_{i+1} \cdots \sigma_{r-1}\left(\alpha_{r}\right) & \text { for } i=0, \ldots, r-2 \\ \alpha_{r} & \text { for } i=r-1\end{cases}
$$

Then $\beta_{0} \prec 0 \prec \beta_{r-1}$, and there is a minimal $s \in\{1, \ldots, r-1\}$ with $0 \prec \beta_{s}$. For this $s$, we have $\sigma_{s}\left(\beta_{s}\right)=\beta_{s-1} \prec 0$. In view of (ii), this shows $\beta_{s}=\alpha_{s}$. By Proposition 6.4.7, for $\sigma:=\sigma_{s+1} \cdots \sigma_{r-1}$, we have

$$
\sigma_{s}=\sigma_{\beta_{s}}=\sigma_{\sigma \alpha_{r}}=\sigma \sigma_{r} \sigma^{-1}=\left(\sigma_{s+1} \cdots \sigma_{r-1}\right) \sigma_{r}\left(\sigma_{r-1} \cdots \sigma_{s+1}\right)
$$

which shows the claim.\\
Proposition 6.4.16. Let $\Delta \subseteq E$ be a root system and $\Pi$ be a basis for $\Delta$.\\
(i) Every $\beta \in \Delta^{+}$can be written in the form $\alpha_{1}+\cdots+\alpha_{m}$ with $\alpha_{j} \in \Pi$ such that $\sum_{j=1}^{k} \alpha_{j} \in \Delta^{+}$for each $k \in\{1, \ldots, m\}$.\\
(ii) Let $\Delta$ be reduced and $\rho:=\frac{1}{2} \sum_{\beta \in \Delta^{+}} \beta$. Then $\sigma_{\alpha}(\rho)=\rho-\alpha$ for all $\alpha \in \Pi$.\\
(iii) Let $\sigma=\sigma_{1} \cdots \sigma_{r}$, where the $\sigma_{j}=\sigma_{\alpha_{j}}$ are reflections associated with the simple roots $\alpha_{j} \in \Pi$, and where $r$ is the minimal number of factors needed to represent $\sigma$ as such a product. Then $\sigma\left(\alpha_{r}\right) \in \Delta^{-}$.

Proof. (i) This immediately follows by induction using Lemma 6.4.15(i).\\
(ii) In view of $\sigma_{\alpha}(\alpha)=-\alpha$, this is a direct consequence of Lemma 6.4.15(ii).\\
(iii) This immediately follows from $\sigma\left(\alpha_{r}\right)=\sigma_{1} \cdots \sigma_{r}\left(\alpha_{r}\right)= \sigma_{1} \cdots \sigma_{r-1}\left(-\alpha_{r}\right)$ and Lemma 6.4.15(iii).

Theorem 6.4.17. Let $\Delta$ be a reduced root system, $W$ be the corresponding Weyl group, and II a basis for $\Delta$.\\
(i) For every regular element $\lambda \in E$, there is a $\sigma \in W$ such that

$$
(\sigma \lambda, \alpha)>0 \quad \text { for } \quad \alpha \in \Pi
$$

i.e., $\sigma(\mathcal{C}(\lambda))=\mathcal{C}(\Pi)$. In particular, $W$ acts transitively on the set of the Weyl chambers.\\
(ii) Let $\Pi^{\prime}$ be another basis for $\Delta$. Then there is a $\sigma \in W$ with $\sigma\left(\Pi^{\prime}\right)=\Pi$, i.e., the Weyl group also acts transitively on the set of the bases.\\
(iii) For every root $\alpha \in \Delta$, there is a $\sigma \in W$ with $\sigma(\alpha) \in \Pi$.\\
(iv) $W$ is generated by the $\sigma_{\alpha}$ with $\alpha \in \Pi$.\\
(v) If $\sigma(\Pi)=\Pi$ for $\sigma \in W$, then $\sigma=\mathbf{1}$.

Proof. Let $W^{\prime}$ be the subgroup of $W$, generated by the $\sigma_{\alpha}$ with $\alpha \in \Pi$. Suppose, (iii) holds for $W^{\prime}$ instead of $W$. Then for $\alpha \in \Delta$, we can find a $\sigma \in W^{\prime}$ with $\sigma \alpha \in \Pi$. Then by $\sigma_{\sigma \alpha}=\sigma \sigma_{\alpha} \sigma^{-1}$ (cf. Proposition 6.4.7), we obtain that $\sigma_{\alpha}=\sigma^{-1} \sigma_{\sigma \alpha} \sigma \in W^{\prime}$. Since $W$ is generated by the $\sigma_{\alpha}$ with $\alpha \in \Delta$, we get $W=W^{\prime}$, and hence (iv). Then (v) is an immediate consequence of Proposition 6.4.16(iii) applied to $\sigma$. Therefore, we see that it suffices to show (i)-(iii) for $W^{\prime}$ instead of $W$.\\
(i) Choose $\sigma \in W^{\prime}$ such that the number ( $\sigma \lambda, \rho$ ) is maximal for the given $\lambda$ and $\rho=\frac{1}{2} \sum_{\beta \in \Delta^{+}} \beta$ as above. For $\alpha \in \Pi$, we have $\sigma_{\alpha} \sigma \in W^{\prime}$, hence, by Proposition 6.4.16(ii), we get

$$
(\sigma \lambda, \rho) \geq\left(\sigma_{\alpha} \sigma \lambda, \rho\right)=\left(\sigma \lambda, \sigma_{\alpha} \rho\right)=(\sigma \lambda, \rho-\alpha)=(\sigma \lambda, \rho)-(\sigma \lambda, \alpha)
$$

This gives $\left(\lambda, \sigma^{-1} \alpha\right)=(\sigma \lambda, \alpha) \geq 0$ for all $\alpha \in \Pi$, hence also for all $\alpha \in \Delta^{+}$. Since $\lambda$ is regular, all inequalities are strict. By Remark 6.4.14, the claim follows.\\
(ii) By (i), this is an immediate implication of Remark 6.4.14.\\
(iii) Because of (ii), it suffices to show that $\alpha$ is an element of some basis. If $\beta \neq \pm \alpha$ is a root, then $\alpha$ and $\beta$ are not proportional since $\Delta$ is reduced. Thus, $\alpha^{\perp} \neq \beta^{\perp}$, and we can find a

$$
\lambda \in \alpha^{\perp} \backslash \bigcup_{\beta \in \Delta \backslash\{ \pm \alpha\}} \beta^{\perp} .
$$

By a small modification of $\lambda$ (adding, e.g., $\varepsilon \alpha$ for a small $\varepsilon>0$ ), we obtain a $\lambda^{\prime} \in E$ with

$$
\left|\left(\lambda^{\prime}, \beta\right)\right|>\left(\lambda^{\prime}, \alpha\right)>0 \quad \text { for } \quad \beta \in \Delta \backslash\{ \pm \alpha\} .
$$

Then $\lambda^{\prime}$ is regular, and we have $\alpha \in \Pi\left(\lambda^{\prime}\right)$ (cf. Theorem 6.4.13).\\
Remark 6.4.18 (The dual root system). If $\Delta \subseteq E$ is a root system, then we put

$$
\check{\Delta}:=\{\check{\alpha}: \alpha \in \Delta\} .
$$

We claim that $\check{\Delta}$ also is a root system, called the dual root system. It is reduced if and only if $\Delta$ is reduced.

To verify (R1) for $\check{\Delta}$ (if $\Delta$ is reduced), we note that $\check{\beta} \in \mathbb{R} \check{\alpha}$ implies $\beta \in \mathbb{R} \alpha$, and hence $\beta= \pm \alpha$, which in turn leads to $\check{\beta}= \pm \check{\alpha}$.

Since $\sigma_{\check{\alpha}}$ is the orthogonal reflection in $\alpha^{\perp}=\check{\alpha}^{\perp}$, we have $\sigma_{\alpha}=\sigma_{\check{\alpha}}$. As $\sigma_{\alpha}$ is an isometry, it satisfies $\sigma_{\alpha}(\check{\beta})=\sigma_{\alpha}(\beta)^{r}$, so that $\check{\Delta}$ satisfies (R2). Finally, we note that for $\alpha, \beta \in \Delta$, we have

$$
(\check{\alpha}, \check{\alpha})=\frac{4}{(\alpha, \alpha)}
$$

so that $(\check{\alpha})^{r}=\alpha$. Therefore,

$$
(\check{\alpha},(\check{\beta}))=(\beta, \check{\alpha}) \in \mathbb{Z},
$$

and we conclude that $\check{\Delta}$ also is a root system.\\
Now let

$$
\Delta^{+}=\Delta^{+}(\lambda)=\{\alpha \in \Delta:(\lambda, \alpha)>0\}
$$

be a positive system of $\Delta$. From the definition of the dual root system, it follows that $\Delta$ and $\Delta$ define the same set of regular elements. Therefore,

$$
\check{\Delta}^{+}:=\{\check{\alpha} \in \check{\Delta}:(\lambda, \check{\alpha})>0\}
$$

also is a positive system of the dual root system $\check{\Delta}$.

The elements $\alpha$ of the basis $\Pi=\Pi(\lambda)$ for the positive system $\Delta^{+}$are uniquely determined by the property that they generate extremal rays of the convex cone

$$
D\left(\Delta^{+}\right):=\sum_{\alpha \in \Delta^{+}} \mathbb{R}^{+} \alpha
$$

and satisfy $\frac{\alpha}{2} \notin \Delta$ (cf. Remark 6.4.5).\\
We define a map $\varphi: \Pi \rightarrow \check{\Delta}$ by $\varphi(\alpha)=\check{\alpha}$ if $2 \alpha \notin \Delta$ and $\varphi(\alpha):=\frac{\check{\alpha}}{2}=(2 \alpha)^{\check{ }}$ otherwise. Then $D\left(\Delta^{+}\right)=D\left(\check{\Delta}^{+}\right)$immediately shows that $\varphi(\Pi)$ is a root basis for $\check{\Delta}$. If, in addition, $\Delta$ is reduced, then

$$
\check{\Pi}:=\{\check{\alpha}: \alpha \in \Pi\}
$$

is a root basis for $\check{\Delta}^{+}$.

\subsubsection*{Notes on Chapter 6}
Cartan subalgebras actually occurred first in the work of W. Killing who classified the finite-dimensional simple complex Lie algebras (cf. [Kil89]). Unfortunately, Killing's work contained some serious gaps concerning the basic properties of Cartan subalgebras which were cleaned up later by Elie Cartan in his thesis [Ca94].

Serre's Theorem on the presentation of semisimple Lie algebras with a splitting Cartan subalgebra can be extended to a construction of a semisimple Lie algebra from an abstract root system $\Delta$ with a basis $\Pi=\left\{\alpha_{1}, \ldots, \alpha_{r}\right\}$. Then we put $a_{i j}:=\alpha_{j}\left(\check{\alpha}_{i}\right)$ and consider the Lie algebra $L(X, R)$ defined by the generators $h_{i}, e_{i}, f_{i}, i=1, \ldots, r$, and the relations

$$
\left[h_{i}, h_{j}\right]=0, \quad\left[h_{i}, e_{j}\right]=a_{i j} e_{j}, \quad\left[h_{i}, f_{j}\right]=-a_{i j} f_{j}, \quad\left[e_{i}, f_{j}\right]=\delta_{i j} h_{i},
$$

and for $i \neq j$,

$$
\left(\operatorname{ad} e_{i}\right)^{1-a_{i j}} e_{j}=0, \quad\left(\operatorname{ad} f_{i}\right)^{1-a_{i j}} f_{j}=0
$$

In this context, the main point is to show that $L(X, R)$ is a semisimple Lie algebra with the Cartan subalgebra $\mathfrak{h}=\operatorname{span}\left\{h_{1}, \ldots, h_{r}\right\}$ and a root system isomorphic to $\Delta$. In the 1960s, this description of the finite-dimensional semisimple Lie algebras was the starting point for the theory of Kac-MoodyLie algebras which are defined by the same set of generators and relations for more general matrices $\left(a_{i j}\right) \in M_{r}(\mathbb{Z})$, called generalized Cartan matrices.

\section*{Representation Theory of Lie Algebras}
Even though representation theory is not in the focus of this book, we provide in the present chapter the basic theory for Lie algebras as it repeatedly plays an important role in structural questions. In this chapter, we first introduce the universal enveloping algebra $\mathcal{U}(\mathfrak{g})$ of a Lie algebra $\mathfrak{g}$. This is a unital associative algebra containing $\mathfrak{g}$ as a Lie subalgebra and is generated by $\mathfrak{g}$. It has the universal property that each representation of $\mathfrak{g}$ extends uniquely to $\mathcal{U}(\mathfrak{g})$, so that any $\mathfrak{g}$-module becomes a $\mathcal{U}(\mathfrak{g})$-module. We may thus translate freely between Lie algebra modules and algebra modules, which is convenient for several representation theoretic constructions. The Poincaré-BirkhoffWitt (PBW) Theorem 7.1.9 provides crucial information on the structure of $\mathcal{U}(\mathfrak{g})$, including the injectivity of the natural map $\mathfrak{g} \rightarrow \mathcal{U}(\mathfrak{g})$.

As a first application of the PBW Theorem, we prove Serre's Theorem 7.2.10 which shows how to construct semisimple Lie algebras from root systems. The second application is the Highest Weight Theorem 7.3.15 providing a classification of irreducible finite-dimensional representations of (split) semisimple Lie algebras. This is the main result of the Cartan-Weyl Theory of simple modules of semisimple complex Lie algebras. In view of Weyl's Theorem that any module over such a Lie algebra is semisimple, the classification of the simple modules provides a complete picture of the finitedimensional representation theory of complex semisimple Lie algebras.

A third application of the PBW Theorem is Ado's Theorem 7.4.1 which says that each finite-dimensional Lie algebra has an injective finite-dimensional representation, i.e., can be viewed as a Lie algebra of matrices. Finally, we introduce basic cohomology theory for Lie algebras which has many applications. Here we use it to describe extensions of Lie algebras.

In general, modules of a Lie algebra $\mathfrak{g}$ are not semisimple, so that submodules need not have module complements. This leads to the concept of a nontrivial extension of modules which is most naturally dealt with in the context of Lie algebra cohomology which we develop in Section 7.5. We also explain how Lie algebra cohomology provides a tool to deal with the nontriviality of other extension problems such as central and abelian extensions of Lie algebras, for which it also provides a classifying parameter space. Several of the structure theoretic results encountered in Section 5.6 have natural\\
cohomological interpretations, such as Weyl's Theorem and Levi's Theorem which are reflected in the two Whitehead Lemmas.

\subsection*{The Universal Enveloping Algebra}
Representing a Lie algebra by linear maps leads to a mapping of the Lie algebra into an associative algebra such that the Lie bracket turns into the commutator bracket. A priori it is not clear that an injective map of this type exists, not even if we allow the associative algebra to be infinite-dimensional. The point of the enveloping algebra $\mathcal{U}(\mathfrak{g})$ of a Lie algebra $\mathfrak{g}$ is that every representation of $\mathfrak{g}$ on $V$ factors through a homomorphism $\mathcal{U}(\mathfrak{g}) \rightarrow \operatorname{End}(V)$ of associative algebras.

Definition 7.1.1. Let $\mathfrak{g}$ be a Lie algebra. A pair ( $\mathcal{U}(\mathfrak{g}), \sigma$ ), consisting of a unital associative algebra $\mathcal{U}(\mathfrak{g})$ and a homomorphism $\sigma: \mathfrak{g} \rightarrow \mathcal{U}(\mathfrak{g})_{L}$ of Lie algebras, is called a (universal) enveloping algebra of $\mathfrak{g}$ if it has the following universal property. For each homomorphism $f: \mathfrak{g} \rightarrow A_{L}$ of $\mathfrak{g}$ into the Lie algebra $A_{L}$, where $A$ is a unital associative algebra, there exists a unique homomorphism $\widetilde{f}: \mathcal{U}(\mathfrak{g}) \rightarrow A$ of unital associative algebras with $\widetilde{f} \circ \sigma=f$.\\
\includegraphics[max width=\textwidth, center]{2025_10_18_3eaf04e77f3cb364fce5g-179}

The universal property determines a universal enveloping algebra uniquely in the following sense:

Lemma 7.1.2 (Uniqueness of the enveloping algebra). If ( $\mathcal{U}(\mathfrak{g}), \sigma$ ) and $(\widetilde{\mathcal{U}}(\mathfrak{g}), \widetilde{\sigma})$ are two enveloping algebras of the Lie algebra $\mathfrak{g}$, then there exists an isomorphism $f: \mathcal{U}(\mathfrak{g}) \rightarrow \widetilde{\mathcal{U}}(\mathfrak{g})$ of unital associative algebras satisfying $f \circ \sigma=\widetilde{\sigma}$.\\
Proof. Since $\widetilde{\sigma}: \mathfrak{g} \rightarrow \widetilde{\mathcal{U}}(\mathfrak{g})_{L}$ is a homomorphism of Lie algebras, the universal property of the pair ( $\mathcal{U}(\mathfrak{g}), \sigma$ ) implies the existence of a unique algebra homomorphism

$$
f: \mathcal{U}(\mathfrak{g}) \rightarrow \widetilde{\mathcal{U}}(\mathfrak{g}) \quad \text { with } \quad f \circ \sigma=\widetilde{\sigma} .
$$

Similarly, the universal property of $(\widetilde{\mathcal{U}}(\mathfrak{g}), \widetilde{\sigma})$ implies the existence of an algebra homomorphism

$$
g: \widetilde{\mathcal{U}}(\mathfrak{g}) \rightarrow \mathcal{U}(\mathfrak{g}) \quad \text { with } \quad g \circ \widetilde{\sigma}=\sigma .
$$

Then $g \circ f: \mathcal{U}(\mathfrak{g}) \rightarrow \mathcal{U}(\mathfrak{g})$ is an algebra homomorphism with $(g \circ f) \circ \sigma=\sigma$, so that the uniqueness part of the universal property of ( $\mathcal{U}(\mathfrak{g}), \sigma$ ) yields $g \circ f= \operatorname{id}_{\mathcal{U}(\mathfrak{g})}$. We likewise get $f \circ g=\operatorname{id}_{\widetilde{\mathcal{U}}(\mathfrak{g})}$, showing that $f$ is an isomorphism of unital algebras. \(\square\)

To prove the existence of an enveloping algebra, we recall some basic algebraic concepts. Let $\mathcal{A}$ be an associative algebra. A subspace $J$ of $\mathcal{A}$ is called an ideal if

$$
\mathcal{A} J \cup J \mathcal{A} \subseteq J .
$$

Let $M$ be a subset of $\mathcal{A}$. Since the intersection of a family of ideals is again an ideal, the intersection $J_{M}$ of all ideals of $\mathcal{A}$ containing $M$ is the smallest ideal of $\mathcal{A}$ containing $M$. It is called the ideal generated by $M$. If $J$ is an ideal of $\mathcal{A}$, then the factor algebra $\mathcal{A} / J$ is the quotient vector space, endowed with the associative multiplication

$$
\left(a_{1}+J\right)\left(a_{2}+J\right):=a_{1} a_{2}+J \quad \text { for } \quad a_{1}, a_{2} \in \mathcal{A} .
$$

\subsubsection*{Existence}
Proposition 7.1.3 (Existence of an enveloping algebra). Each Lie algebra $\mathfrak{g}$ has an enveloping algebra ( $\mathcal{U}(\mathfrak{g}), \sigma$ ).

Proof. Let $\mathcal{T}(\mathfrak{g})$ be the tensor algebra of $\mathfrak{g}$ (Definition B.1.7) and consider the subset

$$
M:=\{x \otimes y-y \otimes x-[x, y] \in \mathcal{T}(\mathfrak{g}): x, y \in \mathfrak{g}\} .
$$

Then

$$
\mathcal{U}(\mathfrak{g}):=\mathcal{T}(\mathfrak{g}) / J_{M}
$$

is a unital associative algebra and

$$
\sigma: \mathfrak{g} \rightarrow \mathcal{U}(\mathfrak{g}), \quad \sigma(x):=x+J_{M},
$$

is a linear map satisfying

$$
\sigma([x, y])=[x, y]+J_{M}=x \otimes y-y \otimes x+J_{M}=\sigma(x) \sigma(y)-\sigma(y) \sigma(x),
$$

so that $\sigma$ is a homomorphism of Lie algebras $\mathfrak{g} \rightarrow \mathcal{U}(\mathfrak{g})_{L}$.\\
To verify the universal property for $(\mathcal{U}(\mathfrak{g}), \sigma)$, let $f: \mathfrak{g} \rightarrow A_{L}$ be a homomorphism of Lie algebras, where $A$ is a unital associative algebra. In view of the universal property of $\mathcal{T}(\mathfrak{g})$ (Lemma B.1.8), there exists an algebra homomorphism $\widehat{f}: \mathcal{T}(\mathfrak{g}) \rightarrow A$ with $\widehat{f}(x)=f(x)$ for all $x \in \mathfrak{g}$. Then $M \subseteq \operatorname{ker} \widehat{f}$, and since $\operatorname{ker} \widehat{f}$ is an ideal of $\mathcal{T}(\mathfrak{g})$, we also have $J_{M} \subseteq \operatorname{ker} \widehat{f}$, so that $\widehat{f}$ factors through an algebra homomorphism

$$
\widetilde{f}: \mathcal{U}(\mathfrak{g}) \rightarrow A \quad \text { with } \quad \widetilde{f} \circ \sigma=f .
$$

To see that $\widetilde{f}$ is unique, it suffices to note that $\sigma(\mathfrak{g})$ and $\mathbf{1}$ generate $\mathcal{U}(\mathfrak{g})$ as an associative algebra because $\mathfrak{g}$ and $\mathbf{1}$ generate $\mathcal{T}(\mathfrak{g})$ as an associative algebra. \(\square\)

Remark 7.1.4. The universal property of ( $\mathcal{U}(\mathfrak{g}), \sigma$ ) implies that each representation ( $\pi, V$ ) of $\mathfrak{g}$ defines a representation $\widetilde{\pi}: \mathcal{U}(\mathfrak{g}) \rightarrow \operatorname{End}(V)$, which is uniquely determined by $\widetilde{\pi} \circ \sigma=\pi$. From the construction of $\mathcal{U}(\mathfrak{g})$, we also know that the algebra $\mathcal{U}(\mathfrak{g})$ is generated by $\sigma(\mathfrak{g})$. This implies that for each $v \in V$, the subspace

$$
\mathcal{U}(\mathfrak{g}) \cdot v \subseteq V
$$

is the smallest subspace containing $v$ and invariant under $\mathfrak{g}$, i.e., the $\mathfrak{g}$ submodule of $V$ generated by $V$. Hence the enveloping algebra provides a tool to understand $\mathfrak{g}$-submodules of a $\mathfrak{g}$-module. But before we are able to use this tool effectively, we need some more information on the structure of $\mathcal{U}(\mathfrak{g})$.

\subsubsection*{The Poincaré-Birkhoff-Witt Theorem}
Note that the canonical map $\sigma: \mathfrak{g} \rightarrow \mathcal{U}(\mathfrak{g})_{L}$ by definition is a homomorphism of Lie algebras. But we do not know yet if it is injective, so that we obtain an embedding of $\mathfrak{g}$ into an associative algebra. Our next goal is the Poincaré-Birkhoff-Witt Theorem which entails in particular that $\sigma$ is injective.

Let $\left\{x_{1}, \ldots, x_{n}\right\}$ be a basis for $\mathfrak{g}$ and set $\xi_{i}:=\sigma\left(x_{i}\right)$. For a finite sequence $I=\left(i_{1}, \ldots, i_{k}\right)$ of natural numbers between 1 and $n$, we set $\xi_{I}:=\xi_{i_{1}} \cdots \xi_{i_{k}}$. If $i \in \mathbb{N}$, then we write $i \leq I$ if $i \leq i_{j}$ for all $j=1, \ldots, k$. We write $\mathcal{U}_{p}(\mathfrak{g}):=\sum_{k \leq p} \sigma(\mathfrak{g})^{k}$ and note that these subspaces are all finite-dimensional and satisfy

$$
\mathcal{U}_{p}(\mathfrak{g}) \mathcal{U}_{q}(\mathfrak{g}) \subseteq \mathcal{U}_{p+q}(\mathfrak{g}) \quad \text { for } \quad p, q \in \mathbb{N}_{0} .
$$

Next we construct a suitable basis for $\mathcal{U}(\mathfrak{g})$.\\
Lemma 7.1.5. Let $y_{1}, \ldots, y_{p} \in \mathfrak{g}$ and $\pi$ be a permutation of $\{1, \ldots, p\}$, then

$$
\sigma\left(y_{1}\right) \cdots \sigma\left(y_{p}\right)-\sigma\left(y_{\pi(1)}\right) \cdots \sigma\left(y_{\pi(p)}\right) \in \mathcal{U}_{p-1}(\mathfrak{g}) .
$$

Proof. Since every permutation is a composition of transpositions of neighboring elements, it suffices to prove the claim for $\pi(j)=j$ for $j \notin\{i, i+1\}$ and $\pi(i)=i+1$. But then we have

$$
\begin{aligned}
& \sigma\left(y_{1}\right) \cdots \sigma\left(y_{p}\right)-\sigma\left(y_{\pi(1)}\right) \cdots \sigma\left(y_{\pi(p)}\right) \\
& \quad=\sigma\left(y_{1}\right) \cdots \sigma\left(y_{i-1}\right)\left(\sigma\left(y_{i}\right) \sigma\left(y_{i+1}\right)-\sigma\left(y_{i+1}\right) \sigma\left(y_{i}\right)\right) \sigma\left(y_{i+2}\right) \cdots \sigma\left(y_{p}\right) \\
& \quad=\sigma\left(y_{1}\right) \cdots \sigma\left(y_{i-1}\right) \sigma\left(\left[y_{i}, y_{i+1}\right]\right) \sigma\left(y_{i+2}\right) \cdots \sigma\left(y_{p}\right) \in \mathcal{U}_{p-1}(\mathfrak{g})
\end{aligned}
$$

Lemma 7.1.6. The vector space $\mathcal{U}_{p}(\mathfrak{g})$ is spanned by the $\xi_{I}$ with increasing sequences I of length less than or equal to $p$.

Proof. It is clear that $\mathcal{U}_{p}(\mathfrak{g})$ is spanned by the $\xi_{I}$ for all sequences $I$ of length less than or equal to $p$. By induction on $p$, we have the claim for $\mathcal{U}_{p-1}(\mathfrak{g})$. But since for a rearrangement $I^{\prime}$ of the sequence $I$, we have $\xi_{I}-\xi_{I^{\prime}} \in \mathcal{U}_{p-1}(\mathfrak{g})$ by Lemma 7.1.5, we also obtain the claim for $\mathcal{U}_{p}(\mathfrak{g})$.

Let $\mathcal{P}:=\mathbb{K}\left[z_{1}, \ldots, z_{n}\right]$ be the associative algebra of polynomials over $\mathbb{K}$ in the commuting variables $z_{1}, \ldots, z_{n}$. For $i \in \mathbb{N} \cup\{0\}$, let $\mathcal{P}_{i}$ be the set of all polynomials of degree less or equal to $i$. As in $\mathcal{U}(\mathfrak{g})$, we write $z_{I}:=z_{i_{1}} \cdots z_{i_{k}}$ for a finite sequence $I$ of natural numbers between 1 and $n$. For the empty sequence, we set $z_{\emptyset}=1$.

Lemma 7.1.7. There exists a $\mathfrak{g}$-module structure on $\mathcal{P}$ with

$$
x_{i} \cdot z_{I}=z_{i} z_{I} \quad \text { for } i \leq I
$$

Proof. By induction on $k \in \mathbb{N}_{0}$, we construct linear maps

$$
\rho_{k}: \mathfrak{g} \rightarrow \operatorname{Hom}\left(\mathcal{P}_{k}, \mathcal{P}_{k+1}\right)
$$

with the following properties:\\
$\left(a_{k}\right) \rho_{k}\left(x_{i}\right) z_{I}=z_{i} z_{I}$ for $i \leq I$ and $z_{I} \in \mathcal{P}_{k}$.\\
$\left(b_{k}\right) \quad \rho_{k}\left(x_{i}\right) z_{I}-z_{i} z_{I} \in \mathcal{P}_{j}$ for $z_{I} \in \mathcal{P}_{j}$ and $j \leq k$.\\
( $c_{k}$ ) $\rho_{k}\left(x_{i}\right) \rho_{k}\left(x_{j}\right) z_{J}-\rho_{k}\left(x_{j}\right) \rho_{k}\left(x_{i}\right) z_{J}=\rho_{k}\left(\left[x_{i}, x_{j}\right]\right) z_{J}$ for $z_{J} \in \mathcal{P}_{k-1}$.\\
$\left.\left(d_{k}\right) \rho_{k}(x)\right|_{\mathcal{P}_{k-1}}=\rho_{k-1}(x)$ for $k>0$ and $x \in \mathfrak{g}$.\\
For $k=0$, we put $\rho_{0}\left(x_{i}\right) 1:=z_{i}$ for $i=1, \ldots, k$. Then $\left(a_{0}\right)$ and $\left(b_{0}\right)$ are satisfied because $x_{1}, \ldots, x_{k}$ is a basis for $\mathfrak{g}$, and ( $c_{0}$ ) and ( $d_{0}$ ) are empty conditions.

Now we assume that we already have $\rho_{k-1}$. In view of $\left(d_{k}\right)$, we have to show that the maps $\rho_{k-1}\left(x_{i}\right)$ can be extended to maps $\mathcal{P}_{k} \rightarrow \mathcal{P}_{k+1}$ so that $\left(a_{k}\right)-\left(c_{k}\right)$ are satisfied. Since $\mathcal{P}$ is commutative, the $z_{I}$ with increasing $I$ form a basis for $\mathcal{P}$. Thus, let $I:=\left(i_{1}, \ldots, i_{k}\right)$ be increasing. For $I_{1}:=\left(i_{2}, \ldots, i_{k}\right)$, by ( $a_{k-1}$ ), we have

$$
z_{I}=z_{i_{1}} z_{I_{1}}=\rho_{k-1}\left(x_{i_{1}}\right) z_{I_{1}}
$$

By $\left(b_{k-1}\right)$,

$$
w(I, i):=\rho_{k-1}\left(x_{i}\right) z_{I_{1}}-z_{i} z_{I_{1}} \in \mathcal{P}_{k-1}
$$

We set

$$
\rho_{k}\left(x_{i}\right) z_{I}:= \begin{cases}z_{i} z_{I} & \text { if } i \leq I \\ z_{i} z_{I}+\rho_{k-1}\left(x_{i_{1}}\right) w(I, i)+\rho_{k-1}\left(\left[x_{i}, x_{i_{1}}\right]\right) z_{I_{1}} & \text { otherwise }\end{cases}
$$

For this definition, $\left(a_{k}\right)$ and $\left(b_{k}\right)$ are obviously satisfied. But we still have to verify $\left(c_{k}\right)$. Two cases occur:

Case 1: $i \neq j$, and one of them is less than $J$.\\
In this case, by $\left[x_{i}, x_{j}\right]=-\left[x_{j}, x_{i}\right]$, we may assume that $j<i$ and $j \leq J$. Then with ( $a_{k-1}$ ) and ( $b_{k-1}$ ), we calculate

$$
\begin{aligned}
& \rho_{k}\left(x_{i}\right) \rho_{k-1}\left(x_{j}\right) z_{J}-\rho_{k}\left(x_{j}\right) \rho_{k-1}\left(x_{i}\right) z_{J} \\
= & \rho_{k}\left(x_{i}\right) z_{j} z_{J}-\rho_{k}\left(x_{j}\right) z_{i} z_{J}-\rho_{k-1}\left(x_{j}\right)\left(\rho_{k-1}\left(x_{i}\right) z_{J}-z_{i} z_{J}\right) \\
= & z_{i} z_{j} z_{J}+\rho_{k-1}\left(x_{j}\right)\left(\rho_{k-1}\left(x_{i}\right) z_{J}-z_{i} z_{J}\right)+\rho_{k-1}\left(\left[x_{i}, x_{j}\right]\right) z_{J} \\
& -z_{i} z_{j} z_{J}-\rho_{k-1}\left(x_{j}\right)\left(\rho_{k-1}\left(x_{i}\right) z_{J}-z_{i} z_{J}\right) \\
= & \rho_{k-1}\left(\left[x_{i}, x_{j}\right]\right) z_{J}=\rho_{k}\left(\left[x_{i}, x_{j}\right]\right) z_{J}
\end{aligned}
$$

Case 2: $J=\left(j_{1}, \ldots, j_{k-1}\right)$ and $j_{1}<i, j$.\\
We set $l:=j_{1}, L:=\left(j_{2}, \ldots, j_{k-1}\right)$, and abbreviate $\rho_{k}\left(x_{i}\right) z_{I}$ by $x_{i}\left(z_{I}\right)$. Then it follows from ( $a_{k-1}$ ) and Case 1 that

$$
x_{j}\left(z_{J}\right)=x_{j}\left(x_{l}\left(z_{L}\right)\right)=x_{l}\left(x_{j}\left(z_{L}\right)\right)+\left[x_{j}, x_{l}\right]\left(z_{L}\right)
$$

and, using ( $c_{k-1}$ ) and Case 1,

$$
\begin{aligned}
& x_{i}\left(x_{j}\left(z_{J}\right)\right)=x_{i}\left(x_{l}\left(x_{j}\left(z_{L}\right)\right)\right)+x_{i}\left(\left[x_{j}, x_{l}\right]\left(z_{L}\right)\right) \\
= & x_{l}\left(x_{i}\left(x_{j}\left(z_{L}\right)\right)\right)+\left[x_{i}, x_{l}\right]\left(x_{j}\left(z_{L}\right)\right)+\left[x_{j}, x_{l}\right]\left(x_{i}\left(z_{L}\right)\right)+\left[x_{i},\left[x_{j}, x_{l}\right]\right]\left(z_{L}\right)
\end{aligned}
$$

Combining this with ( $c_{k-1}$ ) and the Jacobi identity, we obtain

$$
\begin{aligned}
& x_{i}\left(x_{j}\left(z_{J}\right)\right)-x_{j}\left(x_{i}\left(z_{J}\right)\right) \\
= & x_{l}\left(x_{i}\left(x_{j}\left(z_{L}\right)\right)\right)-x_{l}\left(x_{j}\left(x_{i}\left(z_{L}\right)\right)\right)+\left[x_{i},\left[x_{j}, x_{l}\right]\right]\left(z_{L}\right)-\left[x_{j},\left[x_{i}, x_{l}\right]\right]\left(z_{L}\right) \\
= & x_{l}\left(\left[x_{i}, x_{j}\right]\left(z_{L}\right)\right)+\left[x_{i},\left[x_{j}, x_{l}\right]\right]\left(z_{L}\right)+\left[x_{j},\left[x_{l}, x_{i}\right]\right]\left(z_{L}\right) \\
= & {\left[x_{i}, x_{j}\right]\left(x_{l}\left(z_{L}\right)\right)+\left[x_{l},\left[x_{i}, x_{j}\right]\right]\left(z_{L}\right)+\left[x_{i},\left[x_{j}, x_{l}\right]\right]\left(z_{L}\right)+\left[x_{j},\left[x_{l}, x_{i}\right]\right]\left(z_{L}\right) } \\
= & {\left[x_{i}, x_{j}\right]\left(x_{l}\left(z_{L}\right)\right)=\left[x_{i}, x_{j}\right] z_{J} }
\end{aligned}
$$

This completes our induction. In view of ( $d_{k}$ ), we obtain a well-defined map

$$
\rho: \mathfrak{g} \rightarrow \mathfrak{g l}(\mathcal{P}) \quad \text { by }\left.\quad \rho(x)\right|_{\mathcal{P}_{k}}=\rho_{k}(x),
$$

and ( $c_{k}$ ) implies that it is a homomorphism of Lie algebras, so that we obtain on $\mathcal{P}$ a $\mathfrak{g}$-module structure, and $\left(a_{k}\right)$ ensures that it has the required property.

Proposition 7.1.8. The $\xi_{I}$ with increasing I form a basis for $\mathcal{U}(\mathfrak{g})$. In particular, the canonical map $\sigma: \mathfrak{g} \rightarrow \mathcal{U}(\mathfrak{g})$ is injective.

Proof. Let $\rho: \mathfrak{g} \rightarrow \operatorname{End}(\mathcal{P})_{L}$ be the Lie algebra homomorphism defining the module structure constructed in Lemma 7.1.7 and $\widetilde{\rho}: \mathcal{U}(\mathfrak{g}) \rightarrow \operatorname{End}(\mathcal{P})$ the algebra homomorphism determined by $\widetilde{\rho} \circ \sigma=\rho$. Then the property $\rho\left(x_{i}\right) z_{I}= z_{i} z_{I}$ for $i \leq I$ ensures that for $i_{1} \leq \cdots \leq i_{k}$, we have

$$
\widetilde{\rho}\left(\xi_{i_{1}} \cdots \xi_{i_{k}}\right)(\mathbf{1})=z_{i_{1}} \cdots z_{i_{k}}
$$

Then the linear map

$$
\varphi: \mathcal{U}(\mathfrak{g}) \rightarrow \mathcal{P}, \quad \xi \mapsto \widetilde{\rho}(\xi)(\mathbf{1})
$$

maps the set $\mathcal{B}$ of the $\xi_{I}$ with increasing $I$ to the (linearly independent) set of the $z_{I}$ with increasing $I$. Therefore, $\mathcal{B}$ is linearly independent, and the claim now follows from Lemma 7.1.6.

Theorem 7.1.9 (Poincaré-Birkhoff-Witt Theorem (PBW)). Let $\mathfrak{g}$ be a Lie algebra and $\left\{x_{1}, \ldots, x_{n}\right\}$ be a basis for $\mathfrak{g}$. Then

$$
\left\{x_{1}^{\mu_{1}} \cdots x_{n}^{\mu_{n}} \in \mathcal{U}(\mathfrak{g}) \mid \mu_{k} \in \mathbb{N} \cup\{0\}\right\}
$$

is a basis for $\mathcal{U}(\mathfrak{g})$.\\
Proof. We apply Proposition 7.1.8 and identify the $x_{i}$ with the $\xi_{i}=\sigma\left(x_{i}\right)$.

\subsubsection*{Exercises for Section 7.1}
Exercise 7.1.1. Let $\mathfrak{g}$ be a finite-dimensional Lie algebra and $\beta$ be a nondegenerate symmetric bilinear form on $\mathfrak{g}$. Suppose that $x_{1}, \ldots, x_{n}$ is a basis for $\mathfrak{g}$ and let $x^{1}, \ldots, x^{n}$ be the dual basis w.r.t. $\beta$, i.e., $\beta\left(x_{i}, x^{j}\right)=\delta_{i j}$.\\
(i) Show that the Casimir element $\Omega:=\sum_{i=1}^{n} x^{i} x_{i}$ lies in the center of $\mathcal{U}(\mathfrak{g})$.\\
(ii) Show that there is a nondegenerate symmetric invariant bilinear form on the oscillator algebra. Hence such forms do not only exist on semisimple Lie algebras.\\
(iii) Let $\mathfrak{g}=\mathfrak{s o}_{n}(\mathbb{R})$. Show that $\beta(x, y)=-\operatorname{tr}(x y)$ defines an invariant scalar product on $\mathfrak{s o}_{n}(\mathbb{R})$. For an orthonormal basis $x_{1}, \ldots, x_{n}$, we therefore have $\Omega:=\sum_{i=1}^{n} x_{i}^{2} \in Z\left(\mathcal{U}\left(\mathfrak{s o}_{n}(\mathbb{R})\right)\right)$.\\
(iv) Show that the operators of angular momentum

$$
x_{i} \frac{\partial}{\partial x_{j}}-x_{j} \frac{\partial}{\partial x_{i}}: C^{\infty}\left(\mathbb{R}^{n}\right) \rightarrow C^{\infty}\left(\mathbb{R}^{n}\right), \quad i, j=1, \ldots, n
$$

generate a Lie algebra which is isomorphic to $\mathfrak{s o}_{n}(\mathbb{R})$.\\
(v) The Laplace operator $\Delta=\sum_{i=1}^{n} \frac{\partial^{2}}{\partial x_{i}^{2}}$ commutes with the operators of angular momentum.

Exercise 7.1.2. A function $f \in C^{\infty}\left(\mathbb{R}^{n}\right)$ is called harmonic if $\Delta(f)=0$ for the Laplace operator $\Delta=\sum_{i=1}^{n} \frac{\partial^{2}}{\partial x_{i}^{2}}$. Show that the subspace $H \subseteq C^{\infty}\left(\mathbb{R}^{n}\right)$ of the harmonic functions is invariant under the angular momentum operators (cf. Exercise 7.1.1).

\subsection*{Generators and Relations for Semisimple Lie Algebras}
In this section, we shall use the root decomposition of a semisimple Lie algebra to find a description by generators and relations (Serre's Theorem).

\subsubsection*{A Generating Set for Semisimple Lie Algebras}
Proposition 7.2.1. Let $\mathfrak{g}$ be a semisimple Lie algebra and $\mathfrak{h} \subseteq \mathfrak{g}$ a splitting Cartan subalgebra. Fix a positive system $\Delta^{+} \subseteq \Delta$ and let $\Pi \subseteq \Delta^{+}$be the set of simple roots. For each $\alpha \in \Pi$, we fix a corresponding $\mathfrak{s l}_{2}$-triple ( $h_{\alpha}, e_{\alpha}, f_{\alpha}$ ). Then the following assertions hold:\\
(i) The subspace $\mathfrak{n}:=\sum_{\beta \in \Delta^{+}} \mathfrak{g}_{\beta}$ is a nilpotent subalgebra generated by $\left\{e_{\alpha}: \alpha \in \Pi\right\}$.\\
(ii) The Lie algebra $\mathfrak{g}$ is generated by the $\mathfrak{s l}_{2}$-subalgebras $\mathfrak{g}(\alpha), \alpha \in \Pi$, i.e., $\left\{h_{\alpha}, e_{\alpha}, f_{\alpha}: \alpha \in \Pi\right\}$ generates $\mathfrak{g}$. These elements satisfy the relations

$$
\begin{aligned}
{\left[h_{\alpha}, h_{\beta}\right] } & =0, \quad\left[h_{\alpha}, e_{\beta}\right]=\beta(\check{\alpha}) e_{\beta}, \quad\left[h_{\alpha}, f_{\beta}\right]=-\beta(\check{\alpha}) f_{\beta} \\
{\left[e_{\alpha}, f_{\beta}\right] } & =\delta_{\alpha, \beta} h_{\alpha}
\end{aligned}
$$

and for $\alpha \neq \beta$ in $\Pi$, we further have

$$
\left(\operatorname{ad} e_{\alpha}\right)^{1-\beta(\check{\alpha})} e_{\beta}=0, \quad\left(\operatorname{ad} f_{\alpha}\right)^{1-\beta(\check{\alpha})} f_{\beta}=0
$$

The relations (7.1) and (7.2) are called the Serre relations.\\
Proof. (i) If $\beta, \gamma \in \Delta^{+}$, then either $\beta+\gamma \in \Delta^{+}$or $\beta+\gamma$ is not a root. Hence $\mathfrak{n}$ is a subalgebra of $\mathfrak{g}$. Pick $x_{0} \in \mathfrak{h}_{\mathbb{R}}$ with $\Delta^{+}=\left\{\beta \in \Delta: \beta\left(x_{0}\right)>0\right\}$ and let

$$
m:=\min \Delta^{+}\left(x_{0}\right) \quad \text { and } \quad M:=\max \Delta^{+}\left(x_{0}\right)
$$

where $\Delta^{+}\left(x_{0}\right)=\left\{\beta\left(x_{0}\right): \beta \in \Delta^{+}\right\}$. For any $N \in \mathbb{N}$ with $N m>M$ we then have $C^{N}(\mathfrak{n})=\{0\}$, showing that $\mathfrak{n}$ is nilpotent.

In view of Proposition 6.4.16(i), every $\beta \in \Delta^{+}$can be written in the form $\alpha_{1}+\cdots+\alpha_{m}$ with $\alpha_{j} \in \Pi$ such that $\sum_{j=1}^{k} \alpha_{j} \in \Delta^{+}$for each $k \in\{1, \ldots, m\}$. Then Lemma 6.3.5 implies that

$$
\mathfrak{g}_{\beta}=\left[\mathfrak{g}_{\alpha_{m}},\left[\mathfrak{g}_{\alpha_{m-1}},\left[\ldots,\left[\mathfrak{g}_{\alpha_{2}}, \mathfrak{g}_{\alpha_{1}}\right] \ldots\right]\right]\right]
$$

and this proves (i).\\
(ii) From (i) we also derive that $\overline{\mathfrak{n}}:=\sum_{\alpha \in-\Delta^{+}} \mathfrak{g}_{\alpha}$ is generated by $\left\{f_{\beta}: \beta \in \Pi\right\}$. Therefore, the Lie algebra generated by the $\mathfrak{g}(\alpha), \alpha \in \Pi$, contains all root spaces and

$$
\operatorname{span}\left\{h_{\alpha}: \alpha \in \Pi\right\}=\operatorname{span}\left\{h_{\alpha}^{\prime}: \alpha \in \Pi\right\}=\mathfrak{h}
$$

follows from span $\Pi=\mathfrak{h}^{*}$.\\
It remains to verify the Serre relation. Since $h_{\alpha}=\check{\alpha}$, the first three relations are trivial, and the fact that $\alpha-\beta \notin \Delta$ for $\alpha \neq \beta$ in $\Pi$ implies that $\left[e_{\alpha}, f_{\beta}\right]=\{0\}$ in this case.

If $\alpha \neq \beta$, then we consider the $\mathfrak{g}(\alpha)$-submodule of $\mathfrak{g}$ generated by $f_{\beta}$. Since $\left[e_{\alpha}, f_{\beta}\right]=0$, this is a highest weight module $M$ with highest weight $-\beta(\check{\alpha})$, so that Proposition 6.2.4 implies that $\operatorname{dim} M=1-\beta(\check{\alpha})$ and $\left(\operatorname{ad} f_{\alpha}\right)^{\operatorname{dim} M} f_{\beta}=0$.

The first relation in (7.2) is obtained with similar arguments, applied to the $\mathfrak{s} \mathfrak{l}_{2}$-triple $\left(-h_{\alpha}, f_{\alpha}, e_{\alpha}\right)$ and the $\mathfrak{g}(\alpha)$-submodule generated by $e_{\beta}$.

Example 7.2.2. We have seen in Example 6.3.9 how to find a natural root decomposition of the Lie algebra $\mathfrak{s} \mathfrak{l}_{n}(\mathbb{K})$ with respect to the Cartan subalgebra $\mathfrak{h}$ of diagonal matrices. In the root system

$$
\Delta=\left\{\varepsilon_{j}-\varepsilon_{k}: 1 \leq j \neq k \leq n\right\}
$$

the subset

$$
\Delta^{+}=\left\{\varepsilon_{j}-\varepsilon_{k}: 1 \leq j<k \leq n\right\}
$$

is a natural positive system with root basis

$$
\Pi=\left\{\varepsilon_{1}-\varepsilon_{2}, \ldots, \varepsilon_{n-1}-\varepsilon_{n}\right\}
$$

Then

$$
\mathfrak{n}=\sum_{\alpha \in \Delta^{+}} \mathfrak{g}_{\alpha}=\operatorname{span}\left\{E_{j k}: j<k\right\}
$$

is the Lie algebra of strictly upper triangular matrices. It is generated by the root vectors $E_{j, j+1}, j=1, \ldots, n-1$. For each pair of indices $j \neq k$, we have

$$
\mathfrak{g}\left(\varepsilon_{j}-\varepsilon_{k}\right)=\operatorname{span}\left\{E_{j k}, E_{k j}, E_{j j}-E_{k k}\right\} \cong \mathfrak{s} \mathfrak{l}_{2}(\mathbb{K})
$$

and the subalgebras

$$
\mathfrak{g}\left(\varepsilon_{1}-\varepsilon_{2}\right), \quad \ldots, \quad \mathfrak{g}\left(\varepsilon_{n-1}-\varepsilon_{n}\right)
$$

sitting on the diagonal, generate $\mathfrak{s} \mathfrak{l}_{n}(\mathbb{K})$. Moreover, $\mathfrak{s} \mathfrak{l}_{n}(\mathbb{K})$ is also generated by the $2(n-1)$ element: $E_{j, j+1}, E_{j+1, j}, j=1, \ldots, n-1$.

\subsubsection*{Free Lie Algebras}
Free Lie algebras are defined via universal properties. Their existence then has to be proved separately. The point of introducing them is the possibility to describe Lie algebras via generators and relations.\\
Definition 7.2.3. Let $X$ be a set. A pair ( $L, \eta$ ) of a Lie algebra $L$ and a map $\eta: X \rightarrow L$ is said to be a free Lie algebra over $X$ if it has the following universal property. For each map $\alpha: X \rightarrow \mathfrak{g}$ into a Lie algebra $\mathfrak{g}$, there exists a unique morphism $\widetilde{\alpha}: L \rightarrow \mathfrak{g}$ of Lie algebras with $\widetilde{\alpha} \circ \eta=\alpha$.\\
Remark 7.2.4. (a) For two free Lie algebras ( $L_{1}, \eta_{1}$ ) and ( $L_{2}, \eta_{2}$ ) over $X$ there exists a unique isomorphism $\varphi: L_{1} \rightarrow L_{2}$ with $\varphi \circ \eta_{1}=\eta_{2}$. In fact, we choose $\varphi$ as the unique morphism of Lie algebras with $\varphi \circ \eta_{1}=\eta_{2}$. Then the unique morphism of Lie algebras $\psi: L_{2} \rightarrow L_{1}$ with $\psi \circ \eta_{2}=\eta_{1}$ satisfies $\psi \circ \varphi \circ \eta_{1}=\eta_{1}$, so that the uniqueness in the universal property implies $\psi \circ \varphi=\operatorname{id}_{L_{1}}$, and likewise $\varphi \circ \psi=\operatorname{id}_{L_{2}}$.\\
(b) If $(L, \eta)$ is a free Lie algebra over $X$, then $L$ is generated as a Lie algebra by $\eta(X)$. To see this, let $L_{0} \subseteq L$ be the Lie subalgebra generated by $\eta(X)$ and $\eta_{0}: X \rightarrow L_{0}$ be the corestriction of $\eta$. Then ( $L_{0}, \eta_{0}$ ) also has the universal property, so that (a) implies the existence of an isomorphism $\varphi: L_{0} \rightarrow L$ with $\varphi \circ \eta_{0}=\eta$. Since $L_{0}$ is generated by $\eta_{0}(X)=\eta(X)$, the Lie algebra $L$ is generated by $\eta(X)$, which leads to $L=L_{0}$.

Proposition 7.2.5. For each set $X$, there exists a free Lie algebra $(L(X), \eta)$ over $X$.

Proof. Let $X$ be a set. We put $X_{1}=X$ and define inductively

$$
X_{n}:=\bigcup_{p+q=n, p, q \in \mathbb{N}} X_{p} \times X_{q}
$$

The free magma $M(X)$ is defined as the disjoint union

$$
M(X):=\bigcup_{i=1}^{\infty} X_{n}
$$

Elements of the free magma are expressions of the form

$$
\left(\left(\left(\left(x_{1}, x_{2}\right), x_{3}\right), x_{4}\right)\right), \quad\left(\left(x_{1}, x_{2}\right),\left(x_{3}, x_{4}\right)\right), \quad\left(\left(x_{1}, x_{2}\right),\left(\left(x_{3}, x_{4}\right), x_{5}\right)\right) .
$$

One should think of it as a set containing all possible nonassociative products of elements of $X$. The maps $X_{n} \times X_{m} \rightarrow X_{n+m},(x, y) \mapsto x \cdot y:=(x, y)$ can be put together to a multiplication

$$
M(X) \times M(X) \rightarrow M(X), \quad(x, y) \mapsto x \cdot y
$$

The free vector space $A M(X):=\mathbb{K}[M(X)]$, with basis $M(X)$, thus inherits a bilinear multiplication extending the multiplication on $M(X)$.

Let $\mathfrak{n} \subseteq A M(X)$ denote the two sided ideal generated by the expressions of the form

$$
Q(x):=x \cdot x, \quad x \in A M(X)
$$

and

$$
J(x, y, z):=x \cdot(y \cdot z)+y \cdot(z \cdot x)+z \cdot(x \cdot y), \quad x, y, z \in M(X)
$$

We claim that the quotient algebra $L(X):=A M(X) / \mathfrak{n}$ is a free Lie algebra over $X$ with respect to the map $\eta: X \rightarrow L(X), x \mapsto x+\mathfrak{n}$. First, we note that the quotient space $L(X)$ inherits a bilinear multiplication turning it into a Lie algebra. We have to show that for each map $\varphi: X \rightarrow \mathfrak{g}$ into a Lie algebra $\mathfrak{g}$ there exists a unique Lie algebra homomorphism $L(\varphi): L(X) \rightarrow \mathfrak{g}$ with $L(\varphi)(x+\mathfrak{n})=\varphi(x)$ for each $x \in X$. In fact, one inductively defines mappings

$$
\varphi_{n}: X_{n} \rightarrow \mathfrak{g}, \quad(x, y) \mapsto\left[\varphi_{p}(x), \varphi_{q}(x)\right], \quad(x, y) \in X_{p} \times X_{q}, p+q=n
$$

and puts them together to a map $\varphi^{\prime}: M(X) \rightarrow \mathfrak{g}$ satisfying $\varphi(x \cdot y)= [\varphi(x), \varphi(y)]$. This map extends uniquely to a linear map $\varphi^{\prime \prime}: A M(X) \rightarrow \mathfrak{g}$ with the same property. Since $\mathfrak{g}$ is a Lie algebra, the map $\varphi^{\prime \prime}$ factors to a Lie algebra homomorphism $L(\varphi): L(X) \rightarrow \mathfrak{g}$ which is uniquely determined by $L(\varphi)(x+\mathfrak{n})=\varphi(x)$.

Definition 7.2.6. (cf. Exercise 7.2.1) (a) In the following, we denote a free Lie algebra over $X$ by $(L(X), \eta)$.\\
(b) Let $X$ be a set, $R \subseteq L(X)$ a subset and $I_{R} \subseteq L(X)$ the ideal generated by $R$. Then $L(X, R):=L(X) / I_{R}$ is called the Lie algebra defined by the generators $X$ and the relations $R$.

Example 7.2.7. (a) If $X=\{p, q, z\}$ and $R=\{[p, q]-z,[p, z],[q, z]\}$, then $L(X, R)$ is the 3-dimensional Heisenberg-Lie algebra. It is defined by the generators $p, q, z$ and the relations $[p, q]=z,[p, z]=[q, z]=0$.\\
(b) If $X=\{h, e, f\}$ and $R=\{[e, f]-h,[h, e]-2 e,[h, f]+2 f\}$, then $L(X, R) \cong \mathfrak{s l}_{2}(\mathbb{K})$.\\
(c) If $X=\{x\}$ is a one-element set, then $L(X)=\mathbb{C} x$ is one-dimensional. For a two-element set $X=\{x, y\}$, the free Lie algebra $L(X)$ is infinitedimensional (Exercise 7.2.2).

\subsubsection*{Serre's Theorem}
Now we return to the situation of Proposition 7.2.1. Let

$$
X:=\left\{\widehat{h}_{\alpha}, \widehat{e}_{\alpha}, \widehat{f}_{\beta}: \alpha \in \Pi\right\}
$$

and consider in the free Lie algebra $L(X)$ the relations

$$
\begin{aligned}
R & :=\left\{\left[\widehat{h}_{\alpha}, \widehat{h}_{\beta}\right],\left[\widehat{h}_{\alpha}, \widehat{e}_{\beta}\right]-\beta(\check{\alpha}) \widehat{e}_{\beta},\left[\widehat{h}_{\alpha}, \widehat{f}_{\beta}\right]+\beta(\check{\alpha}) \widehat{f}_{\beta},\left[\widehat{e}_{\alpha}, \widehat{f}_{\beta}\right]-\delta_{\alpha, \beta} \widehat{h}_{\alpha}\right\} \\
& \cup\left\{\left(\operatorname{ad} \widehat{e}_{\alpha}\right)^{1-\beta(\check{\alpha})} \widehat{e}_{\beta},\left(\operatorname{ad} \widehat{f}_{\alpha}\right)^{1-\beta(\check{\alpha})} \widehat{f}_{\beta}: \alpha \neq \beta \in \Pi\right\}
\end{aligned}
$$

Then Proposition 7.2.1(ii) implies the existence of a unique surjective Lie algebra homomorphism

$$
q: \widehat{\mathfrak{g}}:=L(X, R) \rightarrow \mathfrak{g} \quad \text { with } \quad q\left(\widehat{h}_{\alpha}\right)=h_{\alpha}, \quad q\left(\widehat{e}_{\alpha}\right)=e_{\alpha} \quad \text { and } \quad q\left(\widehat{f}_{\alpha}\right)=f_{\alpha} .
$$

The goal of this subsection is to prove Serre's Theorem that $q$ is an isomorphism, i.e., that the Lie algebra $\mathfrak{g}$ is defined by the generating set $X$ and the relations $R$. Let

$$
\widehat{\mathfrak{g}}_{+}:=\left\langle\left\{\widehat{e}_{\alpha}: \alpha \in \Pi\right\}\right\rangle_{\text {Lie alg }} \quad \text { and } \quad \widehat{\mathfrak{g}}_{-}:=\left\langle\left\{\widehat{f}_{\alpha}: \alpha \in \Pi\right\}\right\rangle_{\text {Lie alg }}
$$

denote the subalgebras of $\widehat{\mathfrak{g}}$ generated by the elements $\widehat{e}_{\alpha}$, resp., $\widehat{f}_{\alpha}$. Since the elements $h_{\alpha_{2}} \alpha \in \Pi$, are linearly independent in $\mathfrak{h}$, this also holds for the elements $\widehat{h}_{\alpha}$ in $\widehat{\mathfrak{h}}$, so that $\left.q\right|_{\widehat{\mathfrak{h}}}: \widehat{\mathfrak{h}} \rightarrow \mathfrak{h}$ is a linear isomorphism. For $\gamma \in \mathfrak{h}^{*}$, we put $\widehat{\gamma}:=\gamma \circ q \in \widehat{\mathfrak{h}}^{*}$.

Lemma 7.2.8. Let $\mathfrak{b}$ be a subalgebra of a Lie algebra $\mathfrak{g}$ and $M \subseteq \mathfrak{g}$ be a subset. Then $\mathfrak{b}_{M}:=\{x \in \mathfrak{b}:[M, x] \subseteq \mathfrak{b}\}$ is a subalgebra of $\mathfrak{b}$.

Proof. This follows from the following calculation for $m \in M$ and $x, y \in \mathfrak{b}_{M}$ :

$$
[m,[x, y]]=[[m, x], y]+[x,[m, y]] \subseteq[\mathfrak{b}, \mathfrak{b}]+[\mathfrak{b}, \mathfrak{b}] \subseteq \mathfrak{b}
$$

Lemma 7.2.9. The Lie algebra $\widehat{\mathfrak{g}}$ has the following properties:\\
(i) $\widehat{\mathfrak{g}}$ is the direct sum of $\widehat{\mathfrak{h}}$-weight spaces.\\
(ii) $\widehat{\mathfrak{g}}=\widehat{\mathfrak{g}}_{+} \oplus \widehat{\mathfrak{h}} \oplus \widehat{\mathfrak{g}}_{-}$is a direct sum of subalgebras and $\widehat{\mathfrak{h}}=\widehat{\mathfrak{g}}_{0}$. Moreover,

$$
\widehat{\mathfrak{g}}_{ \pm} \subseteq \sum_{\gamma \in \widehat{\Delta}_{ \pm}} \widehat{\mathfrak{g}}_{\gamma} \quad \text { for } \widehat{\Delta}_{ \pm}:= \pm \operatorname{span}_{\mathbb{N}} \widehat{\Pi}
$$

(iii) $\widehat{\mathfrak{g}}_{\widehat{\alpha}}=\mathbb{K} \widehat{e}_{\alpha}$ and $\widehat{\mathfrak{g}}_{-\widehat{\alpha}}=\mathbb{K} \widehat{f}_{\alpha}$ for $\alpha \in \Pi$.\\
(iv) The derivations ad $\widehat{e}_{\alpha}$ and ad $\widehat{f}_{\alpha}$ of $\widehat{\mathfrak{g}}$ are locally nilpotent, i.e., $\widehat{\mathfrak{g}}= \widehat{\mathfrak{g}}^{0}\left(\operatorname{ad} \widehat{e}_{\alpha}\right)=\widehat{\mathfrak{g}}^{0}\left(\operatorname{ad} \widehat{f}_{\alpha}\right)$ for each $\alpha \in \Pi$.\\
(v) The set $\widehat{\Delta}$ of roots of $\widehat{\mathfrak{g}}$ is invariant under the group $\widehat{\mathcal{W}}$ generated by the reflections defined by $\sigma_{\widehat{\alpha}} \widehat{\beta}:=\widehat{\beta}-\widehat{\beta}\left(\widehat{h}_{\alpha}\right) \widehat{\alpha}$. This is a finite group isomorphic to $\mathcal{W}$.\\
(vi) $\widehat{\Delta}=\widehat{\mathcal{W}} \widehat{\Pi}$.

Proof. (i) Let $\widehat{\mathfrak{g}}_{f}:=\sum_{\gamma \in \widehat{\mathfrak{h}}^{*}} \widehat{\mathfrak{g}}_{\gamma}$ be the sum of all $\widehat{\mathfrak{h}}$-weight spaces in $\widehat{\mathfrak{g}}$. In view of $\left[\widehat{\mathfrak{g}}_{\gamma}, \widehat{\mathfrak{g}}_{\delta}\right] \subseteq \widehat{\mathfrak{g}}_{\gamma+\delta}$, this is a subalgebra of $\widehat{\mathfrak{g}}$. We also have $\widehat{\mathfrak{h}} \subseteq \widehat{\mathfrak{g}}_{0}$ and

$$
\widehat{e}_{\alpha} \in \widehat{\mathfrak{g}}_{\widehat{\alpha}} \quad \text { and } \quad \widehat{f}_{\alpha} \in \widehat{\mathfrak{g}}_{-\widehat{\alpha}}
$$

and since $\widehat{\mathfrak{g}}$ is generated by these elements, it follows that $\widehat{\mathfrak{g}}=\widehat{\mathfrak{g}}_{f}=\sum_{\gamma \in \widehat{\mathfrak{h}}^{*}} \widehat{\mathfrak{g}}_{\gamma}$ is a direct sum of $\widehat{\mathfrak{h}}$-weight spaces. The directness of the sum of the weight spaces can be obtained from Lemma 6.1.3 and the observation that each finite subset of $\widehat{\mathfrak{g}}$ is contained in a finite $\widehat{\mathfrak{h}}$-invariant subspace, which follows from $\widehat{\mathfrak{g}}=\widehat{\mathfrak{g}}_{f}$.\\
(ii) Clearly, $\widehat{\mathfrak{h}}+\widehat{\mathfrak{g}}_{ \pm}$are subalgebras of $\widehat{\mathfrak{g}}$ because $\widehat{\mathfrak{g}}_{ \pm}$is generated by $\widehat{\mathfrak{h}}$-eigenvectors. For $\alpha \in \Pi$, we consider the subalgebra

$$
\mathfrak{b}_{\alpha}:=\left\{x \in \widehat{\mathfrak{h}}+\widehat{\mathfrak{g}}_{+}:\left[x, \widehat{f}_{\alpha}\right] \subseteq \widehat{\mathfrak{h}}+\widehat{\mathfrak{g}}_{+}\right\}
$$

(Lemma 7.2.8). For $\beta \in \Pi$, we have $\left[\widehat{f}_{\alpha}, \widehat{e}_{\beta}\right]=0$ for $\alpha \neq \beta$ and $\left[\widehat{f}_{\alpha} . \widehat{e}_{\alpha}\right] \subseteq \widehat{\mathfrak{h}}$. Hence the subalgebra $\mathfrak{b}_{\alpha}$ contains each $\widehat{e}_{\beta}$, and therefore $\widehat{\mathfrak{g}}_{+}$. From that we derive

$$
\left[\widehat{f}_{\alpha}, \widehat{\mathfrak{h}}+\widehat{\mathfrak{g}}_{+}\right] \subseteq \widehat{\mathfrak{h}}+\widehat{\mathfrak{g}}_{+}+\mathbb{K} \widehat{f}_{\alpha},
$$

which in turn implies that each $\widehat{f}_{\alpha}$ normalizes the subspace $\widehat{\mathfrak{g}}_{d}:=\widehat{\mathfrak{g}}_{-}+\widehat{\mathfrak{h}}+\widehat{\mathfrak{g}}_{+}$. Similarly, we obtain $e_{\alpha} \in \mathfrak{n}_{\widehat{\mathfrak{g}}}\left(\widehat{\mathfrak{g}}_{d}\right)$, so that $\widehat{\mathfrak{g}}_{d}$ is a subalgebra of $\widehat{\mathfrak{g}}$. Since it contains all generators, we obtain $\widehat{\mathfrak{g}}=\widehat{\mathfrak{g}}_{d}$.

Next we observe that the set of all $\widehat{\mathfrak{h}}$-weights in $\widehat{\mathfrak{g}}_{ \pm}$is contained in $\widehat{\Delta}_{ \pm}$, so that we obtain in particular that $\widehat{\mathfrak{g}}_{0}=\widehat{\mathfrak{h}}$.\\
(iii) Since $\widehat{\mathfrak{g}}_{+}$is generated by the elements $\widehat{e}_{\alpha} \in \widehat{\mathfrak{g}}_{\widehat{\alpha}}$ and $\widehat{\Pi}$ is linearly independent, the $\widehat{\alpha}$-weight space is spanned by $\widehat{e}_{\alpha}$. Similarly, $\widehat{\mathfrak{g}}_{\widehat{\beta}}=\mathbb{K} \widehat{f}_{\beta}$.\\
(iv) For any derivation $D \in \operatorname{der}(\widehat{\mathfrak{g}})$, the nilspace $\widehat{\mathfrak{g}}^{0}(D)$ is a subalgebra because we have

$$
D^{n}[a, b]=\sum_{k=0}^{n}\binom{n}{k}\left[D^{k} a, D^{n-k} b\right]
$$

For $D=\operatorname{ad} \widehat{e}_{\alpha}$, our relations imply that this subalgebra contains each $\widehat{h}_{\beta}$, each $\widehat{e}_{\beta}$, and also each $\widehat{f}_{\beta}$ because $\left(\operatorname{ad} \widehat{e}_{\alpha}\right)^{3} \widehat{f}_{\alpha}=0$. This leads to $\widehat{\mathfrak{g}}=\widehat{\mathfrak{g}}^{0}\left(\operatorname{ad} \widehat{e}_{\alpha}\right)$, and a similar argument shows that $\widehat{\mathfrak{g}}=\widehat{\mathfrak{g}}^{0}\left(\operatorname{ad} \widehat{f}_{\alpha}\right)$.\\
(v) Let $\widehat{\mathfrak{g}}(\alpha):=\operatorname{span}\left\{\widehat{e}_{\alpha}, \widehat{h}_{\alpha}, \widehat{f}_{\alpha}\right\}$ and note that our relations imply that $\widehat{\mathfrak{g}}(\alpha) \cong \mathfrak{s} \mathfrak{l}_{2}(\mathbb{K})$. Since ad $\widehat{e}_{\alpha}$ and ad $\widehat{f}_{\alpha}$ are locally nilpotent, for each $x \in \widehat{\mathfrak{g}}$, the $\widehat{\mathfrak{g}}(\alpha)$-submodule

$$
U(\widehat{\mathfrak{g}}(\alpha)) x=\sum_{\ell, m, n}\left(\operatorname{ad} \widehat{f}_{\alpha}\right)^{\ell}\left(\operatorname{ad} \widehat{h}_{\alpha}\right)^{m}\left(\operatorname{ad} \widehat{e}_{\alpha}\right)^{n} x
$$

generated by $x$ is finite-dimensional (here we use the Poincaré-Birkhoff-Witt Theorem 7.1.9 below). We therefore obtain a well-defined linear operator

$$
\widehat{\sigma}_{\alpha}:=e^{\widehat{\mathrm{ad}} e_{\alpha}} e^{-\widehat{\mathrm{ad}} f_{\alpha}} e^{\widehat{\mathrm{ad}} e_{\alpha}} \in \operatorname{GL}(\widehat{\mathfrak{g}})
$$

which clearly commutes with $\operatorname{ker} \widehat{\alpha} \subseteq \widehat{\mathfrak{h}}$ and satisfies

$$
\widehat{\sigma}_{\alpha}\left(\widehat{\mathfrak{g}}_{\lambda}\left(\widehat{h}_{\alpha}\right)\right)=\widehat{\mathfrak{g}}_{-\lambda}\left(\widehat{h}_{\alpha}\right)
$$

(cf. Lemma 6.2.9). Next we observe that $\sigma_{\widehat{\alpha}}$ fixes $\widehat{h}_{\alpha}^{\perp} \subseteq \widehat{\mathfrak{h}}^{*}$ pointwise and satisfies $\sigma_{\widehat{\alpha}} \widehat{\alpha}=-\widehat{\alpha}$. We therefore obtain $\widehat{\sigma}_{\alpha}\left(\widehat{\mathfrak{g}}_{\gamma}\right)=\widehat{\mathfrak{g}}_{\sigma_{\alpha} \gamma}$ for each $\gamma \in \widehat{\Delta}$. Clearly, $q \circ \sigma_{\widehat{\alpha}}=\sigma_{\alpha} \circ q$, and since $\left.q\right|_{\widehat{h}}$ is a linear isomorphism to $\mathfrak{h}$, we see that $\widehat{\mathcal{W}}$ is a finite group isomorphic to $\mathcal{W}$.\\
(vi) Since $-\widehat{\Pi} \subseteq \widehat{\mathcal{W}} \widehat{\Pi}$, it suffices to show that each positive root $\gamma \in \widehat{\Delta}_{+}$ lies in $\widehat{\mathcal{W}} \widehat{\Pi}$. For $\gamma=\sum_{\alpha \in \Pi} n_{\alpha} \widehat{\alpha} \in \widehat{\Delta}_{+}$we define its height by

$$
\operatorname{ht}(\gamma):=\sum_{\alpha \in \Pi} n_{\alpha}
$$

Pick $\delta \in \widehat{\mathcal{W}} \gamma \cap \widehat{\Delta}_{+}$with minimal height. If $\delta \notin \widehat{\Pi}$, then $\delta=\sum_{\alpha \in \Pi} n_{\alpha} \widehat{\alpha}$ with $n_{\alpha_{1}}, n_{\alpha_{2}}>0$ for two roots $\alpha_{1} \neq \alpha_{2} \in \Pi$. Then

$$
\sigma_{\widehat{\beta}} \delta=\delta-\delta\left(\widehat{h}_{\beta}\right) \widehat{\beta}=\left(n_{\beta}-\delta\left(\widehat{h}_{\beta}\right)\right) \widehat{\beta}+\sum_{\alpha \in \Pi \backslash \beta} n_{\alpha} \widehat{\alpha} \in \widehat{\Delta}_{+}
$$

because $n_{\alpha_{j}}>0$ for $j=1,2$ and one of these two roots is different from $\beta$. Therefore, the minimality of the height of $\delta$ implies $\delta\left(\widehat{h}_{\beta}\right) \leq 0$ for each\\
$\beta \in \Pi$. For the corresponding linear functional $\bar{\delta} \in \mathfrak{h}^{*}$ with $\delta=\bar{\delta} \circ q$, we then have $\bar{\delta} \in \sum_{\alpha \in \Pi} \mathbb{N}_{0} \alpha$ and $\bar{\delta}(\check{\alpha}) \leq 0$ for each $\alpha \in \Pi$. This leads to

$$
(\bar{\delta}, \bar{\delta})=\sum_{\alpha \in \Pi} n_{\alpha}(\bar{\delta}, \alpha)=\sum_{\alpha \in \Pi} n_{\alpha} \bar{\delta}\left(h_{\alpha}^{\prime}\right) \leq 0
$$

because $h_{\alpha}^{\prime}$ is a positive multiple of $h_{\alpha}$ (Theorem 6.3.4). As the scalar product on $\mathfrak{h}^{*}$ is positive definite, we arrive at a contradiction. This proves that $\delta \in \widehat{\Pi}$, and hence that $\gamma \in \widehat{\mathcal{W}} \widehat{\Pi}$.

Theorem 7.2.10 (Serre's Theorem). The homomorphism

$$
q: \widehat{\mathfrak{g}}=L(X, R) \rightarrow \mathfrak{g}
$$

is an isomorphism of Lie algebras.\\
Proof. We know already that $q$ is surjective, so that it remains to see that it is also injective. Let $\mathfrak{n}:=\operatorname{ker} q$. Then $\mathfrak{n}$ is an ideal of $\widehat{\mathfrak{g}}$, hence in particular invariant under $\widehat{\mathfrak{h}}$, which in turn implies that it is adapted to the $\widehat{\mathfrak{h}}$-weight decomposition of $\widehat{\mathfrak{g}}$ :

$$
\mathfrak{n}=(\mathfrak{n} \cap \widehat{\mathfrak{h}}) \oplus \sum_{\gamma \in \widehat{\Delta}}\left(\mathfrak{n} \cap \widehat{\mathfrak{g}}_{\gamma}\right)
$$

Since $q$ is injective on $\widehat{\mathfrak{h}}+\sum_{\alpha \in \Pi}\left(\widehat{\mathfrak{g}}_{\widehat{\alpha}}+\widehat{\mathfrak{g}}_{-\widehat{\alpha}}\right)$ (Lemma 7.2.9(iii)), we have $\mathfrak{n} \cap \widehat{\mathfrak{g}}_{\gamma}=\{0\}$ for $\gamma \in\{0\} \cup \pm \widehat{\Pi}$.

On the other hand, the invariance of the ideal $\mathfrak{n}$ under $\widehat{\mathfrak{g}}(\alpha)$ implies that it is also invariant under the linear maps $\widehat{\sigma}_{\alpha} \in \mathrm{GL}(\widehat{\mathfrak{g}})$, defined in (7.3), so that

$$
\left\{\gamma \in \widehat{\Delta}: \mathfrak{n} \cap \widehat{\mathfrak{g}}_{\gamma} \neq\{0\}\right\}
$$

is invariant under $\widehat{\mathcal{W}}$ (Lemma 7.2.9). Now $\widehat{\Delta}=\widehat{\mathcal{W}} \widehat{\Pi}$ implies that this set is empty, and this proves that $\mathfrak{n}=\{0\}$.

Corollary 7.2.11. There exists an automorphism $\varphi: \mathfrak{g} \rightarrow \mathfrak{g}$ with $\left.\varphi\right|_{\mathfrak{h}}= -\mathrm{id}_{\mathfrak{h}}$.

Proof. Let $X$ be as above and define the map

$$
\widehat{\varphi}: X \rightarrow \mathfrak{g}, \quad \widehat{\varphi}\left(\widehat{h}_{\alpha}\right):=-h_{\alpha}, \quad \widehat{\varphi}\left(\widehat{e}_{\alpha}\right):=-f_{\alpha}, \quad \widehat{\varphi}\left(\widehat{f}_{\alpha}\right):=-e_{\alpha}
$$

which defines a unique homomorphism $L(X) \rightarrow \mathfrak{g}$, and the Serre relations (7.1) and (7.2) imply that this homomorphism actually factors through a homomorphism $\widetilde{\varphi}: L(X, R) \rightarrow \mathfrak{g}$. Now $\varphi:=\widetilde{\varphi} \circ q^{-1}: \mathfrak{g} \rightarrow \mathfrak{g}$ is an automorphism of $\mathfrak{g}$ with $\left.\varphi\right|_{\mathfrak{h}}=-\mathrm{id}_{\mathfrak{h}}$.

\subsubsection*{Exercises for Section 7.2}
Exercise 7.2.1. Let $\mathfrak{g}$ be a Lie algebra and $\left(b_{j}\right)_{j \in J}$ a basis for $\mathfrak{g}$. Then the Lie bracket on $\mathfrak{g}$ is determined by the numbers $c_{i j}^{k}$ (called structure constants), defined by

$$
\left[b_{i}, b_{j}\right]=\sum_{k \in J} c_{i j}^{k} b_{k}
$$

Show that $\mathfrak{g} \cong L(B, R)$, where

$$
B=\left\{b_{j}: j \in J\right\} \quad \text { and } \quad R=\left\{\left[b_{i}, b_{j}\right]-\sum_{k \in J} c_{i j}^{k} b_{k}: i, j \in J\right\}
$$

Exercise 7.2.2. Let $X=\{x, y\}$ be a two element set. Show that the free Lie algebra $L(X)$ is infinite-dimensional. Hint: Consider the semidirect product $\mathfrak{g}=\mathbb{K}[X] \rtimes_{M} \mathbb{K}$, where $\mathbb{K}[X]$ is considered as an abelian Lie algebra and $M f(X):=X f(X)$ is the multiplication with $X$.

\subsection*{Highest Weight Representations}
We know already from Weyl's Theorem 5.5.21 on Complete Reducibility that any finite-dimensional module over a semisimple Lie algebra $\mathfrak{g}$ is semisimple. This reduces the classification of finite-dimensional modules to the classification of simple ones. In this section, we address this problem for the class of those semisimple Lie algebras which are split, i.e., contain a toral Cartan subalgebra. Note that $\mathfrak{s} \mathfrak{l}_{n}(\mathbb{K})$ and in particular any complex semisimple Lie algebra is split.

Throughout this section, $\mathfrak{g}$ denotes a semisimple Lie algebra and $\mathfrak{h} \subseteq \mathfrak{g}$ a toral Cartan subalgebra. For $\lambda \in \mathfrak{h}^{*}$ and a representation ( $\pi, V$ ) of $\mathfrak{g}$, we write

$$
V_{\lambda}:=\{v \in V:(\forall h \in \mathfrak{h}) \pi(h)(v)=\lambda(h) v\}
$$

for the corresponding weight space in $V$ and $\mathcal{P}(V):=\left\{\lambda \in \mathfrak{h}^{*}: V_{\lambda} \neq\{0\}\right\}$ for the set of $\mathfrak{h}$-weights of $V$. We simply write $\Delta:=\Delta(\mathfrak{g}, \mathfrak{h})$ for the set of roots and $\mathfrak{g}_{\alpha}:=\mathfrak{g}_{\alpha}(\mathfrak{h}), \alpha \in \Delta$, for the root spaces.

Proposition 7.3.1. If $\operatorname{dim} V<\infty$, then $\mathfrak{h}$ acts by diagonalizable operators on $V$, and $V$ is the direct sum of its weight spaces. All weights take real values on $\mathfrak{h}_{\mathbb{R}}$.

Proof. In view of Lemma 6.3.7, $\mathfrak{h}$ is spanned by the coroots $h_{\alpha}=\check{\alpha}$. Therefore, it suffices to see that for each root $\alpha \in \Delta$, the element $h_{\alpha} \in \mathfrak{h}$ is diagonalizable on $V$. Since the $\mathfrak{g}$-representation on $V$ restricts to a representation of

$$
\mathfrak{g}(\alpha)=\mathfrak{g}_{\alpha}+\mathfrak{g}_{-\alpha}+\mathbb{K} h_{\alpha} \cong \mathfrak{s} \mathfrak{l}_{2}(\mathbb{K})
$$

on $V$, it suffices to apply Proposition 6.2.8. Moreover, we see that all eigenvalues of $h_{\alpha}$ are integral, which, since $\mathfrak{h}$ is abelian, implies that each weight takes only real values on the real subspace $\mathfrak{h}_{\mathbb{R}}$.

Definition 7.3.2. Let $\Delta^{+} \subseteq \Delta$ be a positive system of roots. (cf. Theorem 6.4.13). Then

$$
\mathfrak{b}:=\mathfrak{h}+\sum_{\beta \in \Delta^{+}} \mathfrak{g}_{\beta}
$$

is a solvable subalgebra of $\mathfrak{g}$ because it is of the form $\mathfrak{b}=\mathfrak{n} \rtimes \mathfrak{h}$ for a nilpotent Lie algebra $\mathfrak{n}$ (Proposition 7.2.1). Subalgebras of this type are called standard Borel subalgebras with respect to $\mathfrak{h}$.

\subsubsection*{Highest Weights}
Definition 7.3.3. A $\mathfrak{g}$-module $V$ is called a module with highest weight $\lambda \in \mathfrak{h}^{*}$ if there is a $\mathfrak{b}$-invariant line $\mathbb{K} v \in V$ with

$$
h \cdot v=\lambda(h) v \quad \text { for } \quad h \in \mathfrak{h}
$$

and $v$ generates the $\mathfrak{g}$-module $V$ (i.e., $V$ is the smallest submodule containing $v$ ). Then $\lambda$ is called the highest weight and the nonzero elements of the generating line $\mathbb{K} v$ are called highest weight vectors.

Proposition 7.3.4. Each finite-dimensional simple $\mathfrak{g}$-module is a highest weight module.

Proof. Let $V$ be a simple $\mathfrak{g}$-module. In view of Proposition 7.3.1, $V$ is a direct sum of its weight spaces $V=\bigoplus_{\alpha \in \mathcal{P}(V)} V_{\alpha}$. Since $V$ is finite-dimensional, the set $\mathcal{P}(V)$ of weights is finite.

Pick $x_{0} \in \mathfrak{h}_{\mathbb{R}}$ with

$$
\Delta^{+}=\left\{\alpha \in \Delta: \alpha\left(x_{0}\right)>0\right\} .
$$

Then $\mathcal{P}(V)\left(x_{0}\right) \subseteq \mathbb{R}$ (Proposition 7.3.1) and we can pick a $\lambda \in \mathcal{P}(V)$ such that $\lambda\left(x_{0}\right)$ is maximal. Let $v \in V_{\lambda} \backslash\{0\}$. For each $\alpha \in \Delta^{+}$, we then have $\mathfrak{g}_{\alpha} \cdot v \subseteq V_{\lambda+\alpha}$, but the choice of $\lambda$ implies that $V_{\lambda+\alpha}=\{0\}$. Hence $v$ is a $\mathfrak{b}$-eigenvector of weight $\lambda$. Since $V$ is simple, it is generated by $v$.

Remark 7.3.5. If $\mathbb{K}=\mathbb{C}$, then we can also use Lie's Theorem 5.4.8 to see that a simple $\mathfrak{g}$-module $V$ contains a $\mathfrak{b}$-eigenvector, hence is a highest weight module.

For $\beta \in \Delta^{+}=\left\{\beta_{1}, \ldots, \beta_{m}\right\}$, we choose an $\mathfrak{s} \mathfrak{l}_{2}$-triple ( $h_{\beta}, e_{\beta}, f_{\beta}$ ) as in the $\mathfrak{s l}_{2}$-Theorem 6.3.4. As for abstract root systems (see p. 161), we define a partial order $\prec$ on $\mathfrak{h}^{*}$ by

$$
\lambda \prec \mu \quad: \Longleftrightarrow \quad \mu-\lambda \in \mathbb{N}_{0}\left[\Delta^{+}\right]:=\sum_{\beta \in \Delta^{+}} \mathbb{N}_{0} \beta .
$$

Let $\Pi \subseteq \Delta^{+}$be the corresponding set of simple roots (Theorem 6.4.13).\\
The following theorem describes some properties of highest weight modules which are not necessarily finite-dimensional.

Theorem 7.3.6. Let $V$ be a $\mathfrak{g}$-module with highest weight $\lambda$ and $0 \neq v \in V_{\lambda}$ a highest weight vector. Then\\
(i) $\quad V=\operatorname{span}\left\{f_{\beta_{1}}^{i_{1}} \cdots f_{\beta_{m}}^{i_{m}} \cdot v \mid i_{j} \in \mathbb{N}_{0}\right\}$. In particular, $V$ is the direct sum of its weight spaces.\\
(ii) $\mathcal{P}(V) \subseteq \lambda-\mathbb{N}_{0}\left[\Delta^{+}\right]=\lambda-\sum_{\beta \in \Pi} \mathbb{N}_{0} \beta$.\\
(iii) $\operatorname{dim} V_{\mu}<\infty$ for all $\mu \in \mathcal{P}(V)$.\\
(iv) $\operatorname{dim} V_{\lambda}=1$.\\
(v) Every $\mathfrak{g}$-submodule of $V$ is the direct sum of its weight spaces.\\
(vi) $V$ contains exactly one maximal proper $\mathfrak{g}$-submodule $V_{\max }$ and $V / V_{\max }$ is the unique simple quotient module of $V$.\\
(vii) Every nonzero module quotient of $V$ is a module with highest weight $\lambda$.

Proof. (i) Let $\Pi=\left\{\alpha_{1}, \ldots, \alpha_{r}\right\}$. Then

$$
f_{\beta_{1}}, \ldots, f_{\beta_{m}}, h_{\alpha_{1}}, \ldots, h_{\alpha_{r}}, x_{\beta_{1}}, \ldots, x_{\beta_{m}}
$$

is a basis for $\mathfrak{g}$, to which we apply the Poincaré-Birkhoff-Witt Theorem 7.1.9. Then the claim follows from

$$
h_{\alpha_{1}}^{j_{1}} \cdots h_{\alpha_{r}}^{j_{r}} x_{\beta_{1}}^{\ell_{1}} \cdots x_{\beta_{m}}^{\ell_{m}} \cdot v \subseteq \mathbb{K} v
$$

and $V=\mathcal{U}(\mathfrak{g}) \cdot v($ cf. Remark 7.1.4).\\
(ii) In view of (i), $V$ is spanned by vectors of the form

$$
f_{\beta_{1}}^{i_{1}} \cdots f_{\beta_{m}}^{i_{m}} \cdot v
$$

which are weight vectors of weight $\lambda-\sum_{\ell=1}^{m} i_{\ell} \beta_{\ell}$ by Lemma 6.2.3. Assertion (ii) now follows, since the positive roots can be written as sums of simple roots.\\
(iii) For $\mu \in \mathfrak{h}^{*}$, there are only finitely many vectors of the form $f_{\beta_{1}}^{i_{1}} \cdots f_{\beta_{m}}^{i_{m}}$. $v$ for which $\lambda-\sum_{\ell=1}^{m} i_{\ell} \beta_{\ell}$ equals $\mu$. In fact, if $x_{0} \in \mathfrak{h}_{\mathbb{R}}$ satisfies $\beta_{j}\left(x_{0}\right)>0$ for each $j$, then $\mu=\lambda-\sum_{l=1}^{m} i_{\ell} \beta_{\ell}$ yields

$$
\sum_{\ell=1}^{m} i_{\ell} \beta_{\ell}\left(x_{0}\right)=(\lambda-\mu)\left(x_{0}\right)
$$

and there are only finitely many solutions $\left(i_{1}, \ldots, i_{m}\right) \in \mathbb{N}_{0}^{m}$ of this equation.\\
(iv) The equality $\lambda=\lambda-\sum_{l=1}^{m} i_{l} \beta_{l}$ is only possible for $i_{1}=\cdots=i_{m}=0$.\\
(v) Let $W$ be a $\mathfrak{g}$-submodule of $V$. Then Exercise 2.1.1(b) implies that each element of $\mathfrak{h}$ is diagonalizable on $W$, so that $\mathfrak{h}$ is simultaneously diagonalizable on $W$.\\
(vi) By (v), every proper submodule of $V$ is contained in $\sum_{\mu \in \mathcal{P}(V) \backslash\{\lambda\}} V_{\mu}$. Therefore, the union of all proper submodules is still proper, and hence, a maximal proper submodule $V_{\text {max }}$ exists. The quotient module $V / V_{\text {max }}$ is simple since for every nontrivial submodule $W \subseteq V / V_{\text {max }}$ its inverse image $W^{\prime}$\\
in $V$ would be a proper submodule of $V$, strictly containing $V_{\text {max }}$. Conversely, every submodule $W$ of $V$ for which $V / W$ is simple (and nonzero) is a maximal submodule, hence equal to $V_{\text {max }}$.\\
(vii) This is obvious.

Corollary 7.3.7. If $V$ is a simple highest weight module, then $V$ contains only one $\mathfrak{b}$-invariant line.

Proof. Let $\mathbb{K} v$ be a $\mathfrak{b}$-invariant line. Then $v$ is a weight vector for some weight $\mu$ and $v$ generates the simple module $V$ (each simple module is generated by each nonzero element). Hence $\mathcal{P}(V) \subseteq \mu-\mathbb{N}_{0}\left[\Delta^{+}\right]$. If $\lambda$ is the highest weight of $V$, we also have $\mathcal{P}(V) \subseteq \lambda-\mathbb{N}_{0}\left[\Delta^{+}\right]$, which leads to

$$
\lambda \prec \mu \prec \lambda,
$$

and hence to $\lambda=\mu$. Finally, Theorem 7.3.6(iv) implies that $V_{\lambda}$ is onedimensional, which completes the proof.

Proposition 7.3.8. Two simple $\mathfrak{g}$-modules with the same highest weight $\lambda$ are isomorphic.

Proof. Let $V$ and $W$ be two such modules. We choose nonzero elements $v \in V_{\lambda}$ and $w \in W_{\lambda}$. Set $M:=V \oplus W$ and $m:=v+w$. Then $\mathbb{K} m$ is a $\mathfrak{b}$-invariant line and the submodule $M^{\prime}:=\mathcal{U}(\mathfrak{g}) \cdot m$ of $M$ generated by $m$ is a module with highest weight $\lambda$. Let $\operatorname{pr}_{V}: M^{\prime} \rightarrow V$ and $\operatorname{pr}_{W}: M^{\prime} \rightarrow W$ be the canonical projections with respect to the direct sum $V \oplus W$. Then both $\operatorname{pr}_{V}$ and $\operatorname{pr}_{W}$ are homomorphisms of $\mathfrak{g}$-modules. From $\operatorname{pr}_{V}(m)=v$ and $\operatorname{pr}_{W}(m)=w$, we see that $\operatorname{pr}_{V}$ and $\operatorname{pr}_{W}$ are surjective. Therefore, we must have ker $\operatorname{pr}_{V}=M_{\max }^{\prime}=\operatorname{ker}_{W} \operatorname{pr}_{W}$ by Theorem 7.3.6(vi), and this implies $V \cong M^{\prime} / M_{\max }^{\prime} \cong W$.

Definition 7.3.9 (Verma modules). Let $\mathfrak{g}$ be a semisimple Lie algebra, $\mathfrak{h} \subseteq \mathfrak{g}$ a splitting Cartan subalgebra, and $\mathfrak{b}=\mathfrak{h}+\sum_{\beta \in \Delta^{+}} \mathfrak{g}_{\beta}$ the Borel subalgebra of $\mathfrak{g}$ corresponding to the positive system $\Delta^{+}$of $\Delta$. For $\lambda \in \mathfrak{h}^{*}$, we extend $\lambda$ to a linear functional, also called $\lambda$, on $\mathfrak{b}$ vanishing on all root spaces $\mathfrak{g}_{\alpha}$. Then

$$
[\mathfrak{b}, \mathfrak{b}]=\sum_{\beta \in \Delta^{+}} \mathfrak{g}_{\beta} \subseteq \operatorname{ker} \lambda
$$

implies that $\lambda: \mathfrak{b} \rightarrow \mathbb{K} \cong \mathfrak{g} \mathfrak{l}_{1}(\mathbb{K})$ is a homomorphism of Lie algebras, hence defines a one-dimensional representation of $\mathfrak{b}$. We write $\mathbb{K}_{\lambda}$ for the corresponding $\mathfrak{b}$-module. The Lie algebra homomorphism $\lambda$ further extends to an algebra homomorphism $\widetilde{\lambda}: \mathcal{U}(\mathfrak{b}) \rightarrow \mathbb{K}$, turning $\mathbb{K}_{\lambda}$ into a $\mathcal{U}(\mathfrak{b})$-module.

In the following, we shall use the notation

$$
A B:=\operatorname{span}\{a b: a \in A, b \in B\}
$$

for subsets $A, B$ of an associative algebra. In this sense, we write

$$
M(\lambda):=\mathcal{U}(\mathfrak{g}) /(\mathcal{U}(\mathfrak{g})\{b-\lambda(b) \mathbf{1}: b \in \mathfrak{b}\}) .
$$

The module $M(\lambda)$ is called the Verma module of highest weight $\lambda$. We write $[D], D \in \mathcal{U}(\mathfrak{g})$, for its elements. Since $M(\lambda)$ is a quotient by a subspace invariant under the natural (left multiplication) action of $\mathfrak{g}$ on $\mathcal{U}(\mathfrak{g})$, it carries a natural $\mathfrak{g}$-module structure.

It is indeed a highest weight module of highest weight $\lambda$ because $v:=[\mathbf{1}]$ satisfies for $b \in \mathfrak{b}$ :

$$
b \cdot[\mathbf{1}]=[b]=[\lambda(b) \mathbf{1}]=\lambda(b)[\mathbf{1}]
$$

and

$$
\mathcal{U}(\mathfrak{g}) \cdot[\mathbf{1}]=[\mathcal{U}(\mathfrak{g})]=M(\lambda) .
$$

Using Theorem 7.3.6, we conclude that $L(\lambda):=M(\lambda) / M(\lambda)_{\text {max }}$ is a simple highest weight module with highest weight $\lambda$. In particular, such a highest weight module exists for each $\lambda \in \mathfrak{h}^{*}$.

\subsubsection*{Classification of Finite-Dimensional Simple Modules}
We want to characterize the linear functionals on $\mathfrak{h}$ which occur as highest weights of simple $\mathfrak{g}$-modules.

Definition 7.3.10. A linear functional $\lambda \in \mathfrak{h}^{*}$ is said to be integral if

$$
\lambda(\check{\alpha}) \in \mathbb{Z} \quad \text { for } \alpha \in \Delta,
$$

and it is called dominant with respect to the positive system $\Delta^{+}$if

$$
\lambda(\check{\alpha}) \geq 0 \quad \text { for } \alpha \in \Delta^{+} .
$$

We denote the set of all integral functionals on $\mathfrak{h}$ by $\Lambda$, and the set of all dominant integral functionals by $\Lambda^{+}$.

Let $\Pi \subseteq \Delta^{+}$be the set of simple roots, which is a basis for $\mathfrak{h}^{*}$. Since the set $\check{\Pi}=\{\check{\alpha}: \alpha \in \Pi\}$ is a basis for the dual root system $\check{\Delta} \subseteq \mathfrak{h}_{\mathbb{R}}$ (Remark 6.4.18),

$$
\Lambda=\left\{\lambda \in \mathfrak{h}^{*}:(\forall \alpha \in \Pi) \lambda(\check{\alpha}) \in \mathbb{Z}\right\} \cong \mathbb{Z}^{r}
$$

is a lattice in the real vector space $\mathfrak{h}_{\mathbb{R}}^{*}$, called the weight lattice and

$$
\Lambda^{+}=\left\{\lambda \in \mathfrak{h}^{*}:(\forall \alpha \in \Pi) \lambda(\check{\alpha}) \in \mathbb{N}_{0}\right\} \cong \mathbb{N}_{0}^{r} .
$$

Remark 7.3.11. The Weyl group of the root system $\Delta$ is the subgroup $W \subseteq \mathrm{GL}\left(\mathfrak{h}_{\mathbb{R}}^{*}\right)$ generated by the reflections

$$
\sigma_{\alpha}(\lambda)=\lambda-\lambda(\check{\alpha}) \alpha
$$

and this formula immediately implies that the weight lattice $\Lambda$ is invariant under $W$ (Theorem 6.4.17). Moreover, Theorem 6.4.17 implies that for each integral element $\nu \in \mathfrak{h}^{*}$, there exists a $w \in W$ with $w(\nu) \in \Lambda^{+}$.

Proposition 7.3.12. If $V$ is a finite-dimensional module with highest weight $\lambda$, then $\lambda$ is dominant integral.

Proof. Let $\alpha \in \Delta^{+}$and $\mathfrak{g}(\alpha) \cong \mathfrak{s l}_{2}(\mathbb{K})$ be the corresponding 3-dimensional subalgebra of $\mathfrak{g}$. If $v \in V$ is a highest weight vector, then we obtain with Proposition 6.2.4 that $\lambda\left(h_{\alpha}\right)=\lambda(\check{\alpha}) \in \mathbb{N}_{0}$.

Lemma 7.3.13. Let $V$ be a $\mathfrak{g}$-module and $W=W(\Delta)$ the Weyl group of $\Delta$. If, for each $\alpha \in \Pi, V$ is a locally finite $\mathfrak{g}(\alpha)$-module, i.e., a union of finite-dimensional submodules, then each weight $\mu$ of $V$ satisfies

$$
\operatorname{dim} V_{\mu}=\operatorname{dim} V_{w(\mu)}
$$

Proof. Let $\rho: \mathfrak{g} \rightarrow \mathfrak{g} \mathfrak{l}(V)$ denote the representation of $\mathfrak{g}$ on $V$ and ( $h_{\alpha}, e_{\alpha}, f_{\alpha}$ ) an $\mathfrak{s l}_{2}$-triple corresponding to $\alpha \in \Pi$. According to our hypothesis, each $v \in V$ is contained in a finite-dimensional submodule, so that

$$
\Phi_{\alpha} v:=e^{\rho\left(e_{\alpha}\right)} \circ e^{-\rho\left(f_{\alpha}\right)} \circ e^{\rho\left(e_{\alpha}\right)} v
$$

is defined. Moreover, its restriction to each finite-dimensional $\mathfrak{g}(\alpha)$-submodule is invertible, and this implies that $\Phi_{\alpha}$ actually is an element of $\mathrm{GL}(V)$. Let $\mu \in \mathfrak{h}^{*}$ be a weight of $V$. From Lemma 6.2.9, we know that $\Phi_{\alpha}$ maps the $h_{\alpha}$-eigenspace for the eigenvalue $\mu\left(h_{\alpha}\right)$ to the $h_{\alpha}$-eigenspace for the eigenvalue $-\mu\left(h_{\alpha}\right)$. If $h \in \mathfrak{h}$ satisfies $\alpha(h)=0$, then $\rho(h)$ commutes with $\Phi_{\alpha}$, so that $\Phi_{\alpha}$ preserves the $h$-eigenspaces. This implies that

$$
\Phi_{\alpha}\left(V_{\mu}\right)=V_{\sigma_{\alpha}(\mu)}
$$

where $\sigma_{\alpha}(\mu)=\mu-\mu(\check{\alpha}) \alpha$ is the reflection in $W$ corresponding to $\alpha$. Composing the various $\Phi_{\alpha}$ for $\alpha \in \Pi$, we see that the set of weights of $V$ is invariant under the Weyl group $W$ which is generated by $\left\{\sigma_{\alpha}: \alpha \in \Pi\right\}$ (Theorem 6.4.17). Moreover, it is now clear that $\operatorname{dim} V_{\mu}=\operatorname{dim} V_{w(\mu)}$ for all $w \in W$.

Proposition 7.3.14. Let $\mathfrak{g}$ be a semisimple Lie algebra and $\mathfrak{h}$ be a splitting Cartan subalgebra of $\mathfrak{g}$. If $\lambda \in \mathfrak{h}^{*}$ is dominant integral with respect to $\Delta^{+}$, then the simple highest weight module $V=L(\lambda)$ of highest weight $\lambda$ is finitedimensional.

Proof. For each simple root $\alpha \in \Pi$, using Theorem 6.3.4, we choose an $\mathfrak{s} \mathfrak{l}_{2}$-triple $\left(h_{\alpha}, e_{\alpha}, f_{\alpha}\right)$, so that $\mathfrak{g}(\alpha)=\operatorname{span}\left\{h_{\alpha}, e_{\alpha}, f_{\alpha}\right\} \cong \mathfrak{s} \mathfrak{l}_{2}(\mathbb{K})$. Let $v^{+} \in V$ be a highest weight vector. For $\alpha, \beta \in \Pi$, we put $m_{\alpha}:=\lambda\left(h_{\alpha}\right)=\lambda(\check{\alpha}) \in \mathbb{N}_{0}$ and observe that

$$
e_{\beta} f_{\alpha}^{m_{\alpha}+1} v^{+}= \begin{cases}0 & \text { if } \alpha \neq \beta \\ \left(m_{\alpha}+1\right) f_{\alpha}^{m_{\alpha}}\left(h_{\alpha}-m_{\alpha} \mathbf{1}\right) v^{+}=0 & \text { if } \alpha=\beta\end{cases}
$$

Here we use that $\left[e_{\beta}, f_{\alpha}\right] \in \mathfrak{g}_{\beta-\alpha}=\{0\}$ for $\alpha \neq \beta$, and for $\alpha=\beta$ we use the formulas from Lemma 6.2.2(ii) and $e_{\beta} \cdot v^{+}=0$. Since the $e_{\alpha}, \alpha \in \Pi$, generate the subalgebra $\sum_{\beta \in \Delta^{+}} \mathfrak{g}_{\beta}$ (Proposition 7.2.1), the vector $f_{\alpha}^{m_{\alpha}+1} v^{+}$ generates a proper highest weight submodule of $V$ (cf. Theorem 7.3.6 and Corollary 7.3.7), so that the irreducibility of $V$ yields $f_{\alpha}^{m_{\alpha}+1} v^{+}=0$. Therefore,

$$
\operatorname{span}\left\{v^{+}, f_{\alpha} v^{+}, \ldots, f_{\alpha}^{m_{\alpha}} v^{+}\right\}
$$

is a finite-dimensional $\mathfrak{g}(\alpha)$-module with highest weight $\lambda(\check{\alpha})=m_{\alpha}$. Therefore, the subspace $V_{\alpha}^{\prime}$ of $V$ spanned by all finite-dimensional $\mathfrak{g}(\alpha)$-submodules is nonzero.

Let $E \subseteq V$ be a finite-dimensional $\mathfrak{g}(\alpha)$-submodule. Then $\operatorname{span}(\mathfrak{g} \cdot E)$ is finite-dimensional and $\mathfrak{g}(\alpha)$-stable, so it is contained in $V^{\prime}$. This implies that $V_{\alpha}^{\prime}$ is a $\mathfrak{g}$-submodule. Since $V$ is simple, we conclude that $V=V_{\alpha}^{\prime}$. Thus Lemma 7.3.13 is applicable and yields $\operatorname{dim} V_{w \cdot \mu}=\operatorname{dim} V_{\mu}$ for all $w \in W$ as well as the $W$-invariance of $\mathcal{P}(\lambda)$.

It follows from the definition of the ordering $\prec$ that the set of dominant integral elements $\mu \in \mathfrak{h}^{*}$ with $\mu \prec \lambda$ is finite. Since all weights of $V$ are integral, Remark 7.3.11, together with the above, shows that

$$
\mathcal{P}(\lambda) \subseteq W\left(\left\{\mu \in \Lambda^{+} \mid \mu \prec \lambda\right\}\right)
$$

is finite because $W$ is finite and the set $\left\{\mu \in \Lambda^{+} \mid \mu \prec \lambda\right\}$ is finite. As all weight spaces $V_{\mu}$ are finite dimensional by Theorem 7.3.6, this concludes the proof.

Theorem 7.3.15 (Highest Weight Theorem). Let $\mathfrak{g}$ be a split semisimple Lie algebra and $\mathfrak{h}$ a splitting Cartan subalgebra of $\mathfrak{g}$. The assignment $\lambda \mapsto L(\lambda)$ defines a bijection between the set $\Lambda^{+}$of dominant integral functionals and the set of isomorphism classes of finite-dimensional simple $\mathfrak{g}$ modules.

Proof. First, we recall from Proposition 7.3.4 that each finite-dimensional simple $\mathfrak{g}$-module is a highest weight module with some highest weight $\lambda$. In view of Proposition 7.3.12, $\lambda$ is dominant integral, and Proposition 7.3.8 shows that it only depends on the isomorphism class of the module. Finally, Proposition 7.3.14 shows that any dominant integral functional is the highest weight of some finite-dimensional simple $\mathfrak{g}$-module.

\subsubsection*{The Eigenvalue of the Casimir Operator}
In this section, we construct a special element $C_{\mathfrak{g}}$ in the center of the enveloping algebra of $\mathfrak{g}$ and calculate its (scalar) action in a highest weight module. In special cases, this allows us to identify a given simple $\mathfrak{g}$-module.

Definition 7.3.16 (Universal Casimir element). Let $\mathfrak{g}$ be a finitedimensional split semisimple Lie algebra with Cartan-Killing form $\kappa$. As before, we choose for each $\beta \in \Delta^{+}$an $\mathfrak{s l}_{2}$-triple ( $h_{\beta}, e_{\beta}, f_{\beta}$ ) and $e_{\beta}^{*} \in \mathfrak{g}_{-\beta}$, $f_{\beta}^{*} \in \mathfrak{g}_{\beta}$ with

$$
\kappa\left(e_{\beta}, e_{\beta}^{*}\right)=1=\kappa\left(f_{\beta}, f_{\beta}^{*}\right)
$$

We further choose a basis $h_{1}, \ldots, h_{r}$ for $\mathfrak{h}$, and we write $h^{1}, \ldots, h^{r}$ for the dual basis with respect to the nondegenerate restriction of $\kappa$ to $\mathfrak{h} \times \mathfrak{h}$. Then

$$
\left\{h_{i}, e_{\beta}, f_{\beta}: i=1, \ldots, r ; \beta \in \Delta^{+}\right\}
$$

is a basis for $\mathfrak{g}$ and

$$
\left\{h^{i}, e_{\beta}^{*}, f_{\beta}^{*}: i=1, \ldots, r ; \beta \in \Delta^{+}\right\}
$$

is the dual basis with respect to $\kappa$. We therefore obtain a central element of $\mathcal{U}(\mathfrak{g})$ by

$$
C_{\mathfrak{g}}=\sum_{i=1}^{k} h_{i} h^{i}+\sum_{\beta \in \Delta^{+}} e_{\beta} e_{\beta}^{*}+f_{\beta} f_{\beta}^{*}
$$

(Lemma 5.5.16). It is called the universal Casimir element.\\
Lemma 7.3.17. For a positive system $\Delta^{+}$, we put $\rho=\frac{1}{2} \sum_{\alpha \in \Delta^{+}} \alpha$. If ( $\rho_{V}, V$ ) is a highest weight module with highest weight $\lambda$, also considered as a $\mathcal{U}(\mathfrak{g})$-module, then

$$
\rho_{V}\left(C_{\mathfrak{g}}\right)=(\lambda, \lambda+2 \rho) \mathbf{1}=\left(\|\lambda+\rho\|^{2}-\|\rho\|^{2}\right) \mathbf{1}
$$

If $\lambda$ is dominant and nonzero, then $\rho_{V}\left(C_{\mathfrak{g}}\right) \neq 0$.\\
Proof. We write the Casimir element $C_{\mathfrak{g}}$ in the form (7.4) as described in Definition 7.3.16. We compute the action of $C_{\mathfrak{g}}$ on $V$. Let $v^{+}$be a highest weight vector in $V$. Then $e_{\beta} \cdot v^{+}=f_{\beta}^{*} \cdot v^{+}=0$ for each $\beta \in \Delta^{+}$, and $\left[e_{\beta}, e_{\beta}^{*}\right]=h_{\beta}^{\prime}$ (cf. Lemma 6.3.3) implies

$$
e_{\beta} e_{\beta}^{*} \cdot v^{+}=\left[e_{\beta}, e_{\beta}^{*}\right] \cdot v^{+}+e_{\beta}^{*} e_{\beta} \cdot v^{+}=\lambda\left(h_{\beta}^{\prime}\right) v^{+}=(\lambda, \beta) v^{+}
$$

so that $\sum_{\beta \in \Delta^{+}}\left(e_{\beta} e_{\beta}^{*}+f_{\beta} f_{\beta}^{*}\right) \cdot v^{+}=2(\lambda, \rho) v^{+}$. On the other hand, we calculate

$$
\sum_{i=1}^{k} \lambda\left(h_{i}\right) \lambda\left(h^{i}\right)=\lambda\left(\sum_{i=1}^{k} \kappa\left(h_{\lambda}^{\prime}, h^{i}\right) h_{i}\right)=\lambda\left(h_{\lambda}^{\prime}\right)=(\lambda, \lambda) .
$$

Putting these facts together yields

$$
C_{\mathfrak{g}} v^{+}=(\lambda, \lambda+2 \rho) v^{+}=\left(\|\lambda+\rho\|^{2}-\|\rho\|^{2}\right) v^{+} .
$$

Since $C_{\mathfrak{g}}$ is central in $\mathcal{U}(\mathfrak{g})$ (Exercise 7.1.1), $C_{\mathfrak{g}}$ acts by the same scalar on the entire $\mathcal{U}(\mathfrak{g})$-module $V=\mathcal{U}(\mathfrak{g}) v^{+}$.

Finally, we assume that $\lambda$ is dominant and nonzero. Then $\lambda(\check{\alpha}) \geq 0$ for all $\alpha \in \Delta^{+}$implies that $(\lambda, \alpha) \geq 0$, and hence that $(\lambda, \rho) \geq 0$. This leads to $(\lambda, \lambda+2 \rho) \geq(\lambda, \lambda)>0$.

\subsubsection*{Exercises for Section 7.3}
Exercise 7.3.1. A $\mathfrak{g}$-module $V$ is said to be cyclic if it is generated by some element $v \in V$. If $\rho: \mathfrak{g} \rightarrow \mathfrak{g} \mathfrak{l}(V)$ is the module structure and $\widetilde{\rho}: \mathcal{U}(\mathfrak{g}) \rightarrow \operatorname{End}(V)$ the canonical extension, then the annihilator

$$
I:=\operatorname{Ann}_{\mathcal{U}(\mathfrak{g})}(v):=\{D \in \mathcal{U}(\mathfrak{g}): \widetilde{\rho}(D) v=0\}
$$

of $v$ is a left ideal. Show that:\\
(a) If $I \subseteq \mathcal{U}(\mathfrak{g})$ is a left ideal, then the quotient $\mathcal{U}(\mathfrak{g}) / I$ carries a natural $\mathfrak{g}$-module structure, defined by $x \cdot(D+I):=\sigma(x) D+I$, and this $\mathfrak{g}$-module is cyclic.\\
(b) Every cyclic $\mathfrak{g}$-module is isomorphic to one of the form $\mathcal{U}(\mathfrak{g}) / I$, as in (a).

Exercise 7.3.2. Simple $\mathfrak{g}$-modules are particular examples of cyclic $\mathfrak{g}$-modules. Show that:\\
(a) If $I \subseteq \mathcal{U}(\mathfrak{g})$ is a maximal (proper) left ideal, then the quotient $\mathcal{U}(\mathfrak{g}) / I$ is a simple $\mathfrak{g}$-module.\\
(b) Every simple $\mathfrak{g}$-module is isomorphic to one of the form $\mathcal{U}(\mathfrak{g}) / I$, where $I$ is a maximal left ideal in $\mathcal{U}(\mathfrak{g})$.

Exercise 7.3.3. Let $\mathfrak{g}$ be a reductive Lie algebra and $\mathfrak{h} \subseteq \mathfrak{g}$ a toral Cartan subalgebra. Let $V$ be a simple $\mathfrak{g}$-module on which $\mathfrak{h}$ acts by diagonalizable operators (such modules are called weight modules). Identifying the root system of $\mathfrak{g}$ with that of its semisimple commutator algebra, the notion of a positive system makes also sense for $\mathfrak{g}$. Show that, in this sense, for each positive system $\Delta^{+}$of $\Delta, V$ is a highest weight module.

\subsection*{Ado's Theorem}
In this section, we return to the question of when a finite-dimensional Lie algebra can be written as a subalgebra of a finite-dimensional associative algebra equipped with the commutator bracket. Ado's Theorem, in fact, says that this is always the case.

Theorem 7.4.1 (Ado's Theorem). Every finite-dimensional Lie algebra has a faithful, i.e., injective, finite-dimensional representation whose restriction to the maximal nilpotent ideal is nilpotent.

The proof of this theorem will be given below. Ado's Theorem is one of the cornerstones of finite-dimensional Lie theory. For any Lie algebra $\mathfrak{g}$, we have the adjoint representation ad: $\mathfrak{g} \rightarrow \mathfrak{g} \mathfrak{l}(\mathfrak{g})$ which is faithful if and only if $\mathfrak{z}(\mathfrak{g})$ is trivial. Therefore, the main point of Ado's Theorem is to find a representation of $\mathfrak{g}$ which is nontrivial on the center. A crucial first case is to see how to obtain faithful representations of nilpotent Lie algebras.

Proposition 7.4.2. Let $\mathfrak{n}$ be a nilpotent Lie algebra with $C^{k}(\mathfrak{n})=\{0\}$. For $j \in \mathbb{N}_{0}$, let $\mathcal{U}(\mathfrak{n})_{j} \subseteq \mathcal{U}(\mathfrak{n})$ be the ideal spanned by $\sum_{i \geq j} \sigma(\mathfrak{n})^{i}$, where $\sigma: \mathfrak{g} \rightarrow \mathcal{U}(\mathfrak{g})$ is the canonical embedding. Then the $\mathcal{U}(\mathfrak{n}) / \mathcal{U}(\mathfrak{n})_{j}$ are finitedimensional associative unital algebras, and the canonical homomorphism

$$
\sigma_{k}: \mathfrak{n} \rightarrow \mathcal{U}(\mathfrak{n}) / \mathcal{U}(\mathfrak{n})_{k}
$$

is injective.\\
Proof. We have to show that $\mathcal{U}(\mathfrak{n})_{k} \cap \sigma(\mathfrak{n})=\{0\}$.\\
Let $d_{j}:=\operatorname{dim} C^{j}(\mathfrak{n})$ for $j=1, \ldots, k$ and choose a basis $e_{i}$ for $\mathfrak{n}$ with the property that for each $j<k$ the $e_{i}, i \leq d_{j}$, form a basis for $C^{j}(\mathfrak{n})$. For a nonzero element $x \in \mathfrak{n}$, we defined its weight $w(x)$ as the number $h$ with $x \in C^{h}(\mathfrak{n}) \backslash C^{h+1}(\mathfrak{n})$ and $w(0):=\infty$. Since $\left[C^{a}(\mathfrak{n}), C^{b}(\mathfrak{n})\right] \subseteq C^{a+b}(\mathfrak{n})$ (Exercise 5.2.7), we have for any two elements $x, y \in \mathfrak{n}$ the relation

$$
w([x, y]) \geq w(x)+w(y)
$$

Moreover, if $x=\sum_{j} c_{j} e_{j} \in C^{a}(\mathfrak{n})$, then the choice of basis tells us that all $c_{j} e_{j}$ are in $C^{a}(\mathfrak{n})$ as well. For any monomial $e_{i_{1}} e_{i_{2}} \cdots e_{i_{s}} \in \mathcal{U}(\mathfrak{n})$, we define the weight as

$$
w\left(e_{i_{1}} e_{i_{2}} \cdots e_{i_{s}}\right):=\sum_{j=1}^{s} w\left(e_{i_{j}}\right)
$$

Now take a monomial and rewrite it, in the sense of the Poincaré-Birkhoff-Witt Theorem, as a linear combination of monomials of the form

$$
e_{1}^{m_{1}} \cdots e_{j}^{m_{j}}, \quad m_{1}, \ldots, m_{j} \in \mathbb{N}_{0} .
$$

Each time we have to substitute a product $e_{a} e_{b}, b<a$, by $e_{b} e_{a}+\left[e_{a}, e_{b}\right]$, we obtain in the sum a monomial of lesser degree, but when we write $\left[e_{a}, e_{b}\right]= \sum_{j} c_{j} e_{j}$, the weight of each of the resulting monomials is at least as big as the weight of the monomial we started with. This shows that when we write an element of the ideal $\mathcal{U}(\mathfrak{n})_{k}$ in the PBW basis, we have a linear combination of elements of weight at least $k$, but all elements of $\sigma(\mathfrak{n})$ have weight at most $k-1$. This completes the proof.

It follows immediately from the PBW-Theorem that all ideals $\mathcal{U}(\mathfrak{n})_{j}$ have finite codimension. Since left multiplication defines a faithful finitedimensional representation for each unital finite-dimensional associative algebra, the preceding proposition already proves Ado's Theorem for the special class of nilpotent Lie algebras.

Lemma 7.4.3. Let $\mathfrak{g}$ be a Lie algebra and $\sigma: \mathfrak{g} \rightarrow \mathcal{U}(\mathfrak{g})_{L}$ be the canonical homomorphism. For every derivation $D \in \operatorname{der}(\mathfrak{g})$, there exists a unique derivation $\widetilde{D} \in \operatorname{der}(\mathcal{U}(\mathfrak{g}))$ with $\sigma \circ D=\widetilde{D} \circ \sigma$.\\
\includegraphics[max width=\textwidth, center]{2025_10_18_3eaf04e77f3cb364fce5g-202}

Proof. Let $\mathcal{A}:=M_{2}(\mathcal{U}(\mathfrak{g}))$ be the algebra of the $2 \times 2$-matrices with entries in $\mathcal{U}(\mathfrak{g})$. Then the map

$$
\varphi: \mathfrak{g} \rightarrow \mathcal{A}_{L}, \quad \varphi(x)=\left(\begin{array}{cc}
x & D(x) \\
0 & x
\end{array}\right)
$$

is a homomorphism of Lie algebras. From the universal property of $(\mathcal{U}(\mathfrak{g}), \sigma)$, we now obtain a homomorphism $\widetilde{\varphi}: \mathcal{U}(\mathfrak{g}) \rightarrow \mathcal{A}$ of unital algebras with $\widetilde{\varphi} \circ \sigma= \varphi$. Since $\mathcal{U}(\mathfrak{g})$ is generated by $\mathbf{1}$ and $\sigma(\mathfrak{g})$, we have

$$
\widetilde{\varphi}(x)=\left(\begin{array}{cc}
x & \widetilde{D}(x) \\
0 & x
\end{array}\right)
$$

for some linear map $\widetilde{D}: \mathcal{U}(\mathfrak{g}) \rightarrow \mathcal{U}(\mathfrak{g})$, and since $\widetilde{\varphi}$ is an algebra homomorphism, $\widetilde{D}$ is a derivation. It immediately follows from the construction that $\widetilde{D} \circ \sigma=D$. The uniqueness of $\widetilde{D}$ follows from the fact that $\sigma(\mathfrak{g})$ and $\mathbf{1}$ generate $\mathcal{U}(\mathfrak{g})$ because a derivation annihilates the unit and it is determined by its values on a generating set. \(\square\)

Proposition 7.4.4. If $\mathfrak{g}=\mathfrak{n} \rtimes_{\gamma} \mathfrak{l}$ is a semidirect sum of a nilpotent ideal and a Lie algebra $\mathfrak{l}$, then $\mathfrak{g}$ has a finite-dimensional representation which is faithful and nilpotent on $\mathfrak{n}$.

Proof. Suppose that $C^{k}(\mathfrak{n})=\{0\}$ and put $\mathcal{A}:=\mathcal{U}(\mathfrak{n}) / \mathcal{U}(\mathfrak{n})_{k}$. Then Proposition 7.4.2 implies that the natural map $\sigma_{k}: \mathfrak{n} \rightarrow \mathcal{A}$ is injective. Clearly, this defines a nilpotent representation of $\mathfrak{n}$ on $\mathcal{A}$.

In view of the preceding Lemma 7.4.3, the representation $\gamma: \mathfrak{l} \rightarrow \operatorname{der}(\mathfrak{n})$ induces a representation $\widetilde{\gamma}: \mathfrak{l} \rightarrow \operatorname{der}(\mathcal{U}(\mathfrak{n}))$, and it is clear that the derivations $\widetilde{\gamma}(x)$ preserve the ideal $\mathcal{U}(\mathfrak{n})_{k}$, hence induce derivations $\bar{\gamma}(x)$ on $\mathcal{A}$. Write $L_{x}: \mathcal{A} \rightarrow \mathcal{A}, a \mapsto x a$ for the left multiplications on $\mathcal{A}$. Then

$$
\pi: \mathfrak{g}=\mathfrak{n} \rtimes_{\gamma} \mathfrak{l} \rightarrow \mathfrak{g} \mathfrak{l}(\mathcal{A}), \quad(n, x) \mapsto L_{\sigma_{k}(n)}+\bar{\gamma}(x)
$$

defines a representation of $\mathfrak{g}$ on $\mathcal{A}$ whose restriction to $\mathfrak{n}$ is faithful and nilpotent. \(\square\)

Lemma 7.4.5. Let $\mathfrak{n} \subseteq \mathfrak{g}$ be a nilpotent ideal and $x \in \mathfrak{g}$. If there exists an $n \in \mathfrak{n}$ such that $\operatorname{ad}(x+n)$ is nilpotent, then ad $x$ is nilpotent.

Proof. Let $\mathcal{F}=\left(\mathfrak{g}_{0}, \ldots, \mathfrak{g}_{k}\right)$ be a maximal flag in $\mathfrak{g}$ consisting of ideals. This means in particular that the quotient algebras $\mathfrak{q}_{j}:=\mathfrak{g}_{j} / \mathfrak{g}_{j-1}, j=1, \ldots, k$,\\
are simple $\mathfrak{g}$-modules. As $\mathfrak{g}$ is a nilpotent $\mathfrak{n}$-module, the same holds for all the simple quotients $\mathfrak{g}_{j+1} / \mathfrak{g}_{j}$, so that $\mathfrak{n}$ acts trivially on them. This means that $\left[\mathfrak{n}, \mathfrak{g}_{j+1}\right] \subseteq \mathfrak{g}_{j}$, i.e., ad $\mathfrak{n} \subseteq \mathfrak{g}_{n}(\mathcal{F})$. But ad $x$ is nilpotent if and only if all the induced maps $\operatorname{ad}_{\mathfrak{g}_{j+1} / \mathfrak{g}_{j}}(x)$ are nilpotent. In view of

$$
\operatorname{ad}_{\mathfrak{g}_{j+1} / \mathfrak{g}_{j}}(x)=\operatorname{ad}_{\mathfrak{g}_{j+1} / \mathfrak{g}_{j}}(x+n),
$$

all these maps are nilpotent, so that $\operatorname{ad} x$ is nilpotent.\\
Lemma 7.4.6. Let $\mathfrak{r}$ be a solvable Lie algebra and ( $\rho, V$ ) a finite-dimensional representation of $\mathfrak{r}$. Then the set

$$
\mathfrak{n}_{\rho}:=\left\{x \in \mathfrak{r}:(\exists N \in \mathbb{N}) \rho_{V}(x)^{N}=0\right\}
$$

of those elements in $\mathfrak{r}$ for which $\rho_{V}(x)$ is nilpotent is an ideal of $\mathfrak{r}$.\\
Proof. After complexification, we may assume that $\mathbb{K}=\mathbb{C}$. Then we use Lie's Theorem to find a complete flag $\mathcal{F}$ in $V$ with $\rho_{V}(\mathfrak{r}) \subseteq \mathfrak{g}(\mathcal{F})$. Then $\mathfrak{n}_{\rho}=\rho_{V}^{-1}\left(\mathfrak{g}_{n}(\mathcal{F})\right)$ is the inverse image of the ideal $\mathfrak{g}_{n}(\mathcal{F})$ of $\mathfrak{g}(\mathcal{F})$, hence an ideal of $\mathfrak{r}$.

Remark 7.4.7. Applying the preceding lemma to the adjoint representation, we see that for a finite-dimensional Lie algebra $\mathfrak{g}$, the set

$$
\mathfrak{n}_{\mathrm{ad}}:=\left\{x \in \mathfrak{r a d}(\mathfrak{g}):(\exists N \in \mathbb{N})(\operatorname{ad} x)^{N}=0\right\}
$$

is an ideal of $\mathfrak{r}$, whose nilpotency follows from Engel's Theorem. Moreover, $\mathfrak{n}_{\text {ad }}$ contains each nilpotent ideal of $\mathfrak{g}$, hence is the maximal nilpotent ideal of $\mathfrak{g}$.\\
(Proof of Ado's Theorem). We shall prove Ado's Theorem by embedding $\mathfrak{g}$ into a semidirect sum of a nilpotent ideal and a reductive Lie algebra, to which we apply Proposition 7.4.4.

Let $\mathfrak{r}:=\mathfrak{r a d}(\mathfrak{g})$ and $\mathfrak{g}=\mathfrak{r} \rtimes \mathfrak{s}$ be a Levi decomposition of $\mathfrak{g}$. Applying Proposition 5.5.17 to the semisimple $\mathfrak{s}$-module $\mathfrak{r}$, we see that

$$
\mathfrak{r}=[\mathfrak{s}, \mathfrak{r}]+\mathfrak{z}_{\mathfrak{r}}(\mathfrak{s}) .
$$

Let $\mathfrak{h} \subseteq \mathfrak{z}_{\mathfrak{r}}(\mathfrak{s})$ be a Cartan subalgebra. Then $\mathfrak{d}:=(\operatorname{ad} \mathfrak{h})_{s} \subseteq \mathfrak{g} \mathfrak{l}(\mathfrak{g})$ is a commutative subalgebra of derivations of $\mathfrak{g}$ (Lemma 6.1.13). We consider the Lie algebra

$$
\widehat{\mathfrak{g}}:=\mathfrak{g} \rtimes \mathfrak{d} .
$$

Our goal is to show that $\widehat{\mathfrak{g}}$ is the semidirect sum of a reductive Lie algebra and a nilpotent ideal.

Since ad $\mathfrak{h}$ annihilates the subspace $\mathfrak{s}$ of $\mathfrak{g}$, the same is true for $\mathfrak{d}$, so that $\mathfrak{l}:=\mathfrak{s} \oplus \mathfrak{d}$ is a reductive Lie algebra with center $\mathfrak{d}$. We also note that

$$
\widehat{\mathfrak{g}}=(\mathfrak{r} \rtimes \mathfrak{s}) \rtimes \mathfrak{d}=\widehat{\mathfrak{r}} \rtimes \mathfrak{s},
$$

where $\widehat{\mathfrak{r}}:=\mathfrak{r} \rtimes \mathfrak{d}$ is the solvable radical of $\widehat{\mathfrak{g}}$.

Next we recall that the quotient map $\mathfrak{z}_{\mathfrak{r}}(\mathfrak{s}) \rightarrow \mathfrak{z}_{\mathfrak{r}}(\mathfrak{s}) /\left[\mathfrak{z}_{\mathfrak{r}}(\mathfrak{s}), \mathfrak{z}_{\mathfrak{r}}(\mathfrak{s})\right]$ maps $\mathfrak{h}$ surjectively onto the abelian quotient algebra (Proposition 6.1.11(ii)), so that

$$
\mathfrak{z}_{\mathfrak{r}}(\mathfrak{s})=\left[\mathfrak{z}_{\mathfrak{r}}(\mathfrak{s}), \mathfrak{z}_{\mathfrak{r}}(\mathfrak{s})\right]+\mathfrak{h} \subseteq[\mathfrak{r}, \mathfrak{r}]+\mathfrak{h}
$$

and thus

$$
\mathfrak{r} \subseteq[\mathfrak{s}, \mathfrak{r}]+\mathfrak{z}_{\mathfrak{r}}(\mathfrak{s}) \subseteq[\mathfrak{g}, \mathfrak{r}]+\mathfrak{h}
$$

Let $x \in \mathfrak{r}$ and write it as $x=n+h$ with $n \in[\mathfrak{g}, \mathfrak{r}]$ and $h \in \mathfrak{h}$. Then $d:=(\operatorname{ad} h)_{s} \in \mathfrak{d}$, and thus $x-d \in \widehat{\mathfrak{r}}$. Now $[\mathfrak{g}, \mathfrak{r}]$ is a nilpotent ideal of $\widehat{\mathfrak{g}}$ and $\operatorname{ad}_{\widehat{\mathfrak{g}}}(h-d)$ is nilpotent on $\mathfrak{g} \supseteq[\widehat{\mathfrak{g}}, \widehat{\mathfrak{g}}]$, where it coincides with $\operatorname{ad}_{\mathfrak{g}}(h)-d$. We thus conclude with Lemma 7.4.5 that $\operatorname{ad}_{\widehat{\mathfrak{g}}}(x-d)$ is nilpotent, hence contained in the maximal nilpotent ideal

$$
\widehat{\mathfrak{n}}:=\left\{x \in \widehat{\mathfrak{r}}:(\exists N \in \mathbb{N})(\operatorname{ad} x)^{N}=0\right\}
$$

of $\widehat{\mathfrak{g}}$ (Remark 7.4.7). We conclude that $\widehat{\mathfrak{r}} \subseteq \widehat{\mathfrak{n}}+\mathfrak{d}$. On the other hand, $\widehat{\mathfrak{n}} \cap \mathfrak{d}= \{0\}$ follows from the fact that the nonzero elements of $\mathfrak{d}$ act by semisimple endomorphisms on $\mathfrak{g}$. Therefore,

$$
\widehat{\mathfrak{g}}=\widehat{\mathfrak{n}} \rtimes \mathfrak{l}
$$

is a semidirect sum of the nilpotent ideal $\widehat{\mathfrak{n}}$ and the reductive Lie algebra $\mathfrak{l}$. Note that $\mathfrak{z}(\mathfrak{g}) \subseteq \mathfrak{z}(\widehat{\mathfrak{g}}) \subseteq \widehat{\mathfrak{n}}$. Finally, Proposition 7.4.4 provides a representation $\widehat{\pi}: \widehat{\mathfrak{g}} \rightarrow \mathfrak{g} \mathfrak{l}(V)$ which is faithful and nilpotent on $\widehat{\mathfrak{n}}$. Hence the restriction $\pi:=\left.\widehat{\pi}\right|_{\mathfrak{g}}$ is faithful on $\mathfrak{z}(\mathfrak{g})$. Now the direct sum representation $\pi \oplus$ ad of $\mathfrak{g}$ on $V \oplus \mathfrak{g}$ is a faithful finite-dimensional representation of $\mathfrak{g}$. Since the maximal nilpotent ideal $\mathfrak{n}$ of $\mathfrak{g}$ is clearly contained in $\widehat{\mathfrak{n}}$, the representation is nilpotent on $\mathfrak{n}$.

\subsubsection*{Exercises for Section 7.4}
Exercise 7.4.1. Let $\mathfrak{g}$ be a Lie algebra, let $\mathfrak{b}$ be a characteristic ideal in $\mathfrak{g}$ and $\mathfrak{a}$ be a characteristic ideal in $\mathfrak{b}$. Then $\mathfrak{a}$ is a characteristic ideal in $\mathfrak{g}$.

Exercise 7.4.2. Find an example of an ideal of an ideal of a Lie algebra $\mathfrak{g}$ which is not an ideal of $\mathfrak{g}$.

Exercise 7.4.3. Let $\mathfrak{g}$ be a finite-dimensional Lie algebra, let $\mathcal{U}(\mathfrak{g})$ be its enveloping algebra and $\sigma: \mathfrak{g} \rightarrow \mathcal{U}(\mathfrak{g})$ be the canonical embedding. By $\mathcal{U}_{i}(\mathfrak{g})$, we denote the subspace spanned by the products of degree less or equal to $i$ (cf. Lemma 7.1.6). Show:\\
(i) For every automorphism $\alpha$ of the Lie algebra $\mathfrak{g}$, there is precisely one automorphism $\mathcal{U}(\alpha)$ of $\mathcal{U}(\mathfrak{g})$ with $\mathcal{U}(\alpha) \circ \sigma=\sigma \circ \alpha$. This automorphism leaves $\mathcal{U}_{i}(\mathfrak{g})$ invariant for every $i \in \mathbb{N}$.\\
(ii) If $\gamma: \mathbb{R} \rightarrow \operatorname{Aut}(\mathcal{U}(\mathfrak{g}))$ is a one-parameter group of automorphisms of $\mathcal{U}(\mathfrak{g})$ for which

$$
\gamma(t) \mathcal{U}_{i}(\mathfrak{g}) \subseteq \mathcal{U}_{i}(\mathfrak{g}) \quad \text { for } \quad i \in \mathbb{N},
$$

then $D:=\left.\frac{d}{d t}\right|_{t=0} \gamma(t)$ exists and it is a derivation of $\mathcal{U}(\mathfrak{g})$.\\
(iii) Apply (i) and (ii) to obtain a new proof of Lemma 7.4.3.

By (i), the map $\mathcal{U}\left(e^{\operatorname{ad} x}\right)$ is a well defined automorphism of $\mathcal{U}(\mathfrak{g})$ for every $x \in \mathfrak{g}$, and the derivation ad $x$ can directly be continued to $\mathcal{U}(\mathfrak{g})$ by

$$
\operatorname{ad} x(z):=x z-z x \quad \text { for } \quad x \in \mathfrak{g}, z \in \mathcal{U}(\mathfrak{g}) .
$$

(iv) For $z \in \mathcal{U}(\mathfrak{g})$, the following statements are equivalent:\\
(i) $z \in Z(\mathcal{U}(\mathfrak{g}))$.\\
(ii) $\operatorname{ad} x(z)=0$ for all $x \in \mathfrak{g}$.\\
(iii) $\mathcal{U}\left(e^{\operatorname{ad} x}\right) z=z$ for all $x \in \mathfrak{g}$.

Exercise 7.4.4. Let $\alpha: \mathfrak{g} \rightarrow \operatorname{End}(V)$ be a representation of the Lie algebra $\mathfrak{g}$ and $\widetilde{\alpha}: \mathcal{U}(\mathfrak{g}) \rightarrow \operatorname{End}(V)$ be the corresponding representation of the enveloping algebra. Show: Let $z \in Z(\mathcal{U}(\mathfrak{g}))$ and $V_{\alpha}$ be an eigenspace of $\widetilde{\alpha}(z)$, then $V_{\alpha}$ is invariant under $\alpha(\mathfrak{g})$.

\subsection*{Lie Algebra Cohomology}
The cohomology of Lie algebras is the natural tool to understand how we can build new Lie algebras $\widehat{\mathfrak{g}}$ from given Lie algebras $\mathfrak{g}$ and $\mathfrak{n}$ in such a way that $\mathfrak{n} \unlhd \widehat{\mathfrak{g}}$ and $\widehat{\mathfrak{g}} / \mathfrak{n} \cong \mathfrak{g}$. An important special case of this situation arises if $\mathfrak{n}$ is assumed to be abelian, so that $\mathfrak{n}$ is simply a $\mathfrak{g}$-module. We shall see, in particular, how the abelian extensions of Lie algebras can be parameterized by a certain cohomology space.

We shall also deal with the extension problem for $\mathfrak{g}$-modules, i.e., the problem to determine, for a pair ( $V, W$ ) of $\mathfrak{g}$-modules, how many nonisomorphic modules $\widehat{V}$ exist which contain $W$ as a submodule and satisfy $\widehat{V} / W \cong V$.

Throughout this section, $\mathfrak{g}$ denotes a Lie algebra over the field $\mathbb{K}$. We do not have to make any assumption on the dimension of $\mathfrak{g}$ or the nature of the field $\mathbb{K}$.

\subsubsection*{Basic Definitions and Properties}
Definition 7.5.1. Let $V$ and $W$ be vector spaces and $p \in \mathbb{N}$. A multilinear map $f: W^{p} \rightarrow V$ is called alternating if

$$
f\left(w_{\sigma_{1}}, \ldots, w_{\sigma_{p}}\right)=\operatorname{sgn}(\sigma) f\left(w_{1}, \ldots, w_{p}\right)
$$

for $w_{i} \in W$ and $\operatorname{sgn}(\sigma)$ is the sign of the permutation $\sigma \in S_{p}$.

Definition 7.5.2. Let $\mathfrak{g}$ be a Lie algebra and $V$ a $\mathfrak{g}$-module.\\
(a) We write $C^{p}(\mathfrak{g}, V)$ for the space of alternating $p$-linear mappings $\mathfrak{g}^{p} \rightarrow V$ (the $p$-cochains) and put $C^{0}(\mathfrak{g}, V):=V$. We also define

$$
C(\mathfrak{g}, V):=\bigoplus_{k=0}^{\infty} C^{k}(\mathfrak{g}, V) .
$$

On $C^{p}(\mathfrak{g}, V)$, we define the (Chevalley-Eilenberg) differential d by

$$
\begin{aligned}
\mathrm{d} \omega\left(x_{0}, \ldots, x_{p}\right):= & \sum_{j=0}^{p}(-1)^{j} x_{j} \cdot \omega\left(x_{0}, \ldots, \widehat{x}_{j}, \ldots, x_{p}\right) \\
& +\sum_{i<j}(-1)^{i+j} \omega\left(\left[x_{i}, x_{j}\right], x_{0}, \ldots, \widehat{x}_{i}, \ldots, \widehat{x}_{j}, \ldots, x_{p}\right),
\end{aligned}
$$

where $\widehat{x}_{j}$ means that $x_{j}$ is omitted. Observe that the right-hand side defines for each $\omega \in C^{p}(\mathfrak{g}, V)$ an element of $C^{p+1}(\mathfrak{g}, V)$ because it is alternating. To see that this is the case, it suffices to show that it vanishes if $x_{i}=x_{i+1}$ for $i=0, \ldots, p-1$. Since $\omega$ is alternating, we only need to note that for $x_{i}=x_{i+1}$, we have

$$
\begin{aligned}
0= & (-1)^{i} x_{i} \cdot \omega\left(x_{0}, \ldots, \widehat{x}_{i}, x_{i+1} \ldots, x_{p}\right) \\
& +(-1)^{i+1} x_{i} \cdot \omega\left(x_{0}, \ldots, x_{i}, \widehat{x}_{i+1}, \ldots, x_{p}\right), \\
0= & (-1)^{i+j} \omega\left(\left[x_{i}, x_{j}\right], \ldots, \widehat{x}_{i}, \ldots, \widehat{x}_{j}, \ldots, x_{p}\right) \\
& +(-1)^{i+j+1} \omega\left(\left[x_{i+1}, x_{j}\right], \ldots, \widehat{x}_{i+1}, \ldots, \widehat{x}_{j}, \ldots, x_{p}\right), \\
0= & (-1)^{j+i} \omega\left(\left[x_{j}, x_{i}\right], \ldots, \widehat{x}_{j}, \ldots, \widehat{x}_{i}, \ldots, x_{p}\right) \\
& +(-1)^{j+i+1} \omega\left(\left[x_{j}, x_{i+1}\right], \ldots, \widehat{x}_{j}, \ldots, \widehat{x}_{i+1}, \ldots, x_{p}\right) .
\end{aligned}
$$

Putting the differentials on all the spaces $C^{p}(\mathfrak{g}, V)$ together, we obtain a linear map $\mathrm{d}=\mathrm{d}_{\mathfrak{g}}: C(\mathfrak{g}, V) \rightarrow C(\mathfrak{g}, V)$.

The elements of the subspace

$$
Z^{k}(\mathfrak{g}, V):=\operatorname{ker}\left(\left.\mathrm{d}\right|_{C^{k}(\mathfrak{g}, V)}\right)
$$

are called $k$-cocycles, and the elements of the spaces

$$
B^{k}(\mathfrak{g}, V):=\mathrm{d}\left(C^{k-1}(\mathfrak{g}, V)\right) \quad \text { and } \quad B^{0}(\mathfrak{g}, V):=\{0\}
$$

are called $k$-coboundaries. We will see below that $\mathrm{d}^{2}=0$, which implies that $B^{k}(\mathfrak{g}, V) \subseteq Z^{k}(\mathfrak{g}, V)$, so that it makes sense to define the kth cohomology space of $\mathfrak{g}$ with values in the module $V$ :

$$
H^{k}(\mathfrak{g}, V):=Z^{k}(\mathfrak{g}, V) / B^{k}(\mathfrak{g}, V) .
$$

(b) We further define for each $x \in \mathfrak{g}$ and $p>0$ the insertion map or contraction

$$
i_{x}: C^{p}(\mathfrak{g}, V) \rightarrow C^{p-1}(\mathfrak{g}, V), \quad\left(i_{x} \omega\right)\left(x_{1}, \ldots, x_{p-1}\right)=\omega\left(x, x_{1}, \ldots, x_{p-1}\right)
$$

We further define $i_{x}$ to be 0 on $C^{0}(\mathfrak{g}, V)$.\\
Note 7.5.3. The names (co)cycle, (co)boundary, and (co)boundary operator are derived from simplicial homology theory which has large formal similarities with the cohomology theory for Lie algebras. In simplicial homology, one considers simplices, such as triangles or tetrahedra. A triangle is described by its vertices $x_{1}, x_{2}, x_{3}$. Its faces with inherent orientation given by the ordering of the indices are denoted by $x_{12}, x_{23}, x_{31}$. The entire triangle is denoted by $x_{123}$. Chains are formal linear combinations of the type $\sum a_{i j} x_{i j}$. Cycles are chains which have no boundary, for instance, $x_{12}+x_{23}+x_{31}$ or $2 x_{12}+2 x_{23}+2 x_{31}$. For the latter one, the cycle $x_{12}+x_{23}+x_{31}$ is passed through twice. The boundary of $x_{12}$ is $x_{2}-x_{1}$, the boundary of $x_{123}$ is $x_{12}+x_{23}+x_{31}$. Chains which are themselves boundaries of something have no boundary, i.e., they are cycles. This corresponds to the relation $\mathrm{d}^{2}=0$ in Lie algebra cohomology.

Remark 7.5.4. For elements of low degree, we have in particular:

$$
\begin{aligned}
p=0: \quad \mathrm{d} \omega(x)= & x \cdot \omega \\
p=1: \quad \mathrm{d} \omega(x, y)= & x \cdot \omega(y)-y \cdot \omega(x)-\omega([x, y]) \\
p=2: \quad \mathrm{d} \omega(x, y, z)= & x \cdot \omega(y, z)-y \cdot \omega(x, z)+z \cdot \omega(x, y) \\
& \quad-\omega([x, y], z)+\omega([x, z], y)-\omega([y, z], x) \\
= & x \cdot \omega(y, z)+y \cdot \omega(z, x)+z \cdot \omega(x, y) \\
& \quad-\omega([x, y], z)-\omega([y, z], x)-\omega([z, x], y) .
\end{aligned}
$$

Example 7.5.5. (a) This means that

$$
Z^{0}(\mathfrak{g}, V)=V^{\mathfrak{g}}:=\{v \in V: \mathfrak{g} \cdot v=\{0\}\}
$$

is the maximal trivial submodule of $V$. Since $B^{0}(\mathfrak{g}, V)$ is trivial by definition, we obtain

$$
H^{0}(\mathfrak{g}, V)=V^{\mathfrak{g}}
$$

(b) The elements $\alpha \in Z^{1}(\mathfrak{g}, V)$ are also called crossed homomorphisms. They are defined by the condition

$$
\alpha([x, y])=x \cdot \alpha(y)-y \cdot \alpha(x), \quad x, y \in \mathfrak{g} .
$$

The elements $\alpha(x) \cdot v:=x \cdot v$ of the subspace $B^{1}(\mathfrak{g}, V)$ are also called principal crossed homomorphisms. It follows immediately from the definition of a $\mathfrak{g}$-module that each principal crossed homomorphism is a crossed homomorphism.

If $V$ is a trivial module, then it is not hard to compute the cohomology spaces in degree one. In view of $\{0\}=\mathrm{d} V=\mathrm{d} C^{0}(\mathfrak{g}, V)=B^{1}(\mathfrak{g}, V)$, we\\
have $H^{1}(\mathfrak{g}, V)=Z^{1}(\mathfrak{g}, V)$, and the condition that $\alpha: \mathfrak{g} \rightarrow V$ is a crossed homomorphism reduces to $\alpha([x, y])=\{0\}$ for $x, y \in \mathfrak{g}$. This leads to

$$
H^{1}(\mathfrak{g}, V) \cong \operatorname{Hom}(\mathfrak{g} /[\mathfrak{g}, \mathfrak{g}], V) \cong \operatorname{Hom}_{\text {Lie alg }}(\mathfrak{g}, V)
$$

Example 7.5.6. Let $\mathfrak{g}$ be a Lie algebra and $V:=\mathfrak{g}$, considered as a trivial $\mathfrak{g}$-module. Then the Lie bracket

$$
\beta: \mathfrak{g} \times \mathfrak{g} \rightarrow \mathfrak{g}, \quad \beta(x, y):=[x, y]
$$

is a 2 -cocycle. In fact, the two Lie algebra axioms mean that $\beta \in C^{2}(\mathfrak{g}, \mathfrak{g})$ and $\mathrm{d} \beta=0$.

Example 7.5.7. Let $\mathfrak{g}$ be an abelian Lie algebra and $V$ a trivial $\mathfrak{g}$-module. Then $\mathrm{d}=0$, so that $H^{p}(\mathfrak{g}, V)=C^{p}(\mathfrak{g}, V)$ holds for each $p \in \mathbb{N}_{0}$.

Our first goal will be to show that $\mathrm{d}^{2}=0$. This can be proved directly by an lengthy computation (Exercise 7.5.7). We will follow another way which is more conceptual and leads to additional insights and tools which are useful in other situations: Let $\left(\rho_{V}, V\right)$ be a $\mathfrak{g}$-module. Then the representations on the space $\operatorname{Mult}^{p}(\mathfrak{g}, V)$ of $p$-linear maps $\mathfrak{g}^{p} \rightarrow V$ defined by

$$
\rho_{+}(x) \omega:=\rho_{V}(x) \circ \omega
$$

and

$$
\left(\rho_{j}(x) \omega\right)\left(x_{1}, \ldots, x_{p}\right):=-\omega\left(x_{1}, \ldots, x_{j-1}, \operatorname{ad} x\left(x_{j}\right), x_{j+1}, \ldots, x_{p}\right)
$$

commute pairwise. Therefore, the sum of these representations is again a representation, which implies the following lemma, by restricting to the subspace $C^{p}(\mathfrak{g}, V) \subseteq \operatorname{Mult}^{p}(\mathfrak{g}, V)$.\\
Lemma 7.5.8. There exists a representation $\mathcal{L}=\rho_{+}+\rho_{0}=\rho_{+}+\sum_{j=1}^{p} \rho_{j}$ of $\mathfrak{g}$ on $C(\mathfrak{g}, V)$ given on the subspace $C^{p}(\mathfrak{g}, V)$ by

$$
\begin{aligned}
& \left(\mathcal{L}_{x} \omega\right)\left(x_{1}, \ldots, x_{p}\right) \\
& \quad=x \cdot \omega\left(x_{1}, \ldots, x_{p}\right)-\sum_{j=1}^{p} \omega\left(x_{1}, \ldots, x_{j-1},\left[x, x_{j}\right], x_{j+1}, \ldots, x_{p}\right) \\
& \quad=x \cdot \omega\left(x_{1}, \ldots, x_{p}\right)+\sum_{j=1}^{p}(-1)^{j} \omega\left(\left[x, x_{j}\right], x_{1}, \ldots, \widehat{x}_{j}, \ldots, x_{p}\right)
\end{aligned}
$$

Note that $\rho_{0}$ is the representation on $C(\mathfrak{g}, V)$ corresponding to the trivial module structure on $V$.

Lemma 7.5.9 (Cartan Formula). The representation $\mathcal{L}: \mathfrak{g} \rightarrow \mathfrak{g} \mathfrak{l}(C(\mathfrak{g}, V))$ satisfies, for $x \in \mathfrak{g}$, the Cartan formula

$$
\mathcal{L}_{x}=\mathrm{d} \circ i_{x}+i_{x} \circ \mathrm{~d}
$$

Proof. Using the insertion map $i_{x_{0}}$, we can rewrite the formula for the differential as

$$
\begin{aligned}
& \left(i_{x_{0}} \mathrm{~d} \omega\right)\left(x_{1}, \ldots, x_{p}\right) \\
= & x_{0} \cdot \omega\left(x_{1}, \ldots, x_{p}\right)-\sum_{j=1}^{p}(-1)^{j-1} x_{j} \cdot \omega\left(x_{0}, \ldots, \widehat{x}_{j}, \ldots, x_{p}\right) \\
& +\sum_{j=1}^{p}(-1)^{j} \omega\left(\left[x_{0}, x_{j}\right], x_{1}, \ldots, \widehat{x}_{j}, \ldots, x_{p}\right) \\
& +\sum_{1 \leq i<j}(-1)^{i+j} \omega\left(\left[x_{i}, x_{j}\right], x_{0}, \ldots, \widehat{x}_{i}, \ldots, \widehat{x}_{j}, \ldots, x_{p}\right) \\
= & x_{0} \cdot \omega\left(x_{1}, \ldots, x_{p}\right)-\sum_{j=1}^{p} \omega\left(x_{1}, \ldots, x_{j-1},\left[x_{0}, x_{j}\right], x_{j+1}, \ldots, x_{p}\right) \\
& -\sum_{j=1}^{p}(-1)^{j-1} x_{j} \cdot \omega\left(x_{0}, \ldots, \widehat{x}_{j}, \ldots, x_{p}\right) \\
& -\sum_{1 \leq i<j}(-1)^{i+j} \omega\left(x_{0},\left[x_{i}, x_{j}\right], \ldots, \widehat{x}_{i}, \ldots, \widehat{x}_{j}, \ldots, x_{p}\right) \\
= & \left(\mathcal{L}_{x_{0}} \omega\right)\left(x_{1}, \ldots, x_{p}\right)-\mathrm{d}\left(i_{x_{0}} \omega\right)\left(x_{1}, \ldots, x_{p}\right) .
\end{aligned}
$$

This proves the Cartan formula.\\
Lemma 7.5.10. Any two elements $x, y \in \mathfrak{g}$ satisfy $i_{[x, y]}=\left[i_{x}, \mathcal{L}_{y}\right]$.\\
Proof. The explicit formula for $\mathcal{L}_{y}$ (Lemma 7.5.8) yields for $x=x_{1}$ the relation $i_{x} \mathcal{L}_{y}=\mathcal{L}_{y} i_{x}-i_{[y, x]}$.

Lemma 7.5.11. For each $x \in \mathfrak{g}$, the Lie derivative $\mathcal{L}_{x}$ commutes with d .\\
Proof. In view of Lemma 7.5.10, we obtain with the Cartan formula

$$
\begin{aligned}
{\left[\mathcal{L}_{x}, \mathcal{L}_{y}\right] } & =\left[\mathrm{d} \circ i_{x}, \mathcal{L}_{y}\right]+\left[i_{x} \circ \mathrm{~d}, \mathcal{L}_{y}\right] \\
& =\left[\mathrm{d}, \mathcal{L}_{y}\right] \circ i_{x}+\mathrm{d} \circ i_{[x, y]}+i_{[x, y]} \circ \mathrm{d}+i_{x} \circ\left[\mathrm{~d}, \mathcal{L}_{y}\right] \\
& =\left[\mathrm{d}, \mathcal{L}_{y}\right] \circ i_{x}+\mathcal{L}_{[x, y]}+i_{x} \circ\left[\mathrm{~d}, \mathcal{L}_{y}\right]
\end{aligned}
$$

so that the fact that $\mathcal{L}$ is a representation leads to

$$
\left[\mathrm{d}, \mathcal{L}_{y}\right] \circ i_{x}+i_{x} \circ\left[\mathrm{~d}, \mathcal{L}_{y}\right]=0
$$

We now prove by induction over $k$ that $\left[\mathrm{d}, \mathcal{L}_{y}\right]$ vanishes on $C^{k}(\mathfrak{g}, V)$. For $\omega \in C^{0}(\mathfrak{g}, V) \cong V$, we have

$$
\begin{aligned}
& \left(\left[\mathrm{d}, \mathcal{L}_{y}\right] \omega\right)(x)=\mathrm{d}(y \cdot \omega)(x)-(y \cdot(\mathrm{~d} \omega))(x) \\
& =x \cdot(y \cdot \omega)-(y \cdot(x \cdot \omega)-\mathrm{d} \omega([y, x]))=[x, y] \cdot \omega+[y, x] \cdot \omega=0
\end{aligned}
$$

Suppose that $\left[\mathrm{d}, \mathcal{L}_{y}\right] C^{k}(\mathfrak{g}, V)=\{0\}$. Then (7.6) implies that

$$
\begin{aligned}
i_{x}\left[\mathrm{~d}, \mathcal{L}_{y}\right] C^{k+1}(\mathfrak{g}, V) & =-\left[\mathrm{d}, \mathcal{L}_{y}\right] i_{x} C^{k+1}(\mathfrak{g}, V) \\
& \subseteq\left[\mathrm{d}, \mathcal{L}_{y}\right] C^{k}(\mathfrak{g}, V)=\{0\}
\end{aligned}
$$

for each $x \in \mathfrak{g}$. Hence $\left[\mathrm{d}, \mathcal{L}_{y}\right] C^{k+1}(\mathfrak{g}, V)=\{0\}$. By induction, this leads to $\left[\mathrm{d}, \mathcal{L}_{y}\right]=0$ for each $y \in \mathfrak{g}$.

Proposition 7.5.12. $\mathrm{d}^{2}=0$.\\
Proof. We put Lemma 7.5.11 into the Cartan Formula (7.5) and get

$$
0=\left[\mathrm{d}, \mathcal{L}_{x}\right]=\mathrm{d}^{2} \circ i_{x}-i_{x} \circ \mathrm{~d}^{2}
$$

We use this formula to show by induction over $k$ that $\mathrm{d}^{2}$ vanishes on $C^{k}(\mathfrak{g}, V)$. For $\omega \in C^{0}(\mathfrak{g}, V) \cong V$, we have $\mathrm{d} \omega(x)=x \cdot \omega$ and\\
$\mathrm{d}^{2} \omega(x, y)=x \cdot \mathrm{~d} \omega(y)-y \cdot \mathrm{~d} \omega(x)-(\mathrm{d} \omega)([x, y])=x \cdot(y \cdot \omega)-y \cdot(x \cdot \omega)-[x, y] \cdot \omega=0$.\\
If $\mathrm{d}^{2}\left(C^{k}(\mathfrak{g}, V)\right)=\{0\}$, we use (7.7) to see that

$$
i_{x} \mathrm{~d}^{2}\left(C^{k+1}(\mathfrak{g}, V)\right)=\mathrm{d}^{2} i_{x} C^{k+1}(\mathfrak{g}, V) \subseteq \mathrm{d}^{2}\left(C^{k}(\mathfrak{g}, V)\right)=\{0\}
$$

for all $x \in \mathfrak{g}$, and hence that $\mathrm{d}^{2}\left(C^{k+1}(\mathfrak{g}, V)\right)=\{0\}$. By induction on $k$, this proves $\mathrm{d}^{2}=0$.

Since the differential commutes with the action of $\mathfrak{g}$ on the graded vector space $C(\mathfrak{g}, V)$ (Lemma 7.5.11), the space of $k$-cocycles and of $k$-coboundaries is $\mathfrak{g}$-invariant, so that we obtain a natural representation of $\mathfrak{g}$ on the quotient spaces $H^{k}(\mathfrak{g}, V)$.

Lemma 7.5.13. The action of $\mathfrak{g}$ on $H^{k}(\mathfrak{g}, V)$ is trivial, i.e., $\mathcal{L}_{\mathfrak{g}} Z^{k}(\mathfrak{g}, V) \subseteq B^{k}(\mathfrak{g}, V)$.

Proof. In view of Lemma 7.5.9, we have for $\omega \in Z^{k}(\mathfrak{g}, V)$ the relation

$$
\mathcal{L}_{x} \omega=i_{x} \mathrm{~d} \omega+\mathrm{d}\left(i_{x} \omega\right)=\mathrm{d}\left(i_{x} \omega\right) \in B^{k}(\mathfrak{g}, V)
$$

Hence the $\mathfrak{g}$-action induced on the cohomology space $H^{k}(\mathfrak{g}, V)$ is trivial.\\
Remark 7.5.14. Let $M$ be a smooth manifold and $\mathfrak{g}:=\mathcal{V}(M)$ the Lie algebra of smooth vector fields on $M$ (cf. Chapter 9 below). We consider the $\mathfrak{g}$ module $V:=C^{\infty}(M, \mathbb{R})$ of smooth functions on $M$. Then we can identify the space $\Omega^{p}(M, \mathbb{R})$ of alternating $C^{\infty}(M, \mathbb{R})$-multilinear maps $\mathfrak{g}^{p} \rightarrow C^{\infty}(M, \mathbb{R})$ with a subspace of $C^{p}(\mathfrak{g}, V)$. The elements of the space $\Omega^{p}(M, \mathbb{R})$ are called smooth $p$-forms on $M$, and

$$
\Omega(M, \mathbb{R}):=\bigoplus_{p \in \mathbb{N}_{0}} \Omega^{p}(M, \mathbb{R})
$$

is the space of exterior forms on $M$. The restriction of d to these spaces is called the exterior differential. The space $\Omega(M, \mathbb{R})$ is invariant under the differential d and the $\mathfrak{g}$-action given by the Lie derivative. Together with the exterior derivative, the spaces $\Omega^{p}(M, \mathbb{R})$ now form the so called de Rham complex

$$
\Omega^{0}(M, \mathbb{R})=C^{\infty}(M, \mathbb{R}) \xrightarrow{\mathrm{d}} \Omega^{1}(M, \mathbb{R}) \xrightarrow{\mathrm{d}} \Omega^{2}(M, \mathbb{R}) \xrightarrow{\mathrm{d}} \cdots
$$

The cohomology groups of this subcomplex are the de Rham cohomology groups of $M$ :\\
$H_{\mathrm{dR}}^{0}(M, \mathbb{R}):=\operatorname{ker}\left(\left.\mathrm{d}\right|_{\Omega^{0}(M, \mathbb{R})}\right), \quad H_{\mathrm{dR}}^{p}(M, \mathbb{R}):=\frac{\operatorname{ker}\left(\left.\mathrm{d}\right|_{\Omega^{p}(M, \mathbb{R})}\right)}{\mathrm{d}\left(\Omega^{p-1}(M, \mathbb{R})\right)} \quad$ for $\quad p>0$.

\subsubsection*{Extensions and Cocycles}
In this section, we interpret the cohomology spaces in low degrees in terms of extensions of modules and Lie algebras.\\
Definition 7.5.15. (a) Each element $\omega \in Z^{1}(\mathfrak{g}, V)$ defines a $\mathfrak{g}$-module

$$
V_{\omega}:=V \times \mathbb{K} \quad \text { with } \quad x \cdot(v, t)=(x \cdot v+t \omega(x), 0)
$$

Then the inclusion $j: V \hookrightarrow V_{\omega}$ is a module homomorphism, and we thus obtain the short exact sequence

$$
\mathbf{0} \rightarrow V \hookrightarrow V_{\omega} \rightarrow \mathbb{K} \rightarrow \mathbf{0}
$$

of $\mathfrak{g}$-modules, where $\mathbb{K}$ is considered as a trivial $\mathfrak{g}$-module.\\
(b) In the following, we write $\mathfrak{a f f}(V):=V \rtimes \mathfrak{g l}(V)$ for the affine Lie algebra of $V$ with the Lie bracket

$$
\left[(v, A),\left(v^{\prime}, A^{\prime}\right)\right]=\left(A \cdot v^{\prime}-A^{\prime} \cdot v,\left[A, A^{\prime}\right]\right)
$$

An affine representation of a Lie algebra $\mathfrak{g}$ on $V$ is identified with a homomorphism $\pi: \mathfrak{g} \rightarrow \mathfrak{a f f}(V)$. We associate with each pair $(v, A) \in \mathfrak{a f f}(V)$ the affine map $w \mapsto A w+v$. The Lie algebra $\mathfrak{a f f}(V)$ acts linearly on the space $V \times \mathbb{K}$ by $(v, A) \cdot(w, t):=(A w+t v, 0)$.\\
Proposition 7.5.16. Let ( $\rho, V$ ) be a $\mathfrak{g}$-module. An element $\omega \in C^{1}(\mathfrak{g}, V)$ is in $Z^{1}(\mathfrak{g}, V)$ if and only if the map

$$
\rho_{\omega}: \mathfrak{g} \rightarrow \mathfrak{a f f}(V) \cong V \rtimes \mathfrak{g} \mathfrak{l}(V), \quad x \mapsto(\omega(x), \rho(x))
$$

is a homomorphism of Lie algebras.\\
Let $e^{\operatorname{ad} V}:=\mathbf{1}+\operatorname{ad} V \subseteq \operatorname{Aut}(\mathfrak{a f f}(V))$ denote the group of automorphisms defined by the abelian ideal $V \unlhd \mathfrak{a f f}(V)$. Then the space $H^{1}(\mathfrak{g}, V)$ parameterizes the $e^{\operatorname{ad} V}$-conjugacy classes of affine representations of $\mathfrak{g}$ on $V$ whose corresponding linear representation is $\rho$. The coboundaries correspond to those affine representations which are conjugate to a linear representation, i.e., which have a fixed point $p \in V$ in the sense that $\rho(x) p+\omega(x)=0$ for all $x \in \mathfrak{g}$.

Proof. The first assertion is easily checked. For $v \in V$, we consider the automorphism of $\mathfrak{a f f}(V)$ given by $\eta_{v}=e^{\mathrm{ad} v}:=\mathbf{1}+\operatorname{ad} v$. Then $\eta_{v}(w, A)=$ ( $w-A \cdot v, A$ ), so that

$$
\eta_{v} \circ \rho_{\omega}=\rho_{\omega-\mathrm{d} v}
$$

where $(\mathrm{d} v)(x)=x \cdot v$. Thus two affine representations $\rho_{\omega}$ and $\rho_{\omega^{\prime}}$ are conjugate under some $\eta_{v}$ if and only if the cohomology classes of $\omega$ and $\omega^{\prime}$ coincide. In this sense, $H^{1}(\mathfrak{g}, V)$ parameterizes the $e^{\mathrm{ad} V}$-conjugacy classes of affine representations of $\mathfrak{g}$ on $V$ whose corresponding linear representation coincides with $\rho$. The coboundaries correspond to those affine representations which are conjugate to a linear representation. Moreover, it is clear that an affine representation $\rho_{\omega}$ is conjugate to a linear representation, if and only if there exists a fixed point $v \in V$, i.e., $\omega=-\mathrm{d} v$. \(\square\)

Definition 7.5.17. (a) Let $\mathfrak{g}$ and $\mathfrak{n}$ be Lie algebras. The short exact sequence

$$
\mathbf{0} \rightarrow \mathfrak{n} \xrightarrow{\iota} \widehat{\mathfrak{g}} \xrightarrow{q} \mathfrak{g} \rightarrow \mathbf{0}
$$

(this means $\iota$ injective, $q$ surjective, and $\operatorname{im} \iota=\operatorname{ker} q$ ) is called an extension of $\mathfrak{g}$ by $\mathfrak{n}$. If we identify $\mathfrak{n}$ with its image in $\widehat{\mathfrak{g}}$, this means that $\widehat{\mathfrak{g}}$ is a Lie algebra containing $\mathfrak{n}$ as an ideal such that $\widehat{\mathfrak{g}} / \mathfrak{n} \cong \mathfrak{g}$. If $\mathfrak{n}$ is abelian (central) in $\widehat{\mathfrak{g}}$, then the extension is called abelian (central). Two extensions $\mathfrak{n} \hookrightarrow \widehat{\mathfrak{g}}_{1} \rightarrow \mathfrak{g}$ and $\mathfrak{n} \hookrightarrow \widehat{\mathfrak{g}}_{2} \rightarrow \mathfrak{g}$ are called equivalent if there exists a Lie algebra homomorphism $\varphi: \widehat{\mathfrak{g}}_{1} \rightarrow \widehat{\mathfrak{g}}_{2}$ such that the diagram\\
\includegraphics[max width=\textwidth, center]{2025_10_18_3eaf04e77f3cb364fce5g-212}\\
is commutative. It is easy to see that this implies that $\varphi$ is an isomorphism of Lie algebras (Exercise).\\
(b) We call an extension $q: \widehat{\mathfrak{g}} \rightarrow \mathfrak{g}$ with $\operatorname{ker} q=\mathfrak{n}$ trivial, or say that the extension splits, if there exists a Lie algebra homomorphism $\sigma: \mathfrak{g} \rightarrow \widehat{\mathfrak{g}}$ with $q \circ \sigma=\mathrm{id}_{\mathfrak{g}}$. In this case, the map

$$
\mathfrak{n} \rtimes_{\delta} \mathfrak{g} \rightarrow \widehat{\mathfrak{g}}, \quad(n, x) \mapsto n+\sigma(x)
$$

is an isomorphism, where the semidirect sum is defined by the homomorphism

$$
\delta: \mathfrak{g} \rightarrow \operatorname{der}(\mathfrak{n}), \quad \delta(x)(n):=[\sigma(x), n] .
$$

For a trivial central extension, we have $\delta=0$, and therefore $\widehat{\mathfrak{g}} \cong \mathfrak{n} \times \mathfrak{g}$.\\
(c) A particular important case arises if $\mathfrak{n}$ is abelian. Then each Lie algebra extension $q: \widehat{\mathfrak{g}} \rightarrow \mathfrak{g}$ of $\mathfrak{g}$ by $\mathfrak{n}$ leads to a $\mathfrak{g}$-module structure on $\mathfrak{n}$ defined by $q(x) \cdot n:=[x, n]$, which is well defined because $[\mathfrak{n}, \mathfrak{n}]=\{0\}$. It is easy to see that equivalent extensions lead to the same module structure (Exercise). Therefore, it makes sense to write $\operatorname{Ext}_{\rho}(\mathfrak{g}, \mathfrak{n})$ for the set of equivalence classes\\
of extensions of $\mathfrak{g}$ by $\mathfrak{n}$ corresponding to the module structure given by the representation

$$
\rho: \mathfrak{g} \rightarrow \mathfrak{g l}(\mathfrak{n})=\operatorname{der}(\mathfrak{n})
$$

For a $\mathfrak{g}$-module $V$, we also write $\operatorname{Ext}(\mathfrak{g}, V):=\operatorname{Ext}_{\rho_{V}}(\mathfrak{g}, V)$, where $\rho_{V}$ is the representation of $\mathfrak{g}$ on $V$ corresponding to the module structure.

Proposition 7.5.18. For an element $\omega \in C^{2}(\mathfrak{g}, V)$, the formula

$$
\left[(v, x),\left(v^{\prime}, x^{\prime}\right)\right]=\left(x \cdot v^{\prime}-x^{\prime} \cdot v+\omega\left(x, x^{\prime}\right),\left[x, x^{\prime}\right]\right)
$$

defines a Lie bracket on $V \times \mathfrak{g}$ if and only if $\omega \in Z^{2}(\mathfrak{g}, V)$. For a cocycle $\omega \in Z^{2}(\mathfrak{g}, V)$, we write $\mathfrak{g}_{\omega}:=V \oplus_{\omega} \mathfrak{g}$ for the corresponding Lie algebra. Then we obtain for each cocycle $\omega$ an extension of $\mathfrak{g}$ by the abelian ideal $V$ :

$$
\mathbf{0} \rightarrow V \hookrightarrow \mathfrak{g}_{\omega} \rightarrow \mathfrak{g} \rightarrow \mathbf{0}
$$

This extension splits if and only if $\omega$ is a coboundary.\\
The map $Z^{2}(\mathfrak{g}, V) \rightarrow \operatorname{Ext}(\mathfrak{g}, V)$ defined by assigning to $\omega$ the equivalence class of the extension $\mathfrak{g}_{\omega}$ induces a bijection

$$
H^{2}(\mathfrak{g}, V) \rightarrow \operatorname{Ext}(\mathfrak{g}, V)
$$

Therefore, $H^{2}(\mathfrak{g}, V)$ classifies the abelian extensions of $\mathfrak{g}$ by $V$ for which the corresponding representation of $\mathfrak{g}$ on $V$ is given by the module structure on $V$.

Proof. An easy calculation shows that $\mathfrak{g}_{\omega}=V \oplus_{\omega} \mathfrak{g}$ is a Lie algebra if and only if $\omega$ is a 2-cocycle, i.e., an element of $Z^{2}(\mathfrak{g}, V)$.

To see that every abelian Lie algebra extension $q: \widehat{\mathfrak{g}} \rightarrow \mathfrak{g}$ with $\operatorname{ker} q=V$ (as a $\mathfrak{g}$-module) is equivalent to some $\mathfrak{g}_{\omega}$, let $\sigma: \mathfrak{g} \rightarrow \widehat{\mathfrak{g}}$ be a linear map with $q \circ \sigma=\mathrm{id}_{\mathfrak{g}}$. Then the map

$$
V \times \mathfrak{g} \rightarrow \widehat{\mathfrak{g}}, \quad(v, x) \mapsto v+\sigma(x)
$$

is a bijection, and it becomes an isomorphism of Lie algebras if we endow $V \times \mathfrak{g}$ with the bracket of $\mathfrak{g}_{\omega}$ for

$$
\omega(x, y):=[\sigma(x), \sigma(y)]-\sigma([x, y])
$$

This implies that $q: \widehat{\mathfrak{g}} \rightarrow \mathfrak{g}$ is equivalent to $\mathfrak{g}_{\omega}$, and therefore that the map $Z^{2}(\mathfrak{g}, V) \rightarrow \operatorname{Ext}(\mathfrak{g}, V)$ is surjective.

Two Lie algebras $\mathfrak{g}_{\omega}$ and $\mathfrak{g}_{\omega^{\prime}}$ are equivalent as $V$-extensions of $\mathfrak{g}$ if and only if there exists a linear map $\varphi: \mathfrak{g} \rightarrow V$ such that the map

$$
\widetilde{\varphi}: \mathfrak{g}_{\omega}=V \times \mathfrak{g} \rightarrow \mathfrak{g}_{\omega^{\prime}}=V \times \mathfrak{g}, \quad(v, x) \mapsto(v+\varphi(x), x)
$$

is a Lie algebra homomorphism. This means that

$$
\begin{aligned}
\widetilde{\varphi}\left(\left[(v, x),\left(v^{\prime}, x^{\prime}\right)\right]\right) & =\widetilde{\varphi}\left(x \cdot v^{\prime}-x^{\prime} \cdot v+\omega\left(x, x^{\prime}\right),\left[x, x^{\prime}\right]\right) \\
& =\left(x \cdot v^{\prime}-x^{\prime} \cdot v+\omega\left(x, x^{\prime}\right)+\varphi\left(\left[x, x^{\prime}\right]\right),\left[x, x^{\prime}\right]\right)
\end{aligned}
$$

equals

$$
\begin{aligned}
{\left[\widetilde{\varphi}(v, x), \widetilde{\varphi}\left(v^{\prime}, x^{\prime}\right)\right] } & =\left(x \cdot\left(v^{\prime}+\varphi\left(x^{\prime}\right)\right)-x^{\prime} \cdot(v+\varphi(x))+\omega^{\prime}\left(x, x^{\prime}\right),\left[x, x^{\prime}\right]\right) \\
& =\left(x \cdot v^{\prime}-x^{\prime} \cdot v+x \cdot \varphi\left(x^{\prime}\right)-x^{\prime} \cdot \varphi(x)+\omega^{\prime}\left(x, x^{\prime}\right),\left[x, x^{\prime}\right]\right)
\end{aligned}
$$

which is equivalent to

$$
\omega^{\prime}\left(x, x^{\prime}\right)-\omega\left(x, x^{\prime}\right)=\varphi\left(\left[x, x^{\prime}\right]\right)-x \cdot \varphi\left(x^{\prime}\right)+x^{\prime} \cdot \varphi(x)=-(\mathrm{d} \varphi)\left(x, x^{\prime}\right)
$$

Therefore, $\mathfrak{g}_{\omega}$ and $\mathfrak{g}_{\omega^{\prime}}$ are equivalent abelian extensions of $\mathfrak{g}$ if and only if $\omega^{\prime}-\omega$ is a coboundary. Hence the map

$$
Z^{2}(\mathfrak{g}, V) \rightarrow \operatorname{Ext}(\mathfrak{g}, V), \quad \omega \mapsto\left[\mathfrak{g}_{\omega}\right]
$$

induces a bijection $H^{2}(\mathfrak{g}, V) \rightarrow \operatorname{Ext}(\mathfrak{g}, V)$.\\
Example 7.5.19. If $\mathfrak{g}=\mathbb{K}$ is the one-dimensional Lie algebra, then the $\mathfrak{g}$ module structures on a vector space $V$ are in one-to-one correspondence with endomorphisms $D \in \operatorname{End}(V)$.

In this case, $C^{p}(\mathfrak{g}, V)$ vanishes for $p>1$, and we have

$$
C^{0}(\mathfrak{g}, V)=V \quad \text { and } \quad C^{1}(\mathfrak{g}, V) \cong V .
$$

With respect to this identification, the map d corresponds to the endomorphism $D \in \operatorname{End}(V)$, so that

$$
H^{0}(\mathfrak{g}, V)=V^{\mathfrak{g}}=\operatorname{ker} D \quad \text { and } \quad H^{1}(\mathfrak{g}, V) \cong V / D(V)=\operatorname{coker}(D)
$$

Example 7.5.20. Let $\mathfrak{g}:=\mathbb{K}^{2}$ be the abelian two-dimensional Lie algebra with the canonical basis $e_{1}, e_{2}$. Then a $\mathfrak{g}$-module structure on a vector space $V$ corresponds to a pair ( $D_{1}, D_{2}$ ) of commuting endomorphisms of $V$.

We have

$$
C^{1}(\mathfrak{g}, V)=\operatorname{Hom}\left(\mathbb{K}^{2}, V\right) \cong V^{2} \quad \text { and } \quad C^{2}(\mathfrak{g}, V) \cong V
$$

where the isomorphism $C^{2}(\mathfrak{g}, V) \rightarrow V$ is given by $\beta \mapsto \beta\left(e_{1}, e_{2}\right)$. As $\operatorname{dim} \mathfrak{g}=$ 2 , we have $C^{3}(\mathfrak{g}, V)=\{0\}$, so that $C^{2}(\mathfrak{g}, V)=Z^{2}(\mathfrak{g}, V)$. For $(v, w) \in V^{2} \cong C^{1}(\mathfrak{g}, V)$ we have

$$
\mathrm{d}(v, w)\left(e_{1}, e_{2}\right)=e_{1} \cdot w-e_{2} \cdot v=D_{1}(w)-D_{2}(v)
$$

and thus

$$
H^{2}(\mathfrak{g}, V) \cong V /\left(D_{1}(V)+D_{2}(V)\right)
$$

Proposition 7.5.21. If $V=\mathfrak{g}$ with respect to the adjoint representation, then $Z^{1}(\mathfrak{g}, \mathfrak{g}) \cong \operatorname{der} \mathfrak{g}, B^{1}(\mathfrak{g}, \mathfrak{g}) \cong \operatorname{ad} \mathfrak{g}$ and

$$
H^{1}(\mathfrak{g}, \mathfrak{g}) \cong \operatorname{out}(\mathfrak{g}):=\operatorname{der} \mathfrak{g} / \operatorname{ad} \mathfrak{g}
$$

is the space of outer derivations of $\mathfrak{g}$.\\
Proof. Let $V=\mathfrak{g}$ with respect to the adjoint representation. For $c \in C^{1}(\mathfrak{g}, \mathfrak{g})=\operatorname{End}(\mathfrak{g})$, we then have

$$
\mathrm{d} c(x, y)=[x, c(y)]-[y, c(x)]-c([x, y]),
$$

showing that $Z^{1}(\mathfrak{g}, \mathfrak{g})=$ der $\mathfrak{g}$. For $c \in C^{0}(\mathfrak{g}, \mathfrak{g}) \cong \mathfrak{g}$, we have $\mathrm{d} c(x)=[x, c]$, showing that $B^{1}(\mathfrak{g}, \mathfrak{g})=\operatorname{ad} \mathfrak{g}$. \(\square\)

Definition 7.5.22. Let $\mathfrak{g}$ be a Lie algebra and $V$ and $W$ modules of $\mathfrak{g}$. The short exact sequence

$$
\mathbf{0} \rightarrow W \xrightarrow{\iota} \widehat{V} \xrightarrow{q} V \rightarrow \mathbf{0}
$$

(this means $\iota$ injective, $q$ surjective, and $\operatorname{im} \iota=\operatorname{ker} q$ ) is called an extension of $V$ by $W$. If we identify $W$ with its image in $\widehat{V}$, this means that $\widehat{V}$ is a $\mathfrak{g}$-module containing $W$ as a submodule such that $\widehat{V} / W \cong V$.

Two extensions $W \hookrightarrow \widehat{V}_{1} \rightarrow V$ and $W \hookrightarrow \widehat{V}_{2} \rightarrow V$ are called equivalent if there exists a module homomorphism $\varphi: \widehat{V}_{1} \rightarrow \widehat{V}_{2}$ such that the diagram\\
\includegraphics[max width=\textwidth, center]{2025_10_18_3eaf04e77f3cb364fce5g-215}\\
commutes. It is easy to see that this implies that $\varphi$ is an isomorphism of $\mathfrak{g}$-modules (Exercise) and that we thus obtain an equivalence relation on the class of extensions of $V$ by $W$. We write $\operatorname{Ext}(V, W)$ for the set of equivalence classes of module extensions of $V$ by $W$. We call an extension (7.9) trivial, or say that the extension splits, if there exists a module homomorphism $\sigma: V \rightarrow \widehat{V}$ with $q \circ \sigma=\operatorname{id}_{V}$. In this case, the map

$$
W \oplus V \rightarrow \widehat{V}, \quad(w, x) \mapsto \iota(w)+\sigma(x)
$$

is a module isomorphism.\\
The following proposition gives a cohomological interpretation of the set $\operatorname{Ext}(B, A)$ for two $\mathfrak{g}$-modules $B$ and $A$. In particular, it shows that this set carries a natural vector space structure.

Proposition 7.5.23. For $\mathfrak{g}$-modules ( $\rho_{A}, A$ ) and ( $\rho_{B}, B$ ),

$$
\operatorname{Ext}(B, A) \cong H^{1}(\mathfrak{g}, \operatorname{Hom}(B, A)),
$$

where the representation of $\mathfrak{g}$ on $\operatorname{Hom}(B, A)$ is given by

$$
x \cdot \varphi=\rho_{A}(x) \varphi-\varphi \rho_{B}(x) .
$$

Proof. First, we check that a module extension $q: C \rightarrow B$ by the module $A$ can be written as a space $C=A \times B$ on which the $\mathfrak{g}$-module representation is given by

$$
x \cdot(a, b)=(x \cdot a+\omega(x)(b), x \cdot b)
$$

where $\omega \in \operatorname{Hom}(\mathfrak{g}, \operatorname{Hom}(B, A))=C^{1}(\mathfrak{g}, \operatorname{Hom}(B, A))$. To see this, we simply choose a linear map $\sigma: B \rightarrow C$ with $q \circ \sigma=\operatorname{id}_{B}$ and define

$$
\omega(x)(b):=x \cdot \sigma(b)-\sigma(x \cdot b)
$$

Then the linear bijection $A \times B \rightarrow C,(a, b) \mapsto \sigma(b)+a$ is a module isomorphism with respect to the above module structure.

If $\omega \in C^{1}(\mathfrak{g}, \operatorname{Hom}(B, A))$ is given, then the condition that (7.10) defines a $\mathfrak{g}$-module structure on $C$ means that

$$
\begin{aligned}
& ([y, x] \cdot a+\omega([y, x])(b),[y, x] \cdot b) \\
= & {[y, x] \cdot(a, b) \stackrel{!}{=} y \cdot(x \cdot(a, b))-x \cdot(y \cdot(a, b)) } \\
= & (y \cdot(x \cdot a)+y \cdot \omega(x)(b)+\omega(y)(x \cdot b), y \cdot(x \cdot b)) \\
& -(x \cdot(y \cdot a)+x \cdot \omega(y)(b)+\omega(x)(y \cdot b), x \cdot(y \cdot b)) \\
= & ([y, x] \cdot a+y \cdot \omega(x)(b)+\omega(y)(x \cdot b)-x \cdot \omega(y)(b)-\omega(x)(y \cdot b),[y, x] \cdot b) \\
= & ([y, x] \cdot a+(y \cdot \omega(x))(b)-(x \cdot \omega(y))(b),[y, x] \cdot b)
\end{aligned}
$$

This is equivalent to

$$
\omega([y, x])=y \cdot \omega(x)-x \cdot \omega(y)
$$

which in turn means that $\omega \in Z^{1}(\mathfrak{g}, \operatorname{Hom}(B, A))$. The different parameterizations of $C$ as $B \times A$ correspond to linear maps $\sigma: B \rightarrow C$ with $q \circ \sigma=\operatorname{id}_{B}$, where $q: C \rightarrow B$ is the quotient map. In this sense, we have

$$
\omega_{\sigma}(x)(b)=x \cdot \sigma(b)-\sigma(x \cdot b)
$$

For a linear map $\gamma \in \operatorname{Hom}(B, A)$, we therefore have

$$
\omega_{\sigma+\gamma}(x)(b)=\omega_{\sigma}(x)(b)+(x \cdot \gamma)(b)
$$

i.e., $\omega_{\sigma+\gamma}=\omega_{\sigma}+\mathrm{d} \gamma$. We conclude that the different sections lead to cohomologous cocycles, and this observation leads to a bijection $\operatorname{Ext}(B, A) \cong H^{1}(\mathfrak{g}, \operatorname{Hom}(B, A))$.

\subsubsection*{Invariant Volume Forms and Cohomology}
In this subsection, we characterize the existence of volume forms in terms of cohomology. We refer the reader to Appendix B for background material on exterior products of alternating maps.

Proposition 7.5.24. For a finite-dimensional real Lie algebra $\mathfrak{g}$, the following are equivalent:\\
(i) $\quad \operatorname{tr}(\operatorname{ad} x)=0$ for all $x \in \mathfrak{g}$.\\
(ii) $\mathfrak{g}$ carries a volume form $\mu \in C^{n}(\mathfrak{g}, \mathbb{R})$ invariant under the adjoint action.\\
(iii) $H^{n}(\mathfrak{g}, \mathbb{R}) \neq\{0\}$ for $n=\operatorname{dim} \mathfrak{g}$.\\
(iv) $\operatorname{dim} H^{n}(\mathfrak{g}, \mathbb{R})=1$.

Proof. Let $x_{1}, \ldots, x_{n}$ be a basis for $\mathfrak{g}$ and $x_{1}^{*}, \ldots, x_{n}^{*}$ the dual basis for $\mathfrak{g}^{*}$. Then

$$
\mu=x_{1}^{*} \wedge \cdots \wedge x_{n}^{*}
$$

is a nonzero volume form on $\mathfrak{g}$, so that $C^{n}(\mathfrak{g}, \mathbb{R})=\mathbb{R} \mu$ (cf. Definition B.2.27).\\
(i) $\Leftrightarrow$ (ii): For any $x \in \mathfrak{g}$, we have

$$
\mathcal{L}_{x} \mu=-\operatorname{tr}(\operatorname{ad} x) \mu .
$$

In fact, let $\left[x, x_{j}\right]=\sum_{k=1}^{n} a_{k j} x_{k}$. Then

$$
\begin{aligned}
& \left(\mathcal{L}_{x} \mu\right)\left(x_{1}, \ldots, x_{n}\right)=-\sum_{j=1}^{n} \mu\left(x_{1}, \ldots, x_{j-1},\left[x, x_{j}\right], x_{j+1}, \ldots, x_{n}\right) \\
& =-\sum_{j=1}^{n} a_{j j} \mu\left(x_{1}, \ldots, x_{j-1}, x_{j}, x_{j+1}, \ldots, x_{n}\right)=-\sum_{j=1}^{n} a_{j j}=-\operatorname{tr}(\operatorname{ad} x)
\end{aligned}
$$

(iii) $\Leftrightarrow$ (iv) follows from $\operatorname{dim} C^{n}(\mathfrak{g}, \mathbb{R})=1$.\\
(ii) $\Leftrightarrow$ (iii): For each $j$, we have

$$
i_{x_{j}} \mu=(-1)^{j-1} x_{1}^{*} \wedge \cdots x_{j-1}^{*} \wedge x_{j+1}^{*} \cdots \wedge x_{n}^{*}
$$

and these elements form a basis for $C^{n-1}(\mathfrak{g}, \mathbb{R})$. Therefore, $\mathrm{d} \mu=0$ and the Cartan Formula imply that

$$
B^{n}(\mathfrak{g}, \mathbb{R})=\mathrm{d}\left(i_{\mathfrak{g}} \mu\right)=\mathcal{L}_{\mathfrak{g}} \mu
$$

This space vanishes if and only if $\mu$ is invariant if and only if $H^{n}(\mathfrak{g}, \mathbb{R})$ is nonzero.

For reasons that we shall understand later in our discussion of invariant measures on Lie groups, Lie algebras satisfying the equivalent conditions from the preceding proposition are called unimodular.

Example 7.5.25. (a) Each perfect real Lie algebra $\mathfrak{g}$ is unimodular. In fact,

$$
\operatorname{tr} \circ \operatorname{ad}: \mathfrak{g} \rightarrow \mathbb{R}
$$

is a homomorphism of Lie algebras, so that $\mathfrak{g}=[\mathfrak{g}, \mathfrak{g}] \subseteq \operatorname{ker}(\operatorname{tr} \circ \operatorname{ad})$. In particular, each semisimple Lie algebra is unimodular.\\
(b) Each nilpotent Lie algebra is unimodular because ad $x$ is nilpotent for each $x \in \mathfrak{g}$, which implies in particular that $\operatorname{tr}(\operatorname{ad} x)=0$.\\
(c) The 2-dimensional nonabelian Lie algebra $\mathfrak{g}=\operatorname{span}\left\{e_{1}, e_{2}\right\}$ with $\left[e_{1}, e_{2}\right]=e_{2}$ is not unimodular because $\operatorname{tr}\left(\operatorname{ad} e_{1}\right)=1 \neq 0$.

\subsubsection*{Cohomology of Semisimple Lie Algebras}
In this subsection, we discuss the connection between Weyl's and Levi's Theorems and a more general theorem of J.H.C. Whitehead concerning the cohomology of finite-dimensional modules of semisimple Lie algebras.

If $\mathfrak{s}$ is a semisimple finite-dimensional Lie algebra, then Weyl's Theorem states that every finite-dimensional $\mathfrak{g}$-module $V$ is semisimple. This means that every submodule has a module complement, and therefore that all module extensions are trivial. In view of Proposition 7.5.23, this is equivalent to

$$
H^{1}(\mathfrak{g}, \operatorname{Hom}(A, B))=\{0\}
$$

for each pair of finite-dimensional $\mathfrak{g}$-modules $A$ and $B$. For $A=\mathbb{K}$, the trivial module, we have $\operatorname{Hom}(A, B) \cong B$ as $\mathfrak{g}$-modules, and therefore we obtain $H^{1}(\mathfrak{g}, B)=\{0\}$. This argument proves the

Lemma 7.5.26 (First Whitehead Lemma). If $\mathfrak{g}$ is a finite-dimensional semisimple Lie algebra, then each finite-dimensional $\mathfrak{g}$-module $V$ satisfies

$$
H^{1}(\mathfrak{g}, V)=\{0\}
$$

From the above argument, it is easy to see that the First Whitehead Lemma says essentially the same as Weyl's Theorem.

Now let $\omega \in Z^{2}(\mathfrak{g}, V)$ be a 2-cocycle and $\mathfrak{g}_{\omega}:=V \oplus_{\omega} \mathfrak{g}$ the corresponding abelian Lie algebra extension of $\mathfrak{g}$ by $V$. Then $V=\mathfrak{r a d}\left(\mathfrak{g}_{\omega}\right)$ because $V$ is an abelian ideal of $\mathfrak{g}_{\omega}$ and the quotient $\mathfrak{g}_{\omega} / V \cong \mathfrak{g}$ is semisimple. Therefore, Levi's Theorem implies the existence of a Levi complement in $\widehat{\mathfrak{g}}_{\omega}$, which means that there exists a Lie algebra homomorphism $\sigma: \mathfrak{g} \rightarrow \mathfrak{g}_{\omega}$, splitting the exact sequence $V \hookrightarrow \mathfrak{g}_{\omega} \xrightarrow{q} \mathfrak{g}$, i.e., satisfying $q \circ \sigma=\operatorname{id}_{\mathfrak{g}}$. Now Proposition 7.5.18 implies that $\omega$ is a coboundary, which leads to the

Lemma 7.5.27 (Second Whitehead Lemma). If $\mathfrak{g}$ is a finite-dimensional semisimple Lie algebra, then each finite-dimensional $\mathfrak{g}$-module $V$ satisfies

$$
H^{2}(\mathfrak{g}, V)=\{0\}
$$

Again we see from the argument above that the Second Whitehead Lemma is essentially the case of Levi's Theorem where $\alpha: \mathfrak{g} \rightarrow \mathfrak{s}$ is a surjective homomorphism onto a semisimple Lie algebra $\mathfrak{s}$ with abelian kernel, but this was the crucial case in the proof of Levi's Theorem.

We now aim at more general results concerning also the higher degree cohomology of modules of semisimple Lie algebras.

Lemma 7.5.28. If $\omega \in C^{p}(\mathfrak{g}, V)^{\mathfrak{g}}$ is a $\mathfrak{g}$-invariant element with respect to the action $\rho$ from Lemma 7.5.8, and $\mathrm{d}^{0}$ is the Chevalley-Eilenberg differential with respect to the trivial $\mathfrak{g}$-module structure on $V$, then $\mathrm{d} \omega=-\mathrm{d}^{0} \omega$.

Proof. First, we note that for $x_{0}, x_{1}, \ldots, x_{p} \in \mathfrak{g}$, the invariance of $\omega$ implies that

$$
\begin{aligned}
& x_{i} \cdot \omega\left(x_{0}, \ldots, \widehat{x}_{i}, \ldots, x_{p}\right) \\
= & \omega\left(\left[x_{i}, x_{0}\right], x_{1}, \ldots, \widehat{x}_{i}, \ldots, x_{p}\right)+\cdots+\omega\left(x_{0}, \ldots, \widehat{x}_{i}, \ldots,\left[x_{i}, x_{p}\right]\right) \\
= & \sum_{j=0}^{i-1}(-1)^{j} \omega\left(\left[x_{i}, x_{j}\right], \ldots, \widehat{x}_{j}, \ldots, \widehat{x}_{i}, \ldots, x_{p}\right) \\
& +\sum_{j=i+1}^{p}(-1)^{j+1} \omega\left(\left[x_{i}, x_{j}\right], \ldots, \widehat{x}_{i}, \ldots, \widehat{x}_{j}, \ldots, x_{p}\right)
\end{aligned}
$$

This leads to

$$
\begin{aligned}
& \sum_{i=0}^{p}(-1)^{i} x_{i} \cdot \omega\left(x_{0}, \ldots, \widehat{x}_{i}, \ldots, x_{p}\right) \\
= & \sum_{j<i}(-1)^{j+i} \omega\left(\left[x_{i}, x_{j}\right], \ldots, \widehat{x}_{j}, \ldots, \widehat{x}_{i}, \ldots, x_{p}\right) \\
& +\sum_{j>i}(-1)^{i+j+1} \omega\left(\left[x_{i}, x_{j}\right], \ldots, \widehat{x}_{i}, \ldots, \widehat{x}_{j}, \ldots, x_{p}\right) \\
= & \sum_{j<i}(-1)^{j+i} \omega\left(\left[x_{i}, x_{j}\right], \ldots, \widehat{x}_{j}, \ldots, \widehat{x}_{i}, \ldots, x_{p}\right) \\
& +\sum_{j<i}(-1)^{i+j+1} \omega\left(\left[x_{j}, x_{i}\right], \ldots, \widehat{x}_{j}, \ldots, \widehat{x}_{i}, \ldots, x_{p}\right) \\
= & \sum_{j<i}(-1)^{j+i} \omega\left(\left[x_{i}, x_{j}\right], \ldots, \widehat{x}_{j}, \ldots, \widehat{x}_{i}, \ldots, x_{p}\right) \\
& +\sum_{j<i}(-1)^{i+j} \omega\left(\left[x_{i}, x_{j}\right], \ldots, \widehat{x}_{j}, \ldots, \widehat{x}_{i}, \ldots, x_{p}\right) \\
= & 2 \sum_{j<i}(-1)^{i+j} \omega\left(\left[x_{i}, x_{j}\right], \ldots, \widehat{x}_{j}, \ldots, \widehat{x}_{i}, \ldots, x_{p}\right) \\
= & -2\left(\mathrm{~d}^{0} \omega\right)\left(x_{0}, \ldots, x_{p}\right) .
\end{aligned}
$$

We conclude that $\mathrm{d} \omega=-2 \mathrm{~d}^{0} \omega+\mathrm{d}^{0} \omega=-\mathrm{d}^{0} \omega$.\\
Lemma 7.5.29. If $V$ is a finite-dimensional module over the semisimple Lie algebra $\mathfrak{g}$, then

$$
H^{p}(\mathfrak{g}, V) \cong Z^{p}(\mathfrak{g}, V)^{\mathfrak{g}} / B^{p}(\mathfrak{g}, V)^{\mathfrak{g}}
$$

i.e., each cohomology class can be represented by an invariant cocycle.

Proof. From Proposition 5.5.17, we obtain for each $p \in \mathbb{N}_{0}$ the decomposition

$$
Z^{p}(\mathfrak{g}, V)=Z^{p}(\mathfrak{g}, V)^{\mathfrak{g}} \oplus \mathfrak{g} \cdot Z^{p}(\mathfrak{g}, V)
$$

Since $\mathfrak{g}$ acts trivially on the quotient space $H^{p}(\mathfrak{g}, V)$, we have $\mathfrak{g} \cdot Z^{p}(\mathfrak{g}, V) \subseteq B^{p}(\mathfrak{g}, V)$, and the assertion follows.

Proposition 7.5.30. If $V$ is a trivial $\mathfrak{g}$-module, then

$$
H^{p}(\mathfrak{g}, V) \cong Z^{p}(\mathfrak{g}, V)^{\mathfrak{g}}=C^{p}(\mathfrak{g}, V)^{\mathfrak{g}} \quad \text { for each } p \in \mathbb{N}_{0}
$$

Proof. The preceding Lemma 7.5.29 shows that each cohomology class is represented by an invariant cocycle. Further, Lemma 7.5.28 and the triviality of $V$ implies that each $\mathfrak{g}$-invariant cochain $\omega$ satisfies $\mathrm{d}^{0} \omega=\mathrm{d} \omega=-\mathrm{d}^{0} \omega$, so that $\omega$ is a cocycle. This proves $C^{p}(\mathfrak{g}, V)^{\mathfrak{g}}=Z^{p}(\mathfrak{g}, V)^{\mathfrak{g}}$.

To see that $B^{p}(\mathfrak{g}, V)^{\mathfrak{g}}$ vanishes, we note that we have the direct sum decomposition

$$
C^{p-1}(\mathfrak{g}, V)=C^{p-1}(\mathfrak{g}, V)^{\mathfrak{g}} \oplus \mathfrak{g} \cdot C^{p-1}(\mathfrak{g}, V)
$$

and likewise

$$
B^{p}(\mathfrak{g}, V)=B^{p}(\mathfrak{g}, V)^{\mathfrak{g}} \oplus \mathfrak{g} \cdot B^{p}(\mathfrak{g}, V)
$$

so that

$$
B^{p}(\mathfrak{g}, V)^{\mathfrak{g}}=\mathrm{d}\left(C^{p-1}(\mathfrak{g}, V)^{\mathfrak{g}}\right)=\{0\}
$$

since Lemma 7.5.11 shows that $\mathrm{d}\left(\mathfrak{g} \cdot C^{p-1}(\mathfrak{g}, V)\right) \subseteq \mathfrak{g} \cdot \mathrm{d}\left(C^{p-1}(\mathfrak{g}, V)\right)$. Therefore, $H^{p}(\mathfrak{g}, V)=Z^{p}(\mathfrak{g}, V)$, and the proof is complete.

Example 7.5.31. For each semisimple Lie algebra $\mathfrak{g}$, we have

$$
H^{3}(\mathfrak{g}, \mathbb{K}) \neq\{0\}
$$

In fact, the invariance of the Cartan-Killing form $\kappa$ implies that

$$
\Gamma(\kappa)(x, y, z):=\kappa([x, y], z)
$$

is an invariant alternating 3 -form on $\mathfrak{g}$, hence a nonzero element of $C^{3}(\mathfrak{g}, \mathbb{K})^{\mathfrak{g}} \cong H^{3}(\mathfrak{g}, \mathbb{K})$.

Lemma 7.5.32. If $V$ is a real module over the real Lie algebra $\mathfrak{g}$ and $V_{\mathbb{C}}$ the corresponding complex $\mathfrak{g}_{\mathbb{C}}$-module, then

$$
H^{p}(\mathfrak{g}, V)_{\mathbb{C}} \cong H^{p}\left(\mathfrak{g}_{\mathbb{C}}, V_{\mathbb{C}}\right) \quad \text { for } \quad p \in \mathbb{N}_{0}
$$

Proof. Since any alternating $p$-form $\omega \in C^{p}(\mathfrak{g}, V)$ extends uniquely to a complex $p$-linear map $\omega_{\mathbb{C}} \in C^{p}\left(\mathfrak{g}_{\mathbb{C}}, V_{\mathbb{C}}\right)$, we have an embedding

$$
C^{p}(\mathfrak{g}, V) \hookrightarrow C^{p}\left(\mathfrak{g}_{\mathbb{C}}, V_{\mathbb{C}}\right)
$$

The image of this map is the set of all elements $\eta \in C^{p}\left(\mathfrak{g}_{\mathbb{C}}, V_{\mathbb{C}}\right)$, whose values on $\mathfrak{g}^{p}$ lie in the real subspace $V \cong 1 \otimes V \subseteq V_{\mathbb{C}}$. Since any $\eta \in C^{p}\left(\mathfrak{g}_{\mathbb{C}}, V_{\mathbb{C}}\right)$ can be written in a unique fashion as a $\eta_{1}+i \eta_{2}$ with $\eta_{j}\left(\mathfrak{g}^{p}\right) \subseteq V$, we see that

$$
C^{p}\left(\mathfrak{g}_{\mathbb{C}}, V_{\mathbb{C}}\right) \cong C^{p}(\mathfrak{g}, V)_{\mathbb{C}} .
$$

The subspaces of cocycles and coboundaries inherit the corresponding property, and from that the assertion follows.

Theorem 7.5.33 (Whitehead's Vanishing Theorem). If $V$ is a finitedimensional module over the semisimple Lie algebra $\mathfrak{g}$ with $V^{\mathfrak{g}}=\{0\}$, then $H^{p}(\mathfrak{g}, V)=\{0\}$ for any $p \in \mathbb{N}_{0}$.

Proof. In view of Weyl's Theorem, the condition $V^{\mathfrak{g}}=\{0\}$ implies that $V$ is a direct sum of simple nontrivial $\mathfrak{g}$-modules $V=\bigoplus_{i=1}^{k} V_{i}$. Then

$$
H^{p}(\mathfrak{g}, V) \cong \bigoplus_{i=1}^{k} H^{p}\left(\mathfrak{g}, V_{i}\right)
$$

(Exercise), so that we may w.lo.g. assume that $V$ is a simple nontrivial module. By Lemma 7.5.32, we may further assume that $\mathfrak{g}$ and $V$ are complex.

In view of Lemma 7.5.29, we have to show that $Z^{p}(\mathfrak{g}, V)^{\mathfrak{g}} \subseteq B^{p}(\mathfrak{g}, V)^{\mathfrak{g}}$. For $p=0$, this follows from our assumption $Z^{0}(\mathfrak{g}, V)=V^{\mathfrak{g}}=\{0\}$. Let $x_{1}, \ldots, x_{m}$ be a basis for $\mathfrak{g}$ and $x^{1}, \ldots, x^{m}$ the dual basis with respect to the Cartan-Killing form. For $p>0$, we then define a linear map

$$
\Gamma: C^{p}(\mathfrak{g}, V) \rightarrow C^{p-1}(\mathfrak{g}, V), \quad \Gamma(\omega):=\sum_{j=1}^{m} \rho_{V}\left(x_{j}\right) \circ i\left(x^{j}\right) \omega .
$$

First, we show that $\Gamma$ is $\mathfrak{g}$-equivariant. For $x \in \mathfrak{g}$, we have

$$
\begin{aligned}
x \cdot \Gamma(\omega) & =\sum_{j=1}^{m} x \cdot\left(\rho_{V}\left(x_{j}\right) \circ i\left(x^{j}\right) \omega\right) \\
& =\sum_{j=1}^{m} \rho_{V}\left(\left[x, x_{j}\right]\right) \circ i\left(x^{j}\right) \omega+\rho_{V}\left(x_{j}\right) \circ i\left(\left[x, x^{j}\right]\right) \omega+\Gamma(x \cdot \omega) .
\end{aligned}
$$

If

$$
\operatorname{ad} x\left(x_{j}\right)=\sum_{k=1}^{n} a_{k j} x_{k} \quad \text { and } \quad \operatorname{ad} x\left(x^{j}\right)=\sum_{k=1}^{m} a^{k j} x^{k}
$$

then $a_{k j}=\kappa\left(\left[x, x_{j}\right], x^{k}\right)=-\kappa\left(x_{j},\left[x, x^{k}\right]\right)=-a^{j k}$, and this leads to

$$
\begin{aligned}
& \sum_{j=1}^{m} \rho_{V}\left(\left[x, x_{j}\right]\right) \circ i\left(x^{j}\right) \omega+\rho_{V}\left(x_{j}\right) \circ i\left(\left[x, x^{j}\right]\right) \omega \\
& =\sum_{j, k=1}^{m} a_{k j} \rho_{V}\left(x_{k}\right) \circ i\left(x^{j}\right) \omega+a^{k j} \rho_{V}\left(x_{j}\right) \circ i\left(x^{k}\right) \omega \\
& =\sum_{j, k=1}^{m} a_{k j} \rho_{V}\left(x_{k}\right) \circ i\left(x^{j}\right) \omega+a^{j k} \rho_{V}\left(x_{k}\right) \circ i\left(x^{j}\right) \omega=0 .
\end{aligned}
$$

We conclude that $\Gamma$ is $\mathfrak{g}$-equivariant, hence maps $\mathfrak{g}$-invariant cochains to $\mathfrak{g}$ invariant cochains.

For $\omega \in Z^{p}(\mathfrak{g}, V)^{\mathfrak{g}}$, we thus have $\Gamma(\omega) \in C^{p-1}(\mathfrak{g}, V)^{\mathfrak{g}}$. With Lemma 7.5.28, we obtain $\mathrm{d}^{0} \omega=-\mathrm{d} \omega=0$. We also note that the representation $\mathcal{L}$ of $\mathfrak{g}$ on $C(\mathfrak{g}, V)$ can be written as $\mathcal{L}=\rho_{0}+\rho_{+}$, where $\rho_{0}$ is the representation corresponding to the trivial $\mathfrak{g}$-module structure on $V$. Since $\omega$ is $\mathcal{L}_{\mathfrak{g}}$-invariant, we therefore have

$$
\rho_{0}(x) \omega=-\rho_{V}(x) \circ \omega .
$$

Next we calculate with the Cartan formulas for the trivial $\mathfrak{g}$-module $V$ and the corresponding representation $\rho_{0}$ of $\mathfrak{g}$ on $C(\mathfrak{g}, V)$ :

$$
\begin{aligned}
& \mathrm{d}(\Gamma(\omega))=-\mathrm{d}^{0}(\Gamma(\omega))=-\sum_{j=1}^{m} \rho_{V}\left(x_{j}\right) \circ \mathrm{d}^{0}\left(i_{x^{j}} \omega\right) \\
& =-\sum_{j=1}^{m} \rho_{V}\left(x_{j}\right) \circ\left(\rho_{0}\left(x^{j}\right) \omega-i_{x^{j}} \mathrm{~d}^{0} \omega\right)=\sum_{j=1}^{m} \rho_{V}\left(x_{j}\right) \rho_{V}\left(x^{j}\right) \circ \omega \\
& =\rho_{V}\left(C_{\mathfrak{g}}\right) \circ \omega
\end{aligned}
$$

where $C_{\mathfrak{g}} \in \mathcal{U}(\mathfrak{g})$ is the universal Casimir operator of $\mathfrak{g}$. Since $V$ is a simple $\mathfrak{g}$-module, $\rho_{V}\left(C_{\mathfrak{g}}\right)$ is a nonzero multiple of $\operatorname{id}_{V}$ (Lemma 7.3.17), so that $\omega$ is a coboundary. \(\square\)

\subsubsection*{Cohomology of Nilpotent Lie Algebras}
We conclude this section with a short subsection in which we show that a certain vanishing result for the cohomology of nilpotent Lie algebras implies that all Cartan subalgebras of a solvable Lie algebra are conjugate.

Proposition 7.5.34. Let $V$ be a finite-dimensional module over the nilpotent Lie algebra $\mathfrak{h}$. If $V^{0}(\mathfrak{h})=\{0\}$, then

$$
H^{p}(\mathfrak{h}, V)=\{0\} \quad \text { for } \quad p \in \mathbb{N}_{0} .
$$

Proof. If $V$ and $\mathfrak{h}$ are real, then $V_{\mathbb{C}}$ carries a natural $\mathfrak{h}_{\mathbb{C}}$-module structure with $V^{0}\left(\mathfrak{h}_{\mathbb{C}}\right)=V^{0}(\mathfrak{h})_{\mathbb{C}}=\{0\}$ (Exercise). In view of Lemma 7.5.32, we may\\
therefore assume that $V$ and $\mathfrak{h}$ are complex, so that our assumption implies that

$$
V=\bigoplus_{0 \neq \lambda \in \mathfrak{h}^{*}} V^{\lambda}(\mathfrak{h})
$$

(Lemma 6.1.3). We therefore have

$$
H^{p}(\mathfrak{h}, V)=\bigoplus_{0 \neq \lambda \in \mathfrak{h}^{*}} H^{p}\left(\mathfrak{h}, V^{\lambda}(\mathfrak{h})\right)
$$

so that we may assume that $V=V^{\lambda}(\mathfrak{h})$ for some nonzero $\lambda \in \mathfrak{h}^{*}$.\\
Then the action of $\mathfrak{h}$ on $C^{p}(\mathfrak{g}, V)$ is given by

$$
\mathcal{L}_{x} \omega=\rho_{V}(x) \circ \omega+\rho_{0}(x) \omega=\rho_{+}(x) \omega+\rho_{0}(x) \omega
$$

where $\rho_{0}$ is the action corresponding to the trivial module structure on $V$. The operator $\rho_{0}(x)$ is on the larger space $\operatorname{Mult}^{p}(\mathfrak{g}, V)$ a sum of $p+1$ pairwise commuting linear maps, all of which are nilpotent (cf. Lemma 7.5.8). These in turn commute with $\rho_{+}(x)$, which is a sum of $\lambda(x) \mathbf{1}$ and the nilpotent linear map $\rho_{V}(x)-\lambda(x) \mathbf{1}$. This implies that

$$
C^{p}(\mathfrak{h}, V)=C^{p}(\mathfrak{h}, V)^{\lambda}(\mathfrak{h})
$$

under the action of $\mathfrak{h}$. We therefore have

$$
Z^{p}(\mathfrak{h}, V)=Z^{p}(\mathfrak{h}, V)^{\lambda}(\mathfrak{h}) \subseteq \operatorname{span}\left(\mathcal{L}_{\mathfrak{h}} Z^{p}(\mathfrak{h}, V)\right) \subseteq B^{p}(\mathfrak{h}, V)
$$

which implies that $H^{p}(\mathfrak{h}, V)$ vanishes.\\
Theorem 7.5.35 (Conjugacy Theorem for Cartan Subalgebras of Solvable Lie Algebras). Let $\mathfrak{h}$ and $\mathfrak{h}^{\prime}$ be Cartan subalgebras of the solvable Lie algebra $\mathfrak{g}$. Then there exists an element $x \in C^{\infty}(\mathfrak{g})=\bigcap_{n \in \mathbb{N}} C^{n}(\mathfrak{g})$ with $e^{\text {ad } x} \mathfrak{h}=\mathfrak{h}^{\prime}$.

Proof. We use induction on the dimension of $\mathfrak{g}$. The case $\mathfrak{g}=\{0\}$ is trivial. Now assume that $\mathfrak{g} \neq\{0\}$. The last nonzero term of the derived series of $\mathfrak{g}$ is an abelian ideal of $\mathfrak{g}$, so that $\mathfrak{g}$ possesses a nonzero minimal abelian ideal $\mathfrak{n}$. Let $\varphi: \mathfrak{g} \rightarrow \mathfrak{g} / \mathfrak{n}$ denote the quotient homomorphism. Then $\varphi(\mathfrak{h})$ and $\varphi\left(\mathfrak{h}^{\prime}\right)$ are Cartan subalgebras of $\mathfrak{g} / \mathfrak{n}$ (Proposition 6.1.11), so that the induction hypothesis implies the existence of an element $y=\varphi(x) \in C^{\infty}(\mathfrak{g} / \mathfrak{n})=\varphi\left(C^{\infty}(\mathfrak{g})\right)$ with $e^{\operatorname{ad} y} \varphi(\mathfrak{h})=\varphi\left(\mathfrak{h}^{\prime}\right)$. Replacing $\mathfrak{h}$ by $e^{\operatorname{ad} x} \mathfrak{h}$, we may therefore assume that $\varphi(\mathfrak{h})=\varphi\left(\mathfrak{h}^{\prime}\right)$, i.e.,

$$
\mathfrak{h}+\mathfrak{n}=\mathfrak{h}^{\prime}+\mathfrak{n}
$$

Then $\mathfrak{h}$ and $\mathfrak{h}^{\prime}$ are Cartan subalgebras of the subalgebra $\mathfrak{h}+\mathfrak{n}$ (Lemma 6.1.11(iii)). If $\mathfrak{h}+\mathfrak{n} \neq \mathfrak{g}$, the assertion follows from the induction hypothesis. We may thus assume that $\mathfrak{g}=\mathfrak{h}+\mathfrak{n}=\mathfrak{h}^{\prime}+\mathfrak{n}$.

In view of the minimality of $\mathfrak{n}$, we either have $\mathfrak{n}=[\mathfrak{g}, \mathfrak{n}]$ or $[\mathfrak{g}, \mathfrak{n}]=\{0\}$. In the latter case, $\mathfrak{n}$ is central in $\mathfrak{g}$, so that $\mathfrak{n} \subseteq \mathfrak{h} \cap \mathfrak{h}^{\prime}$, so that $\mathfrak{h}=\mathfrak{h}+\mathfrak{n}= \mathfrak{h}^{\prime}+\mathfrak{n}=\mathfrak{h}^{\prime}$. This leaves the first case $\mathfrak{n}=[\mathfrak{g}, \mathfrak{n}]$. In this case, $\mathfrak{n} \subseteq C^{\infty}(\mathfrak{g})$ and $\mathfrak{n}$ is a simple $\mathfrak{g}$-module. As $\mathfrak{n}$ is abelian and $\mathfrak{g}=\mathfrak{n}+\mathfrak{h}$, the $\mathfrak{g}$-action on $\mathfrak{n}$ factors through an action of $\mathfrak{h}$, so that $\mathfrak{n}$ is a simple $\mathfrak{h}$-module. If $\mathfrak{n} \cap \mathfrak{h} \neq\{0\}$, the simplicity of the $\mathfrak{h}$-module $\mathfrak{n}$ yields $\mathfrak{n} \subseteq \mathfrak{h}$, and therefore $\mathfrak{h}=\mathfrak{g}$, which leads to $\mathfrak{h}=\mathfrak{h}^{\prime}$.

We may therefore assume that $\mathfrak{h} \cap \mathfrak{n}=\{0\}$, so that $\mathfrak{g} \cong \mathfrak{n} \rtimes \mathfrak{h}$ is a semidirect sum. Since $\mathfrak{h}$ and $\mathfrak{h}^{\prime}$ have the same dimension, we also have $\mathfrak{h}^{\prime} \cap \mathfrak{n}=\{0\}$, and there exists a linear map $f: \mathfrak{h} \rightarrow \mathfrak{n}$ for which

$$
\mathfrak{h}^{\prime}=\{h+f(h): h \in \mathfrak{h}\}
$$

is the corresponding graph. Since $\mathfrak{h}^{\prime}$ is a subalgebra, we have

$$
\left[h+f(h), h^{\prime}+f\left(h^{\prime}\right)\right]=\left[h, h^{\prime}\right]+\left[h, f\left(h^{\prime}\right)\right]-\left[h^{\prime}, f(h)\right] \in \mathfrak{h}^{\prime},
$$

showing that

$$
f\left(\left[h, h^{\prime}\right]\right)=\left[h, f\left(h^{\prime}\right)\right]-\left[h^{\prime}, f(h)\right]
$$

i.e., $f \in Z^{1}(\mathfrak{h}, \mathfrak{n})$. As $\mathfrak{n}$ is a simple nontrivial $\mathfrak{h}$-module, we have $\mathfrak{n}^{0}(\mathfrak{h})=\{0\}$, and Proposition 7.5.34 implies that $H^{1}(\mathfrak{h}, \mathfrak{n})=\{0\}$. Hence there exists a $v \in \mathfrak{n}$ with $f(h)=[h, v]$ for all $h \in \mathfrak{h}$, and then

$$
e^{\operatorname{ad} v} \mathfrak{h}=(\mathbf{1}+\operatorname{ad} v)(\mathfrak{h})=\{h+[x, h]: h \in \mathfrak{h}\}=\mathfrak{h}^{\prime} .
$$

\subsubsection*{Exercises for Section 7.5}
Exercise 7.5.1. (a) A central extension $q: \widehat{\mathfrak{g}} \rightarrow \mathfrak{g}$ with kernel $\mathfrak{z}$ is trivial if and only if $\mathfrak{z} \cap[\widehat{\mathfrak{g}}, \widehat{\mathfrak{g}}]=\{0\}$.\\
(b) For a vector space $V$, the central extension

$$
\mathbb{K} \mathbf{1} \hookrightarrow \mathfrak{g} \mathfrak{l}(V) \rightarrow \mathfrak{p} \mathfrak{g} \mathfrak{l}(V):=\mathfrak{g} \mathfrak{l}(V) / \mathbb{K} \mathbf{1}
$$

is trivial if and only if $\operatorname{dim} V<\infty$ and the characteristic $\operatorname{char}(\mathbb{K})$ of the field $\mathbb{K}$ is either 0 or does not divide $\operatorname{dim}(V)$. If this is not the case, then there exist endomorphisms $P, Q \in \operatorname{End}(V)$ with $[P, Q]=\mathbf{1}$.

Exercise 7.5.2. (a) Let $V$ be a module over the Lie algebra $\mathfrak{g}$ and $V^{*}$ be the dual space. Then $V^{*}$ becomes a $\mathfrak{g}$-module with

$$
(x \cdot f)(v):=-f(x \cdot v) \quad \text { for } \quad x \in \mathfrak{g}, f \in V^{*}, v \in V
$$

(b) If $V_{1}, \ldots, V_{n}$ and $W$ are $\mathfrak{g}$-modules, then the space $L$ of $n$-linear maps $V_{1} \times \cdots \times V_{n} \rightarrow W$ carries a $\mathfrak{g}$-module structure defined by

$$
(x \cdot f)\left(v_{1}, \ldots, v_{n}\right):=x \cdot f\left(v_{1}, \ldots, v_{n}\right)-\sum_{i=1}^{n} f\left(v_{1}, \ldots, v_{i-1}, x \cdot v_{i}, v_{i+1}, \ldots, v_{n}\right)
$$

Exercise 7.5.3. Let $V_{1}, \ldots, V_{n}$ be modules of the Lie algebra $\mathfrak{g}$. Then the tensor product $E:=V_{1} \otimes \cdots \otimes V_{n}$ becomes a $\mathfrak{g}$-module with\\
$x \cdot\left(v_{1} \otimes \cdots \otimes v_{n}\right):=x \cdot v_{1} \otimes \cdots \otimes v_{n}+v_{1} \otimes x \cdot v_{2} \otimes \cdots \otimes v_{n}+v_{1} \otimes \cdots \otimes x \cdot v_{n}$.

Exercise 7.5.4. Let $V$ and $W$ be two finite-dimensional $\mathfrak{g}$-modules and consider the linear isomorphism\\
$\Phi: V^{*} \otimes W \rightarrow \operatorname{Hom}(V, W), \quad \Phi(f \otimes w)(v):=f(v) w, \quad f \in V^{*}, w \in W, v \in V$.\\
Show that $\Phi$ is an isomorphism of $\mathfrak{g}$-modules if the $\mathfrak{g}$-module structure on $\operatorname{Hom}(V, W)$ is the one defined in Proposition 7.5.23, and the $\mathfrak{g}$-module structure on the tensor product is given by Exercises 7.5.2 and 7.5.3.

Exercise 7.5.5. $\mathbb{R}^{2}$ becomes an $\mathfrak{s l}_{2}(\mathbb{R})$-module by $x \cdot v:=x(v)$. Show:\\
(i) $\mathbb{R}^{2}$ is a simple $\mathfrak{s l}_{2}(\mathbb{R})$-module.\\
(ii) The $\mathfrak{s} \mathfrak{l}_{2}(\mathbb{R})$-module $\mathbb{R}^{2} \otimes \mathbb{R}^{2}$ splits into the direct sum of two simple $\mathfrak{s} \mathfrak{l}_{2}(\mathbb{R})$ modules.

Exercise 7.5.6. Let $\kappa: \mathfrak{g} \times \mathfrak{g} \rightarrow \mathbb{R}$ be an invariant bilinear form and put

$$
\Gamma(\kappa)(x, y, z):=\kappa([x, y], z)
$$

Show that:\\
(a) If $\kappa$ is skew-symmetric, then $\Gamma(\kappa)=0$.\\
(b) If $\kappa$ is symmetric, then $\Gamma(\kappa) \in Z^{3}(\mathfrak{g}, \mathbb{K})$ is a 3-cocycle.

Exercise 7.5.7. Show directly by computation that the Chevalley-Eilenberg differential d: $C(\mathfrak{g}, V) \rightarrow C(\mathfrak{g}, V)$ satisfies $\mathrm{d}^{2}=0$.

\subsection*{General Extensions of Lie Algebras}
In this section, we discuss a method to classify extensions of a Lie algebra $\mathfrak{g}$ by a Lie algebra $\mathfrak{n}$ in terms of data associated to these two Lie algebras. This generalizes the classification of abelian extension of $\mathfrak{g}$ by a $\mathfrak{g}$-module $V$ in terms of the second cohomology space $H^{2}(\mathfrak{g}, V)$.

Throughout this section, we shall denote an extension of $\mathfrak{g}$ by $\mathfrak{n}$ by the short exact sequence

$$
\mathbf{0} \rightarrow \mathfrak{n} \xrightarrow{\iota} \widehat{\mathfrak{g}} \xrightarrow{q} \mathfrak{g} \rightarrow \mathbf{0}
$$

where we identify $\mathfrak{n}$ with the ideal $\iota(\mathfrak{n})$ of $\widehat{\mathfrak{g}}$ to simplify notation.

\subsection*{1 $\mathfrak{g}$-Kernels}
One of the major difficulties of extensions by nonabelian Lie algebras is that $\mathfrak{n}$ acts nontrivially on itself by the adjoint action of $\widehat{\mathfrak{g}}$, so that the action of\\
$\widehat{\mathfrak{g}}$ on $\mathfrak{n}$ does not factor through an action of $\mathfrak{g}$. To overcome this problem, we introduce the concept of a $\mathfrak{g}$-kernel:

Definition 7.6.1. Recall the Lie algebra out( $\mathfrak{n}):=\operatorname{der}(\mathfrak{n}) / \operatorname{ad}(\mathfrak{n})$ of outer derivations of $\mathfrak{n}$ and write $[D]$ for the image of a derivation $D$ in out $(\mathfrak{n})$. A $\mathfrak{g}$-kernel for $\mathfrak{n}$ is a homomorphism

$$
s: \mathfrak{g} \rightarrow \operatorname{out}(\mathfrak{n})
$$

of Lie algebras. As ad $(\mathfrak{n})$ acts trivially on the center $\mathfrak{z}(\mathfrak{n})$, we have a natural out $(\mathfrak{n})$-module structure on $\mathfrak{z}(\mathfrak{n})$, so that each $\mathfrak{g}$-kernel $s: \mathfrak{g} \rightarrow$ out $(\mathfrak{n})$ defines in particular a $\mathfrak{g}$-module structure on $\mathfrak{z}(\mathfrak{n})$.

Remark 7.6.2. If $\mathfrak{n}$ is abelian, then a $\mathfrak{g}$-kernel for $\mathfrak{n}$ is simply a homomorphism $s: \mathfrak{g} \rightarrow \operatorname{der}(\mathfrak{n})$ since $\operatorname{ad}(\mathfrak{n})=\{0\}$.

Lemma 7.6.3. For each extension $q: \widehat{\mathfrak{g}} \rightarrow \mathfrak{g}$ of $\mathfrak{g}$ by the Lie algebra $\mathfrak{n}$, the adjoint action of $\widehat{\mathfrak{g}}$ on the ideal $\mathfrak{n}$ induces a homomorphism

$$
s: \mathfrak{g} \rightarrow \operatorname{out}(\mathfrak{n}), \quad q(x) \mapsto\left[\left.\operatorname{ad}_{\widehat{\mathfrak{g}}}(x)\right|_{\mathfrak{n}}\right]
$$

called the corresponding $\mathfrak{g}$-kernel. Equivalent extensions define the same $\mathfrak{g}$ kernel.

Proof. Since the adjoint representation

$$
\operatorname{ad}_{\mathfrak{n}}: \widehat{\mathfrak{g}} \rightarrow \operatorname{der}(\mathfrak{n}),\left.\quad x \mapsto \operatorname{ad}(x)\right|_{\mathfrak{n}}
$$

maps $\mathfrak{n}$ onto $\operatorname{ad}(\mathfrak{n})$, it factors through a homomorphism $s: \mathfrak{g} \cong \widehat{\mathfrak{g}} / \mathfrak{n} \rightarrow \operatorname{out}(\mathfrak{n})$, satisfying $s(q(x))=\left[\left.\operatorname{ad}_{\widehat{\mathfrak{g}}}(x)\right|_{\mathfrak{n}}\right]$.

If $\varphi: \widehat{\mathfrak{g}}_{1} \rightarrow \widehat{\mathfrak{g}}_{2}$ is an equivalence of extensions of $\mathfrak{g}$ by $\mathfrak{n}$, inducing the identity on $\mathfrak{n}$, then $q_{2} \circ \varphi=q_{1}$, and we have for each $x \in \widehat{\mathfrak{g}}_{1}$ the relation

$$
\left.\operatorname{ad}_{\widehat{\mathfrak{g}}_{1}}(x)\right|_{\mathfrak{n}}=\left.\varphi \circ \operatorname{ad}_{\widehat{\mathfrak{g}}_{1}}(x)\right|_{\mathfrak{n}}=\left.\operatorname{ad}_{\widehat{\mathfrak{g}}_{2}}(\varphi(x))\right|_{\mathfrak{n}} .
$$

Therefore, the corresponding $\mathfrak{g}$-kernels satisfy

$$
s_{1}\left(q_{1}(x)\right)=\left[\left.\operatorname{ad}_{\widehat{\mathfrak{g}}_{1}}(x)\right|_{\mathfrak{n}}\right]=\left[\left.\operatorname{ad}_{\widehat{\mathfrak{g}}_{2}}(\varphi(x))\right|_{\mathfrak{n}}\right]=s_{2}\left(q_{2}(\varphi(x))\right)=s_{2}\left(q_{1}(x)\right) .
$$

This proves that $s_{2}=s_{1}$.\\
For a given $\mathfrak{g}$-kernel $s$ for $\mathfrak{n}$, we write $\operatorname{Ext}(\mathfrak{g}, \mathfrak{n})_{s}$ for the set of equivalence classes of extensions of $\mathfrak{g}$ by $\mathfrak{n}$ for which $s$ is the corresponding $\mathfrak{g}$-kernel in the sense of Lemma 7.6.3. The classification of all extensions of $\mathfrak{g}$ by $\mathfrak{n}$ now decomposes into two different problems:\\
(P1) For a fixed $\mathfrak{g}$-kernel $s$ we have to parameterize the set $\operatorname{Ext}(\mathfrak{g}, \mathfrak{n})_{s}$.\\
(P2) We need a method to decide for which $\mathfrak{g}$-kernels the set $\operatorname{Ext}(\mathfrak{g}, \mathfrak{n})_{s}$ is nonempty.

Let $q: \widehat{\mathfrak{g}} \rightarrow \mathfrak{g}$ be an extension corresponding to the $\mathfrak{g}$-kernel $s$. To obtain a suitable description of $\widehat{\mathfrak{g}}$ in terms of $\mathfrak{g}$ and $\mathfrak{n}$, we choose a linear section $\sigma: \mathfrak{g} \rightarrow \widehat{\mathfrak{g}}$ of $q$. Then the linear map

$$
\Phi: \mathfrak{n} \times \mathfrak{g} \rightarrow \widehat{\mathfrak{g}}, \quad(n, x) \mapsto n+\sigma(x)
$$

is an isomorphism of vector spaces. To express the Lie bracket of $\widehat{\mathfrak{g}}$ in terms of the corresponding product coordinates, we define the linear map

$$
S: \mathfrak{g} \rightarrow \operatorname{der} \mathfrak{n}, \quad S(x):=\operatorname{ad}_{\mathfrak{n}}(\sigma(x)):=\left.(\operatorname{ad} \sigma(x))\right|_{\mathfrak{n}}
$$

and the alternating bilinear map

$$
\omega: \mathfrak{g} \times \mathfrak{g} \rightarrow \mathfrak{n}, \quad \omega(x, y):=[\sigma(x), \sigma(y)]-\sigma([x, y]) .
$$

Then $\Phi$ is an isomorphism of Lie algebras if we endow $\mathfrak{n} \times \mathfrak{g}$ with the Lie bracket

$$
\left[(n, x),\left(n^{\prime}, x^{\prime}\right)\right]:=\left(\left[n, n^{\prime}\right]+S(x) n^{\prime}-S\left(x^{\prime}\right) n+\omega\left(x, x^{\prime}\right),\left[x, x^{\prime}\right]\right)
$$

Remark 7.6.4. It is easy to verify directly that if $\mathfrak{n}$ and $\mathfrak{g}$ are Lie algebras, $\omega \in C^{2}(\mathfrak{g}, \mathfrak{n})$ and $S \in C^{1}(\mathfrak{g}, \operatorname{der}(\mathfrak{n}))$, then (7.11) defines a Lie bracket on $\mathfrak{n} \times \mathfrak{g}$ if and only if

$$
[S(x), S(y)]-S([x, y])=\operatorname{ad}(\omega(x, y)) \quad \text { for } \quad x, y \in \mathfrak{g}
$$

and

$$
\left(\mathrm{d}_{S} \omega\right)(x, y, z):=\sum_{\text {cycl. }} S(x) \omega(y, z)-\omega([x, y], z)=0
$$

Definition 7.6.5. If $\mathfrak{g}$ and $\mathfrak{n}$ are Lie algebras, then a factor system for ( $\mathfrak{g}, \mathfrak{n}$ ) is a pair ( $S, \omega$ ) with $S \in C^{1}\left(\mathfrak{g}\right.$, der $\mathfrak{n}$ ) and $\omega \in C^{2}(\mathfrak{g}, \mathfrak{n})$, satisfying (7.12) and (7.13).

We write $\widehat{\mathfrak{g}}:=\mathfrak{n} \times{ }_{(S, \omega)} \mathfrak{g}$ for the corresponding Lie algebra and note that $q: \widehat{\mathfrak{g}} \rightarrow \mathfrak{g}, q(n, x)=x$, is a surjective homomorphism with kernel $\mathfrak{n} \cong \mathfrak{n} \times\{0\}$. The corresponding $\mathfrak{g}$-kernel is given by $s: \mathfrak{g} \rightarrow$ out $(\mathfrak{n}), s(x)=[S(x)]$. In order to keep track of the $\mathfrak{g}$-action on $\mathfrak{z}(\mathfrak{n})$ induced by $s$, we write $Z^{k}(\mathfrak{g}, \mathfrak{z}(\mathfrak{n}))_{s}, B^{k}(\mathfrak{g}, \mathfrak{z}(\mathfrak{n}))_{s}$, and $H^{k}(\mathfrak{g}, \mathfrak{z}(\mathfrak{n}))_{s}$ instead of $Z^{k}(\mathfrak{g}, \mathfrak{z}(\mathfrak{n})), B^{k}(\mathfrak{g}, \mathfrak{z}(\mathfrak{n}))$, and $H^{k}(\mathfrak{g}, \mathfrak{z}(\mathfrak{n}))$.

Theorem 7.6.6. Let $s: \mathfrak{g} \rightarrow \operatorname{out}(\mathfrak{n})$ be a $\mathfrak{g}$-kernel for $\mathfrak{n}$. If $\operatorname{Ext}(\mathfrak{g}, \mathfrak{n})_{s} \neq \emptyset$ and $S \in C^{1}(\mathfrak{g}$, der $\mathfrak{n}$ ) satisfies $s(x)=[S(x)]$ for each $x \in \mathfrak{g}$, then each extension of $\mathfrak{g}$ by $\mathfrak{n}$ corresponding to $s$ is equivalent to some $\mathfrak{n} \times(S, \omega) \mathfrak{g}$, where ( $S, \omega$ ) is a factor system. For any such factor system, we obtain a bijection

$$
\Gamma: H^{2}(\mathfrak{g}, \mathfrak{z}(\mathfrak{n}))_{s} \rightarrow \operatorname{Ext}(\mathfrak{g}, \mathfrak{n})_{s}, \quad[\eta] \mapsto[\mathfrak{n} \times(S, \omega+\eta) \mathfrak{g}] .
$$

Proof. Let $q: \widehat{\mathfrak{g}} \rightarrow \mathfrak{g}$ be any extension of $\mathfrak{g}$ by $\mathfrak{n}$ corresponding to $s$. If $\sigma: \mathfrak{g} \rightarrow \widehat{\mathfrak{g}}$ is a linear section, then we have $[S(x)]=s(x)=\left[\operatorname{ad}_{\mathfrak{n}}(\sigma(x))\right]$, so that there exists a linear map $\gamma: \mathfrak{g} \rightarrow \mathfrak{n}$ with $S(x)=\operatorname{ad}_{\mathfrak{n}}(\sigma(x))+\operatorname{ad}(\gamma(x))$ for each $x \in \mathfrak{g}$. Then the new section $\widetilde{\sigma}:=\sigma+\gamma$ satisfies $S(x)=\operatorname{ad}_{\mathfrak{n}}(\widetilde{\sigma}(x))$, so that

$$
\omega(x, y):=[\widetilde{\sigma}(x), \widetilde{\sigma}(y)]-\widetilde{\sigma}([x, y])
$$

leads to a factor system ( $S, \omega$ ) with $\widehat{\mathfrak{g}} \cong \mathfrak{n} \times{ }_{(S, \omega)} \mathfrak{g}$.\\
For any other factor system ( $S, \widetilde{\omega}$ ), we have $\eta:=\widetilde{\omega}-\omega \in C^{2}(\mathfrak{g}, \mathfrak{z}(\mathfrak{n}))$ with

$$
\mathrm{d}_{\mathfrak{g}} \eta=\mathrm{d}_{S} \widetilde{\omega}-\mathrm{d}_{S} \omega=0
$$

Here we use that $\mathfrak{g}$-module structure on $\mathfrak{z}(\mathfrak{n})$ is given by $x \cdot z=S(x) z$, so that the corresponding Chevalley-Eilenberg differential $\mathrm{d}_{\mathfrak{g}}$ coincides with $\mathrm{d}_{S}$ on $\mathfrak{z}(\mathfrak{n})$-valued cochains. We conclude that the map

$$
Z^{2}(\mathfrak{g}, \mathfrak{z}(\mathfrak{n}))_{s} \rightarrow \operatorname{Ext}(\mathfrak{g}, \mathfrak{n})_{s}, \quad[\eta] \mapsto\left[\mathfrak{n} \times_{(S, \omega+\eta)} \mathfrak{g}\right]
$$

is surjective.\\
If $\varphi: \mathfrak{n} \times{ }_{(S, \omega)} \mathfrak{g} \rightarrow \mathfrak{n} \times{ }_{\left(S, \omega^{\prime}\right)} \mathfrak{g}$ is an equivalence of extensions, then $\varphi(n, x)= (n+\gamma(x), x)$ for some linear map $\gamma: \mathfrak{g} \rightarrow \mathfrak{n}$. In view of

$$
\begin{aligned}
(S(x) n+[\gamma(x), n], 0) & =[(\gamma(x), x),(n, 0)]=[\varphi(0, x),(n, 0)]=\varphi(S(x) n, 0) \\
& =(S(x) n, 0)
\end{aligned}
$$

we have $\gamma(\mathfrak{g}) \subseteq \mathfrak{z}(\mathfrak{n})$. Further,

$$
\begin{aligned}
(\omega(x, y)+\gamma([x, y]),[x, y]) & =\varphi([(0, x),(0, y)])=[(\gamma(x), x),(\gamma(y), y)] \\
& =\left(S(x) \gamma(y)-S(y) \gamma(x)+\omega^{\prime}(x, y),[x, y]\right)
\end{aligned}
$$

implies that $\omega=\mathrm{d}_{\mathfrak{g}} \gamma+\omega^{\prime}$. If, conversely, $\omega-\omega^{\prime}=\mathrm{d}_{\mathfrak{g}} \gamma$ for some $\gamma \in C^{1}(\mathfrak{g}, \mathfrak{z}(\mathfrak{n}))$, then $\varphi(n, x)=(n+\gamma(x), x)$ defines an equivalence of extensions

$$
\varphi: \mathfrak{n} \times{ }_{(S, \omega)} \mathfrak{g} \rightarrow \mathfrak{n} \times{ }_{\left(S, \omega^{\prime}\right)} \mathfrak{g},
$$

as a straight forward calculation shows. We conclude that the extensions defined by two factor systems ( $S, \omega$ ) and ( $S, \omega^{\prime}$ ) are equivalent if and only if $\omega-\omega^{\prime} \in B^{2}(\mathfrak{g}, \mathfrak{z}(\mathfrak{n}))$, and this implies the theorem.

Remark 7.6.7. The preceding theorem shows that the set $\operatorname{Ext}(\mathfrak{g}, \mathfrak{n})_{s}$, if nonempty, has the structure of an affine space whose translation group is $H^{2}(\mathfrak{g}, \mathfrak{z}(\mathfrak{n}))_{s} \cong \operatorname{Ext}(\mathfrak{g}, \mathfrak{z}(\mathfrak{n}))_{s}$ (cf. Proposition 7.5.18).\\
Remark 7.6.8. The group $H^{2}(\mathfrak{g}, \mathfrak{z}(\mathfrak{n}))_{s}$ very much depends on the $\mathfrak{g}$-kernel $s$. Let $\mathfrak{g}=\mathbb{K}^{2}$ and $\mathfrak{n}=\mathbb{K}$. Then $C^{2}(\mathfrak{g}, \mathfrak{z}(\mathfrak{n}))$ is 1-dimensional. Further, $\operatorname{dim} \mathfrak{g}=2$ implies $C^{3}(\mathfrak{g}, \mathfrak{z}(\mathfrak{n}))=\{0\}$, so that each 2-cochain is a cocycle. Since $B^{2}(\mathfrak{g}, \mathfrak{z}(\mathfrak{n}))_{s}$ vanishes if the module $\mathfrak{z}(\mathfrak{n})$ is trivial and coincides with $Z^{2}(\mathfrak{g}, \mathfrak{z}(\mathfrak{n}))_{s}$ otherwise, we have

$$
H^{2}(\mathfrak{g}, \mathfrak{z}(\mathfrak{n}))_{s} \cong \begin{cases}\mathbb{K} & \text { for } \mathfrak{g} \cdot \mathfrak{z}(\mathfrak{n})=\{0\} \\ \{0\} & \text { for } \mathfrak{g} \cdot \mathfrak{z}(\mathfrak{n}) \neq\{0\}\end{cases}
$$

Example 7.6.9. If $\mathfrak{g}$ is finite-dimensional semisimple, then the Second Whitehead Lemma 7.5.27 implies that $H^{2}(\mathfrak{g}, \mathfrak{z}(\mathfrak{n}))_{s}=\{0\}$ for each $\mathfrak{g}$-kernel $s: \mathfrak{g} \rightarrow \operatorname{out}(\mathfrak{n})$. Further, Corollary 5.6.10 implies that $s$ lifts to a homomorphism $S: \mathfrak{g} \rightarrow \operatorname{der}(\mathfrak{n})$. Therefore, any extension $\widehat{\mathfrak{g}}$ of $\mathfrak{g}$ by a finite-dimensional Lie algebra $\mathfrak{n}$ is a semidirect product $\mathfrak{n} \rtimes_{S} \mathfrak{g}$, and all the sets $\operatorname{Ext}(\mathfrak{g}, \mathfrak{n})_{s}$ consist of only one element.

Remark 7.6.10. If $\mathfrak{z}(\mathfrak{n})=\{0\}$, then $H^{2}(\mathfrak{g}, \mathfrak{z}(\mathfrak{n}))_{s}=\{0\}$ for each $\mathfrak{g}$-kernel $s$, so that all sets $\operatorname{Ext}(\mathfrak{g}, \mathfrak{n})_{s}$ contain at most one element (Theorem 7.6.6). On the other hand, we obtain for each $\mathfrak{g}$-kernel $s: \mathfrak{g} \rightarrow \operatorname{out}(\mathfrak{n})$ the pullback extension

$$
\widehat{\mathfrak{g}}:=s^{*} \operatorname{der}(\mathfrak{n}):=\{(x, D) \in \mathfrak{g} \times \operatorname{der}(\mathfrak{n}):[D]=s(x)\}
$$

of $\mathfrak{g}$ by $\operatorname{ad}(\mathfrak{n}) \cong \mathfrak{n}$. Since this extension obviously corresponds to the $\mathfrak{g}$-kernel $s$, it follows that $\operatorname{Ext}(\mathfrak{g}, \mathfrak{n})_{s} \neq \emptyset$ for any $\mathfrak{g}$-kernel and further that $\operatorname{Ext}(\mathfrak{g}, \mathfrak{n})_{s}= \left\{\left[s^{*} \operatorname{der}(\mathfrak{n})\right]\right\}$.

Remark 7.6.11 (Split extensions). An extension $q: \widehat{\mathfrak{g}} \rightarrow \mathfrak{g}$ of $\mathfrak{g}$ by $\mathfrak{n}$ splits if and only if there exists a homomorphism $\sigma: \mathfrak{g} \rightarrow \widehat{\mathfrak{g}}$ with $q \circ \sigma=\operatorname{id}_{\mathfrak{g}}$. Then $S(x):=\left.\operatorname{ad}(\sigma(x))\right|_{\mathfrak{n}}$ defines a homomorphism $S: \mathfrak{g} \rightarrow \operatorname{der}(\mathfrak{n})$ with

$$
\widehat{\mathfrak{g}} \cong \mathfrak{n} \rtimes_{S} \mathfrak{g} \cong \mathfrak{n} \times_{(S, 0)} \mathfrak{g}
$$

(a) For a $\mathfrak{g}$-kernel $s: \mathfrak{g} \rightarrow \operatorname{out}(\mathfrak{n})$, the existence of a split extension in the set $\operatorname{Ext}(\mathfrak{g}, \mathfrak{n})_{s}$ is therefore equivalent to the existence of a homomorphism $S: \mathfrak{g} \rightarrow \operatorname{der}(\mathfrak{n})$ with $[S(x)]=s(x)$ for each $x \in \mathfrak{g}$. Such a lift always exists if the extension $\operatorname{der}(\mathfrak{n})$ of out $(\mathfrak{n})$ by ad $\mathfrak{n}$ splits, but if this is not the case, then $\mathfrak{g}=\operatorname{out}(\mathfrak{n})$ and $s=\operatorname{id}_{\mathfrak{g}}$ is a $\mathfrak{g}$-kernel for which a homomorphic lift does not exist.\\
(b) If $\mathfrak{n}$ is abelian, then each set $\operatorname{Ext}(\mathfrak{g}, \mathfrak{n})_{s}$ contains exactly one split extension, namely $\mathfrak{n} \rtimes_{s} \mathfrak{g}$, but in general $\operatorname{Ext}(\mathfrak{g}, \mathfrak{n})_{s}$ may contain different classes of split extensions. To understand this phenomenon, consider two homomorphic lifts $S, S^{\prime}: \mathfrak{g} \rightarrow \operatorname{der}(\mathfrak{n})$ of the $\mathfrak{g}$-kernel $s$. Then $\gamma:=S^{\prime}-S: \mathfrak{g} \rightarrow \operatorname{ad} \mathfrak{n}$ is a linear map with $S^{\prime}=S+\gamma$, and the requirement that $S^{\prime}$ also is a homomorphism of Lie algebras is equivalent to $\gamma$ being a crossed homomorphism, i.e.,

$$
[S(x), \gamma(y)]-[S(y), \gamma(x)]-\gamma([x, y])+[\gamma(x), \gamma(y)]=0 \quad \text { for } \quad x, y \in \mathfrak{g} .
$$

A particularly simple case arises for $s=0$. Then $\operatorname{Ext}(\mathfrak{g}, \mathfrak{n})_{s}$ contains the class of the split extension $\mathfrak{n} \oplus \mathfrak{g}$ (Lie algebra direct sum). For the lift $S=0$, the other lifts simply correspond to homomorphisms $\gamma: \mathfrak{g} \rightarrow$ ad $\mathfrak{n}$. For any such homomorphism, there exists a linear lift $\widetilde{\gamma}: \mathfrak{g} \rightarrow \mathfrak{n}$ with ad $\circ \widetilde{\gamma}=\gamma$, but in general $\widetilde{\gamma}$ is not a homomorphism of Lie algebras. Then

$$
\omega(x, y):=[\widetilde{\gamma}(x), \widetilde{\gamma}(y)]-\widetilde{\gamma}([x, y])
$$

is a Lie algebra cocycle in $Z^{2}(\mathfrak{g}, \mathfrak{z}(\mathfrak{n}))$, corresponding to the central extension

$$
\gamma^{*} \mathfrak{n}=\{(x, n) \in \mathfrak{g} \times \mathfrak{n}: \gamma(x)=\operatorname{ad} n\}
$$

of $\mathfrak{g}$ by $\mathfrak{z}(\mathfrak{n})$. We thus obtain a well defined map

$$
\operatorname{Hom}(\mathfrak{g}, \operatorname{ad} \mathfrak{n}) \rightarrow H^{2}(\mathfrak{g}, \mathfrak{z}(\mathfrak{n})), \quad \gamma \mapsto\left[\gamma^{*} \mathfrak{n}\right]=[\omega]
$$

We are interested in a criterion for the split extension $\mathfrak{n} \rtimes_{\gamma} \mathfrak{g}$ to be equivalent to the trivial one.

For any lift $\widetilde{\gamma}$ as above, the map

$$
\varphi: \mathfrak{n} \times{ }_{(\gamma, \omega)} \mathfrak{g} \rightarrow \mathfrak{n} \oplus \mathfrak{g}, \quad(n, x) \mapsto(n+\widetilde{\gamma}(x), x)
$$

is a homomorphism of Lie algebras, hence an equivalence of extensions because we have for the $\mathfrak{n}$-components:

$$
\begin{aligned}
{\left[n+\widetilde{\gamma}(x), n^{\prime}+\widetilde{\gamma}\left(x^{\prime}\right)\right] } & =\left[n, n^{\prime}\right]+\gamma(x) n^{\prime}-\gamma\left(x^{\prime}\right) n+\left[\widetilde{\gamma}(x), \widetilde{\gamma}\left(x^{\prime}\right)\right] \\
& =\left(\left[n, n^{\prime}\right]+\gamma(x) n^{\prime}-\gamma\left(x^{\prime}\right) n+\omega\left(x, x^{\prime}\right)\right)+\widetilde{\gamma}\left(\left[x, x^{\prime}\right]\right)
\end{aligned}
$$

With Theorem 7.6.6, we therefore see that $\mathfrak{n} \rtimes_{\gamma} \mathfrak{g}$ is equivalent to the trivial extension if and only if the cohomology class $[\omega] \in H^{2}(\mathfrak{g}, \mathfrak{z}(\mathfrak{n}))$ vanishes, which in turn is equivalent to the existence of a homomorphic lift $\widetilde{\gamma}$ of $\gamma$.

If $\mathfrak{g}$ and ad $\mathfrak{n}$ are abelian, then we have a bracket map

$$
\eta: \operatorname{ad} \mathfrak{n} \times \operatorname{ad} \mathfrak{n} \rightarrow \mathfrak{z}(\mathfrak{n}), \quad \eta\left(\operatorname{ad} n, \operatorname{ad} n^{\prime}\right):=\left[n, n^{\prime}\right]
$$

and for a homomorphism $\gamma: \mathfrak{g} \rightarrow$ ad $\mathfrak{n}$ the extension $\gamma^{*} \mathfrak{n}$ is trivial if and only if $\gamma^{*} \eta$ vanishes, i.e., $\gamma(\mathfrak{g})$ is isotropic for $\eta$.

Example 7.6.12. We have seen above, that the theory of extensions of a Lie algebra $\mathfrak{g}$ by a Lie algebra $\mathfrak{n}$ simplifies significantly if the extension der( $\mathfrak{n}$ ) of out $(\mathfrak{n})$ by ad $(\mathfrak{n})$ splits. Here we describe a series of finite-dimensional nilpotent Lie algebras, for which this is not the case.

For each $n \geq 3$, we consider the ( $n+1$ )-dimensional filiform Lie algebra $L_{n}$ with a basis $e_{0}, \ldots, e_{n}$, where all nonzero brackets between basis elements are

$$
\left[e_{0}, e_{i}\right]=e_{i+1}, \quad i=1, \ldots, n-1
$$

Clearly, $L_{n}$ is generated as a Lie algebra by $e_{0}$ and $e_{1}$, and $L_{n} \cong V \rtimes_{A} \mathbb{K}$, where $V=\operatorname{span}\left\{e_{1}, \ldots, e_{n}\right\}$ and $A \in \mathfrak{g} \mathfrak{l}(V)$ is the nilpotent shift corresponding to the action of ad $e_{0}$ on $V$. Further, $\mathfrak{z}\left(L_{n}\right)=\mathbb{K} e_{n}$, so that $\operatorname{ad}\left(L_{n}\right) \cong L_{n-1}$ is an $n$-dimensional Lie algebra. One easily verifies that the following maps define derivations of $L_{n}$ :

$$
h_{k}:=\left(\operatorname{ad} e_{0}\right)^{k} \quad \text { and } \quad k=2, \ldots, n-1,
$$

$$
t_{1}\left(e_{i}\right):=\left\{\begin{array}{ll}
0 & \text { for } i=0, \\
e_{i} & \text { for } i>0,
\end{array} \quad t_{2}\left(e_{i}\right):= \begin{cases}e_{0} & \text { for } i=0 \\
(i-1) e_{i} & \text { for } i>0\end{cases}\right.
$$

and

$$
t_{3}\left(e_{i}\right):= \begin{cases}e_{1} & \text { for } i=0 \\ 0 & \text { for } i>0\end{cases}
$$

(cf. [GK96]). For the brackets of these derivations, we find

$$
\begin{aligned}
{\left[t_{1}, t_{2}\right] } & =0, \quad\left[t_{1}, t_{3}\right]=t_{3}, \quad\left[t_{2}, t_{3}\right]=-t_{3}, \quad\left[h_{i}, h_{j}\right]=0, \\
{\left[t_{1}, h_{k}\right] } & =0, \quad\left[t_{2}, h_{k}\right]=k h_{k}, \quad \text { and } \quad\left[t_{3}, h_{k}\right]=\operatorname{ad}\left(e_{k}\right) .
\end{aligned}
$$

Considering the action on the hyperplane ideal $V$, we see that $e_{0}, h_{2}, \ldots, h_{k-1}$, $t_{1}$, and $t_{2}$ lead to a linearly independent set of endomorphisms of $V$. Taking also the values on $e_{0}$ into account, we see that the $t_{i}, i=1,2,3, h_{k}$, $k=2, \ldots, n-1$, and $\operatorname{ad}\left(e_{i}\right), i=0, \ldots, n-1$, are linearly independent. Therefore,

$$
\operatorname{span}\left\{\left[h_{k}\right],\left[t_{i}\right]: i=1,2,3 ; k=2, \ldots, n-1\right\}
$$

is an ( $n+1$ )-dimensional Lie subalgebra of $\operatorname{out}\left(L_{n}\right)$ with the bracket relations

$$
\left[\left[t_{1}\right],\left[h_{k}\right]\right]=0, \quad\left[\left[t_{2}\right],\left[h_{k}\right]\right]=k\left[h_{k}\right], \quad \text { and } \quad\left[\left[t_{3}\right],\left[h_{k}\right]\right]=0
$$

We consider the 3-dimensional Lie subalgebra

$$
\mathfrak{g}:=\operatorname{span}\left\{\left[t_{1}\right],\left[t_{3}\right],\left[h_{2}\right]\right\} \subseteq \operatorname{out}\left(L_{n}\right)
$$

and claim that the Lie algebra extension

$$
\mathbf{0} \rightarrow \operatorname{ad}(\mathfrak{n}) \subseteq \widehat{\mathfrak{g}}:=\operatorname{ad}(\mathfrak{n})+\operatorname{span}\left\{t_{1}, t_{3}, h_{2}\right\} \rightarrow \mathfrak{g} \rightarrow \mathbf{0}
$$

does not split. This implies in particular that $\operatorname{der}(\mathfrak{n})$ does not split as an extension of out $(\mathfrak{n})$.

Suppose that $\widehat{\mathfrak{g}}$ does split, i.e., that there are 3 elements $z_{1}, z_{3}, x_{2} \in L_{n}$ such that the derivations

$$
\widetilde{t}_{1}:=t_{1}+\operatorname{ad}\left(z_{1}\right), \quad \widetilde{t}_{3}:=t_{3}+\operatorname{ad}\left(z_{3}\right), \quad \text { and } \quad \widetilde{h}_{2}:=h_{2}+\operatorname{ad}\left(x_{2}\right)
$$

satisfy

$$
\left[\widetilde{t}_{1}, \widetilde{t}_{3}\right]=\widetilde{t}_{3} \quad \text { and } \quad\left[\widetilde{t}_{1}, \widetilde{h}_{2}\right]=\left[\widetilde{t}_{3}, \widetilde{h}_{2}\right]=0
$$

The first relation is equivalent to

$$
\operatorname{ad}\left(t_{1} z_{3}-t_{3} z_{1}+\left[z_{1}, z_{3}\right]\right)=\operatorname{ad}\left(z_{3}\right)
$$

Writing $z_{3}=z_{3}^{0} e_{0}+z_{3}^{\prime}$ with $z_{3}^{\prime} \in V$, this leads to

$$
\operatorname{ad}\left(z_{3}^{\prime}-z_{1}^{0} e_{1}+\left[z_{1}, z_{3}\right]\right)=\operatorname{ad}\left(z_{3}^{0} e_{0}+z_{3}^{\prime}\right)
$$

and therefore to

$$
\operatorname{ad}\left(-z_{1}^{0} e_{1}+\left[z_{1}, z_{3}\right]\right)=\operatorname{ad}\left(z_{3}^{0} e_{0}\right)
$$

Applying this identity to $e_{1}$, we get $z_{3}^{0}=0$, so that $z_{3}=z_{3}^{\prime} \in V$.\\
Next we analyze the relation $\left[\widetilde{t}_{3}, \widetilde{h}_{2}\right]=0$, which is equivalent to

$$
0=\operatorname{ad}\left(e_{2}+t_{3} x_{2}-h_{2} z_{3}+\left[z_{3}, x_{2}\right]\right)
$$

Since $e_{2}-h_{2} z_{3}+\left[z_{3}, x_{2}\right] \in \operatorname{span}\left\{e_{2}, \ldots, e_{n}\right\}$ and $t_{3} x_{2} \in \mathbb{K} e_{1}$, it follows that $t_{3} x_{2}=0$ and hence $x_{2} \in V$. Then $\left[z_{3}, x_{2}\right]=0$, and we derive $\operatorname{ad}\left(e_{2}-h_{2} z_{3}\right)=$ 0 , which leads to the contradiction $e_{2} \in h_{2}(V)+\mathbb{K} e_{n} \subseteq \operatorname{span}\left\{e_{3}, \ldots, e_{n}\right\}$.

\subsubsection*{Integrability of $\mathfrak{g}$-Kernels}
We call a $\mathfrak{g}$-kernel $s: \mathfrak{g} \rightarrow$ out $(\mathfrak{n})$ integrable if there exists a factor system $(S, \omega)$ with $s(x)=[S(x)]$ for each $x \in \mathfrak{g}$ which is equivalent to $\operatorname{Ext}(\mathfrak{g}, \mathfrak{n})_{s} \neq \emptyset$ (cf. Theorem 7.6.6).

Remark 7.6.13. Let

$$
Q_{\mathfrak{n}}: \operatorname{der}(\mathfrak{n}) \rightarrow \operatorname{out}(\mathfrak{n}):=\operatorname{der}(\mathfrak{n}) / \operatorname{ad} \mathfrak{n}
$$

denote the quotient homomorphism and $s: \mathfrak{g} \rightarrow \operatorname{out}(\mathfrak{n})$. Then there exists a linear map $S: \mathfrak{g} \rightarrow$ der $\mathfrak{n}$ with $Q_{\mathfrak{n}} \circ S=s$. Since $s$ and $Q_{\mathfrak{n}}$ are homomorphisms,

$$
[S(x), S(y)]-S([x, y]) \in \operatorname{ad}(\mathfrak{n})
$$

for $x, y \in \mathfrak{g}$, and from that we derive the existence of an alternating bilinear map $\omega \in C^{2}(\mathfrak{g}, \mathfrak{n})$ satisfying (7.12). Indeed, for any linear section $\alpha: \operatorname{ad}(\mathfrak{n}) \rightarrow \mathfrak{n}$, we may put $\omega(x, y):=\alpha([S(x), S(y)]-S([x, y]))$. In general, the relation $\mathrm{d}_{S} \omega=0$ will not be satisfied.

Definition 7.6.14. For the following calculations, we introduce a product structure on $C(\mathfrak{g}, \mathfrak{n})$, defined for $\alpha \in C^{p}(\mathfrak{g}, \mathfrak{n})$ and $\beta \in C^{q}(\mathfrak{g}, \mathfrak{n})$ by

$$
\begin{aligned}
& {[\alpha, \beta]\left(x_{1}, \ldots, x_{p+q}\right)} \\
& \quad:=\sum_{\sigma \in \operatorname{Sh}(p, q)} \operatorname{sgn}(\sigma)\left[\alpha\left(x_{\sigma(1)}, \ldots, x_{\sigma(p)}\right), \beta\left(x_{\sigma(p+1)}, \ldots, x_{\sigma(p+q)}\right)\right]
\end{aligned}
$$

where $\operatorname{Sh}(p, q)$ denotes the set of all ( $p, q$ )-shuffles in $S_{p+q}$, i.e., all permutations with

$$
\sigma(1)<\cdots<\sigma(p) \quad \text { and } \quad \sigma(p+1)<\cdots<\sigma(p+q)
$$

For $p=1$ and $q=2$, we then have

$$
[\alpha, \alpha](x, y)=[\alpha(x), \alpha(y)]-[\alpha(y), \alpha(x)]=2[\alpha(x), \alpha(y)]
$$

and

$$
[\alpha, \beta](x, y, z)=\sum_{\text {cycl. }}[\alpha(x), \beta(y, z)]
$$

as well as

$$
[\beta, \alpha](x, y, z)=\sum_{\text {cycl. }}[\beta(x, y), \alpha(z)]=-[\alpha, \beta](x, y, z)
$$

In particular, the Jacobi identity implies that

$$
[\alpha,[\alpha, \alpha]]=0
$$

We also put for $\alpha \in C^{1}(\mathfrak{g}, \mathfrak{n})$ :

$$
\left(\mathrm{d}_{S} \alpha\right)(x, y):=S(x) \alpha(y)-S(y) \alpha(x)-\alpha([x, y])
$$

Lemma 7.6.15. Every $\gamma \in C^{1}(\mathfrak{g}, \mathfrak{n})$ satisfies $\frac{1}{2} \mathrm{~d}_{S}[\gamma, \gamma]=\left[\mathrm{d}_{S} \gamma, \gamma\right]= -\left[\gamma, \mathrm{d}_{S} \gamma\right]$.\\
Proof. We calculate

$$
\begin{aligned}
& \frac{1}{2}\left(\mathrm{~d}_{S}[\gamma, \gamma]\right)(x, y, z) \\
& =\sum_{\text {cycl. }} S(x)[\gamma(y), \gamma(z)]-[\gamma([x, y]), \gamma(z)] \\
& =\sum_{\text {cycl. }}[S(x) \gamma(y), \gamma(z)]+[\gamma(y), S(x) \gamma(z)]-[\gamma([x, y]), \gamma(z)] \\
& =\sum_{\text {cycl. }}[S(x) \gamma(y), \gamma(z)]+[\gamma(z), S(y) \gamma(x)]-[\gamma([x, y]), \gamma(z)] \\
& =\sum_{\text {cycl. }}\left[\left(\mathrm{d}_{S} \gamma\right)(x, y), \gamma(z)\right]=\left[\mathrm{d}_{S} \gamma, \gamma\right](x, y, z)
\end{aligned}
$$

Lemma 7.6.16. If $S: \mathfrak{g} \rightarrow \operatorname{der}(\mathfrak{n})$ is a linear map and $\omega \in C^{2}(\mathfrak{g}, \mathfrak{n})$ satisfies (7.12), then $\mathrm{d}_{S} \omega \in Z^{3}(\mathfrak{g}, \mathfrak{z}(\mathfrak{n}))_{s}$, where $s=Q_{\mathfrak{n}} \circ S$.

Proof. First, we show that $\operatorname{im}\left(\mathrm{d}_{S} \omega\right) \subseteq \mathfrak{z}(\mathfrak{n})$ :

$$
\begin{aligned}
\operatorname{ad} & \left(\mathrm{d}_{S} \omega(x, y, z)\right) \\
= & \sum_{\text {cycl. }} \operatorname{ad}(S(x) \omega(y, z))-\operatorname{ad}(\omega([x, y], z)) \\
= & \sum_{\text {cycl. }}[S(x), \operatorname{ad}(\omega(y, z))]-[S([x, y]), S(z)]-S([[x, y], z]) \\
= & \sum_{\text {cycl. }}[S(x),[S(y), S(z)]]-[S(x), S([y, z])]-[S([x, y]), S(z)] \\
& -S([[x, y], z]) \\
= & \sum_{\text {cycl. }}[S(x),[S(y), S(z)]]-S([[x, y], z])=0
\end{aligned}
$$

where the last equality follows from the Jacobi identities in $\operatorname{der}(\mathfrak{n})$ and $\mathfrak{g}$. We conclude that $\operatorname{im}\left(\mathrm{d}_{S} \omega\right) \subseteq \mathfrak{z}(\mathfrak{n})$.

It remains to see that $\mathrm{d}_{S} \omega$ is a 3-cocycle. Let

$$
\widetilde{\mathfrak{g}}:=s^{*} \operatorname{der}(\mathfrak{n})=\{(x, D) \in \mathfrak{g} \times \operatorname{der}(\mathfrak{n}):[D]=s(x)\}
$$

and note that $q(x, D):=x$ defines a Lie algebra extension of $\mathfrak{g}$ by ad $\mathfrak{n}$. Further, $\rho(x, D):=D$ defines an action of $\widetilde{\mathfrak{g}}$ on $\mathfrak{n}$ by derivations, satisfying

$$
[\rho(x, D)]=s(x)=s(q(x, D))
$$

Since $\mathfrak{z}(\mathfrak{n})$ is $\widetilde{\mathfrak{g}}$-invariant and the action of $\widetilde{\mathfrak{g}}$ on $\mathfrak{z}(\mathfrak{n})$ factors through the given action of $\mathfrak{g}$ on $\mathfrak{z}(\mathfrak{n})$, it suffices to see that the pullback $q^{*}\left(\mathrm{~d}_{S} \omega\right) \in C^{3}(\widetilde{\mathfrak{g}}, \mathfrak{z}(\mathfrak{n}))$ is a cocycle. For $\widetilde{S}(x, D):=S(x)$ and $\widetilde{\omega}:=q^{*} \omega$, we have

$$
\begin{aligned}
\left(q^{*}\left(\mathrm{~d}_{S} \omega\right)\right)(x, y, z) & =\sum_{\text {cycl. }} S(q(x)) \omega(q(y), q(z))-\omega([q(x), q(y)], q(z)) \\
& =\sum_{\text {cycl. }} \widetilde{S}(x) \widetilde{\omega}(y, z)-\widetilde{\omega}([x, y], z)=\left(\mathrm{d}_{\widetilde{S}} \widetilde{\omega}\right)(x, y, z)
\end{aligned}
$$

Since $\widetilde{S}(x, D)-D=S(x)-D \in \operatorname{ad}(\mathfrak{n})$ for each $x \in \mathfrak{g}$, there exists a linear map $\beta: \widetilde{\mathfrak{g}} \rightarrow \mathfrak{n}$ with

$$
\widetilde{S}=\rho+\operatorname{ad} \circ \beta: \widetilde{\mathfrak{g}} \rightarrow \operatorname{der}(\mathfrak{n})
$$

With respect to the $\widetilde{\mathfrak{g}}$-module structure on $\mathfrak{n}$, we thus have

$$
\mathrm{d}_{\widetilde{S}} \widetilde{\omega}=\mathrm{d}_{\widetilde{\mathfrak{g}}} \widetilde{\omega}+[\beta, \widetilde{\omega}]
$$

Since

$$
\begin{aligned}
\operatorname{ad}(\widetilde{\omega}(x, y)) & =[\widetilde{S}(x), \widetilde{S}(y)]-\widetilde{S}([x, y]) \\
& =[\rho(x)+\operatorname{ad} \beta(x), \rho(y)+\operatorname{ad} \beta(y)]-\rho([x, y])-\operatorname{ad}(\beta([x, y])) \\
& =\operatorname{ad}(\rho(x) \beta(y)-\rho(y) \beta(x)+[\beta(x), \beta(y)]-\beta([x, y])) \\
& =\operatorname{ad} \circ\left(\mathrm{d}_{\widetilde{\mathfrak{g}}} \beta+\frac{1}{2}[\beta, \beta]\right)(x, y),
\end{aligned}
$$

we further get with Lemma 7.6.15 and $[\beta,[\beta, \beta]]=0$ :

$$
[\beta, \widetilde{\omega}]=\left[\beta, \mathrm{d}_{\widetilde{\mathfrak{g}}} \beta+\frac{1}{2}[\beta, \beta]\right]=\left[\beta, \mathrm{d}_{\widetilde{\mathfrak{g}}} \beta\right]=-2 \mathrm{~d}_{\mathfrak{g}}[\beta, \beta]
$$

This proves that $\mathrm{d}_{\widetilde{S}} \widetilde{\omega}=\mathrm{d}_{\mathfrak{g}}(\widetilde{\omega}-2[\beta, \beta])$ is an element of $B^{3}(\widetilde{\mathfrak{g}}, \mathfrak{n})_{s}$, hence in particular a cocycle, and therefore $\mathrm{d}_{S} \omega$ also is a cocycle.

Lemma 7.6.17. The cohomology class $\chi(s):=\left[\mathrm{d}_{S} \omega\right] \in H^{3}(\mathfrak{g}, \mathfrak{z}(\mathfrak{n}))_{s}$ only depends on $s$, as long as $s(x)=[S(x)]$ holds for each $x \in \mathfrak{g}$.

Proof. If $S$ is fixed and $\omega^{\prime} \in C^{2}(\mathfrak{g}, \mathfrak{n})$ also satisfies (7.12), then $\omega^{\prime}-\omega \in C^{2}(\mathfrak{g}, \mathfrak{z}(\mathfrak{n}))$, so that

$$
\mathrm{d}_{S} \omega^{\prime}=\mathrm{d}_{S} \omega+\mathrm{d}_{S}\left(\omega^{\prime}-\omega\right)=\mathrm{d}_{S} \omega+\mathrm{d}_{\mathfrak{g}}\left(\omega^{\prime}-\omega\right) \in \mathrm{d}_{S} \omega+B^{3}(\mathfrak{g}, \mathfrak{z}(\mathfrak{n}))_{s}
$$

and thus $\left[\mathrm{d}_{S} \omega\right]=\left[\mathrm{d}_{S} \omega^{\prime}\right]$.\\
If $S^{\prime}: \mathfrak{g} \rightarrow \operatorname{der}(\mathfrak{n})$ is another linear map with $\left[S^{\prime}(x)\right]=s(x)$ for each $x \in \mathfrak{g}$, then $S^{\prime}=S+$ ad $\circ \gamma$ for some linear map $\gamma: \mathfrak{g} \rightarrow \mathfrak{n}$. A direct calculation shows that

$$
\omega^{\prime}:=\omega+\mathrm{d}_{S} \gamma+\frac{1}{2}[\gamma, \gamma]
$$

satisfies

$$
\operatorname{ad}\left(\omega^{\prime}(x, y)\right)=\left[S^{\prime}(x), S^{\prime}(y)\right]-S^{\prime}([x, y])
$$

We claim that $\mathrm{d}_{S^{\prime}} \omega^{\prime}=\mathrm{d}_{S} \omega$, and this will complete the proof. In fact, we obtain with Lemma 7.6.15

$$
\begin{aligned}
\mathrm{d}_{S^{\prime}} \omega^{\prime} & =\mathrm{d}_{S} \omega^{\prime}+\left[\gamma, \omega^{\prime}\right]=\mathrm{d}_{S} \omega+\mathrm{d}_{S}^{2} \gamma+\frac{1}{2} \mathrm{~d}_{S}[\gamma, \gamma]-\left[\omega^{\prime}, \gamma\right] \\
& =\mathrm{d}_{S} \omega+\mathrm{d}_{S}^{2} \gamma+\left[\mathrm{d}_{S} \gamma, \gamma\right]-\left[\omega^{\prime}, \gamma\right]
\end{aligned}
$$

Further,

$$
\begin{aligned}
\left(\mathrm{d}_{S}^{2} \gamma\right)(x, y, z)=\sum_{\text {cycl. }} S(x)\left(\mathrm{d}_{S} \gamma(y, z)\right)-\mathrm{d}_{S} \gamma([x, y], z) \\
=\sum_{\text {cycl. }} S(x)(S(y) \gamma(z)-S(z) \gamma(y)-\gamma([y, z])) \\
-S([x, y]) \gamma(z)+S(z) \gamma([x, y])+\gamma([[x, y], z]) \\
=\sum_{\text {cycl. }} S(x) S(y) \gamma(z)-S(y) S(x) \gamma(z)-S(x) \gamma([y, z]) \\
-S([x, y]) \gamma(z)+S(x) \gamma([y, z]) \\
=\sum_{\text {cycl. }}[\omega(x, y), \gamma(z)]-S(x) \gamma([y, z])+S(x) \gamma([y, z]) \\
=[\omega, \gamma](x, y, z) .
\end{aligned}
$$

We thus obtain

$$
\mathrm{d}_{S^{\prime}} \omega^{\prime}=\mathrm{d}_{S} \omega+\left[\omega+\mathrm{d}_{S} \gamma-\omega^{\prime}, \gamma\right]=\mathrm{d}_{S} \omega-\frac{1}{2}[[\gamma, \gamma], \gamma]=\mathrm{d}_{S} \omega
$$

Theorem 7.6.18. For a $\mathfrak{g}$-kernel $s: \mathfrak{g} \rightarrow$ out $(\mathfrak{n})$, the set $\operatorname{Ext}(\mathfrak{g}, \mathfrak{n})_{s}$ is nonempty if and only if the cohomology class $\chi(s)$ vanishes.

Proof. The set $\operatorname{Ext}(\mathfrak{g}, \mathfrak{n})_{s}$ is nonempty if and only if there exists a factor system ( $S, \omega$ ) corresponding to $s$ with $\mathrm{d}_{S} \omega=0$. If this is the case, then\\
$\chi(s)=\left[\mathrm{d}_{S} \omega\right]=0$. If, conversely, $\chi(s)=0$ and $\beta \in C^{2}(\mathfrak{g}, \mathfrak{z}(\mathfrak{n}))$ satisfies $\mathrm{d}_{S} \omega=\mathrm{d}_{\mathfrak{g}} \beta$, then $\widetilde{\omega}:=\omega-\beta$ also satisfies

$$
\operatorname{ad}(\widetilde{\omega}(x, y))=\operatorname{ad}(\omega(x, y))=[S(x), S(y)]-S([x, y])
$$

for $x, y \in \mathfrak{g}$, and further $\mathrm{d}_{S} \widetilde{\omega}=\mathrm{d}_{S} \omega-\mathrm{d}_{S} \beta=\mathrm{d}_{S} \omega-\mathrm{d}_{\mathfrak{g}} \beta=0$. Hence $(S, \widetilde{\omega})$ is a factor system corresponding to $s$.

\subsubsection*{Exercises for Section 7.6}
Definition 7.6.19. A Lie superalgebra (over a field $\mathbb{K}$ with $2,3 \in \mathbb{K}^{\times}$) is a $\mathbb{Z} / 2 \mathbb{Z}$-graded vector space $\mathfrak{g}=\mathfrak{g}_{\overline{0}} \oplus \mathfrak{g}_{\overline{1}}$ with a bilinear map $[., \cdot]$ satisfying\\
(LS1) $[\alpha, \beta]=(-1)^{p q+1}[\beta, \alpha]$ for $x \in \mathfrak{g}_{p}$ and $y \in \mathfrak{g}_{q}$.\\
(LS2) $(-1)^{p r}[[\alpha, \beta], \gamma]+(-1)^{q p}[[\beta, \gamma], \alpha]+(-1)^{q r}[[\gamma, \alpha], \beta]=0$ for $\alpha \in \mathfrak{g}_{p}$, $\beta \in \mathfrak{g}_{q}$ and $\gamma \in \mathfrak{g}_{r}$.

Note that (LS1) implies that $[\alpha, \alpha]=0=[\beta,[\beta, \beta]]$ for $\alpha \in \mathfrak{g}_{\overline{0}}$ and $\beta \in \mathfrak{g}_{\overline{1}}$.\\
Exercise 7.6.1. Let $\mathfrak{g}$ and $\mathfrak{n}$ be Lie algebras. Then the bracket defined in Definition 7.6.14 defines on the $\mathbb{Z} / 2 \mathbb{Z}$-graded vector space

$$
C(\mathfrak{g}, \mathfrak{n}):=\bigoplus_{p \in \mathbb{N}_{0}} C^{p}(\mathfrak{g}, \mathfrak{n}):=C^{\text {even }}(\mathfrak{g}, \mathfrak{n}) \oplus C^{\text {odd }}(\mathfrak{g}, \mathfrak{n})
$$

the structure of a Lie superalgebra.

\subsubsection*{Notes on Chapter 7}
The finite dimension of the Lie algebra $\mathfrak{g}$ was not essential for the proof of the Poincaré-Birkhoff-Witt Theorem 7.1.9. Only slight changes yield the same theorem for any Lie algebra (cf. [Bou89, Chapter 1]).

The origin of Lie algebra cohomology lies in the work of É. Cartan who reduced the determination of the rational (singular) cohomology of compact Lie groups to a purely algebraic problem which later led to the invention of Lie algebra cohomology by C. Chevalley and S. Eilenberg [CE48]. Computational methods to evaluate Lie algebra cohomology spaces have been developed in [Kos50], where the cohomology of semisimple Lie algebras is studied in detail.

The proof of Ado's Theorem we present in Section 7.4 follows an idea of Y.A. Neretin [Ner02].

In [Ho54a], G. Hochschild shows that, for each $\mathfrak{g}$-module $V$ of a Lie algebra $\mathfrak{g}$, each element of $H^{3}(\mathfrak{g}, V)$ arises as an obstruction class for a homomorphism $s: \mathfrak{g} \rightarrow \operatorname{out}(\mathfrak{n})$, where $\mathfrak{n}$ is a Lie algebra with $V=\mathfrak{z}(\mathfrak{n})$. In [Ho54b], he analyzes for a finite-dimensional Lie algebra $\mathfrak{g}$ and a finite-dimensional $\mathfrak{g}$-module $V$, the question of the existence of a finite-dimensional Lie algebra $\mathfrak{n}$ with the above properties. In this case, the answer is affirmative if $\mathfrak{g}$ is solvable, but if $\mathfrak{g}$ is semisimple, then all obstructions for a homomorphism $s: \mathfrak{g} \rightarrow \operatorname{out}(\mathfrak{n})$\\
are trivial because $s$ lifts to a homomorphism $S: \mathfrak{g} \rightarrow$ der $\mathfrak{n}$ by Levi's Theorem 5.6.6. The general result is that a cohomology class $[\omega] \in H^{3}(\mathfrak{g}, V)$ arises as an obstruction $\chi(s)$ of a $\mathfrak{g}$-kernel $s$ if and only if its restriction to a Levi complement $\mathfrak{s}$ in $\mathfrak{g}$ vanishes.

\section*{Part III}
\section*{Manifolds and Lie Groups}
\section*{Smooth Manifolds}
Even though it is possible to prove that each Lie group, up to covering, is isomorphic to a linear Lie group of the type discussed in Part I, the natural setting for Lie groups is the category of smooth manifolds in which Lie groups can be viewed as the group objects. Thus we will use linear Lie groups rather as a source of examples and will start building the theory of Lie groups from scratch in Chapter 9 by defining them as groups which are smooth manifolds for which the group operations are smooth.

We start in the present chapter by reviewing same basic features of differential analysis on open domains in $\mathbb{R}^{n}$. Then we introduce smooth manifolds and their tangent bundles, which allows us to define the Lie algebra vector fields, thus preparing the grounds for the definition of the Lie algebra of a general Lie group. Relating vector fields to ordinary differential equations leads to integral curves and local flows. These are used later on to define the exponential function of a Lie group. We conclude the chapter with a discussion of various concepts of submanifolds which naturally play are role in the study of subgroups of Lie groups.

In basic calculus courses, one mostly deals with (differentiable) functions on open subsets of $\mathbb{R}^{n}$, but as soon as one wants to solve equations of the form $f(x)=y$, where $f: U \rightarrow \mathbb{R}^{m}$ is a differentiable function and $U$ is open in $\mathbb{R}^{n}$, one observes that the set $f^{-1}(y)$ of solutions behaves in a much more complicated manner than one is used to from linear algebra, where $f$ is linear and $f^{-1}(y)$ is the intersection of $U$ with an affine subspace. One way to approach differentiable manifolds is to think of them as the natural objects arising as solutions of nonlinear equations as above (under some nondegeneracy condition on $f$, made precise by the Implicit Function Theorem). For submanifolds of $\mathbb{R}^{n}$, this is a quite natural approach which immediately leads to the method of Lagrange multipliers to deal with extrema of differentiable functions under differentiable constraints. This is the external perspective on differentiable manifolds, which has the serious disadvantage that it depends very much on the surrounding space $\mathbb{R}^{n}$.

It is much more natural to adopt a more intrinsic perspective: an $n$ dimensional manifold is a topological space which locally looks like $\mathbb{R}^{n}$. More precisely, it arises by gluing open subsets of $\mathbb{R}^{n}$ in a smooth (differentiable) way. Below we shall make this more precise.

The theory of smooth manifolds has three levels:\\
(1) The infinitesimal level, where one deals with tangent spaces, tangent vectors and differentials of maps,\\
(2) the local level, which is analysis on open subsets of $\mathbb{R}^{n}$, and\\
(3) the global level, where one studies the global behavior of manifolds and other related structures.

These three levels are already visible in one-variable calculus: Suppose we are interested in the global maximum of a differentiable function $f: \mathbb{R} \rightarrow \mathbb{R}$ which is a question about the global behavior of this function. The necessary condition $f^{\prime}\left(x_{0}\right)=0$ belongs to the infinitesimal level because it says something about the behavior of $f$ infinitesimally close to the point $x_{0}$. The sufficient criterion for a local maximum: $f^{\prime}\left(x_{0}\right)=0, f^{\prime \prime}\left(x_{0}\right)<0$ provides information on the local level. Of course, this is far from being the whole story and one really has to study global properties of $f$, such as $\lim _{x \rightarrow \pm \infty} f(x)=0$, to guarantee the existence of global maxima.

\subsection*{Smooth Maps in Several Variables}
First, we recall some facts and definitions from calculus of several variables, formulated in a way that will be convenient for us in the following.

Definition 8.1.1 (Differentiable maps). (a) Let $n, m \in \mathbb{N}$ and $U \subseteq \mathbb{R}^{n}$ be an open subset. A function $f: U \rightarrow \mathbb{R}^{m}$ is called differentiable at $x \in U$ if there exists a linear map $L \in \operatorname{Hom}\left(\mathbb{R}^{n}, \mathbb{R}^{m}\right)$ such that for one norm (and hence for all norms) on $\mathbb{R}^{n}$ we have

$$
\lim _{h \rightarrow 0} \frac{f(x+h)-f(x)-L(h)}{\|h\|}=0
$$

If $f$ is differentiable at $x$, then for each $h \in \mathbb{R}^{n}$ we have

$$
\lim _{t \rightarrow 0} \frac{1}{t}(f(x+t h)-f(x))=\lim _{t \rightarrow 0} \frac{1}{t} L(t h)=L(h),
$$

so that $L(h)$ is the directional derivative of $f$ at $x$ in the direction $h$. It follows in particular that condition (8.1) determines the linear map $L$ uniquely. We therefore write

$$
\mathrm{d} f(x)(h):=\lim _{t \rightarrow 0} \frac{1}{t}(f(x+t h)-f(x))=L(h)
$$

and call the linear map $\mathrm{d} f(x)$ the differential of $f$ at $x$.\\
(b) Let $e_{1}, \ldots, e_{n}$ denote the canonical basis vectors in $\mathbb{R}^{n}$. Then

$$
\frac{\partial f}{\partial x_{i}}(x):=\mathrm{d} f(x)\left(e_{i}\right)
$$

is called the $i$ th partial derivative of $f$ at $x$. If $f$ is differentiable at each $x \in U$, then the partial derivatives are functions

$$
\frac{\partial f}{\partial x_{i}}: U \rightarrow \mathbb{R}^{m}
$$

and we say that $f$ is continuously differentiable, or a $C^{1}$-map, if all its partial derivatives are continuous. For $k \geq 2$, the map $f$ is said to be a $C^{k}$-map if it is $C^{1}$ and all its partial derivatives are $C^{k-1}$-maps. We say that $f$ is smooth, or a $C^{\infty}$-map, if it is $C^{k}$ for each $k \in \mathbb{N}$. We denote the space of $C^{k}$-maps $U \rightarrow \mathbb{R}^{m}$ by $C^{k}\left(U, \mathbb{R}^{m}\right)$.\\
(c) If $I \subseteq \mathbb{R}$ is an interval and $\gamma: I \rightarrow \mathbb{R}^{n}$ is a differentiable curve, we also write

$$
\dot{\gamma}(t)=\gamma^{\prime}(t)=\lim _{h \rightarrow 0} \frac{\gamma(t+h)-\gamma(t)}{h}
$$

This is related to the notation from above by

$$
\gamma^{\prime}(t)=\mathrm{d} \gamma(t)\left(e_{1}\right)
$$

where $e_{1}=1 \in \mathbb{R}$ is the canonical basis vector.\\
Definition 8.1.2. Let $U \subseteq \mathbb{R}^{n}$ and $V \subseteq \mathbb{R}^{m}$ be open subsets. A map $f: U \rightarrow V$ is called $C^{k}$ if it is $C^{k}$ as a map $U \rightarrow \mathbb{R}^{m}$.

For $n \geq 1$ a $C^{k}$-map $f: U \rightarrow V$ is called a $C^{k}$-diffeomorphism if there exists a $C^{k}$-map $g: V \rightarrow U$ with

$$
f \circ g=\mathrm{id}_{V} \quad \text { and } \quad g \circ f=\mathrm{id}_{U}
$$

Obviously, this is equivalent to $f$ being bijective and $f^{-1}$ being a $C^{k}$-map. Whenever such a diffeomorphism exists, we say that the domains $U$ and $V$ are $C^{k}$-diffeomorphic. For $k=0$, we thus obtain the notion of a homeomorphism.

Theorem 8.1.3 (Chain Rule). Let $U \subseteq \mathbb{R}^{n}$ and $V \subseteq \mathbb{R}^{m}$ be open subsets. Further, let $f: U \rightarrow V$ be a $C^{k}$-map and $g: V \rightarrow \mathbb{R}^{d}$ a $C^{k}$-map. Then $g \circ f$ is a $C^{k}$-map, and for each $x \in U$ we have in $\operatorname{Hom}\left(\mathbb{R}^{n}, \mathbb{R}^{d}\right)$ :

$$
\mathrm{d}(g \circ f)(x)=\mathrm{d} g(f(x)) \circ \mathrm{d} f(x)
$$

The Chain Rule is an important tool which permits us to "linearize" nonlinear information. The following proposition is an example.

Proposition 8.1.4 (Invariance of the Dimension). If the nonempty open subsets $U \subseteq \mathbb{R}^{n}$ and $V \subseteq \mathbb{R}^{m}$ are $C^{1}$-diffeomorphic, then $m=n$.

Proof. Let $f: U \rightarrow V$ be a $C^{1}$-diffeomorphism and $g: V \rightarrow U$ its inverse. Pick $x \in U$. Then the Chain Rule implies that

$$
\mathrm{id}_{\mathbb{R}^{n}}=\mathrm{d}(g \circ f)(x)=\mathrm{d} g(f(x)) \circ \mathrm{d} f(x)
$$

and

$$
\mathrm{id}_{\mathbb{R}^{m}}=\mathrm{d}(f \circ g)(f(x))=\mathrm{d} f(x) \circ \mathrm{d} g(f(x))
$$

so that $\mathrm{d} f(x): \mathbb{R}^{n} \rightarrow \mathbb{R}^{m}$ is a linear isomorphism. This implies that $m=n$.

Theorem 8.1.5 (Inverse Function Theorem). Let $U \subseteq \mathbb{R}^{n}$ be an open subset, $x_{0} \in U, k \in \mathbb{N} \cup\{\infty\}$, and $f: U \rightarrow \mathbb{R}^{n}$ a $C^{k}$-map for which the linear map $\mathrm{d} f\left(x_{0}\right)$ is invertible. Then there exists an open neighborhood $V$ of $x_{0}$ in $U$ for which $\left.f\right|_{V}: V \rightarrow f(V)$ is a $C^{k}$-diffeomorphism onto an open subset of $\mathbb{R}^{n}$.

Corollary 8.1.6. Let $U \subseteq \mathbb{R}^{n}$ be an open subset and $f: U \rightarrow \mathbb{R}^{n}$ be an injective $C^{k}$-map ( $k \geq 1$ ) for which $\mathrm{d} f(x)$ is invertible for each $x \in U$. Then $f(U)$ is open and $f: U \rightarrow f(U)$ is a $C^{k}$-diffeomorphism.

Proof. First, we use the Inverse Function Theorem to see that for each $x \in U$ the image $f(U)$ contains a neighborhood of $f(x)$, so that $f(U)$ is an open subset of $\mathbb{R}^{n}$. Since $f$ is injective, the inverse function $g=f^{-1}: f(U) \rightarrow U$ exists. Now we apply the Inverse Function Theorem again to see that for each $x \in U$ there exists a neighborhood of $f(x)$ in $f(U)$ on which $g$ is $C^{k}$. Therefore, $g$ is a $C^{k}$-map, and this means that $f$ is a $C^{k}$-diffeomorphism.

Example 8.1.7. That the injectivity assumption in Corollary 8.1.6 is crucial is shown by the following example, which is a real description of the complex exponential function. We consider the smooth map

$$
f: \mathbb{R}^{2} \rightarrow \mathbb{R}^{2}, \quad f\left(x_{1}, x_{2}\right)=\left(e^{x_{1}} \cos x_{2}, e^{x_{1}} \sin x_{2}\right)
$$

Then the matrix of $\mathrm{d} f(x)$ with respect to the canonical basis is

$$
[\mathrm{d} f(x)]=\left(\begin{array}{cc}
e^{x_{1}} \cos x_{2} & -e^{x_{1}} \sin x_{2} \\
e^{x_{1}} \sin x_{2} & e^{x_{1}} \cos x_{2}
\end{array}\right)
$$

Its determinant is $e^{2 x_{1}} \neq 0$, so that $\mathrm{d} f(x)$ is invertible for each $x \in \mathbb{R}^{2}$.\\
Polar coordinates immediately show that $f\left(\mathbb{R}^{2}\right)=\mathbb{R}^{2} \backslash\{(0,0)\}$, which is an open subset of $\mathbb{R}^{2}$, but the map $f$ is not injective because it is $2 \pi$-periodic in $x_{2}$ :

$$
f\left(x_{1}, x_{2}+2 \pi\right)=f\left(x_{1}, x_{2}\right)
$$

Therefore, the Inverse Function Theorem applies to each $x \in \mathbb{R}^{2}$, but $f$ is not a global diffeomorphism.

Remark 8.1.8. The best way to understand the Implicit Function Theorem is to consider the linear case first. Let $g: \mathbb{R}^{m} \times \mathbb{R}^{n} \rightarrow \mathbb{R}^{m}$ be a linear map. We are interested in conditions under which the equation $g(x, y)=0$ can be solved for $x$, i.e., there is a function $f: \mathbb{R}^{n} \rightarrow \mathbb{R}^{m}$ such that $g(x, y)=0$ is equivalent to $x=f(y)$.

Since we are dealing with linear maps, there are matrices $A \in M_{m}(\mathbb{R})$ and $B \in M_{m, n}(\mathbb{R})$ with

$$
g(x, y)=A x+B y \quad \text { for } \quad x \in \mathbb{R}^{m}, y \in \mathbb{R}^{n} .
$$

The unique solvability of the equation $g(x, y)=0$ for $x$ is equivalent to the unique solvability of the equation $A x=-B y$, which is equivalent to the invertibility of the matrix $A$. If $A \in \mathrm{GL}_{m}(\mathbb{R})$, we thus obtain the linear function

$$
f: \mathbb{R}^{n} \rightarrow \mathbb{R}^{m}, \quad f(y)=-A^{-1} B y
$$

for which $x=f(y)$ is equivalent to $g(x, y)=0$.\\
Theorem 8.1.9 (Implicit Function Theorem). Let $U \subseteq \mathbb{R}^{m} \times \mathbb{R}^{n}$ be an open subset and $g: U \rightarrow \mathbb{R}^{m}$ be a $C^{k}$-function, $k \in \mathbb{N} \cup\{\infty\}$. Further, let $\left(x_{0}, y_{0}\right) \in U$ with $g\left(x_{0}, y_{0}\right)=0$ such that the linear map

$$
\mathrm{d}_{1} g\left(x_{0}, y_{0}\right): \mathbb{R}^{m} \rightarrow \mathbb{R}^{m}, \quad v \mapsto \mathrm{~d} g\left(x_{0}, y_{0}\right)(v, 0)
$$

is invertible. Then there exist open neighborhoods $V_{1}$ of $x_{0}$ in $\mathbb{R}^{m}$ and $V_{2}$ of $y_{0}$ in $\mathbb{R}^{n}$ with $V_{1} \times V_{2} \subseteq U$, and a $C^{k}$-function $f: V_{2} \rightarrow V_{1}$ with $f\left(y_{0}\right)=x_{0}$ such that

$$
\left\{(x, y) \in V_{1} \times V_{2}: g(x, y)=0\right\}=\left\{(f(y), y): y \in V_{2}\right\}
$$

Definition 8.1.10 (Higher derivatives). For $k \geq 2$, a $C^{k}$-map $f: U \rightarrow \mathbb{R}^{m}$, and $U \subseteq \mathbb{R}^{n}$ open, higher derivatives are defined inductively by

$$
\begin{aligned}
& \mathrm{d}^{k} f(x)\left(h_{1}, \ldots, h_{k}\right) \\
& :=\lim _{t \rightarrow 0} \frac{1}{t}\left(\mathrm{~d}^{k-1} f\left(x+t h_{k}\right)\left(h_{1}, \ldots, h_{k-1}\right)-\mathrm{d}^{k-1} f(x)\left(h_{1}, \ldots, h_{k-1}\right)\right)
\end{aligned}
$$

We thus obtain continuous maps

$$
\mathrm{d}^{k} f: U \times\left(\mathbb{R}^{n}\right)^{k} \rightarrow \mathbb{R}^{m}
$$

In terms of concrete coordinates and the canonical basis $e_{1}, \ldots, e_{n}$ for $\mathbb{R}^{n}$, we then have

$$
\mathrm{d}^{k} f(x)\left(e_{i_{1}}, \ldots, e_{i_{k}}\right)=\frac{\partial^{k} f}{\partial x_{i_{k}} \cdots \partial x_{i_{1}}}(x)
$$

Let $V$ and $W$ be vector spaces. We recall that a map $\beta: V^{k} \rightarrow W$ is called $k$-linear if all the maps

$$
V \rightarrow W, \quad v \mapsto \beta\left(v_{1}, \ldots, v_{j-1}, v, v_{j+1}, \ldots, v_{k}\right)
$$

are linear. It is said to be symmetric if

$$
\beta\left(v_{\sigma(1)}, \ldots, v_{\sigma(k)}\right)=\beta\left(v_{1}, \ldots, v_{k}\right)
$$

holds for all permutations $\sigma \in S_{k}$.

Proposition 8.1.11. If $f \in C^{k}\left(U, \mathbb{R}^{m}\right)$ and $k \geq 2$, then the functions $\left(h_{1}, \ldots, h_{k}\right) \mapsto \mathrm{d}^{k} f(x)\left(h_{1}, \ldots, h_{k}\right), x \in U$, are symmetric $k$-linear maps.\\
Proof. From the definition, it follows inductively that $\left(\mathrm{d}^{k} f\right)(x)$ is linear in each argument $h_{i}$ because, if all other arguments are fixed, it is the differential of a $C^{1}$-function.

To verify the symmetry of $\left(\mathrm{d}^{k} f\right)(x)$, we may also proceed by induction. It suffices to show that for $h_{1}, \ldots, h_{k-2}$ fixed, the map

$$
\beta(v, w):=\left(\mathrm{d}^{k} f(x)\right)\left(h_{1}, \ldots, h_{k-2}, v, w\right)
$$

is symmetric. This map is the second derivative $\mathrm{d}^{2} F(x)$ of the function

$$
F(x):=\left(\mathrm{d}^{k-2} f\right)(x)\left(h_{1}, \ldots, h_{k-2}\right) .
$$

We may therefore assume that $k=2$.\\
In view of the bilinearity, it suffices to observe that the Schwarz Lemma implies

$$
\left(\mathrm{d}^{2} F\right)(x)\left(e_{j}, e_{i}\right)=\left(\frac{\partial^{2}}{\partial x_{i} x_{j}} F\right)(x)=\left(\frac{\partial^{2}}{\partial x_{j} x_{i}} F\right)(x)=\left(\mathrm{d}^{2} F\right)(x)\left(e_{i}, e_{j}\right)
$$

Theorem 8.1.12 (Taylor's Theorem). Let $U \subseteq \mathbb{R}^{n}$ be open and $f: U \rightarrow \mathbb{R}^{m}$ of class $C^{k+1}$. If $x+[0,1] h \subseteq U$, then we have the Taylor Formula

$$
\begin{aligned}
f(x+h)= & f(x)+\mathrm{d} f(x)(h)+\cdots+\frac{1}{k!} \mathrm{d}^{k} f(x)(h, \ldots, h) \\
& +\frac{1}{k!} \int_{0}^{1}(1-t)^{k}\left(\mathrm{~d}^{k+1} f(x+t h)\right)(h, \ldots, h) d t
\end{aligned}
$$

Proof. For each $i \in\{1, \ldots, m\}$, we consider the $C^{k+1}$-maps

$$
F:[0,1] \rightarrow \mathbb{R}, \quad F(t):=f_{i}(x+t h) \quad \text { with } F^{(k)}(t)=\mathrm{d}^{k} f_{i}(x+t h)(h, \ldots, h)
$$

and apply the Taylor Formula for functions $[0,1] \rightarrow \mathbb{R}$ to get

$$
F(1)=F(0)+\cdots+\frac{F^{(k)}(0)}{k!}+\frac{1}{k!} \int_{0}^{1}(1-t)^{k} F^{(k+1)}(t) d t
$$

\subsection*{Smooth Manifolds and Smooth Maps}
Before we turn to the concept of a smooth manifold, we recall the concept of a Hausdorff space. We assume, however, some familiarity with basic topological constructions and concepts, such as the quotient topology. A topological space ( $X, \tau$ ) is called a Hausdorff space if for two different points $x, y \in X$ there exist disjoint open subsets $O_{x}, O_{y}$ with $x \in O_{x}$ and $y \in O_{y}$. Recall that each metric space ( $X, d$ ) is Hausdorff.

Definition 8.2.1. Let $M$ be a topological space.\\
(a) A pair ( $\varphi, U$ ) consisting of an open subset $U \subseteq M$ and a homeomorphism $\varphi: U \rightarrow \varphi(U) \subseteq \mathbb{R}^{n}$ of $U$ onto an open subset of $\mathbb{R}^{n}$ is called an $n$-dimensional chart of $M$.\\
(b) Two $n$-dimensional charts ( $\varphi, U$ ) and ( $\psi, V$ ) of $M$ are said to be $C^{k}$-compatible ( $k \in \mathbb{N} \cup\{\infty\}$ ) if $U \cap V=\emptyset$ or the map

$$
\left.\psi \circ \varphi^{-1}\right|_{\varphi(U \cap V)}: \varphi(U \cap V) \rightarrow \psi(U \cap V)
$$

is a $C^{k}$-diffeomorphism. Since $\varphi: U \rightarrow \varphi(U)$ is a homeomorphism onto an open subset of $\mathbb{R}^{n}, \varphi(U \cap V)$ is an open subset of $\varphi(U)$ and hence of $\mathbb{R}^{n}$.\\
(c) An $n$-dimensional $C^{k}$-atlas of $M$ is a family $\mathcal{A}:=\left(\varphi_{i}, U_{i}\right)_{i \in I}$ of $n$ dimensional charts of $M$ with the following properties:\\
(A1) $\bigcup_{i \in I} U_{i}=M$, i.e., $\left(U_{i}\right)_{i \in I}$ is an open covering of $M$.\\
(A2) All charts $\left(\varphi_{i}, U_{i}\right), i \in I$, are pairwise $C^{k}$-compatible. For $U_{i j}:=U_{i} \cap U_{j}$, this means that all maps

$$
\varphi_{j i}:=\left.\varphi_{j} \circ \varphi_{i}^{-1}\right|_{\varphi_{i}\left(U_{i j}\right)}: \varphi_{i}\left(U_{i j}\right) \rightarrow \varphi_{j}\left(U_{i j}\right)
$$

are $C^{k}$-maps because $\varphi_{j i}^{-1}=\varphi_{i j}$.\\
(d) A chart ( $\varphi, U$ ) is called compatible with a $C^{k}$-atlas $\left(\varphi_{i}, U_{i}\right)_{i \in I}$ if it is $C^{k}$-compatible with all charts of the atlas $\mathcal{A}$. A $C^{k}$-atlas $\mathcal{A}$ is called maximal if it contains all charts compatible with it. A maximal $C^{k}$-atlas is also called a $C^{k}$-differentiable structure on $M$. For $k=\infty$, we also call it a smooth structure.

Remark 8.2.2. (a) In Definition 8.2.1(b), we required that the map

$$
\left.\psi \circ \varphi^{-1}\right|_{\varphi(U \cap V)}: \varphi(U \cap V) \rightarrow \psi(U \cap V)
$$

were a $C^{k}$-diffeomorphism. Since $\varphi$ and $\psi$ are homeomorphisms, this map is always a homeomorphism between open subsets of $\mathbb{R}^{n}$. The differentiability is an additional requirement.\\
(b) For $M=\mathbb{R}$, the maps ( $M, \varphi$ ) and ( $M, \psi$ ) with $\varphi(x)=x$ and $\psi(x)=x^{3}$ are 1 -dimensional charts. These charts are not $C^{1}$-compatible: the map

$$
\psi \circ \varphi^{-1}: \mathbb{R} \rightarrow \mathbb{R}, \quad x \mapsto x^{3}
$$

is smooth, but not a diffeomorphism, since its inverse $\varphi \circ \psi^{-1}$ is not differentiable.\\
(c) Every atlas $\mathcal{A}$ is contained in a unique maximal atlas: We simply add all charts compatible with $\mathcal{A}$, and thus obtain a maximal atlas. This atlas is unique (Exercise 8.2.2).

Definition 8.2.3. An $n$-dimensional $C^{k}$-manifold is a pair ( $M, \mathcal{A}$ ) of a Hausdorff space $M$ and a maximal $n$-dimensional $C^{k}$-atlas $\mathcal{A}$ for $M$. For $k=\infty$, we call it a smooth manifold.

To specify a manifold structure, it suffices to specify a $C^{k}$-atlas $\mathcal{A}$ because this atlas is contained in a unique maximal one (Exercise 8.2.2). In the following, we shall never describe a maximal atlas. We shall always try to keep the number of charts as small as possible. For simplicity, we always assume in the following that $k=\infty$.

\subsubsection*{Examples}
Example 8.2.4 (Open subsets of $\mathbb{R}^{n}$ ). Let $U \subseteq \mathbb{R}^{n}$ be an open subset. Then $U$ is a Hausdorff space with respect to the induced topology. The inclusion map $\varphi: U \rightarrow \mathbb{R}^{n}$ defines a chart ( $\varphi, U$ ) which already defines a smooth atlas of $U$, turning $U$ into an $n$-dimensional smooth manifold.

Example 8.2.5 (The $n$-dimensional sphere). We consider the unit sphere

$$
\mathbb{S}^{n}:=\left\{\left(x_{0}, \ldots, x_{n}\right) \in \mathbb{R}^{n+1}: x_{0}^{2}+x_{1}^{2}+\cdots+x_{n}^{2}=1\right\}
$$

in $\mathbb{R}^{n}$, endowed with the subspace topology, turning it into a compact space.\\
(a) To specify a smooth manifold structure on $\mathbb{S}^{n}$, we consider the open subsets

$$
U_{i}^{\varepsilon}:=\left\{x \in \mathbb{S}^{n}: \varepsilon x_{i}>0\right\}, \quad i=0, \ldots, n, \quad \varepsilon \in\{ \pm 1\}
$$

These $2(n+1)$ subsets form a covering of $\mathbb{S}^{n}$. We have homeomorphisms

$$
\varphi_{i}^{\varepsilon}: U_{i}^{\varepsilon} \rightarrow B:=\left\{x \in \mathbb{R}^{n}:\|x\|_{2}<1\right\}
$$

onto the open unit ball in $\mathbb{R}^{n}$, given by

$$
\varphi_{i}^{\varepsilon}(x)=\left(x_{0}, x_{1}, \ldots, x_{i-1}, x_{i+1}, \ldots, x_{n}\right)
$$

and with continuous inverse map

$$
\left(y_{1}, \ldots, y_{n}\right) \mapsto\left(y_{1}, \ldots, y_{i}, \varepsilon \sqrt{1-\|y\|_{2}^{2}}, y_{i+1}, \ldots, y_{n}\right)
$$

This leads to charts $\left(\varphi_{i}^{\varepsilon}, U_{i}^{\varepsilon}\right)$ of $\mathbb{S}^{n}$.\\
It is easy to see that these charts are pairwise compatible. We have $\varphi_{i}^{\varepsilon} \circ \left(\varphi_{i}^{\varepsilon^{\prime}}\right)^{-1}=\operatorname{id}_{B}$, and for $i<j$, we have

$$
\varphi_{i}^{\varepsilon} \circ\left(\varphi_{j}^{\varepsilon^{\prime}}\right)^{-1}(y)=\left(y_{1}, \ldots, y_{i}, y_{i+2}, \ldots, y_{j}, \varepsilon^{\prime} \sqrt{1-\|y\|_{2}^{2}}, y_{j+1}, \ldots, y_{n}\right),
$$

which is a smooth map

$$
\varphi_{j}^{\varepsilon^{\prime}}\left(U_{i}^{\varepsilon} \cap U_{j}^{\varepsilon^{\prime}}\right) \rightarrow \varphi_{i}^{\varepsilon}\left(U_{i}^{\varepsilon} \cap U_{j}^{\varepsilon^{\prime}}\right)
$$

(b) There is another atlas of $\mathbb{S}^{n}$ consisting only of two charts, where the maps are slightly more complicated.

We call the unit vector $e_{0}:=(1,0, \ldots, 0)$ the north pole of the sphere and $-e_{0}$ the south pole. We then have the corresponding stereographic projection maps

$$
\varphi_{+}: U_{+}:=\mathbb{S}^{n} \backslash\left\{e_{0}\right\} \rightarrow \mathbb{R}^{n}, \quad\left(y_{0}, y\right) \mapsto \frac{1}{1-y_{0}} y
$$

and

$$
\varphi_{-}: U_{-}:=\mathbb{S}^{n} \backslash\left\{-e_{0}\right\} \rightarrow \mathbb{R}^{n}, \quad\left(y_{0}, y\right) \mapsto \frac{1}{1+y_{0}} y
$$

Both maps are bijective with inverse maps

$$
\varphi_{ \pm}^{-1}(x)=\left( \pm \frac{\|x\|_{2}^{2}-1}{\|x\|_{2}^{2}+1}, \frac{2 x}{1+\|x\|_{2}^{2}}\right)
$$

(Exercise 8.2.8). This implies that ( $\varphi_{+}, U_{+}$) and ( $\varphi_{-}, U_{-}$) are charts of $\mathbb{S}^{n}$. That both are smoothly compatible, hence form a smooth atlas, follows from

$$
\left(\varphi_{+} \circ \varphi_{-}^{-1}\right)(x)=\left(\varphi_{-} \circ \varphi_{+}^{-1}\right)(x)=\frac{x}{\|x\|^{2}}, \quad x \in \mathbb{R}^{n} \backslash\{0\}
$$

which is the inversion at the unit sphere.\\
Example 8.2.6. Let $E$ be an $n$-dimensional real vector space. We know from Linear Algebra that $E$ is isomorphic to $\mathbb{R}^{n}$, and that for each ordered basis $B:=\left(b_{1}, \ldots, b_{n}\right)$ for $E$, the linear map

$$
\varphi_{B}: \mathbb{R}^{n} \rightarrow E, \quad x=\left(x_{1}, \ldots, x_{n}\right) \mapsto \sum_{j=1}^{n} x_{j} b_{j}
$$

is a linear isomorphism. Using such a linear isomorphism $\varphi_{B}$, we define a topology on $E$ in such a way that $\varphi_{B}$ is a homeomorphism, i.e., $O \subseteq E$ is open if and only if $\varphi_{B}^{-1}(O)$ is open in $\mathbb{R}^{n}$.

For any other choice of a basis $C=\left(c_{1}, \ldots, c_{m}\right)$ in $E$ we recall from linear algebra that $m=n$ and that the map

$$
\varphi_{C}^{-1} \circ \varphi_{B}: \mathbb{R}^{n} \rightarrow \mathbb{R}^{n}
$$

is a linear isomorphism, hence a homeomorphism. This implies that for a subset $O \subseteq E$ the condition that $\varphi_{B}^{-1}(O)$ is open is equivalent to $\varphi_{C}^{-1}(O)= \varphi_{C}^{-1} \circ \varphi_{B} \circ \varphi_{B}^{-1}(O)$ being open. We conclude that the topology introduced on $E$ by $\varphi_{B}$ does not depend on the choice of a basis.

We thus obtain on $E$ a natural topology for which it is homeomorphic to $\mathbb{R}^{n}$, hence in particular a Hausdorff space. From each coordinate map $\kappa_{B}:=\varphi_{B}^{-1}$ we obtain a chart ( $\kappa_{B}, E$ ) which already defines an atlas of $E$. We thus obtain on $E$ the structure of an $n$-dimensional smooth manifold. That all these charts are compatible follows from the smoothness of the linear maps $\kappa_{C} \circ \kappa_{B}^{-1}=\varphi_{C}^{-1} \circ \varphi_{B}: \mathbb{R}^{n} \rightarrow \mathbb{R}^{n}$.

Example 8.2.7 (Submanifolds of $\mathbb{R}^{n}$ ). A subset $M \subseteq \mathbb{R}^{n}$ is called a $d$-dimensional submanifold if for each $p \in M$ there exists an open neighborhood $U$ of $p$ in $\mathbb{R}^{n}$ and a diffeomorphism

$$
\varphi: U \rightarrow \varphi(U) \subseteq \mathbb{R}^{n}
$$

onto an open subset $\varphi(U)$ with

$$
\varphi(U \cap M)=\varphi(U) \cap\left(\mathbb{R}^{d} \times\{0\}\right) .
$$

Whenever this condition is satisfied, we call ( $\varphi, U$ ) a submanifold chart.\\
A submanifold of codimension 1, i.e., $\operatorname{dim} M=n-1$, is called a smooth hypersurface.

We claim that $M$ carries a natural $d$-dimensional manifold structure when endowed with the topology inherited from $\mathbb{R}^{n}$, which obviously turns it into a Hausdorff space.

In fact, for each submanifold chart ( $\varphi, U$ ), we obtain a $d$-dimensional chart

$$
\left(\left.\varphi\right|_{U \cap M}, U \cap M\right),
$$

where we have identified $\mathbb{R}^{d}$ with $\mathbb{R}^{d} \times\{0\}$. For two such charts coming from ( $\varphi, U$ ) and ( $\psi, V$ ), we have

$$
\left.\psi \circ \varphi^{-1}\right|_{\varphi(U \cap V \cap M)}=\left.\left(\left.\psi\right|_{V \cap M}\right) \circ\left(\left.\varphi\right|_{U \cap M}\right)^{-1}\right|_{\varphi(U \cap V \cap M)},
$$

which is a smooth map onto an open subset of $\mathbb{R}^{d}$. We thus obtain a smooth atlas of $M$.

The following proposition provides a particularly handy criterion to verify that the set of solutions of a nonlinear equation is a submanifold.

Definition 8.2.8. Let $f: U \rightarrow \mathbb{R}^{m}$ be a $C^{1}$-map. We call $y \in \mathbb{R}^{m}$ a regular value of $f$ if for each $x \in U$ with $f(x)=y$ the differential $\mathrm{d} f(x)$ is surjective. Otherwise $y$ is called a singular value of $f$. Note that, in particular, each $y \in \mathbb{R}^{m} \backslash f(U)$ is a regular value.

Proposition 8.2.9 (Regular Value Theorem-Local Version). Let $U \subseteq \mathbb{R}^{n}$ be an open subset, $f: U \rightarrow \mathbb{R}^{m}$ a smooth map and $y \in \mathbb{R}^{m}$ a regular value of $f$. Then $M:=f^{-1}(y)$ is an ( $n-m$ )-dimensional submanifold of $\mathbb{R}^{n}$, hence in particular a smooth manifold.

Proof. Let $d:=n-m$ and observe that $d \geq 0$ because $\mathrm{d} f(x): \mathbb{R}^{n} \rightarrow \mathbb{R}^{m}$ is surjective for each $x \in M$. We have to show that for each $x_{0} \in M$ there exists an open neighborhood $V$ of $x_{0}$ in $\mathbb{R}^{n}$ and a diffeomorphism

$$
\varphi: V \rightarrow \varphi(V) \subseteq \mathbb{R}^{n}
$$

with

$$
\varphi(V \cap M)=\varphi(V) \cap\left(\mathbb{R}^{d} \times\{0\}\right) .
$$

After a permutation of the coordinates, we may w.l.o.g. assume that the vectors

$$
\mathrm{d} f\left(x_{0}\right)\left(e_{d+1}\right), \ldots, \mathrm{d} f\left(x_{0}\right)\left(e_{n}\right)
$$

form a basis for $\mathbb{R}^{m}$. Then we consider the map

$$
\varphi: U \rightarrow \mathbb{R}^{n}, \quad x=\left(x_{1}, \ldots, x_{n}\right) \mapsto\left(x_{1}, \ldots, x_{d}, f_{1}(x)-y_{1}, \ldots, f_{m}(x)-y_{m}\right) .
$$

In view of

$$
\mathrm{d} \varphi\left(x_{0}\right)\left(e_{j}\right)= \begin{cases}\left(e_{j}, \mathrm{~d} f\left(x_{0}\right) e_{j}\right) & \text { for } j \leq d \\ \mathrm{~d} f\left(x_{0}\right)\left(e_{j}\right) & \text { for } j>d\end{cases}
$$

it follows that the linear map $\mathrm{d} \varphi\left(x_{0}\right): \mathbb{R}^{n} \rightarrow \mathbb{R}^{n}$ is invertible. Hence the Inverse Function Theorem implies the existence of an open neighborhood $V \subseteq U$ of $x_{0}$ for which $\left.\varphi\right|_{V}: V \rightarrow \varphi(V)$ is a diffeomorphism onto an open subset of $\mathbb{R}^{n}$.

Since

$$
M=\left\{p \in U: \varphi(p)=\left(\varphi_{1}(p), \ldots, \varphi_{d}(p), 0, \ldots, 0\right)\right\}=\varphi^{-1}\left(\mathbb{R}^{d} \times\{0\}\right),
$$

it follows that $\varphi(M \cap V)=\varphi(V) \cap\left(\mathbb{R}^{d} \times\{0\}\right)$.\\
Example 8.2.10. The preceding proposition is particularly easy to apply for hypersurfaces, i.e., to the case $m=1$. Then $f: U \rightarrow \mathbb{R}$ is a smooth function and the condition that $\mathrm{d} f(x)$ is surjective simply means that $\mathrm{d} f(x) \neq 0$, i.e., that there exists some $j$ with $\frac{\partial f}{\partial x_{j}}(x) \neq 0$.\\
(a) Let $A=A^{\top} \in M_{n}(\mathbb{R})$ be a symmetric matrix and

$$
f(x):=x^{\top} A x=\sum_{i, j=1}^{n} a_{i j} x_{i} x_{j}
$$

the corresponding quadratic form. We want to show that the corresponding quadric

$$
Q:=\left\{x \in \mathbb{R}^{n}: f(x)=1\right\}
$$

is a submanifold of $\mathbb{R}^{n}$. To verify the criterion from Proposition 8.2.9, we assume that $f(x)=1$ and note that

$$
\mathrm{d} f(x) v=v^{\top} A x+x^{\top} A v=2 v^{\top} A x
$$

(Exercise; use Exercise 8.2.11). Therefore, $\mathrm{d} f(x)=0$ is equivalent to $A x=0$, which is never the case if $x^{\top} A x=1$. We conclude that all level surfaces of $f$ are smooth hypersurfaces of $\mathbb{R}^{n}$.

For $A=E_{n}$ (the identity matrix), we obtain the ( $n-1$ )-dimensional unit sphere $Q=\mathbb{S}^{n-1}$.

For $A=\operatorname{diag}\left(\lambda_{1}, \ldots, \lambda_{n}\right)$ and nonzero $\lambda_{i}$ we obtain the hyperboloids

$$
Q=\left\{x \in \mathbb{R}^{n}: \sum_{i=1}^{n} \lambda_{i} x_{i}^{2}=1\right\}
$$

which degenerate to hyperbolic cylinders if some $\lambda_{i}$ vanish.\\
(b) For singular values, the level sets may or may not be submanifolds: For the quadratic form

$$
f: \mathbb{R}^{2} \rightarrow \mathbb{R}, \quad f\left(x_{1}, x_{2}\right)=x_{1} x_{2}
$$

the value 0 is singular because $f(0,0)=0$ and $\mathrm{d} f(0,0)=0$. The inverse image is

$$
f^{-1}(0)=(\mathbb{R} \times\{0\}) \cup(\{0\} \times \mathbb{R})
$$

which is not a submanifold of $\mathbb{R}^{2}$ (Exercise).\\
For the quadratic form

$$
f: \mathbb{R}^{2} \rightarrow \mathbb{R}, \quad f\left(x_{1}, x_{2}\right)=x_{1}^{2}+x_{2}^{2}
$$

the value 0 is singular because $f(0,0)=0$ and $\mathrm{d} f(0,0)=0$. The inverse image is

$$
f^{-1}(0)=\{(0,0)\}
$$

which is a zero-dimensional submanifold of $\mathbb{R}^{2}$.\\
For the quadratic form

$$
f: \mathbb{R} \rightarrow \mathbb{R}, \quad f(x)=x^{2}
$$

the value 0 is singular because $f(0)=0$ and $f^{\prime}(0)=0$. The inverse image is

$$
f^{-1}(0)=\{0\}
$$

which is a submanifold of $\mathbb{R}$.\\
(c) On $M_{n}(\mathbb{R}) \cong \mathbb{R}^{n^{2}}$, we consider the quadratic function

$$
f: M_{n}(\mathbb{R}) \rightarrow \operatorname{Sym}_{n}(\mathbb{R}):=\left\{A \in M_{n}(\mathbb{R}): A^{\top}=A\right\}, \quad X \mapsto X X^{\top}
$$

Then

$$
f^{-1}(\mathbf{1})=\mathrm{O}_{n}(\mathbb{R}):=\left\{g \in \mathrm{GL}_{n}(\mathbb{R}): g^{\top}=g^{-1}\right\}
$$

is the orthogonal group.\\
To see that this is a submanifold of $M_{n}(\mathbb{R})$, we note that

$$
\mathrm{d} f(X)(Y)=X Y^{\top}+Y X^{\top}
$$

(Exercise 8.2.11). If $f(X)=\mathbf{1}$, we have $X^{\top}=X^{-1}$, so that for any $Z \in \operatorname{Sym}_{n}(\mathbb{R})$ the matrix $Y:=\frac{1}{2} Z X$ satisfies

$$
X Y^{\top}+Y X^{\top}=\frac{1}{2}\left(X X^{\top} Z+Z X X^{\top}\right)=Z .
$$

Therefore, $\mathrm{d} f(X)$ is surjective for each orthogonal matrix $X$, and Proposition 8.2.9 implies that $\mathrm{O}_{n}(\mathbb{R})$ is a submanifold of $M_{n}(\mathbb{R})$ of dimension

$$
d=n^{2}-\operatorname{dim}\left(\operatorname{Sym}_{n}(\mathbb{R})\right)=n^{2}-\frac{n(n+1)}{2}=\frac{n(n-1)}{2} .
$$

Remark 8.2.11 (The gluing picture). Let $M$ be an $n$-dimensional manifold with an atlas $\mathcal{A}=\left(\varphi_{i}, U_{i}\right)_{i \in A}$ and $V_{i}:=\varphi_{i}\left(U_{i}\right)$ the corresponding open subsets of $\mathbb{R}^{n}$.

Note that we have used the topology of $M$ to define the notion of a chart. We now explain how the topological space $M$ can be reconstructed from the atlas $\mathcal{A}$. First, we consider the set

$$
S:=\bigcup_{i \in I}\{i\} \times V_{i},
$$

which we consider as the disjoint union of the open subset $V_{i} \subseteq \mathbb{R}^{n}$. We endow $S$ with the topology of the disjoint sum, i.e., a subset $O \subseteq S$ is open if and only if all its intersections with the subsets $\{i\} \times V_{i} \cong V_{i}$ are open.

Then we consider the surjective map

$$
\Phi: S \rightarrow M, \quad(i, x) \mapsto \varphi_{i}^{-1}(x) .
$$

On each subset $\{i\} \times V_{i}$, this map is a homeomorphism onto $U_{i}$. Hence $\Phi$ is continuous, surjective and open, which implies that it is a quotient map, i.e., that the topology on $M$ coincides with the quotient topology on $S / \sim$, where

$$
(i, x) \sim(j, y) \quad \Longleftrightarrow \quad \varphi_{i}\left(\varphi_{j}^{-1}(y)\right)=x .
$$

In this sense, we can think of $M$ as obtained by gluing of the patches $U_{i} \cong V_{i}$, where $U_{i j}=U_{i} \cap U_{j}$ and $x_{i} \in \varphi_{i}\left(U_{i j}\right) \subseteq V_{i}$ is identified with the point $x_{j}=\varphi_{j}\left(\varphi_{i}^{-1}\left(x_{i}\right)\right) \in V_{j}$.

Example 8.2.12. We discuss an example of a "non-Hausdorff manifold". We endow the set $S:=(\{1\} \times \mathbb{R}) \cup(\{2\} \times \mathbb{R})$ with the disjoint sum topology and define an equivalence relation on $S$ by

$$
(1, x) \sim(2, y) \quad \Longleftrightarrow \quad x=y \neq 0 .
$$

If $[i, x]$ denotes the class of $(i, x)$ we see that all classes except $[1,0]$ and $[2,0]$ contain 2 points. The topological quotient space

$$
M:=S / \sim=\{[1, x]: x \in \mathbb{R}\} \cup\{[2,0]\}=\{[2, x]: x \in \mathbb{R}\} \cup\{[1,0]\}
$$

is the union of a real line with an extra point, but the two points $[1,0]$ and [2, 0] have no disjoint open neighborhoods.

The subsets $U_{j}:=\{[j, x]: x \in \mathbb{R}\}, j=1,2$, of $M$ are open, and the maps

$$
\varphi_{j}: U_{j} \rightarrow \mathbb{R}, \quad[j, x] \mapsto x,
$$

are homeomorphisms defining a smooth atlas on $M$ (Exercise 8.2.12).

Remark 8.2.13. One can also define manifold structures on sets carrying no a priori topology. To this end, one proceeds as follows. We start with a set $M$. An $n$-dimensional chart of $M$ is a pair ( $\varphi, U$ ), where $U \subseteq M$ is a subset and $\varphi: U \rightarrow \mathbb{R}^{n}$ an injection with an open image. Two charts $(\varphi, U)$ and ( $\psi, V$ ) are called $C^{k}$-compatible if both $\varphi(U \cap V)$ and $\psi(U \cap V)$ are open and

$$
\psi^{-1} \circ \varphi: \varphi(U \cap V) \rightarrow \psi(U \cap V)
$$

is a $C^{k}$-diffeomorphism. The notion of a $C^{k}$-atlas is introduced as in Definition 8.2.1(c).

Let $\mathcal{A}=\left(\varphi_{i}, U_{i}\right)_{i \in I}$ be an $n$-dimensional $C^{k}$-atlas on $M$. We call a subset $O \subseteq M$ open if for each $i \in I$ the subset $\varphi_{i}\left(O \cap U_{i}\right) \subseteq \mathbb{R}^{n}$ is open. It is easy to see that we thus obtain a topology on $M$. For any subset $O_{j} \subseteq U_{j}$, we have

$$
\varphi_{i}\left(O_{j} \cap U_{i}\right)=\left(\varphi_{i} \circ \varphi_{j}^{-1}\right)\left(\varphi_{j}\left(O_{j}\right) \cap \varphi_{j}\left(U_{i} \cap U_{j}\right)\right),
$$

showing that $O_{j}$ is open in $M$ if and only if $\varphi_{j}\left(O_{j}\right)$ is open in $\mathbb{R}^{n}$. Hence each chart ( $\varphi_{j}, U_{j}$ ) defines a homeomorphism $U_{j} \rightarrow \varphi\left(U_{j}\right)$. As Example 8.2.12 shows, the topology on $M$ need not be Hausdorff, but if it is, then the arguments above show that the pair ( $M, \mathcal{A}$ ) is an $n$-dimensional $C^{k}$-manifold.

A sufficient condition for the Hausdorff property to hold is that for any pair $x \neq y$ in $M$ either\\
(H1) there exists an $i \in I$ with $x, y \in U_{i}$, or\\
(H2) there exist $i, j \in I$ with $x \in U_{i}, y \in U_{j}$ and $U_{i} \cap U_{j}=\emptyset$.\\
In fact, in the second case, $x$ and $y$ are separated by the open sets $U_{i}$ and $U_{j}$; and in the first case, the fact that $U_{i}$ is Hausdorff implies the existence of disjoint open subsets of $U_{i}$, and hence of $M$, separating $x$ and $y$.

As an application of the preceding remark, we discuss the Graßmann manifolds of $k$-dimensional subspaces of an $n$-dimensional vector space.

Example 8.2.14 (Graßmannians). Let $V=\mathbb{R}^{n}$ and write $\operatorname{Gr}_{k}(V)= \operatorname{Gr}_{k}\left(\mathbb{R}^{n}\right)$ for the set of all $k$-dimensional linear subspaces of $V$. We now use the approach described in Remark 8.2.13 to define on $\operatorname{Gr}_{k}\left(\mathbb{R}^{n}\right)$ the structure of a smooth manifold of dimension $k(n-k)$.

For an ( $n-k$ )-dimensional subspace $F \subseteq V$, we consider the subset

$$
U_{F}:=\left\{E \in \operatorname{Gr}_{k}(V): F \oplus E=V\right\}
$$

of all subspaces complementing $F$. Fixing an element $E_{0} \in U_{F}$, we have $V=F \oplus E_{0}$, and each $E \in U_{F}$ can be written as the graph

$$
\Gamma(f)=\left\{(x, f(x)): x \in E_{0}\right\} \quad \text { of some } \quad f \in \operatorname{Hom}\left(E_{0}, F\right)
$$

In terms of the projections $\operatorname{pr}_{E_{0}}^{F}: V \rightarrow E_{0}$ and $\operatorname{pr}_{F}^{E_{0}}: V \rightarrow F$ along the decomposition $V=F \oplus E_{0}$, the linear map $f$ can be expressed as

$$
f=\operatorname{pr}_{F}^{E_{0}} \circ\left(\left.\operatorname{pr}_{E_{0}}^{F}\right|_{E}\right)^{-1}
$$

Therefore,

$$
\varphi_{F, E_{0}}: U_{F} \rightarrow \operatorname{Hom}\left(E_{0}, F\right),\left.\quad E \mapsto \operatorname{pr}_{F}^{E_{0}}\right|_{E} \circ\left(\left.\operatorname{pr}_{E_{0}}^{F}\right|_{E}\right)^{-1}
$$

is a bijection. Replacing $E_{0}$ by another subspace $E_{1} \in U_{F}$, the relation

$$
\Gamma(f)=\Gamma(\widetilde{f})=E \quad \text { for } \quad f \in \operatorname{Hom}\left(E_{0}, F\right), \widetilde{f} \in \operatorname{Hom}\left(E_{1}, F\right),
$$

leads to the relation $x+f(x)=y+\widetilde{f}(y)$, where $x \in E_{0}$ and $y=\operatorname{pr}_{E_{1}}^{F}(x) \in E_{1}$. But then

$$
\widetilde{f}(y)=f(x)+x-y=f\left(\left(\left.\operatorname{pr}_{E_{1}}^{F}\right|_{E_{0}}\right)^{-1}(y)\right)+\left(\left.\operatorname{pr}_{E_{1}}^{F}\right|_{E_{0}}\right)^{-1}(y)-y .
$$

This means that

$$
\varphi_{F, E_{1}}(E)=\widetilde{f}=f \circ\left(\left.\operatorname{pr}_{E_{1}}^{F}\right|_{E_{0}}\right)^{-1}+\varphi_{F, E_{1}}\left(E_{0}\right) .
$$

Therefore, the transition between the charts $\varphi_{F, E_{1}}$ and $\varphi_{F, E_{0}}$ is given by an invertible affine map.

If two sets $U_{F}$ and $U_{F^{\prime}}$ intersect nontrivially, then the set of those $f \in \operatorname{Hom}\left(E_{0}, F\right)$ for which $F^{\prime} \oplus \Gamma(f)=V$ is an open subset of $\operatorname{Hom}\left(E_{0}, F\right)$. In fact, if $b_{1}, \ldots, b_{k}$ is a basis for $E_{0}$ and $c_{1}^{\prime}, \ldots, c_{n-k}^{\prime}$ is a basis for $F^{\prime}$, then the condition on $f$ is that

$$
\operatorname{det}\left(f\left(b_{1}\right), \ldots, f\left(b_{k}\right), c_{1}^{\prime}, \ldots, c_{n-k}^{\prime}\right) \neq 0
$$

which specifies an open subset of $\operatorname{Hom}\left(E_{0}, F\right) \cong \mathbb{R}^{k(n-k)}$.\\
We further note that, for $E_{0} \in U_{F} \cap U_{F^{\prime}}$, we have

$$
\begin{aligned}
\varphi_{F^{\prime}, E_{0}}(E) & =\operatorname{pr}_{F^{\prime}}^{E_{0}} \circ\left(\left.\operatorname{pr}_{E_{0}}^{F^{\prime}}\right|_{E}\right)^{-1}=\left(\left.\operatorname{pr}_{F^{\prime}}^{E}\right|_{F}\right) \circ \operatorname{pr}_{F}^{E_{0}} \circ\left(\left.\operatorname{pr}_{E_{0}}^{F}\right|_{E}\right)^{-1} \\
& =\left(\left.\operatorname{pr}_{F^{\prime}}^{E}\right|_{F}\right) \circ \varphi_{F, E_{0}}(E) .
\end{aligned}
$$

Hence all these coordinate changes are given by restrictions of invertible linear maps. This proves that the collection $\left(\varphi_{F, E_{0}}, U_{F}\right)_{F \oplus E_{0}=V}$ yields a smooth $k(n-k)$-dimensional atlas of $\operatorname{Gr}_{k}(V)$ if we identify all the spaces $\operatorname{Hom}\left(E_{0}, F\right)$ with $\mathbb{R}^{n(k-n)}$ by some linear isomorphism.

In view of Remark 8.2.13, it now suffices to verify that the topology on $\operatorname{Gr}_{k}(V)$ defined by our smooth atlas is Hausdorff and (H1) shows that it suffices to find for $E_{0}, E_{1} \in \operatorname{Gr}_{k}(V)$ a subspace $F$ with $E_{0}, E_{1} \in U_{F}$.

Writing $V=F_{1} \oplus\left(E_{0} \cap E_{1}\right)$ for some subspace $F_{1}$, we have to find a subspace $F_{2}$ of $F_{1} \cap\left(E_{0}+E_{1}\right)$ complementing $E_{0} \cap F_{1}$ and $E_{1} \cap F_{1}$. Since the latter two spaces intersect trivially, we have (Exercise)

$$
F_{1} \cap\left(E_{0}+E_{1}\right)=\left(E_{0} \cap F_{1}\right) \oplus\left(E_{1} \cap F_{1}\right) .
$$

Because of $\left(E_{0} \cap F_{1}\right) \oplus\left(E_{0} \cap E_{1}\right)=E_{0}$, using standard dimension formulas from linear algebra, one obtains a linear isomorphism

$$
\varphi: E_{0} \cap F_{1} \rightarrow E_{1} \cap F_{1}
$$

Then its graph $\Gamma(\varphi)$ is a linear subspace of $F_{1} \cap\left(E_{0}+E_{1}\right)$ complementing both subspaces $E_{i} \cap F_{1}$.\\
Example 8.2.15 (Projective space). As an important special case of a Graßmann manifold we obtain the real projective space

$$
\mathbb{P}\left(\mathbb{R}^{n}\right):=\operatorname{Gr}_{1}\left(\mathbb{R}^{n}\right)
$$

of all 1 -dimensional subspaces of $\mathbb{R}^{n}$. It is a smooth manifold of dimension $n-1$. For $n=2$, this space is called the projective line and for $n=3$ it is called the projective plane (it is a 2 -dimensional manifold, thus also called a surface).

We write $[x]:=\mathbb{R} x \in \mathbb{P}\left(\mathbb{R}^{n}\right)$ for the subspace generated by a nonzero element $x \in \mathbb{R}^{n}$. Since each one-dimensional subspace $\mathbb{R} x$ of $\mathbb{R}^{n}$ is not contained in some of the hyperplanes $F_{i}:=\operatorname{span}\left\{e_{1}, \ldots, e_{i-1}, e_{i+1}, \ldots, e_{n}\right\}$, the manifold $\mathbb{P}\left(\mathbb{R}^{n}\right)$ is covered by the $n$ open sets

$$
U_{i}:=U_{F_{i}}=\left\{[x]: x_{i} \neq 0\right\}, \quad i=1, \ldots, n
$$

The simplest charts are obtained by picking the complement $E_{0}=\mathbb{R} e_{i}$ of the hyperplane $F_{i}$ and identifying $\operatorname{Hom}\left(E_{0}, F_{i}\right) \cong F_{i} \cong \mathbb{R}^{n-1}$. For any $[x] \in U_{i}$, we then have

$$
\mathbb{R} x=\Gamma(f) \quad \text { with } \quad f\left(e_{i}\right)=\left(\frac{x_{1}}{x_{i}}, \ldots, \frac{x_{i-1}}{x_{i}}, \frac{x_{i+1}}{x_{i}}, \ldots, \frac{x_{n}}{x_{i}}\right)
$$

which leads to the coordinates

$$
\varphi_{i}([x])=\left(\frac{x_{1}}{x_{i}}, \ldots, \frac{x_{i-1}}{x_{i}}, \frac{x_{i+1}}{x_{i}}, \ldots, \frac{x_{n}}{x_{i}}\right)
$$

These charts are called homogeneous coordinates. They play a fundamental role in projective geometry.\\
Example 8.2.16. We discuss an example of a manifold which is not separable. In particular, its topology has no countable basis. We define a new topology on $\mathbb{R}^{2}$ by defining $O \subseteq \mathbb{R}^{2}$ to be open if and only if for each $y \in \mathbb{R}$ the set

$$
O^{y}:=\{x \in \mathbb{R}:(x, y) \in O\}
$$

is open in $\mathbb{R}$. This defines a topology on $\mathbb{R}^{2}$ for which all sets $U_{y}:=\mathbb{R} \times\{y\}$ are open and the maps

$$
\varphi_{y}: U_{y} \rightarrow \mathbb{R}, \quad(x, y) \mapsto x
$$

are homeomorphisms. We thus obtain a smooth 1 -dimensional manifold structure on $\mathbb{R}^{2}$ for which it has uncountably many connected components, namely the subsets $U_{y}, y \in \mathbb{R}$.

Remark 8.2.17 (Coordinates versus parameterizations). (a) Let $(\varphi, U)$ be an $n$-dimensional chart of the smooth manifold $M$. Then $\varphi: U \rightarrow \mathbb{R}^{n}$ has $n$ components $\varphi_{1}, \ldots, \varphi_{n}$ which we consider as coordinate functions on $U$. Sometimes it is convenient to write $x_{i}(p):=\varphi_{i}(p)$ for $p \in U$, so that $\left(x_{1}(p), \ldots, x_{n}(p)\right)$ are the coordinates of $p \in U$ w.r.t. the chart $(\varphi, U)$.

If we have another chart ( $\psi, V$ ) of $M$ with $U \cap V \neq \emptyset$, then any $p \in U \cap V$ has a second tuple of coordinates, $x_{i}^{\prime}(p):=\psi_{i}(p)$, given by the components of $\psi$. Now the change of coordinates is given by

$$
x^{\prime}(x)=\psi\left(\varphi^{-1}(x)\right) \quad \text { and } \quad x\left(x^{\prime}\right)=\varphi\left(\psi^{-1}\left(x^{\prime}\right)\right) .
$$

In this sense, the maps $\psi \circ \varphi^{-1}$ and $\varphi \circ \psi^{-1}$ describe how we translate between the $x$-coordinates and the $x^{\prime}$-coordinates.\\
(b) Instead of putting the focus on coordinates, which are functions on open subsets of the manifold, one can also parameterize open subset of $M$. This is done by maps $\varphi: V \rightarrow M$, where $V$ is an open subsets of some $\mathbb{R}^{n}$ and ( $\varphi^{-1}, \varphi(V)$ ) is a chart of $M$. Then the point $p \in M$ corresponding to the parameter values $\left(x_{1}, \ldots, x_{n}\right) \in V$ is $p=\varphi(x)$. In this picture, the lines

$$
t \mapsto \varphi\left(x_{1}, \ldots, x_{i-1}, t, x_{i+1}, \ldots, x_{n}\right)
$$

are curves on $M$, called the parameter lines.

\subsubsection*{New Manifolds from Old Ones}
Definition 8.2.18 (Open subsets are manifolds). Let $M$ be a smooth manifold and $N \subseteq M$ an open subset. Then $N$ carries a natural smooth manifold structure.

Let $\mathcal{A}=\left(\varphi_{i}, U_{i}\right)_{i \in I}$ be an atlas of $M$. Then $V_{i}:=U_{i} \cap N$ and $\psi_{i}:=\left.\varphi_{i}\right|_{V_{i}}$ define a smooth atlas $\mathcal{B}:=\left(\psi_{i}, V_{i}\right)_{i \in I}$ of $N$ (Exercise).\\
Definition 8.2.19 (Products of manifolds). Let $M$ and $N$ be smooth manifolds of dimensions $d$, resp., $k$, and

$$
M \times N=\{(m, n): m \in M, n \in N\}
$$

the product set, which we endow with the product topology.\\
We show that $M \times N$ carries a natural structure of a smooth ( $d+k$ )-dimensional manifold. Let $\mathcal{A}=\left(\varphi_{i}, U_{i}\right)_{i \in I}$ be an atlas of $M$ and $\mathcal{B}=\left(\psi_{j}, V_{j}\right)_{j \in J}$ an atlas of $N$. Then the product sets $W_{i j}:=U_{i} \times V_{j}$ are open in $M \times N$ and the maps

$$
\gamma_{i j}:=\varphi_{i} \times \psi_{j}: U_{i} \times V_{j} \rightarrow \mathbb{R}^{d} \times \mathbb{R}^{k} \cong \mathbb{R}^{d+k}, \quad(x, y) \mapsto\left(\varphi_{i}(x), \psi_{j}(y)\right)
$$

are homeomorphisms onto open subsets of $\mathbb{R}^{d+k}$. On $\gamma_{i^{\prime} j^{\prime}}\left(W_{i j} \cap W_{i^{\prime} j^{\prime}}\right)$, we have

$$
\gamma_{i j} \circ \gamma_{i^{\prime} j^{\prime}}^{-1}=\left(\varphi_{i} \circ \varphi_{i^{\prime}}^{-1}\right) \times\left(\psi_{j} \circ \psi_{j^{\prime}}^{-1}\right),
$$

which is a smooth map. Therefore, $\left(\varphi_{i j}, W_{i j}\right)_{(i, j) \in I \times J}$ is a smooth atlas on $M \times N$.

Definition 8.2.20. Let $M$ and $N$ be smooth manifolds.\\
(a) A continuous map $f: M \rightarrow N$ is called smooth, if for each chart $(\varphi, U)$ of $M$ and each chart ( $\psi, V$ ) of $N$ the map

$$
\psi \circ f \circ \varphi^{-1}: \varphi\left(f^{-1}(V) \cap U\right) \rightarrow \psi(V)
$$

is smooth. Note that $\varphi\left(f^{-1}(V) \cap U\right)$ is open because $f$ is continuous.\\
We write $C^{\infty}(M, N)$ for the set of smooth maps $M \rightarrow N$ and abbreviate $C^{\infty}(M, \mathbb{R})$ by $C^{\infty}(M)$.\\
(b) A map $f: M \rightarrow N$ is called a diffeomorphism, or a smooth isomorphism, if there exists a smooth map $g: N \rightarrow M$ with

$$
f \circ g=\operatorname{id}_{N} \quad \text { and } \quad g \circ f=\operatorname{id}_{M}
$$

This condition obviously is equivalent to $f$ being bijective and its inverse $f^{-1}$ being a smooth map.

We write $\operatorname{Diff}(M)$ for the set of all diffeomorphisms of $M$.\\
Lemma 8.2.21. Compositions of smooth maps are smooth. In particular, the set $\operatorname{Diff}(M)$ is a group (with respect to composition) for each smooth manifold $M .^{1}$

Proof. Let $f: M \rightarrow N$ and $g: N \rightarrow L$ be smooth maps. Pick charts $(\varphi, U)$ of $M$ and $(\gamma, W)$ of $L$. To see that the map $\gamma \circ(g \circ f) \circ \varphi^{-1}$ is smooth on $\varphi\left((g \circ f)^{-1}(W)\right)$, we have to show that each element $x=\varphi(p)$ in this set has a neighborhood on which it is smooth. Let $q:=f(p)$ and note that $g(q) \in W$. We choose a chart ( $\psi, V$ ) of $N$ with $q \in V$. We then have

$$
\gamma \circ(g \circ f) \circ \varphi^{-1}=\left(\gamma \circ g \circ \psi^{-1}\right) \circ\left(\psi \circ f \circ \varphi^{-1}\right)
$$

on the open neighborhood $\varphi\left(f^{-1}(V) \cap(g \circ f)^{-1}(W)\right)$ of $x$. Since compositions of smooth maps on open domains in $\mathbb{R}^{n}$ are smooth by the Chain Rule (Theorem 8.1.3), $\gamma \circ(g \circ f) \circ \varphi^{-1}$ is smooth on $\varphi\left((g \circ f)^{-1}(W)\right)$. This proves that $g \circ f: M \rightarrow L$ is a smooth map.

Remark 8.2.22. (a) If $I \subseteq \mathbb{R}$ is an open interval, then a smooth map $\gamma: I \rightarrow M$ is called a smooth curve.

For a not necessarily open interval $I \subseteq \mathbb{R}$, a map $\gamma: I \rightarrow \mathbb{R}^{n}$ is called smooth if all derivatives $\gamma^{(k)}$ exist at all points of $I$ and define continuous functions $I \rightarrow \mathbb{R}^{n}$. Based on this generalization of smoothness for curves, a curve $\gamma: I \rightarrow M$ is said to be smooth, if for each chart $(\varphi, U)$ of $M$ the curves

\footnotetext{${ }^{1}$ For each manifold $M$ the identity $\operatorname{id}_{M}: M \rightarrow M$ is a smooth map, so that this lemma leads to the "category of smooth manifolds". The objects of this category are smooth manifolds and the morphisms are the smooth maps. In the following, we shall use consistently category theoretical language, but we shall not go into the formal details of category theory.
}
$$
\varphi \circ \gamma: \gamma^{-1}(U) \rightarrow \mathbb{R}^{n}
$$
are smooth.\\
A curve $\gamma:[a, b] \rightarrow M$ is called piecewise smooth if $\gamma$ is continuous and there exists a subdivision $x_{0}=a<x_{1}<\cdots<x_{N}=b$ such that $\left.\gamma\right|_{\left[x_{i}, x_{i+1}\right]}$ is smooth for $i=0, \ldots, N-1$.\\
(b) Smoothness of maps $f: M \rightarrow \mathbb{R}^{n}$ can be checked more easily. Since the identity is a chart of $\mathbb{R}^{n}$, the smoothness condition simply means that for each chart ( $\varphi, U$ ) of $M$ the map
$$
f \circ \varphi^{-1}: \varphi\left(f^{-1}(V) \cap U\right) \rightarrow \mathbb{R}^{n}
$$
is smooth.\\
(c) If $U$ is an open subset of $\mathbb{R}^{n}$, then a map $f: U \rightarrow M$ to a smooth $m$-dimensional manifold $M$ is smooth if and only if for each chart ( $\varphi, V$ ) of $M$ the map
$$
\varphi \circ f: f^{-1}(V) \rightarrow \mathbb{R}^{n}
$$
is smooth.\\
(d) Any chart ( $\varphi, U$ ) of a smooth $n$-dimensional manifold $M$ defines a diffeomorphism $U \rightarrow \varphi(U) \subseteq \mathbb{R}^{n}$, when $U$ is endowed with the canonical manifold structure as an open subset of $M$.

In fact, by definition, we may use ( $\varphi, U$ ) as an atlas of $U$. Then the smoothness of $\varphi$ is equivalent to the smoothness of the map $\varphi \circ \varphi^{-1}=\operatorname{id}_{\varphi(U)}$, which is trivial. Likewise, the smoothness of $\varphi^{-1}: \varphi(U) \rightarrow U$ is equivalent to the smoothness of $\varphi \circ \varphi^{-1}=\operatorname{id}_{\varphi(U)}$.

\subsubsection*{Exercises for Section 8.2}
Exercise 8.2.1. Let $M:=\mathbb{R}$, endowed with its standard topology. Show that $C^{k}$-compatibility of 1 -dimensional charts is not an equivalence relation.

Exercise 8.2.2. Show that each $n$-dimensional $C^{k}$-atlas is contained in a unique maximal one.

Exercise 8.2.3. Let $M_{i}, i=1, \ldots, n$, be smooth manifolds of dimension $d_{i}$. Show that the product space $M:=M_{1} \times \cdots \times M_{n}$ carries the structure of a ( $d_{1}+\cdots+d_{n}$ )-dimensional manifold.

Exercise 8.2.4 (Relaxation of the smoothness definition). Let $M$ and $N$ be smooth manifolds. Show that a map $f: M \rightarrow N$ is smooth if and only if for each point $x \in M$ there exists a chart ( $\varphi, U$ ) of $M$ with $x \in U$ and a chart ( $\psi, V$ ) of $N$ with $f(x) \in V$ such that the map

$$
\psi \circ f \circ \varphi^{-1}: \varphi\left(f^{-1}(V)\right) \rightarrow \psi(V)
$$

is smooth.

Exercise 8.2.5. Show that the set $A:=C^{\infty}(M, \mathbb{R})$ of smooth real-valued functions on $M$ is a real algebra. If $g \in A$ is nonzero and $U:=g^{-1}\left(\mathbb{R}^{\times}\right)$, then $\frac{1}{g} \in C^{\infty}(U, \mathbb{R})$.

Exercise 8.2.6. Let $f_{1}: M_{1} \rightarrow N_{1}$ and $f_{2}: M_{2} \rightarrow N_{2}$ be smooth maps. Show that the map

$$
f_{1} \times f_{2}: M_{1} \times M_{2} \rightarrow N_{1} \times N_{2}, \quad(x, y) \mapsto\left(f_{1}(x), f_{2}(y)\right)
$$

is smooth.\\
Exercise 8.2.7. Let $f_{1}: M \rightarrow N_{1}$ and $f_{2}: M \rightarrow N_{2}$ be smooth maps. Show that the map

$$
\left(f_{1}, f_{2}\right): M \rightarrow N_{1} \times N_{2}, \quad x \mapsto\left(f_{1}(x), f_{2}(x)\right)
$$

is smooth.\\
Exercise 8.2.8. (a) Verify the details in Example 8.2.5, where we describe an atlas of $\mathbb{S}^{n}$ by stereographic projections.\\
(b) Show that the two atlases of $\mathbb{S}^{n}$ constructed in Example 8.2.5 and the atlas obtained from the realization of $\mathbb{S}^{n}$ as a quadric in $\mathbb{R}^{n+1}$ define the same differentiable structure.

Exercise 8.2.9. Let $N$ be an open subset of the smooth manifold $M$. Show that if $\mathcal{A}=\left(\varphi_{i}, U_{i}\right)_{i \in I}$ is a smooth atlas of $M, V_{i}:=U_{i} \cap N$ and $\psi_{i}:=\left.\varphi_{i}\right|_{V_{i}}$, then $\mathcal{B}:=\left(\psi_{i}, V_{i}\right)_{i \in I}$ is a smooth atlas of $N$.

Exercise 8.2.10. Smoothness is a local property: Show that a map $f: M \rightarrow N$ between smooth manifolds is smooth if and only if for each $p \in M$ there is an open neighborhood $U$ such that $\left.f\right|_{U}$ is smooth.

Exercise 8.2.11. Let $V_{1}, \ldots, V_{k}$ and $V$ be finite-dimensional real vector space and

$$
\beta: V_{1} \times \cdots \times V_{k} \rightarrow V
$$

be a $k$-linear map. Show that $\beta$ is smooth with

$$
\mathrm{d} \beta\left(x_{1}, \ldots, x_{k}\right)\left(h_{1}, \ldots, h_{k}\right)=\sum_{j=1}^{k} \beta\left(x_{1}, \ldots, x_{j-1}, h_{j}, x_{j+1}, \ldots, x_{k}\right) .
$$

Exercise 8.2.12. Show that the space $M$ defined in Example 8.2.12 is not Hausdorff, but that the two maps $\varphi_{j}([j, x]):=x, j=1,2$, define a smooth atlas of $M$.

Exercise 8.2.13. A map $f: X \rightarrow Y$ between topological spaces is called a quotient map if a subset $O \subseteq Y$ is open if and only if $f^{-1}(O)$ is open. Show that:\\
(1) If $f_{1}: X_{1} \rightarrow Y_{1}$ and $f_{2}: X_{2} \rightarrow Y_{2}$ are open quotient maps, then the cartesian product

$$
f_{1} \times f_{2}: X_{1} \times X_{2} \rightarrow Y_{1} \times Y_{2}, \quad\left(x_{1}, x_{2}\right) \mapsto\left(f_{1}\left(x_{1}\right), f_{2}\left(x_{2}\right)\right)
$$

is a quotient map.\\
(2) If $f: X \rightarrow Y$ is a quotient map and we define on $X$ an equivalence relation by $x \sim y$ if $f(x)=f(y)$, then the map $\bar{f}: X / \sim \rightarrow Y$ is a homeomorphism if $X / \sim$ is endowed with the quotient topology.\\
(3) The map $q: \mathbb{R}^{n} \rightarrow \mathbb{T}^{n}, x \mapsto\left(e^{2 \pi i x_{j}}\right)_{j=1, \ldots, n}$ is a quotient map.\\
(4) The map $\bar{q}: \mathbb{R}^{n} / \mathbb{Z}^{n} \rightarrow \mathbb{T}^{n},[x] \mapsto\left(e^{2 \pi i x_{j}}\right)_{j=1, \ldots, n}$ is a homeomorphism.

Exercise 8.2.14. Let $M$ be a compact smooth manifold containing at least two points. Then each atlas of $M$ contains at least two charts. In particular, the atlas of $\mathbb{S}^{n}$ obtained from stereographic projections is minimal.

Exercise 8.2.15. Let $X$ and $Y$ be topological spaces and $q: X \rightarrow Y$ a quotient map, i.e., $q$ is surjective and $O \subseteq Y$ is open if and only if $q^{-1}(O)$ is open in $X$. Show that a map $f: Y \rightarrow Z$ ( $Z$ a topological space) is continuous if and only if the map $f \circ q: X \rightarrow Z$ is continuous.

Exercise 8.2.16. Let $M$ and $B$ be smooth manifolds. A smooth map $\pi$ : $M \rightarrow B$ is said to define a (locally trivial) fiber bundle with a typical fiber $F$ over the base manifold $B$ if each $b_{0} \in B$ has an open neighborhood $U$ for which there exists a diffeomorphism

$$
\varphi: \pi^{-1}(U) \rightarrow U \times F
$$

satisfying $\operatorname{pr}_{U} \circ \varphi=\pi$, where $\operatorname{pr}_{U}: U \times F \rightarrow U,(u, f) \mapsto u$ is the projection onto the first factor. Then the pair ( $\varphi, U$ ) is called a local trivialization.

Show that:\\
(1) If $(\varphi, U),(\psi, V)$ are local trivializations, then

$$
\varphi \circ \psi^{-1}(b, f)=\left(b, g_{\varphi \psi}(b)(f)\right)
$$

holds for a function $g_{\varphi \psi}: U \cap V \rightarrow \operatorname{Diff}(F)$.\\
(2) If ( $\gamma, W$ ) is another local trivialization, then

$$
g_{\varphi \varphi}=\operatorname{id}_{F} \quad \text { and } \quad g_{\varphi \psi} g_{\psi \gamma}=g_{\varphi \gamma} \quad \text { on } \quad U \cap V \cap W .
$$

Exercise 8.2.17. Show that a function $f: \mathbb{R} \rightarrow \mathbb{R}$ is a diffeomorphism if and only if either\\
(1) $f^{\prime}>0$ and $\lim _{x \rightarrow \pm \infty} f(x)= \pm \infty$.\\
(2) $f^{\prime}<0$ and $\lim _{x \rightarrow \pm \infty} f(x)=\mp \infty$.

Exercise 8.2.18. Let $B \in \mathrm{GL}_{n}(\mathbb{R})$ be an invertible matrix which is symmetric or skew-symmetric. Show that:\\
(1) $G:=\left\{g \in \mathrm{GL}_{n}(\mathbb{R}): g^{\top} B g=B\right\}$ is a subgroup of $\mathrm{GL}_{n}(\mathbb{R})$.\\
(2) If $B=B^{\top}$, then $B$ is a regular value of the smooth function

$$
f: M_{n}(\mathbb{R}) \rightarrow \operatorname{Sym}_{n}(\mathbb{R}), \quad x \mapsto x^{\top} B x .
$$

(3) If $B=-B^{\top}$, then $B$ is a regular value of the smooth function

$$
f: M_{n}(\mathbb{R}) \rightarrow \operatorname{Skew}_{n}(\mathbb{R}):=\left\{A \in M_{n}(\mathbb{R}): A^{\top}=-A\right\}, \quad x \mapsto x^{\top} B x
$$

(4) $G$ is a submanifold of $M_{n}(\mathbb{R})$.\\
(5) For $B=I_{p, q}:=\operatorname{diag}(\underbrace{1, \ldots, 1}_{p}, \underbrace{-1, \ldots,-1}_{q})$, the indefinite orthogonal group

$$
\mathrm{O}_{p, q}(\mathbb{R}):=\left\{g \in \mathrm{GL}_{n}(\mathbb{R}): g^{\top} I_{p, q} g=I_{p, q}\right\}
$$

is a submanifold on $M_{n}(\mathbb{R})$ of dimension $\frac{n(n-1)}{2}$.\\
(6) For $J=\left(\begin{array}{cc}\mathbf{0} & I \\ -I & \mathbf{0}\end{array}\right) \in M_{2}\left(M_{n}(\mathbb{R})\right) \cong M_{2 n}(\mathbb{R})$, the symplectic group

$$
\mathrm{Sp}_{2 n}(\mathbb{R}):=\left\{g \in \mathrm{GL}_{2 n}(\mathbb{R}): g^{\top} J g=J\right\}
$$

is a submanifold on $M_{2 n}(\mathbb{R})$ of dimension $n(2 n+1)$.

\subsection*{The Tangent Bundle}
The real strength of the theory of smooth manifolds is due to the fact that it permits us to analyze differentiable structures in terms of their derivatives. To model these derivatives appropriately, we introduce the tangent bundle $T M$ of a smooth manifold, tangent maps of smooth maps, and smooth vector fields.

We start with the definition of a tangent vector of a smooth manifold. The subtle point of this definition is that tangent vectors and the vector space structure can only be defined rather indirectly. The most straightforward way is to construct tangent vectors as "tangents" to smooth curves.

\subsubsection*{Tangent Vectors and Tangent Maps}
Definition 8.3.1. Let $M$ be a smooth manifold, $p \in M$, and $(\varphi, U)$ a chart of $M$ with $p \in U$. Let $\gamma: I \rightarrow M$ be a smooth curve, where $I \subseteq \mathbb{R}$ is an interval containing 0 and $\gamma(0)=p$. We call two such curves $\gamma_{i}: I_{i} \rightarrow M$, $i=1,2$, equivalent, denoted $\gamma_{1} \sim \gamma_{2}$, if

$$
\left(\varphi \circ \gamma_{1}\right)^{\prime}(0)=\left(\varphi \circ \gamma_{2}\right)^{\prime}(0)
$$

Clearly, this defines an equivalence relation. The equivalence classes are called tangent vectors at $p$. We write $T_{p}(M)$ for the set of all tangent vectors at $p$ and $[\gamma] \in T_{p}(M)$ for the equivalence class of the curve $\gamma$. The disjoint union

$$
T(M):=\coprod_{p \in M} T_{p}(M)
$$

is called the tangent bundle of $M$ and we write $\pi_{T M}: T M \rightarrow M$ for the projection, mapping $T_{p}(M)$ to $\{p\}$.

Remark 8.3.2. (a) The equivalence relation defining tangent vectors does not depend on the chart ( $\varphi, U$ ). If ( $\psi, V$ ) is a second chart with $p \in V$ and $\gamma: I \rightarrow M$ a smooth curve with $\gamma(0)=p$, then

$$
(\psi \circ \gamma)^{\prime}(0)=\mathrm{d}\left(\psi \circ \varphi^{-1}\right)(\varphi(p))(\varphi \circ \gamma)^{\prime}(0)
$$

so that we obtain the same equivalence relation on curves through $p$.\\
(b) If $U \subseteq \mathbb{R}^{n}$ is an open subset and $p \in U$, then each smooth curve $\gamma: I \rightarrow U$ with $\gamma(0)=p$ is equivalent to the curve $\eta_{v}(t):=p+t v$ for $v=\gamma^{\prime}(0)$. Hence each equivalence class contains exactly one curve $\eta_{v}$. We may therefore think of a tangent vector at $p \in U$ as a vector $v \in \mathbb{R}^{n}$ attached to the point $p$, and the map

$$
\mathbb{R}^{n} \rightarrow T_{p}(U), \quad v \mapsto\left[\eta_{v}\right]
$$

is a bijection. In this sense, we identify all tangent spaces $T_{p}(U)$ with $\mathbb{R}^{n}$, so that we obtain a bijection

$$
T(U) \cong U \times \mathbb{R}^{n}
$$

As an open subset of the product space $T\left(\mathbb{R}^{n}\right) \cong \mathbb{R}^{2 n}$, the tangent bundle $T(U)$ inherits a natural manifold structure.\\
(c) For each $p \in M$ and any chart ( $\varphi, U$ ) with $p \in U$, the map

$$
T_{p}(\varphi): T_{p}(M) \rightarrow \mathbb{R}^{n}, \quad[\gamma] \mapsto(\varphi \circ \gamma)^{\prime}(0)
$$

is well-defined by the definition of the equivalence relation. Moreover, the curve

$$
\gamma(t):=\varphi^{-1}(\varphi(p)+t v),
$$

which is smooth and defined on some neighborhood of 0 , satisfies $(\varphi \circ \gamma)^{\prime}(0)=v$. Hence $T_{p}(\varphi)$ is a bijection.

Definition 8.3.3. (a) Each tangent space $T_{p}(M)$ carries the unique structure of an $n$-dimensional vector space with the property that for each chart ( $\varphi, U$ ) of $M$ with $p \in U$, the map

$$
T_{p}(\varphi): T_{p}(M) \rightarrow \mathbb{R}^{n}, \quad[\gamma] \mapsto(\varphi \circ \gamma)^{\prime}(0)
$$

is a linear isomorphism.\\
In fact, since $T_{p}(\varphi)$ is a bijection, we may define a vector space structure on $T_{p}(M)$ by

$$
v+w:=T_{p}(\varphi)^{-1}\left(T_{p}(\varphi) v+T_{p}(\varphi) w\right) \quad \text { and } \quad \lambda v:=T_{p}(\varphi)^{-1}\left(\lambda T_{p}(\varphi) v\right)
$$

for $\lambda \in \mathbb{R}, v, w \in T_{p}(M)$. For any other chart ( $\psi, V$ ) with $p \in V$, we then have

$$
T_{p}(\psi)=\mathrm{d}\left(\psi \circ \varphi^{-1}\right)(\varphi(p)) \circ T_{p}(\varphi)
$$

and since $\mathrm{d}\left(\psi \circ \varphi^{-1}\right)(\varphi(p))$ is a linear automorphism of $\mathbb{R}^{n}$, the vector space structure on $T_{p}(M)$ does not depend on the chart we use for its definition.\\
(b) If $f: M \rightarrow N$ is a smooth map and $p \in M$, then we obtain a linear map

$$
T_{p}(f): T_{p}(M) \rightarrow T_{f(p)}(N), \quad[\gamma] \mapsto[f \circ \gamma]
$$

In fact, we only have to observe that for any chart $(\varphi, U)$ of $N$ with $f(p) \in U$ and any chart ( $\psi, V$ ) of $M$ with $p \in V$, we have

$$
\begin{aligned}
T_{f(p)}(\varphi)[f \circ \gamma] & =(\varphi \circ f \circ \gamma)^{\prime}(0)=\mathrm{d}\left(\varphi \circ f \circ \psi^{-1}\right)(\psi(p))(\psi \circ \gamma)^{\prime}(0) \\
& =\mathrm{d}\left(\varphi \circ f \circ \psi^{-1}\right)(\psi(p)) T_{p}(\psi)[\gamma] .
\end{aligned}
$$

This relation shows that $T_{p}(f)$ is well-defined, and a linear map.\\
The collection of all these maps defines a map

$$
T(f): T(M) \rightarrow T(N) \quad \text { with } \quad T_{p}(f)=\left.T(f)\right|_{T_{p}(M)}, p \in M
$$

It is called the tangent map of $f$.\\
(c) If $M \subseteq \mathbb{R}$ is an open subset, then $f: M \rightarrow N$ is a smooth curve in $N$, and its tangent vector is $f^{\prime}(t):=T_{t}(f)(1)$, where $1 \in T_{t}(\mathbb{R}) \cong \mathbb{R}$ is considered as a tangent vector.\\
(d) If $N$ is a vector space, then we identify $T(N)$ in a natural way with $N \times N$. Accordingly, we have

$$
T_{p}(f)(v)=(f(p), \mathrm{d} f(p) v)
$$

for a map $\mathrm{d} f: T(M) \rightarrow N$ with $\mathrm{d} f(p):=\left.\mathrm{d} f\right|_{T_{p}(M)}$.\\
Remark 8.3.4. (a) For an open subset $U \subseteq \mathbb{R}^{n}$ and $p \in U$, the vector space structure on $T_{p}(U)=\{p\} \times \mathbb{R}^{n}$ is simply given by

$$
(p, v)+(p, w):=(p, v+w) \quad \text { and } \quad \lambda(p, v):=(p, \lambda v)
$$

for $v, w \in \mathbb{R}^{n}$ and $\lambda \in \mathbb{R}$.\\
(b) If $f: U \rightarrow V$ is a smooth map between open subsets $U \subseteq \mathbb{R}^{n}$ and $V \subseteq \mathbb{R}^{m}, p \in U$, and $\eta_{v}(t)=p+t v$, then the tangent map satisfies

$$
T(f)(p, v)=\left[f \circ \eta_{v}\right]=\left(f \circ \eta_{v}\right)^{\prime}(0)=\left(f(p), \mathrm{d} f(p) \eta_{v}^{\prime}(0)\right)=(f(p), \mathrm{d} f(p) v)
$$

The main difference to the map $\mathrm{d} f$ is the bookkeeping; here we keep track of what happens to the point $p$ and the tangent vector $v$. We may also write

$$
T(f)=\left(f \circ \pi_{T U}, \mathrm{~d} f\right): T U \cong U \times \mathbb{R}^{n} \rightarrow T V \cong V \times \mathbb{R}^{n}
$$

where $\pi_{T U}: T U \rightarrow U,(p, v) \mapsto p$, is the projection map.\\
(c) If ( $\varphi, U$ ) is a chart of $M$ and $p \in U$, then we identify $T(\varphi(U))$ with $\varphi(U) \times \mathbb{R}^{n}$ and obtain for $[\gamma] \in T_{p}(M)$ :

$$
T(\varphi)([\gamma])=(\varphi(p),[\varphi \circ \gamma])=\left(\varphi(p),(\varphi \circ \gamma)^{\prime}(0)\right)
$$

which is consistent with our previously introduced notation $T_{p}(\varphi)$ (Remark 8.3.2).

Lemma 8.3.5 (Chain Rule for Tangent Maps). For smooth maps $f: M \rightarrow N$ and $g: N \rightarrow L$, the tangent maps satisfy

$$
T(g \circ f)=T(g) \circ T(f)
$$

Proof. We recall from Lemma 8.2.21 that $g \circ f: M \rightarrow L$ is a smooth map, so that $T(g \circ f)$ is defined. For $p \in M$ and $[\gamma] \in T_{p}(M)$, we further have

$$
T_{p}(g \circ f)[\gamma]=[g \circ f \circ \gamma]=T_{f(p)}(g)[f \circ \gamma]=T_{f(p)}(g) T_{p}(f)[\gamma]
$$

Since $p$ was arbitrary, this implies the lemma.\\
So far we only considered the tangent bundle $T(M)$ of a smooth manifold $M$ as a set, but this set also carries a natural topology and a smooth manifold structure.

Definition 8.3.6 (Manifold structure on $T(M)$ ). Let $M$ be a smooth manifold. First, we introduce a topology on $T(M)$.

For each chart ( $\varphi, U$ ) of $M$, we have a tangent map

$$
T(\varphi): T(U) \rightarrow T(\varphi(U)) \cong \varphi(U) \times \mathbb{R}^{n}
$$

where we consider $T(U)=\bigcup_{p \in U} T_{p}(M)$ as a subset of $T(M)$. We define a topology on $T(M)$ by declaring a subset $O \subseteq T(M)$ to be open if for each chart ( $\varphi, U$ ) of $M$, the set $T(\varphi)(O \cap T(U))$ is an open subset of $T(\varphi(U))$. It is easy to see that this defines indeed a Hausdorff topology on $T(M)$ for which all the subsets $T(U)$ are open and the maps $T(\varphi)$ are homeomorphisms onto open subsets of $\mathbb{R}^{2 n}$ (Exercise 8.3.1; see also Remark 8.2.13).

Since for two charts $(\varphi, U),(\psi, V)$ of $M$, the map

$$
T\left(\varphi \circ \psi^{-1}\right)=T(\varphi) \circ T(\psi)^{-1}: T(\psi(V)) \rightarrow T(\varphi(U))
$$

is smooth, for each atlas $\mathcal{A}$ of $M$, the collection $(T(\varphi), T(U))_{(\varphi, U) \in \mathcal{A}}$ is a smooth atlas of $T(M)$. We thus obtain on $T(M)$ the structure of a smooth manifold.

Lemma 8.3.7. If $f: M \rightarrow N$ is a smooth map, then its tangent map $T(f)$ is smooth.

Proof. Let $p \in M$ and choose charts ( $\varphi, U$ ) and ( $\psi, V$ ) of $M$, resp., $N$, with $p \in U$ and $f(p) \in V$. Then the map

$$
T(\psi) \circ T(f) \circ T(\varphi)^{-1}=T\left(\psi \circ f \circ \varphi^{-1}\right): T\left(\varphi\left(f^{-1}(V) \cap U\right)\right) \rightarrow T(V)
$$

is smooth, and this implies that $T(f)$ is a smooth map.\\
Remark 8.3.8. For smooth manifolds $M_{1}, \ldots, M_{n}$, the projection maps

$$
\pi_{i}: M_{1} \times \cdots \times M_{n} \rightarrow M_{i}, \quad\left(p_{1}, \ldots, p_{n}\right) \mapsto p_{i}
$$

induce a diffeomorphism

$$
\left(T\left(\pi_{1}\right), \ldots, T\left(\pi_{n}\right)\right): T\left(M_{1} \times \cdots \times M_{n}\right) \rightarrow T M_{1} \times \cdots \times T M_{n}
$$

(Exercise 8.3.2).\\
Definition 8.3.9. Let $M$ be a differentiable manifold and $V$ a finite-dimensional vector space. A $V$-valued Pfaffian form, also simply called a 1 -form, on $M$ is a smooth map $\omega: T M \rightarrow V$ whose restrictions $\omega_{p}: T_{p}(M) \rightarrow V$ are linear for each $p \in M$. The set of all $V$-valued 1-forms on $M$ is denoted by $\Omega^{1}(M, V)$, and for $V=\mathbb{R}$ we abbreviate $\Omega^{1}(M):=\Omega^{1}(M, \mathbb{R})$.

Recall the bundle projection $\pi_{T M}: T M \rightarrow M$. Then for any smooth function $f: M \rightarrow \mathbb{R}$, the pointwise product

$$
f \cdot \omega:=\left(f \circ \pi_{T M}\right) \cdot \omega: T M \rightarrow V
$$

also is a $V$-valued 1-form. Combined with the obvious vector space structure on $\Omega^{1}(M, V)$, defined by pointwise addition and scalar multiplication, we thus obtain on $\Omega^{1}(M, V)$ the structure of a $C^{\infty}(M)$-module.

Example 8.3.10. If $f: M \rightarrow V$ is a smooth function, then $T f: T M \rightarrow T V$ is also smooth. Identifying $T V$ in the canonical fashion with $V \times V$ and writing $\operatorname{pr}_{2}: T V \rightarrow V,(x, v) \mapsto v$ for the projection, the map

$$
\mathrm{d} f:=\mathrm{pr}_{2} \circ T f: T M \rightarrow V
$$

is a $V$-valued 1-form, and the tangent map $T f$ can now be written $T f(v)= (f(p), \mathrm{d} f(p) v)$ for $v \in T_{p}(M)$.

Remark 8.3.11. Now let $\omega \in \Omega^{1}(M)$ and $X \in \mathcal{V}(M)$. For every $p \in M$, we can apply $\omega_{p}$ to $X(p)$. We thus get a smooth function $\omega(X)=\omega \circ X: M \rightarrow \mathbb{R}$.

Definition 8.3.12. Let $\omega \in \Omega^{1}(M, V)$ be a $V$-valued 1 -form and $f: N \rightarrow M$ a smooth map. Then we obtain a $V$-valued 1 -form

$$
f^{*} \omega:=\omega \circ T f: T N \rightarrow V
$$

called the pull-back of $\omega$ to $N$.

Remark 8.3.13. We clearly have the following rules for pull-backs of 1 -forms. For $\omega \in \Omega^{1}(M, V)$ and smooth maps $\psi: W \rightarrow N$ and $\varphi: N \rightarrow M$, we have

$$
\operatorname{id}_{M}^{*} \omega=\omega \quad \text { and } \quad(\varphi \circ \psi)^{*} \omega=\psi^{*}\left(\varphi^{*} \omega\right)
$$

Moreover,

$$
\varphi^{*}: \Omega^{1}(M, V) \rightarrow \Omega^{1}(N, V)
$$

is a linear map satisfying $\varphi^{*}(f \cdot \omega)=(f \circ \varphi) \cdot \varphi^{*} \omega$ for $f \in C^{\infty}(M)$.\\
Definition 8.3.14. Let $f: M_{1} \rightarrow M_{2}$ be a smooth map and $m \in M_{1}$. The map $f$ is called submersive at $m$ if the differential $T_{m}(f)$ is surjective. Otherwise $m$ is called a critical point of $f$.

The map $f$ is said to be a submersion if $T_{m}(f)$ is surjective for each $m \in M_{1}$.

Lemma 8.3.15. If $f: M_{1} \rightarrow M_{2}$ is a smooth map which is submersive at $m \in M_{1}$, then there exists an open neighborhood $U \subseteq M_{2}$ of $p:=f(m)$ and a smooth map $\sigma: U \rightarrow M_{1}$ with $f \circ \sigma=\operatorname{id}_{U}$ and $\sigma(p)=m$.

Proof. Since the assertion is purely local, we may w.l.o.g. assume that $M_{1}$ is an open subset of $\mathbb{R}^{n}$ and $M_{2}$ is an open subset of $\mathbb{R}^{k}$. That $f$ is submersive at $m$ means that $\mathrm{d} f(m): \mathbb{R}^{n} \rightarrow \mathbb{R}^{k}$ is surjective. Let $V \subseteq \mathbb{R}^{n}$ be a linear subspace for which $\left.\mathrm{d} f(m)\right|_{V}: V \rightarrow \mathbb{R}^{k}$ is a linear isomorphism. Then $F:=\left.f\right|_{M_{1} \cap V}: M_{1} \cap V \rightarrow M_{2}$ is a smooth map whose differential $T_{m}(F)$ is invertible. Hence the Inverse Function Theorem implies the existence of a smooth inverse $\sigma$, defined on an open neighborhood of $p=F(m)$.

Proposition 8.3.16 (Universal property of submersions). Suppose that $f: M_{1} \rightarrow M_{2}$ is a surjective submersion and $N$ a smooth manifold. Then a map $h: M_{2} \rightarrow N$ is smooth if and only if the map $h \circ f: M_{1} \rightarrow N$ is smooth. In particular, for each smooth map $g: M_{1} \rightarrow N$ which is constant on all fibers of $f$, there exists a unique smooth map $h: M_{2} \rightarrow N$ with $h \circ f=g$.

Proof. If $h$ is smooth, then also the composition $h \circ f$ is smooth. Assume, conversely, that $h \circ f$ is smooth. To see that $h$ is smooth, pick $p \in M_{2}$ and an open neighborhood $U \subseteq M_{2}$ of $p$ on which there exists a smooth section $\sigma: U \rightarrow M_{1}$ with $f \circ \sigma=\operatorname{id}_{U}$ (Lemma 8.3.15). Then $\left.h\right|_{U}=h \circ f \circ \sigma$ is smooth. Hence $h$ is smooth on a neighborhood of $p$, and since $p \in M_{2}$ was arbitrary, $h$ is smooth.

The second assertion immediately follows from the first one if we define the map $h: M_{2} \rightarrow N$ by $h(f(x)):=g(x)$, which works if and only if $g$ is constant on the fibers of $f$.

Corollary 8.3.17. If $f: M_{1} \rightarrow M_{2}$ is a bijective submersion, then $f$ is a diffeomorphism.

Proof. Apply the preceding proposition with $N:=M_{1}$ and $g=\operatorname{id}_{M_{1}}$.

Remark 8.3.18. The smooth map $f: M_{1}:=\mathbb{R} \rightarrow M_{2}:=\mathbb{R}, x \mapsto x^{3}$ is submersive at all points $x \neq 0$. The map $g=\operatorname{id}_{\mathbb{R}}: M_{1}=\mathbb{R} \rightarrow N:=\mathbb{R}$ is smooth and bijective, hence constant on the fibers of $f$, but the map $\bar{g}: M_{2} \rightarrow N, x \mapsto x^{\frac{1}{3}}$ is not smooth at 0 . This shows that the assumption in Proposition 8.3.16 that $f$ is submersive is really needed.

\subsubsection*{Algebraic Tangent Spaces}
An alternative algebraic approach to the tangent space is to consider not the velocity vectors of curves, but to study their actions on smooth functions: Let $v \in T_{p}(M)$ and $U$ be an open neighborhood of $p$ in $M$. For a smooth function $f: U \rightarrow \mathbb{R}$, we then put (cf. Definition 8.3.3(d))

$$
v(f):=\mathrm{d} f(p) v
$$

From the product rule (see Exercise 8.3.7), one immediately derives

$$
v(f \cdot g)=v(f) g(p)+f(p) v(g)
$$

We also observe that $v(f)=v(g)$ if $f$ and $g$ coincide on a neighborhood of $p$. This motivates the following definition:

Definition 8.3.19. Let $M$ be a differentiable manifold. For $p \in M$, let

$$
C^{\infty}(M, p)=\coprod_{p \in U} C^{\infty}(U)
$$

be the set of smooth functions which are defined on an open subset $U$ containing $p$. On $C^{\infty}(M, p)$, we consider the equivalence relation defined by $f \sim g$ if there exists an open neighborhood $U$ of $p$ in $M$ such that $\left.f\right|_{U}=\left.g\right|_{U}$. We denote the equivalence class of $f \in C^{\infty}(M, p)$ with respect to $\sim$ by $f_{p}$. It is called the germ of $f$ at $p$. The set $C^{\infty}(M)_{p}:=C^{\infty}(M, p) / \sim$ of equivalence classes inherits an algebra structure (given by pointwise addition and multiplication of functions)

$$
f_{p}+g_{p}:=(f+g)_{p} \quad \text { and } \quad f_{p} g_{p}:=(f g)_{p} .
$$

A linear functional $v: C^{\infty}(M)_{p} \rightarrow \mathbb{R}$ is called a derivation at $p$ if

$$
v\left(f_{p} \cdot g_{p}\right)=v\left(f_{p}\right) \cdot g(p)+f(p) \cdot v\left(g_{p}\right)
$$

for $f_{p}, g_{p} \in C^{\infty}(M)_{p}$. The algebraic tangent space $T_{p}^{\mathrm{alg}}(M)$ of $M$ at $p$ is defined to be the set of all derivations $v: C^{\infty}(M)_{p} \rightarrow \mathbb{R}$.

Note that we have assigned a derivation to every $v \in T_{p}(M)$, but we have not assigned a tangent vector to every derivation. This is nevertheless possible, and will be done in Proposition 8.3.24.

Proposition 8.3.20. Let $M$ be a differentiable manifold. Then for $p \in M$ the following is true:\\
(i) $T_{p}^{\mathrm{alg}}(M)$ is a vector space with respect to the pointwise operations.\\
(ii) If $f_{p} \in C^{\infty}(M)_{p}$ is such that $f$ is constant on a neighborhood of $p$, then $v\left(f_{p}\right)=0$ for all $v \in T_{p}^{\mathrm{alg}}(M)$.\\
Proof. The first statement is obvious, and for the second, it suffices to prove that $v\left(1_{p}\right)=0$ for all $v \in T_{p}^{\mathrm{alg}}(M)$, where 1 is the constant function with value 1 . For this, we compute:

$$
v\left(1_{p}\right)=v\left(1_{p} \cdot 1_{p}\right)=v\left(1_{p}\right) \cdot 1+v\left(1_{p}\right) \cdot 1=2 v\left(1_{p}\right)
$$

which implies $v\left(1_{p}\right)=0$.\\
Note that at this point it is not clear what the size of the tangent space is. We cannot even say whether it is finite-dimensional. In order to clarify the structure of $T_{p}^{\mathrm{alg}}(M)$, we describe it in local coordinates.\\
Definition 8.3.21. Let $\varphi: U \rightarrow V \subset \mathbb{R}^{n}$ be a chart with $p \in U$. We define $n$ elements $\left.\frac{\partial}{\partial x_{j}}\right|_{p} \in T_{p}^{\mathrm{alg}}(M)$ for $j=1, \ldots, n$ by

$$
\left.\frac{\partial}{\partial x_{j}}\right|_{p}\left(f_{p}\right):=\frac{\partial\left(f \circ \varphi^{-1}\right)}{\partial x_{j}}(\varphi(p)),
$$

where the points in $\mathbb{R}^{n}$ are denoted by $x=\left(x_{1}, \ldots, x_{n}\right)$. Note that this definition depends on the choice of the chart.

It turns out that the $\left.\frac{\partial}{\partial x_{j}}\right|_{p}$ form a basis for $T_{p} M$, which we call the $\varphi$-basis for $T_{p}^{\mathrm{alg}}(M)$. To prove this, we start with a lemma:\\
Lemma 8.3.22 (Hadamard). Let $B$ be an open ball in $\mathbb{R}^{n}$ with center a and $f \in C^{\infty}(B)$. Then there exist smooth functions $g_{1}, \ldots, g_{n} \in C^{\infty}(B)$ such that\\
(i) $f(x)=f(a)+\sum_{j=1}^{n}\left(x_{j}-a_{j}\right) g_{j}(x)$ for all $x \in B$.\\
(ii) $g_{j}(a)=\frac{\partial f}{\partial x_{j}}(a)$.

Proof. For fixed $x \in B$, we consider the function

$$
\xi:[-1,1] \rightarrow \mathbb{R}, \quad \xi(t):=f(a+t(x-a))
$$

Then, for $g_{j}(x):=\int_{0}^{1} \frac{\partial f}{\partial x_{j}}(a+t(x-a)) d t$, we have

$$
\begin{aligned}
f(x) & =\xi(0)+\int_{0}^{1} \xi^{\prime}(t) d t=\xi(0)+\int_{0}^{1} \sum_{j=1}^{n} \frac{\partial f}{\partial x_{j}}(a+t(x-a)) \cdot\left(x_{j}-a_{j}\right) d t \\
& =f(a)+\sum_{j=1}^{n}\left(x_{j}-a_{j}\right)\left(\int_{0}^{1} \frac{\partial f}{\partial x_{j}}(a+t(x-a)) d t\right) \\
& =f(a)+\sum_{j=1}^{n}\left(x_{j}-a_{j}\right) g_{j}(x)
\end{aligned}
$$

Since $f \in C^{\infty}(B)$ by assumption, we also have $g_{j} \in C^{\infty}(B)$ (Exercise).

Proposition 8.3.23. Let $M$ be a smooth manifold and $(\varphi, U)$ be a chart on $M$ with coordinate functions $x_{j}: U \rightarrow \mathbb{R}$ given by $\varphi(q)=\left(x_{1}(q), \ldots, x_{n}(q)\right)$, see Remark 8.2.17. Then the tangent vectors $\left.\frac{\partial}{\partial x_{j}}\right|_{p}, j=1, \ldots, n$, form a basis for $T_{p}^{\mathrm{alg}}(M)$ and for every $v \in T_{p}^{\mathrm{alg}}(M)$, the formula holds

$$
v=\left.\sum_{j=1}^{n} v\left(\left(x_{j}\right)_{p}\right) \cdot \frac{\partial}{\partial x_{j}}\right|_{p}
$$

Proof. Fix $v \in T_{p}^{\mathrm{alg}}(M)$ and $f \in C^{\infty}(\widetilde{V})$, where $\widetilde{V} \subseteq U$ is an open subset containing $p$. Let $V=\varphi(\widetilde{V})$ and put $F=f \circ \varphi^{-1}: V \rightarrow \mathbb{R}$. Shrinking the neighborhoods, we may assume that $V$ is an open ball with center $\varphi(p)$. We apply Lemma 8.3.22 to $F$ and get functions $g_{j}: V \rightarrow \mathbb{R}$ satisfying $F(x)= F(a)+\sum_{j=1}^{n}\left(x_{j}-a_{j}\right) g_{j}(x)$. Then, using Proposition 8.3.20, we compute

$$
\begin{aligned}
v\left(f_{p}\right) & =v\left((F \circ \varphi)_{p}\right)=v\left(F(\varphi(p))+\sum_{j=1}^{n}\left(x_{j}-x_{j}(p)\right)_{p}\left(g_{j} \circ \varphi\right)_{p}\right) \\
& =v\left(f(p) \cdot 1_{p}\right)+v\left(\sum_{j=1}^{n}\left(x_{j}-x_{j}(p) \cdot 1\right)_{p} \cdot\left(g_{j} \circ \varphi\right)_{p}\right) \\
& =0+\sum_{j=1}^{n}\left(v\left(\left(x_{j}\right)_{p}\right) \cdot g_{j}(\varphi(p))+v\left(\left(g_{j} \circ \varphi\right)_{p}\right) \cdot\left(x_{j}(p)-x_{j}(p) \cdot 1\right)\right) \\
& =\sum_{j=1}^{n} v\left(\left(x_{j}\right)_{p}\right) \cdot \frac{\partial\left(f \circ \varphi^{-1}\right)}{\partial x_{j}}(\varphi(p))=\left.\sum_{j=1}^{n} v\left(\left(x_{j}\right)_{p}\right) \cdot \frac{\partial}{\partial x_{j}}\right|_{p}\left(f_{p}\right) .
\end{aligned}
$$

Hence, every derivation $v \in T_{p}^{\mathrm{alg}}(M)$ is a linear combination of the $\left.\frac{\partial}{\partial x_{j}}\right|_{p}, j= 1, \ldots, n$. It remains to show that the $\left.\frac{\partial}{\partial x_{j}}\right|_{p}$ for $j=1, \ldots, n$ are linearly independent. But this is immediate from $\left.\frac{\partial}{\partial x_{j}}\right|_{p}\left(\left(x_{j}\right)_{p}\right)=\delta_{i j}$.

Proposition 8.3.24. Let $M$ be a smooth manifold and $p \in M$. Then

$$
\Gamma: T_{p}(M) \rightarrow T_{p}^{\mathrm{alg}}(M), \quad \Gamma(v)\left(f_{p}\right):=v\left(f_{p}\right)=\mathrm{d} f(p) v
$$

is a linear isomorphism.\\
Proof. It follows from (8.3) and the subsequent calculation that $\Gamma$ is a linear map into the algebraic tangent space.

Consider a chart $(\varphi, U)$ with $p \in U$. In view of Proposition 8.3.23, the condition $\Gamma(v)=0$ is equivalent to $0=v\left(\left(x_{j}\right)_{p}\right)=\mathrm{d} x_{j}(p) v$ for all coordinate functions, which in turn is equivalent to $T_{p}(\varphi) v=0$, i.e., to $v=0$. Thus $\Gamma$ is injective. Since both spaces are of the same dimension, $\Gamma$ is a linear bijection.

In view of Proposition 8.3.24, we will no longer distinguish between the geometric and the algebraic tangent spaces.

Recall that with any differential map $f: M \rightarrow N$ between two manifolds we associate a linear map $T_{p}(f): T_{p} M \rightarrow T_{f(p)} N$ as derivative.\\
\includegraphics[max width=\textwidth, center]{2025_10_18_3eaf04e77f3cb364fce5g-270}\\
\includegraphics[max width=\textwidth, center]{2025_10_18_3eaf04e77f3cb364fce5g-270(1)}

Remark 8.3.25. For charts ( $\varphi, U$ ) of $M$ and ( $\psi, V$ ) of $N$, we obtained distinguished bases for $T_{p} M$ and $T_{f(p)} N$. We want to describe the derivative of $f$ with respect to these bases.\\
\includegraphics[max width=\textwidth, center]{2025_10_18_3eaf04e77f3cb364fce5g-270(2)}

Let $F:=\psi \circ f \circ \varphi^{-1}:: \varphi\left(U \cap f^{-1}(V)\right) \rightarrow \psi(V)$ and write $F_{i}$ for its components.\\
\includegraphics[max width=\textwidth, center]{2025_10_18_3eaf04e77f3cb364fce5g-270(3)}

Then we find (exercise)

$$
T_{p}(f)\left(\left.\frac{\partial}{\partial x_{j}}\right|_{p}\right)=\left.\sum_{i=1}^{\operatorname{dim} N} \frac{\partial F_{i}(\varphi(p))}{\partial x_{j}} \frac{\partial}{\partial y_{i}}\right|_{f(p)}
$$

Hence the Jacobian of $F$ at $\varphi(p)$ is exactly the matrix of $T_{p}(f)$ with respect to the bases of $T_{p} M$ and $T_{f(p)} N$, associated with $\varphi$ and $\psi$, respectively.

\subsubsection*{Exercises for Section 8.3}
Exercise 8.3.1. Let $M$ be a smooth manifold. We call a subset $O \subseteq T(M)$ open if for each chart ( $\varphi, U$ ) of $M$, the set $T(\varphi)(O \cap T(U))$ is an open subset of $T(\varphi(U))$. Show that:\\
(1) This defines a topology on $T(M)$.\\
(2) All subsets $T(U)$ are open (Remark 8.3.4(b)).\\
(3) The maps $T(\varphi): T U \rightarrow T(\varphi(U)) \cong \varphi(U) \times \mathbb{R}^{n}$ are homeomorphisms onto open subsets of $\mathbb{R}^{2 n} \cong T\left(\mathbb{R}^{n}\right)$.\\
(4) The projection $\pi_{T M}: T(M) \rightarrow M$ is continuous.\\
(5) $T(M)$ is Hausdorff.

Exercise 8.3.2. For smooth manifolds $M_{1}, \ldots, M_{n}$, the projection maps

$$
\pi_{i}: M_{1} \times \cdots \times M_{n} \rightarrow M_{i}, \quad\left(p_{1}, \ldots, p_{n}\right) \mapsto p_{i}
$$

induce a diffeomorphism

$$
\left(T\left(\pi_{1}\right), \ldots, T\left(\pi_{n}\right)\right): T\left(M_{1} \times \cdots \times M_{n}\right) \rightarrow T M_{1} \times \cdots \times T M_{n}
$$

Exercise 8.3.3. Let $N$ and $M_{1}, \ldots, M_{n}$ be a smooth manifolds. Show that a map

$$
f: N \rightarrow M_{1} \times \cdots \times M_{n}
$$

is smooth if and only if all its component functions $f_{i}: N \rightarrow M_{i}$ are smooth.\\
Exercise 8.3.4. Let $f: M \rightarrow N$ be a smooth map between manifolds, $\pi_{T M}: T M \rightarrow M$ the tangent bundle projection and $\sigma_{M}: M \rightarrow T M$ the zero section. Show that for each smooth map $f: M \rightarrow N$ we have

$$
\pi_{T N} \circ T f=f \circ \pi_{T M} \quad \text { and } \quad \sigma_{N} \circ f=T f \circ \sigma_{M} .
$$

Exercise 8.3.5 (Inverse Function Theorem for manifolds). Let $f$ : $M \rightarrow N$ be a smooth map and $p \in M$ such that $T_{p}(f): T_{p}(M) \rightarrow T_{p}(N)$ is a linear isomorphism. Show that there exists an open neighborhood $U$ of $p$ in $M$ such that the restriction $\left.f\right|_{U}: U \rightarrow f(U)$ is a diffeomorphism onto an open subset of $N$.

Exercise 8.3.6. Let $\mathbb{P}_{n}(\mathbb{R})$ be the $n$-dimensional projective space over $\mathbb{R}$. For a point $x \in \mathbb{P}_{n}(\mathbb{R})$, let $L_{x}$ be the line in $\mathbb{R}^{n+1}$ represented by $x$. Show that $L:=\coprod_{x \in M} L_{x}$ carries the structure of a line bundle, i.e., a vector bundle of rank 1.

Exercise 8.3.7. Let $\mu: E \times F \rightarrow W$ be a bilinear map and $M$ a smooth manifold. For $f \in C^{\infty}(M, E), g \in C^{\infty}(M, F)$ and $p \in M$ set $h(p):=\mu(f(p), g(p))$. Show that $h$ is smooth with

$$
T(h) v=\mu(T(f) v, g(p))+\mu(f(p), T(g) v) \quad \text { for } v \in T_{p}(M)
$$

\subsection*{Vector Fields}
Vector fields are maps which associate with each point in a manifold a tangent vector at this point. They can be interpreted as a geometric way to formulate first order differential equations on a manifold, a point of view we will elaborate on in Section 8.5. First, we introduce the Lie algebra structure on the space of smooth vector fields.

\subsubsection*{The Lie Algebra of Vector Fields}
Definition 8.4.1. (a) Let $\pi_{T M}: T M \rightarrow M$ denote the canonical projection mapping $T_{p}(M)$ to $p$. A (smooth) vector field $X$ on $M$ is a smooth section of the tangent bundle $T M$, i.e., a smooth map $X: M \rightarrow T M$ with $\pi_{T M} \circ X= \operatorname{id}_{M}$. We write $\mathcal{V}(M)$ for the space of all vector fields on $M$.\\
(b) If $U \subseteq M$ is an open subset and $f \in C^{\infty}(U, V)$ is a smooth function on $U$ with values in some finite-dimensional vector space $V$ and $X \in \mathcal{V}(M)$, then we obtain a smooth function on $U$ via

$$
\mathcal{L}_{X} f:=\left.\mathrm{d} f \circ X\right|_{U}: U \rightarrow T U \rightarrow V
$$

We thus obtain for each $X \in \mathcal{V}(M)$ a linear operator $\mathcal{L}_{X}$ on $C^{\infty}(U, V)$. The function $\mathcal{L}_{X} f$ is also called the Lie derivative of $f$ with respect to $X$.

Remark 8.4.2. (a) If $U$ is an open subset of $\mathbb{R}^{n}$, then $T U=U \times \mathbb{R}^{n}$ with the bundle projection

$$
\pi_{T U}: U \times \mathbb{R}^{n} \rightarrow U, \quad(x, v) \mapsto x
$$

Therefore, each smooth vector field is of the form $X(x)=(x, \widetilde{X}(x))$ for some smooth function $\widetilde{X}: U \rightarrow \mathbb{R}^{n}$, and we may thus identify $\mathcal{V}(U)$ with the space $C^{\infty}\left(U, \mathbb{R}^{n}\right)$ of smooth $\mathbb{R}^{n}$-valued functions on $U$.\\
(b) The space $\mathcal{V}(M)$ carries a natural vector space structure given by

$$
(X+Y)(p):=X(p)+Y(p), \quad(\lambda X)(p):=\lambda X(p)
$$

(Exercise 8.4.1).\\
More generally, we can multiply vector fields with smooth functions

$$
(f X)(p):=f(p) X(p), \quad f \in C^{\infty}(M, \mathbb{R}), X \in \mathcal{V}(M)
$$

Before we turn to the Lie bracket on the space $\mathcal{V}(M)$ of smooth vector fields on a manifold $M$, we take a closer look at the local level.

Lemma 8.4.3. Let $U \subseteq \mathbb{R}^{n}$ be an open subset. Then we obtain a Lie bracket on the space $C^{\infty}\left(U, \mathbb{R}^{n}\right)$ by

$$
[X, Y](p):=\mathrm{d} Y(p) X(p)-\mathrm{d} X(p) Y(p) \quad \text { for } \quad p \in U
$$

With respect to this Lie bracket, the map

$$
C^{\infty}\left(U, \mathbb{R}^{n}\right) \rightarrow \operatorname{End}\left(C^{\infty}(U, \mathbb{R})\right), \quad X \mapsto \mathcal{L}_{X}, \quad \mathcal{L}_{X}(f)(p):=\mathrm{d} f(p) X(p)
$$

is an injective homomorphism of Lie algebras, i.e., $\mathcal{L}_{[X, Y]}=\left[\mathcal{L}_{X}, \mathcal{L}_{Y}\right]$.\\
Proof. (L1) and (L2) are obvious from the definition. To verify the Jacobi identity, we first observe that the map $X \mapsto \mathcal{L}_{X}$ is injective. In fact, if $\mathcal{L}_{X}=0$, then we have for each linear function $f: \mathbb{R}^{n} \rightarrow \mathbb{R}$ the relation $0=\left(\mathcal{L}_{X} f\right)(p)= \mathrm{d} f(p) X(p)=f(X(p))$, and therefore $X(p)=0$.

Next we observe that

$$
\begin{aligned}
\mathcal{L}_{X} \mathcal{L}_{Y}(f)(p) & =\mathrm{d}\left(\mathcal{L}_{Y} f\right)(p) X(p)=\mathrm{d}(\mathrm{~d} f \circ Y)(p) X(p) \\
& =\left(\mathrm{d}^{2} f\right)(p)(X(p), Y(p))+\mathrm{d} f(p) \mathrm{d} Y(p) X(p)
\end{aligned}
$$

so that the Schwarz Lemma implies $\mathcal{L}_{X} \mathcal{L}_{Y} f-\mathcal{L}_{Y} \mathcal{L}_{X} f=\mathcal{L}_{[X, Y]} f$. Since $\operatorname{End}\left(C^{\infty}(U, \mathbb{R})\right)$ is a Lie algebra with respect to the commutator bracket (Lemma 4.1.2) and $\mathcal{L}_{[X, Y]}=\left[\mathcal{L}_{X}, \mathcal{L}_{Y}\right]$, the Jacobi identity in $\operatorname{End}\left(C^{\infty}(U, \mathbb{R})\right.$ ) implies the Jacobi identity in $C^{\infty}\left(U, \mathbb{R}^{n}\right)$.

Remark 8.4.4. For any open subset $U \subseteq \mathbb{R}^{n}$, the map

$$
\mathcal{V}(U) \rightarrow C^{\infty}\left(U, \mathbb{R}^{n}\right), \quad X \rightarrow \widetilde{X}
$$

with $X(p)=(p, \widetilde{X}(p))$ is a linear isomorphism. We use this map to transfer the Lie bracket on $C^{\infty}\left(U, \mathbb{R}^{n}\right)$, defined in Lemma 8.4.3, to a Lie bracket on $\mathcal{V}(U)$, determined by

$$
[X, Y](p):=[\widetilde{X}, \widetilde{Y}](p)=\mathrm{d} \widetilde{Y}(p) \widetilde{X}(p)-\mathrm{d} \widetilde{X}(p) \widetilde{Y}(p)
$$

Our goal is to use the Lie brackets on the space $\mathcal{V}(U)$ and local charts to define a Lie bracket on $\mathcal{V}(M)$. The following lemma will be needed to ensure consistency in this process.

First, we introduce the concept or related vector fields. If $\varphi: M \rightarrow N$ is a smooth map, then we call two vector fields $X^{\prime} \in \mathcal{V}(M)$ and $X \in \mathcal{V}(N) \varphi$-related if

$$
X \circ \varphi=T \varphi \circ X^{\prime}: M \rightarrow T N
$$

With respect to the pullback map

$$
\varphi^{*}: C^{\infty}(N, \mathbb{R}) \rightarrow C^{\infty}(M, \mathbb{R}), \quad f \mapsto f \circ \varphi
$$

the $\varphi$-relatedness of $X$ and $X^{\prime}$ implies that\\
$\mathcal{L}_{X^{\prime}}\left(\varphi^{*} f\right)=\mathcal{L}_{X^{\prime}}(f \circ \varphi)=\mathrm{d}(f \circ \varphi) \circ X^{\prime}=\mathrm{d} f \circ T \varphi \circ X^{\prime}=\mathrm{d} f \circ X \circ \varphi=\varphi^{*}\left(\mathcal{L}_{X} f\right)$, i.e.,

$$
\mathcal{L}_{X^{\prime}} \circ \varphi^{*}=\varphi^{*} \circ \mathcal{L}_{X}
$$

Lemma 8.4.5. Let $M \subseteq \mathbb{R}^{n}$ and $N \subseteq \mathbb{R}^{m}$ be open subsets. Suppose that $X^{\prime}$, resp., $Y^{\prime} \in \mathcal{V}(M)$, is $\varphi$-related to $X$, resp., $Y \in \mathcal{V}(N)$. Then $\left[X^{\prime}, Y^{\prime}\right]$ is $\varphi$-related to $[X, Y]$.

Proof. In view of (8.5), we have

$$
\mathcal{L}_{X^{\prime}} \circ \varphi^{*}=\varphi^{*} \circ \mathcal{L}_{X} \quad \text { and } \quad \mathcal{L}_{Y^{\prime}} \circ \varphi^{*}=\varphi^{*} \circ \mathcal{L}_{Y}
$$

as linear maps $C^{\infty}(N, \mathbb{R}) \rightarrow C^{\infty}(M, \mathbb{R})$. Therefore,

$$
\begin{aligned}
{\left[\mathcal{L}_{X^{\prime}}, \mathcal{L}_{Y^{\prime}}\right] \circ \varphi^{*} } & =\mathcal{L}_{X^{\prime}} \circ \mathcal{L}_{Y^{\prime}} \circ \varphi^{*}-\mathcal{L}_{Y^{\prime}} \circ \mathcal{L}_{X^{\prime}} \circ \varphi^{*} \\
& =\varphi^{*} \circ \mathcal{L}_{X} \circ \mathcal{L}_{Y}-\varphi^{*} \circ \mathcal{L}_{Y} \circ \mathcal{L}_{X}=\varphi^{*} \circ\left[\mathcal{L}_{X}, \mathcal{L}_{Y}\right]
\end{aligned}
$$

For any $f \in C^{\infty}(N, \mathbb{R})$, we thus obtain

$$
\mathrm{d} f \circ T \varphi \circ\left[X^{\prime}, Y^{\prime}\right]=\mathcal{L}_{\left[X^{\prime}, Y^{\prime}\right]}(f \circ \varphi)=\left(\mathcal{L}_{[X, Y]} f\right) \circ \varphi=\mathrm{d} f \circ[X, Y] \circ \varphi
$$

Since each linear functional on the space $T_{x}(N) \cong \mathbb{R}^{m}$ is of the form $\mathrm{d} f(x)$ for some linear map $f: \mathbb{R}^{m} \rightarrow \mathbb{R}$, the assertion follows.

Proposition 8.4.6. For a vector field $X \in \mathcal{V}(M)$ and a chart $(\varphi, U)$ of $M$, we write $X_{\varphi}:=T \varphi \circ X \circ \varphi^{-1}$ for the corresponding vector field on the open subset $\varphi(U) \subseteq \mathbb{R}^{n}$.

For $X, Y \in \mathcal{V}(M)$, there exists a vector field $[X, Y] \in \mathcal{V}(M)$ which is uniquely determined by the property that for each chart ( $\varphi, U$ ) of $M$, the following equation holds

$$
[X, Y]_{\varphi}=\left[X_{\varphi}, Y_{\varphi}\right]
$$

Proof. If ( $\varphi, U$ ) and ( $\psi, V$ ) are charts of $M$, the vector fields $X_{\varphi}$ on $\varphi(U)$ and $X_{\psi}$ on $\psi(V)$ are ( $\psi \circ \varphi^{-1}$ )-related. Therefore, Lemma 8.4.5 implies that $\left[X_{\varphi}, Y_{\varphi}\right]$ is $\left(\psi \circ \varphi^{-1}\right)$-related to $\left[X_{\psi}, Y_{\psi}\right]$. This in turn is equivalent to

$$
T(\varphi)^{-1} \circ\left[X_{\varphi}, Y_{\varphi}\right] \circ \varphi=T(\psi)^{-1} \circ\left[X_{\psi}, Y_{\psi}\right] \circ \psi
$$

which is an identity of vector fields on the open subset $U \cap V$.\\
Hence there exists a unique vector field $[X, Y] \in \mathcal{V}(M)$, satisfying

$$
\left.[X, Y]\right|_{U}=T(\varphi)^{-1} \circ\left[X_{\varphi}, Y_{\varphi}\right] \circ \varphi
$$

for each chart $(\varphi, U)$, i.e., $[X, Y]_{\varphi}=\left[X_{\varphi}, Y_{\varphi}\right]$ on $\varphi(U)$.

Lemma 8.4.7. For $f \in C^{\infty}(M, \mathbb{R})$ and $X, Y \in \mathcal{V}(M)$, the following equation holds

$$
\mathcal{L}_{[X, Y]} f=\mathcal{L}_{X}\left(\mathcal{L}_{Y} f\right)-\mathcal{L}_{Y}\left(\mathcal{L}_{X} f\right)
$$

Proof. It suffices to show that this relation holds on $U$ for any chart $(\varphi, U)$ of $M$. For $f_{\varphi}:=f \circ \varphi^{-1}$, we then obtain with (8.5)

$$
\begin{aligned}
\mathcal{L}_{[X, Y]} f & =\mathrm{d} f \circ[X, Y]=\mathrm{d} f \circ T\left(\varphi^{-1}\right) \circ[X, Y]_{\varphi} \circ \varphi \\
& =\mathrm{d} f_{\varphi} \circ\left[X_{\varphi}, Y_{\varphi}\right] \circ \varphi=\varphi^{*}\left(\mathcal{L}_{\left[X_{\varphi}, Y_{\varphi}\right]} f_{\varphi}\right) \\
& =\varphi^{*}\left(\mathcal{L}_{\left(X_{\varphi}\right)} \mathcal{L}_{\left(Y_{\varphi}\right)} f_{\varphi}-\mathcal{L}_{\left(Y_{\varphi}\right)} \mathcal{L}_{\left(X_{\varphi}\right)} f_{\varphi}\right) \\
& =\mathcal{L}_{X}\left(\mathcal{L}_{Y} f\right)-\mathcal{L}_{Y}\left(\mathcal{L}_{X} f\right),
\end{aligned}
$$

because $\varphi^{*} f_{\varphi}=f$.\\
Proposition 8.4.8. ( $\mathcal{V}(M),[\cdot, \cdot]$ ) is a Lie algebra.\\
Proof. Clearly, (L1) and (L2) are satisfied. To verify the Jacobi identity, let $X, Y, Z \in \mathcal{V}(M)$ and $(\varphi, U)$ be a chart of $M$. For the vector field $J(X, Y, Z):= \sum_{\text {cycl. }}[X,[Y, Z]] \in \mathcal{V}(M)$, we then obtain from the definition of the bracket, Remark 8.4.4, and Proposition 8.4.6:

$$
J(X, Y, Z)_{\varphi}=J\left(X_{\varphi}, Y_{\varphi}, Z_{\varphi}\right)=0
$$

because $[\cdot, \cdot]$ is a Lie bracket on $\mathcal{V}(\varphi(U))$. This means that $J(X, Y, Z)$ vanishes on $U$, but since the chart ( $\varphi, U$ ) was arbitrary, $J(X, Y, Z)=0$.

We shall see later that the following lemma is an extremely important tool.

Lemma 8.4.9 (Related Vector Field Lemma). Let $M$ and $N$ be smooth manifolds, $\varphi: M \rightarrow N$ a smooth map, $X, Y \in \mathcal{V}(N)$, and $X^{\prime}, Y^{\prime} \in \mathcal{V}(M)$. If $X^{\prime}$ is $\varphi$-related to $X$ and $Y^{\prime}$ is $\varphi$-related to $Y$, then the Lie bracket $\left[X^{\prime}, Y^{\prime}\right]$ is $\varphi$-related to $[X, Y]$.

Proof. We have to show that for each $p \in M$ we have

$$
[X, Y](\varphi(p))=T_{p}(\varphi)\left[X^{\prime}, Y^{\prime}\right](p)
$$

Let ( $\rho, U$ ) be a chart of $M$ with $p \in U$ and $(\psi, V)$ a chart of $N$ with $\varphi(p) \in V$. Then the vector fields $X_{\rho}^{\prime}$ and $X_{\psi}$ are $\psi \circ \varphi \circ \rho^{-1}$-related on the domain $\rho\left(\varphi^{-1}(V) \cap U\right)$ :

$$
\begin{aligned}
& T\left(\psi \circ \varphi \circ \rho^{-1}\right) X_{\rho}^{\prime}=T\left(\psi \circ \varphi \circ \rho^{-1}\right)\left(T(\rho) \circ X^{\prime} \circ \rho^{-1}\right) \\
& =T(\psi) \circ T(\varphi) \circ X^{\prime} \circ \rho^{-1}=T(\psi) \circ X \circ \varphi \circ \rho^{-1}=X_{\psi} \circ\left(\psi \circ \varphi \circ \rho^{-1}\right)
\end{aligned}
$$

and the same holds for the vector fields $Y_{\rho}^{\prime}$ and $Y_{\psi}$, hence for their Lie brackets (Lemma 8.4.5).

Now the definition of the Lie bracket on $\mathcal{V}(N)$ and $\mathcal{V}(M)$ implies that

$$
\begin{aligned}
& T(\psi) \circ T(\varphi) \circ\left[X^{\prime}, Y^{\prime}\right]=T\left(\psi \circ \varphi \circ \rho^{-1}\right) \circ\left[X^{\prime}, Y^{\prime}\right]_{\rho} \circ \rho \\
& =T\left(\psi \circ \varphi \circ \rho^{-1}\right) \circ\left[X_{\rho}^{\prime}, Y_{\rho}^{\prime}\right] \circ \rho=\left[X_{\psi}, Y_{\psi}\right] \circ \psi \circ \varphi \circ \rho^{-1} \circ \rho \\
& =\left[X_{\psi}, Y_{\psi}\right] \circ \psi \circ \varphi=[X, Y]_{\psi} \circ \psi \circ \varphi=T(\psi) \circ[X, Y] \circ \varphi,
\end{aligned}
$$

and since $T(\psi)$ is injective, the assertion follows.\\
Example 8.4.10. Let $(\varphi, U)$ be a chart of $M$ and $x_{1}, \ldots, x_{n}: U \rightarrow \mathbb{R}$ the corresponding coordinate functions. Then we obtain on $U$ the vector fields $X_{j}, j=1, \ldots, n$, defined by

$$
X_{j}(p):=T_{p}(\varphi)^{-1} e_{j}:=\frac{\partial}{\partial x_{j}}(p):=\left.\frac{\partial}{\partial x_{j}}\right|_{p}
$$

where $e_{1}, \ldots, e_{n}$ is the standard basis for $\mathbb{R}^{n}$. We call these vector fields the $\varphi$-basic vector fields on $U$. The expression basic vector field is doubly justified. On the one hand, $\left(X_{1}(p), \ldots, X_{n}(p)\right)$ is a basis for $T_{p}(M)$ for every $p \in U$. On the other hand, the definition shows that every $X \in \mathcal{V}(U)$ can be written as

$$
X=\sum_{j=1}^{n} a_{j} \cdot X_{j} \quad \text { for } \quad a_{j} \in C^{\infty}(U)
$$

Since basic vector fields are $\varphi$-related with the constant vector fields on $\mathbb{R}^{n}$, they commute (Related Vector Field Lemma 8.4.9), i.e., $\left[X_{j}, X_{k}\right]=0$.

Remark 8.4.11. Let $M$ be a smooth manifold of dimension $n,(\varphi, U)$ be a chart on $M$, and $x_{j}: U \rightarrow \mathbb{R}$ be the $j$ th component of $\varphi$. Then, for $p \in U$, the differentials $\left(\mathrm{d} x_{j}(p)\right)_{j=1, \ldots, n}$ form the dual basis to $\left(\left.\frac{\partial}{\partial x_{j}}\right|_{p}\right)_{j=1, \ldots, n}$ for $T_{p} M^{*}$. Indeed, $x_{j} \circ \varphi^{-1}: \varphi(U) \rightarrow \mathbb{R}$ is simply the projection to the $j$ th coordinate, so

$$
\mathrm{d} x_{i}(p)\left(\left.\frac{\partial}{\partial x_{j}}\right|_{p}\right)=\frac{\partial x_{i}}{\partial x_{j}}(p)=\delta_{i j}= \begin{cases}1 & \text { for } i=j \\ 0 & \text { for } i \neq j\end{cases}
$$

(Kronecker delta).\\
Example 8.4.12. Let $(\varphi, U)$ be a chart of $M$ and $x_{1}, \ldots, x_{n}: U \rightarrow \mathbb{R}$ the corresponding coordinate functions. Then we obtain on $U$ the 1 -forms $\mathrm{d} x_{j}$, $j=1, \ldots, n$. We call these 1 -forms the $\varphi$-basic forms on $U$. As in the case of $\varphi$-basic vector fields the expression basic form is doubly justified. On the one hand, $\left(\mathrm{d} x_{1}(p), \ldots, \mathrm{d} x_{n}(p)\right)$ is a basis for the dual space $T_{p}(M)^{*}$ for every $p \in U$. On the other hand, setting $a_{j}:=\omega\left(\frac{\partial}{\partial x_{j}}\right)$, the formula $\mathrm{d} x_{k}\left(\frac{\partial}{\partial x_{j}}\right)=\delta_{k j}$ implies that every $\omega \in \Omega^{1}(U)$ can be written as

$$
\omega=\sum_{j=1}^{n} a_{j} \cdot \mathrm{~d} x_{j} \quad \text { for } \quad a_{j} \in C^{\infty}(U)
$$

\subsubsection*{Vector Fields as Derivations}
Let $M$ be a smooth manifold with tangent bundle $T M$ and $\pi_{T M}: T M \rightarrow M$ the canonical projection. Recall from Definition 8.4.1 that a vector field on $M$ is a smooth section of the tangent bundle and that it induces a map

$$
\mathcal{L}_{X}: C^{\infty}(M) \rightarrow C^{\infty}(M), \quad f \mapsto \mathcal{L}_{X} f=\mathrm{d} f \circ X
$$

Then the Product Rule $\mathrm{d}(f g)=f \mathrm{~d} g+g \mathrm{~d} f$ implies

$$
\mathcal{L}_{X}(f \cdot g)=\mathcal{L}_{X} f \cdot g+f \cdot \mathcal{L}_{X} g
$$

Definition 8.4.13. Let $A$ be an associative algebra over $\mathbb{K}$. A linear map $D: A \rightarrow A$ is called a derivation of $A$, if it satisfies the following property:

$$
D(f \cdot g)=D(f) \cdot g+f \cdot D(g) \quad \text { for } f, g \in A
$$

The set of all derivations of $A$ is denoted by $\operatorname{der}(A)$.\\
Remark 8.4.14. The discussion above shows that for a vector field $X \in \mathcal{V}(M)$ the map $\mathcal{L}_{X}$ is a derivation of the algebra $C^{\infty}(M)$. We thus obtain a map

$$
\mathcal{V}(M) \rightarrow \operatorname{der}\left(C^{\infty}(M)\right), \quad X \mapsto \mathcal{L}_{X}
$$

Below, we will show that this map is a bijection, and after that we will also write $X f$ for $\mathcal{L}_{X} f$. Note that the set $\operatorname{der}\left(C^{\infty}(M)\right)$ carries several algebraic structures. It is obvious that we can add derivations and multiply them with scalars. Derivations can also be multiplied with functions by

$$
(g \cdot D)(f):=g \cdot(D(f)) \quad \forall f, g \in C^{\infty}(M)
$$

Moreover, it is easy to check directly that $\operatorname{der}\left(C^{\infty}(M)\right)$ is closed under the commutator bracket (this also follows from Lemma 8.4.7 once $\operatorname{der}\left(C^{\infty}(M)\right.$ ) is identified with $\mathcal{V}(M)$ ). One thus obtains that $\operatorname{der}\left(C^{\infty}(M)\right)$ is at the same time a Lie algebra and a $C^{\infty}(M)$-module (cf. Exercise 5.1.4).

Lemma 8.4.15. Let $M$ be a smooth manifold.\\
(i) Let $C$ be a compact subset of $M$ and $V \supseteq C$ be an open subset. Then there exists a smooth function $f$ with $\left.f\right|_{C}=1$ and $\left.f\right|_{M \backslash V}=0$.\\
(ii) For each $p \in M$, the map $C^{\infty}(M) \rightarrow C^{\infty}(M)_{p}, f \mapsto f_{p}$ is surjective.

Proof. To prove (i), we first claim that for each $p \in M$ and each open subset $B$ containing $p$, there exists a smooth function $f \in C^{\infty}(B)$ with $f=0$ on $M \backslash B$ and $f=1$ on a neighborhood $A$ of $p$.

Since this is local property, it suffice s to prove it for balls in $\mathbb{R}^{n}$. So let $A=B_{a}(0)$ and $B=B_{b}(0)$ be the open balls of radius $a$ and $b$ around 0 , where $0<a<b$. Consider the function $f: \mathbb{R} \rightarrow \mathbb{R}$ which is defined by

$$
f(x)= \begin{cases}\exp \left(\frac{1}{x-b}-\frac{1}{x-a}\right) & \text { for } a<x<b \\ 0 & \text { otherwise }\end{cases}
$$

It is elementary to check (Exercise) that $f$ and the function $F: \mathbb{R} \rightarrow \mathbb{R}$ defined by

$$
F(x)=\frac{\int_{x}^{b} f(t) d t}{\int_{a}^{b} f(t) d t}
$$

are smooth.\\
\includegraphics[max width=\textwidth, center]{2025_10_18_3eaf04e77f3cb364fce5g-278(2)}\\
\includegraphics[max width=\textwidth, center]{2025_10_18_3eaf04e77f3cb364fce5g-278}

We define our desired function by

$$
\zeta\left(x_{1}, \ldots, x_{n}\right):=F\left(\sqrt{x_{1}^{2}+\cdots+x_{n}^{2}}\right) .
$$

\begin{center}
\includegraphics[max width=\textwidth]{2025_10_18_3eaf04e77f3cb364fce5g-278(1)}
\end{center}

One immediately checks that $\left.\zeta\right|_{A} \equiv 1$ and $\left.\zeta\right|_{\mathbb{R}^{n} \backslash B} \equiv 0$, so that the claim is proved.

Now cover $C$ by finitely many open sets $U_{i} \subseteq V$ such that one has smooth functions $f_{i} \in C^{\infty}(M)$ with $f_{i}=0$ on $M \backslash V$ and $f_{i}=1$ on $U_{i}$. Then\\
$f:=1-\prod_{i}\left(1-f_{i}\right)$ has the desired properties which completes the proof of (i).

Finally, let $U$ be an open neighborhood of $p \in M$ and $h \in C^{\infty}(U)$. Further, let $W$ be an open neighborhood of $p$ with compact closure $\bar{W} \subseteq U$. Let $f \in C^{\infty}(M)$ be a smooth function vanishing on $M \backslash W$ and which is 1 on a neighborhood of $p$. Then

$$
H(x):= \begin{cases}f(x) h(x) & \text { for } x \in U \\ 0 & \text { for } x \notin \bar{W}\end{cases}
$$

defines a smooth function on $M$ whose germ at $p$ is $H_{p}=h_{p}$.\\
Lemma 8.4.16. The map $\mathcal{V}(M) \rightarrow \operatorname{der}\left(C^{\infty}(M)\right), X \mapsto \mathcal{L}_{X}$ is injective.\\
Proof. Let $X \in \mathcal{V}(M)$ with $\mathcal{L}_{X}=0$. Then by (8.7) we have

$$
X(p)\left(f_{p}\right)=\left(\mathcal{L}_{X} f\right)(p)=0
$$

for all $f \in C^{\infty}(M)$ and all $p \in M$. Since each germ at $p$ comes from a global smooth function on $M$ by Lemma 8.4.15, this implies $X(p)=0$ for each $p \in M$, hence $X=0$.

Lemma 8.4.17. Let $M$ be a smooth manifold and $D \in \operatorname{der}\left(C^{\infty}(M)\right)$.\\
(i) $D(f)=0$ for all constant functions $f \in C^{\infty}(M)$.\\
(ii) Let $V \subset M$ be an open subset and $f \in C^{\infty}(M)$ with $\left.f\right|_{V}=0$. Then $\left.D(f)\right|_{V}=0$ as well.

Proof. Just as Proposition 8.3.20, the assertion (i) is a direct consequence of the defining equation of a derivation. To show (ii), pick $p \in V$. By Lemma 8.4.15, there exists a function $g \in C^{\infty}(M)$ which is 1 outside $V$ and which vanishes at $p$. Then $g \cdot f=f$, and therefore,

$$
D(f)(p)=D(f \cdot g)(p)=f(p)(D(g)(p))+g(p)(D(f)(p))=0
$$

Since $p \in V$ was chosen arbitrarily, the lemma is proved.\\
To finish the program outlined in Remark 8.4.14, we still have to show that the assignment $X \mapsto \mathcal{L}_{X}$ is also surjective. So, let $D \in \operatorname{der}\left(C^{\infty}(M)\right)$ be given. By Lemma 8.4.17, $D(f)(p)$ only depends on $f_{p}$, so, in view of Lemma 8.4.16, we can define an element $X_{D}(p) \in T_{p}(M)$ via

$$
X_{D}(p)\left(f_{p}\right):=D(f)(p)
$$

for all $f \in C^{\infty}(M)$ and for all $p \in M$.

All that remains to show to complete the proof of the bijectivity of $X \mapsto \mathcal{L}_{X}$ is the smoothness of $X_{D}$ at each $p \in M$. This is a local property, so it suffices to verify the smoothness of $X_{D}$ on some neighborhood of $p$. So let ( $\varphi, U$ ) be a chart of $M$ with $p \in U$. Let $\widehat{x}_{j} \in C^{\infty}(M)$ be functions which coincide with $x_{j}$ on a neighborhood $V$ of $p$ (Lemma 8.4.15). Then Proposition 8.3.23 implies for $q \in V$

$$
X_{D}(q)=\left.\sum_{j=1}^{n} X_{D}(q)\left(\left(x_{j}\right)_{q}\right) \cdot \frac{\partial}{\partial x_{j}}\right|_{q}=\left.\sum_{j=1}^{n} D\left(\widehat{x}_{j}\right)(q) \cdot \frac{\partial}{\partial x_{j}}\right|_{q}
$$

Therefore, $X_{D}$ is smooth on $V$. This completes the proof of the following theorem.

Theorem 8.4.18. Let $M$ be a smooth manifold. Then the map

$$
\mathcal{V}(M) \rightarrow \operatorname{der}\left(C^{\infty}(M)\right), \quad X \mapsto \mathcal{L}_{X}
$$

defined by $\mathcal{L}_{X}(f)=\mathrm{d} f \circ X$, is a bijection.\\
Theorem 8.4.18 is interesting not only in its own right, but also because of the methods that were used to prove it. We collect some of the facts we essentially derived in the course of the proof:

Remark 8.4.19. Let $M$ be a smooth manifold and $p \in M$.\\
(i) Let $C \subseteq M$ be a compact subset, $U \supseteq C$ an open neighborhood, and $X \in \mathcal{V}(U)$ a vector field on $U$. Let $W \subseteq U$ be a compact neighborhood of $C$ in $U$ and $f \in C^{\infty}(M)$ vanishing on $M \backslash W$ with $\left.f\right|_{C}=1$ (Lemma 8.4.15(i)). Then $f X$ defines a smooth vector field on $M$ which coincides with $X$ on $C$.\\
(ii) For any chart $(\varphi, U)$ with $p \in U$ a vector field $X \in \operatorname{der}\left(C^{\infty}(M)\right) \cong \mathcal{V}(M)$ can be written on $U$ as

$$
\left.X\right|_{U}=\sum_{j=1}^{n} a_{j} \cdot X_{j}
$$

with $a_{j} \in C^{\infty}(U)$ and $X_{j} \in \mathcal{V}(U)$, given by $\left.\frac{\partial}{\partial x_{j}}\right|_{q}$ for all $q \in U$.\\
(iii) For every $v \in T_{p} M$, there exists a vector field $X \in \mathcal{V}(M)$ with $X(p)=v$.

\subsubsection*{Exercises for Section 8.4}
Exercise 8.4.1. Let $M$ be a smooth manifold and $q: \mathbb{V} \rightarrow M$ a smooth $\mathbb{K}$-vector bundle. Show that

$$
\left(s_{1}+s_{2}\right)(p):=s_{1}(p)+s_{2}(p), \quad(\lambda s)(p):=\lambda s(p), \quad \lambda \in \mathbb{K}
$$

defines on the space $\Gamma \mathbb{V}$ of smooth sections of $\mathbb{V}$ the structure of a $\mathbb{K}$-vector space. Show also that the multiplication with smooth $\mathbb{K}$-valued functions defined by

$$
(f s)(p):=f(p) s(p)
$$

satisfies for $s, s_{1}, s_{2} \in \Gamma \mathbb{V}$ and $f, g \in C^{\infty}(M, \mathbb{K})$ :\\
(1) $f\left(s_{1}+s_{2}\right)=f s_{1}+f s_{2}$.\\
(2) $f(\lambda s)=\lambda \cdot f s=(\lambda f) s$ for $\lambda \in \mathbb{K}$.\\
(3) $(f+g) s=f s+g s$.\\
(4) $f(g s)=(f g) s$.

Exercise 8.4.2. Let $M$ be a smooth manifold, $X, Y \in \mathcal{V}(M)$ and $f, g \in C^{\infty}(M, \mathbb{R})$. Show that\\
(1) $\mathcal{L}_{X}(f \cdot g)=\mathcal{L}_{X}(f) \cdot g+f \cdot \mathcal{L}_{X}(g)$, i.e., the map $f \mapsto \mathcal{L}_{X}(f)$ is a derivation.\\
(2) $\mathcal{L}_{f X}(g)=f \cdot \mathcal{L}_{X}(g)$.

Exercise 8.4.3. Let $M$ be a smooth manifold of dimension $n$ and $p \in M$. Pick a vector $v_{\varphi} \in \mathbb{R}^{n}$ for each chart ( $\varphi, U$ ) with $p \in U$. We call the family $\left(v_{\varphi}\right)$ of all such vectors a tangent vector (physicists' definition) if

$$
v_{\psi}=\mathrm{d}\left(\psi \circ \varphi^{-1}\right)_{\varphi(p)} v_{\varphi}
$$

It is obvious that these tangent vectors form a vector space isomorphic to $\mathbb{R}^{n}$. Fix a chart ( $\varphi, U$ ) with $p \in U$. Consider the map

$$
\mathbb{R}^{n} \rightarrow T_{p} M,\left.\quad v_{\varphi} \mapsto \sum_{j=1}^{n}\left(v_{\varphi}\right)_{j} \frac{\partial}{\partial x_{j}}\right|_{p}
$$

and show that it defines a linear isomorphism between the physicists and the algebraic version of the tangent space.

Exercise 8.4.4. Let $M$ be a smooth manifold, and let ( $\varphi, U$ ) and ( $\psi, V$ ) be charts on $M$. Further, let $p \in U \cap V$ and $v \in T_{p} M$. Express $v$ in the $\varphi$-basis for $T_{p} M$ as well as in the $\psi$-basis for $T_{p} M$. How are the coefficients related?\\
Exercise 8.4.5. Let $f: M \rightarrow N$ be a smooth map between smooth manifolds. Pick charts ( $\varphi, U$ ) and ( $\psi, V$ ) of $M$ and $N$, respectively, with $p \in U$ and $f(p) \in V$. Show that the derivative $T_{p} f: T_{p} M \rightarrow T_{f(p)} N$ of $f$ at $p$ in the physicists' definition of tangent space is given by

$$
\left(T_{p} f\left(v_{\varphi}\right)\right)_{\psi}=\mathrm{d}\left(\psi \circ f \circ \varphi^{-1}\right)_{\varphi(p)}\left(v_{\varphi}\right)
$$

Exercise 8.4.6. Show that $\operatorname{der}\left(C^{\infty}(M)\right)$ is a $C^{\infty}(M)$-module with respect to the multiplication (8.9). This means that for all $D, E \in \operatorname{der}\left(C^{\infty}(M)\right)$, $g, f \in C^{\infty}(M)$, and $r, s \in \mathbb{R}$ the following properties are satisfied:

$$
\begin{aligned}
g \cdot(r D+s E) & =r(g \cdot D)+s(g \cdot E) \\
g \cdot(f \cdot D) & =(g \cdot f) \cdot D \\
(g+f) \cdot D & =g \cdot D+f \cdot D
\end{aligned}
$$

\subsection*{Integral Curves and Local Flows}
In this section, we turn to the geometric nature of vector fields as infinitesimal generators of local flows on manifolds. This provides a natural perspective on (autonomous) ordinary differential equations.

\subsubsection*{Integral Curves}
Throughout this subsection, $M$ denotes an $n$-dimensional smooth manifold.\\
Definition 8.5.1. Let $X \in \mathcal{V}(M)$ and $I \subseteq \mathbb{R}$ an open interval containing 0 . A differentiable map $\gamma: I \rightarrow M$ is called an integral curve of $X$ if

$$
\gamma^{\prime}(t)=X(\gamma(t)) \quad \text { for each } \quad t \in I .
$$

Note that the preceding equation implies that $\gamma^{\prime}$ is continuous and further that if $\gamma$ is $C^{k}$, then $\gamma^{\prime}$ is also $C^{k}$. Therefore, integral curves are automatically smooth.

If $J \supseteq I$ is an interval containing $I$, then an integral curve $\eta: J \rightarrow M$ is called an extension of $\gamma$ if $\left.\eta\right|_{I}=\gamma$. An integral curve $\gamma$ is said to be maximal if it has no proper extension.

Remark 8.5.2. (a) If $U \subseteq \mathbb{R}^{n}$ is an open subset of $\mathbb{R}^{n}$, then we write a vector field $X \in \mathcal{V}(U)$ as $X(x)=(x, F(x))$, where $F: U \rightarrow \mathbb{R}^{n}$ is a smooth function. A curve $\gamma: I \rightarrow U$ is an integral curve of $X$ if and only if it satisfies the ordinary differential equation

$$
\gamma^{\prime}(t)=F(\gamma(t)) \quad \text { for all } \quad t \in I
$$

(b) If ( $\varphi, U$ ) is a chart of the manifold $M$ and $X \in \mathcal{V}(M)$, then a curve $\gamma: I \rightarrow M$ is an integral curve of $X$ if and only if the curve $\eta:=\varphi \circ \gamma$ is an integral curve of the vector field $X_{\varphi}:=T(\varphi) \circ X \circ \varphi^{-1} \in \mathcal{V}(\varphi(U))$ because

$$
X_{\varphi}(\eta(t))=T_{\gamma(t)}(\varphi) X(\gamma(t)) \quad \text { and } \quad \eta^{\prime}(t)=T_{\gamma(t)}(\varphi) \gamma^{\prime}(t)
$$

Remark 8.5.3. A curve $\gamma: I \rightarrow M$ is an integral curve of $X$ if and only if $\widetilde{\gamma}(t):=\gamma(-t)$ is an integral curve of the vector field $-X$.

More generally, for $a, b \in \mathbb{R}$, the curve $\eta(t):=\gamma(a t+b)$ is an integral curve of the vector field $a X$.

Definition 8.5.4. Let $a<b \in[-\infty, \infty]$. For a continuous curve $\gamma:] a, b[\rightarrow M$ we say that

$$
\lim _{t \rightarrow b} \gamma(t)=\infty
$$

if for each compact subset $K \subseteq M$ there exists a $c<b$ with $\gamma(t) \notin K$ for $t>c$. Similarly, we define

$$
\lim _{t \rightarrow a} \gamma(t)=\infty
$$

Theorem 8.5.5 (Existence and Uniqueness of Integral Curves). Let $X \in \mathcal{V}(M)$ and $p \in M$. Then there exists a unique maximal integral curve $\gamma_{p}: I_{p} \rightarrow M$ with $\gamma_{p}(0)=p$. If $a:=\inf I_{p}>-\infty$, then $\lim _{t \rightarrow a} \gamma_{p}(t)=\infty$ and if $b:=\sup I_{p}<\infty$, then $\lim _{t \rightarrow b} \gamma_{p}(t)=\infty$.

Proof. We have seen in Remark 8.5.2 that in local charts, integral curves are solutions of an ordinary differential equation with a smooth right-hand side. We now reduce the proof to the Local Existence- and Uniqueness Theorem for ODEs.

Uniqueness: Let $\gamma, \eta: I \rightarrow M$ be two integral curves of $X$ with $\gamma(0)= \eta(0)=p$. The continuity of the curves implies that

$$
0 \in J:=\{t \in I: \gamma(t)=\eta(t)\}
$$

is a closed subset of $I$. In view of the Local Uniqueness Theorem for ODEs, for each $t_{0} \in J$ there exists an $\varepsilon>0$ with $\left[t_{0}, t_{0}+\varepsilon\right] \subseteq J$, and likewise $\left[t_{0}-\varepsilon, t_{0}\right] \subseteq J$ (Remark 8.5.3). Therefore, $J$ is also open. Now the connectedness of $I$ implies $I=J$, so that $\gamma=\eta$.

Existence: The Local Existence Theorem implies the existence of some integral curve $\gamma: I \rightarrow M$ on some open interval containing 0 . For any other integral curve $\eta: J \rightarrow M$, the intersection $I \cap J$ is an interval containing 0 , so that the uniqueness assertion implies that $\eta=\gamma$ on $I \cap J$.

Let $I_{p} \subseteq \mathbb{R}$ be the union of all open intervals $I_{j}$ containing 0 on which there exists an integral curve $\gamma_{j}: I_{j} \rightarrow M$ of $X$ with $\gamma_{j}(0)=p$. Then the preceding argument shows that

$$
\gamma(t):=\gamma_{j}(t) \quad \text { for } \quad t \in I_{j}
$$

defines an integral curve of $X$ on $I_{p}$, which is maximal by definition. The uniqueness of the maximal integral curve also follows from its definition.

Limit condition: Suppose that $b:=\sup I_{p}<\infty$. If $\lim _{t \rightarrow b} \gamma(t)=\infty$ does not hold, then there exists a compact subset $K \subseteq M$ and a sequence $t_{m} \in I_{p}$ with $t_{m} \rightarrow b$ and $\gamma\left(t_{m}\right) \in K$. As $K$ can be covered with finitely many closed subsets homeomorphic to a closed subset of a ball in $\mathbb{R}^{n}$, after passing to a suitable subsequence, we may w.l.o.g. assume that $K$ itself is homeomorphic to a compact subset of $\mathbb{R}^{n}$. Then a subsequence of $\left(\gamma\left(t_{m}\right)\right)_{m \in \mathbb{N}}$ converges, and we may replace the original sequence by this subsequence, hence assume that $q:=\lim _{m \rightarrow \infty} \gamma\left(t_{m}\right)$ exists.

The Local Existence Theorem for ODEs implies the existence of a compact neighborhood $V \subseteq M$ of $q$ and $\varepsilon>0$ such that the initial value problem

$$
\eta(0)=x, \quad \eta^{\prime}=X \circ \eta
$$

has a solution on $[-\varepsilon, \varepsilon]$ for each $x \in V$. Pick $m \in \mathbb{N}$ with $t_{m}>b-\varepsilon$ and $\gamma\left(t_{m}\right) \in V$. Further let $\eta:[-\varepsilon, \varepsilon] \rightarrow M$ be an integral curve with $\eta(0)= \gamma\left(t_{m}\right)$. Then

$$
\gamma(t):=\eta\left(t-t_{m}\right) \quad \text { for } \quad t \in\left[t_{m}-\varepsilon, t_{m}+\varepsilon\right]
$$

defines an extension of $\gamma$ to the interval $\left.I_{p} \cup\right] t_{m}, t_{m}+\varepsilon[$ strictly containing $] a, b\left[\right.$, hence contradicting the maximality of $I_{p}$. This proves that $\lim _{t \rightarrow b} \gamma(t)=\infty$. Replacing $X$ by $-X$, we also obtain $\lim _{t \rightarrow a} \gamma(t)=\infty$.

If $q=\gamma_{p}(t)$ is a point on the unique maximal integral curve of $X$ through $p \in M$, then $I_{q}=I_{p}-t$ and

$$
\gamma_{q}(s):=\gamma_{p}(t+s)
$$

is the unique maximal integral curve through $q$. Here $I_{p}$ is the domain of definition of the maximal integral curve through $p$ and $I_{q}$ is the domain of definition of the maximal integral curve through $q$.\\
Example 8.5.6. (a) On $M=\mathbb{R}$, we consider the vector field $X$ given by the function $F(s)=1+s^{2}$, i.e., $X(s)=\left(s, 1+s^{2}\right)$. The corresponding ODE is

$$
\gamma^{\prime}(s)=X(\gamma(s))=1+\gamma(s)^{2}
$$

For $\gamma(0)=0$, the function $\gamma(s):=\tan (s)$ on $I:=]-\frac{\pi}{2}, \frac{\pi}{2}[$ is the unique maximal solution because

$$
\lim _{t \rightarrow \frac{\pi}{2}} \tan (t)=\infty \quad \text { and } \quad \lim _{t \rightarrow-\frac{\pi}{2}} \tan (t)=-\infty
$$

(b) Let $M:=]-1,1[$ and $X(s)=(s, 1)$, so that the corresponding ODE is $\gamma^{\prime}(s)=1$. Then the unique maximal solution is

$$
\gamma(s)=s, \quad I=]-1,1[
$$

Note that we also have in this case

$$
\lim _{s \rightarrow \pm 1} \gamma(s)=\infty
$$

if we consider $\gamma$ as a curve in the noncompact manifold $M$.\\
For $M=\mathbb{R}$, the same vector field has the maximal integral curve

$$
\gamma(s)=s, \quad I=\mathbb{R}
$$

(c) For $M=\mathbb{R}$ and $X(s)=(s,-s)$, the differential equation is $\gamma^{\prime}(t)= -\gamma(t)$, so that we obtain the maximal integral curves $\gamma(t)=\gamma_{0} e^{-t}$. For $\gamma_{0}=0$, this curve is constant; and for $\gamma_{0} \neq 0$, we have $\lim _{t \rightarrow \infty} \gamma(t)=0$, hence $\lim _{t \rightarrow \infty} \gamma(t) \neq \infty$. This shows that maximal integral curves do not always leave every compact subset of $M$ if they are defined on an interval unbounded from above.

The preceding example shows in particular that the global existence of integral curves can also be destroyed by deleting parts of the manifold $M$, i.e., by considering $M^{\prime}:=M \backslash K$ for some closed subset $K \subseteq M$.\\
Definition 8.5.7. A vector field $X \in \mathcal{V}(M)$ is said to be complete if all its maximal integral curves are defined on all of $\mathbb{R}$.\\
Corollary 8.5.8. All vector fields on a compact manifold $M$ are complete.

\subsubsection*{Local Flows}
Definition 8.5.9. Let $M$ be a smooth manifold. A local flow on $M$ is a smooth map

$$
\Phi: U \rightarrow M
$$

where $U \subseteq \mathbb{R} \times M$ is an open subset containing $\{0\} \times M$, such that for each $x \in M$ the intersection $I_{x}:=U \cap(\mathbb{R} \times\{x\})$ is an interval containing 0 and

$$
\Phi(0, x)=x \quad \text { and } \quad \Phi(t, \Phi(s, x))=\Phi(t+s, x)
$$

hold for all $t, s, x$ for which both sides are defined. The maps

$$
\alpha_{x}: I_{x} \rightarrow M, \quad t \mapsto \Phi(t, x)
$$

are called the flow lines. The flow $\Phi$ is said to be global if $U=\mathbb{R} \times M$.\\
Lemma 8.5.10. If $\Phi: U \rightarrow M$ is a local flow, then

$$
X^{\Phi}(x):=\left.\frac{d}{d t}\right|_{t=0} \Phi(t, x)=\alpha_{x}^{\prime}(0)
$$

defines a smooth vector field.\\
It is called the velocity field or the infinitesimal generator of the local flow $\Phi$.

Lemma 8.5.11. If $\Phi: U \rightarrow M$ is a local flow on $M$, then the flow lines are integral curves of the vector field $X^{\Phi}$. In particular, the local flow $\Phi$ is uniquely determined by the vector field $X^{\Phi}$.

Proof. Let $\alpha_{x}: I_{x} \rightarrow M$ be a flow line and $s \in I_{x}$. For sufficiently small $t \in \mathbb{R}$, we then have

$$
\alpha_{x}(s+t)=\Phi(s+t, x)=\Phi(t, \Phi(s, x))=\Phi\left(t, \alpha_{x}(s)\right)
$$

so that taking derivatives in $t=0$ leads to $\alpha_{x}^{\prime}(s)=X^{\Phi}\left(\alpha_{x}(s)\right)$.\\
That $\Phi$ is uniquely determined by the vector field $X^{\Phi}$ follows from the uniqueness of integral curves (Theorem 8.5.5).

Theorem 8.5.12. Each smooth vector field $X$ is the velocity field of a unique local flow defined by

$$
\mathcal{D}_{X}:=\bigcup_{x \in M} I_{x} \times\{x\} \quad \text { and } \quad \Phi(t, x):=\gamma_{x}(t) \quad \text { for } \quad(t, x) \in \mathcal{D}_{X}
$$

where $\gamma_{x}: I_{x} \rightarrow M$ is the unique maximal integral curve through $x \in M$.

Proof. If $(s, x),(t, \Phi(s, x))$ and $(s+t, x) \in \mathcal{D}_{X}$, the relation

$$
\Phi(s+t, x)=\Phi(t, \Phi(s, x)) \quad \text { and } \quad I_{\Phi(s, x)}=I_{\gamma_{x}(s)}=I_{x}-s
$$

follow from the fact that both curves

$$
t \mapsto \Phi(t+s, x)=\gamma_{x}(t+s) \quad \text { and } \quad t \mapsto \Phi(t, \Phi(s, x))=\gamma_{\Phi(s, x)}(t)
$$

are integral curves of $X$ with the initial value $\Phi(s, x)$, hence coincide.\\
We claim that all maps

$$
\Phi_{t}: M_{t}:=\left\{x \in M:(t, x) \in \mathcal{D}_{X}\right\} \rightarrow M, \quad x \mapsto \Phi(t, x)
$$

are injective. In fact, if $p:=\Phi_{t}(x)=\Phi_{t}(y)$, then $\gamma_{x}(t)=\gamma_{y}(t)$, and on $[0, t]$ the curves $s \mapsto \gamma_{x}(t-s), \gamma_{y}(t-s)$ are integral curves of $-X$, starting in $p$. Hence the Uniqueness Theorem 8.5.5 implies that they coincide for $s=t$, which mans that $x=\gamma_{x}(0)=\gamma_{y}(0)=y$. From this argument, it further follows that $\Phi_{t}\left(M_{t}\right)=M_{-t}$ and $\Phi_{t}^{-1}=\Phi_{-t}$.

It remains to show that $\mathcal{D}_{X}$ is open and $\Phi$ smooth. The local Existence Theorem provides for each $x \in M$ an open neighborhood $U_{x}$ diffeomorphic to a cube and some $\varepsilon_{x}>0$, as well as a smooth map

$$
\left.\varphi_{x}:\right]-\varepsilon_{x}, \varepsilon_{x}\left[\times U_{x} \rightarrow M, \quad \varphi_{x}(t, y)=\gamma_{y}(t)=\Phi(t, y)\right.
$$

Hence $]-\varepsilon_{x}, \varepsilon_{x}\left[\times U_{x} \subseteq \mathcal{D}_{X}\right.$, and the restriction of $\Phi$ to this set is smooth. Therefore, $\Phi$ is smooth on a neighborhood of $\{0\} \times M$ in $\mathcal{D}_{X}$.

Now let $J_{x}$ be the set of all $t \in\left[0, \infty\left[\right.\right.$ for which $\mathcal{D}_{X}$ contains a neighborhood of $[0, t] \times\{x\}$ on which $\Phi$ is smooth. The interval $J_{x}$ is open in $\mathbb{R}^{+}:=\left[0, \infty\left[\right.\right.$ by definition. We claim that $J_{x}=I_{x} \cap \mathbb{R}^{+}$. This entails that $\mathcal{D}_{X}$ is open because the same argument applies to $\left.\left.I_{x} \cap\right]-\infty, 0\right]$.

We assume the contrary and find a minimal $\tau \in I_{x} \cap \mathbb{R}^{+} \backslash J_{x}$ because this interval is closed. Put $p:=\Phi(\tau, x)$ and pick a product set $I \times W \subseteq \mathcal{D}_{X}$, where $W$ is an open neighborhood of $p$ and $I=]-2 \varepsilon, 2 \varepsilon[$ a 0 -neighborhood such that $2 \varepsilon<\tau$ and $\Phi: I \times W \rightarrow M$ is smooth. By assumption, there exists an open neighborhood $V$ of $x$ such that $\Phi$ is smooth on $[0, \tau-\varepsilon] \times V \subseteq \mathcal{D}_{X}$. Then $\Phi_{\tau-\varepsilon}$ is smooth on $V$ and

$$
V^{\prime}:=\Phi_{\tau-\varepsilon}^{-1}\left(\Phi_{\varepsilon}^{-1}(W)\right) \cap V
$$

is a neighborhood of $[0, \tau+\varepsilon] \times\{x\}$ in $\mathcal{D}_{X}$. Further,

$$
V^{\prime}=\Phi_{\tau-\varepsilon}^{-1}\left(\Phi_{\varepsilon}^{-1}(W)\right) \cap V=\Phi_{\tau}^{-1}(W) \cap V,
$$

and $\Phi$ is smooth on $V^{\prime}$ because it is a composition of smooth maps:

$$
] \tau-2 \varepsilon, \tau+2 \varepsilon\left[\times V^{\prime} \rightarrow M, \quad(t, y) \mapsto \Phi(t-\tau, \Phi(\varepsilon, \Phi(\tau-\varepsilon, y))) .\right.
$$

We thus arrive at the contradiction $\tau \in J_{x}$.\\
This completes the proof of the openness of $\mathcal{D}_{X}$ and the smoothness of $\Phi$. The uniqueness of the flow follows from the uniqueness of the integral curves.

Remark 8.5.13. Let $X \in \mathcal{V}(M)$ be a complete vector field. If

$$
\Phi^{X}: \mathbb{R} \times M \rightarrow M
$$

is the corresponding global flow, then the maps $\Phi_{t}^{X}: x \mapsto \Phi^{X}(t, x)$ satisfy\\
(A1) $\Phi_{0}^{X}=\mathrm{id}_{M}$.\\
(A2) $\Phi_{t+s}^{X}=\Phi_{t}^{X} \circ \Phi_{s}^{X}$ for $t, s \in \mathbb{R}$.\\
It follows in particular that $\Phi_{t}^{X} \in \operatorname{Diff}(M)$ with $\left(\Phi_{t}^{X}\right)^{-1}=\Phi_{-t}^{X}$, so that we obtain a group homomorphism

$$
\gamma_{X}: \mathbb{R} \rightarrow \operatorname{Diff}(M), \quad t \mapsto \Phi_{t}^{X}
$$

With respect to the terminology introduced below, (A1) and (A2) mean that $\Phi^{X}$ defines a smooth action of $\mathbb{R}$ on $M$. As $\Phi^{X}$ is determined by the vector field $X$, we call $X$ the infinitesimal generator of this action. In this sense, the smooth $\mathbb{R}$-actions on a manifold $M$ are in one-to-one correspondence with the complete vector fields on $M$.

Remark 8.5.14. Let $\Phi^{X}: \mathcal{D}_{X} \rightarrow M$ be the maximal local flow of a vector field $X$ on $M$. Let $M_{t}=\left\{x \in M:(t, x) \in \mathcal{D}_{X}\right\}$, and observe that this is an open subset of $M$. We have already seen in the proof of Theorem 8.5.12 above that all the smooth maps $\Phi_{t}^{X}: M_{t} \rightarrow M$ are injective with $\Phi_{t}^{X}\left(M_{t}\right)=M_{-t}$ and $\left(\Phi_{t}^{X}\right)^{-1}=\Phi_{-t}^{X}$ on the image. It follows in particular that $\Phi_{t}^{X}\left(M_{t}\right)=M_{-t}$ is open and that

$$
\Phi_{t}^{X}: M_{t} \rightarrow M_{-t}
$$

is a diffeomorphism whose inverse is $\Phi_{-t}^{X}$.\\
Proposition 8.5.15 Smooth Dependence Theorem. Let $M$ be a smooth manifold, $V$ a finite-dimensional vector space, $V_{1} \subseteq V$ an open subset, and $\Psi: V_{1} \rightarrow \mathcal{V}(M)$ a map for which the map

$$
\widehat{\Psi}: V_{1} \times M \rightarrow T(M), \quad(v, p) \mapsto \Psi(v)(p)
$$

is smooth (the vector field $\Psi(v)$ depends smoothly on $v$ ). Then there exist for each $\left(p_{0}, v_{0}\right) \in M \times V_{1}$ an open neighborhood $U$ of $p_{0}$ in $M$, an open interval $I \subseteq \mathbb{R}$ containing 0 , an open neighborhood $W$ of $v_{0}$ in $V_{1}$, and a smooth map

$$
\Phi: I \times U \times W \rightarrow M
$$

such that for each $(p, v) \in U \times W$ the curve

$$
\Phi_{p}^{v}: I \rightarrow M, \quad t \mapsto \Phi(t, p, v)
$$

is an integral curve of the vector field $\Psi(v)$ with $\Phi_{p}^{v}(0)=p$.

Proof. The parameters do not cause any additional problems, which can be seen by the following trick: On the product manifold $N:=V_{1} \times M$, we consider the smooth vector field $Y$ given by

$$
Y(v, p):=(0, \Psi(v)(p))
$$

Then the integral curves of $Y$ are of the form

$$
\gamma(t)=\left(v, \gamma_{v}(t)\right)
$$

where $\gamma_{v}$ is an integral curve of the smooth vector field $\Psi(v)$ on $M$. Therefore, the assertion is an immediate consequence on the smoothness of the local flow of $Y$ on $V_{1} \times M$ (Theorem 8.5.12).

\subsubsection*{Lie Derivatives}
We take a closer look at the interaction of local flows and vector fields. It will turn out that this leads to a new concept of a directional derivative which works for general tensor fields.

Let $X \in \mathcal{V}(M)$ and $\Phi^{X}: \mathcal{D}_{X} \rightarrow M$ its maximal local flow. For $f \in C^{\infty}(M)$ and $t \in \mathbb{R}$, we set

$$
\left(\Phi_{t}^{X}\right)^{*} f:=f \circ \Phi_{t}^{X} \in C^{\infty}\left(M_{t}\right) .
$$

Then we find

$$
\lim _{t \rightarrow 0} \frac{1}{t}\left(\left(\Phi_{t}^{X}\right)^{*} f-f\right)=\mathrm{d} f(X)=\mathcal{L}_{X} f \in C^{\infty}(M)
$$

For a second vector field $Y \in \mathcal{V}(M)$, we define a smooth vector field on the open subset $M_{-t} \subseteq M$ by

$$
\left(\Phi_{t}^{X}\right)_{*} Y:=T\left(\Phi_{t}^{X}\right) \circ Y \circ \Phi_{-t}^{X}=T\left(\Phi_{t}^{X}\right) \circ Y \circ\left(\Phi_{t}^{X}\right)^{-1}
$$

(cf. Remark 8.5.14) and define the Lie derivative by

$$
\mathcal{L}_{X} Y:=\lim _{t \rightarrow 0} \frac{1}{t}\left(\left(\Phi_{-t}^{X}\right)_{*} Y-Y\right)=\left.\frac{d}{d t}\right|_{t=0}\left(\Phi_{-t}^{X}\right)_{*} Y
$$

which is defined on all of $M$ since for each $p \in M$ the vector $\left(\left(\Phi_{t}^{X}\right)_{*} Y\right)(p)$ is defined for sufficiently small $t$ and depends smoothly on $t$.

Theorem 8.5.16. $\mathcal{L}_{X} Y=[X, Y]$ for $X, Y \in \mathcal{V}(M)$.\\
Proof. Fix $p \in M$. It suffices to show that $\mathcal{L}_{X} Y$ and $[X, Y]$ coincide at $p$. We may therefore work in a local chart, hence assume that $M=U$ is an open subset of $\mathbb{R}^{n}$.

Identifying vector fields with smooth $\mathbb{R}^{n}$-valued functions, we then have

$$
[X, Y](x)=\mathrm{d} Y(x) X(x)-\mathrm{d} X(x) Y(x), \quad x \in U
$$

On the other hand,

$$
\begin{aligned}
\left(\left(\Phi_{-t}^{X}\right)_{*} Y\right)(x) & =T\left(\Phi_{-t}^{X}\right) \circ Y \circ \Phi_{t}^{X}(x) \\
& =\mathrm{d}\left(\Phi_{-t}^{X}\right)\left(\Phi_{t}^{X}(x)\right) Y\left(\Phi_{t}^{X}(x)\right)=\left(\mathrm{d}\left(\Phi_{t}^{X}\right)(x)\right)^{-1} Y\left(\Phi_{t}^{X}(x)\right) .
\end{aligned}
$$

To calculate the derivative of this expression with respect to $t$, we first observe that it does not matter if we first take derivatives with respect to $t$ and then with respect to $x$, or vice versa. This leads to

$$
\left.\frac{d}{d t}\right|_{t=0} \mathrm{~d}\left(\Phi_{t}^{X}\right)(x)=\mathrm{d}\left(\left.\frac{d}{d t}\right|_{t=0} \Phi_{t}^{X}\right)(x)=\mathrm{d} X(x) .
$$

Next we note that for any smooth curve $\alpha:[-\varepsilon, \varepsilon] \rightarrow \mathrm{GL}_{n}(\mathbb{R})$ with $\alpha(0)=\mathbf{1}$ we have

$$
\left(\alpha^{-1}\right)^{\prime}(t)=-\alpha(t)^{-1} \alpha^{\prime}(t) \alpha(t)^{-1}
$$

and in particular $\left(\alpha^{-1}\right)^{\prime}(0)=-\alpha^{\prime}(0)$. Combining all this, we obtain with the Product Rule

$$
\mathcal{L}_{X}(Y)(x)=-\mathrm{d} X(x) Y(x)+\mathrm{d} Y(x) X(x)=[X, Y](x) .
$$

Corollary 8.5.17. If $X, Y \in \mathcal{V}(M)$ are complete vector fields, then their global flows $\Phi^{X}, \Phi^{Y}: \mathbb{R} \rightarrow \operatorname{Diff}(M)$ commute if and only if $X$ and $Y$ commute, i.e., $[X, Y]=0$.

Proof. (1) Suppose first that $\Phi^{X}$ and $\Phi^{Y}$ commute, i.e.,

$$
\Phi^{X}(t) \circ \Phi^{Y}(s)=\Phi^{Y}(s) \circ \Phi^{X}(t) \quad \text { for } t, s \in \mathbb{R}
$$

Let $p \in M$ and $\gamma_{p}(s):=\Phi_{s}^{Y}(p)$ the $Y$-integral curve through $p$. We then have

$$
\gamma_{p}(s)=\Phi_{s}^{Y}(p)=\Phi_{t}^{X} \circ \Phi_{s}^{Y} \circ \Phi_{-t}^{X}(p),
$$

and passing to the derivative at $s=0$ yields

$$
Y(p)=\gamma_{p}^{\prime}(0)=T\left(\Phi_{t}^{X}\right) Y\left(\Phi_{-t}^{X}(p)\right)=\left(\left(\Phi_{t}^{X}\right)_{*} Y\right)(p) .
$$

Passing now to the derivative at $t=0$, we arrive at $[X, Y]=\mathcal{L}_{X}(Y)=0$.\\
(2) Now we assume $[X, Y]=0$. First, we show that $\left(\Phi_{t}^{X}\right)_{*} Y=Y$ holds for all $t \in \mathbb{R}$. For $t, s \in \mathbb{R}$ we have

$$
\left(\Phi_{t+s}^{X}\right)_{*} Y=\left(\Phi_{t}^{X}\right)_{*}\left(\Phi_{s}^{X}\right)_{*} Y,
$$

so that

$$
\frac{d}{d t}\left(\Phi_{t}^{X}\right)_{*} Y=-\left(\Phi_{t}^{X}\right)_{*} \mathcal{L}_{X}(Y)=0
$$

for each $t \in \mathbb{R}$. Since for each $p \in M$ the curve

$$
\mathbb{R} \rightarrow T_{p}(M), \quad t \mapsto\left(\left(\Phi_{t}^{X}\right)_{*} Y\right)(p)
$$

is smooth and its derivative vanishes, it is constant $Y(p)$. This shows that $\left(\Phi_{t}^{X}\right)_{*} Y=Y$ for each $t \in \mathbb{R}$.

For $\gamma(s):=\Phi_{t}^{X} \Phi_{s}^{Y}(p)$, we now have $\gamma(0)=\Phi_{t}^{X}(p)$ and

$$
\gamma^{\prime}(s)=T\left(\Phi_{t}^{X}\right) \circ Y\left(\Phi_{s}^{Y}(p)\right)=Y\left(\Phi_{t}^{X} \Phi_{s}^{Y}(p)\right)=Y(\gamma(s))
$$

so that $\gamma$ is an integral curve of $Y$. We conclude that $\gamma(s)=\Phi_{s}^{Y}\left(\Phi_{t}^{X}(p)\right)$, and this means that the flows of $X$ and $Y$ commute.

Remark 8.5.18. Let $X, Y \in \mathcal{V}(M)$ be two complete vector fields with $\Phi^{X}$ and $\Phi^{Y}$ as their global flows. We then consider the commutator map

$$
F: \mathbb{R}^{2} \rightarrow \operatorname{Diff}(M), \quad(t, s) \mapsto \Phi_{t}^{X} \circ \Phi_{s}^{Y} \circ \Phi_{-t}^{X} \circ \Phi_{-s}^{Y}
$$

We know from Corollary 8.5.17 that it vanishes if and only if $[X, Y]=0$, but there is also a more direct way from $F$ to the Lie bracket. In fact, we first observe that

$$
\frac{\partial F}{\partial s}(t, 0)=\left(\Phi_{t}^{X}\right)_{*} Y-Y
$$

and hence that

$$
\frac{\partial^{2} F}{\partial t \partial s}(0,0)=[Y, X]
$$

Here we use that if $I \subseteq \mathbb{R}$ is an interval and

$$
\alpha: I \rightarrow \operatorname{Diff}(M) \quad \text { and } \quad \beta: I \rightarrow \operatorname{Diff}(M)
$$

are maps for which\\
$\widehat{\alpha}: I \times M \rightarrow M, \quad(t, x) \mapsto \alpha(t)(x) \quad$ and $\quad \widehat{\beta}: I \times M \rightarrow M, \quad(t, x) \mapsto \beta(t)(x)$\\
are smooth, then the curve $\gamma(t):=\alpha(t) \circ \beta(t)$ also has this property (by the Chain Rule), and if $\alpha(0)=\beta(0)=\mathrm{id}_{M}$, then $\gamma$ satisfies

$$
\gamma^{\prime}(0)=\alpha^{\prime}(0) \circ \beta(0)+T(\alpha(0)) \circ \beta^{\prime}(0)=\alpha^{\prime}(0)+\beta^{\prime}(0)
$$

\subsubsection*{Exercises for Section 8.5}
Exercise 8.5.1. Let $M:=\mathbb{R}^{n}$. For a matrix $A \in M_{n}(\mathbb{R})$, we consider the linear vector field $X_{A}(x):=A x$. Determine the maximal flow $\Phi^{X}$ of this vector field.

Exercise 8.5.2. Let $M$ be a smooth manifold and $Y \in \mathcal{V}(M)$ a smooth vector field on $M$. Suppose that $Y$ generates a local flow which is defined on an entire box of the form $[-\varepsilon, \varepsilon] \times M$. Show that this implies the completeness of $X$.

\subsection*{Submanifolds}
In this section, we describe subsets of a smooth manifold which themselves carry manifold structures. In applications, there occur subsets with various degrees of compatibility of their smooth structures with the ambient manifold. These give rise to different concepts of submanifolds.

Definition 8.6.1 (Submanifolds). (a) Let $M$ be a smooth $n$-dimensional manifold. A subset $S \subseteq M$ is called a $d$-dimensional submanifold if for each $p \in S$ there exists a chart $(\varphi, U)$ of $M$ with $p \in U$ such that

$$
\varphi(U \cap S)=\varphi(U) \cap\left(\mathbb{R}^{d} \times\{0\}\right)
$$

A submanifold of codimension 1, i.e., $\operatorname{dim} S=n-1$, is called a smooth hypersurface.\\
(b) An immersed submanifold of $M$ is a subset $S \subseteq M$, endowed with a smooth manifold structure such that the inclusion map $i_{S}: S \rightarrow M$ is a smooth immersion, i.e., its tangent map $T\left(i_{S}\right)$ is injective on each tangent space of $S$.\\
(c) An initial submanifold of $M$ is an immersed submanifold ${ }^{2}$ such that for each other smooth manifold $N$ a map $f: N \rightarrow S$ is smooth if and only if $i_{S} \circ f: N \rightarrow M$ is smooth. The latter condition means that a map into $S$ is smooth if and only if it is smooth when considered as a map into $M$.

Remark 8.6.2. For any two initial submanifold structures $\iota_{1}: S_{1} \rightarrow M$ and $\iota_{2}: S_{2} \rightarrow M$ on a subset $S=\iota_{1}\left(S_{1}\right)=\iota_{2}\left(S_{2}\right)$ of a smooth manifold $M$, there exists a diffeomorphism $\varphi: S_{2} \rightarrow S_{1}$ with $\iota_{1} \circ \varphi=\iota_{2}$. In particular, the smooth manifold structures on $M$ coincide.

In fact, by the definition of an initial submanifold, the maps

$$
\iota_{1}^{-1} \circ \iota_{2}: S_{2} \rightarrow S_{1} \quad \text { and } \quad \iota_{2}^{-1} \circ \iota_{1}: S_{1} \rightarrow S_{2}
$$

are smooth. Since they are mutually inverse, $\varphi:=\iota_{1}^{-1} \circ \iota_{2}: S_{2} \rightarrow S_{1}$ is a diffeomorphism with the required properties.

We shall see below that submanifolds are initial submanifolds. Later (see Example 9.3.12) we will see an example of an initial submanifold which is not a submanifold. The following example shows that not all immersed submanifolds are initial.

Example 8.6.3 (Figure Eight). Consider the immersed submanifold $S \subseteq \mathbb{R}^{2}$ defined by the immersion $\left.\varphi:\right] 0,2 \pi\left[\rightarrow \mathbb{R}^{2}, t \mapsto\left(\sin ^{3} t, \sin t \cos t\right)\right.$. Define a map $\psi:] 0,2 \pi[\rightarrow S$ via

\footnotetext{${ }^{2}$ Note that the assumption that $M$ be immersed is not redundant (see [KM97], §27.11).
}
$$
\psi(\varphi(t)):= \begin{cases}\varphi(t+\pi) & \text { for } t<\pi \\ 0 & \text { for } t=\pi \\ \varphi(t-\pi) & \text { for } t>\pi\end{cases}
$$

Then we have $\psi(t)=\left(-\sin ^{3} t, \sin t \cos t\right)$, so that $\left.i_{S} \circ \psi:\right] 0,2 \pi\left[\rightarrow \mathbb{R}^{2}\right.$ is smooth. But the limit $\lim _{t \rightarrow \pi} \psi(t)$ does not exist in $S$, so $\psi$ is not even continuous.

Remark 8.6.4. (a) Any discrete subset $S$ of $M$ is a 0 -dimensional submanifold.\\
(b) If $n=\operatorname{dim} M$, any open subset $S \subseteq M$ is an $n$-dimensional submanifold. If, conversely, $S \subseteq M$ is an $n$-dimensional submanifold, then the definition immediately shows that $S$ is open.

Lemma 8.6.5. Any submanifold $S$ of a manifold $M$ has a natural manifold structure, turning it into an initial submanifold.\\
Proof. (a) We endow $S$ with the subspace topology inherited from $M$, which turns it into a Hausdorff space. For each chart ( $\varphi, U$ ) satisfying (8.12), we obtain a $d$-dimensional chart

$$
\left(\left.\varphi\right|_{U \cap S}, U \cap S\right)
$$

of $S$. For two such charts coming from $(\varphi, U)$ and $(\psi, V)$, we have

$$
\left.\psi \circ \varphi^{-1}\right|_{\varphi(U \cap V \cap S)}=\left.\left(\left.\psi\right|_{V \cap S}\right) \circ\left(\left.\varphi\right|_{U \cap S}\right)^{-1}\right|_{\varphi(U \cap V \cap S)},
$$

which is a smooth map onto an open subset of $\mathbb{R}^{d}$. We thus obtain a smooth atlas on $S$.\\
(b) To see that $i_{S}$ is smooth, let $p \in S$ and $(\varphi, U)$ be a chart satisfying (8.12). Then

$$
\varphi \circ i_{S} \circ\left(\left.\varphi\right|_{S \cap U}\right)^{-1}: \varphi(U) \cap\left(\mathbb{R}^{d} \times\{0\}\right) \rightarrow \varphi(U) \subseteq \mathbb{R}^{n}
$$

is the inclusion map, hence smooth. This implies that $i_{S}$ is smooth.\\
(c) If $f: N \rightarrow S$ is smooth, then the composition $i_{S} \circ f$ is smooth (Lemma 8.2.21). Suppose, conversely, that $i_{S} \circ f$ is smooth. Let $p \in N$ and choose a chart ( $\varphi, U$ ) of $M$ satisfying (8.12) with $f(p) \in U$. Then the map

$$
\left.\varphi \circ i_{S} \circ f\right|_{f^{-1}(U)}: f^{-1}(U) \rightarrow \varphi(U) \subseteq \mathbb{R}^{n}
$$

is smooth, but its values lie in

$$
\varphi(U \cap S)=\varphi(U) \cap\left(\mathbb{R}^{d} \times\{0\}\right)
$$

Therefore, $\left.\varphi \circ i_{S} \circ f\right|_{f^{-1}(U)}$ is also smooth as a map into $\mathbb{R}^{d}$, which means that

$$
\left.\left.\varphi\right|_{U \cap S} \circ f\right|_{f^{-1}(U)}: f^{-1}(U) \rightarrow \varphi(U \cap S) \subseteq \mathbb{R}^{d}
$$

is smooth, and hence that $f$ is smooth as a map $N \rightarrow S$.

Remark 8.6.6 (Tangent spaces of submanifolds). From the construction of the manifold structure on $S$, it follows that for each $p \in S$ and each chart ( $\varphi, U$ ) satisfying (8.12), we may identify the tangent space $T_{p}(S)$ with the subspace $T_{p}(\varphi)^{-1}\left(\mathbb{R}^{d}\right)$ mapped by $T_{p}(\varphi)$ onto the subspace $\mathbb{R}^{d}$ of $\mathbb{R}^{n}$.

Lemma 8.6.7. $A$ subset $S$ of a smooth manifold $M$ carries at most one initial submanifold structure, i.e., for any two smooth manifold structures $S_{1}, S_{2}$ on the same set $S$ the map $\varphi=\operatorname{id}_{S}: S_{1} \rightarrow S_{2}$ is a diffeomorphism.

Proof. Since $i_{S_{2}} \circ \varphi=i_{S_{1}}: S_{1} \rightarrow M$ is smooth, the map $\varphi$ is smooth. Likewise, we see that the inverse map $\varphi^{-1}: S_{2} \rightarrow S_{1}$ is smooth, showing that $\varphi$ is a diffeomorphism $S_{1} \rightarrow S_{2}$.

Definition 8.6.8. Let $f: M \rightarrow N$ be a smooth map. We call $n \in N$ a regular value of $f$ if for each $x \in f^{-1}(n)$ the tangent map $T_{x}(f): T_{x}(M) \rightarrow T_{n}(N)$ is surjective. Otherwise $n$ is called a singular value of $f$. Note that, in particular, each $n \in N \backslash f(M)$ is a regular value.

We are now ready to prove a manifold version of the fact that inverse images of regular values are submanifolds.

Theorem 8.6.9 (Regular Value Theorem-Global Version). Let $M$ and $N$ be smooth manifolds of dimensions $m$, resp., $n$, and $f: M \rightarrow N$ a smooth map. If $y \in N$ is a regular value of $f$, then $S:=f^{-1}(y)$ is a submanifold of $M$ of dimension $(m-n)$.

Proof. We may assume that $y \in f(M)$ and note that then $d:=m-n \geq 0$ because $T_{x}(f): T_{x}(M) \cong \mathbb{R}^{m} \rightarrow \mathbb{R}^{n} \cong T_{f(x)}(N)$ is surjective for some $x \in M$.

Let $p \in S$ and choose charts $(\varphi, U)$ of $M$ with $p \in U$ and $(\psi, V)$ of $N$ with $f(p) \in V$. Then the map

$$
F:=\psi \circ f \circ \varphi^{-1}: \varphi\left(U \cap f^{-1}(V)\right) \rightarrow \mathbb{R}^{n}
$$

is a smooth map, and for each $x \in F^{-1}(\psi(y))=\varphi(S \cap U)$ the linear map

$$
T_{x}(F)=T_{f \circ \varphi^{-1}(x)}(\psi) \circ T_{\varphi^{-1}(x)}(f) \circ T_{x}\left(\varphi^{-1}\right): \mathbb{R}^{m} \rightarrow \mathbb{R}^{n}
$$

is surjective. Therefore, Proposition 8.2.9 implies the existence of an open subset $U^{\prime} \subseteq \varphi\left(U \cap f^{-1}(V)\right)$ containing $\varphi(p)$ and a diffeomorphism $\gamma: U^{\prime} \rightarrow \gamma\left(U^{\prime}\right) \subseteq \mathbb{R}^{m}$ with

$$
\gamma\left(U^{\prime} \cap \varphi(U \cap S)\right)=\left(\mathbb{R}^{d} \times\{0\}\right) \cap \gamma\left(U^{\prime}\right) .
$$

Then $\left(\gamma \circ \varphi, \varphi^{-1}\left(U^{\prime}\right)\right)$ is a chart of $M$ with

$$
(\gamma \circ \varphi)\left(S \cap \varphi^{-1}\left(U^{\prime}\right)\right)=\gamma\left(\varphi(S \cap U) \cap U^{\prime}\right)=\left(\mathbb{R}^{d} \times\{0\}\right) \cap \gamma\left(U^{\prime}\right) .
$$

This shows that $S$ is a $d$-dimensional submanifold of $M$.

Remark 8.6.10. If $S \subseteq M$ is a submanifold, then we may identify the tangent spaces $T_{p}(S)$ with the subspaces $\operatorname{im}\left(T_{p}\left(i_{S}\right)\right)$ of $T_{p}(M)$, where $i_{S}: S \rightarrow M$ is the smooth inclusion map (cf. Remark 8.6.6). If, in addition, $S=f^{-1}(y)$ for some regular value $y$ of the smooth map $f: M \rightarrow N$, then we have

$$
T_{p}(S)=\operatorname{ker} T_{p}(f) \quad \text { for } \quad p \in S
$$

To verify this relation, we recall that we know already that

$$
\operatorname{dim} S=n-m=\operatorname{dim} T_{p}(S)
$$

On the other hand, $f \circ i_{S}=y: S \rightarrow N$ is the constant map, so that $T_{p}\left(f \circ i_{S}\right)= T_{p}(f) \circ T_{p}\left(i_{S}\right)=0$, which leads to $T_{p}(S) \subseteq \operatorname{ker} T_{p}(f)$. Since $T_{p}(f)$ is surjective by assumption, equality follows by comparing dimensions.

The following theorem which we mention without proof implies in particular the existence of regular values for surjective smooth maps (cf. [La99, Thm. XVI.1.4]).

Theorem 8.6.11 (Sard's Theorem). Let $M_{1}$ and $M_{2}$ be smooth second countable manifolds, $f: M_{1} \rightarrow M_{2}$ a smooth map, and $M_{1}^{c}$ the set of critical points of $f$. Then $f\left(M_{1}^{c}\right)$ is a set of measure zero in $M_{2}$, i.e., for each chart ( $\varphi, U$ ) of $M_{2}$ the set $\varphi\left(U \cap f\left(M_{1}^{c}\right)\right.$ ) is of Lebesgue measure zero.

Here we recall that a topological space $X$ is said to be second countable if its topology has a countable basis, i.e., there exists a countable family $\left(O_{n}\right)_{n \in \mathbb{N}}$ of open subsets such that any open subset is the union of some of the $O_{n}$.

\subsubsection*{Notes on Chapter 8}
The notion of a smooth manifold is more subtle than one may think on the surface. One of these subtleties arises from the fact that a topological space may carry different smooth manifold structures which are not diffeomorphic. Important examples of low dimension are the 7 -sphere $\mathbb{S}^{7}$ and $\mathbb{R}^{4}$. Actually, 4 is the only dimension $n$ for which $\mathbb{R}^{n}$ carries two nondiffeomorphic smooth structures. At this point, it is instructive to observe that two smooth structures might be diffeomorphic without having the same maximal atlas: The charts ( $\varphi, \mathbb{R}$ ) and ( $\psi, \mathbb{R}$ ) on $\mathbb{R}$ given by

$$
\varphi(x)=x \quad \text { and } \quad \psi(x)=x^{3}
$$

define two different smooth manifold structures $\mathbb{R}_{\varphi}$ and $\mathbb{R}_{\psi}$, but the map

$$
\gamma: \mathbb{R}_{\psi} \rightarrow \mathbb{R}_{\varphi}, \quad x \mapsto x^{3}
$$

is a diffeomorphism.

Later we shall see that there are also purely topological subtleties due to the fact that the topology might be "too large". The regularity assumption which is needed in many situations is the paracompactness of the underlying Hausdorff space.

Vector fields and their zeros play an important role in the topology of manifolds. To each manifold $M$ we associate the maximal number $\alpha(M)=k$ for which there exist smooth vector fields $X_{1}, \ldots, X_{k} \in \mathcal{V}(M)$ which are linearly independent at each point of $M$. A manifold is called parallelizable if $\alpha(M)=\operatorname{dim} M$ (which is the maximal value). Clearly, $\alpha\left(\mathbb{R}^{n}\right)=n$, so that $\mathbb{R}^{n}$ is parallelizable, but it is a deep theorem that the $n$-sphere $\mathbb{S}^{n}$ is only parallelizable if $n=0,1,3$, or 7 . This in turn has important applications on the existence of real division algebras, namely that they only exist in dimensions $1,2,4$, or 8 (this is the famous 1-2-4-8-Theorem). Another important result in topology is $\alpha\left(\mathbb{S}^{2}\right)=0$, i.e., each vector field on the 2 -sphere has a zero (Hairy Ball Theorem).

As we have seen in Theorem 8.5.16, the Lie bracket on the space $\mathcal{V}(M)$ of vector fields is closely related to the commutator in the group $\operatorname{Diff}(M)$ of diffeomorphisms of $M$. This fact is part of a more general correspondence in the theory of Lie groups which associates to each Lie group $G$ a Lie algebra $\mathbf{L}(G)$ given by a suitable bracket on the tangent space $T_{\mathbf{1}}(G)$ which is defined in terms of the Lie bracket of vector fields.

\section*{Basic Lie Theory}
This chapter is devoted to the subject proper of this book: Lie groups, defined as smooth manifolds with a group structure such that all structure maps (multiplication and inversion) are smooth. Here we use vector fields to build the key tools of Lie theory. The Lie functor which associates a Lie algebra with a Lie group and the exponential function from the Lie algebra to the Lie group. They provide the means to translate global problems into infinitesimal ones and to lift infinitesimal solutions to local ones. Passing from the local to the global level usually requires tools from covering theory, resp., topology. In the process, we introduce smooth group actions and the adjoint representation, and provide a number of topological facts about Lie groups.

Further, we prove the Baker-Campbell-Dynkin-Hausdorff (BCDH) formula which expresses the group multiplication locally in terms of a universal power series involving Lie brackets. As a first set of applications of the translation mechanisms, we identify the Lie group structures of closed subgroups of Lie groups and show how to construct Lie groups from local and infinitesimal data. The key result is the Integral Subgroup Theorem 9.4.8 describing the subgroup generated by the exponential image of a Lie subalgebra as a Lie group. Combined with Ado's Theorem 7.4.1, it yields Lie's Third Theorem 9.4.11 saying that each finite-dimensional real Lie algebra is the Lie algebra of a Lie group. Then we systematically study coverings of Lie groups, thus providing the means to classify Lie groups with a given Lie algebra. Finally, we prove Yamabe's Theorem 9.6.1 asserting that any arcwise connected subgroup of a Lie group carries a natural Lie group structure and allows equipping any subgroup of a Lie group with a Lie group structure.

\subsection*{Lie Groups and Their Lie Algebras}
In the context of smooth manifolds, the natural class of groups are those endowed with a manifold structure compatible with the group structure. Such groups will be called Lie groups.

\subsubsection*{Lie Groups, First Examples, and the Tangent Group}
Definition 9.1.1. A Lie group is a group $G$ endowed with the structure of a smooth manifold such that the group operations

$$
m_{G}: G \times G \rightarrow G, \quad(x, y) \mapsto x y \quad \text { and } \quad \iota_{G}: G \rightarrow G, \quad x \mapsto x^{-1}
$$

are smooth.\\
Throughout this section, $G$ denotes a Lie group with multiplication map $m_{G}: G \times G \rightarrow G,(x, y) \mapsto x y$, inversion map $\iota_{G}: G \rightarrow G, x \mapsto x^{-1}$, and neutral element 1. For $g \in G$, we write $\lambda_{g}: G \rightarrow G, x \mapsto g x$ for the left multiplication map, $\rho_{g}: G \rightarrow G, x \mapsto x g$ for the right multiplication map, and $c_{g}: G \rightarrow G, x \mapsto g x g^{-1}$ for the conjugation with $g$. A morphism of Lie groups is a smooth homomorphism of Lie groups $\varphi: G_{1} \rightarrow G_{2}$.

Remark 9.1.2. All maps $\lambda_{g}, \rho_{g}$, and $c_{g}$ are smooth. Moreover, they are bijective with $\lambda_{g^{-1}}=\lambda_{g}^{-1}, \rho_{g^{-1}}=\rho_{g}^{-1}$, and $c_{g^{-1}}=c_{g}^{-1}$, so that they are diffeomorphisms of $G$ onto itself.

In addition, the maps $c_{g}$ are automorphisms of $G$, so that we obtain a group homomorphism

$$
C: G \rightarrow \operatorname{Aut}(G), \quad g \mapsto c_{g},
$$

where $\operatorname{Aut}(G)$ stands for the group of automorphisms of the Lie group $G$, i.e., the group automorphisms which are diffeomorphisms. The automorphisms of the form $c_{g}$ are called inner automorphisms of $G$. The group of inner automorphisms of $G$ is denoted by $\operatorname{Inn}(G)$.

One can show that the requirement of $\iota_{G}$ being smooth is redundant (Exercise 9.1.4).

Example 9.1.3. We consider the additive group $G:=\left(\mathbb{R}^{n},+\right)$, endowed with the natural $n$-dimensional manifold structure. A corresponding chart is given by $\left(\mathrm{id}_{\mathbb{R}^{n}}, \mathbb{R}^{n}\right)$, which shows that the corresponding product manifold structure on $\mathbb{R}^{n} \times \mathbb{R}^{n}$ is given by the chart $\left(\mathrm{id}_{\mathbb{R}^{n}} \times \mathrm{id}_{\mathbb{R}^{n}}, \mathbb{R}^{n} \times \mathbb{R}^{n}\right)=\left(\mathrm{id}_{\mathbb{R}^{2 n}}, \mathbb{R}^{2 n}\right)$, hence coincides with the natural manifold structure on $\mathbb{R}^{2 n}$. Therefore, the smoothness of addition and inversion follows from the smoothness of the maps

$$
\mathbb{R}^{2 n} \rightarrow \mathbb{R}^{n}, \quad(x, y) \mapsto x+y \quad \text { and } \quad \mathbb{R}^{n} \rightarrow \mathbb{R}^{n}, \quad x \mapsto-x .
$$

Example 9.1.4. Next we consider the group $G:=\mathrm{GL}_{n}(\mathbb{R})$ of invertible $(n \times n)$-matrices. If $\operatorname{det}: M_{n}(\mathbb{R}) \rightarrow \mathbb{R}$ denotes the determinant function

$$
\operatorname{det}(A)=\sum_{\sigma \in S_{n}} \operatorname{sgn}(\sigma) a_{1, \sigma(1)} \cdots a_{n, \sigma(n)}
$$

then det is a polynomial, hence in particular continuous, and therefore $\mathrm{GL}_{n}(\mathbb{R})=\operatorname{det}^{-1}\left(\mathbb{R}^{\times}\right)$is an open subset of $M_{n}(\mathbb{R}) \cong \mathbb{R}^{n^{2}}$. Hence $G$ carries a natural manifold structure.

We claim that $G$ is a Lie group. The smoothness of the multiplication map follows directly from the smoothness of the bilinear multiplication map

$$
M_{n}(\mathbb{R}) \times M_{n}(\mathbb{R}) \rightarrow M_{n}(\mathbb{R}), \quad(A, B) \mapsto\left(\sum_{k=1}^{n} a_{i k} b_{k j}\right)_{i, j=1, \ldots, n}
$$

which is given in each component by a polynomial function in the $2 n^{2}$ variables $a_{i j}$ and $b_{i j}$ (cf. Exercise 8.2.11 on the smoothness of bilinear maps).

The smoothness of the inversion map follows from Cramer's Rule

$$
g^{-1}=\frac{1}{\operatorname{det} g}\left(b_{i j}\right), \quad b_{i j}=(-1)^{i+j} \operatorname{det}\left(G_{j i}\right)
$$

where $G_{i j} \in M_{n-1}(\mathbb{R})$ is the matrix obtained by erasing the $i$ th row and the $j$ th column in $g$.

Example 9.1.5. (a) (The circle group) We have already seen how to endow the circle

$$
\mathbb{S}^{1}:=\left\{(x, y) \in \mathbb{R}^{2}: x^{2}+y^{2}=1\right\}
$$

with a manifold structure (Example 8.2.5). Identifying it with the unit circle

$$
\mathbb{T}:=\{z \in \mathbb{C}:|z|=1\}
$$

in $\mathbb{C}$, it also inherits a group structure, given by

$$
(x, y) \cdot\left(x^{\prime}, y^{\prime}\right):=\left(x x^{\prime}-y y^{\prime}, x y^{\prime}+x^{\prime} y\right) \quad \text { and } \quad(x, y)^{-1}=(x,-y) .
$$

With these explicit formulas, it is easy to verify that $\mathbb{T}$ is a Lie group (Exercise 9.1.1).\\
(b) (The $n$-dimensional torus) In view of (a), we have a natural manifold structure on the $n$-dimensional torus $\mathbb{T}^{n}:=\left(\mathbb{S}^{1}\right)^{n}$. The corresponding direct product group structure

$$
\left(t_{1}, \ldots, t_{n}\right)\left(s_{1}, \ldots, s_{n}\right):=\left(t_{1} s_{1}, \ldots, t_{n} s_{n}\right)
$$

turns $\mathbb{T}^{n}$ into a Lie group (Exercise 9.1.2).\\
Lemma 9.1.6. (a) As usual, we identify $T(G \times G)$ with $T(G) \times T(G)$. Then the tangent map\\
$T\left(m_{G}\right): T(G \times G) \cong T(G) \times T(G) \rightarrow T(G), \quad(v, w) \mapsto v \cdot w:=T m_{G}(v, w)$\\
defines a Lie group structure on $T(G)$ with identity element $0_{\mathbf{1}} \in T_{\mathbf{1}}(G)$ and inversion $T\left(\iota_{G}\right)$. The canonical projection $\pi_{T(G)}: T(G) \rightarrow G$ is a morphism\\
of Lie groups with kernel ( $T_{\mathbf{1}}(G),+$ ) and the zero section $\sigma: G \rightarrow T(G), g \mapsto 0_{g} \in T_{g}(G)$ is a homomorphism of Lie groups with $\pi_{T(G)} \circ \sigma=\operatorname{id}_{G}$.\\
(b) The map

$$
\Phi: G \times T_{\mathbf{1}}(G) \rightarrow T(G), \quad(g, x) \mapsto g \cdot x:=0_{g} \cdot x=T\left(\lambda_{g}\right) x
$$

is a diffeomorphism.\\
Proof. (a) Since the multiplication map $m_{G}: G \times G \rightarrow G$ is smooth, the same holds for its tangent map

$$
T m_{G}: T(G \times G) \cong T(G) \times T(G) \rightarrow T(G)
$$

Let $\varepsilon_{G}: G \rightarrow G, g \mapsto \mathbf{1}$ be the constant homomorphism. Then the group axioms for $G$ are encoded in the relations\\
(1) $m_{G} \circ\left(m_{G} \times \mathrm{id}_{G}\right)=m_{G} \circ\left(\mathrm{id}_{G} \times m_{G}\right)$ (associativity),\\
(2) $m_{G} \circ\left(\iota_{G}, \mathrm{id}_{G}\right)=m_{G} \circ\left(\mathrm{id}_{G}, \iota_{G}\right)=\varepsilon_{G}$ (inversion), and\\
(3) $m_{G} \circ\left(\varepsilon_{G}, \mathrm{id}_{G}\right)=m_{G} \circ\left(\mathrm{id}_{G}, \varepsilon_{G}\right)=\mathrm{id}_{G}$ (unit element).

Using the functoriality (cf. Lemma 8.3.5) of $T$ and its compatibility with products, we see that these properties carry over to the corresponding maps on $T(G)$ :\\
(1) $T\left(m_{G}\right) \circ T\left(m_{G} \times \operatorname{id}_{G}\right)=T\left(m_{G}\right) \circ\left(T\left(m_{G}\right) \times \operatorname{id}_{T(G)}\right) =T\left(m_{G}\right) \circ\left(\mathrm{id}_{T(G)} \times T\left(m_{G}\right)\right)$ (associativity),\\
(2) $T\left(m_{G}\right) \circ\left(T\left(\iota_{G}\right), \mathrm{id}_{T(G)}\right)=T\left(m_{G}\right) \circ\left(\mathrm{id}_{T(G)}, T\left(\iota_{G}\right)\right)=T\left(\varepsilon_{G}\right)$ (inversion), and\\
(3) $T\left(m_{G}\right) \circ\left(T\left(\varepsilon_{G}\right), \mathrm{id}_{T(G)}\right)=T\left(m_{G}\right) \circ\left(\mathrm{id}_{T(G)}, T\left(\varepsilon_{G}\right)\right)=\mathrm{id}_{T(G)}$ (unit element).

Here we only have to observe that the tangent map $T\left(\varepsilon_{G}\right)$ maps each $v \in T(G)$ to $0_{\mathbf{1}} \in T_{\mathbf{1}}(G)$, which is the neutral element of $T(G)$. We conclude that $T(G)$ is a Lie group with multiplication $T\left(m_{G}\right)$, inversion $T\left(\iota_{G}\right)$, and unit element $0_{1} \in T_{1}(G)$.

The definition of the tangent map implies that the zero section $\sigma: G \rightarrow T(G)$ satisfies

$$
T m_{G} \circ(\sigma \times \sigma)=\sigma \circ m_{G}, \quad T m_{G}\left(0_{g}, 0_{h}\right)=0_{m_{G}(g, h)}=0_{g h}
$$

which means that it is a morphism of Lie groups. That $\pi_{T(G)}$ also is a morphism of Lie groups follows likewise from the relation

$$
\pi_{T(G)} \circ T m_{G}=m_{G} \circ\left(\pi_{T(G)} \times \pi_{T(G)}\right)
$$

which also is an immediate consequence of the definition of the tangent map $T_{m_{G}}$ : it maps $T_{g}(G) \times T_{h}(G)$ into $T_{g h}(G)$.

For $v \in T_{g}(G)$ and $w \in T_{h}(G)$ the linearity of $T_{(g, h)}\left(m_{G}\right)$ implies that

$$
\begin{aligned}
T m_{G}(v, w) & =T_{(g, h)}\left(m_{G}\right)(v, w)=T_{(g, h)}\left(m_{G}\right)(v, 0)+T_{(g, h)}\left(m_{G}\right)(0, w) \\
& =T_{g}\left(\rho_{h}\right) v+T_{h}\left(\lambda_{g}\right) w
\end{aligned}
$$

and in particular $T_{(\mathbf{1}, \mathbf{1})}\left(m_{G}\right)(v, w)=v+w$, so that the multiplication on the normal subgroup $\operatorname{ker} \pi_{T(G)}=T_{\mathbf{1}}(G)$ is simply given by addition.\\
(b) The smoothness of $\Phi$ follows from the smoothness of the multiplication of $T(G)$ and the smoothness of the zero section $\sigma: G \rightarrow T(G), g \mapsto 0_{g}$. That $\Phi$ is a diffeomorphism follows from the following explicit formula for its inverse: $\Phi^{-1}(v)=\left(\pi_{T(G)}(v), \pi_{T(G)}(v)^{-1} . v\right)$, so that its smoothness follows from the smoothness of $\pi_{T(G)}$ (its first component), and the smoothness of the multiplication on $T(G)$. \(\square\)

\subsubsection*{The Lie Functor}
The Lie functor assigns a Lie algebra to each Lie group and a Lie algebra homomorphism to each morphism of Lie groups. It is the key tool to translate Lie group problems into problems in linear algebra.

Definition 9.1.7 (The Lie algebra of $G$ ). A vector field $X \in \mathcal{V}(G)$ is called left invariant if

$$
X=\left(\lambda_{g}\right)_{*} X:=T\left(\lambda_{g}\right) \circ X \circ \lambda_{g}^{-1}
$$

holds for each $g \in G$, i.e., $\left(\lambda_{g}\right)_{*} X=X$. We write $\mathcal{V}(G)^{l}$ for the set of left invariant vector fields in $\mathcal{V}(G)$. Clearly, $\mathcal{V}(G)^{l}$ is a linear subspace of $\mathcal{V}(G)$.

Writing the left invariance as $X=T\left(\lambda_{g}\right) \circ X \circ \lambda_{g}^{-1}$, we see that it means that $X$ is $\lambda_{g}$-related to itself (cf. Exercise 9.1.6). Therefore, the Related Vector Field Lemma 8.4.9 implies that if $X$ and $Y$ are left-invariant, their Lie bracket $[X, Y]$ is also $\lambda_{g}$-related to itself for each $g \in G$, hence left invariant. We conclude that the vector space $\mathcal{V}(G)^{l}$ is a Lie subalgebra of ( $\mathcal{V}(G),[\cdot, \cdot]$ ).

Next we observe that the left invariance of a vector field $X$ implies that for each $g \in G$ we have $X(g)=g . X(\mathbf{1})$ (Lemma 9.1.6(b)), so that $X$ is completely determined by its value $X(\mathbf{1}) \in T_{\mathbf{1}}(G)$. Conversely, for each $x \in T_{\mathbf{1}}(G)$, we obtain a left invariant vector field $x_{l} \in \mathcal{V}(G)^{l}$ with $x_{l}(\mathbf{1})=x$ by $x_{l}(g):=g . x$. That this vector field is indeed left invariant follows from

$$
x_{l} \circ \lambda_{h}(g)=x_{l}(h g)=(h g) \cdot x=h \cdot(g \cdot x)=T\left(\lambda_{h}\right) x_{l}(g)
$$

for all $h, g \in G$. Hence

$$
T_{\mathbf{1}}(G) \rightarrow \mathcal{V}(G)^{l}, \quad x \mapsto x_{l}
$$

is a linear bijection. We thus obtain a Lie bracket $[\cdot, \cdot]$ on $T_{\mathbf{1}}(G)$ satisfying

$$
[x, y]_{l}=\left[x_{l}, y_{l}\right] \quad \text { for all } \quad x, y \in T_{\mathbf{1}}(G) .
$$

The Lie algebra

$$
\mathbf{L}(G):=\left(T_{\mathbf{1}}(G),[\cdot, \cdot]\right) \cong \mathcal{V}(G)^{l}
$$

is called the Lie algebra of $G$.

Proposition 9.1.8 (Functoriality of the Lie algebra). If $\varphi: G \rightarrow H$ is a morphism of Lie groups, then the tangent map

$$
\mathbf{L}(\varphi):=T_{\mathbf{1}}(\varphi): \mathbf{L}(G) \rightarrow \mathbf{L}(H)
$$

is a homomorphism of Lie algebras.\\
Proof. Let $x, y \in \mathbf{L}(G)$ and let $x_{l}, y_{l}$ be the corresponding left invariant vector fields. Then $\varphi \circ \lambda_{g}=\lambda_{\varphi(g)} \circ \varphi$ for each $g \in G$ implies that

$$
T(\varphi) \circ T\left(\lambda_{g}\right)=T\left(\lambda_{\varphi(g)}\right) \circ T(\varphi)
$$

and applying this relation to $x, y \in T_{\mathbf{1}}(G)$, we get

$$
T \varphi \circ x_{l}=(\mathbf{L}(\varphi) x)_{l} \circ \varphi \quad \text { and } \quad T \varphi \circ y_{l}=(\mathbf{L}(\varphi) y)_{l} \circ \varphi
$$

i.e., $x_{l}$ is $\varphi$-related to $\left(\mathbf{L}(\varphi) x_{l}\right.$ and $y_{l}$ is $\varphi$-related to $(\mathbf{L}(\varphi) y)_{l}$. Therefore, the Related Vector Field Lemma implies that

$$
T \varphi \circ\left[x_{l}, y_{l}\right]=\left[(\mathbf{L}(\varphi) x)_{l},(\mathbf{L}(\varphi) y)_{l}\right] \circ \varphi
$$

Evaluating at $\mathbf{1}$, we obtain $\mathbf{L}(\varphi)[x, y]=[\mathbf{L}(\varphi)(x), \mathbf{L}(\varphi)(y)]$, showing that $\mathbf{L}(\varphi)$ is a homomorphism of Lie algebras.

Remark 9.1.9. We obviously have $\mathbf{L}\left(\mathrm{id}_{G}\right)=\mathrm{id}_{\mathbf{L}(G)}$, and for two morphisms $\varphi_{1}: G_{1} \rightarrow G_{2}$ and $\varphi_{2}: G_{2} \rightarrow G_{3}$ of Lie groups, we obtain

$$
\mathbf{L}\left(\varphi_{2} \circ \varphi_{1}\right)=\mathbf{L}\left(\varphi_{2}\right) \circ \mathbf{L}\left(\varphi_{1}\right)
$$

from the Chain Rule:

$$
T_{\mathbf{1}}\left(\varphi_{2} \circ \varphi_{1}\right)=T_{\varphi_{\mathbf{1}}(\mathbf{1})}\left(\varphi_{2}\right) \circ T_{\mathbf{1}}\left(\varphi_{1}\right)=T_{\mathbf{1}}\left(\varphi_{2}\right) \circ T_{\mathbf{1}}\left(\varphi_{1}\right)
$$

The preceding lemma implies that the assignments $G \mapsto \mathbf{L}(G)$ and $\varphi \mapsto \mathbf{L}(\varphi)$ define a functor, called the Lie functor,

$$
\mathbf{L}: \underline{\text { LieGrp }} \rightarrow \underline{\text { LieAlg }}
$$

from the category LieGrp of Lie groups to the category LieAlg of (finitedimensional) Lie algebras.

Corollary 9.1.10. For each isomorphism of Lie groups $\varphi: G \rightarrow H$, the map $\mathbf{L}(\varphi)$ is an isomorphism of Lie algebras, and for each $x \in \mathbf{L}(G)$, the following equation holds

$$
\varphi_{*} x_{l}:=T(\varphi) \circ x_{l} \circ \varphi^{-1}=\left(\mathbf{L}(\varphi) x_{l} .\right.
$$

Proof. Let $\psi: H \rightarrow G$ be the inverse of $\varphi$. Then $\varphi \circ \psi=\operatorname{id}_{H}$ and $\psi \circ \varphi=\operatorname{id}_{G}$ lead to $\mathbf{L}(\varphi) \circ \mathbf{L}(\psi)=\operatorname{id}_{\mathbf{L}(H)}$ and $\mathbf{L}(\psi) \circ \mathbf{L}(\varphi)=\operatorname{id}_{\mathbf{L}(G)}$ (Remark 9.1.9). Further (9.3) follows from (9.2) in the proof of Proposition 9.1.8.

\subsubsection*{Smooth Actions of Lie Groups}
We already encountered smooth flows on manifolds in Chapter 8. These can be viewed as actions of the one-dimensional Lie group ( $\mathbb{R},+$ ). In particular, we have seen that these actions are in one-to-one correspondence with complete vector fields, which is the corresponding Lie algebra picture. Now we describe the corresponding concept for general Lie groups.

Definition 9.1.11. Let $M$ be a smooth manifold and $G$ a Lie group. A (smooth) action of $G$ on $M$ is a smooth map

$$
\sigma: G \times M \rightarrow M
$$

with the following properties:\\
(A1) $\sigma(\mathbf{1}, m)=m$ for all $m \in M$.\\
(A2) $\sigma\left(g_{1}, \sigma\left(g_{2}, m\right)\right)=\sigma\left(g_{1} g_{2}, m\right)$ for $g_{1}, g_{2} \in G$ and $m \in M$.\\
We also write

$$
g . m:=\sigma(g, m), \quad \sigma_{g}(m):=\sigma(g, m), \quad \sigma^{m}(g):=\sigma(g, m)=g . m .
$$

The map $\sigma^{m}$ is called the orbit map.\\
For each smooth action $\sigma$, the map

$$
\widehat{\sigma}: G \rightarrow \operatorname{Diff}(M), \quad g \mapsto \sigma_{g}
$$

is a group homomorphism, and any homomorphism $\gamma: G \rightarrow \operatorname{Diff}(M)$ for which the map

$$
\sigma_{\gamma}: G \times M \rightarrow M, \quad(g, m) \mapsto \gamma(g)(m)
$$

is smooth defines a smooth action of $G$ on $M$.\\
Remark 9.1.12. What we call an action is sometimes called a left action. Likewise one defines a right action as a smooth map $\sigma_{R}: M \times G \rightarrow M$ with

$$
\sigma_{R}(m, \mathbf{1})=m, \quad \sigma_{R}\left(\sigma_{R}\left(m, g_{1}\right), g_{2}\right)=\sigma_{R}\left(m, g_{1} g_{2}\right) .
$$

For $m . g:=\sigma_{R}(m, g)$, this takes the form

$$
m \cdot\left(g_{1} g_{2}\right)=\left(m \cdot g_{1}\right) \cdot g_{2}
$$

of an associativity condition.\\
If $\sigma_{R}$ is a smooth right action of $G$ on $M$, then

$$
\sigma_{L}(g, m):=\sigma_{R}\left(m, g^{-1}\right)
$$

defines a smooth left action of $G$ on $M$. Conversely, if $\sigma_{L}$ is a smooth left action, then

$$
\sigma_{R}(m, g):=\sigma_{L}\left(g^{-1}, m\right)
$$

defines a smooth left action. This translation is one-to-one, so that we may freely pass from one type of action to the other.

Examples 9.1.13. (a) If $X \in \mathcal{V}(M)$ is a complete vector field (cf. Definition 8.5.7) and $\Phi: \mathbb{R} \times M \rightarrow M$ its global flow, then $\Phi$ defines a smooth action of $G=(\mathbb{R},+)$ on $M$.\\
(b) If $G$ is a Lie group, then the multiplication map $\sigma:=m_{G}: G \times G \rightarrow G$ defines a smooth left action of $G$ on itself. In this case, the $\left(m_{G}\right)_{g}=\lambda_{g}$ are the left multiplications.

The multiplication map also defines a smooth right action of $G$ on itself. The corresponding left action is

$$
\sigma: G \times G \rightarrow G, \quad(g, h) \mapsto h g^{-1} \quad \text { with } \quad \sigma_{g}=\rho_{g}^{-1} .
$$

There is a third action of $G$ on itself, the conjugation action:

$$
\sigma: G \times G \rightarrow G, \quad(g, h) \mapsto g h g^{-1} \quad \text { with } \quad \sigma_{g}=c_{g} .
$$

(c) We have a natural smooth action of the Lie group $\mathrm{GL}_{n}(\mathbb{R})$ on $\mathbb{R}^{n}$ :

$$
\sigma: \mathrm{GL}_{n}(\mathbb{R}) \times \mathbb{R}^{n} \rightarrow \mathbb{R}^{n}, \quad \sigma(g, x):=g x
$$

We further have an action of $\mathrm{GL}_{n}(\mathbb{R})$ on $M_{n}(\mathbb{R})$ :

$$
\sigma: \mathrm{GL}_{n}(\mathbb{R}) \times M_{n}(\mathbb{R}) \rightarrow M_{n}(\mathbb{R}), \quad \sigma(g, A)=g A g^{-1}
$$

(d) On the set $M_{p, q}(\mathbb{R})$ of $(p \times q)$-matrices we have an action of the direct product Lie group $G:=\mathrm{GL}_{p}(\mathbb{R}) \times \mathrm{GL}_{q}(\mathbb{R})$ by $\sigma((g, h), A):=g A h^{-1}$.

The following proposition generalizes the passage from flows of vector fields to actions of general Lie groups.

Proposition 9.1.14. Let $G$ be a Lie group and $\sigma: G \times M \rightarrow M$ a smooth action of $G$ on $M$. Then the assignment

$$
\dot{\sigma}: \mathbf{L}(G) \rightarrow \mathcal{V}(M), \quad \dot{\sigma}(x)(m):=\mathbf{L}(\sigma)(x)(m):=-T_{\mathbf{1}}\left(\sigma^{m}\right)(x)
$$

is a homomorphism of Lie algebras.\\
Proof. First, we observe that for each $x \in \mathbf{L}(G)$ the map $\mathbf{L}(\sigma)(x)$ defines a smooth map $M \rightarrow T(M)$, and since $\mathbf{L}(\sigma)(x)(m) \in T_{\sigma(\mathbf{1}, m)}(M)=T_{m}(M)$, it is a smooth vector field on $M$.

To see that $\dot{\sigma}$ is a homomorphism of Lie algebras, we pick $m \in M$ and write

$$
\varphi^{m}:=\sigma^{m} \circ \iota_{G}: G \rightarrow M, \quad g \mapsto g^{-1} . m
$$

for the reversed orbit map. Then

$$
\varphi^{m}(g h)=(g h)^{-1} \cdot m=h^{-1} \cdot\left(g^{-1} \cdot m\right)=\varphi^{g^{-1} \cdot m}(h)
$$

which can be written as

$$
\varphi^{m} \circ \lambda_{g}=\varphi^{g^{-1} \cdot m}
$$

Taking the differential at $\mathbf{1} \in G$, we obtain for each $x \in \mathbf{L}(G)=T_{\mathbf{1}}(G)$ :

$$
\begin{aligned}
T_{g}\left(\varphi^{m}\right) x_{l}(g) & =T_{g}\left(\varphi^{m}\right) T_{\mathbf{1}}\left(\lambda_{g}\right) x=T_{\mathbf{1}}\left(\varphi^{m} \circ \lambda_{g}\right) x=T_{\mathbf{1}}\left(\varphi^{g^{-1} . m}\right) x \\
& =T_{\mathbf{1}}\left(\sigma^{g^{-1} . m}\right) T_{\mathbf{1}}\left(\iota_{G}\right) x=-T_{\mathbf{1}}\left(\sigma^{\varphi^{m}(g)}\right) x=\dot{\sigma}(x)\left(\varphi^{m}(g)\right)
\end{aligned}
$$

This means that the left invariant vector field $x_{l}$ on $G$ is $\varphi^{m}$-related to the vector field $\mathbf{L}(\sigma)(x)$ on $M$. Therefore, the Related Vector Field Lemma 8.4.9 implies that for $x, y \in \mathbf{L}(G)$ the vector field $\left[x_{l}, y_{l}\right]$ is $\varphi^{m}$-related to $[\mathbf{L}(\sigma)(x)$, $\mathbf{L}(\sigma)(y)]$, which leads for each $m \in M$ to

$$
\begin{aligned}
\mathbf{L}(\sigma)([x, y])(m) & =T_{\mathbf{1}}\left(\varphi^{m}\right)[x, y]_{l}(\mathbf{1})=T_{\mathbf{1}}\left(\varphi^{m}\right)\left[x_{l}, y_{l}\right](\mathbf{1}) \\
& =[\mathbf{L}(\sigma)(x), \mathbf{L}(\sigma)(y)]\left(\varphi^{m}(\mathbf{1})\right)=[\mathbf{L}(\sigma)(x), \mathbf{L}(\sigma)(y)](m)
\end{aligned}
$$

\subsubsection*{Basic Topology of Lie Groups}
In this subsection, we collect some basic topological properties of Lie groups.\\
Proposition 9.1.15. The topology of a Lie group $G$ has the following properties:\\
(i) $G$ is a locally compact space, i.e., each neighborhood of an element of $g$ contains a compact one.\\
(ii) The identity component $G_{0}$ of $G$ is an open normal subgroup which coincides with the arc-component of $\mathbf{1}$.\\
(iii) For a subgroup $H$ of $G$, the following are equivalent:\\
(a) $H$ is a neighborhood of $\mathbf{1}$.\\
(b) $H$ is open.\\
(c) $H$ is open and closed.\\
(d) $H$ contains $G_{0}$.\\
(iv) If the set $\pi_{0}(G):=G / G_{0}$ of connected components of $G$ is countable, then, in addition, the following statements hold:\\
(a) $G$ is countable at infinity, i.e., a countable union of compact subsets.\\
(b) For each $\mathbf{1}$-neighborhood $U$ in $G$, there exists a sequence $\left(g_{n}\right)_{n \in \mathbb{N}}$ in $G$ with $G=\bigcup_{n \in \mathbb{N}} g_{n} U$.\\
(c) $G$ is second countable, i.e., the topology of $G$ has a countable basis.\\
(d) If $\left(U_{i}\right)_{i \in I}$ is a pairwise disjoint collection of open subsets of $G$, then $I$ is countable.

Proof. (i) This is true for any smooth $n$-dimensional manifold $M$. If $m \in M$, $V$ is a neighborhood of $m$ and ( $\varphi, U$ ) is a chart with $m \in M$, then $\varphi(U \cap V)$ is a neighborhood of $\varphi(m)$ in $\mathbb{R}^{n}$. If $B \subseteq \varphi(U \cap V)$ is a closed ball around $\varphi(m)$, which is compact due to the Heine-Borel Theorem, its inverse image $\varphi^{-1}(B)$ is a compact neighborhood of $m$, contained in $V$. Here we use that $M$ is Hausdorff to see that $\varphi^{-1}(B)$ is compact.\\
(ii) Since $G$ is a smooth manifold, each point has an open neighborhood $U$ homeomorphic to an open ball in some $\mathbb{R}^{n}$. Then $U$ is in particular arcwise\\
connected. This implies that the arc-components of $G$ are open, hence that they coincide with the connected components.

To see that the identity component $G_{0}$ of $G$ is a subgroup, we first note that $G_{0} G_{0}$ is the image of the connected set $G_{0} \times G_{0}$ under the multiplication map, hence connected. Since it contains 1 , we find $G_{0} G_{0} \subseteq G_{0}$. Similarly, we see that the inversion preserves $G_{0}$, i.e., $G_{0}^{-1} \subseteq G_{0}$, showing that $G_{0}$ is a subgroup of $G$. Each conjugation $c_{g}(x):=g x g^{-1}$ fixes the identity element $\mathbf{1}$, hence maps the identity component $G_{0}$ into itself. Thus $G_{0}$ is normal.\\
(iii) (a) $\Rightarrow$ (b): If $H$ is a neighborhood of $\mathbf{1}$, then each coset $g H$ is a neighborhood of $g$ because the left multiplication maps $\lambda_{g}: G \rightarrow G$ are homeomorphisms. Hence all left cosets of $H$ are open. In particular, $H$ is open.\\
(b) $\Rightarrow$ (c): If $H$ is an open subgroup, then its complement is the union of all cosets $g H, g \notin H$, hence also open. Therefore, $H$ is also closed.\\
(c) $\Rightarrow$ (d): If $H$ is open and closed, then the connectedness of $G_{0}$ implies $G_{0} \subseteq H$.\\
(d) $\Rightarrow$ (a) is trivial.\\
(iv) (a) In view of (i), there exists a compact identity neighborhood $U$ in $G$. Replacing $U$ by $U \cap U^{-1}$, we may w.l.o.g. assume that $U=U^{-1}$. Then each set

$$
U^{n}:=\left\{u_{1} \cdots u_{n}: u_{i} \in U\right\}
$$

is also compact because it is the image of the compact topological product space $U^{\times n}$ under the $n$-fold multiplication map which is continuous.

Now $H:=\bigcup_{n \in \mathbb{N}} U^{n}$ is a subgroup of $G$ which is a $\mathbf{1}$-neighborhood, and (iii) implies $G_{0} \subseteq H$. Hence the set of $H$-cosets is countable, and since each coset $g H$ is a union of the countably many compact subsets $g U^{n}$, we see that $G$ also is a countable union of compact subsets.\\
(b) In view of (a), we have $G=\bigcup_{n \in \mathbb{N}} K_{n}$, where each $K_{n}$ is a compact subset. For each $n$, the open sets $k U^{\circ}, k \in K_{n}$, cover the compact set $K_{n}$, so that there exist finitely many $k_{n, 1}, \ldots, k_{n, m_{n}}$ with $K_{n} \subseteq \bigcup_{j=1}^{m_{n}} k_{n, j} U$. Then $G \subseteq \bigcup_{n \in \mathbb{N}} \bigcup_{j=1}^{m_{n}} k_{n, j} U$.\\
(c) Let $\left(U_{n}\right)_{n \in \mathbb{N}}$ be a countable basis of open 1-neighborhoods, we may take $U_{n}=\varphi\left(\frac{1}{n} B\right)$, where $B \subseteq \mathbf{L}(G)$ is an open ball with respect to some norm and $\varphi: B \rightarrow G$ is a diffeomorphism onto an open subset of $G$ with $\varphi(0)=\mathbf{1}$. In view of (b), there exists for each $n \in \mathbb{N}$ a sequence $\left(g_{n, k}\right)_{k \in \mathbb{N}}$ in $G$ with $G=\bigcup_{k \in \mathbb{N}} g_{n, k} U_{n}$. We claim that $\left\{g_{n, k} U_{n}: n, k \in \mathbb{N}\right\}$ is a basis for the topology of $G$. In fact, if $O \subseteq G$ is an open subset and $g \in O$, then there exists some $n$ with $g U_{n} \subseteq O$. Next we pick $m$ such that $U_{m}^{-1} U_{m} \subseteq U_{n}$ and some $k \in \mathbb{N}$ with $g \in g_{m, k} U_{m}$. Then $g_{m, k} U_{m} \subseteq g U_{m}^{-1} U_{m} \subseteq g U_{n} \subseteq O$, and this proves our claim.\\
(d) follows immediately from (c).

\subsubsection*{Exercises for Section 9.1}
Exercise 9.1.1. Show that the natural group structure on $\mathbb{T} \cong \mathbb{S}^{1} \subseteq \mathbb{C}^{\times}$ turns it into a Lie group.

Exercise 9.1.2. Let $G_{1}, \ldots, G_{n}$ be Lie groups and $G:=G_{1} \times \cdots \times G_{n}$, endowed with the direct product group structure

$$
\left(g_{1}, \ldots, g_{n}\right)\left(g_{1}^{\prime}, \ldots, g_{n}^{\prime}\right):=\left(g_{1} g_{1}^{\prime}, \ldots, g_{n} g_{n}^{\prime}\right)
$$

and the product manifold structure. Show that $G$ is a Lie group with

$$
\mathbf{L}(G) \cong \mathbf{L}\left(G_{1}\right) \times \cdots \times \mathbf{L}\left(G_{n}\right)
$$

Exercise 9.1.3. Let $V$ and $W$ be finite-dimensional real vector spaces and $\beta: V \times V \rightarrow W$ a bilinear map. Show that $G:=W \times V$ is a Lie group with respect to

$$
(w, v)\left(w^{\prime}, v^{\prime}\right):=\left(w+w^{\prime}+\beta\left(v, v^{\prime}\right), v+v^{\prime}\right)
$$

For $(w, v) \in \mathbf{L}(G) \cong T_{(0,0)}(G)$, find a formula for the corresponding left invariant vector field $(w, v)_{l}$, considered as a smooth function $G \rightarrow W \times V$.

Exercise 9.1.4 (Automatic smoothness of the inversion). Let $G$ be an $n$-dimensional smooth manifold, endowed with a group structure for which the multiplication map $m_{G}$ is smooth. Show that:\\
(1) $T_{(g, h)}\left(m_{G}\right)=T_{g}\left(\rho_{h}\right)+T_{h}\left(\lambda_{g}\right)$ for $\lambda_{g}(x)=g x$ and $\rho_{h}(x)=x h$.\\
(2) $T_{(\mathbf{1}, \mathbf{1})}\left(m_{G}\right)(v, w)=v+w$.\\
(3) The inverse map $\iota_{G}: G \rightarrow G, g \mapsto g^{-1}$ is smooth if it is smooth in a neighborhood of $\mathbf{1}$.\\
(4) The inverse map $\iota_{G}$ is smooth.

Exercise 9.1.5. Let $A$ be a finite-dimensional unital real algebra and $A^{\times}$its group of units. We write $\lambda_{a}(b):=a b$ for the left multiplication with $a \in A$. Show that:\\
(1) $A^{\times}=\left\{a \in A: \operatorname{det}\left(\lambda_{a}\right) \neq 0\right\}$.\\
(2) $A^{\times}$is an open subset of $A$ and with respect to the corresponding manifold structure it is a Lie group.\\
(3) Identifying vector fields on the open subset $A^{\times}$with smooth $A$-valued functions, a vector field $X \in \mathcal{V}\left(A^{\times}\right) \cong C^{\infty}\left(A^{\times}, A\right)$ is left invariant if and only if there exists an element $x \in A$ with $X(a)=a x$ for $a \in A^{\times}$.

Exercise 9.1.6. Let $G$ be a Lie group and $X$ a vector field on $G$, viewed as a derivation of $C^{\infty}(G)$. Show that $X$ is left invariant if and only if

$$
(\mathrm{id} \otimes X) \circ m_{G}^{*}=m_{G}^{*} \circ X
$$

where $m_{G}^{*} f(g, h)=f(g h)$ and $(\mathrm{id} \otimes X) F(g, h)=\left(X F_{g}\right)(h)$ for $f \in C^{\infty}(G)$, $F \in C^{\infty}(G \times G)$ and $F_{g}(h)=F(g, h)$.

Exercise 9.1.7. Let $G=\mathbb{R}^{n}$. Show that the vector fields $X_{i}:=\frac{\partial}{\partial x_{i}}$ form a basis for the space of (left) invariant vector fields.

Exercise 9.1.8. Consider the three-dimensional Heisenberg group

$$
G=\left\{\left(\begin{array}{ccc}
1 & x & y \\
0 & 1 & z \\
0 & 0 & 1
\end{array}\right): x, y, z \in \mathbb{R}\right\}
$$

Determine the space of (left) invariant vector fields in the coordinates ( $x, y, z$ ).

\subsection*{The Exponential Function of a Lie Group}
In the preceding section, we have introduced the Lie functor which assigns to a Lie group $G$ its Lie algebra $\mathbf{L}(G)$ and to a morphism $\varphi$ of Lie groups its tangent morphism $\mathbf{L}(\varphi)$ of Lie algebras. In this section, we introduce a key tool of Lie theory which will allow us to also go in the opposite direction: the exponential function $\exp _{G}: \mathbf{L}(G) \rightarrow G$. It is a natural generalization of the matrix exponential map, which is obtained for $G=\mathrm{GL}_{n}(\mathbb{R})$ and its Lie algebra $\mathbf{L}(G)=\mathfrak{g} \mathfrak{l}_{n}(\mathbb{R})$. We conclude this section with a discussion of the naturality of the exponential function (Proposition 9.2.10) and the Lie group versions of the Trotter Product Formula and the Commutator Formula.

\subsubsection*{Basic Properties of the Exponential Function}
Proposition 9.2.1. Each left invariant vector field $X$ on $G$ is complete.\\
Proof. Let $g \in G$ and $\gamma: I \rightarrow G$ be the unique maximal integral curve (cf. Theorem 8.5.5) of $X \in \mathcal{V}(G)^{l}$ with $\gamma(0)=g$.

For each $h \in G$, we have $\left(\lambda_{h}\right)_{*} X=X$, which implies that $\eta:=\lambda_{h} \circ \gamma$ is also an integral curve of $X$. Put $h=\gamma(s) g^{-1}$ for some $s>0$. Then

$$
\eta(0)=\left(\lambda_{h} \circ \gamma\right)(0)=h \gamma(0)=h g=\gamma(s)
$$

and the uniqueness of integral curves implies that $\gamma(t+s)=\eta(t)$ for all $t$ in the interval $I \cap(I-s)$ which is nonempty because it contains 0 . In view of the maximality of $I$, it now follows that $I-s \subseteq I$, and hence that $I-n s \subseteq I$ for each $n \in \mathbb{N}$, so that the interval $I$ is unbounded from below. Applying the same argument to some $s<0$, we see that $I$ is also unbounded from above. Hence $I=\mathbb{R}$, which means that $X$ is complete.

Definition 9.2.2. We now define the exponential function

$$
\exp _{G}: \mathbf{L}(G) \rightarrow G, \quad \exp _{G}(x):=\gamma_{x}(1)
$$

where $\gamma_{x}: \mathbb{R} \rightarrow G$ is the unique maximal integral curve of the left invariant vector field $x_{l}$, satisfying $\gamma_{x}(0)=\mathbf{1}$. This means that $\gamma_{x}$ is the unique solution of the initial value problem

$$
\gamma(0)=\mathbf{1}, \quad \gamma^{\prime}(t)=x_{l}(\gamma(t))=\gamma(t) . x \quad \text { for all } \quad t \in \mathbb{R} .
$$

Example 9.2.3. (a) Let $G:=(V,+)$ be the additive group of a finitedimensional vector space. The left invariant vector fields on $V$ are given by

$$
x_{l}(w):=\left.\frac{d}{d t}\right|_{t=0} w+t x=x
$$

so that they are simply the constant vector fields. Hence (cf. Lemma 8.4.3)

$$
\left[x_{l}, y_{l}\right](0)=\mathrm{d} y_{l}\left(x_{l}(0)\right)-\mathrm{d} x_{l}\left(y_{l}(0)\right)=\mathrm{d} y_{l}(x)-\mathrm{d} x_{l}(y)=0 .
$$

Therefore, $\mathbf{L}(V)$ is an abelian Lie algebra.\\
For each $x \in V$, the flow of $x_{l}$ is given by $\Phi^{x_{l}}(t, v)=v+t x$, so that

$$
\exp _{V}(x)=\Phi^{x_{l}}(1,0)=x, \quad \text { i.e., } \quad \exp _{V}=\operatorname{id}_{V}
$$

(b) Now let $G:=\mathrm{GL}_{n}(\mathbb{R})$ be the Lie group of invertible ( $n \times n$ )-matrices, which inherits its manifold structure from the embedding as an open subset of the vector space $M_{n}(\mathbb{R})$.

The left invariant vector field $A_{l}$ corresponding to a matrix $A$ is given by

$$
A_{l}(g)=T_{\mathbf{1}}\left(\lambda_{g}\right) A=g A
$$

because $\lambda_{g}(h)=g h$ extends to a linear endomorphism of $M_{n}(\mathbb{R})$. The unique solution $\gamma_{A}: \mathbb{R} \rightarrow \mathrm{GL}_{n}(\mathbb{R})$ of the initial value problem

$$
\gamma(0)=\mathbf{1}, \quad \gamma^{\prime}(t)=A_{l}(\gamma(t))=\gamma(t) A
$$

is the curve describing the fundamental system of the linear differential equation defined by the matrix $A$ :

$$
\gamma_{A}(t)=e^{t A}=\sum_{k=0}^{\infty} \frac{1}{k!} t^{k} A^{k}
$$

It follows that $\exp _{G}(A)=e^{A}$ is the matrix exponential function.\\
The Lie algebra $\mathbf{L}(G)$ of $G$ is determined from

$$
\begin{aligned}
{[A, B] } & =\left[A_{l}, B_{l}\right](\mathbf{1})=\mathrm{d} B_{l}(\mathbf{1}) A_{l}(\mathbf{1})-\mathrm{d} A_{l}(\mathbf{1}) B_{l}(\mathbf{1}) \\
& =\mathrm{d} B_{l}(\mathbf{1}) A-\mathrm{d} A_{l}(\mathbf{1}) B=A B-B A
\end{aligned}
$$

Therefore, the Lie bracket on $\mathbf{L}(G)=T_{\mathbf{1}}(G) \cong M_{n}(\mathbb{R})$ is given by the commutator bracket. This Lie algebra is denoted $\mathfrak{g} \mathfrak{l}_{n}(\mathbb{R})$, to express that it is the Lie algebra of $\mathrm{GL}_{n}(\mathbb{R})$.\\
(c) If $V$ is a finite-dimensional real vector space, then $V \cong \mathbb{R}^{n}$, so that we can immediately use (b) to see that $\mathrm{GL}(V)$ is a Lie group with Lie algebra $\mathfrak{g} \mathfrak{l}(V):=(\operatorname{End}(V),[\cdot, \cdot])$ and exponential function

$$
\exp _{\mathrm{GL}(V)}(A)=\sum_{k=0}^{\infty} \frac{A^{k}}{k!}
$$

Lemma 9.2.4. (a) For each $x \in \mathbf{L}(G)$, the curve $\gamma_{x}: \mathbb{R} \rightarrow G$ is a smooth homomorphism of Lie groups with $\gamma_{x}^{\prime}(0)=x$.\\
(b) The global flow of the left invariant vector field $x_{l}$ is given by

$$
\Phi(t, g)=g \gamma_{x}(t)=g \exp _{G}(t x)
$$

(c) If $\gamma: \mathbb{R} \rightarrow G$ is a smooth homomorphism of Lie groups and $x:=\gamma^{\prime}(0)$, then $\gamma=\gamma_{x}$. In particular, the map

$$
\operatorname{Hom}(\mathbb{R}, G) \rightarrow \mathbf{L}(G), \quad \gamma \mapsto \gamma^{\prime}(0)
$$

is a bijection, where $\operatorname{Hom}(\mathbb{R}, G)$ stands for the set of morphisms, i.e., smooth homomorphisms, of Lie groups $\mathbb{R} \rightarrow G$.

Proof. (a), (b) Since $\gamma_{x}$ is an integral curve of the smooth vector field $x_{l}$, it is a smooth curve. Hence the smoothness of the multiplication in $G$ implies that $\Phi(t, g):=g \gamma_{x}(t)$ defines a smooth map $\mathbb{R} \times G \rightarrow G$. In view of the left invariance of $x_{l}$, we have for each $g \in G$ and $\Phi^{g}(t):=\Phi(t, g)$ the relation

$$
\left(\Phi^{g}\right)^{\prime}(t)=T\left(\lambda_{g}\right) \gamma_{x}^{\prime}(t)=T\left(\lambda_{g}\right) x_{l}\left(\gamma_{x}(t)\right)=x_{l}\left(g \gamma_{x}(t)\right)=x_{l}\left(\Phi^{g}(t)\right)
$$

Therefore, $\Phi^{g}$ is an integral curve of $x_{l}$ with $\Phi^{g}(0)=g$, and this proves that $\Phi$ is the unique maximal flow of the complete vector field $x_{l}$.

In particular, we obtain for $t, s \in \mathbb{R}$ :

$$
\gamma_{x}(t+s)=\Phi(t+s, \mathbf{1})=\Phi(t, \Phi(s, \mathbf{1}))=\Phi(s, \mathbf{1}) \gamma_{x}(t)=\gamma_{x}(s) \gamma_{x}(t)
$$

Hence $\gamma_{x}$ is a group homomorphism $(\mathbb{R},+) \rightarrow G$.\\
(c) If $\gamma:(\mathbb{R},+) \rightarrow G$ is a smooth group homomorphism, then

$$
\Phi(t, g):=g \gamma(t)
$$

defines a global flow on $G$ whose infinitesimal generator is the vector field given by

$$
X(g)=\left.\frac{d}{d t}\right|_{t=0} \Phi(t, g)=T\left(\lambda_{g}\right) \gamma^{\prime}(0)
$$

We conclude that $X=x_{l}$ for $x=\gamma^{\prime}(0)$, so that $X$ is a left invariant vector field. Since $\gamma$ is its unique integral curve through 0 , it follows that $\gamma=\gamma_{x}$. In view of (a), this proves (c).

Proposition 9.2.5. For a Lie group $G$, the exponential function

$$
\exp _{G}: \mathbf{L}(G) \rightarrow G
$$

is smooth and satisfies

$$
T_{0}\left(\exp _{G}\right)=\mathrm{id}_{\mathbf{L}(G)}
$$

In particular, $\exp _{G}$ is a local diffeomorphism at 0 in the sense that it maps some 0 -neighborhood in $\mathbf{L}(G)$ diffeomorphically onto some $\mathbf{1}$-neighborhood in $G$.

Proof. Let $n \in \mathbb{N}$. In view of Lemma 9.2.4(c), we have

$$
\exp _{G}(n x)=\gamma_{x}(n)=\gamma_{x}(1)^{n}=\exp _{G}(x)^{n}
$$

for each $x \in \mathbf{L}(G)$. Since the $n$-fold multiplication map

$$
G^{n} \rightarrow G, \quad\left(g_{1}, \ldots, g_{n}\right) \mapsto g_{1} \cdots g_{n}
$$

is smooth, the $n$th power map $G \rightarrow G, g \mapsto g^{n}$ is smooth. Therefore, it suffices to verify the smoothness of $\exp _{G}$ in some 0 -neighborhood $W$. Then (9.5) immediately implies smoothness in $n W$ for each $n$, and hence on all of $\mathbf{L}(G)$.

The map $\Psi: \mathbf{L}(G) \rightarrow \mathcal{V}(G), x \mapsto x_{l}$ satisfies the assumptions of Proposition 8.5.15 because the map

$$
\mathbf{L}(G) \times G \rightarrow T(G), \quad(x, g) \mapsto x_{l}(g)=g \cdot x
$$

is smooth (Lemma 9.1.6). In the terminology of Proposition 8.5.15, it now follows that the map

$$
\Phi: \mathbb{R} \times \mathbf{L}(G) \times G \rightarrow G, \quad(t, x, g) \mapsto g \gamma_{x}(t)=g \exp _{G}(t x)
$$

is smooth on a neighborhood of $(0,0, \mathbf{1})$. In particular, for some $t>0$, the map $x \mapsto \exp _{G}(t x)$ is smooth on a 0-neighborhood of $\mathbf{L}(G)$, and this proves that $\exp _{G}$ is smooth in some 0-neighborhood.

Finally, we observe that

$$
T_{0}\left(\exp _{G}\right)(x)=\left.\frac{d}{d t}\right|_{t=0} \exp _{G}(t x)=\gamma_{x}^{\prime}(0)=x
$$

so that $T_{0}\left(\exp _{G}\right)=\operatorname{id}_{\mathbf{L}(G)}$.\\
Lemma 9.2.6 (Canonical Coordinates). Let $G$ be a Lie group and $b_{1}$, $\ldots, b_{n}$ be a basis for its Lie algebra $\mathbf{L}(G)$. Then the following maps restrict to diffeomorphisms of some 0 -neighborhood in $\mathbb{R}^{n}$ to some open $\mathbf{1}$-neighborhood in $G$ :\\
(i) $x \mapsto \exp _{G}\left(x_{1} b_{1}+\cdots+x_{n} b_{n}\right)$ (Canonical coordinates of the first kind).\\
(ii) $x \mapsto \exp _{G}\left(x_{1} b_{1}\right) \cdots \exp _{G}\left(x_{n} b_{n}\right)$ (Canonical coordinates of the second kind).

Proof. (i) This is immediate from Proposition 9.2.5.\\
(ii) In view of Proposition 9.2.5, $T_{0}\left(\exp _{G}\right)=\operatorname{id}_{\mathbf{L}(G)}$, and further $T_{\mathbf{1}}\left(m_{G}\right)(x, y)=x+y$ by Lemma 9.1.6. Therefore,

$$
\Phi: \mathbb{R}^{n} \rightarrow G, \quad x \mapsto \exp _{G}\left(x_{1} b_{1}\right) \cdots \exp _{G}\left(x_{n} b_{n}\right)
$$

satisfies $T_{0}(\Phi)(x)=\sum_{i=1}^{n} x_{i} b_{i}$. Hence the claim follows from the Inverse Function Theorem.

Lemma 9.2.7. If $\sigma: G \times M \rightarrow M$ is a smooth action and $x \in \mathbf{L}(G)$, then the global flow of the vector field $\dot{\sigma}(x)$ is given by $\Phi^{x}(t, m)=\exp _{G}(-t x) . m$. In particular,

$$
\dot{\sigma}(x)(m)=\left.\frac{d}{d t}\right|_{t=0} \exp _{G}(-t x) . m
$$

Proof. In the proof of Proposition 9.1.14, we have seen that

$$
T_{g}\left(\varphi^{m}\right) x_{l}(g)=\dot{\sigma}(x)\left(\varphi^{m}(g)\right)
$$

holds for the map $\varphi^{m}(g)=g^{-1} . m$. In view of Proposition 9.2.5, this yields

$$
\left.\frac{d}{d t}\right|_{t=0} \exp _{G}(-t x) \cdot m=T_{\mathbf{1}}\left(\varphi^{m}\right) T_{0}\left(\exp _{G}\right) x=T_{\mathbf{1}}\left(\varphi^{m}\right) x=\dot{\sigma}(x)(m)
$$

and hence proves the lemma.\\
Lemma 9.2.8. If $x, y \in \mathbf{L}(G)$ commute, i.e., $[x, y]=0$, then

$$
\exp _{G}(x+y)=\exp _{G}(x) \exp _{G}(y)
$$

Proof. If $x$ and $y$ commute, then the corresponding left invariant vector fields commute, and Corollary 8.5.17 implies that their flows commute. We conclude that for all $t, s \in \mathbb{R}$ we have

$$
\exp _{G}(t x) \exp _{G}(s y)=\exp _{G}(s y) \exp _{G}(t x)
$$

Therefore,

$$
\gamma(t):=\exp _{G}(t x) \exp _{G}(t y)
$$

is a smooth group homomorphism. In view of

$$
\gamma^{\prime}(0)=T_{(\mathbf{1}, \mathbf{1})}\left(m_{G}\right)(x, y)=x+y
$$

(Lemma 9.1.6), Lemma 9.2.4(c) leads to $\gamma(t)=\exp _{G}(t(x+y))$, and for $t=1$ we obtain the lemma.

Lemma 9.2.9. The subgroup $\left\langle\exp _{G}(\mathbf{L}(G))\right\rangle$ of $G$ generated by $\exp _{G}(\mathbf{L}(G))$ coincides with the identity component $G_{0}$ of $G$, i.e., the connected component containing 1.

Proof. Since $\exp _{G}$ is a local diffeomorphism at 0 (Proposition 9.2.5), the Inverse Function Theorem (see Exercise 8.3.5) implies that $\exp _{G}(\mathbf{L}(G))$ is a neighborhood of $\mathbf{1}$. We conclude that the subgroup $H:=\left\langle\exp _{G}(\mathbf{L}(G))\right\rangle$ generated by the exponential image is a $\mathbf{1}$-neighborhood, hence contains $G_{0}$ (Proposition 9.1.15(iii)(d)). On the other hand, $\exp _{G}$ is continuous, so that it maps the connected space $\mathbf{L}(G)$ into the identity component $G_{0}$ of $G$, which leads to $H \subseteq G_{0}$, and hence to equality. \(\square\)

\subsubsection*{Naturality of the Exponential Function}
In this subsection, we study how the exponential function is related to the Lie functor.

Proposition 9.2.10. Let $\varphi: G_{1} \rightarrow G_{2}$ be a morphism of Lie groups and $\mathbf{L}(\varphi): \mathbf{L}\left(G_{1}\right) \rightarrow \mathbf{L}\left(G_{2}\right)$ its differential at $\mathbf{1}$. Then

$$
\exp _{G_{2}} \circ \mathbf{L}(\varphi)=\varphi \circ \exp _{G_{1}},
$$

i.e., the following diagram commutes\\
\includegraphics[max width=\textwidth, center]{2025_10_18_3eaf04e77f3cb364fce5g-312}

Proof. For $x \in \mathbf{L}\left(G_{1}\right)$, we consider the smooth homomorphism

$$
\gamma_{x} \in \operatorname{Hom}\left(\mathbb{R}, G_{1}\right), \quad \gamma_{x}(t)=\exp _{G_{1}}(t x) .
$$

According to Lemma 9.2.4, we have

$$
\varphi \circ \gamma_{x}(t)=\exp _{G_{2}}(t y)
$$

for $y=\left(\varphi \circ \gamma_{x}\right)^{\prime}(0)=\mathbf{L}(\varphi) x$, because $\varphi \circ \gamma_{x}: \mathbb{R} \rightarrow G_{2}$ is a smooth group homomorphism. For $t=1$, we obtain in particular

$$
\exp _{G_{2}}(\mathbf{L}(\varphi) x)=\varphi\left(\exp _{G_{1}}(x)\right),
$$

which we had to show. \(\square\)

Corollary 9.2.11. Let $G_{1}$ and $G_{2}$ be Lie groups and $\varphi: G_{1} \rightarrow G_{2}$ be a group homomorphism. Then the following are equivalent:\\
(a) $\varphi$ is smooth in an identity neighborhood of $G_{1}$.\\
(b) $\varphi$ is smooth.\\
(c) There exists a linear map $\psi: \mathbf{L}\left(G_{1}\right) \rightarrow \mathbf{L}\left(G_{2}\right)$ satisfying

$$
\exp _{G_{2}} \circ \psi=\varphi \circ \exp _{G_{1}}
$$

Proof. (a) $\Rightarrow$ (b): Let $U$ be an open $\mathbf{1}$-neighborhood of $G_{1}$ such that $\left.\varphi\right|_{U}$ is smooth. Since each left translation $\lambda_{g}$ is a diffeomorphism, $\lambda_{g}(U)=g U$ is an open neighborhood of $g$, and we have

$$
\varphi(g x)=\varphi(g) \varphi(x), \quad \text { i.e., } \quad \varphi \circ \lambda_{g}=\lambda_{\varphi(g)} \circ \varphi .
$$

Hence the smoothness of $\varphi$ on $U$ implies the smoothness of $\varphi$ on $g U$, and therefore that $\varphi$ is smooth.\\
(b) $\Rightarrow$ (c): If $\varphi$ is smooth, then $\psi:=\mathbf{L}(\varphi)$ satisfies (9.8).\\
(c) $\Rightarrow$ (a): If $\psi$ is a linear map satisfying (9.8), then the fact that the exponential functions $\exp _{G_{1}}$ and $\exp _{G_{2}}$ are local diffeomorphisms and the smoothness of the linear map $\psi$ imply (a).

Corollary 9.2.12. If $\varphi_{1}, \varphi_{2}: G_{1} \rightarrow G_{2}$ are morphisms of Lie groups with $\mathbf{L}\left(\varphi_{1}\right)=\mathbf{L}\left(\varphi_{2}\right)$, then $\varphi_{1}$ and $\varphi_{2}$ coincide on the identity component of $G_{1}$.

Proof. In view of Proposition 9.2.10, we have for $x \in \mathbf{L}\left(G_{1}\right)$ :

$$
\varphi_{1}\left(\exp _{G_{1}}(x)\right)=\exp _{G_{2}}\left(\mathbf{L}\left(\varphi_{1}\right) x\right)=\exp _{G_{2}}\left(\mathbf{L}\left(\varphi_{2}\right) x\right)=\varphi_{2}\left(\exp _{G_{1}}(x)\right)
$$

so that $\varphi_{1}$ and $\varphi_{2}$ coincide on the image of $\exp _{G_{1}}$, hence on the subgroup generated by this set. Now the assertion follows from Lemma 9.2.9.

Proposition 9.2.13. For a morphism $\varphi: G_{1} \rightarrow G_{2}$ of Lie groups, the following assertions hold:\\
(1) $\operatorname{ker} \mathbf{L}(\varphi)=\left\{x \in \mathbf{L}\left(G_{1}\right): \exp _{G_{1}}(\mathbb{R} x) \subseteq \operatorname{ker} \varphi\right\}$.\\
(2) $\varphi$ is an open map if and only if $\mathbf{L}(\varphi)$ is surjective.\\
(3) If $\mathbf{L}(\varphi)$ is a linear isomorphism and $\varphi$ is bijective, then $\varphi$ is an isomorphism of Lie groups.

Proof. (1) The condition $x \in \operatorname{ker} \mathbf{L}(\varphi)$ is equivalent to

$$
\{\boldsymbol{1}\}=\exp _{G_{2}}(\mathbb{R} \mathbf{L}(\varphi) x)=\varphi\left(\exp _{G_{1}}(\mathbb{R} x)\right)
$$

(2) Suppose first that $\varphi$ is an open map. Since $\exp _{G_{i}}, i=1,2$, are local diffeomorphisms,

$$
\exp _{G_{2}} \circ \mathbf{L}(\varphi)=\varphi \circ \exp _{G_{1}}
$$

implies that there exists some 0 -neighborhood in $\mathbf{L}\left(G_{1}\right)$ on which $\mathbf{L}(\varphi)$ is an open map, hence that $\mathbf{L}(\varphi)$ is surjective.

If, conversely, $\mathbf{L}(\varphi)$ is surjective, then $\mathbf{L}(\varphi)$ is an open map, so that the relation (9.9) implies that there exists an open $\mathbf{1}$-neighborhood $U_{1}$ in $G_{1}$ such that $\left.\varphi\right|_{U_{1}}$ is an open map. We claim that this implies that $\varphi$ is an open map.

In fact, suppose that $O \subseteq G_{1}$ is open and $g \in O$. Then there exists an open 1-neighborhood $U_{2}$ of $G_{1}$ with $g U_{2} \subseteq O$ and $U_{2} \subseteq U_{1}$. Then

$$
\varphi(O) \supseteq \varphi\left(g U_{2}\right)=\varphi(g) \varphi\left(U_{2}\right)
$$

and since $\varphi\left(U_{2}\right)$ is open in $G_{2}$, we see that $\varphi(O)$ is a neighborhood of $\varphi(g)$, hence that $\varphi(O)$ is open because $g \in O$ was arbitrary.\\
(3) From the relation $\exp _{G_{2}} \circ \mathbf{L}(\varphi)=\varphi \circ \exp _{G_{1}}$ and the bijectivity of $\varphi$, we derive that the group homomorphism $\varphi^{-1}$ satisfies

$$
\varphi^{-1} \circ \exp _{G_{2}}=\exp _{G_{1}} \circ \mathbf{L}(\varphi)^{-1}
$$

so that Corollary 9.2.11 implies that $\varphi^{-1}$ is also smooth.\\
Proposition 9.2.14. Let $G$ be a Lie group with Lie algebra $\mathbf{L}(G)$. Then for $x, y \in \mathbf{L}(G)$ the following equations hold:\\
(1) (Product Formula)

$$
\exp _{G}(x+y)=\lim _{k \rightarrow \infty}\left(\exp _{G}\left(\frac{1}{k} x\right) \exp _{G}\left(\frac{1}{k} y\right)\right)^{k}
$$

(2) (Commutator Formula)

$$
\exp _{G}([x, y])=\lim _{k \rightarrow \infty}\left(\exp _{G}\left(\frac{1}{k} x\right) \exp _{G}\left(\frac{1}{k} y\right) \exp _{G}\left(-\frac{1}{k} x\right) \exp _{G}\left(-\frac{1}{k} y\right)\right)^{k^{2}}
$$

Proof. Let $U \subseteq \mathbf{L}(G)$ be an open 0-neighborhood for which

$$
\exp _{U}:=\left.\exp _{G}\right|_{U}: U \rightarrow \exp _{G}(U)
$$

is a diffeomorphism onto an open subset of $G$. Put

$$
U^{1}:=\left\{(x, y) \in U \times U: \exp _{G}(x) \exp _{G}(y) \in \exp _{G}(U)\right\}
$$

and observe that this is an open subset of $U \times U$ containing $(0,0)$ because $\exp _{G}(U)$ is open and $\exp _{G}$ is continuous.

For $(x, y) \in U_{1}$, we then define

$$
x * y:=\exp _{U}^{-1}\left(\exp _{G}(x) \exp _{G}(y)\right)
$$

and thus obtain a smooth map

$$
m: U^{1} \rightarrow \mathbf{L}(G), \quad(x, y) \mapsto x * y .
$$

(1) In view of $m(0, x)=m(x, 0)=x$, we have

$$
\mathrm{d} m(0,0)(x, y)=\mathrm{d} m(0,0)(x, 0)+\mathrm{d} m(0,0)(0, y)=x+y
$$

This implies that

$$
\begin{aligned}
\lim _{k \rightarrow \infty} k \cdot\left(\frac{1}{k} x * \frac{1}{k} y\right) & =\lim _{k \rightarrow \infty} k \cdot\left(m\left(\frac{1}{k} x, \frac{1}{k} y\right)-m(0,0)\right) \\
& =\mathrm{d} m(0,0)(x, y)=x+y
\end{aligned}
$$

Applying $\exp _{G}$, it follows that

$$
\begin{aligned}
\exp _{G}(x+y) & =\lim _{k \rightarrow \infty} \exp _{G}\left(k \cdot\left(\frac{1}{k} x * \frac{1}{k} y\right)\right)=\lim _{k \rightarrow \infty} \exp _{G}\left(\frac{1}{k} x * \frac{1}{k} y\right)^{k} \\
& =\lim _{k \rightarrow \infty}\left(\exp _{G}\left(\frac{1}{k} x\right) \exp _{G}\left(\frac{1}{k} y\right)\right)^{k}
\end{aligned}
$$

(2) Now let $x_{l}^{*}:=T\left(\exp _{U}\right)^{-1} \circ x_{l} \circ \exp _{U}$ be the smooth vector field on $U$ corresponding to the left invariant vector field $x_{l}$ on $\exp _{G}(U)$. Then $x_{l}^{*}$ and $x_{l}$ are $\exp _{U}$-related, so that the Related Vector Field Lemma 8.4.9 leads to

$$
\left[x_{l}^{*}, y_{l}^{*}\right](0)=T_{0}\left(\exp _{U}\right)\left[x_{l}^{*}, y_{l}^{*}\right](0)=\left[x_{l}, y_{l}\right]\left(\exp _{G}(0)\right)=\left[x_{l}, y_{l}\right](\mathbf{1})=[x, y]
$$

The local flow of $x_{l}$ through a point $\exp _{U}(y)$ is given by the curve $t \mapsto \exp _{G}(y) \exp _{G}(t x)$, which implies that the integral curve of $x_{l}^{*}$ through $y \in U$ is given for $t$ close to 0 by $\Phi_{t}^{x_{l}^{*}}(y)=y * t x$. We therefore obtain with Remark 8.5.18 (on commutators of flows):

$$
\begin{aligned}
\left.\frac{\partial^{2}}{\partial s \partial t}\right|_{t=s=0} t x * s y *(-t x) *(-s y) & =\left.\frac{\partial^{2}}{\partial s \partial t}\right|_{t=s=0} \Phi_{-s}^{y_{l}^{*}} \circ \Phi_{-t}^{x_{l}^{*}} \circ \Phi_{s}^{y_{l}^{*}} \circ \Phi_{t}^{x_{l}^{*}}(0) \\
& =\left[-x_{l}^{*},-y_{l}^{*}\right](0)=\left[x_{l}^{*}, y_{l}^{*}\right](0)=[x, y] .
\end{aligned}
$$

Note that $F(t, s):=t x * s y *(-t x) *(-s y)$ vanishes for $t=0$ and $s=0$.\\
For $f(t):=F(t, t)$, we have

$$
\begin{gathered}
f^{\prime}(t)=\frac{\partial F}{\partial t}(t, t)+\frac{\partial F}{\partial s}(t, t) \\
f^{\prime \prime}(t)=\frac{\partial^{2} F}{\partial t^{2}}(t, t)+2 \frac{\partial^{2} F}{\partial t \partial s}(t, t)+\frac{\partial^{2} F}{\partial s^{2}}(t, t)
\end{gathered}
$$

and

$$
\frac{\partial^{2} F}{\partial t^{2}}(0,0)=\frac{\partial^{2} F}{\partial s^{2}}(0,0)=f^{\prime}(0)=0
$$

Therefore,

$$
\frac{1}{2} f^{\prime \prime}(0)=\frac{\partial^{2} F}{\partial t \partial s}(0,0)=[x, y]
$$

leads to

$$
\lim _{k \rightarrow \infty} k^{2}\left(\frac{1}{k} x * \frac{1}{k} y *\left(-\frac{1}{k} x\right) *\left(-\frac{1}{k} y\right)\right)=\lim _{k \rightarrow \infty} k^{2} f\left(\frac{1}{k}\right)=\frac{1}{2} f^{\prime \prime}(0)=[x, y] .
$$

Applying the exponential function, we obtain the Commutator Formula.

Theorem 9.2.15 (One-parameter Group Theorem). Let $G$ be a Lie group. For each $x \in \mathfrak{g}:=\mathbf{L}(G)$, the map $\gamma_{x}:(\mathbb{R},+) \rightarrow G, t \mapsto \exp _{G}(t x)$ is a smooth group homomorphism. Conversely, every continuous one-parameter group $\gamma: \mathbb{R} \rightarrow G$ is of this form.

Proof. The first assertion is an immediate consequence of Lemma 9.2.4(c). It therefore remains to show that each continuous one-parameter group $\gamma$ of $G$ is a $\gamma_{x}$ for some $x \in \mathfrak{g}$. Let $U=-U$ be a convex 0-neighborhood in $\mathfrak{g}$ for which $\left.\exp _{G}\right|_{U}$ is a diffeomorphism onto an open subset of $G$ and put $U_{1}:=\frac{1}{2} U$. Since $\gamma$ is continuous at 0 , there exists an $\varepsilon>0$ such that $\gamma([-\varepsilon, \varepsilon]) \subseteq \exp _{G}\left(U_{1}\right)$. Then $\alpha(t):=\left(\left.\exp _{G}\right|_{U}\right)^{-1}(\gamma(t))$ defines a continuous curve $\alpha:[-\varepsilon, \varepsilon] \rightarrow U_{1}$ with $\exp (\alpha(t))=\gamma(t)$ for $|t| \leq \varepsilon$. With the same arguments as in the proof of Theorem 3.2.6, we see that $\alpha(t)=t x$ for some $x \in \mathfrak{g}$. Hence $\gamma(t)=\exp _{G}(t x)$ for $|t| \leq \varepsilon$, but then $\gamma(n t)=\exp _{G}(n t x)$ for $n \in \mathbb{N}$ leads to $\gamma(t)=\exp _{G}(t x)$ for each $t \in \mathbb{R}$.

Theorem 9.2.16 (Automatic Smoothness Theorem). Each continuous homomorphism $\varphi: G \rightarrow H$ of Lie groups is smooth.

Proof. From Theorem 9.2.15, we know that the map

$$
\mathbf{L}(G) \rightarrow \operatorname{Hom}_{c}(\mathbb{R}, G), \quad x \mapsto \gamma_{x}, \quad \gamma_{x}(t):=\exp _{G}(t x)
$$

is a bijection, where $\operatorname{Hom}_{c}(\mathbb{R}, G)$ denotes the set of all continuous oneparameter groups of $G$. For $x \in \mathbf{L}\left(G_{1}\right)$, we consider the continuous homomorphism $\varphi \circ \gamma_{x} \in \operatorname{Hom}_{c}\left(\mathbb{R}, G_{2}\right)$. Since this one-parameter group is smooth (Theorem 9.2.15), it is of the form

$$
\varphi \circ \gamma_{x}(t)=\exp _{G_{2}}(t y)
$$

for $y=\left(\varphi \circ \gamma_{x}\right)^{\prime}(0) \in \mathbf{L}\left(G_{2}\right)$. We define a map $\mathbf{L}(\varphi): \mathbf{L}\left(G_{1}\right) \rightarrow \mathbf{L}\left(G_{2}\right)$ by $\mathbf{L}(\varphi) x:=\left(\varphi \circ \gamma_{x}\right)^{\prime}(0)$. For $t=1$, we then obtain

$$
\exp _{G_{2}} \circ \mathbf{L}(\varphi)=\varphi \circ \exp _{G_{1}}: \mathbf{L}\left(G_{1}\right) \rightarrow G_{2}
$$

Next we show that $\mathbf{L}(\varphi)$ is a linear map. Our definition immediately shows that $\mathbf{L}(\varphi) \lambda x=\lambda \mathbf{L}(\varphi) x$ for each $x \in \mathbf{L}\left(G_{1}\right)$. Further, the Product Formula (Proposition 9.2.14) yields

$$
\begin{aligned}
& \exp _{G_{2}}(\mathbf{L}(\varphi)(x+y))=\varphi\left(\exp _{G_{1}}(x+y)\right) \\
& =\lim _{k \rightarrow \infty} \varphi\left(\exp _{G_{1}}\left(\frac{1}{k} x\right) \exp _{G_{1}}\left(\frac{1}{k} y\right)\right)^{k} \\
& =\lim _{k \rightarrow \infty}\left(\exp _{G_{2}}\left(\frac{1}{k} \mathbf{L}(\varphi) x\right) \exp _{G_{2}}\left(\frac{1}{k} \mathbf{L}(\varphi) y\right)\right)^{k} \\
& =\exp _{G_{2}}(\mathbf{L}(\varphi) x+\mathbf{L}(\varphi) y)
\end{aligned}
$$

This proves that $\mathbf{L}(\varphi)(x+y)=\mathbf{L}(\varphi) x+\mathbf{L}(\varphi) y$, so that $\mathbf{L}(\varphi)$ is indeed a linear map, hence in particular smooth. Since both maps $\exp _{G_{i}}$ are local diffeomorphisms, (9.10) implies that $\varphi$ is smooth in an identity neighborhood of $G_{1}$, hence smooth by Corollary 9.2.11.

Corollary 9.2.17. A topological group $G$ carries at most one Lie group structure.

Proof. If $G_{1}$ and $G_{2}$ are two Lie groups which are isomorphic as topological groups, then the Automatic Smoothness Theorem applies to each topological isomorphism $\varphi: G_{1} \rightarrow G_{2}$ and shows that $\varphi$ is smooth. It likewise applies to $\varphi^{-1}$, so that $\varphi$ is an isomorphism of Lie groups.

\subsubsection*{The Adjoint Representation}
The Lie functor associates linear automorphisms of the Lie algebra with conjugations on the Lie group. The resulting representation of the Lie group is called the adjoint representation. Its interplay with the exponential function will be important in the entire theory.

Definition 9.2.18. (a) We know that for each finite-dimensional vector space $V$, the group $\mathrm{GL}(V)$ carries a natural Lie group structure. For a Lie group $G$, a smooth homomorphism $\pi: G \rightarrow \mathrm{GL}(V)$ is a called a representation of $G$ on $V$ (cf. Exercise 9.2.3).

Any representation defines a smooth action of $G$ on $V$ via

$$
\sigma(g, v):=\pi(g)(v)
$$

In this sense, representations are the same as linear actions, i.e., actions on vector spaces for which the $\sigma_{g}$ are linear.\\
(b) If $\mathfrak{g}$ is a Lie algebra, then a homomorphism of Lie algebras $\varphi: \mathfrak{g} \rightarrow \mathfrak{g} \mathfrak{l}(V)$ is called a representation of $\mathfrak{g}$ on $V$ (cf. Definition 5.1.4).

As a consequence of Proposition 9.1.8, we obtain\\
Proposition 9.2.19. If $\varphi: G \rightarrow \mathrm{GL}(V)$ is a representation of $G$, then $\mathbf{L}(\varphi): \mathbf{L}(G) \rightarrow \mathfrak{g} \mathfrak{l}(V)$ is a representation of its Lie algebra $\mathbf{L}(G)$.

The representation $\mathbf{L}(\varphi)$ obtained in Proposition 9.2.19 from the group representation $\varphi$ is called the derived representation. This is motivated by the fact that for each $x \in \mathbf{L}(G)$ we have

$$
\mathbf{L}(\varphi)(x)=\left.\frac{d}{d t}\right|_{t=0} e^{t \mathbf{L}(\varphi) x}=\left.\frac{d}{d t}\right|_{t=0} \varphi\left(\exp _{G} t x\right)
$$

Let $G$ be a Lie group and $\mathbf{L}(G)$ its Lie algebra. For $g \in G$, we recall the conjugation automorphism $c_{g} \in \operatorname{Aut}(G), c_{g}(x)=g x g^{-1}$, and define

$$
\operatorname{Ad}(g):=\mathbf{L}\left(c_{g}\right) \in \operatorname{Aut}(\mathbf{L}(G)) .
$$

Then

$$
\operatorname{Ad}\left(g_{1} g_{2}\right)=\mathbf{L}\left(c_{g_{1} g_{2}}\right)=\mathbf{L}\left(c_{g_{1}}\right) \circ \mathbf{L}\left(c_{g_{2}}\right)=\operatorname{Ad}\left(g_{1}\right) \operatorname{Ad}\left(g_{2}\right)
$$

shows that $\mathrm{Ad}: G \rightarrow \operatorname{Aut}(\mathbf{L}(G))$ is a group homomorphism. It is called the adjoint representation. To see that it is smooth, we observe that for each $x \in \mathbf{L}(G)$ we have

$$
\operatorname{Ad}(g) x=T_{\mathbf{1}}\left(c_{g}\right) x=T_{\mathbf{1}}\left(\lambda_{g} \circ \rho_{g^{-1}}\right) x=T_{g^{-1}}\left(\lambda_{g}\right) T_{\mathbf{1}}\left(\rho_{g^{-1}}\right) x=0_{g} \cdot x \cdot 0_{g^{-1}}
$$

in the Lie group $T(G)$ (Lemma 9.1.6). Since the multiplication in $T(G)$ is smooth, the representation Ad of $G$ on $\mathbf{L}(G)$ is smooth (cf. Exercise 9.2.3), and

$$
\mathbf{L}(\mathrm{Ad}): \mathbf{L}(G) \rightarrow \mathfrak{g} \mathfrak{l}(\mathbf{L}(G))
$$

is a representation of $\mathbf{L}(G)$ on $\mathbf{L}(G)$. The following lemma gives a formula for this representation.\\
Lemma 9.2.20. $\mathbf{L}(\mathrm{Ad})=$ ad, i.e., $\mathbf{L}(\mathrm{Ad})(x)(y)=[x, y]$.\\
Proof. Let $x, y \in \mathbf{L}(G)$ and $x_{l}, y_{l}$ be the corresponding left invariant vector fields. Corollary 9.1.10 implies for $g \in G$ the relation

$$
\left(c_{g}\right)_{*} y_{l}=\left(\mathbf{L}\left(c_{g}\right) y\right)_{l}=(\operatorname{Ad}(g) y)_{l}
$$

On the other hand, the left invariance of $y_{l}$ leads to

$$
\left(c_{g}\right)_{*} y_{l}=\left(\rho_{g}^{-1} \circ \lambda_{g}\right)_{*} y_{l}=\left(\rho_{g}^{-1}\right)_{*}\left(\lambda_{g}\right)_{*} y_{l}=\left(\rho_{g}^{-1}\right)_{*} y_{l}
$$

Next we observe that $\Phi_{t}^{x_{l}}=\rho_{\exp _{G}(t x)}$ is the flow of the vector field $x_{l}$ (Lemma 9.2.4), so that Theorem 8.5.16 implies that

$$
\begin{aligned}
{\left[x_{l}, y_{l}\right] } & =\mathcal{L}_{x_{l}} y_{l}=\left.\frac{d}{d t}\right|_{t=0}\left(\Phi_{-t}^{x_{l}}\right)_{*} y_{l}=\left.\frac{d}{d t}\right|_{t=0}\left(c_{\exp _{G}(t x)}\right)_{*} y_{l} \\
& =\left.\frac{d}{d t}\right|_{t=0}\left(\operatorname{Ad}\left(\exp _{G}(t x)\right) y\right)_{l}
\end{aligned}
$$

Evaluating at 1, we get

$$
[x, y]=\left[x_{l}, y_{l}\right](\mathbf{1})=\left.\frac{d}{d t}\right|_{t=0} \operatorname{Ad}\left(\exp _{G}(t x)\right) y=\mathbf{L}(\operatorname{Ad})(x)(y)
$$

Combining Proposition 9.2.10 with Lemma 9.2.20, we obtain the important formula

$$
\text { Ado } \exp _{G}=\exp _{\mathrm{Aut}(\mathbf{L}(G))} \circ \mathrm{ad}
$$

i.e.,

$$
\operatorname{Ad}\left(\exp _{G}(x)\right)=e^{\operatorname{ad} x} \quad \text { for } \quad x \in \mathbf{L}(G)
$$

Lemma 9.2.21. For a Lie group $G$, the kernel of the adjoint representation $\mathrm{Ad}: G \rightarrow \operatorname{Aut}(\mathbf{L}(G))$, is given by

$$
Z_{G}\left(G_{0}\right):=\left\{g \in G:\left(\forall x \in G_{0}\right) g x=x g\right\}
$$

where $G_{0}$ is the connected component of the identity in $G$. If, in addition, $G$ is connected, then

$$
\operatorname{ker} \mathrm{Ad}=Z(G)
$$

Proof. Since $G_{0}$ is connected, the automorphism $\left.c_{g}\right|_{G_{0}}$ of $G_{0}$ is trivial if and only if $\mathbf{L}\left(c_{g}\right)=\operatorname{Ad}(g)$ is trivial. This implies the lemma.

\subsubsection*{Semidirect Products}
The easiest way to construct a new Lie group from two given Lie groups $G$ and $H$, is to endow the product manifold $G \times H$ with the multiplication

$$
\left(g_{1}, h_{1}\right)\left(g_{2}, h_{2}\right):=\left(g_{1} g_{2}, h_{1} h_{2}\right)
$$

The resulting group is called the direct product of the Lie groups $G$ and $H$. Here $G$ and $H$ can be identified with normal subgroups of $G \times H$ for which the multiplication map

$$
(G \times\{\mathbf{1}\}) \times(\{\mathbf{1}\} \times H) \rightarrow G \times H, \quad((g, \mathbf{1}),(\mathbf{1}, h)) \mapsto(g, \mathbf{1})(\mathbf{1}, h)=(g, h)
$$

is a diffeomorphism. Relaxing this condition in the sense that only one factor is assumed to be normal, leads to the concept of a semidirect product of Lie groups, introduced below.

Definition 9.2.22. Let $N$ and $G$ be Lie groups and $\alpha: G \rightarrow \operatorname{Aut}(N)$ be a group homomorphism defining a smooth action $(g, n) \mapsto \alpha(g)(n)$ of $G$ on $N$. Then the product manifold $N \times G$ is a group with respect to the product (cf. Lemma 2.2.3)

$$
(n, g)\left(n^{\prime}, g^{\prime}\right):=\left(n \alpha(g)\left(n^{\prime}\right), g g^{\prime}\right)
$$

and the inversion

$$
(n, g)^{-1}=\left(\alpha\left(g^{-1}\right)\left(n^{-1}\right), g^{-1}\right)
$$

Since multiplication and inversion are smooth, this group is a Lie group, called the semidirect product of $N$ and $G$ with respect to $\alpha$. It is denoted by $N \rtimes_{\alpha} G$.

On the manifold $G \times N$, we also obtain a Lie group structure by

$$
(g, n)\left(g^{\prime}, n^{\prime}\right):=\left(g g^{\prime}, \alpha\left(g^{\prime}\right)^{-1}(n) n^{\prime}\right)
$$

and this Lie group is denoted $G \ltimes_{\alpha} N$. It is easy to verify that the map

$$
\Phi: N \rtimes_{\alpha} G \rightarrow G \ltimes_{\alpha} N, \quad(n, g) \mapsto\left(g, \alpha(g)^{-1}(n)\right)
$$

is an isomorphism of Lie groups.\\
Remark 9.2.23. If $\widehat{G}:=N \rtimes_{\alpha} G$ is a semidirect product, then

$$
\pi: \widehat{G} \rightarrow G, \quad(n, g) \mapsto g, \quad \sigma: G \rightarrow \widehat{G}, \quad g \mapsto(\mathbf{1}, g)
$$

and $\iota: N \rightarrow \widehat{G}, n \mapsto(n, \mathbf{1})$ are morphisms of Lie groups with $\pi \circ \sigma=\operatorname{id}_{G}$ and $\iota$ is an isomorphism of $N$ onto the submanifold $\operatorname{ker} \pi$ of $\widehat{G}$ (cf. Remark 8.6.10).

Example 9.2.24. Let $G$ be a Lie group and $T(G)$ its tangent Lie group (Lemma 9.1.6). We have already seen that the map $G \times \mathbf{L}(G) \rightarrow T G,(g, x) \mapsto g \cdot x=0_{g} \cdot x$ is a diffeomorphism, and for similar reasons, the map $\mathbf{L}(G) \times G \rightarrow T G,(x, g) \mapsto x . g:=x \cdot 0_{g}$ is a diffeomorphism. In these coordinates, the multiplication is given by

$$
(x . g) \cdot\left(x^{\prime} \cdot g^{\prime}\right)=x \cdot 0_{g} \cdot x^{\prime} \cdot 0_{g^{\prime}}=x \cdot \operatorname{Ad}(g) x^{\prime} \cdot 0_{g} \cdot 0_{g^{\prime}}=\left(x+\operatorname{Ad}(g) x^{\prime}\right) \cdot g g^{\prime}
$$

This shows that the tangent bundle is a semidirect product

$$
T G \cong \mathbf{L}(G) \rtimes_{\mathrm{Ad}} G
$$

Similarly, the calculation

$$
(g \cdot x) \cdot\left(g^{\prime} \cdot x^{\prime}\right)=0_{g g^{\prime}} \cdot \operatorname{Ad}\left(g^{\prime}\right)^{-1} x \cdot x^{\prime}=\left(g g^{\prime}, \operatorname{Ad}\left(g^{\prime}\right)^{-1} x+x^{\prime}\right)
$$

shows that also

$$
T G \cong G \ltimes_{\mathrm{Ad}} \mathbf{L}(G)
$$

Proposition 9.2.25. The Lie algebra of the semidirect product group $N \rtimes_{\alpha} G$ is given by

$$
\mathbf{L}\left(N \rtimes_{\alpha} G\right) \cong \mathbf{L}(N) \rtimes_{\beta} \mathbf{L}(G)
$$

where $\beta: \mathbf{L}(G) \rightarrow \operatorname{der}(\mathbf{L}(N))$ is the derived representation of $\mathbf{L}(G)$ on $\mathbf{L}(N)$ corresponding to the representation of $G$ on $\mathbf{L}(N)$ given by $g . x:=\mathbf{L}(\alpha(g)) x$.

Proof. We identify $\mathbf{L}(N)$, resp., $\mathbf{L}(G)$, with a subspace of

$$
T_{\mathbf{1}}(N) \oplus T_{\mathbf{1}}(G)=T_{(\mathbf{1}, \mathbf{1})}(N \times G) \cong \mathbf{L}\left(N \rtimes_{\alpha} G\right)
$$

Since $N$ and $G$ are subgroups, the functoriality of $\mathbf{L}$ implies that $\mathbf{L}(G)$ and $\mathbf{L}(N)$ are Lie subalgebras of $\mathbf{L}\left(N \rtimes_{\alpha} G\right)$. The normal subgroup $N$ is the kernel of the projection $\pi: N \rtimes_{\alpha} G \rightarrow G$, so that our identification shows that $\mathbf{L}(N)=\operatorname{ker} \mathbf{L}(\pi)$ is an ideal of $\mathbf{L}\left(N \rtimes_{\alpha} G\right)$. This already implies

$$
\mathbf{L}\left(N \rtimes_{\alpha} G\right) \cong \mathbf{L}(N) \rtimes_{\beta} \mathbf{L}(G)
$$

for the homomorphism $\beta: \mathbf{L}(G) \rightarrow \operatorname{der}(\mathbf{L}(N))$, given by

$$
(\beta(x)(y), 0)=[(0, x),(y, 0)]
$$

To determine $\beta$ in terms of $\alpha$, we note that the smooth action of $G$ on $N$ by automorphisms induces a smooth action of $G$ on the tangent bundle $T(N)$, hence in particular on $T_{\mathbf{1}}(N) \cong \mathbf{L}(N)$. We thus obtain a representation $\pi: G \rightarrow \operatorname{Aut}(\mathbf{L}(N))$. In $N \rtimes_{\alpha} G$, we have $(\mathbf{1}, g)(n, \mathbf{1})(\mathbf{1}, g)^{-1}=(\alpha(g)(n), \mathbf{1})$, so that

$$
\pi(g) y=\mathbf{L}(\alpha(g)) y=\operatorname{Ad}(\mathbf{1}, g)(y, 0)
$$

Now Lemma 9.2.20 immediately shows that $\mathbf{L}(\pi) x=\operatorname{ad}(0, x)=\beta(x)$.

To form semidirect products, we need smooth actions of a Lie group $G$ on a Lie group $N$. The following lemma is a useful tool to verify smoothness in this context.

Lemma 9.2.26. Let $G$ and $N$ be Lie groups and $\alpha: G \rightarrow \operatorname{Aut}(N)$ be a group homomorphism.\\
(i) The action $\sigma(g, n):=\alpha(g)(n)$ of $G$ on $N$ is smooth if and only if\\
(a) $\mathbf{L} \circ \alpha: G \rightarrow \mathrm{Aut}(\mathbf{L}(N))$ is smooth as a map into $\mathrm{GL}(\mathbf{L}(N))$.\\
(b) All orbit maps $\sigma^{n}: G \rightarrow N, g \mapsto \alpha(g)(n)$ are smooth.\\
(ii) If $N$ is connected, then (a) implies (b).

Proof. (i) If $\sigma$ is smooth, then (b) is clearly satisfied. Moreover,

$$
\exp _{N} \circ \mathbf{L}(\alpha(g))=\alpha(g) \circ \exp _{N}
$$

holds for each $g \in G$, so that the map

$$
G \times \mathbf{L}(N) \rightarrow \mathbf{L}(N), \quad(g, x) \mapsto \mathbf{L}(\alpha(g)) x
$$

is smooth in an open neighborhood of $(\mathbf{1}, 0)$. This implies that $\mathbf{L} \circ \alpha: G \rightarrow \mathrm{GL}(\mathbf{L}(N))$ is smooth in an open neighborhood of $\mathbf{1}$, hence smooth because it is a group homomorphism.

Next we assume that (a) and (b) are satisfied. We have to show that for any pair $\left(g_{1}, n_{1}\right) \in G \times N$ the expression

$$
\sigma\left(g_{1} g_{2}, n_{1} n_{2}\right)=\alpha\left(g_{1}\right)\left(\alpha\left(g_{2}\right)\left(n_{1}\right) \alpha\left(g_{2}\right)\left(n_{2}\right)\right)
$$

is a smooth function of ( $g_{2}, n_{2}$ ) in a neighborhood of ( $\mathbf{1}, \mathbf{1}$ ). In view of the smoothness of $\alpha\left(g_{1}\right)$ and (b), it suffices to prove that $\sigma$ is smooth in an open neighborhood of $(\mathbf{1}, \mathbf{1})$. According to (a), this follows from

$$
\sigma\left(g, \exp _{N} x\right)=\alpha(g)\left(\exp _{N}(x)\right)=\exp _{N}(\mathbf{L}(\alpha(g)) x)
$$

(ii) Since $G$ acts by group automorphisms, the set

$$
S:=\left\{n \in N: \sigma^{n} \in C^{\infty}(G, N)\right\}
$$

is a subgroup. In view of (9.12), (a) implies that $S$ contains the image of $\exp _{N}$, hence all of $N$ (Lemma 9.2.9).

\subsubsection*{The Baker-Campbell-Dynkin-Hausdorff Formula}
In this subsection, we show that the formula

$$
\exp _{G}(x * y)=\exp _{G} x \exp _{G} y
$$

where $x * y$ for sufficiently small elements $x, y \in \mathfrak{g}=\mathbf{L}(G)$ is given by the Hausdorff series (cf. Proposition 3.4.5), also holds for the exponential function of a general Lie group $G$ with Lie algebra $\mathfrak{g}$.

The Maurer-Cartan form $\kappa_{G} \in \Omega^{1}(G, \mathfrak{g})$ is the unique left invariant 1form on $G$ with $\kappa_{G, \mathbf{1}}=\operatorname{id}_{\mathfrak{g}}$, i.e., $\kappa_{G}(v)=g^{-1} . v$ for $v \in T_{g}(G)$. The logarithmic derivative of a smooth function $f: M \rightarrow G$ is now defined as the pull-back of the Maurer-Cartan form:

$$
\delta(f):=f^{*} \kappa_{G} \in \Omega^{1}(M, \mathfrak{g})
$$

For each $v \in T_{m}(M)$, we then have

$$
\delta(f)_{m} v=f(m)^{-1} \cdot T_{m}(f) v
$$

If $\alpha \in \Omega^{1}(M, \mathfrak{g})$ is a Lie algebra-valued 1-form and $f: M \rightarrow G$ a smooth function, then we define $\operatorname{Ad}(f) \alpha$ pointwise by $(\operatorname{Ad}(f) \alpha)_{m}:=\operatorname{Ad}(f(m)) \alpha_{m}$.

Lemma 9.2.27. For two smooth maps $f, h: M \rightarrow G$, the logarithmic derivative of the pointwise products $f h$ and $f h^{-1}$ is given by the\\
(1) Product Rule: $\delta(f h)=\delta(h)+\operatorname{Ad}\left(h^{-1}\right) \delta(f)$, and the\\
(2) Quotient Rule: $\delta\left(f h^{-1}\right)=\operatorname{Ad}(h)(\delta(f)-\delta(h))$.

Proof. Writing $f g=m_{G} \circ(f, g)$, we obtain from

$$
T_{(a, b)}\left(m_{G}\right)(v, w)=v \cdot b+a \cdot w
$$

for $a, b \in G$ and $v, w \in \mathbf{L}(G) \subseteq T G$ the relation

$$
T(f h)=T\left(m_{G}\right) \circ(T(f), T(h))=T(f) \cdot h+f \cdot T(h): T(M) \rightarrow T(G)
$$

where $f \cdot T(h)$, resp., $T(f) \cdot h$, refers to the pointwise product in the group $T(G)$ containing $G$ as the zero section (Lemma 9.1.6). This immediately leads to the Product Rule

$$
\begin{aligned}
\delta(f h) & =(f h)^{-1} \cdot(T(f) \cdot h+f \cdot T(h))=h^{-1} \cdot(\delta(f) \cdot h)+\delta(h) \\
& =\operatorname{Ad}(h)^{-1} \delta(f)+\delta(h)
\end{aligned}
$$

For $h=f^{-1}$, we then obtain

$$
0=\delta\left(f f^{-1}\right)=\operatorname{Ad}(f) \delta(f)+\delta\left(f^{-1}\right)
$$

hence $\delta\left(f^{-1}\right)=-\operatorname{Ad}(f) \delta(f)$. This in turn leads to

$$
\delta\left(f h^{-1}\right)=\operatorname{Ad}(h) \delta(f)+\delta\left(h^{-1}\right)=\operatorname{Ad}(h) \delta(f)-\operatorname{Ad}(h) \delta(h)
$$

which is the Quotient Rule.\\
For any $g \in G$ and a smooth function $f: M \rightarrow G$, the function $\lambda_{g} \circ f$ has the same logarithmic derivative as $f$ because $\left(\lambda_{g} \circ f\right)^{*} \kappa_{G}=f^{*} \lambda_{g}^{*} \kappa_{G}=f^{*} \kappa_{G}$. The following lemma provides a converse, which is a very convenient tool.

\section*{Lemma 9.2.28 (Uniqueness Lemma for Logarithmic Derivatives).}
Let $G$ be a Lie group and $M$ a connected manifold. If $f, h \in C^{\infty}(M, G)$ satisfy $\delta(f)=\delta(h)$, then $h=\lambda_{g} \circ f$ holds for some $g \in G$.

Proof. The Quotient Rule leads to

$$
\delta\left(h f^{-1}\right)=\delta\left(f^{-1}\right)+\operatorname{Ad}(f) \delta(h)=-\operatorname{Ad}(f) \delta(f)+\operatorname{Ad}(f) \delta(f)=0 .
$$

Hence $T_{m}\left(h f^{-1}\right)=0$ in each $m \in M$, so that $h \cdot f^{-1}$ is constant equal to some $g \in G$. This means that $h=\lambda_{g} \circ f$.

Proposition 9.2.29. The logarithmic derivative of $\exp _{G}$ is given by

$$
\delta\left(\exp _{G}\right)(x)=\frac{\mathbf{1}-e^{-\operatorname{ad} x}}{\operatorname{ad} x}: \mathfrak{g} \rightarrow \mathfrak{g}
$$

where the fraction on the right means $\Phi(\operatorname{ad} x)$ for the entire function

$$
\Phi(z):=\frac{1-e^{-z}}{z}=\sum_{k=1}^{\infty} \frac{(-z)^{k-1}}{k!}
$$

The series $\Phi(\operatorname{ad} x)$ converges for each $x \in \mathfrak{g}$.\\
Proof. Fix $t, s \in \mathbb{R}$. Then the smooth functions $f, f_{t}, f_{s}: \mathbf{L}(G) \rightarrow G$, given by

$$
f(x):=\exp _{G}((t+s) x), \quad f_{t}(x):=\exp _{G}(t x) \quad \text { and } \quad f_{s}(x):=\exp _{G}(s x)
$$

satisfy $f=f_{t} f_{s}$ pointwise on $\mathbf{L}(G)$. The Product Rule (Lemma 9.2.27) therefore implies that

$$
\delta(f)=\delta\left(f_{s}\right)+\operatorname{Ad}\left(f_{s}\right)^{-1} \delta\left(f_{t}\right)
$$

For the smooth curve $\psi: \mathbb{R} \rightarrow \mathbf{L}(G), \psi(t):=\delta\left(\exp _{G}\right)_{t x}(t y)$, we now obtain

$$
\begin{aligned}
\psi(t+s) & =\delta(f)_{x}(y)=\delta\left(f_{s}\right)_{x}(y)+\operatorname{Ad}\left(f_{s}\right)^{-1} \delta\left(f_{t}\right)_{x}(y) \\
& =\psi(s)+\operatorname{Ad}\left(\exp _{G}(-s x)\right) \psi(t)
\end{aligned}
$$

We have $\psi(0)=0$ and

$$
\psi^{\prime}(0)=\lim _{t \rightarrow 0} \delta\left(\exp _{G}\right)_{t x}(y)=\delta\left(\exp _{G}\right)_{0}(y)=y
$$

so that taking derivatives with respect to $t$ at 0 leads to

$$
\psi^{\prime}(s)=\operatorname{Ad}\left(\exp _{G}(-s x)\right) y=e^{-\operatorname{ad}(s x)} y
$$

Now the assertion follows by integration from

$$
\delta\left(\exp _{G}\right)_{x}(y)=\psi(1)=\int_{0}^{1} \psi^{\prime}(s) d s
$$

and $\int_{0}^{1} e^{-s \operatorname{ad} x} d s=\sum_{k=0}^{\infty} \frac{(-\operatorname{ad} x)^{k}}{(k+1)!}=\Phi(\operatorname{ad} x)$, which we saw already in the proof of Proposition 3.4.2.

Definition 9.2.30. We have seen in Proposition 9.2.29 that

$$
\delta\left(\exp _{G}\right)(x)=\frac{\mathbf{1}-e^{-\mathrm{ad} x}}{\operatorname{ad} x}
$$

From Exercise 9.2.13, we now derive that

$$
\operatorname{sing} \exp (\mathfrak{g}):=\{x \in \mathfrak{g}: \operatorname{Spec}(\operatorname{ad} x) \cap 2 \pi i \mathbb{Z} \nsubseteq\{0\}\}
$$

is the set of singular points of $\exp _{G}$, and

$$
\operatorname{regexp}(\mathfrak{g}):=\{x \in \mathfrak{g}: \operatorname{Spec}(\operatorname{ad} x) \cap 2 \pi i \mathbb{Z} \subseteq\{0\}\}
$$

is the set of regular points of the exponential function. We note in particular that these sets only depend on the Lie algebra $\mathfrak{g}$, and not on the group $G$.

Lemma 9.2.31. Let $G$ be a Lie group with Lie algebra $\mathfrak{g}$. Then the following assertions hold for $x \in \mathfrak{g}$ :\\
(a) If $x \in \operatorname{regexp}(\mathfrak{g})$ and $\exp _{G} x=\exp _{G} y$ for some $y \in \mathfrak{g}$, then

$$
[x, y]=0 \quad \text { and } \quad \exp _{G}(x-y)=\mathbf{1}
$$

(b) If $x \in \operatorname{sing} \exp (\mathfrak{g})$, then $\exp _{G}$ is not injective on any neighborhood of $x$.

Proof. (a) The whole one-parameter group $\exp _{G}(\mathbb{R} y)$ commutes with $\exp _{G} y=\exp _{G} x$, which leads to

$$
\exp _{G}(x)=c_{\exp _{G} t y}\left(\exp _{G} x\right)=\exp _{G}\left(e^{t \operatorname{ad} y} x\right) \quad \text { for } \quad t \in \mathbb{R}
$$

Consequently,

$$
\left.\frac{d}{d t}\right|_{t=0} e^{t \operatorname{ad} y} x=[y, x] \in \operatorname{ker} T_{x}(\exp )=\{0\}
$$

and thus $[x, y]=0$. This in turn leads to

$$
\exp _{G}(x-y)=\exp _{G}(x) \exp _{G}(-y)=\mathbf{1}
$$

(Lemma 9.2.8).\\
(b) There exists some $0 \neq y \in \mathfrak{g}$ with $T_{x}(\exp ) y=0$, resp., $\Phi(\operatorname{ad} x) y=0$ (Proposition 9.2.29). We consider the one-parameter group

$$
\alpha: \mathbb{R} \rightarrow \operatorname{Aut}(G), \quad t \mapsto c_{\exp _{G} t y}
$$

of automorphisms of $G$. It is generated by the vector field

$$
\mathcal{X}(g)=\left.\frac{d}{d t}\right|_{t=0} \exp _{G}(t y) g \exp _{G}(-t y)
$$

which has at $\exp _{G} x$ the value

$$
\begin{aligned}
\mathcal{X}(\exp x) & =\left.\frac{d}{d t}\right|_{t=0} c_{\exp _{G}(t y)} \exp _{G}(x)=\left.\frac{d}{d t}\right|_{t=0} \exp _{G}\left(e^{\operatorname{ad} t y} x\right)=T_{x}(\exp )[y, x] \\
& =T_{\mathbf{1}}\left(\lambda_{\exp _{G} x}\right) \Phi(\operatorname{ad} x)(-\operatorname{ad} x) y=T_{\mathbf{1}}\left(\lambda_{\exp _{G} x}\right)(-\operatorname{ad} x) \Phi(\operatorname{ad} x) y=0
\end{aligned}
$$

(Proposition 9.2.29). This implies that $\alpha(t)\left(\exp _{G} x\right)=\exp _{G}\left(e^{\mathrm{ad} t y} x\right)= \exp _{G} x$ for every $t \in \mathbb{R}$. In view of $y \neq 0$ and $\Phi(\operatorname{ad} x) y=0$, we have $[y, x] \neq 0$ (Exercise 9.2.12), so that the curve $e^{\text {ad } t y} x$ is not constant. Hence $\exp _{G}$ is not injective on any neighborhood of $x$.

Let $U \subseteq \mathfrak{g}$ be a convex 0-neighborhood for which $\left.\exp _{G}\right|_{U}$ is a diffeomorphism onto an open subset of $G$ and $V \subseteq U$ a smaller convex open 0neighborhood with $\exp _{G} V \exp _{G} V \subseteq \exp _{G} U$. Put $\log _{U}:=\left(\left.\exp _{G}\right|_{U}\right)^{-1}$ and define for $x, y \in V$

$$
x * y:=\log _{U}\left(\exp _{G} x \exp _{G} y\right),
$$

which defines a smooth map $V \times V \rightarrow U$. Fix $x, y \in V$. Then the smooth curve $F(t):=x * t y$ satisfies $\exp _{G} F(t)=\exp _{G}(x) \exp _{G}(t y)$, so that the logarithmic derivative of this curve is

$$
\delta\left(\exp _{G}\right)_{F(t)} F^{\prime}(t)=y .
$$

We now choose $U$ so small that the power series $\Psi(z)=\frac{z \log z}{z-1}$ from Lemma 3.4.3 satisfies

$$
\Psi\left(e^{\operatorname{ad} z}\right) \Phi(\operatorname{ad} z)=\operatorname{id}_{\mathfrak{g}} \quad \text { for } \quad z \in U
$$

(cf. Proposition 3.1.6). For $z=F(t)$, we then arrive with Proposition 9.2.29 at

$$
F^{\prime}(t)=\Psi\left(e^{\operatorname{ad} F(t)}\right) y
$$

Now the same arguments as in Propositions 3.4.4 and 3.4.5 imply that

$$
x * y=F(1)=x+y+\frac{1}{2}[x, y]+\cdots
$$

is given by the convergent Hausdorff series:\\
Proposition 9.2.32. If $G$ is a Lie group, then there exists a convex 0neighborhood $V \subseteq \mathfrak{g}$ such that for $x, y \in V$ the Hausdorff series

$$
\begin{aligned}
x * y & :=x+\sum_{\substack{k, m \geq 0 \\
p_{i}+q_{i}>0}} \frac{(-1)^{k}}{(k+1)\left(q_{1}+\cdots+q_{k}+1\right)} \\
& \cdot \frac{(\operatorname{ad} x)^{p_{1}}(\operatorname{ad} y)^{q_{1}} \cdots(\operatorname{ad} x)^{p_{k}}(\operatorname{ad} y)^{q_{k}}(\operatorname{ad} x)^{m}}{p_{1}!q_{1}!\cdots p_{k}!q_{k}!m!} y
\end{aligned}
$$

converges and satisfies

$$
\exp _{G}(x * y)=\exp _{G}(x) \exp _{G}(y)
$$

\subsubsection*{Exercises for Section 9.2}
Exercise 9.2.1. Let $G$ be a connected Lie group and $x \in \mathfrak{g}=\mathbf{L}(G)$. Show that the corresponding left invariant vector field $x_{l} \in \mathcal{V}(G)$ is biinvariant, i.e., also invariant under all right multiplications, if and only if $x \in \mathfrak{z}(\mathfrak{g})$.

Exercise 9.2.2. Let $f_{1}, f_{2}: G \rightarrow H$ be two group homomorphisms. Show that the pointwise product

$$
f_{1} f_{2}: G \rightarrow H, \quad g \mapsto f_{1}(g) f_{2}(g)
$$

is a homomorphism if and only if $f_{1}(G)$ commutes with $f_{2}(G)$.\\
Exercise 9.2.3. Let $M$ be a manifold and $V$ a finite-dimensional vector space with a basis $\left(b_{1}, \ldots, b_{n}\right)$. Let $f: M \rightarrow \mathrm{GL}(V)$ be a map. Show that the following are equivalent:\\
(1) $f$ is smooth.\\
(2) For each $v \in V$, the map $f_{v}: M \rightarrow V, m \mapsto f(m) v$ is smooth.\\
(3) For each $i$, the map $f: M \rightarrow V, m \mapsto f(m) b_{i}$ is smooth.

Exercise 9.2.4. A vector field $X$ on a Lie group $G$ is called right invariant if for each $g \in G$ the vector field $\left(\rho_{g}\right)_{*} X=T\left(\rho_{g}\right) \circ X \circ \rho_{g}^{-1}$ coincides with $X$. We write $\mathcal{V}(G)^{r}$ for the set of right invariant vector fields on $G$. Show that:\\
(1) The evaluation map $\mathrm{ev}_{\mathbf{1}}: \mathcal{V}(G)^{r} \rightarrow T_{\mathbf{1}}(G)$ is a linear isomorphism.\\
(2) If $X$ is right invariant, then there exists a unique $x \in T_{\mathbf{1}}(G)$ such that $X(g)=x_{r}(g):=T_{\mathbf{1}}\left(\rho_{g}\right) x=x \cdot 0_{g}($ w.r.t. multiplication in $T(G))$.\\
(3) If $X$ is right invariant, then $\widetilde{X}:=\left(\iota_{G}\right)_{*} X:=T\left(\iota_{G}\right) \circ X \circ \iota_{G}^{-1}$ is left invariant and vice versa.\\
(4) Show that $\left(\iota_{G}\right)_{*} x_{r}=-x_{l}$ and $\left[x_{r}, y_{r}\right]=-[x, y]_{r}$ for $x, y \in T_{\mathbf{1}}(G)$.\\
(5) Show that each right invariant vector field is complete and determine its flow.

Exercise 9.2.5. Let $M$ be a smooth manifold, $\varphi \in \operatorname{Diff}(M)$ and $X \in \mathcal{V}(M)$. Show that the following are equivalent:\\
(1) $\varphi$ commutes with the flow maps $\Phi_{t}^{X}: M_{t} \rightarrow M$ of $X$, i.e., each set $M_{t}$ is $\varphi$-invariant and $\Phi_{t}^{X} \circ \varphi=\varphi \circ \Phi_{t}^{X}$ holds on $M_{t}$.\\
(2) For each integral curve $\gamma: I \rightarrow M$ of $X$ the curve $\varphi \circ \gamma$ is also an integral curve of $X$.\\
(3) $X=\varphi_{*} X=T(\varphi) \circ X \circ \varphi^{-1}$, i.e., $X$ is $\varphi$-invariant.

Exercise 9.2.6. Let $G$ be a Lie group. Show that any map $\varphi: G \rightarrow G$ commuting with all left multiplications $\lambda_{g}, g \in G$, is a right multiplication.

Exercise 9.2.7. Let $X, Y \in \mathcal{V}(M)$ be two commuting complete vector fields, i.e., $[X, Y]=0$. Show that the vector field $X+Y$ is complete and that its flow is given by

$$
\Phi_{t}^{X+Y}=\Phi_{t}^{X} \circ \Phi_{t}^{Y} \quad \text { for all } \quad t \in \mathbb{R}
$$

Exercise 9.2.8. Let $V$ be a finite-dimensional vector space and $\mu_{t}(v):=t v$ for $t \in \mathbb{R}^{\times}$. Show that:\\
(1) A vector field $X \in \mathcal{V}(V)$ is linear if and only if $\left(\mu_{t}\right)_{*} X=X$ holds for all $t \in \mathbb{R}^{\times}$.\\
(2) A diffeomorphism $\varphi \in \operatorname{Diff}(V)$ is linear if and only if it commutes with all the maps $\mu_{t}, t \in \mathbb{R}^{\times}$.

Exercise 9.2.9. Let $G$ be a connected Lie group and $H$ a Lie group. For a smooth map $f: G \rightarrow H$, the following are equivalent:\\
(1) $f$ is a homomorphism.\\
(2) $f(\mathbf{1})=\mathbf{1}$ and there exists a homomorphism $\psi: \mathbf{L}(G) \rightarrow \mathbf{L}(H)$ of Lie algebras with $\delta(f)=\psi \circ \kappa_{G}$.

Exercise 9.2.10. Let $G$ be a connected Lie group, $H$ a Lie group and $\alpha: G \rightarrow H$ a homomorphism defining a smooth action of $G$ on $H$. For a smooth map $f: G \rightarrow H$, the following are equivalent:\\
(1) $f$ is a crossed homomorphism, i.e.,

$$
f(x y)=f(x) \cdot \alpha(x)(f(y)) \quad \text { for } \quad x, y \in G .
$$

(2) $f(\mathbf{1})=\mathbf{1}$ and $\delta(f) \in \Omega^{1}(G, \mathbf{L}(H))$ is equivariant, i.e.,

$$
\lambda_{g}^{*} \delta(f)=\mathbf{L}(\alpha(g)) \circ \delta(f)
$$

for each $g \in G$.\\
(3) $\left(f, \mathrm{id}_{G}\right): G \rightarrow H \rtimes_{\alpha} G$ is a homomorphism.

Exercise 9.2.11. No one-parameter group $\gamma: \mathbb{R} \rightarrow \mathrm{SU}_{2}(\mathbb{C})$ is injective, in particular, the image of $\gamma(\mathbb{R})$ is a circle group.

Exercise 9.2.12. (i) Let $A$ be a semisimple endomorphism of the complex vector space $V$ and let $h(z):=\sum_{n=0}^{\infty} a_{n} z^{n}$ be a complex power series converging on $\mathbb{C}$. We define $h(A):=\sum_{n=0}^{\infty} a_{n} A^{n}$. Then

$$
\operatorname{ker} h(A)=\bigoplus_{z \in h^{-1}(0) \cap \operatorname{Spec}(A)} \operatorname{ker}(A-z 1)
$$

(ii) Let $A$ be a semisimple endomorphism of the real vector space $V$, and let $A_{\mathbb{C}}: V_{\mathbb{C}} \rightarrow V_{\mathbb{C}}$ be its complex linear extension. Then $\operatorname{Spec}(A):= \operatorname{Spec}\left(A_{\mathbb{C}}\right)$ decomposes into the subsets

$$
S_{\mathrm{re}}:=\operatorname{Spec}(A) \cap \mathbb{R} \quad \text { and } \quad S_{\mathrm{im}}:=\operatorname{Spec}(A) \backslash S_{\mathrm{re}} .
$$

Now let $h$ be as above and assume, in addition, that $h(\bar{z})=\overline{h(z)}$. Then $h(A) V \subseteq V$ and

$$
\begin{aligned}
\operatorname{ker} h(A) & =\bigoplus_{z \in h^{-1}(0) \cap S_{\mathrm{re}}} \operatorname{ker}(A-z 1) \\
& \oplus \bigoplus_{x+i y \in h^{-1}(0) \cap S_{\mathrm{im}}, y>0} \operatorname{ker}\left(A^{2}-2 x A+\left(x^{2}+y^{2}\right)\right)
\end{aligned}
$$

Exercise 9.2.13. Let $A \in \operatorname{End}(V)$, where $V$ is a finite dimensional real vector space and

$$
f(z):=\frac{1-e^{-z}}{z}=\sum_{k=0}^{\infty} \frac{(-1)^{k} z^{k}}{(k+1)!}
$$

Show that $f(A)$ is invertible if and only if

$$
\operatorname{Spec}(A) \cap 2 \pi i \mathbb{Z} \subseteq\{0\}
$$

\subsection*{Closed Subgroups of Lie Groups and Their Lie Algebras}
In this section, we show that closed subgroups of Lie groups are always Lie groups and that, for a closed subgroup $H$ of $G$, its Lie algebra can be computed as

$$
\mathbf{L}(H)=\left\{x \in \mathbf{L}(G): \exp _{G}(\mathbb{R} x) \subseteq H\right\}
$$

This makes it particularly easy to verify that concrete groups of matrices are Lie groups and to determine their algebras.

\subsubsection*{The Lie Algebra of a Closed Subgroup}
Definition 9.3.1. Let $G$ be a Lie group and $H \leq G$ a closed subgroup. We define the set

$$
\mathbf{L}^{e}(H):=\left\{x \in \mathbf{L}(G): \exp _{G}(\mathbb{R} x) \subseteq H\right\}
$$

and observe that $\mathbb{R} \mathbf{L}^{e}(H) \subseteq \mathbf{L}^{e}(H)$ follows immediately from the definition.\\
Note that, for each $x \in \mathbf{L}(G)$, the set

$$
\left\{t \in \mathbb{R}: \exp _{G}(t x) \in H\right\}=\gamma_{x}^{-1}(H)
$$

is a closed subgroup of $\mathbb{R}$, hence either discrete cyclic or equal to $\mathbb{R}$ (cf. Exercise 9.3.4).

Example 9.3.2. We consider the Lie group $G:=\mathbb{R} \times \mathbb{T}$ (the cylinder) with Lie algebra $\mathbf{L}(G) \cong \mathbb{R}^{2}$ (Exercise 9.1.2) and the exponential function

$$
\exp _{G}(x, y)=\left(x, e^{2 \pi i y}\right)
$$

For the closed subgroup $H:=\mathbb{R} \times\{\mathbf{1}\}$, we then see that $(x, y) \in \mathbf{L}^{e}(H)$ is equivalent to $y=0$, but $\exp _{G}^{-1}(H)=\mathbb{R} \times \mathbb{Z}$.

Proposition 9.3.3. If $H \leq G$ is a closed subgroup of the Lie group $G$, then $\mathbf{L}^{e}(H)$ is a real Lie subalgebra of $\mathbf{L}(G)$.

Proof. Let $x, y \in \mathbf{L}^{e}(H)$. For $k \in \mathbb{N}$, we then have $\exp _{G} \frac{1}{k} x, \exp _{G} \frac{1}{k} y \in H$, and with the Product Formula (Proposition 9.2.14), we get

$$
\exp _{G}(x+y)=\lim _{k \rightarrow \infty}\left(\exp _{G} \frac{x}{k} \exp _{G} \frac{y}{k}\right)^{k} \in H
$$

because $H$ is closed. Therefore, $\exp _{G}(x+y) \in H$, and $\mathbb{R} \mathbf{L}^{e}(H)=\mathbf{L}^{e}(H)$ now implies $\exp _{G}(\mathbb{R}(x+y)) \subseteq H$, hence $x+y \in \mathbf{L}^{e}(H)$.

Similarly, we use the Commutator Formula to get

$$
\exp _{G}[x, y]=\lim _{k \rightarrow \infty}\left(\exp _{G} \frac{x}{k} \exp _{G} \frac{y}{k} \exp _{G}-\frac{x}{k} \exp _{G}-\frac{y}{k}\right)^{k^{2}} \in H
$$

hence $\exp _{G}([x, y]) \in H$, and $\mathbb{R} \mathbf{L}^{e}(H)=\mathbf{L}^{e}(H)$ yields $[x, y] \in \mathbf{L}^{e}(H)$.

\subsubsection*{The Closed Subgroup Theorem and Its Consequences}
As we shall see below in the Initial Subgroup Theorem 9.6.13, $\mathbf{L}^{e}(H)$ is a Lie algebra for any subgroup $H$ of $G$. We now address more detailed information on closed subgroups of Lie groups. We start with three key lemmas providing the main information for the proof of the Closed Subgroup Theorem.

Lemma 9.3.4. Let $W \subseteq \mathbf{L}(G)$ be an open 0 -neighborhood for which $\left.\exp _{G}\right|_{W}$ is a diffeomorphism and $\log _{W}: \exp _{G}(W) \rightarrow W$ its inverse function. Further, let $H \subseteq G$ be a closed subgroup and $\left(g_{k}\right)_{k \in \mathbb{N}}$ be a sequence in $H \cap \exp _{G}(W)$ with $g_{k} \neq \mathbf{1}$ for all $k \in \mathbb{N}$ and $g_{k} \rightarrow \mathbf{1}$. We put $y_{k}:=\log _{W} g_{k}$ and fix a norm $\|\cdot\|$ on $\mathbf{L}(G)$. Then every cluster point of the sequence $\left\{\frac{y_{k}}{\left\|y_{k}\right\|}: k \in \mathbb{N}\right\}$ is contained in $\mathbf{L}^{e}(H)$.

Proof. Let $x$ be such a cluster point. Replacing the original sequence by a subsequence, we may assume that

$$
x_{k}:=\frac{y_{k}}{\left\|y_{k}\right\|} \rightarrow x \in \mathbf{L}(G)
$$

Note that this implies $\|x\|=1$. Let $t \in \mathbb{R}$ and put $p_{k}:=\frac{t}{\left\|y_{k}\right\|}$. Then $t x_{k}= p_{k} y_{k}$ and $y_{k} \rightarrow \log _{W} \mathbf{1}=0$, so that

$$
\exp _{G}(t x)=\lim _{k \rightarrow \infty} \exp _{G}\left(t x_{k}\right)=\lim _{k \rightarrow \infty} \exp _{G}\left(p_{k} y_{k}\right)
$$

and

$$
\exp _{G}\left(p_{k} y_{k}\right)=\exp _{G}\left(y_{k}\right)^{\left[p_{k}\right]} \exp _{G}\left(\left(p_{k}-\left[p_{k}\right]\right) y_{k}\right)
$$

where $\left[p_{k}\right]=\max \left\{l \in \mathbb{Z}: l \leq p_{k}\right\}$ is the Gauß function. We then have

$$
\left\|\left(p_{k}-\left[p_{k}\right]\right) y_{k}\right\| \leq\left\|y_{k}\right\| \rightarrow 0
$$

and

$$
\exp _{G}(t x)=\lim _{k \rightarrow \infty}\left(\exp _{G} y_{k}\right)^{\left[p_{k}\right]}=\lim _{k \rightarrow \infty} g_{k}^{\left[p_{k}\right]} \in H
$$

because $H$ is closed. This implies $x \in \mathbf{L}^{e}(H)$.\\
Lemma 9.3.5. Let $H \subseteq G$ be a closed subgroup and $E \subseteq \mathbf{L}(G)$ be a vector subspace complementing $\mathbf{L}^{e}(H)$. Then there exists a 0 -neighborhood $U_{E} \subseteq E$ with

$$
H \cap \exp _{G}\left(U_{E}\right)=\{\mathbf{1}\}
$$

Proof. We argue by contradiction. If a neighborhood $U_{E}$ with the required properties does not exist, then for each compact convex 0 -neighborhood $V_{E} \subseteq E$ we have for each $k \in \mathbb{N}$ :

$$
\left(\exp _{G} \frac{1}{k} V_{E}\right) \cap H \neq\{\boldsymbol{1}\}
$$

For each $k \in \mathbb{N}$, we therefore find $y_{k} \in V_{E}$ with $\mathbf{1} \neq g_{k}:=\exp _{G}\left(\frac{y_{k}}{k}\right) \in H$. Now the compactness of $V_{E}$ implies that the sequence $\left(y_{k}\right)_{k \in \mathbb{N}}$ is bounded, so that $\frac{y_{k}}{k} \rightarrow 0$, which implies $g_{k} \rightarrow \mathbf{1}$. Now let $x \in E$ be a cluster point of the sequence $\frac{y_{k}}{\left\|y_{k}\right\|}$ which lies in the compact set $S_{E}:=\{z \in E:\|z\|=1\}$. According to Lemma 9.3.4, we have $x \in \mathbf{L}^{e}(H) \cap E=\{0\}$ because $g_{k} \in H \cap W$ for $k$ sufficiently large. We arrive at a contradiction to $\|x\|=1$. This proves the lemma.

Lemma 9.3.6. Let $E, F \subseteq \mathbf{L}(G)$ be vector subspaces with $E \oplus F=\mathbf{L}(G)$. Then the map

$$
\Phi: E \times F \rightarrow G, \quad(x, y) \mapsto \exp _{G}(x) \exp _{G}(y),
$$

restricts to a diffeomorphism of a neighborhood of $(0,0)$ to an open 1neighborhood in $G$.

Proof. The Chain Rule implies that

$$
\begin{aligned}
T_{(0,0)}(\Phi)(x, y) & =T_{(\mathbf{1}, \mathbf{1})}\left(m_{G}\right) \circ\left(\left.T_{0}\left(\exp _{G}\right)\right|_{E},\left.T_{0}\left(\exp _{G}\right)\right|_{F}\right)(x, y) \\
& =T_{(\mathbf{1}, \mathbf{1})}\left(m_{G}\right)(x, y)=x+y
\end{aligned}
$$

Since the addition map $E \times F \rightarrow \mathbf{L}(G) \cong T_{\mathbf{1}}(G)$ is bijective, the Inverse Function Theorem implies that $\Phi$ restricts to a diffeomorphism of an open neighborhood of ( 0,0 ) in $E \times F$ onto an open neighborhood of $\mathbf{1}$ in $G$.

Theorem 9.3.7 (Closed Subgroup Theorem). Let $H$ be a closed subgroup of the Lie group $G$. Then the following assertions hold:\\
(i) Each 0 -neighborhood in $\mathbf{L}^{e}(H)$ contains an open 0 -neighborhood $V$ such that $\left.\exp _{G}\right|_{V}: V \rightarrow \exp _{G}(V)$ is a homeomorphism onto an open subset of $H$.\\
(ii) $H$ is a submanifold of $G$ and $m_{H}:=\left.m_{G}\right|_{H \times H}$ induces a Lie group structure on $H$ such that the inclusion map $\iota_{H}: H \rightarrow G$ is a morphism of Lie groups for which $\mathbf{L}\left(\iota_{H}\right): \mathbf{L}(H) \rightarrow \mathbf{L}(G)$ is an isomorphism of $\mathbf{L}(H)$ onto $\mathbf{L}^{e}(H)$.\\
(iii) Let $E \subseteq \mathbf{L}(G)$ be a vector space complement of $\mathbf{L}^{e}(H)$. Then there exists an open 0 -neighborhood $V_{E} \subseteq E$ such that

$$
\varphi: V_{E} \times H \rightarrow \exp _{G}\left(V_{E}\right) H, \quad(x, h) \mapsto \exp _{G}(x) h
$$

is a diffeomorphism onto an open subset of $G$.\\
In view of (ii) above, we shall always identify $\mathbf{L}(H)$ with the subalgebra $\mathbf{L}^{e}(H)$ if $H$ is a closed subgroup of $G$.

Proof. (i) Let $E \subseteq \mathbf{L}(G)$ be a vector space complement of the subspace $\mathbf{L}^{e}(H)$ of $\mathbf{L}(G)$ and define

$$
\Phi: E \times \mathbf{L}^{e}(H) \rightarrow G, \quad(x, y) \mapsto \exp _{G} x \exp _{G} y
$$

According to Lemma 9.3.6, there exist open 0-neighborhoods $U_{E} \subseteq E$ and $U_{H} \subseteq \mathbf{L}^{e}(H)$ such that

$$
\Phi_{1}:=\left.\Phi\right|_{U_{E} \times U_{H}}: U_{E} \times U_{H} \rightarrow \exp _{G}\left(U_{E}\right) \exp _{G}\left(U_{H}\right)
$$

is a diffeomorphism onto an open 1-neighborhood in $G$. In view of Lemma 9.3.5, we may even choose $U_{E}$ so small that $\exp _{G}\left(U_{E}\right) \cap H=\{\mathbf{1}\}$.

Since $\exp _{G}\left(U_{H}\right) \subseteq H$, the condition

$$
g=\exp _{G} x \exp _{G} y \in H \cap\left(\exp _{G}\left(U_{E}\right) \exp _{G}\left(U_{H}\right)\right)
$$

implies $\exp _{G} x=g\left(\exp _{G} y\right)^{-1} \in H \cap \exp _{G} U_{E}=\{\boldsymbol{1}\}$. Therefore,

$$
H \supseteq \exp _{G}\left(U_{H}\right)=H \cap\left(\exp _{G}\left(U_{E}\right) \exp _{G}\left(U_{H}\right)\right)
$$

is an open $\mathbf{1}$-neighborhood in $H$. This proves (i).\\
(ii) Let $\Phi_{1}, U_{E}$ and $U_{H}$ be as in (i). For $h \in H$, the set $U_{h}:=\lambda_{h}\left(\operatorname{im}\left(\Phi_{1}\right)\right)= h \operatorname{im}\left(\Phi_{1}\right)$ is an open neighborhood of $h$ in $G$. Moreover, the map

$$
\varphi_{h}: U_{h} \rightarrow E \oplus \mathbf{L}^{e}(H)=\mathbf{L}(G), \quad x \mapsto \Phi_{1}^{-1}\left(h^{-1} x\right)
$$

is a diffeomorphism onto the open subset $U_{E} \times U_{H}$ of $\mathbf{L}(G)$, and we have

$$
\begin{aligned}
\varphi_{h}\left(U_{h} \cap H\right) & =\varphi_{h}\left(h \operatorname{im}\left(\Phi_{1}\right) \cap H\right)=\varphi_{h}\left(h\left(\operatorname{im}\left(\Phi_{1}\right) \cap H\right)\right) \\
& =\varphi_{h}\left(h \exp _{G}\left(U_{H}\right)\right)=\{0\} \times U_{H}=\left(U_{E} \times U_{H}\right) \cap\left(\{0\} \times \mathbf{L}^{e}(H)\right) .
\end{aligned}
$$

Therefore, the family $\left(\varphi_{h}, U_{h}\right)_{h \in H}$ provides a submanifold atlas for $H$ in $G$. This defines a manifold structure on $H$ for which $\left.\exp _{G}\right|_{U_{H}}$ is a local chart (see Lemma 8.6.5).

The map $m_{H}: H \times H \rightarrow H$ is a restriction of the multiplication map $m_{G}$ of $G$, hence smooth as a map $H \times H \rightarrow G$, and since $H$ is an initial submanifold of $G$, Lemma 8.6.5 implies that $m_{H}$ is smooth. With a similar argument we see that the inversion $\iota_{H}$ of $H$ is smooth. Therefore, $H$ is a Lie group, and the inclusion map $\iota_{H}: H \rightarrow G$ a smooth homomorphism. The corresponding morphism of Lie algebras $\mathbf{L}\left(\iota_{H}\right): \mathbf{L}(H) \rightarrow \mathbf{L}(G)$ is injective, and from $\exp _{G} \circ \mathbf{L}\left(\iota_{H}\right)=\iota_{H} \circ \exp _{H}$ it follows that its image consists of the set $\mathbf{L}^{e}(H)$ of all elements $x \in \mathbf{L}(G)$ with $\exp _{G}(\mathbb{R} x) \subseteq H$ because each element of $L^{e}(H)$ defines a smooth one-parameter group of $H$ (cf. Lemma 9.2.4).\\
(iii) Let $E$ be as in the proof of (i) and consider the smooth map

$$
\Psi: E \times H \rightarrow G, \quad(x, h) \mapsto \exp _{G}(x) h
$$

where $H$ carries the submanifold structure from (ii). Since $\exp _{H}: \mathbf{L}^{e}(H) \rightarrow H$ is a local diffeomorphism at 0 , the proof of (i) implies the existence of a 0 neighborhood $U_{E} \subseteq E$ and a 1-neighborhood $V_{H} \subseteq H$ such that

$$
\Psi_{1}:=\left.\Psi\right|_{U_{E} \times V_{H}}: U_{E} \times V_{H} \rightarrow \exp _{G}\left(U_{E}\right) V_{H}
$$

is a diffeomorphism onto an open subset of $G$. We further recall from Lemma 9.3.5, that we may assume, in addition, that

$$
\exp _{G}\left(U_{E}\right) \cap H=\{\boldsymbol{1}\}
$$

We now pick a small symmetric 0-neighborhood $V_{E}=-V_{E} \subseteq U_{E}$ such that $\exp _{G}\left(V_{E}\right) \exp _{G}\left(V_{E}\right) \subseteq \exp _{G}\left(U_{E}\right) V_{H}$. Its existence follows from the continuity of the multiplication in $G$. We claim that the map

$$
\varphi:=\left.\Psi\right|_{V_{E} \times H}: V_{E} \times H \rightarrow \exp _{G}\left(V_{E}\right) H
$$

is a diffeomorphism onto an open subset of $G$. To this end, we first observe that

$$
\varphi \circ\left(\operatorname{id}_{V_{E}} \times \rho_{h}\right)=\rho_{h} \circ \varphi \quad \text { for each } \quad h \in H
$$

i.e., $\varphi\left(x, h^{\prime} h\right)=\varphi\left(x, h^{\prime}\right) h$, so that

$$
T_{(x, h)}(\varphi) \circ\left(\operatorname{id}_{E} \times T_{\mathbf{1}}\left(\rho_{h}\right)\right)=T_{\varphi(x, \mathbf{1})}\left(\rho_{h}\right) \circ T_{(x, \mathbf{1})}(\varphi)
$$

Since $T_{(x, \mathbf{1})}(\varphi)=T_{(x, \mathbf{1})}(\Psi)$ is invertible for each $x \in V_{E}, T_{(x, h)}(\varphi)$ is invertible for each $(x, h) \in V_{E} \times H$. This implies that $\varphi$ is a local diffeomorphism at each point $(x, h)$. To see that $\varphi$ is injective, we observe that

$$
\exp _{G}(x) h=\varphi(x, h)=\varphi\left(x^{\prime}, h^{\prime}\right)=\exp _{G}\left(x^{\prime}\right) h^{\prime}
$$

implies that\\
$\exp _{G}(x)^{-1} \exp _{G}\left(x^{\prime}\right)=h\left(h^{\prime}\right)^{-1} \in \exp _{G}\left(V_{E}\right)^{2} \cap H \subseteq\left(\exp _{G}\left(U_{E}\right) V_{H}\right) \cap H=V_{H}$,\\
where we have used (9.13). We thus obtain $\exp _{G}\left(x^{\prime}\right) \in \exp _{G}(x) V_{H}$, so that the injectivity of $\Psi_{1}$ yields $x=x^{\prime}$, which in turn leads to $h=h^{\prime}$. This proves that $\varphi$ is injective and a local diffeomorphism, hence a diffeomorphism.

Example 9.3.8. We take a closer look at closed subgroups of the Lie group $(V,+)$, where $V$ is a finite-dimensional vector space. From Example 9.2.3, we know that $\exp _{V}=\operatorname{id}_{V}$. Let $H \subseteq V$ be a closed subgroup. Then

$$
\mathbf{L}(H)=\{x \in V: \mathbb{R} x \subseteq H\} \subseteq H
$$

is the largest vector subspace contained in $H$. Let $E \subseteq V$ be a vector space complement for $\mathbf{L}(H)$. Then $V \cong \mathbf{L}(H) \times E$, and we derive from $\mathbf{L}(H) \subseteq H$ that

$$
H \cong \mathbf{L}(H) \times(E \cap H)
$$

Lemma 9.3.5 implies the existence of some 0-neighborhood $U_{E} \subseteq E$ with $U_{E} \cap H=\{0\}$, hence that $H \cap E$ is discrete because 0 is an isolated point of $H \cap E$. Now Exercise 9.3.4 implies the existence of linearly independent elements $f_{1}, \ldots, f_{k} \in E$ with

$$
E \cap H=\mathbb{Z} f_{1}+\cdots+\mathbb{Z} f_{k}
$$

We conclude that

$$
H \cong \mathbf{L}(H) \times \mathbb{Z}^{k} \cong \mathbb{R}^{d} \times \mathbb{Z}^{k} \quad \text { for } d=\operatorname{dim} \mathbf{L}(H)
$$

Note that $\mathbf{L}(H)$ coincides with the connected component $H_{0}$ of 0 in $H$.\\
In view of Corollary 9.2.17, we may think of Lie groups as a special class of topological groups. We may therefore ask which subgroups of a Lie group $G$ are Lie groups with respect to the subspace topology:

Proposition 9.3.9. A subgroup of a Lie group is a Lie group with respect to the induced topology if and only if it is closed.

Proof. If $H$ is closed, then the Closed Subgroup Theorem 9.3.7 implies that $H$ is a submanifold of $G$ which is a Lie group.

Suppose, conversely, that $H$ is a Lie group. Since the inclusion map $\iota: H \rightarrow G$ is assumed to be a topological embedding, it is in particular continuous, hence smooth by the Automatic Smoothness Theorem 9.2.16 and since $\iota$ is injective, the same holds for $\mathbf{L}(\iota): \mathbf{L}(H) \rightarrow \mathbf{L}(G)$ (Proposition 9.2.13).

Let $V \subseteq \mathbf{L}(G)$ be an open convex 0-neighborhood for which $\left.\exp _{G}\right|_{V}$ is a diffeomorphism onto an open subset of $G$. Since $\mathbf{L}(\iota)$ is continuous, there exists a 0-neighborhood $U \subseteq \mathbf{L}(H)$ such that $\left.\exp _{H}\right|_{U}$ is a diffeomorphism onto an open subset of $H$ and $\mathbf{L}(\iota) U \subseteq V$. Since $H$ is a topological subgroup of $G$, there exists an open subset $U_{\mathfrak{g}} \subseteq V$ with

$$
\exp _{G}\left(U_{\mathfrak{g}}\right) \cap H=\iota\left(\exp _{H}(U)\right)=\exp _{G}(\mathbf{L}(\iota) U)
$$

Since $\mathbf{L}(\iota) U$ is locally closed in $\mathbf{L}(G)$, it now follows that $H$ is locally closed in $G$, hence closed by Exercise 9.3.3.

Definition 9.3.10. Let $G$ be a Lie group. A Lie subgroup of $G$ is a closed subgroup $H$ together with its Lie group structure provided by Proposition 9.3.9.

\subsubsection*{Examples}
Example 9.3.11 (Closed Subgroups of $\mathbb{T}$ ). Let $H \subseteq \mathbb{T} \subseteq\left(\mathbb{C}^{\times}, \cdot\right)$ be a closed proper (= different from $\mathbb{T}$ ) subgroup. Since $\operatorname{dim} \mathbb{T}=1$, it follows that $\mathbf{L}(H)=\{0\}$, so that the Identity Neighborhood Theorem implies that $H$ is discrete, hence finite because $\mathbb{T}$ is compact.

If $q: \mathbb{R} \rightarrow \mathbb{T}$ is the covering projection, $q^{-1}(H)$ is a closed proper subgroup of $\mathbb{R}$, hence cyclic (this is a very simple case of Exercise 9.3.4), which implies that $H=q\left(q^{-1}(H)\right)$ is also cyclic. Therefore, $H$ is one of the groups

$$
C_{n}=\left\{z \in \mathbb{T}: z^{n}=1\right\}
$$

of $n$th roots of unity.\\
Example 9.3.12 (Subgroups of $\mathbb{T}^{2}$ ). (a) Let $H \subseteq \mathbb{T}^{2}$ be a closed proper subgroup. Then $\mathbf{L}(H) \neq \mathbf{L}\left(\mathbb{T}^{2}\right)$ implies $\operatorname{dim} H<\operatorname{dim} \mathbb{T}^{2}=2$. Further, $H$ is compact, so that the group $\pi_{0}(H)$ of connected components of $H$ is finite.

If $\operatorname{dim} H=0$, then $H$ is finite, and for $n:=|H|$ it is contained in a subgroup of the form $C_{n} \times C_{n}$, where $C_{n} \subseteq \mathbb{T}$ is the subgroup of $n$th roots of unity (cf. Example 9.3.11).

If $\operatorname{dim} H=1$, then $H_{0}$ is a compact connected 1-dimensional Lie group, hence isomorphic to $\mathbb{T}$ (Exercise 9.3.5). Therefore, $H_{0}=\exp _{\mathbb{T}^{2}}(\mathbb{R} x)$ for some $x \in \mathbf{L}(H)$ with $\exp _{\mathbb{T}^{2}}(x)=\left(e^{2 \pi i x_{1}}, e^{2 \pi i x_{2}}\right)=(1,1)$, which is equivalent to $x \in \mathbb{Z}^{2}$. We conclude that the Lie algebras of the closed subgroups are of the form $\mathbf{L}(H)=\mathbb{R} x$ for some $x \in \mathbb{Z}^{2}$.\\
(b) For each $\theta \in \mathbb{R} \backslash \mathbb{Q}$, the image of the 1-parameter group

$$
\gamma: \mathbb{R} \rightarrow \mathbb{T}^{2}, \quad t \mapsto\left(e^{i \theta t}, e^{i t}\right)
$$

is not closed because $\gamma$ is injective. Hence the closure of $\gamma(\mathbb{R})$ is a closed subgroup of dimension at least 2 , which shows that $\gamma(\mathbb{R})$ is dense in $\mathbb{T}^{2}$. The subgroup $\gamma(\mathbb{R})$ is called a dense wind. We leave it as an exercise to the reader to verify that the dense wind is an initial submanifold of $\mathbb{T}^{2}$.

As we shall see later, in many situations it is important to have some information on the center of (simply) connected Lie groups. Below we shall use Lemma 9.2.21 to determine the kernel of the adjoint representation for various Lie groups. For that we have to know their center.

Example 9.3.13. (a) Let $\mathbb{K} \in\{\mathbb{R}, \mathbb{C}\}$. First, we recall from Proposition 2.1.10 that $Z\left(\mathrm{GL}_{n}(\mathbb{K})\right)=\mathbb{K}^{\times} \mathbf{1}$ and from Exercise 2.2.14(v) that

$$
Z\left(\mathrm{SL}_{n}(\mathbb{K})\right)=\left\{z \mathbf{1}: z \in \mathbb{K}^{\times}, z^{n}=1\right\}
$$

In particular,

$$
Z\left(\mathrm{SL}_{n}(\mathbb{C})\right)=\left\{z \mathbf{1}: z^{n}=1\right\} \cong C_{n}
$$

and

$$
Z\left(\mathrm{SL}_{n}(\mathbb{R})\right)= \begin{cases}\mathbf{1} & \text { for } n \in 2 \mathbb{N}_{0}+1 \\ \{ \pm \mathbf{1}\} & \text { for } n \in 2 \mathbb{N}\end{cases}
$$

(b) For $g \in Z\left(\mathrm{SU}_{n}(\mathbb{C})\right)=$ ker Ad we likewise have $g x=x g$ for all $x \in \mathfrak{s u}_{n}(\mathbb{C})$. From

$$
\mathfrak{g} \mathfrak{l}_{n}(\mathbb{C})=\mathfrak{u}_{n}(\mathbb{C})+i \mathfrak{u}_{n}(\mathbb{C})=\mathfrak{s} \mathfrak{u}_{n}(\mathbb{C})+i \mathfrak{s} \mathfrak{u}_{n}(\mathbb{C})+\mathbb{C} \mathbf{1}
$$

we derive that $g \in Z\left(\mathrm{GL}_{n}(\mathbb{C})\right)=\mathbb{C}^{\times} \mathbf{1}$. From that we immediately get

$$
Z\left(\mathrm{SU}_{n}(\mathbb{C})\right)=\left\{z \mathbf{1}: z^{n}=1\right\} \cong C_{n}
$$

and similarly we obtain

$$
Z\left(\mathrm{U}_{n}(\mathbb{C})\right)=\{z \mathbf{1}:|z|=1\} \cong \mathbb{T}
$$

(c) (cf. also Exercise 2.2.16) Next we show that

$$
Z\left(\mathrm{O}_{n}(\mathbb{R})\right)=\{ \pm \mathbf{1}\} \quad \text { and } \quad Z\left(\mathrm{SO}_{n}(\mathbb{R})\right)= \begin{cases}\mathrm{SO}_{2}(\mathbb{R}) & \text { for } n=2 \\ \mathbf{1} & \text { for } n \in 2 \mathbb{N}+1 \\ \{ \pm \mathbf{1}\} & \text { for } n \in 2 \mathbb{N}+2\end{cases}
$$

If $g \in Z\left(\mathrm{O}_{n}(\mathbb{R})\right)$, then $g$ commutes with each orthogonal reflection

$$
\sigma_{v}: \mathbb{R}^{n} \rightarrow \mathbb{R}^{n}, \quad w \mapsto w-2\langle v, w\rangle v
$$

in the hyperplane $v^{\perp}$, where $v$ is a unit vector. Since $\mathbb{R} v$ is the -1-eigenspace of $\sigma_{v}$, this space is invariant under $g$ (Exercise 2.1.1). This implies that for each $v \in \mathbb{R}^{n}$ we have $g . v \in \mathbb{R} v$ which by an elementary argument leads to $g \in \mathbb{R}^{\times} \mathbf{1}$ (Exercise 9.3.11). We conclude that

$$
Z\left(\mathrm{O}_{n}(\mathbb{R})\right)=\mathrm{O}_{n}(\mathbb{R}) \cap \mathbb{R}^{\times} \mathbf{1}=\{ \pm \mathbf{1}\}
$$

To determine the center of $\mathrm{SO}_{n}(\mathbb{R})$, we consider for orthogonal unit vectors $v_{1}, v_{2}$ the map $\sigma_{v_{1}, v_{2}}:=\sigma_{v_{1}} \sigma_{v_{2}} \in \mathrm{SO}_{n}(\mathbb{R})$ (a reflection in the subspace $\left.v_{1}^{\top} \cap v_{2}^{\top}\right)$. Since an element $g \in Z\left(\mathrm{SO}_{n}(\mathbb{R})\right)$ commutes with $\sigma_{v_{1}, v_{2}}$, it leaves the plane $\mathbb{R} v_{1}+\mathbb{R} v_{2}=\operatorname{ker}\left(\sigma_{v_{1}, v_{2}}+\mathbf{1}\right)$ invariant. If a linear map preserves all two-dimensional planes and $n \geq 3$, then it preserves all one-dimensional subspaces. As above, we get $g \in \mathbb{R}^{\times} \mathbf{1}$, which in turn leads to

$$
Z\left(\mathrm{SO}_{n}(\mathbb{R})\right)=\mathrm{SO}_{n}(\mathbb{R}) \cap \mathbb{R}^{\times} \mathbf{1}
$$

and the assertion follows.

\subsubsection*{Exercises for Section 9.3}
Exercise 9.3.1. If $\left(H_{j}\right)_{j \in J}$ is a family of subgroups of the Lie group $G$, then $\mathbf{L}\left(\bigcap_{j \in J} H_{j}\right)=\bigcap_{j \in J} \mathbf{L}\left(H_{j}\right)$.

Exercise 9.3.2. Let $\varphi: G \rightarrow H$ be a morphism of Lie groups. Show that

$$
\mathbf{L}(\operatorname{ker} \varphi)=\operatorname{ker} \mathbf{L}(\varphi)
$$

Exercise 9.3.3. (a) Show that each submanifold $S$ of a manifold $M$ is locally closed, i.e., for each point $s \in S$ there exists an open neighborhood $U$ of $s$ in $M$ such that $U \cap S$ is closed.\\
(b) Show that any locally closed subgroup $H$ of a Lie group $G$ is closed.

Exercise 9.3.4. Let $D \subseteq \mathbb{R}^{n}$ be a discrete subgroup. Then there exist linearly independent elements $v_{1}, \ldots, v_{k} \in \mathbb{R}^{n}$ with $D=\sum_{i=1}^{k} \mathbb{Z} v_{i}$.

Exercise 9.3.5 (Connected abelian Lie groups). Let $A$ be a connected abelian Lie group. Show that\\
(1) $\exp _{A}:(\mathbf{L}(A),+) \rightarrow A$ is a morphism of Lie groups.\\
(2) $\exp _{A}$ is surjective.\\
(3) $\Gamma_{A}:=\operatorname{ker} \exp _{A}$ is a discrete subgroup of $(\mathbf{L}(A),+)$.\\
(4) $\mathbf{L}(A) / \Gamma_{A} \cong \mathbb{R}^{k} \times \mathbb{T}^{m}$ for some $k, m \geq 0$. In particular, it is a Lie group and the quotient map $q_{A}: \mathbf{L}(A) \rightarrow \mathbf{L}(A) / \Gamma_{A}$ is a smooth map.\\
(5) $\exp _{A}$ factors through a diffeomorphism $\varphi: \mathbf{L}(A) / \Gamma_{A} \rightarrow A$.\\
(6) $A \cong \mathbb{R}^{k} \times \mathbb{T}^{m}$ as Lie groups.

Exercise 9.3.6 (Divisible groups). An abelian group $D$ is called divisible if for each $d \in D$ and $n \in \mathbb{N}$ there exists an $a \in D$ with $a^{n}=d$. Show that:\\
(1)* If $G$ is an abelian group, $H$ a subgroup and $f: H \rightarrow D$ a homomorphism into an abelian divisible group $D$, then there exists an extension of $f$ to a homomorphism $\widetilde{f}: G \rightarrow D$.\\
(2) If $G$ is an abelian group and $D$ a divisible subgroup, then $G \cong D \times H$ for some subgroup $H$ of $G$.

Exercise 9.3.7 (Nonconnected abelian Lie groups). Let $A$ be an abelian Lie group. Show that:\\
(1) The identity component of $A_{0}$ is isomorphic to $\mathbb{R}^{k} \times \mathbb{T}^{m}$ for some $k, m \in \mathbb{N}_{0}$.\\
(2) $A_{0}$ is divisible.\\
(3) $A \cong A_{0} \times \pi_{0}(A)$, where $\pi_{0}(A):=A / A_{0}$.\\
(4) There exists a discrete abelian group $D$ with $A \cong \mathbb{R}^{k} \times \mathbb{T}^{m} \times D$.

Exercise 9.3.8. If $q: G \rightarrow H$ is a surjective open morphism of topological groups, then the induced map $G / \operatorname{ker} q \rightarrow H$ is an isomorphism of topological groups, where $G / \operatorname{ker} q$ is endowed with the quotient topology.

Exercise 9.3.9. If $G$ is a topological group and $\mathbf{1} \in U \subseteq G$ a connected subset. Then all the sets $U^{n}:=U \cdots U$ are connected, and so is their union $\bigcup_{n} U^{n}$.

Exercise 9.3.10. Let $G$ be a topological group. Then for each open subset $O \subseteq G$ and for each subset $S \subseteq G$, the product sets

$$
O S=\{g h: g \in O, h \in S\} \quad \text { and } \quad S O=\{h g: g \in O, h \in S\}
$$

are open.\\
Exercise 9.3.11. Let $V$ be a $\mathbb{K}$-vector space and $A \in \operatorname{End}(V)$ with $A v \in \mathbb{K} v$ for all $v \in V$. Show that $A \in \mathbb{K} \operatorname{id}_{V}$.

\subsection*{Constructing Lie Group Structures on Groups}
In this subsection, we describe some methods to construct Lie group structures on groups, starting from a manifold structure on some "identity neighborhood" for which the group operations are smooth close to $\mathbf{1}$.

\subsubsection*{Group Topologies from Local Data}
The following lemma describes how to construct a group topology on a group $G$, i.e., a Hausdorff topology for which the group multiplication and the inversion are continuous, from a filter basis of subsets which then becomes a filter basis of identity neighborhoods for the group topology.

Definition 9.4.1. Let $X$ be a set. A set $\mathcal{F} \subseteq \mathbb{P}(X)$ of subsets of $X$ is called a filter basis if the following conditions are satisfied:\\
$(\mathrm{F} 1) \mathcal{F} \neq \emptyset$;\\
(F2) Each set $F \in \mathcal{F}$ is nonempty;\\
(F3) $A, B \in \mathcal{F} \Rightarrow(\exists C \in \mathcal{F}) C \subseteq A \cap B$.\\
Lemma 9.4.2. Let $G$ be a group and $\mathcal{F}$ a filter basis of subsets of $G$ satisfying $\bigcap \mathcal{F}=\{\mathbf{1}\}$ and\\
(U1) $(\forall U \in \mathcal{F})(\exists V \in \mathcal{F}) V V \subseteq U$.\\
(U2) $(\forall U \in \mathcal{F})(\exists V \in \mathcal{F}) V^{-1} \subseteq U$.\\
(U3) $(\forall U \in \mathcal{F})(\forall g \in G)(\exists V \in \mathcal{F}) g V g^{-1} \subseteq U$.\\
Then there exists a unique group topology on $G$ such that $\mathcal{F}$ is a basis of $\mathbf{1}$-neighborhoods in $G$. This topology is given by

$$
\{U \subseteq G:(\forall g \in U)(\exists V \in \mathcal{F}) g V \subseteq U\}
$$

Proof. Let

$$
\tau:=\{U \subseteq G:(\forall g \in U)(\exists V \in \mathcal{F}) g V \subseteq U\} .
$$

First, we show that $\tau$ is a topology. Clearly, $\emptyset, G \in \tau$. Let $\left(U_{j}\right)_{j \in J}$ be a family of elements of $\tau$ and $U:=\bigcup_{j \in J} U_{j}$. For each $g \in U$, there exist a $j_{0} \in J$ with $g \in U_{j_{0}}$ and a $V \in \mathcal{F}$ with $g V \subseteq U_{j_{0}} \subseteq U$. Thus $U \in \tau$ and we see that $\tau$ is stable under arbitrary unions.

If $U_{1}, U_{2} \in \tau$ and $g \in U_{1} \cap U_{2}$, there exist $V_{1}, V_{2} \in \mathcal{F}$ with $g V_{i} \subseteq U_{i}$. Since $\mathcal{F}$ is a filter basis, there exists $V_{3} \in \mathcal{F}$ with $V_{3} \subseteq V_{1} \cap V_{2}$, and then $g V_{3} \subseteq U_{1} \cap U_{2}$. We conclude that $U_{1} \cap U_{2} \in \tau$, and hence that $\tau$ is a topology on $G$.

We claim that the interior $U^{\circ}$ of a subset $U \subseteq G$ is given by

$$
U_{1}:=\{u \in U:(\exists V \in \mathcal{F}) u V \subseteq U\} .
$$

In fact, if there exists a $V \in \mathcal{F}$ with $u V \subseteq U$, then we pick a $W \in \mathcal{F}$ with $W W \subseteq V$ and obtain $u W W \subseteq U$, so that $u W \subseteq U_{1}$. Hence $U_{1} \in \tau$, i.e., $U_{1}$ is open, and it clearly is the largest open subset contained in $U$, i.e., $U_{1}=U^{\circ}$. It follows in particular that $U$ is a neighborhood of $g$ if and only if $g \in U^{\circ}$, and we see in particular that $\mathcal{F}$ is a neighborhood basis at $\mathbf{1}$. The property $\bigcap \mathcal{F}=\{\mathbf{1}\}$ implies that for $x \neq y$ there exists $U \in \mathcal{F}$ with $y^{-1} x \notin U$. For $V \in \mathcal{F}$ with $V V \subseteq U$ and $W \in \mathcal{F}$ with $W^{-1} \subseteq V$ we then obtain $y^{-1} x \notin V W^{-1}$, i.e., $x W \cap y V=\emptyset$. Thus ( $G, \tau$ ) is a Hausdorff space.

To see that $G$ is a topological group, we have to verify that the map

$$
f: G \times G \rightarrow G, \quad(x, y) \mapsto x y^{-1}
$$

is continuous. So let $x, y \in G, U \in \mathcal{F}$ and pick $V \in \mathcal{F}$ with $y V y^{-1} \subseteq U$ and $W \in \mathcal{F}$ with $W W^{-1} \subseteq V$. Then

$$
f(x W, y W)=x W W^{-1} y^{-1}=x y^{-1} y\left(W W^{-1}\right) y^{-1} \subseteq x y^{-1} y V y^{-1} \subseteq x y^{-1} U
$$

implies that $f$ is continuous at ( $x, y$ ). \(\square\)

Before we turn to Lie group structures, it is illuminating to first consider the topological variant of Lemma 9.4.3 below.

Lemma 9.4.3. Let $G$ be a group and $U=U^{-1}$ a symmetric subset containing 1. We further assume that $U$ carries a Hausdorff topology for which\\
(T1) $D:=\{(x, y) \in U \times U: x y \in U\}$ is an open subset of $U \times U$ and the group multiplication $m_{U}: D \rightarrow U,(x, y) \mapsto x y$ is continuous,\\
(T2) the inversion map $\iota_{U}: U \rightarrow U, u \mapsto u^{-1}$ is continuous, and\\
(T3) for each $g \in G$, there exists an open $\mathbf{1}$-neighborhood $U_{g}$ in $U$ with $c_{g}\left(U_{g}\right) \subseteq U$ such that the conjugation map $c_{g}: U_{g} \rightarrow U, x \mapsto g x g^{-1}$ is continuous.

Then there exists a unique group topology on $G$ for which the inclusion map $U \hookrightarrow G$ is a homeomorphism onto an open subset of $G$.

If, in addition, $U$ generates $G$, then (T1) and (T2) imply (T3).\\
Proof. First, we consider the filter basis $\mathcal{F}$ of $\mathbf{1}$-neighborhoods in $U$. Then (T1) implies (U1), (T2) implies (U2), and (T3) implies (U3). Moreover, the assumption that $U$ is Hausdorff implies that $\bigcap \mathcal{F}=\{\mathbf{1}\}$. Therefore, Lemma 9.4.2 implies that $G$ carries a unique structure of a (Hausdorff) topological group for which $\mathcal{F}$ is a basis of $\mathbf{1}$-neighborhoods.

We claim that the inclusion map $U \rightarrow G$ is an open embedding. So let $x \in U$. Then

$$
U_{x}:=U \cap x^{-1} U=\{y \in U:(x, y) \in D\}
$$

is open in $U$ and $\lambda_{x}$ restricts to a continuous map $U_{x} \rightarrow U$ with image $U_{x^{-1}}$. Its inverse is also continuous. Hence $\lambda_{x}^{U}: U_{x} \rightarrow U_{x^{-1}}$ is a homeomorphism. We conclude that the sets of the form $x V$, where $V$ a neighborhood of $\mathbf{1}$, form a basis of neighborhoods of $x \in U$. Hence the inclusion map $U \hookrightarrow G$ is an open embedding.

Suppose, in addition, that $G$ is generated by $U$. For each $g \in U$, there exists an open 1-neighborhood $U_{g}$ with $g U_{g} \times\left\{g^{-1}\right\} \subseteq D$. Then $c_{g}\left(U_{g}\right) \subseteq U$, and the continuity of $m_{U}$ implies that $\left.c_{g}\right|_{U_{g}}: U_{g} \rightarrow U$ is continuous.

Hence, for each $g \in U$, the conjugation $c_{g}$ is continuous in a neighborhood of $\mathbf{1}$. Since the set of all these $g$ is a submonoid of $G$ containing $U$, it contains $U^{n}$ for each $n \in \mathbb{N}$, hence all of $G$ because $G$ is generated by $U=U^{-1}$. Therefore, (T3) follows from (T1) and (T2).

\subsubsection*{Lie Group Structures from Local Data}
The following theorem, the smooth version of the preceding lemma, is an important tool to construct Lie group structures on groups. It is an important supplement to the Closed Subgroup Theorem 9.3.7.

Theorem 9.4.4. Let $G$ be a group and $U=U^{-1}$ a symmetric subset containing 1. We further assume that $U$ is a smooth manifold and that\\
(L1) $D:=\{(x, y) \in U \times U: x y \in U\}$ is an open subset and the multiplication $m_{U}: D \rightarrow U,(x, y) \mapsto x y$ is smooth,\\
(L2) the inversion map $\iota_{U}: U \rightarrow U, u \mapsto u^{-1}$ is smooth, and\\
(L3) for each $g \in G$ there exists an open $\mathbf{1}$-neighborhood $U_{g} \subseteq U$ with $c_{g}\left(U_{g}\right) \subseteq U$ and such that the conjugation map $c_{g}: U_{g} \rightarrow U, x \mapsto g x g^{-1}$ is smooth.

Then there exists a unique structure of a Lie group on $G$ such that the inclusion map $U \hookrightarrow G$ is a diffeomorphism onto an open subset of $G$.

If, in addition, $U$ generates $G$, then (L1) and (L2) imply (L3).

Proof. From the preceding Lemma 9.4.3, we immediately obtain the unique group topology on $G$ for which the inclusion map $U \hookrightarrow G$ is an open embedding.

Now we turn to the manifold structure. Let $V=V^{-1} \subseteq U$ be an open 1-neighborhood with $V V \times V V \subseteq D$, for which there exists a chart $(\varphi, V)$ of $U$. For $g \in G$, we consider the maps

$$
\varphi_{g}: g V \rightarrow E, \quad \varphi_{g}(x)=\varphi\left(g^{-1} x\right)
$$

which are homeomorphisms of $g V$ onto $\varphi(V) \subseteq E$. We claim that $\left(\varphi_{g}, g V\right)_{g \in G}$ is a smooth atlas of $G$.

Let $g_{1}, g_{2} \in G$ and put $W:=g_{1} V \cap g_{2} V$. If $W \neq \emptyset$, then $g_{2}^{-1} g_{1} \in V V^{-1}= V V$. The smoothness of the map

$$
\psi:=\left.\varphi_{g_{2}} \circ \varphi_{g_{1}}^{-1}\right|_{\varphi_{g_{1}}(W)}: \varphi_{g_{1}}(W) \rightarrow \varphi_{g_{2}}(W)
$$

given by

$$
\psi(x)=\varphi_{g_{2}}\left(\varphi_{g_{1}}^{-1}(x)\right)=\varphi_{g_{2}}\left(g_{1} \varphi^{-1}(x)\right)=\varphi\left(g_{2}^{-1} g_{1} \varphi^{-1}(x)\right)
$$

follows from the smoothness of the multiplication $V V \times V V \rightarrow U$. This proves that the charts $\left(\varphi_{g}, g U\right)_{g \in G}$ form a smooth atlas of $G$. Moreover, the construction implies that all left translations of $G$ are smooth maps.

The construction also shows that for each $g \in G$ the conjugation map $c_{g}: G \rightarrow G$ is smooth in a neighborhood of $\mathbf{1}$. Since all left translations are smooth, and

$$
c_{g} \circ \lambda_{x}=\lambda_{c_{g}(x)} \circ c_{g}
$$

the smoothness of $c_{g}$ in a neighborhood of $x \in G$ follows. Therefore, all conjugations and hence also all right multiplications are smooth. The smoothness of the inversion follows from its smoothness on $V$ and the fact that left and right multiplications are smooth. Finally, the smoothness of the multiplication follows from the smoothness in $\mathbf{1} \times \mathbf{1}$ because

$$
g_{1} x g_{2} y=g_{1} g_{2} c_{g_{2}^{-1}}(x) y
$$

Next we show that the inclusion $U \hookrightarrow G$ of $U$ is a diffeomorphism. So let $x \in U$ and recall the open set $U_{x}=U \cap x^{-1} U$. Then $\lambda_{x}$ restricts to a smooth map $U_{x} \rightarrow U$ with image $U_{x^{-1}}$. Its inverse is also smooth. Hence $\lambda_{x}^{U}: U_{x} \rightarrow U_{x^{-1}}$ is a diffeomorphism. Since $\lambda_{x}: G \rightarrow G$ also is a diffeomorphism, it follows that the inclusion $\lambda_{x} \circ \lambda_{x^{-1}}^{U}: U_{x^{-1}} \rightarrow G$ is a diffeomorphism. As $x$ was arbitrary, the inclusion of $U$ in $G$ is a diffeomorphic embedding.

The uniqueness of the Lie group structure is clear because each locally diffeomorphic bijective homomorphism between Lie groups is a diffeomorphism (Proposition 9.2.13(3)).

Finally, we assume that $G$ is generated by $U$. We show that in this case (L3) is a consequence of (L1) and (L2); the argument is similar to the topological case. Indeed, for each $g \in U$, there exists an open 1-neighborhood $U_{g}$\\
with $g U_{g} \times\left\{g^{-1}\right\} \subseteq D$. Then $c_{g}\left(U_{g}\right) \subseteq U$, and the smoothness of $m_{U}$ implies that $\left.c_{g}\right|_{U_{g}}: U_{g} \rightarrow U$ is smooth. Hence, for each $g \in U$, the conjugation $c_{g}$ is smooth in a neighborhood of $\mathbf{1}$. Since the set of all these $g$ is a submonoid of $G$ containing $U$, it contains $U^{n}$ for each $n \in \mathbb{N}$, hence all of $G$ because $G$ is generated by $U=U^{-1}$. Therefore, (L3) is satisfied.

Corollary 9.4.5. Let $G$ be a group and $N \unlhd G$ a normal subgroup of $G$ that carries a Lie group structure. Then there exists a Lie group structure on $G$ for which $N$ is an open subgroup if and only if for each $g \in G$ the restriction $\left.c_{g}\right|_{N}$ is a smooth automorphism of $N$.

Proof. If $N$ is an open normal subgroup of the Lie group $G$, then clearly all inner automorphisms of $G$ restrict to smooth automorphisms of $N$.

Suppose, conversely, that $N$ is a normal subgroup of the group $G$ which is a Lie group and that all inner automorphisms of $G$ restrict to smooth automorphisms of $N$. Then we can apply Theorem 9.4.4 with $U=N$ and obtain a Lie group structure on $G$ for which the inclusion $N \rightarrow G$ is a diffeomorphism onto an open subgroup of $G$.

Corollary 9.4.6. Let $\varphi: G \rightarrow H$ be a continuous homomorphism of topological groups which is a covering map. If $G$ or $H$ is a Lie group, then the other group carries a unique Lie group structure for which $\varphi$ is a morphism of Lie groups which is a local diffeomorphism.

Proof. Since $\operatorname{ker} \varphi$ is discrete, there exists an open symmetric identity neighborhood $U_{G} \subseteq G$ for which $U_{G}^{3}:=U_{G} U_{G} U_{G}$ intersects $\operatorname{ker}(\varphi)$ in $\{\mathbf{1}\}$. For $x, y \in U_{G}$ with $\varphi(x)=\varphi(y)$, we then have $x^{-1} y \in U_{G}^{2} \cap \operatorname{ker}(\varphi)=\{\mathbf{1}\}$, so that $\left.\varphi\right|_{U_{G}}$ is injective. Since $\varphi$ is an open map, this implies that $\left.\varphi\right|_{U_{G}}$ is a homeomorphism onto an open subset $U_{H}:=\varphi\left(U_{G}\right)$ of $H$.

Suppose first that $G$ is a Lie group. Then we apply Theorem 9.4.4 to $U_{H}$, endowed with the manifold structure for which $\left.\varphi\right|_{U_{G}}$ is a diffeomorphism. Then (L2) follows from $\varphi(x)^{-1}=\varphi\left(x^{-1}\right)$. To verify the smoothness of the multiplication map

$$
m_{U_{H}}: D_{H}:=\left\{(a, b) \in U_{H} \times U_{H}: a b \in U_{H}\right\} \rightarrow U_{H}
$$

we first observe that, if $x, y \in U_{G}$ satisfy $(\varphi(x), \varphi(y)) \in D_{H}$, i.e., $\varphi(x y) \in U_{H}$, then there exists a $z \in U_{G}$ with $\varphi(x y)=\varphi(z)$, and $x y z^{-1} \in U_{G}{ }^{3} \cap \operatorname{ker}(\varphi)= \{1\}$ yields $z=x y$. We thus have $D_{H}=(\varphi \times \varphi)\left(D_{G}\right)$ for

$$
D_{G}:=\left\{(x, y) \in U_{G} \times U_{G}: x y \in U_{G}\right\}
$$

and the smoothness of $m_{H}$ follows from the smoothness of the multiplication $m_{U_{G}}: D_{G} \rightarrow U_{G}$ and

$$
m_{U_{H}} \circ(\varphi \times \varphi)=\varphi \circ m_{U_{G}}
$$

To verify (L3), we note that the surjectivity of $\varphi$ implies that for each $h \in H$ there is an element $g \in G$ with $\varphi(g)=h$. Now we choose an open 1-neighborhood $U_{g} \subseteq U_{G}$ with $c_{g}\left(U_{g}\right) \subseteq U_{G}$ and put $U_{h}:=\varphi\left(U_{g}\right)$.

If, conversely, $H$ is a Lie group, then we apply Theorem 9.4.4 to $U_{G}$, endowed with the manifold structure for which $\left.\varphi\right|_{U_{G}}$ is a diffeomorphism onto $U_{H}$. Again, (L2) follows right away, and (L1) follows from $(\varphi \times \varphi)\left(D_{G}\right) \subseteq D_{H}$ and the smoothness of

$$
m_{U_{H}} \circ(\varphi \times \varphi)=\varphi \circ m_{U_{G}} .
$$

For (L3), we choose $U_{g}$ as any open 1-neighborhood in $U_{G}$ with $c_{g}\left(U_{g}\right) \subseteq U$. Then the smoothness of $\left.c_{g}\right|_{U_{g}}$ follows from the smoothness of $\varphi \circ c_{g}= c_{\varphi(g)} \circ \varphi$.

Corollary 9.4.7. If $G$ is a connected Lie group and $q_{G}: \widetilde{G} \rightarrow G$ its universal covering space, then $\widetilde{G}$ carries a unique Lie group structure for which $q_{G}$ is a smooth covering map. We call this Lie group the simply connected covering group of $G$.

Proof. We first have to construct a (topological) group structure on the universal covering space $\widetilde{G}$. Pick an element $\widetilde{\mathbf{1}} \in q_{G}^{-1}(\mathbf{1})$. Then the multiplication map $m_{G}: G \times G \rightarrow G$ lifts uniquely to a continuous map $\widetilde{m}_{G}: \widetilde{G} \times \widetilde{G} \rightarrow \widetilde{G}$ with $\widetilde{m}_{G}(\widetilde{\mathbf{1}}, \widetilde{\mathbf{1}})=\widetilde{\mathbf{1}}$. To see that the multiplication map $\widetilde{m}_{G}$ is associative, we observe that

$$
\begin{aligned}
& q_{G} \circ \widetilde{m}_{G} \circ\left(\operatorname{id}_{\widetilde{G}} \times \widetilde{m}_{G}\right)=m_{G} \circ\left(q_{G} \times q_{G}\right) \circ\left(\operatorname{id}_{\widetilde{G}} \times \widetilde{m}_{G}\right) \\
& =m_{G} \circ\left(\operatorname{id}_{G} \times m_{G}\right) \circ\left(q_{G} \times q_{G} \times q_{G}\right)=m_{G} \circ\left(m_{G} \times \operatorname{id}_{G}\right) \circ\left(q_{G} \times q_{G} \times q_{G}\right) \\
& =q_{G} \circ \widetilde{m}_{G} \circ\left(\widetilde{m}_{G} \times \operatorname{id}_{\widetilde{G}}\right)
\end{aligned}
$$

so that the two continuous maps

$$
\widetilde{m}_{G} \circ\left(\mathrm{id}_{\widetilde{G}} \times \widetilde{m}_{G}\right), \quad \widetilde{m}_{G} \circ\left(\widetilde{m}_{G} \times \mathrm{id}_{\widetilde{G}}\right): \widetilde{G}^{3} \rightarrow G
$$

are lifts of the same map $G^{3} \rightarrow G$ and both map $(\widetilde{\mathbf{1}}, \widetilde{\mathbf{1}}, \widetilde{\mathbf{1}})$ to $\widetilde{\mathbf{1}}$. Hence the uniqueness of lifts implies that $\widetilde{m}_{G}$ is associative. We likewise obtain that the unique lift $\widetilde{\iota}_{G}: \widetilde{G} \rightarrow \widetilde{G}$ of the inversion map $\iota_{G}: G \rightarrow G$ with $\widetilde{\iota}_{G}(\widetilde{\mathbf{1}})=\widetilde{\mathbf{1}}$ satisfies

$$
\widetilde{m}_{G} \circ\left(\iota_{G}, \mathrm{id}_{\widetilde{G}}\right)=\widetilde{\mathbf{1}}=\widetilde{m}_{G} \circ\left(\mathrm{id}_{\widetilde{G}}, \iota_{G}\right) .
$$

Finally, $\lambda_{\widetilde{\mathbf{1}}}$ lifts $\lambda_{\mathbf{1}}=\operatorname{id}_{G}$, so that $\lambda_{\widetilde{\mathbf{1}}}(\widetilde{\mathbf{1}})=\widetilde{\mathbf{1}}$ leads to $\lambda_{\widetilde{\mathbf{1}}}=\operatorname{id}_{\widetilde{G}}$, and likewise one shows that $\rho_{\widetilde{\mathbf{1}}}=\operatorname{id}_{\widetilde{G}}$, so that $\widetilde{\mathbf{1}}$ is a neutral element for the multiplication on $\widetilde{G}$. Therefore, $\widetilde{m}_{G}$ defines on $\widetilde{G}$ a topological group structure such that $q_{G}: \widetilde{G} \rightarrow G$ is a covering morphism of topological groups. Now Corollary 9.4.6 applies.

\subsubsection*{Existence of Lie Groups for a Given Lie Algebra}
Theorem 9.4.8 (Integral Subgroup Theorem). Let $G$ be a Lie group with Lie algebra $\mathfrak{g}$ and $\mathfrak{h} \subseteq \mathfrak{g}$ a Lie subalgebra. Then the subgroup $H:=\langle\exp \mathfrak{h}\rangle$ of $G$ generated by $\exp _{G} \mathfrak{h}$ carries a Lie group structure with the following properties:\\
(a) There exists an open 0 -neighborhood $W \subseteq \mathfrak{h}$ on which the Hausdorff series converges, and

$$
\exp _{\mathfrak{h}}: \mathfrak{h} \rightarrow H, \quad x \mapsto \exp _{G} x
$$

maps $W$ diffeomorphism onto its open image in $H$ and satisfies

$$
\exp _{\mathfrak{h}}(x * y)=\exp _{\mathfrak{h}}(x) \exp _{\mathfrak{h}}(y) \quad \text { for } \quad x, y \in W
$$

(b) The inclusion $i_{H}: H \rightarrow G$ is a smooth morphism of Lie groups and $\mathbf{L}\left(i_{H}\right): \mathbf{L}(H) \rightarrow \mathfrak{h}$ an isomorphism of Lie algebras. These two properties determine the Lie group structure on $H$ uniquely.\\
(c) If $H \subseteq H_{1}$ for some subgroup $H_{1}$ for which $H_{1} / H$ is countable, then $\mathfrak{h}=\left\{x \in \mathbf{L}(G): \exp _{G}(\mathbb{R} x) \subseteq H_{1}\right\}$. In particular,

$$
\mathfrak{h}=\mathbf{L}^{e}(H):=\left\{x \in \mathbf{L}(G): \exp _{G}(\mathbb{R} x) \subseteq H\right\}
$$

(d) $H$ is connected.\\
(e) $H$ is closed in $G$ if and only if $i_{H}$ is a topological embedding.

Proof. (a) Let $V \subseteq \mathfrak{g}$ be an open convex symmetric 0 -neighborhood for which the Hausdorff series for $x * y$ converges for $x, y \in V$ and satisfies

$$
\exp _{G}(x * y)=\exp _{G}(x) \exp _{G}(y)
$$

(Proposition 9.2.32). We further assume that $\left.\exp _{G}\right|_{V}$ is a diffeomorphism onto an open subset of $G$.

Put $W:=V \cap \mathfrak{h}$. Then $x * y \in \mathfrak{h}$ for $x, y \in W$ because each summand in the Hausdorff series is an iterated Lie bracket. Further, $x * y$ defines a smooth function $W \times W \rightarrow \mathfrak{h}$ because it is the restriction of a smooth function $V \times V \rightarrow \mathfrak{g}$. We consider the subset $U:=\exp _{G}(W) \subseteq H$. From $W=-W$, we derive $U=U^{-1}$ and note that $\varphi:=\left.\exp _{G}\right|_{W}$ is injective. We may thus endow $U \subseteq H$ with the manifold structure turning $\varphi$ into a diffeomorphism.

Then

$$
\widetilde{D}=\{(x, y) \in W \times W: x * y \in W\}
$$

is an open subset of $W \times W$ on which the *-multiplication is smooth, so that the multiplication $D \rightarrow U$ is also smooth. We further observe that

$$
\exp _{G}(-x)=\exp _{G}(x)^{-1}
$$

from which it follows that the inversion on $U$ is smooth. Since $U$ generates $H$, (L3) follows from (L1) and (L2). Therefore, $U$ satisfies all assumptions of

Theorem 9.4.4, so that we obtain a Lie group structure on $H$ for which $\varphi$, resp., $\exp _{\mathfrak{h}}$, induces a local diffeomorphism at 0 .\\
(b) Since the map $i_{H} \circ \exp _{\mathfrak{h}}: \mathfrak{h} \rightarrow G$ is smooth and $\exp _{\mathfrak{h}}$ is a local diffeomorphism at 0 , the inclusion $i_{H}: H \rightarrow G$ is smooth. Now

$$
\mathbf{L}\left(i_{H}\right): \mathbf{L}(H) \rightarrow \mathbf{L}(G)
$$

is injective, and by construction, its image contains $\mathfrak{h}$ because each element of $x$ generates a one-parameter group of $H$. As $\operatorname{dim} \mathbf{L}(H)=\operatorname{dim} H=\operatorname{dim} \mathfrak{h}$, we have $\mathbf{L}\left(\iota_{H}\right) \mathbf{L}(H)=\mathfrak{h}$.

If $\widehat{H}$ denotes another Lie group structure on the subgroup $H$ for which $i_{\widehat{H}}: \widehat{H} \rightarrow G$ is smooth and $\mathbf{L}\left(i_{\widehat{H}}\right): \mathbf{L}(\widehat{H}) \rightarrow \mathfrak{h} \subseteq \mathbf{L}(G)$ is an isomorphism of Lie algebras, then the relation

$$
i_{\widehat{H}} \circ \exp _{\widehat{H}}=\exp _{G} \circ \mathbf{L}\left(\iota_{\widehat{H}}\right)=\exp _{\mathfrak{h}} \circ \mathbf{L}\left(\iota_{\widehat{H}}\right)
$$

and (a) imply that the identical map

$$
j:=\iota_{H}^{-1} \circ \iota_{\widehat{H}}: \widehat{H} \rightarrow H
$$

is a bijective morphism of connected Lie groups for which $\mathbf{L}(j)$ is an isomorphism, hence an isomorphism (Proposition 9.2.13(3)).\\
(c) Clearly, $\exp _{G}(\mathfrak{h}) \subseteq H \subseteq H_{1}$, so that it remains to show that any $x \in \mathbf{L}(G)$ satisfying $\exp _{G}(\mathbb{R} x) \subseteq H_{1}$ is contained in $\mathfrak{h}$. Let $U_{H} \subseteq H$ be a 1neighborhood with respect to the Lie group topology for which $U_{H}^{-1} U_{H} \subseteq U$, where $U=\exp _{G}(W)$ is chosen as in (a). In view of Proposition 9.1.15(iv) and the countability of the set $H_{1} / H$ of $H$-cosets in $H_{1}$, there exists a sequence $\left(h_{n}\right)_{n \in \mathbb{N}}$ with $H_{1}=\bigcup_{n \in \mathbb{N}} h_{n} U_{H}$. Choose $\varepsilon>0$ with $\exp _{G}([-2 \varepsilon, 2 \varepsilon] x) \subseteq \exp _{G}(V)$. Since the interval $[0, \varepsilon]$ is uncountable, there exists some $n$ for which $\left\{t \in[0, \varepsilon]: \exp _{G}(t x) \in h_{n} U_{H}\right\}$ is uncountable. In particular, there are two such elements $t_{1}<t_{2}$ and $u_{1}, u_{2} \in U_{H}$ with

$$
\exp _{G}\left(t_{1} x\right)=h_{n} u_{1} \neq \exp _{G}\left(t_{2} x\right)=h_{n} u_{2}
$$

Then

$$
\begin{aligned}
\exp _{G}\left(\left(-t_{1}+t_{2}\right) x\right) & =\exp _{G}\left(t_{1} x\right)^{-1} \exp _{G}\left(t_{2} x\right) \\
& =u_{1}^{-1} h_{n}^{-1} h_{n} u_{2} \in U_{H}^{-1} U_{H} \subseteq U=\exp _{G}(W)
\end{aligned}
$$

In particular, there is a nonzero $y \in W \subseteq \mathfrak{h}$ with

$$
\exp _{G} y=\exp _{G}\left(\left(-t_{1}+t_{2}\right) x\right)
$$

Since $y$ and $\left(t_{2}-t_{1}\right) x$ are contained in $V$, it follows that $x=\frac{1}{t_{2}-t_{1}} y \in \mathfrak{h}$.\\
(d) Since $H$ is generated by $\exp _{H}(\mathbf{L}(H))$, this follows from Lemma 9.2.9.\\
(e) If $H$ is closed, then the Closed Subgroup Theorem 9.3.7 shows that $H$ is a Lie group with respect to the subspace topology, so that the uniqueness part of (b) implies that $i_{H}$ is a topological embedding.

If, conversely, $i_{H}$ is a topological embedding, then the subgroup $H$ of $G$ is locally compact, hence in particular locally closed, and therefore closed by Exercise 9.3.3.

Remark 9.4.9. Example 9.3.12, the dense wind in the 2 -torus, shows that we cannot expect that the group $H=\left\langle\exp _{G} \mathfrak{h}\right\rangle$ is closed in $G$ or that the inclusion map $H \rightarrow G$ (which is a smooth homomorphism) is a topological embedding.

Definition 9.4.10. Let $G$ be a Lie group. An integral subgroup $H$ of $G$ is a subgroup that is generated by $\exp \mathfrak{h}$ for a subalgebra $\mathfrak{h}$ of the Lie algebra $\mathfrak{g}$ of $G$.

The Integral Subgroup Theorem implies in particular that each Lie subalgebra $\mathfrak{h}$ of the Lie algebra $\mathbf{L}(G)$ of a Lie group $G$ is integrable in the sense that it is the Lie algebra of some Lie group $H$.

Combining this with Ado's Theorem on the existence of faithful linear representations of a Lie algebra, we obtain one of the cornerstones of the theory of Lie groups:

Theorem 9.4.11 (Lie's Third Theorem). Each finite-dimensional Lie algebra $\mathfrak{g}$ is the Lie algebra of a connected Lie group $G$.

Proof. Ado's Theorem implies that $\mathfrak{g}$ is isomorphic to a subalgebra of some $\mathfrak{g} \mathfrak{l}_{n}(\mathbb{R})$, so that the assertion follows directly from the Integral Subgroup Theorem.

\subsubsection*{Exercises for Section 9.4}
Exercise 9.4.1. Let $\varphi: G \rightarrow H$ be a surjective morphism of topological groups. Show that the following conditions are equivalent:\\
(1) $\varphi$ is open with discrete kernel.\\
(2) $\varphi$ is a covering in the topological sense, i.e., each $h \in H$ has an open neighborhood $U$ such that $\varphi^{-1}(U)=\bigcup_{i \in I} U_{i}$ is a disjoint union of open subsets $U_{i}$ for which all restrictions $\left.\varphi\right|_{U_{i}}: U_{i} \rightarrow U$ are homeomorphisms.

Exercise 9.4.2 (Refining Lemma 9.4.3). Show that the conclusion of Lemma 9.4.3 is still valid if the assumption (T1) is weakened as follows: There exists an open subset $D \subseteq U \times U$ with $x y \in U$ for all $(x, y) \in D$, containing all pairs $\left(x, x^{-1}\right),(x, \mathbf{1}),(\mathbf{1}, x)$ for $x \in U$, such that the group multiplication $m: D \rightarrow U$ is continuous.

Exercise 9.4.3. Let $G$ be an abelian group and $N \leq G$ a subgroup carrying a Lie group structure. Then there exists a unique Lie group structure on $G$ for which $N$ is an open subgroup.

Exercise 9.4.4. Let $G$ be a connected topological group and $\Gamma \unlhd G$ a discrete normal subgroup. Then $\Gamma$ is central.

Exercise 9.4.5. Let $X$ be a topological space and $\left(X_{i}\right)_{i \in I}$ connected subspaces of $X$ with $X=\bigcup_{i \in I} X_{i}$. If $\bigcap_{i \in I} X_{i} \neq \emptyset$, then $X$ is connected.

Exercise 9.4.6. We consider the simply connected covering group $G:= \widetilde{\mathrm{SL}}_{2}(\mathbb{R})$ with $\mathbf{L}(G)=\mathfrak{s l}_{2}(\mathbb{R})$ and we write $q: G \rightarrow \mathrm{SL}_{2}(\mathbb{R})$ for the covering homomorphism. The map

$$
\alpha: \mathbb{R} \rightarrow G, \quad t \mapsto \exp _{G}(t 2 \pi u), \quad u=:\left(\begin{array}{cc}
0 & 1 \\
-1 & 0
\end{array}\right)
$$

is injective.

\subsection*{Covering Theory for Lie Groups}
In this section, we eventually turn to applications of covering theory to Lie groups. Our goal is to see to which extent the Lie algebra and the fundamental group determine a connected Lie group. From the preceding section, we know that each connected Lie group $G$ has a simply connected covering group $\widetilde{G}$ which also carries a Lie group structure. As we shall see below, the kernel of the covering morphism $q_{G}: \widetilde{G} \rightarrow G$ can be identified with the fundamental group $\pi_{1}(G)$. Since $\mathbf{L}\left(q_{G}\right)$ is an isomorphism of Lie algebras, we then have $\mathbf{L}(G) \cong \mathbf{L}(\widetilde{G})$. We further prove the Monodromy Principle which implies that any Lie algebra morphism $\mathbf{L}(G) \rightarrow \mathbf{L}(H)$ can be integrated to a group homomorphism, provided $G$ is 1 -connected, i.e., connected and simply connected. From that we shall derive that the Lie algebra $\mathbf{L}(G)$ determines the corresponding simply connected group up to isomorphy.

\subsubsection*{Simply Connected Coverings of Lie Groups}
Proposition 9.5.1. A surjective morphism $\varphi: G \rightarrow H$ of Lie groups is a covering if and only if $\mathbf{L}(\varphi): \mathbf{L}(G) \rightarrow \mathbf{L}(H)$ is a linear isomorphism.

Proof. If $\varphi$ is a covering, then it is an open homomorphism with discrete kernel (Exercise 9.4.1), so that $\mathbf{L}(\operatorname{ker} \varphi)=\{0\}$, and Proposition 9.2.13 implies that $\mathbf{L}(\varphi)$ is bijective, hence an isomorphism of Lie algebras.

If, conversely, $\mathbf{L}(\varphi)$ is bijective, then Proposition 9.2.13 implies that

$$
\mathbf{L}(\operatorname{ker} \varphi)=\operatorname{ker} \mathbf{L}(\varphi)=\{0\}
$$

and the Closed Subgroup Theorem 9.3.7 shows that $\operatorname{ker} \varphi$ is discrete. Since $\mathbf{L}(\varphi)$ is surjective, Proposition 9.2.13 implies that $\varphi$ is an open map. Finally, Exercise 9.4.1 shows that $\varphi$ is a covering.

Proposition 9.5.2. For a covering $q: G_{1} \rightarrow G_{2}$ of connected Lie groups, the following equalities hold

$$
q\left(Z\left(G_{1}\right)\right)=Z\left(G_{2}\right) \quad \text { and } \quad Z\left(G_{1}\right)=q^{-1}\left(Z\left(G_{2}\right)\right) .
$$

Proof. Since $q$ is a covering, $\mathbf{L}(q): \mathbf{L}\left(G_{1}\right) \rightarrow \mathbf{L}\left(G_{2}\right)$ is an isomorphism of Lie algebras, and the adjoint representations satisfy

$$
\operatorname{Ad}_{G_{2}}\left(q\left(g_{1}\right)\right) \circ \mathbf{L}(q)=\mathbf{L}(q) \circ \operatorname{Ad}_{G_{1}}\left(g_{1}\right)
$$

Hence

$$
Z\left(G_{1}\right)=\operatorname{ker} \operatorname{Ad}_{G_{1}}=q^{-1} \operatorname{ker} \operatorname{Ad}_{G_{2}}=q^{-1}\left(Z\left(G_{2}\right)\right)
$$

Now the claim follows from the surjectivity of $q$.\\
Theorem 9.5.3 (Lifting Theorem for Groups). Let $q: G \rightarrow H$ be a covering morphism of Lie groups. If $f: L \rightarrow H$ is a morphism of Lie groups, where $L$ is 1 -connected, then there exists a unique lift $\widehat{f}: L \rightarrow G$ which is a morphism of Lie groups.

Proof. Since Lie groups are locally arcwise connected, the Lifting Theorem A.2.9 implies the existence of a unique lift $\widetilde{f}$ with $\widetilde{f}\left(\mathbf{1}_{L}\right)=\mathbf{1}_{G}$. Then

$$
m_{G} \circ(\tilde{f} \times \tilde{f}): L \times L \rightarrow G
$$

is the unique lift of $m_{H} \circ(f \times f): L \times L \rightarrow H$ mapping $\left(\mathbf{1}_{L}, \mathbf{1}_{L}\right)$ to $\mathbf{1}_{G}$. We also have

$$
q \circ \widetilde{f} \circ m_{L}=f \circ m_{L}=m_{H} \circ(f \times f)
$$

so that $\widetilde{f} \circ m_{L}$ is another lift mapping $\left(\mathbf{1}_{L}, \mathbf{1}_{L}\right)$ to $\mathbf{1}_{G}$. Therefore,

$$
\tilde{f} \circ m_{L}=m_{G} \circ(\tilde{f} \times \tilde{f})
$$

which means that $\widetilde{f}$ is a group homomorphism.\\
Since $q$ is a local diffeomorphism and $\widetilde{f}$ is a continuous lift of $f$, it is also smooth in an identity neighborhood of $L$, hence smooth by Corollary 9.2.11.

Theorem 9.5.4. Let $G$ be a connected Lie group and $q_{G}: \widetilde{G} \rightarrow G$ a universal covering homomorphism. Then $\operatorname{ker} q_{G} \cong \pi_{1}(G)$ is a discrete central subgroup and $G \cong \widetilde{G} / \operatorname{ker} q_{G}$.

Moreover, for any discrete central subgroup $\Gamma \subseteq \widetilde{G}$, the group $\widetilde{G} / \Gamma$ is a connected Lie group with the same universal covering group as $G$. We thus obtain a bijection from the set of all $\operatorname{Aut}(\widetilde{G})$-orbits in the set of discrete central subgroups of $\widetilde{G}$ onto the set of isomorphy classes of connected Lie groups whose universal covering group is isomorphic to $\widetilde{G}$.

Proof. First, we note that $\operatorname{ker} q_{G}$ is a discrete normal subgroup of the connected Lie group $\widetilde{G}$, hence central by Exercise 9.4.4. Left multiplications by elements of $\operatorname{ker} q_{G}$ lead to deck transformations of the covering $\widetilde{G} \rightarrow G$, and this group of deck transformations acts transitively on the fiber $\operatorname{ker} q_{G}$ of $\mathbf{1}$. Proposition A.2.15 now shows that

$$
\pi_{1}(G) \cong \operatorname{ker} q_{G}
$$

as groups. Since $q_{G}: \widetilde{G} \rightarrow G$ is open and surjective, we have $G \cong \widetilde{G} / \operatorname{ker} q_{G}$ as topological groups (Exercise 9.3.8), hence as Lie groups (Theorem 9.2.16).

If, conversely, $\Gamma \subseteq \widetilde{G}$ is a discrete central subgroup, then the topological quotient group $\widetilde{G} / \Gamma$ is a Lie group (Corollary 9.4.6) whose universal covering group is $\widetilde{G}$. Two such groups $\widetilde{G} / \Gamma_{1}$ and $\widetilde{G} / \Gamma_{2}$ are isomorphic if and only if there exists a Lie group automorphism $\varphi \in \operatorname{Aut}(\widetilde{G})$ with $\varphi\left(\Gamma_{1}\right)=\Gamma_{2}$ (Theorem 9.5.3). Therefore, the isomorphism classes of Lie groups with the same universal covering group $G$ are parameterized by the orbits of the group $\operatorname{Aut}(\widetilde{G})$ in the set of discrete central subgroups of $\widetilde{G}$. \(\square\)

Remark 9.5.5. (a) Since the normal subgroup $\operatorname{Inn}(\widetilde{G}):=\left\{c_{g}: g \in \widetilde{G}\right\}$ of inner automorphisms acts trivially on the center of $\widetilde{G}$, the action of $\operatorname{Aut}(\widetilde{G})$ on the set of all discrete normal subgroups factors through an action of the $\operatorname{group} \operatorname{Out}(\widetilde{G}):=\operatorname{Aut}(\widetilde{G}) / \operatorname{Inn}(\widetilde{G})$.\\
(b) Since each automorphism $\varphi \in \operatorname{Aut}(G)$ lifts to a unique automor$\operatorname{phism} \widetilde{\varphi} \in \operatorname{Aut}(\widetilde{G})$ (Theorem 9.5.3), we have a natural embedding $\operatorname{Aut}(G) \hookrightarrow \operatorname{Aut}(\widetilde{G})$, and the image of this homomorphism consists of the stabilizer of the subgroup ker $q_{G} \subseteq Z(\widetilde{G})$ in $\operatorname{Aut}(\widetilde{G})$.

Example 9.5.6 (Connected abelian Lie groups). Let $A$ be a connected abelian Lie group and $\exp _{A}: \mathbf{L}(A) \rightarrow A$ its exponential function. Then $\exp _{A}$ is a morphism of Lie groups with $\mathbf{L}\left(\exp _{A}\right)=\operatorname{id}_{\mathbf{L}(A)}$, hence a covering morphism. Since $\mathbf{L}(A)$ is simply connected, we have $\mathbf{L}(A) \cong \widetilde{A}$ and ker $\exp _{A} \cong \pi_{1}(A)$ is the fundamental group of $A$.

As special cases we obtain in particular the finite-dimensional tori

$$
\mathbb{T}^{d} \cong \mathbb{R}^{d} / \mathbb{Z}^{d} \quad \text { with } \quad \pi_{1}\left(\mathbb{T}^{n}\right) \cong \mathbb{Z}^{n}
$$

If we want to classify all connected abelian Lie groups $A$ of dimension $n$, we can now proceed as follows. First, we note that $\widetilde{A} \cong \mathbf{L}(A) \cong\left(\mathbb{R}^{n},+\right)$ as abelian Lie groups. Then $\operatorname{Aut}(\widetilde{A}) \cong \mathrm{GL}_{n}(\mathbb{R})$ follows from the Automatic Smoothness Theorem 9.2.16. Further, Exercise 9.3.4 implies that the discrete subgroup $\pi_{1}(A)$ of $\widetilde{A} \cong \mathbb{R}^{n}$ can be mapped by some $\varphi \in \mathrm{GL}_{n}(\mathbb{R})$ onto

$$
\mathbb{Z}^{k} \cong \mathbb{Z}^{k} \times\{0\} \subseteq \mathbb{R}^{k} \times \mathbb{R}^{n-k} \cong \mathbb{R}^{n}
$$

Therefore,

$$
A \cong \mathbb{R}^{n} / \mathbb{Z}^{k} \cong \mathbb{T}^{k} \times \mathbb{R}^{n-k}
$$

and it is clear that the number $k$ is an isomorphy invariant of the Lie group $A$, namely, the rank of its fundamental subgroup. Therefore, connected abelian Lie groups $A$ are determined up to isomorphism by the pair ( $n, k$ ), where $n=\operatorname{dim} A$ and $k=\operatorname{rk} \pi_{1}(A)$. The case where $n=k$ gives the compact connected abelian Lie groups. The above argument shows that such groups are always of the form $A \cong \mathbf{L}(A) / \Gamma$, where $\Gamma$ is a discrete subgroup of $\mathbf{L}(T)$ generated by a basis for $\mathbf{L}(T)$. Such discrete subgroups are called lattices.

Example 9.5.7. We show that

$$
\pi_{1}\left(\mathrm{SO}_{3}(\mathbb{R})\right) \cong C_{2}=\{ \pm 1\}
$$

by constructing a surjective homomorphism

$$
\varphi: \mathrm{SU}_{2}(\mathbb{C}) \rightarrow \mathrm{SO}_{3}(\mathbb{R})
$$

with $\operatorname{ker} \varphi=\{ \pm \mathbf{1}\}$. Since $\mathrm{SU}_{2}(\mathbb{C})$ is homeomorphic to $\mathbb{S}^{3}$, it is simply connected (Exercise A.1.3), so that we then obtain $\pi_{1}\left(\mathrm{SO}_{3}(\mathbb{R})\right) \cong C_{2}$ (Theorem 9.5.4).

We consider

$$
\mathfrak{s} \mathfrak{u}_{2}(\mathbb{C})=\left\{x \in \mathfrak{g l}_{2}(\mathbb{C}): x^{*}=-x, \operatorname{tr} x=0\right\}=\left\{\left(\begin{array}{cc}
a i & b \\
-\bar{b} & -a i
\end{array}\right): b \in \mathbb{C}, a \in \mathbb{R}\right\}
$$

and observe that this is a three-dimensional real subspace of $\mathfrak{g} \mathfrak{l}_{2}(\mathbb{C})$. We obtain on $E:=\mathfrak{s u}_{2}(\mathbb{C})$ the structure of a euclidian vector space by the scalar product

$$
\beta(x, y):=\operatorname{tr}\left(x y^{*}\right)=-\operatorname{tr}(x y) .
$$

Now we consider the adjoint representation

$$
\mathrm{Ad}: \mathrm{SU}_{2}(\mathbb{C}) \rightarrow \mathrm{GL}(E), \quad \mathrm{Ad}(g)(x)=g x g^{-1} .
$$

Then we have for $x, y \in E$ and $g \in \mathrm{SU}_{2}(\mathbb{C})$ the relation

$$
\begin{aligned}
\beta(\operatorname{Ad}(g) x, \operatorname{Ad}(g) y) & =\operatorname{tr}\left(g x g^{-1}\left(g y g^{-1}\right)^{*}\right)=\operatorname{tr}\left(g x g^{-1}\left(g^{-1}\right)^{*} y^{*} g^{*}\right) \\
& =\operatorname{tr}\left(g x g^{-1} g y^{*} g^{-1}\right)=\operatorname{tr}\left(x y^{*}\right)=\beta(x, y)
\end{aligned}
$$

This means that

$$
\operatorname{Ad}\left(\mathrm{SU}_{2}(\mathbb{C})\right) \subseteq \mathrm{O}(E, \beta) \cong \mathrm{O}_{3}(\mathbb{R})
$$

Since $\mathrm{SU}_{2}(\mathbb{C})$ is connected, we further obtain $\operatorname{Ad}\left(\mathrm{SU}_{2}(\mathbb{C})\right) \subseteq \mathrm{SO}(E, \beta) \cong \mathrm{SO}_{3}(\mathbb{R})$, the identity component of $\mathrm{O}(E, \beta)$.

The derived representation is given by

$$
\operatorname{ad}: \mathfrak{s u}_{2}(\mathbb{C}) \rightarrow \mathfrak{s o}(E, \beta) \cong \mathfrak{s o}_{3}(\mathbb{R}), \quad \operatorname{ad}(x)(y)=[x, y]
$$

If $\operatorname{ad} x=0$, then $\operatorname{ad} x(i \mathbf{1})=0$ implies that $\operatorname{ad} x\left(\mathfrak{u}_{2}(\mathbb{C})\right)=\{0\}$, so that $\operatorname{ad} x\left(\mathfrak{g} \mathfrak{l}_{2}(\mathbb{C})\right)=\{0\}$ follows from $\mathfrak{g} \mathfrak{l}_{2}(\mathbb{C})=\mathfrak{u}_{2}(\mathbb{C})+i \mathfrak{u}_{2}(\mathbb{C})$. This implies that $x \in \mathbb{C} 1$, so that $\operatorname{tr} x=0$ leads to $x=0$. Hence ad is injective, and we conclude with $\operatorname{dim} \mathfrak{s o}(E, \beta)=\operatorname{dim} \mathfrak{s o}_{3}(\mathbb{R})=3$ that

$$
\operatorname{ad}\left(\mathfrak{s} \mathfrak{u}_{2}(\mathbb{C})\right)=\mathfrak{s} \mathfrak{o}(E, \beta)
$$

Therefore,

$$
\operatorname{Im} \operatorname{Ad}=\langle\exp \mathfrak{s o}(E, \beta)\rangle=\operatorname{SO}(E, \beta)_{0}=\operatorname{SO}(E, \beta)
$$

We thus obtain a surjective homomorphism

$$
\varphi: \mathrm{SU}_{2}(\mathbb{C}) \rightarrow \mathrm{SO}_{3}(\mathbb{R})
$$

Since $\mathrm{SU}_{2}(\mathbb{C})$ is compact, the quotient group $\mathrm{SU}_{2}(\mathbb{C}) / \operatorname{ker} \varphi$ is also compact, and the induced bijective morphism $\mathrm{SU}_{2}(\mathbb{C}) / \operatorname{ker} \varphi \rightarrow \mathrm{SO}_{3}(\mathbb{R})$ is a homeomorphism, hence an isomorphism of topological groups.

We further have

$$
\operatorname{ker} \varphi=Z\left(\mathrm{SU}_{2}(\mathbb{C})\right)=\{ \pm \mathbf{1}\}
$$

(Exercise 9.3.8), so that

$$
\widetilde{\mathrm{SO}}_{3}(\mathbb{R}) \cong \mathrm{SU}_{2}(\mathbb{C}) \quad \text { and } \quad \pi_{1}\left(\mathrm{SO}_{3}(\mathbb{R})\right) \cong C_{2}
$$

\subsubsection*{The Monodromy Principle and Its Applications}
To round off the picture, we still have to provide the link between Lie algebras and covering groups. The main point is that, in general, one cannot integrate morphisms of Lie algebras $\mathbf{L}(G) \rightarrow \mathbf{L}(H)$ to morphisms of the corresponding groups $G \rightarrow H$ if $G$ is not simply connected.

Proposition 9.5.8 (Monodromy Principle). Let $G$ be a connected simply connected Lie group and $H$ a group. Let $V$ be an open symmetric connected identity neighborhood in $G$ and $f: V \rightarrow H$ a function with

$$
f(x y)=f(x) f(y) \quad \text { for } \quad x, y, x y \in V
$$

Then there exists a unique group homomorphism extending $f$. If, in addition, $H$ is a Lie group and $f$ is smooth, then its extension is also smooth.

Proof. We consider the group $G \times H$ and the subgroup $S \subseteq G \times H$ generated by the subset $U:=\{(x, f(x)): x \in V\}$. We endow $U$ with the topology for which $x \mapsto(x, f(x)), V \rightarrow U$ is a homeomorphism. Note that $f(\mathbf{1})^{2}=f\left(\mathbf{1}^{2}\right)= f(\mathbf{1})$ implies $f(\mathbf{1})=\mathbf{1}$, which further leads to $\mathbf{1}=f\left(x x^{-1}\right)=f(x) f\left(x^{-1}\right)$, so that $f\left(x^{-1}\right)=f(x)^{-1}$. Hence $U=U^{-1}$.

We now apply Lemma 9.4.3 because $S$ is generated by $U$, and (T1)-(T2) directly follow from the corresponding properties of $V$ and $(x, f(x))(y, f(y))=$\\
$(x y, f(x y))$ for $x, y, x y \in V$. This leads to a group topology on $S$, for which $S$ is a connected topological group. Indeed, its connectedness follows from $S=\bigcup_{n \in \mathbb{N}} U^{n}$ and the connectedness of all sets $U^{n}$ (Exercise 9.4.5). The projection $p_{G}: G \times H \rightarrow G$ induces a covering homomorphism $q: S \rightarrow G$ because its restriction to the open 1-neighborhood $U$ is a homeomorphism (Exercise A.2.2(c)), and the connectedness of $S$ and the simple connectedness of $G$ imply that $q$ is a homeomorphism (Corollary A.2.8). Now $F:=p_{H} \circ q^{-1}: G \rightarrow H$ provides the required extension of $f$. In fact, for $x \in U$ we have $q^{-1}(x)=(x, f(x))$, and therefore $F(x)=f(x)$.

If, in addition, $H$ is Lie and $f$ is smooth, then the smoothness of the extension follows directly from Corollary 9.2.11.

Theorem 9.5.9 (Integrability Theorem for Lie Algebra Homomorphisms). Let $G$ be a connected simply connected Lie group, $H$ a Lie group, and $\psi: \mathbf{L}(G) \rightarrow \mathbf{L}(H)$ a Lie algebra morphism. Then there exists a unique morphism $\varphi: G \rightarrow H$ with $\mathbf{L}(\varphi)=\psi$.

Proof. Let $U \subseteq \mathbf{L}(G)$ be an open connected symmetric 0-neighborhood on which the BCH -product is defined and satisfies $\exp _{G}(x * y)=\exp _{G}(x) \exp _{G}(y)$ and $\exp _{H}(\psi(x) * \psi(y))=\exp _{H}(\psi(x)) \exp _{H}(\psi(y))$ for $x, y \in U$ (Theorem 9.4.8). Assume further that $\left.\exp _{G}\right|_{U}$ is a homeomorphism onto an open subset of $G$ (cf. Proposition 9.2.5).

The continuity of $\psi$ and the fact that $\psi$ is a Lie algebra homomorphism imply that for $x, y \in U$ the element $\psi(x * y)$ coincides with the convergent Hausdorff series $\psi(x) * \psi(y)$. We define

$$
f: \exp _{G}(U) \rightarrow H, \quad f\left(\exp _{G}(x)\right):=\exp _{H}(\psi(x)) .
$$

For $x, y, x * y \in U$, we then obtain

$$
\begin{aligned}
& f\left(\exp _{G}(x) \exp _{G}(y)\right)=f\left(\exp _{G}(x * y)\right)=\exp _{H}(\psi(x * y)) \\
= & \exp _{H}(\psi(x) * \psi(y))=\exp _{H}(\psi(x)) \exp _{H}(\psi(y))=f\left(\exp _{G}(x)\right) f\left(\exp _{G}(y)\right)
\end{aligned}
$$

Then $f: \exp (U) \rightarrow H$ satisfies the assumptions of Proposition 9.5.8, and we see that $f$ extends uniquely to a group homomorphism $\varphi: G \rightarrow H$. Since $\exp _{G}$ is a local diffeomorphism, $f$ is smooth in a $\mathbf{1}$-neighborhood, and therefore $\varphi$ is smooth. We finally observe that $\varphi$ is uniquely determined by $\mathbf{L}(\varphi)=\psi$ because $G$ is connected (Corollary 9.2.12).

The following corollary can be viewed as an integrability condition for $\psi$.\\
Corollary 9.5.10. If $G$ is a connected Lie group and $H$ is a Lie group, then for a Lie algebra morphism $\psi: \mathbf{L}(G) \rightarrow \mathbf{L}(H)$ there exists a morphism $\varphi: G \rightarrow H$ with $\mathbf{L}(\varphi)=\psi$ if and only if $\pi_{1}(G) \subseteq \operatorname{ker} \widetilde{\varphi}$, where $\pi_{1}(G)$ is identified with the kernel of the universal covering map $q_{G}: \widetilde{G} \rightarrow G$ and $\widetilde{\varphi}: \widetilde{G} \rightarrow H$ is the unique morphism with $\mathbf{L}(\widetilde{\varphi})=\psi \circ \mathbf{L}\left(q_{G}\right)$.

Proof. If $\varphi$ exists, then

$$
\left(\varphi \circ q_{G}\right) \circ \exp _{\widetilde{G}}=\varphi \circ \exp _{G} \circ \mathbf{L}\left(q_{G}\right)=\exp _{H} \circ \psi \circ \mathbf{L}\left(q_{G}\right)
$$

and the uniqueness of $\widetilde{\varphi}$ imply that $\widetilde{\varphi}=\varphi \circ q_{G}$, and hence that $\pi_{1}(G)= \operatorname{ker} q_{G} \subseteq \operatorname{ker} \widetilde{\varphi}$.

If, conversely, $\operatorname{ker} q_{G} \subseteq \operatorname{ker} \widetilde{\varphi}$, then $\varphi\left(q_{G}(g)\right):=\widetilde{\varphi}(g)$ defines a continuous morphism $G \cong \widetilde{G} / \operatorname{ker} q_{G} \rightarrow H$ with $\varphi \circ q_{G}=\widetilde{\varphi}$ (Exercise 9.3.8) and

$$
\varphi \circ \exp _{G} \circ \mathbf{L}\left(q_{G}\right)=\varphi \circ q_{G} \circ \exp _{\widetilde{G}}=\widetilde{\varphi} \circ \exp _{\widetilde{G}}=\exp _{H} \circ \psi \circ \mathbf{L}\left(q_{G}\right) .
$$

We recall that a Lie group $G$ is called 1 -connected if it is connected and simply connected.

Corollary 9.5.11. If $G$ is a 1 -connected Lie group with Lie algebra $\mathfrak{g}$, then the map

$$
\mathbf{L}: \operatorname{Aut}(G) \rightarrow \operatorname{Aut}(\mathfrak{g})
$$

is an isomorphism of groups. As a closed subgroup of $\operatorname{GL}(\mathfrak{g})$, the group $\operatorname{Aut}(\mathfrak{g})$ carries a natural Lie group structure, and we endow $\operatorname{Aut}(G)$ with the Lie group structure for which $\mathbf{L}$ is an isomorphism of Lie groups. Then the action of $\operatorname{Aut}(G)$ on $G$ is smooth.

Proof. First, we recall from Corollary 9.1.10 that for each automorphism $\varphi \in \operatorname{Aut}(G)$ the endomorphism $\mathbf{L}(\varphi)$ of $\mathfrak{g}$ also is an automorphism. That $\mathbf{L}$ is injective follows from the connectedness of $G$ (Corollary 9.2.12) and that $\mathbf{L}$ is surjective from the Integrability Theorem 9.5.9.

If we endow $\operatorname{Aut}(G)$ with the Lie group structure for which $\mathbf{L}$ is an isomorphism, then the smoothness of the action of $\operatorname{Aut}(\mathfrak{g})$ on $\mathfrak{g}$ and Lemma 9.2.26(i) imply that the action of $\operatorname{Aut}(G)$ on $G$ is smooth.

Remark 9.5.12. With the Integrability Theorem 9.5.9 we can give a second proof of Lie's Third Theorem which does not require Ado's Theorem.

From the structure theory of Lie algebras, we know that $\mathfrak{g}$ is a semidirect sum $\mathfrak{g}=\mathfrak{r} \rtimes_{\beta} \mathfrak{s}$, where $\mathfrak{r}$ is solvable and $\mathfrak{s}$ is semisimple. Since $\mathfrak{s}$ is semisimple, $\mathfrak{z}(\mathfrak{s})=\{0\}$ and the adjoint representation of $\mathfrak{s}$ is faithful. Therefore, the Integral Subgroup Theorem 9.4.8 implies the existence of a connected Lie group $S$ with $\mathbf{L}(S)=\mathfrak{s}$. Now Theorem 9.5.9 yields a homomorphism $\alpha: \widetilde{S} \rightarrow \operatorname{Aut}(\mathfrak{r})$ of the simply connected covering group $\widetilde{S}$ of $S$ integrating the adjoint representation of $\mathfrak{s}$ on the ideal $\mathfrak{r}$ of $\mathfrak{g}$.

Next we recall that the solvable Lie algebra $\mathfrak{r}$ can be written as an iterated semidirect sum

$$
\mathfrak{r}=\mathfrak{r}_{1} \rtimes_{\beta_{1}} \mathbb{R}, \quad \mathfrak{r}_{1}=\mathfrak{r}_{2} \rtimes_{\beta_{2}} \mathbb{R}, \quad \text { etc. }
$$

If $r=\operatorname{dim} \mathfrak{r}$, it follows inductively that there exists a simply connected Lie group $R_{j}$, diffeomorphic to $\mathbb{R}^{r-j}$, with $\mathbf{L}\left(R_{j}\right)=\mathfrak{r}_{j}, j=r, r-1, \ldots, 1$. In view of $\operatorname{Aut}\left(R_{j}\right) \cong \operatorname{Aut}\left(\mathfrak{r}_{j}\right)$, the action $\beta_{j}$ of $\mathbb{R}$ on $\mathfrak{r}_{j}$ integrates to a smooth action\\
$\alpha_{j}: \mathbb{R} \rightarrow \operatorname{Aut}\left(R_{j}\right)$ (Corollary 9.5.11), so that $R_{j-1}:=R_{j} \rtimes_{\alpha_{j}} \mathbb{R}$ is a Lie group diffeomorphic to $\mathbb{R}^{k-j+1}$ with Lie algebra $\mathfrak{r}_{j-1}$ (Proposition 9.2.25). For $j=1$, we obtain a simply connected Lie group $R$ with $\mathbf{L}(R)=\mathfrak{r}$. Again we use Corollary 9.5.11 to see that the smooth action of $\widetilde{S}$ on $\mathfrak{r}$ induces a smooth action of $\widetilde{S}$ on $R$, so that the semidirect product $G:=R \rtimes \widetilde{S}$ is a simply connected Lie group with Lie algebra $\mathfrak{r} \rtimes \mathfrak{s}=\mathfrak{g}$ (Proposition 9.2.25).

\subsubsection*{Classification of Lie Groups with Given Lie Algebra}
Let $G$ and $H$ be linear Lie groups. If $\varphi: G \rightarrow H$ is an isomorphism, then the functoriality of $\mathbf{L}$ directly implies that $\mathbf{L}(\varphi): \mathbf{L}(G) \rightarrow \mathbf{L}(H)$ is an isomorphism. In fact, if $\psi: H \rightarrow G$ is a morphism with $\varphi \circ \psi=\operatorname{id}_{H}$ and $\psi \circ \varphi=\operatorname{id}_{G}$, then

$$
\mathrm{id}_{\mathbf{L}(H)}=\mathbf{L}\left(\mathrm{id}_{H}\right)=\mathbf{L}(\varphi \circ \psi)=\mathbf{L}(\varphi) \circ \mathbf{L}(\psi)
$$

and likewise $\mathbf{L}(\psi) \circ \mathbf{L}(\varphi)=\operatorname{id}_{\mathbf{L}(G)}$.\\
In this subsection, we ask to which extent a Lie group $G$ is determined by its Lie algebra $\mathbf{L}(G)$.

Theorem 9.5.13. Two connected Lie groups $G$ and $H$ have isomorphic Lie algebras if and only if their universal covering groups $\widetilde{G}$ and $\widetilde{H}$ are isomorphic.

Proof. If $\widetilde{G}$ and $\widetilde{H}$ are isomorphic, then we clearly have

$$
\mathbf{L}(G) \cong \mathbf{L}(\widetilde{G}) \cong \mathbf{L}(\widetilde{H}) \cong \mathbf{L}(H)
$$

(cf. Proposition 9.5.1).\\
Conversely, let $\psi: \mathbf{L}(G) \rightarrow \mathbf{L}(H)$ be an isomorphism. Using Theorem 9.5.9, we obtain a unique morphism $\varphi: \widetilde{G} \rightarrow \widetilde{H}$ with $\mathbf{L}(\varphi)=\psi$ and also a unique morphism $\widehat{\varphi}: \widetilde{H} \rightarrow \widetilde{G}$ with $\mathbf{L}(\widehat{\varphi})=\psi^{-1}$. Then $\mathbf{L}(\varphi \circ \widehat{\varphi})=\operatorname{id}_{\mathbf{L}(\widetilde{G})}$ implies $\varphi \circ \widehat{\varphi}=\operatorname{id}_{\widetilde{G}}$, and likewise $\widehat{\varphi} \circ \varphi=\operatorname{id}_{\widetilde{H}}$. Therefore, $\widetilde{G}$ and $\widetilde{H}$ are isomorphic Lie groups.

Combining the preceding theorem with Theorem 9.5.4, we obtain:\\
Corollary 9.5.14. Let $G$ be a connected Lie group and $q: \widetilde{G} \rightarrow G$ the universal covering morphism of connected Lie groups. Then for each discrete central subgroup $\Gamma \subseteq \widetilde{G}$, the group $\widetilde{G} / \Gamma$ is a connected Lie group with $\mathbf{L}(\widetilde{G} / \Gamma) \cong \mathbf{L}(G)$ and, conversely, each Lie group with the same Lie algebra as $G$ is isomorphic to some quotient $\widetilde{G} / \Gamma$.

Example 9.5.15. We now describe a pair of nonisomorphic Lie groups with isomorphic Lie algebras and isomorphic fundamental groups.

Let

$$
\widetilde{G}:=\mathrm{SU}_{2}(\mathbb{C}) \times \mathrm{SU}_{2}(\mathbb{C})
$$

whose center is $C_{2} \times C_{2}$,

$$
G:=\widetilde{G} /\left(C_{2} \times\{\mathbf{1}\}\right) \cong \mathrm{SO}_{3}(\mathbb{R}) \times \mathrm{SU}_{2}(\mathbb{C})
$$

and

$$
H:=\widetilde{G} /\{(\mathbf{1}, \mathbf{1}),(-\mathbf{1},-\mathbf{1})\} \cong \mathrm{SO}_{4}(\mathbb{R})
$$

where the latter isomorphy follows from Proposition 9.5.21 below. Then $\pi_{1}(G) \cong \pi_{1}(H) \cong C_{2}$, but there is no automorphism of $\widetilde{G}$ mapping $\pi_{1}(G)$ to $\pi_{1}(H)$.

Indeed, one can show that the two direct factors are the only nontrivial connected normal subgroups of $\widetilde{G}$, so that each automorphism of $\widetilde{G}$ either preserves both or exchanges them. Since $\pi_{1}(H)$ is not contained in any of them, it cannot be mapped to $\pi_{1}(G)$ by an automorphism of $\widetilde{G}$.\\
Examples 9.5.16. Here are some examples of pairs of linear Lie groups with isomorphic Lie algebras:\\
(1) $G=\mathrm{SO}_{3}(\mathbb{R})$ and $\widetilde{G} \cong \mathrm{SU}_{2}(\mathbb{C})$ (Example 9.5.7).\\
(2) $G=\mathrm{SO}_{2,1}(\mathbb{R})_{0}$ and $H=\mathrm{SL}_{2}(\mathbb{R})$ : In this case, we actually have a covering morphism $\varphi: H \rightarrow G$ coming from the adjoint representation

$$
\mathrm{Ad}: \mathrm{SL}_{2}(\mathbb{R}) \rightarrow \mathrm{GL}(\mathbf{L}(H)) \cong \mathrm{GL}_{3}(\mathbb{R})
$$

On $\mathbf{L}(H)=\mathfrak{s} \mathfrak{l}_{2}(\mathbb{R})$, we consider the symmetric bilinear form given by $\beta(x, y):=\frac{1}{2} \operatorname{tr}(x y)$ and the basis

$$
e_{1}:=\left(\begin{array}{cc}
1 & 0 \\
0 & -1
\end{array}\right), \quad e_{2}:=\left(\begin{array}{cc}
0 & 1 \\
1 & 0
\end{array}\right), \quad e_{3}:=\left(\begin{array}{cc}
0 & 1 \\
-1 & 0
\end{array}\right) .
$$

Then the matrix $B$ of $\beta$ with respect to this basis is

$$
B:=\left(\begin{array}{ccc}
1 & 0 & 0 \\
0 & 1 & 0 \\
0 & 0 & -1
\end{array}\right)
$$

One easily verifies that

$$
\operatorname{Im} \mathrm{Ad} \subseteq \mathrm{O}(\mathbf{L}(H), \beta) \cong \mathrm{O}_{2,1}(\mathbb{R})
$$

and since ad: $\mathbf{L}(H) \rightarrow \mathfrak{o}_{2,1}(\mathbb{R})$ is injective between spaces of the same dimension 3 (Exercise), it is bijective. Therefore, im $\operatorname{Ad}=\left\langle\exp \mathfrak{o}_{2,1}(\mathbb{R})\right\rangle= \mathrm{SO}_{2,1}(\mathbb{R})_{0}$ and Proposition 9.5.1 imply that

$$
\mathrm{Ad}: \mathrm{SL}_{2}(\mathbb{R}) \rightarrow \mathrm{SO}_{2,1}(\mathbb{R})_{0}
$$

is a covering morphism. Its kernel is given by $Z\left(\mathrm{SL}_{2}(\mathbb{R})\right)=\{ \pm \mathbf{1}\}$.\\
From the polar decomposition, one derives that both groups are homeomorphic to $\mathbb{T} \times \mathbb{R}^{2}$, and topologically the map Ad is like $(z, x, y) \mapsto\left(z^{2}, x, y\right)$, a two-fold covering.\\
(3) $G=\mathrm{SL}_{2}(\mathbb{C})$ and $H=\mathrm{SO}_{3,1}(\mathbb{R})_{0}$ : This will be explained in detail and greater generality in Example 17.2.6.

Example 9.5.17. Let $G=\mathrm{SL}_{2}(\mathbb{R})$ and $H=\mathrm{SO}_{2,1}(\mathbb{R})_{0}$ and recall that $\widetilde{G} \cong \widetilde{H}$ follows from $\mathfrak{s} \mathfrak{l}_{2}(\mathbb{R}) \cong \mathfrak{s o}_{2,1}(\mathbb{R})$ (cf. Example 9.5.16).

We further have $q_{G}(Z(\widetilde{G})) \subseteq Z(G)=\{ \pm \mathbf{1}\}$ and $\pi_{1}(G)=\operatorname{ker} q_{G} \subseteq Z(\widetilde{G})$ (cf. Proposition 9.5.2). Likewise $q_{H}(Z(\widetilde{G})) \subseteq Z(H)=\{\boldsymbol{1}\}$ implies

$$
Z(\widetilde{G}) \cong \pi_{1}(H) \cong \pi_{1}\left(\mathrm{O}_{2}(\mathbb{R}) \times \mathrm{O}_{1}(\mathbb{R})\right) \cong \mathbb{Z}
$$

where the latter is a consequence of the polar decomposition. This implies that $Z(\widetilde{G}) \cong \mathbb{Z}$, where

$$
\pi_{1}(G) \cong 2 \mathbb{Z} \quad \text { and } \quad \pi_{1}(H) \cong \mathbb{Z}=Z(\widetilde{G})
$$

Therefore, $G$ and $H$ are not isomorphic, but they have isomorphic Lie algebras and isomorphic fundamental groups.

\subsubsection*{Nonlinear Lie Groups}
We have already seen how to describe all connected Lie groups with a given Lie algebra. Determining all such groups which are, in addition, linear turns out to be a much more subtle enterprise. If $\widetilde{G}$ is a simply connected group with a given Lie algebra, it means determining which of the groups $\widetilde{G} / D$ are linear. As the following examples show, the answer to this problem is not easy. In fact, a complete answer requires detailed knowledge of the structure of finite-dimensional Lie algebras.

Example 9.5.18. We show that the universal covering group $G:=\widetilde{\mathrm{SL}}_{2}(\mathbb{R})$ of $\mathrm{SL}_{2}(\mathbb{R})$ is not a linear Lie group. Moreover, we show that every continuous homomorphism $\varphi: G \rightarrow \mathrm{GL}_{n}(\mathbb{R})$ satisfies $D:=\pi_{1}\left(\mathrm{SL}_{2}(\mathbb{R})\right) \subseteq \operatorname{ker} \varphi$, hence factors through $G / D \cong \mathrm{SL}_{2}(\mathbb{R})$.

We consider the Lie algebra homomorphism $\mathbf{L}(\varphi): \mathfrak{s l}_{2}(\mathbb{R}) \rightarrow \mathfrak{g l} \mathfrak{l}_{n}(\mathbb{R})$. Then it is easy to see that

$$
\mathbf{L}(\varphi)_{\mathbb{C}}(x+i y):=\mathbf{L}(\varphi) x+i \mathbf{L}(\varphi) y
$$

defines an extension of $\mathbf{L}(\varphi)$ to a complex linear Lie algebra homomorphism

$$
\mathbf{L}(\varphi)_{\mathbb{C}}: \mathfrak{s} \mathfrak{l}_{2}(\mathbb{C}) \rightarrow \mathfrak{g} \mathfrak{l}_{n}(\mathbb{C})
$$

Since the group $\mathrm{SL}_{2}(\mathbb{C})$ is simply connected, there exists a unique group homomorphism $\psi: \mathrm{SL}_{2}(\mathbb{C}) \rightarrow \mathrm{GL}_{n}(\mathbb{C})$ with $\mathbf{L}(\psi)=\mathbf{L}(\varphi)_{\mathbb{C}}$.

Let $\alpha: G \rightarrow G / D \cong \mathrm{SL}_{2}(\mathbb{R}) \rightarrow \mathrm{SL}_{2}(\mathbb{C})$ be the canonical morphism. Then

$$
\mathbf{L}(\varphi)=\mathbf{L}(\varphi)_{\mathbb{C}} \circ \mathbf{L}(\alpha)=\mathbf{L}(\psi) \circ \mathbf{L}(\alpha)=\mathbf{L}(\psi \circ \alpha)
$$

implies $\varphi=\psi \circ \alpha$. We conclude that $\operatorname{ker} \varphi \supseteq \operatorname{ker} \alpha=D$. Therefore, $G$ has no faithful linear representation.

Lemma 9.5.19. If $A$ is a Banach algebra with unit $\mathbf{1}$ and $p, q \in A$ with $[p, q]=\lambda \mathbf{1}$, then $\lambda=0$.

Proof. ([Wie49]) By induction we obtain

$$
\left[p, q^{n}\right]=\lambda n q^{n-1} \quad \text { for } \quad n \in \mathbb{N} .
$$

In fact,

$$
\left[p, q^{n+1}\right]=[p, q] q^{n}+q\left[p, q^{n}\right]=\lambda q^{n}+\lambda n q^{n}=\lambda(n+1) q^{n} .
$$

Therefore,

$$
|\lambda| n\left\|q^{n-1}\right\| \leq 2\|p\|\left\|q^{n}\right\| \leq 2\|p\|\|q\|\left\|q^{n-1}\right\|
$$

for each $n \in \mathbb{N}$, which leads to

$$
(|\lambda| n-2\|p\|\|q\|)\left\|q^{n-1}\right\| \leq 0 .
$$

If $\lambda \neq 0$, then we obtain for sufficiently large $n$ that $q^{n-1}=0$. For $n>1$, we derive from (9.15) that $q^{n-2}=0$. Inductively we arrive at the contradiction $q=0$. \(\square\)

If $A$ is a finite-dimensional algebra, we may w.l.o.g. assume that it is a subalgebra of some matrix algebra $M_{n}(\mathbb{K})$, and then $[p, q]=\lambda \mathbf{1}$ implies

$$
n \lambda=\operatorname{tr}(\lambda \mathbf{1})=\operatorname{tr}([p, q])=0
$$

so that $\lambda=0$.\\
Example 9.5.20. We consider the three-dimensional Heisenberg group

$$
\begin{aligned}
& G=\left\{\left(\begin{array}{ccc}
1 & x & y \\
0 & 1 & z \\
0 & 0 & 1
\end{array}\right): x, y, z \in \mathbb{R}\right\} \\
& \text { with } \quad \mathbf{L}(G)=\left\{\left(\begin{array}{ccc}
0 & x & y \\
0 & 0 & z \\
0 & 0 & 0
\end{array}\right): x, y, z \in \mathbb{R}\right\} .
\end{aligned}
$$

Note that $\exp _{G}: \mathbf{L}(G) \rightarrow G$ is a diffeomorphism whose inverse is given by

$$
\log (g)=(g-\mathbf{1})-\frac{1}{2}(g-\mathbf{1})^{2}
$$

(Proposition 3.3.3). Let

$$
z:=\left(\begin{array}{lll}
0 & 0 & 1 \\
0 & 0 & 0 \\
0 & 0 & 0
\end{array}\right), \quad p:=\left(\begin{array}{lll}
0 & 1 & 0 \\
0 & 0 & 0 \\
0 & 0 & 0
\end{array}\right), \quad \text { and } \quad q:=\left(\begin{array}{lll}
0 & 0 & 0 \\
0 & 0 & 1 \\
0 & 0 & 0
\end{array}\right) .
$$

Then $[p, q]=z,[p, z]=[q, z]=0, \exp \mathbb{R} z=\mathbf{1}+\mathbb{R} z \subseteq Z(G)$, and $D:=\exp (\mathbb{Z} z)$ is a discrete central subgroup of $G$. We claim that the group $G / D$ is not a linear Lie group. This will be verified by showing that each homomorphism $\alpha: G \rightarrow \mathrm{GL}_{n}(\mathbb{C})$ with $D \subseteq \operatorname{ker} \alpha$ satisfies $\exp (\mathbb{R} z) \subseteq \operatorname{ker} \alpha$.

The map $\mathbf{L}(\alpha): \mathbf{L}(G) \rightarrow \mathfrak{g l}_{n}(\mathbb{C})$ is a Lie algebra homomorphism and we obtain linear maps

$$
P:=\mathbf{L}(\alpha)(p), \quad Q:=\mathbf{L}(\alpha)(q), \quad \text { and } \quad Z:=\mathbf{L}(\alpha)(z)
$$

with $[P, Q]=Z$. Now $\exp _{G} z \in D=\operatorname{ker} \alpha$ implies that $e^{Z}=\alpha(\exp z)=\mathbf{1}$, and hence that $Z$ is diagonalizable with all eigenvalues contained in $2 \pi i \mathbb{Z}$ (Exercise 3.2.12). Let $V_{\lambda}:=\operatorname{ker}(Z-\lambda \mathbf{1})$. Since $z$ is central in $\mathbf{L}(G)$, the space $V_{\lambda}$ is invariant under $G$ (Exercise 2.1.1), hence also under $\mathbf{L}(G)$ (Exercise 4.2.4). Therefore, the restrictions $P_{\lambda}:=\left.P\right|_{V_{\lambda}}$ and $Q_{\lambda}:=\left.Q\right|_{V_{\lambda}}$ satisfy $\left[P_{\lambda}, Q_{\lambda}\right]=\lambda$ id in the Banach algebra End $\left(V_{\lambda}\right)$. In view of the preceding lemma, we have $\lambda=0$. Therefore, the diagonalizability of $Z$ entails that $Z=0$ and hence that $\mathbb{R} z \subseteq \operatorname{ker} \mathbf{L}(\alpha)$. It follows in particular that the group $G / D$ has no faithful linear representation.

\subsubsection*{The Quaternions, $\mathrm{SU}_{2}(\mathbb{C})$ and $\mathrm{SO}_{4}(\mathbb{R})$}
In this subsection, we shall use the quaternion algebra $\mathbb{H}$ (Section 2.3) to get some more information on the structure of the group $\mathrm{SO}_{4}(\mathbb{R})$. Here the idea is to identify $\mathbb{R}^{4}$ with $\mathbb{H}$.

Proposition 9.5.21. There exists a covering homomorphism

$$
\varphi: \mathrm{SU}_{2}(\mathbb{C}) \times \mathrm{SU}_{2}(\mathbb{C}) \rightarrow \mathrm{SO}_{4}(\mathbb{R}) \subseteq \mathrm{GL}(\mathbb{H}), \quad \varphi(a, b) . x=a x b^{-1}
$$

This homomorphism is a universal covering with $\operatorname{ker} \varphi=\{ \pm(\mathbf{1}, \mathbf{1})\}$.\\
Proof. Since $|a|=|b|=1$, all the maps $\varphi(a, b): \mathbb{H} \rightarrow \mathbb{H}$ are orthogonal, so that $\varphi$ is a homomorphism

$$
\mathrm{SU}_{2}(\mathbb{C}) \times \mathrm{SU}_{2}(\mathbb{C}) \rightarrow \mathrm{O}_{4}(\mathbb{R})
$$

Since $\mathrm{SU}_{2}(\mathbb{C}) \times \mathrm{SU}_{2}(\mathbb{C})$ is connected, it further follows that $\operatorname{im}(\varphi) \subseteq \mathrm{SO}_{4}(\mathbb{R})$.\\
To determine the kernel of $\varphi$, suppose that $\varphi(a, b)=\operatorname{id}_{\mathbb{H}}$. Then $a x b^{-1}=x$ for all $x \in \mathbb{H}$. For $x=b$, we obtain in particular $a=b$. Hence $a x=x a$ for all $x \in \mathbb{H}$. With $x=I$ and $x=J$, this leads to $a \in \mathbb{R} \mathbf{1}$, and hence to $(a, b) \in\{ \pm(\mathbf{1}, \mathbf{1})\}$. This proves the assertion on $\operatorname{ker} \varphi$.

The derived representation is given by

$$
\mathbf{L}(\varphi): \mathfrak{s u}_{2}(\mathbb{C}) \times \mathfrak{s u}_{2}(\mathbb{C}) \rightarrow \mathfrak{s o}_{4}(\mathbb{R}), \quad \mathbf{L}(\varphi)(x, y)(z)=x z-z y
$$

Since $\operatorname{ker} \varphi$ is discrete, it follows that $\operatorname{ker} \mathbf{L}(\varphi) \subseteq \mathbf{L}(\operatorname{ker} \varphi)=\{0\}$. Hence $\mathbf{L}(\varphi)$ is injective. Next $\operatorname{dim} \mathfrak{s} \mathfrak{o}_{4}(\mathbb{R})=6=2 \operatorname{dim} \mathfrak{s} \mathfrak{u}_{2}(\mathbb{C})$ shows that $\mathbf{L}(\varphi)$ is surjective, and we conclude that

$$
\operatorname{im}(\varphi)=\langle\operatorname{expim} \mathbf{L}(\varphi)\rangle=\mathrm{SO}_{4}(\mathbb{R}) .
$$

Therefore, $\varphi$ is a covering morphism (Proposition 9.5.1). Since $\mathrm{SU}_{2}(\mathbb{C})$ is simply connected, $\widetilde{\mathrm{SO}}_{4}(\mathbb{R}) \cong \mathrm{SU}_{2}(\mathbb{C})^{2}$.

Let $G:=\mathrm{SU}_{2}(\mathbb{C})^{2}$. We have just seen that this is the universal covering group of $\mathrm{SO}_{4}(\mathbb{R})$. On the other hand, $\mathrm{SU}_{2}(\mathbb{C}) \cong \widetilde{\mathrm{SO}}_{3}(\mathbb{R})$. From $Z\left(\mathrm{SU}_{2}(\mathbb{C})\right)= \{ \pm \mathbf{1}\}$, we derive that

$$
Z(G)=\{(\mathbf{1}, \mathbf{1}),(\mathbf{1},-\mathbf{1}),(-\mathbf{1}, \mathbf{1}),(-\mathbf{1},-\mathbf{1})\} \cong C_{2}^{2}
$$

We have

$$
G / Z(G) \cong \mathrm{SO}_{3}(\mathbb{R}) \times \mathrm{SO}_{3}(\mathbb{R})
$$

and therefore,

$$
\mathrm{SO}_{4}(\mathbb{R}) /\{ \pm \mathbf{1}\} \cong G / Z(G) \cong \mathrm{SO}_{3}(\mathbb{R}) \times \mathrm{SO}_{3}(\mathbb{R})
$$

The group $\mathrm{SO}_{4}(\mathbb{R})$ is a twofold covering group of $\mathrm{SO}_{3}(\mathbb{R})^{2}$.

\subsubsection*{Exercises for Section 9.5}
Exercise 9.5.1. Let $G$ be a connected linear Lie group. Show that

$$
\mathbf{L}(Z(G))=\mathfrak{z}(\mathbf{L}(G)):=\{x \in \mathbf{L}(G):(\forall y \in \mathbf{L}(G))[x, y]=0\} .
$$

Exercise 9.5.2. Let $q_{G}: \widetilde{G} \rightarrow G$ be a simply connected covering of the connected Lie group $G$.\\
(1) Show that each automorphism $\varphi \in \operatorname{Aut}(G)$ has a unique lift $\widetilde{\varphi} \in \operatorname{Aut}(\widetilde{G})$.\\
(2) Derive from (1) that $\operatorname{Aut}(G) \cong\left\{\widetilde{\varphi} \in \operatorname{Aut}(\widetilde{G}): \widetilde{\varphi}\left(\pi_{1}(G)\right)=\pi_{1}(G)\right\}$.\\
(3) Show that, in general, $\left\{\widetilde{\varphi} \in \operatorname{Aut}(\widetilde{G}): \widetilde{\varphi}\left(\pi_{1}(G)\right) \subseteq \pi_{1}(G)\right\}$ is not a subgroup of $\operatorname{Aut}(\widetilde{G})$.

\subsection*{Arcwise Connected Subgroups and Initial Subgroups}
The central result of this section is Yamabe's theorem characterizing integral subgroups of a Lie group by the purely topological property of being arcwise connected. Groups which are not arcwise connected can then be dealt with using the concept of initial subgroups.

\subsubsection*{Yamabe's Theorem}
Theorem 9.6.1 (Yamabe). A subgroup $H$ of a connected Lie group $G$ is arcwise connected if and only if it is an integral subgroup. More precisely, $H$\\
is of the form $\left\langle\exp _{G} \mathfrak{h}\right\rangle$ for the Lie subalgebra $\mathfrak{h}$ of $\mathbf{L}(G)$, determined by

$$
\mathfrak{h}=\left\{x \in \mathbf{L}(G): \exp _{G}(\mathbb{R} x) \subseteq H\right\}
$$

Clearly, each integral subgroup is arcwise connected, so that the main point of the proof is the converse. For that it is decisive to find the Lie algebra $\mathfrak{h}$ associated with an arcwise connected subgroup $H$. This is our first target.

Definition 9.6.2. Let $H \subseteq G$ be a subgroup and $x \in \mathbf{L}(G)$. We call $x H$-attainable if for each 1-neighborhood $U$ in $G$ there exists a continuous path $\gamma:[0,1] \rightarrow H$ with $\gamma(0)=\mathbf{1}$ and

$$
\gamma(t) \in \exp _{G}(t x) U
$$

for $t \in[0,1]$. We write $A(H) \subseteq \mathbf{L}(G)$ for the set of $H$-attainable elements of $\mathbf{L}(G)$.

Remark 9.6.3. For each $x \in A(H)$, each neighborhood of $\exp _{G}(x)$ is of the form $\exp _{G}(x) U, U$ an identity neighborhood of $G$, and all these sets contain elements of $H$. This implies that $\exp _{G}(x) \in \bar{H}$.

Lemma 9.6.4. Let ( $X, d$ ) be a compact metric space, $G$ a topological group and $f: X \rightarrow G$ a continuous map. Then the following assertions hold:\\
(a) $f$ is uniformly continuous in the sense that for each $\mathbf{1}$-neighborhood $U \subseteq G$ there exists a $\delta>0$ with

$$
f(y) \in f(x) U \quad \text { for } \quad d(x, y)<\delta
$$

(b) If $\left(\gamma_{n}\right)$ is a sequence of paths $\gamma_{n}:[0,1] \rightarrow X$ converging uniformly to $\gamma:[0,1] \rightarrow X$, then for each $\mathbf{1}$-neighborhood $U$ in $G$, the relation

$$
f\left(\gamma_{n}(t)\right) \in f(\gamma(t)) U
$$

holds for all $t \in[0,1]$ if $n$ is sufficiently large.\\
(c) Let $K \subseteq G$ be a compact subset and ( $X, d$ ) a metric space. Then each continuous function $f: K \rightarrow X$ is uniformly continuous in the sense that for each $\varepsilon>0$, there exists a $\mathbf{1}$-neighborhood $U$ in $G$ such that

$$
d(f(x), f(y))<\varepsilon \quad \text { holds for } \quad x, y \in K, y \in x U
$$

Proof. (a) First, we pick a $\mathbf{1}$-neighborhood $V$ in $G$ with $V^{-1} V \subseteq U$. Since $f$ is continuous, there exists for each $x \in X$ a $\delta_{x}$ such that $f(y) \in f(x) V$ whenever $d(x, y)<\delta_{x}$. Now the open balls $U_{x}:=\left\{y \in X: 2 d(x, y)<\delta_{x}\right\}$ form an open cover of $X$, and the compactness implies the existence of a finite subcover. Hence there exist $x_{1}, \ldots, x_{n} \in X$ with $X \subseteq \bigcup_{i=1}^{n} U_{x_{i}}$. Let $\delta:=\frac{1}{2} \min \left\{\delta_{x_{i}}: i=1, \ldots, n\right\}$ and $x, y \in X$ with $d(x, y)<\delta$. Then there exists an $x_{i}$ with $2 d\left(x, x_{i}\right)<\delta_{x_{i}}$, and then we also have

$$
d\left(y, x_{i}\right) \leq d(y, x)+d\left(x, x_{i}\right)<\delta+\frac{1}{2} \delta_{x_{i}} \leq \frac{1}{2} \delta_{x_{i}}+\frac{1}{2} \delta_{x_{i}}=\delta_{x_{i}} .
$$

Hence $f(x), f(y) \in f\left(x_{i}\right) V$, and thus $f(x)^{-1} f(y) \in V^{-1} V \subseteq U$.\\
(b) is an immediate consequence of (a).\\
(c) Since $f$ is continuous, we find for each $x \in K$ a $\mathbf{1}$-neighborhood $U_{x}$ in $G$ with

$$
f\left(K \cap x U_{x}\right) \subseteq U_{\frac{\varepsilon}{2}}(f(x)):=\left\{z \in X: d(z, f(x))<\frac{\varepsilon}{2}\right\} .
$$

We pick for each $x \in K$ another 1-neighborhood $V_{x}$ in $G$ with $V_{x} V_{x} \subseteq U_{x}$. Since $K$ is compact, it is covered by finitely many $x_{i} V_{x_{i}}, i=1, \ldots, n$. Let $U:=\bigcap_{i=1}^{n} V_{x_{i}}$.

Now let $x, y \in K$ with $y \in x U$. Then $x \in x_{i} V_{x_{i}}$ for some $i$, and therefore $y \in x_{i} V_{x_{i}} U \subseteq x_{i} V_{x_{i}} V_{x_{i}} \subseteq x_{i} U_{x_{i}}$. Therefore, $f(x), f(y) \in U_{\frac{\varepsilon}{2}}\left(f\left(x_{i}\right)\right)$, and thus $d(f(x), f(y))<\varepsilon$. \(\square\)

Remark 9.6.5. (a) Let $\gamma:[0,1] \rightarrow G$ be a $C^{1}$-curve with $\gamma(0)=\mathbf{1}$ and $x:=\gamma^{\prime}(0)$.

Further, let $U_{\mathfrak{g}} \subseteq \mathfrak{g}=\mathbf{L}(G)$ be an open 0-neighborhood for which $\left.\exp _{G}\right|_{U_{\mathfrak{g}}}$ is a diffeomorphism onto the open 1-neighborhood $U_{G}=\exp _{G}\left(U_{\mathfrak{g}}\right)$ in $G$. We assume that $\operatorname{im}(\gamma) \subseteq U_{G}$, so that

$$
\beta:=\left(\left.\exp _{G}\right|_{U_{\mathfrak{g}}}\right)^{-1} \circ \gamma:[0,1] \rightarrow \mathfrak{g}
$$

is a $C^{1}$-curve starting at 0 . We consider the sequence of curves

$$
\beta_{n}(t):=n \beta\left(\frac{t}{n}\right)=\int_{0}^{t} \beta^{\prime}\left(\frac{s}{n}\right) d s, \quad 0 \leq t \leq 1,
$$

which converges uniformly to the curve $\beta(t):=t x$. Since $\exp _{G}$ is uniformly continuous on each compact neighborhood of $\operatorname{im}(\beta)$, Lemma 9.6.4 now implies that the sequence of curves

$$
\gamma_{n}(t):=\left(\exp _{G} \circ \beta_{n}\right)(t)=\gamma\left(\frac{t}{n}\right)^{n}
$$

has the property that for each $\mathbf{1}$-neighborhood $U$ at $\mathbf{1}$, the relation

$$
\gamma_{n}(t) \in \exp _{G}(t x) U
$$

holds for all $t \in[0,1]$ whenever $n$ is sufficiently large.\\
(b) If $H$ is a subgroup of $G$ and $\operatorname{im}(\gamma) \subseteq H$, then $\operatorname{im}\left(\gamma_{n}\right) \subseteq H$ for each $n$, so that (a) implies that $x=\gamma^{\prime}(0) \in A(H)$.

Lemma 9.6.6. $A(H)$ is a Lie subalgebra of $\mathbf{L}(G)$.

Proof. We split the proof into several steps.\\
Step 1: $x \in A(H) \Rightarrow-x \in A(H)$ :\\
Let $U$ be a $\mathbf{1}$-neighborhood in $G$. Since $[0,1]$ is compact and the conjugation action of $G$ on itself is continuous, we find with Exercise 9.6.2 a 1-neighborhood $V$ of $\mathbf{1}$ in $G$ with $V=V^{-1}$ and

$$
\exp _{G}(t x) V \exp _{G}(-t x) \subseteq U \quad \text { for } \quad t \in[0,1]
$$

Since $x \in A(H)$, there exists a continuous path $\gamma:[0,1] \rightarrow H$ starting at $\mathbf{1}$ with $\gamma(t) \in \exp _{G}(t x) V$ for $t \in[0,1]$. We put $\widetilde{\gamma}(t)=\gamma(t)^{-1}$. Then we find with (9.16)

$$
\widetilde{\gamma}(t) \in V^{-1} \exp _{G}(-t x) \subseteq \exp _{G}(-t x) U
$$

Now $\widetilde{\gamma}(t) \in H$ implies $-x \in A(H)$.\\
Step 2: $\mathbb{R} x \subseteq A(H)$ for each $x \in A(H)$ :\\
It is immediately clear from the definition that $s x \in A(H)$ for each $s \in [0,1]$. In view of Step 1, we further get $[-1,1] x \subseteq A(H)$ for $x \in A(H)$. Hence it suffices to show that $k x \in A(H)$ holds for $k \in \mathbb{N}$ and $x \in A(H)$. Let $U \subseteq G$ be a $\mathbf{1}$-neighborhood. As above, we find a $\mathbf{1}$-neighborhood $V$ in $G$ with

$$
\prod_{j=1}^{k+1}\left(\exp _{G}\left(t_{j} x\right) V\right) \subseteq \exp _{G}\left(\sum_{j=1}^{k+1} t_{j} x\right) U \quad \text { for } \quad t_{j} \in[0,1]
$$

To see that such a $V$ exists, we apply Exercise 9.6.2 to the continuous map

$$
\begin{aligned}
\mathbb{R}^{k+1} \times G^{k+1} & \rightarrow G \\
\left(t_{1}, \ldots, t_{k+1}, g_{1}, \ldots, g_{k+1}\right) & \mapsto \exp _{G}\left(\sum_{j=1}^{k+1} t_{j} x\right)^{-1} \prod_{j=1}^{k+1}\left(\exp _{G}\left(t_{j} x\right) g_{j}\right),
\end{aligned}
$$

mapping $[0,1]^{k+1} \times\{\boldsymbol{1}\}^{k+1}$ to $\boldsymbol{1}$.\\
Now let $\gamma:[0,1] \rightarrow H$ be a continuous path starting at $\mathbf{1}$ with $\gamma(t) \in \exp _{G}(t x) V$ for $t \in[0,1]$. For $t \in[0,1]$, we put

$$
\widetilde{\gamma}(t)=\gamma(1)^{[k t]} \gamma(k t-[k t]),
$$

where $[k t]=\max \{b \in \mathbb{N}: b \leq k t\}$. Then $\widetilde{\gamma}$ is continuous, and (9.17) leads to

$$
\widetilde{\gamma}(t) \in\left(\exp _{G}(x) V\right)^{[k t]} \exp _{G}((k t-[k t]) x) V \subseteq \exp _{G}(t k x) U
$$

for $t \in[0,1]$, and hence to $k x \in A(H)$.\\
Step 3: $x+y \in A(H)$ for $x, y \in A(H)$ :\\
Let $\beta(t):=\exp _{G}(t x) \exp _{G}(t y)$. Then $\beta:[0,1] \rightarrow G$ is a smooth curve with $\beta(0)=1$ and $\beta^{\prime}(0)=x+y$. In view of Remark 9.6.5, there exists for each 1-neighborhood $V$ in $G$ an $N \in \mathbb{N}$ such that

$$
\beta\left(\frac{t}{k}\right)^{k}=\left(\left(\exp _{G} \frac{t x}{k}\right)\left(\exp _{G} \frac{t y}{k}\right)\right)^{k} \in \exp _{G}(t(x+y)) V
$$

for all $t \in[0,1]$ and $k>N$. For any such $k$, there exists a 1-neighborhood $W$ in $G$ with

$$
\left(\left(\exp _{G} \frac{t x}{k}\right) W\left(\exp _{G} \frac{t y}{k}\right) W\right)^{k} \subseteq\left(\left(\exp _{G} \frac{t x}{k}\right)\left(\exp _{G} \frac{t y}{k}\right)\right)^{k} V
$$

for all $t \in[0,1]$.\\
Now let $U$ be a $\mathbf{1}$-neighborhood in $G$, choose $V$ such that $V V \subseteq U, W$ as above. Further, let $\gamma_{x}, \gamma_{y}:[0,1] \rightarrow H$ be continuous curves starting at $\mathbf{1}$ with $\gamma_{x}(t) \in\left(\exp _{G} \frac{t x}{k}\right) W$ and $\gamma_{y}(t) \in\left(\exp _{G} \frac{t y}{k}\right) W$ for $t \in[0,1]$. The existence of such paths is due to $\frac{x}{k}, \frac{y}{k} \in A(H)$. Put

$$
\widetilde{\gamma}(t):=\left(\gamma_{x}(t) \gamma_{y}(t)\right)^{k},
$$

so that we obtain with (9.18) and (9.19)

$$
\begin{aligned}
\widetilde{\gamma}(t) & \in\left(\left(\exp _{G} \frac{t x}{k}\right) W\left(\exp _{G} \frac{t y}{k}\right) W\right)^{k} \subseteq\left(\left(\exp _{G} \frac{t x}{k}\right)\left(\exp _{G} \frac{t y}{k}\right)\right)^{k} V \\
& \subseteq \exp _{G}(t(x+y)) V V \subseteq \exp _{G}(t(x+y)) U
\end{aligned}
$$

This implies that $x+y \in A(H)$.\\
Step 4: $\operatorname{Ad}(h) x \in A(H)$ for $h \in H$ and $x \in A(H)$ :\\
Let $h \in H$ and $x \in A(H)$. For an identity neighborhood $U$ in $G$, we then find an identity neighborhood $U_{h}$ of $G$ with $c_{h}\left(U_{h}\right) \subseteq U$. If $\gamma_{x}:[0,1] \rightarrow H$ satisfies $\gamma_{x}(t) \in \exp _{G}(t x) U$ for each $t \in[0,1]$, we then obtain

$$
h \gamma_{x}(t) h^{-1} \in c_{h}\left(\exp _{G}(t x)\right) c_{h}(U) \subseteq \exp _{G}(t \operatorname{Ad}(h) x) U
$$

so that $\operatorname{Ad}(h) x \in A(H)$.\\
Step 5: $[x, y] \in A(H)$ for $x, y \in A(H)$ :\\
The normalizer $N:=\{g \in G: \operatorname{Ad}(g) A(H)=A(H)\}$ of the subspace $A(H) \subseteq \mathbf{L}(G)$ in $G$ is a closed subgroup, and we know from Step 4 that it contains $H$. In view of Remark 9.6.3, it also contains $\exp _{G}(A(H))$, so that $e^{t \operatorname{ad} x} y \in A(H)$ for $x, y \in A(H)$ and $t \in \mathbb{R}$. Taking derivatives at $t=0$, we obtain $[x, y] \in A(H)$.

Lemma 9.6.7. If $H$ is an arcwise connected topological subgroup of $G$, then for each identity neighborhood $U$ of $\mathbf{1}$ in $G$, the arc-component $U_{a}$ of $U \cap H$ generates $H$.

Proof. Let $H_{U} \subseteq H$ denote the subgroup generated by $U_{a}$. Further, let $h \in H$ and $\gamma:[0,1] \rightarrow H$ be a continuous path from $\mathbf{1}$ to $h$. We consider the set

$$
S:=\left\{t \in[0,1]: \gamma(t) \in H_{U}\right\}
$$

Then the subset $\gamma^{-1}(U)$ of $S$ is a neighborhood of 0 . For any $t_{0} \in S$, there exists an $\varepsilon>0$ with

$$
\gamma\left(t_{0}\right)^{-1} \gamma(t) \in U_{a}
$$

for $\left|t-t_{0}\right| \leq \varepsilon$. This implies that $S$ is an open subset of $[0,1]$. If $t_{1} \in[0,1] \backslash S$, then we also find an $\varepsilon>0$ with

$$
\gamma\left(t_{1}\right)^{-1} \gamma(t) \in U_{a} \subseteq H_{U}
$$

for $\left|t-t_{1}\right| \leq \varepsilon$, so that $\gamma\left(t_{1}\right) \notin H_{U}$ leads to $\gamma(t) \notin H_{U}$ for $\left|t-t_{1}\right| \leq \varepsilon$. Thus $[0,1] \backslash S$ is also open in [ 0,1 ], and now the connectedness of [ 0,1 ] implies that $S=[0,1]$, hence $h=\gamma(1) \in H_{U}$.

It remains to show that the integral subgroup $\left\langle\exp _{G} A(H)\right\rangle$ coincides with $H$. In the following lemma, we show one inclusion.

Lemma 9.6.8. If $H$ is an arcwise connected subgroup of the Lie group $G$, then $H \subseteq\left\langle\exp _{G}(A(H))\right\rangle$.

Proof. Let $\left(r_{m}\right)_{m \in \mathbb{N}}$ be any sequence of positive real numbers converging to 0 with $r_{m+1}<r_{m}$ for $m \in \mathbb{N}$. We pick a norm $\|\cdot\|$ on $\mathfrak{g}$ and put

$$
B_{m}:=\left\{x \in \mathfrak{g}:\|x\|<r_{m}\right\} \quad \text { and } \quad U_{m}:=\exp _{G}\left(B_{m}\right)
$$

Let $\mathfrak{p} \subseteq \mathfrak{g}=\mathbf{L}(G)$ be a vector space complement of the Lie subalgebra $A(H)$. Then there exist open convex symmetric 0-neighborhoods $V_{m} \subseteq A(H)$ and $W_{m} \subseteq \mathfrak{p}$ such that the map

$$
\psi_{m}: V_{m} \times W_{m} \rightarrow U_{m}, \quad(x, y) \mapsto \exp _{G}(x) \exp _{G}(y)
$$

is a diffeomorphism onto an open subset $U_{m}^{\prime}:=\exp _{G}\left(V_{m}\right) \exp _{G}\left(W_{m}\right)$ of $U_{m}$ (cf. Lemma 9.3.6). An easy induction shows that we may choose the sets $V_{m}$ and $W_{m}$ in such a way that $V_{m+1} \subseteq V_{m}$ and $W_{m+1} \subseteq W_{m}$ for each $m \in \mathbb{N}$.

Let $H_{m}$ denote the arc-component of $H \cap U_{m}^{\prime}$ containing 1 . Then $H_{m}$ generates $H$ by Lemma 9.6.7. Therefore, it remains to show that $H_{m} \subseteq \exp _{G}(A(H))$ for some $m \in \mathbb{N}$. If this is not the case, then there exists a sequence

$$
h_{m}=\exp _{G}\left(v_{m}\right) \exp _{G}\left(w_{m}\right) \in H_{m} \quad \text { with } \quad v_{m} \in V_{m}, 0 \neq w_{m} \in W_{m} .
$$

Then the sequence $\widetilde{w}_{m}:=\frac{w_{m}}{\left\|w_{m}\right\|}$ has a cluster point, and by passing to a suitable subsequence of $\left(r_{m}\right)_{m \in \mathbb{N}}$, we may assume that the sequence $\left(\widetilde{w}_{m}\right)_{m \in \mathbb{N}}$ converges to some $w \in \mathfrak{p}$ with $\|w\|=1$.

To arrive at a contradiction, we claim that $w \in A(H)$. To this end, we fix some $m \in \mathbb{N}$ and let $U^{(m)}$ be a $\mathbf{1}$-neighborhood in $G$. Then there exists a smaller 1-neighborhood $U_{m}^{\prime \prime} \subseteq U^{(m)} \cap U_{m}$ with

$$
h_{m}^{-1} U_{m}^{\prime \prime} h_{m} \subseteq U^{(m)}
$$

Further, $-v_{m} \in A(H)$ implies the existence of a continuous path $\gamma_{m}:[0,1] \rightarrow H$ starting at $\mathbf{1}$ with $\gamma_{m}(t) \in \exp _{G}\left(-t v_{m}\right) U_{m}^{\prime \prime}$ for $t \in[0,1]$. Since $H_{m}$ is\\
arcwise connected, there also exists a continuous path $\eta_{m}:[0,1] \rightarrow H_{m}$ from 1 to $h_{m}$. Now

$$
\widetilde{\gamma}(t):=\gamma_{m}(t) \eta_{m}(t)
$$

satisfies $\widetilde{\gamma}_{m}(0)=\mathbf{1}$ and

$$
\widetilde{\gamma}_{m}(1) \in \exp _{G}\left(-v_{m}\right) U_{m}^{\prime \prime} h_{m} \subseteq \exp _{G}\left(-v_{m}\right) h_{m} U^{(m)}=\left(\exp _{G} w_{m}\right) U^{(m)} .
$$

Moreover,

$$
\widetilde{\gamma}_{m}(t) \in\left(\exp _{G} V_{m}\right) U_{m}^{\prime \prime} H_{m} \subseteq U_{m}^{3} .
$$

If $m$ is sufficiently large, then we have a smooth inverse $\log : U_{m}^{\prime} \rightarrow \mathfrak{g}$ of the exponential function, so that we may put

$$
z_{m}:=\log \left(\widetilde{\gamma}_{m}(1)\right) .
$$

Then (9.21) yields

$$
z_{m}=w_{m} * w \quad \text { with some } \quad w \in V^{(m)}:=\log \left(U^{(m)}\right) .
$$

Next we choose $U^{(m)}=\exp _{G}\left(V^{(m)}\right)$ so small that

$$
\left\|w_{m} * w-w_{m}\right\| \leq\left\|w_{m}\right\|^{2} \leq r_{m} \quad \text { for all } \quad w \in V^{(m)}
$$

and

$$
w_{m} * V^{(m)} \subseteq B_{m} .
$$

Then

$$
\lim _{m \rightarrow \infty} \frac{z_{m}}{\left\|w_{m}\right\|}=\lim _{m \rightarrow \infty}\left(\frac{z_{m}-w_{m}}{\left\|w_{m}\right\|}+\frac{w_{m}}{\left\|w_{m}\right\|}\right)=w .
$$

For $p_{m}:=\left\|w_{m}\right\|^{-1}$ and

$$
\widehat{\gamma}_{m}(t):=\widetilde{\gamma}_{m}(1)^{\left[t p_{m}\right]} \widetilde{\gamma}_{m}\left(t p_{m}-\left[t p_{m}\right]\right),
$$

we also have

$$
\begin{aligned}
\widehat{\gamma}_{m}(t) & =\exp _{G}\left(\left[t p_{m}\right] z_{m}\right) \widetilde{\gamma}_{m}\left(t p_{m}-\left[t p_{m}\right]\right) \\
& =\exp _{G}\left(t p_{m} z_{m}\right) \exp _{G}\left(\left(\left[t p_{m}\right]-t p_{m}\right) z_{m}\right) \widetilde{\gamma}_{m}\left(t p_{m}-\left[t p_{m}\right]\right) \\
& \in \exp _{G}\left(t p_{m} z_{m}\right) U_{m} U_{m}^{3} \subseteq \exp _{G}\left(t p_{m} z_{m}\right) U_{m}^{4}
\end{aligned}
$$

Finally, let $U$ be any $\mathbf{1}$-neighborhood in $G$. Then there exists a $k \in \mathbb{N}$ with $U_{k}^{5} \subseteq U$, and the sequence $\exp _{G}\left(t p_{m} z_{m}\right)$ converges for $m \rightarrow \infty$ uniformly in $t \in[0,1]$ to $\exp _{G}(t w)$ (Lemma 9.6.4(b)). Hence there exists an $m>k$ with

$$
\exp _{G}\left(t p_{m} z_{m}\right) \subseteq\left(\exp _{G} t w\right) U_{k} \quad \text { for all } t \in[0,1]
$$

Then we finally arrive at

$$
\widehat{\gamma}_{m}(t) \in \exp _{G}\left(t p_{m} z_{m}\right) U_{m}^{4} \subseteq\left(\exp _{G} t w\right) U_{k} U_{m}^{4} \subseteq\left(\exp _{G} t w\right) U_{k}^{5} \subseteq \exp _{G}(t w) U,
$$

and thus $w \in A(H)$; contradicting $w \in \mathfrak{p} \backslash\{0\}$. \(\square\)

To prove also the converse of Lemma 9.6.8, we need a corollary to Brouwer's Fixed Point Theorem, saying that each continuous selfmap of the closed unit ball in $\mathbb{R}^{m}$ (with respect to any norm) has a fixed point (cf. [Hir76, p. 73]).

Lemma 9.6.9. Let $\|x\|:=\sqrt{\sum_{i} x_{i}^{2}}$ denote the euclidian norm on $\mathbb{R}^{n}$. If $f:[-1,1]^{m} \rightarrow \mathbb{R}^{m}$ is a continuous map with

$$
\|f(x)-x\| \leq \frac{1}{2} \quad \text { for } \quad x \in[-1,1]^{m}
$$

then $\left\{x \in \mathbb{R}^{m}:\|x\|<\frac{1}{2}\right\} \subseteq f\left([-1,1]^{m}\right)$.\\
Proof. Let $y \in \mathbb{R}^{m}$ with $\|y\|_{\infty} \leq\|y\|<\frac{1}{2}$ and put $g(x):=x-f(x)+y$. Then

$$
g:[-1,1]^{m} \rightarrow[-1,1]^{m}
$$

is a continuous map and Brouwer's Fixed Point Theorem implies that $g$ has some fixed point $x_{o}$. Then $f\left(x_{o}\right)=x_{o}-g\left(x_{o}\right)+y=y$.

Lemma 9.6.10. If $H$ is an arcwise connected subgroup of the Lie group $G$, then $\exp _{G}(A(H)) \subseteq H$.

Proof. It suffices to show that $H$ contains an open 1-neighborhood in $H^{\#}:= \left\langle\exp _{G}(A(H))\right\rangle$ with respect to the intrinsic Lie group topology of this group (Integral Subgroup Theorem 9.4.8). We choose a basis $x_{1}, \ldots, x_{m}$ of $A(H)$, for which the map

$$
\varphi:]-2,2\left[{ }^{m} \rightarrow U, \quad t=\left(t_{1}, \ldots, t_{m}\right) \mapsto \prod_{j=1}^{m} \exp _{G}\left(t_{j} x_{j}\right)\right.
$$

is a homeomorphism onto an open subset $U$ of $H^{\#}$. The existence of such a basis follows from the Inverse Function Theorem (cf. Lemma 9.3.6). In view of the compactness of $[-1,1]^{m}$, there exists a compact $\mathbf{1}$-neighborhood $U_{1}$ in $H^{\sharp}$ with\\
(a) $\varphi\left([-1,1]^{m}\right) U_{1} \subseteq U$,\\
(b) $\|s-t\| \leq \frac{1}{2}$ if $\left.s \in\right]-2,2\left[^{m}, t \in[-1,1]^{m}\right.$ and $\varphi(s) \in \varphi(t) U_{1}$ (Lemma 9.6.4(c)).\\
We further find a 1-neighborhood $U_{2}$ in $H^{\#}$ with

$$
\prod_{j=1}^{m}\left(\exp _{G}\left(t_{j} x_{j}\right) U_{2}\right) \subseteq \varphi\left(t_{1}, \ldots, t_{m}\right) U_{1} \quad \text { for } \quad t_{1}, \ldots, t_{m} \in[-1,1]
$$

Since $\pm x_{j} \in A(H)$, there exist continuous curves $\gamma_{j}:[-1,1] \rightarrow H$ with $\gamma_{j}(0)=\mathbf{1}$ and $\gamma_{j}(t) \in\left(\exp _{G}\left(t x_{j}\right)\right) U_{2}$ for all $t$. We then have

$$
\prod_{j=1}^{m} \gamma_{j}\left(t_{j}\right) \in \prod_{j=1}^{m}\left(\exp _{G}\left(t_{j} x_{j}\right) U_{2}\right) \subseteq \varphi\left(t_{1}, \ldots, t_{m}\right) U_{1} \subseteq U .
$$

Using $\varphi$, we may thus write

$$
\prod_{j=1}^{m} \gamma_{j}\left(t_{j}\right)=\prod_{j=1}^{m} \exp _{G}\left(f_{j}\left(t_{1}, \ldots, t_{m}\right) x_{j}\right),
$$

and obtain a continuous function

$$
f:[-1,1]^{m} \rightarrow \mathbb{R}^{m}, \quad\left(t_{1}, \ldots, t_{m}\right) \mapsto\left(f_{1}\left(t_{1}, \ldots, t_{m}\right), \ldots, f_{m}\left(t_{1}, \ldots, t_{m}\right)\right) .
$$

Note that

$$
\varphi\left(f_{1}\left(t_{1}, \ldots, t_{m}\right), \ldots, f_{m}\left(t_{1}, \ldots, t_{m}\right)\right)=\prod_{j=1}^{m} \gamma_{j}\left(t_{j}\right) \in \varphi\left(t_{1}, \ldots, t_{m}\right) U_{1} \subseteq U,
$$

so that, in view of (b), $\|f(t)-t\| \leq \frac{1}{2}$. Finally, we apply Lemma 9.6.9 and see that $f\left([-1,1]^{m}\right)$ contains a neighborhood of 0 in $]-2,2\left[^{m}\right.$. Therefore,

$$
\gamma_{1}([-1,1]) \cdots \gamma_{m}([-1,1])
$$

contains a $\mathbf{1}$-neighborhood of $H^{\#}$, and this completes the proof.\\
Proof. (of Theorem 9.6.1) Combining Lemmas 9.6.8 and 9.6.10, we see that $H=\left\langle\exp _{G}(A(H))\right\rangle$ is the integral subgroup of $G$ corresponding to the Lie subalgebra $A(H)$ of $\mathbf{L}(G)$. The definition of $A(H)$ implies in particular that

$$
\left\{x \in \mathbf{L}(G): \exp _{G}(\mathbb{R} x) \subseteq H\right\} \subseteq A(H),
$$

which leads to the equality

$$
A(H)=\left\{x \in \mathbf{L}(G): \exp _{G}(\mathbb{R} x) \subseteq H\right\} .
$$

\subsubsection*{Initial Lie Subgroups}
Definition 9.6.11. An injective morphism $\iota: H \rightarrow G$ of Lie groups is called an initial Lie subgroup if $\mathbf{L}(\iota): \mathbf{L}(H) \rightarrow \mathbf{L}(G)$ is injective, and for each smooth map $f: M \rightarrow G$ from a smooth manifold $M$ to $G$ with $\operatorname{im}(f) \subseteq H$, the corresponding map $\iota^{-1} \circ f: M \rightarrow H$ is smooth.

The following lemma shows that the existence of an initial Lie group structure only depends on the subgroup $H$, considered as a subset of $G$.

Lemma 9.6.12. Any subgroup $H$ of a Lie group $G$ carries at most one structure of an initial Lie subgroup.

Proof. If $\iota^{\prime}: H^{\prime} \hookrightarrow G$ is another initial Lie subgroup with the same range as $\iota: H \rightarrow G$, then $\iota^{-1} \circ \iota^{\prime}: H^{\prime} \rightarrow H$ and $\iota^{\prime-1} \circ \iota: H \rightarrow H^{\prime}$ are smooth morphisms of Lie groups, so that $H$ and $H^{\prime}$ are isomorphic.

Theorem 9.6.13 (Initial Subgroup Theorem). Each subgroup $H$ of a Lie group $G$ carries an initial Lie subgroup structure for which the identity component $H_{0}$ coincides with the arc-component of $H$ with respect to the subspace topology and

$$
\mathbf{L}(H) \cong\left\{x \in \mathbf{L}(G): \exp _{G}(\mathbb{R} x) \subseteq H\right\}
$$

Proof. Let $H_{a} \subseteq H$ be the arc-component of $H$, viewed as a topological subgroup of $G$. According to Yamabe's Theorem 9.6.1, $H_{a}$ is an integral subgroup with Lie algebra

$$
\mathfrak{h}=\left\{x \in \mathbf{L}(G): \exp _{G}(\mathbb{R} x) \subseteq H_{a}\right\}=\left\{x \in \mathbf{L}(G): \exp _{G}(\mathbb{R} x) \subseteq H\right\}
$$

Let $H_{a}^{L}$ denote the group $H_{a}$, endowed with its intrinsic Lie group topology for which $\exp =\left.\exp _{G}\right|_{\mathfrak{h}}: \mathfrak{h} \rightarrow H_{a}^{L}$ is a local diffeomorphism in 0 . Then $H \subseteq\{g \in G: \operatorname{Ad}(g) \mathfrak{h}=\mathfrak{h}\}$ and $c_{g} \circ \exp =\left.\exp \circ \operatorname{Ad}(g)\right|_{\mathfrak{h}}$ imply that for each $h \in H$, the conjugation $c_{h}$ defines a smooth automorphisms on $H_{a}^{L}$, so that $H$ carries a Lie group structure for which $H_{a}^{L}$ is an open subgroup (Corollary 9.4.5). Let $H^{L}$ denote this Lie group. Now the inclusion map $\iota: H^{L} \rightarrow G$ is an immersion whose differential at $\mathbf{1}$ is the inclusion $\mathfrak{h} \hookrightarrow \mathfrak{g}$.

We claim that $\iota: H^{L} \rightarrow G$ is an initial Lie subgroup. In fact, let $f: M \rightarrow G$ be a smooth map from the smooth manifold $M$ to $G$ with $f(M) \subseteq H$. We have to show that $f$ is smooth, i.e., that $f$ is smooth in a neighborhood of each point $m \in M$. Replacing $f$ by $f \cdot f(m)^{-1}$ and observing that the group operations in $H$ and $G$ are smooth, we may w.lo.g. assume that $f(m)=\mathbf{1}$.

Let $U \subseteq \mathfrak{h}$ be an open 0-neighborhood, $\mathfrak{m} \subseteq \mathfrak{g}$ be a vector space complement to $\mathfrak{h}$, and $V \subseteq \mathfrak{m}$ an open 0-neighborhood for which the map

$$
\Phi: U \times V \rightarrow G, \quad(x, y) \mapsto \exp _{G} x \exp _{G} y
$$

is a diffeomorphism onto an open subset of $G$. Then

$$
H_{a} \cap\left(\exp _{G} U \exp _{G} V\right)=\bigcup_{y \in V, \exp _{G}} y \in H_{a} \exp _{G} U \exp _{G} y
$$

where the union on the right is disjoint because $\varphi$ is bijective, and each set $\exp _{G} U \exp _{G} y$ contained in $H_{a}$ also is an open subset of $H^{L}$. Since the topology of $H_{a}^{L}$ is second countable (Proposition 9.1.15(iv)), the set $\exp _{G}^{-1}\left(H_{a}\right) \cap V$ is countable. Every smooth arc $\gamma: I \rightarrow \exp _{G} U \exp _{G} V$ is of the form $\gamma(t)=\exp _{G} \alpha(t) \exp _{G} \beta(t)$ with smooth $\operatorname{arcs} \alpha: I \rightarrow U$ and $\beta: I \rightarrow V$, and for every smooth arc contained in $H_{a}$, the countability of $\beta(I)$ implies that $\beta$ is constant.

We conclude that if $W \subseteq M$ is an open connected neighborhood of $m$ with $f(W) \subseteq \exp _{G} U \exp _{G} V$, then $f(W) \subseteq \exp _{G} U$. Then the map

$$
\left.\left.\exp _{G}\right|_{U} ^{-1} \circ f\right|_{W}: W \rightarrow \mathfrak{h}
$$

is smooth, so that the corresponding map

$$
\left.\iota^{-1} \circ f\right|_{W}=\left.\left.\exp _{H^{L}} \circ \exp _{G}\right|_{U} ^{-1} \circ f\right|_{W}: W \rightarrow H^{L}
$$

is also smooth. This proves that the map $\iota^{-1} \circ f: M \rightarrow H^{L}$ is smooth, and hence that $\iota: H^{L} \rightarrow G$ is an initial Lie subgroup of $G$. \(\square\)

\subsubsection*{Exercises for Section 9.6}
Exercise 9.6.1. Let $X$ and $Y$ be topological spaces and $K_{X} \subseteq X$ and $K_{Y} \subseteq Y$ compact subsets. If $V \subseteq X \times Y$ is an open subset containing $K_{X} \times K_{Y}$, then there exist open subsets $U_{X} \subseteq X$ and $U_{Y} \subseteq Y$ with $U_{X} \times U_{Y} \subseteq U$.

Exercise 9.6.2. Let $X, Y$, and $Z$ be topological spaces and $f: X \times Y \rightarrow Z$ a continuous map. If $K_{X} \subseteq X$ and $K_{Y} \subseteq Y$ are compact and $U \subseteq Z$ open with $f\left(K_{X} \times K_{Y}\right) \subseteq U$, then there exist open subsets $U_{X} \subseteq X$ and $U_{Y} \subseteq Y$ with $f\left(U_{X} \times U_{Y}\right) \subseteq U$.

Exercise 9.6.3. Construct a subgroup of $\mathbb{T}^{2}=\mathbb{R}^{2} / \mathbb{Z}^{2}$ which is connected but not arcwise connected.

\subsubsection*{Notes on Chapter 9}
In a way this chapter describes the essence of Lie theory, namely the translation process between global and infinitesimal objects. A subtle point in this matter is the correspondence of subobjects. Given a Lie group $G$ with Lie algebra $\mathbf{L}(G)$, Lie subalgebras generate subgroups of $G$, via integral manifolds or the exponential function, but unless they are closed, these groups are Lie groups only if one changes the topology. On the other hand, the Initial Subgroup Theorem 9.6.13 shows that any subgroup $H$ can be given a topology such that it becomes a Lie group whose connected component arises in this way. Thus in defining the concept of a Lie subgroup one has to decide whether one wants to include the topology into the structure or not. If not, any subgroup is a Lie subgroup and the concept is superfluous. If one makes the topology part of the structure, only closed subgroups will be Lie subgroups, which explains our Definition 9.3.10. Of course, this does not mean that the nonclosed subgroups associated with Lie subgroups are not important. So they also get names. In the literature, arcwise connected subgroups often go, for instance, under the name analytic subgroups. Singling out this class of subgroups is, of course, motivated by Yamabe's Theorem 9.6.1. We prefer to call them integral subgroups since they systematically arise as integral manifolds.

If one wants to characterize the Lie group structure of general subgroups, in view of Lemma 9.6.12 and the Initial Subgroup Theorem 9.6.13, the initial submanifold property is the crucial one. Our terminology concerning different kinds of subgroups basically follows Bourbaki (cf. [Bou89], Chapter 3).

In the 1890s, Sophus Lie developed his theory of differentiable groups (called continuous groups at the time when the concept of a topological space was not yet developed) to study symmetries of differential equations. The infinitesimal and local theory was worked out and further developed by Friedrich Engel (1861-1941), Wilhelm Killing (1847-1923), and Élie Car$\tan (1869-1951)$. The global theory of Lie groups was initiated by Hermann Weyl (1885-1955) in the seminal series [We25] of papers from 1925-1926, motivated by representation theory and harmonic analysis. It inspired Élie Cartan to reconsider his approach to Lie groups and eventually led to him to deep insights into the topology of Lie groups as well as to the invention of symmetric spaces (see [Ca30, He01]). An important role in the circulation of Lie theoretic ideas was played by Claude Chevalley (1909-1984) through his papers and in particular his book [Ch46] (see also [BG81]). Later developments were often motivated by either representation theory or harmonic analysis, pioneered by Harish-Chandra (1923-1983) and Israil M. Gelfand (1913-2009).

\section*{Smooth Actions of Lie Groups}
In many areas of mathematics, Lie groups appear naturally as symmetry groups. Examples are groups of isometries of Riemannian manifolds, groups of holomorphic automorphisms of complex domains, or groups of canonical transformations in hamiltonian mechanics. In all these cases, one considers group actions on manifolds by smooth maps. Even though the focus of this book is the geometry and structure theory of Lie groups rather than their applications, we have to study the concept of a smooth action of a Lie group on a manifold in some detail since it is an essential tool in the smooth versions of group theoretic considerations like the study of quotient groups and conjugacy classes.

In the first section, we collect some basic facts on orbits and stabilizers for smooth actions. Section 10.1.3 is devoted to the structure of homogeneous spaces. Its results are used to provide orbits of smooth actions with manifold structures. We also introduce frame bundles which turn out to be very useful in the description of bundles with symmetries. With this preparation, it is easy to describe general tensor and density bundles as well as the corresponding sections. Using (invariant) densities and differential forms, one can then quickly describe a basic integration theory for manifolds together with its interaction with symmetries. In particular, we find the Haar measure on a Lie group and its basic properties, which allows us to use averaging techniques in Lie theory.

A smooth action of a Lie group on a manifold gives rise to a representation of the Lie algebra by vector fields on the manifold. In Section 10.5, we give a proof of Palais' Theorem which gives conditions under which, conversely, a finite-dimensional Lie algebra of vector fields can be integrated to a smooth group action.

\subsection*{Homogeneous Spaces}
In this section, we begin our study of the structure of smooth actions of Lie groups on smooth manifolds (cf. Definition 9.1.11) with some elementary observations on orbits and stabilizers.

\subsubsection*{Orbits and Stabilizers}
Definition 10.1.1. Let $\sigma: G \times M \rightarrow M,(g, m) \mapsto g . m$ be a group action. In the following, we write $\sigma_{g}(m):=\sigma(g, m)$ for the corresponding diffeomorphisms of $M$ and $\sigma^{m}(g):=\sigma(g, m)$ for the orbit map of $m$. For $m \in M$, the set

$$
\mathcal{O}_{m}:=G . m:=\{g . m: g \in G\}=\{\sigma(g, m): g \in G\}
$$

is called the orbit of $m$. The action is said to be transitive if there exists only one orbit, i.e., for $x, y \in M$, there exists a $g \in G$ with $y=g . x$. We write $M / G:=\left\{\mathcal{O}_{m}: m \in M\right\}$ for the set of $G$-orbits on $M$.

Remark 10.1.2. If $\sigma: G \times X \rightarrow X$ is an action of $G$ on $X$, then the orbits form a partition of $X$ (Exercise 10.1.1).

A subset $R \subseteq X$ is called a set of representatives for the action if each $G$-orbit in $X$ meets $R$ exactly once:

$$
(\forall x \in X) \quad\left|R \cap \mathcal{O}_{x}\right|=1
$$

Examples 10.1.3. (1) We consider the action of the circle group

$$
\mathbb{T}=\left\{z \in \mathbb{C}^{\times}:|z|=1\right\}
$$

on $\mathbb{C}$ by

$$
\sigma: \mathbb{T} \times \mathbb{C} \rightarrow \mathbb{C}, \quad t . z=t z
$$

The orbits of this action are concentric circles:

$$
\mathcal{O}_{z}=\{t z: t \in \mathbb{T}\}=\{w \in \mathbb{C}:|w|=|z|\}
$$

A set of representatives is given by

$$
R:=[0, \infty[=\{r \in \mathbb{R}: r \geq 0\}
$$

(2) For $\mathbb{K} \in\{\mathbb{R}, \mathbb{C}\}$ and the action

$$
\sigma: \mathrm{GL}_{n}(\mathbb{K}) \times \mathbb{K}^{n} \rightarrow \mathbb{K}^{n}, \quad(g, x) \mapsto g x,
$$

there are only two orbits:

$$
\mathcal{O}_{0}=\{0\} \quad \text { and } \quad \mathcal{O}_{x}=\mathbb{K}^{n} \backslash\{0\} \quad \text { for } \quad x \neq 0
$$

Each non-zero vector $x \in \mathbb{K}^{n}$ can be complemented to a basis for $\mathbb{K}^{n}$, hence arises as a first column of an invertible matrix $g$. Then $g e_{1}=x$ implies that $\mathcal{O}_{x}=\mathcal{O}_{e_{1}}$. We conclude that $\mathbb{K}^{n} \backslash\{0\}=\mathcal{O}_{e_{1}}$.\\
(3) For the conjugation action

$$
\mathrm{GL}_{n}(\mathbb{K}) \times M_{n}(\mathbb{K}) \rightarrow M_{n}(\mathbb{K}), \quad(g, A) \mapsto g A g^{-1}
$$

the orbits are the similarity classes of matrices $\mathcal{O}_{A}=\left\{g A g^{-1}: g \in \mathrm{GL}_{n}(\mathbb{K})\right\}$.

Definition 10.1.4. Let $\sigma: G \times M \rightarrow M,(g, m) \mapsto g . m$ be an action of the group $G$ on $M$. For $m \in M$, the subset

$$
G_{m}:=\{g \in G: g . m=m\}
$$

is called the stabilizer of $m$. For $g \in G$, we write

$$
\operatorname{Fix}(g):=M^{g}:=\{m \in M: g \cdot m=m\}
$$

for the set of fixed points of $g$ in $M$. We then have

$$
m \in M^{g} \Longleftrightarrow g \in G_{m} .
$$

For a subset $S \subseteq M$, we write

$$
G_{S}:=\bigcap_{m \in S} G_{m}=\{g \in G:(\forall m \in S) g \cdot m=m\},
$$

and for $H \subseteq G$ we write

$$
M^{H}:=\{m \in M:(\forall h \in H) h . m=m\}
$$

for the set of points in $M$ fixed by $H$.\\
Lemma 10.1.5. For each smooth action $\sigma: G \times M \rightarrow M$, the following assertions hold:\\
(i) For each $m \in M$, the stabilizer $G_{m}$ of $m$ is a closed subgroup of $G$.\\
(ii) For $m \in M$ and $g \in G$, the following equality holds: $G_{g . m}=g G_{m} g^{-1}$.\\
(iii) If $S \subseteq M$ is a $G$-invariant subset, then $G_{S} \unlhd G$ is a normal subgroup.

Proof. (i) That $G_{m}$ is a subgroup is a trivial consequence of the action axioms. Its closedness follows from the continuity of the orbit map $\sigma^{m}$ and the closedness of the points of $M$.\\
(ii) If $h \in G_{m}$, then

$$
g h g^{-1} \cdot(g \cdot m)=g h g^{-1} g \cdot m=g h \cdot m=g \cdot m,
$$

hence $g h g^{-1} \in G_{g . m}$ and thus $g G_{m} g^{-1} \subseteq G_{g . m}$. Similarly, we get $g^{-1} G_{g . m} g \subseteq G_{g^{-1} .(g . m)}=G_{m}$, and therefore $g G_{m} g^{-1}=G_{g . m}$.\\
(iii) follows directly from (ii).

The normal subgroup $G_{M}$ consisting of all elements of $G$ which do not move any element of $M$ is called the effectivity kernel of the action. It is the kernel of the corresponding homomorphism $G \rightarrow \operatorname{Diff}(M), g \mapsto \sigma_{g}$.

Proposition 10.1.6. Let $\sigma: G \times M \rightarrow M$ be a smooth action of $G$ on $M$. Recall the associated Lie algebra homomorphism $\mathbf{L}(\sigma): \mathbf{L}(G) \rightarrow \mathcal{V}(M)$ from Proposition 9.1.14. Then the following assertions hold:\\
(i) $m \in M^{G} \Rightarrow \mathbf{L}(\sigma)(x)(m)=0$ for each $x \in \mathbf{L}(G)$. The converse holds if $G$ is connected.\\
(ii) If $\mathbf{L}(\sigma)(\mathbf{L}(G))(m)=T_{m}(M)$, then the orbit $\mathcal{O}_{m}$ of $m$ is open.

Proof. (i) Suppose first that $m$ is a fixed point and let $x \in \mathbf{L}(G)$. Then

$$
\mathbf{L}(\sigma)(x)(m)=\left.\frac{d}{d t}\right|_{t=0} \exp _{G}(-t x) \cdot m=\left.\frac{d}{d t}\right|_{t=0} m=0 .
$$

If, conversely, all vector fields $\mathbf{L}(\sigma)(x)$ vanish at $m$, then $m$ is a fixed point for all flows generated by these vector fields, which leads to $\exp _{G}(x) . m=m$ for each $x \in \mathbf{L}(G)$. This means that $G_{m} \supseteq\left\langle\exp _{G} \mathbf{L}(G)\right\rangle$, which in turn is the identity component of $G$ (Lemma 9.2.9). If $G$ is connected, we get $G=G_{m}$.\\
(ii) Since the linear map $T_{\mathbf{1}}\left(\sigma^{m}\right): \mathbf{L}(G) \rightarrow T_{m}(M), x \mapsto \mathbf{L}(\sigma)(x)(m)$ is surjective, the Implicit Function Theorem 8.1.9 implies that $G . m=\sigma^{m}(G)$ is a neighborhood of $m$. Since all maps $\sigma_{g}$ are diffeomorphisms of $M, \sigma_{g}(G . m)= g G . m=G . m$ also is a neighborhood of $g . m$, so that $G . m$ is open.

Corollary 10.1.7. For each $m \in M$,

$$
\mathbf{L}\left(G_{m}\right)=\{x \in \mathbf{L}(G): \mathbf{L}(\sigma)(x)(m)=0\}
$$

The preceding proposition shows in particular that the orbit $\mathcal{O}_{m}$ is a submanifold if $m$ is a fixed point (zero-dimensional case) and if $\mathcal{O}_{m}$ is open. Our next goal is to show that orbits of smooth actions always carry a natural manifold structure. This leads us to the geometry of homogeneous spaces in the next section.

\subsubsection*{Exercises for Section 10.1}
Exercise 10.1.1. Show that the orbits of a group action $\sigma: G \times M \rightarrow M$ form a partition of $M$.

Exercise 10.1.2. Show that the following maps define group actions and determine their orbits by naming a representative for each orbit ( $\mathbb{K}=\mathbb{R}, \mathbb{C}$ ):\\
(a) $\mathrm{GL}_{n}(\mathbb{K}) \times \mathbb{K}^{n} \rightarrow \mathbb{K}^{n},(g, v) \mapsto g v$;\\
(b) $\mathrm{O}_{n}(\mathbb{R}) \times \mathbb{R}^{n} \rightarrow \mathbb{R}^{n},(g, v) \mapsto g v$;\\
(c) $\mathrm{GL}_{n}(\mathbb{C}) \times M_{n}(\mathbb{C}) \rightarrow M_{n}(\mathbb{C}),(g, x) \mapsto g x g^{-1}$;\\
(d) $\mathrm{U}_{n}(\mathbb{K}) \times \operatorname{Herm}_{n}(\mathbb{K}) \rightarrow \operatorname{Herm}_{n}(\mathbb{K}),(g, x) \mapsto g x g^{-1}$;\\
(e) $\mathrm{GL}_{n}(\mathbb{K}) \times \operatorname{Herm}_{n}(\mathbb{K}) \rightarrow \operatorname{Herm}_{n}(\mathbb{K}),(g, x) \mapsto g x g^{*} ;$\\
(f) $\mathrm{O}_{n}(\mathbb{R}) \times\left(\mathbb{R}^{n} \times \mathbb{R}^{n}\right) \rightarrow \mathbb{R}^{n} \times \mathbb{R}^{n},(g,(x, y)) \mapsto(g x, g y)$.

Exercise 10.1.3. For a complex number $\lambda \in \mathbb{C}$, consider the smooth action

$$
\sigma: \mathbb{R} \times \mathbb{C} \rightarrow \mathbb{C}, \quad \sigma(t, z):=e^{t \lambda} z
$$

(1) Sketch the orbits of this action in dependence of $\lambda$.\\
(2) Under which conditions are there compact orbits?\\
(3) Describe the corresponding vector field.

Exercise 10.1.4. For complex numbers $\lambda_{1}, \lambda_{2} \in \mathbb{C}$, consider the smooth action

$$
\sigma: \mathbb{R} \times \mathbb{C}^{2} \rightarrow \mathbb{C}^{2}, \quad \sigma\left(t,\left(z_{1}, z_{2}\right)\right):=\left(e^{t \lambda_{1}} z_{1}, e^{t \lambda_{2}} z_{2}\right)
$$

(1) For which pairs ( $\lambda_{1}, \lambda_{2}$ ) are there bounded orbits?\\
(2) For which pairs ( $\lambda_{1}, \lambda_{2}$ ) are there compact orbits?\\
(3) Describe a situation where the closure of some orbit is compact, but the orbit itself is not.

Exercise 10.1.5. Let $G$ be a Lie group. For $x \in \mathfrak{g}$, let $x_{l}$ be the left invariant vector field on $G$ with $x_{l}(\mathbf{1})=x$ and $x_{r}$ the right invariant vector field on $G$ with $x_{r}(\mathbf{1})=x$. Show that\\
(i) the flow $\Phi_{t}^{x_{l}}: G \rightarrow G$ of $x_{l}$ is given by

$$
\Phi_{t}^{x_{l}}(g)=g \exp _{G}(t x)=\exp _{G}(t \operatorname{Ad}(g) x) g,
$$

whereas the flow $\Phi_{t}^{x_{r}}: G \rightarrow G$ of $x_{r}$ is given by $\Phi_{t}^{x_{r}}(g)=\exp _{G}(t x) g$,\\
(ii) $x_{r}(g)=T\left(c_{g}\right) x_{l}(g)$,\\
(iii) for the (left) actions $\sigma^{l}: G \times G \rightarrow G,(g, h) \mapsto \sigma^{l}(g, h)=\lambda_{g}(h)=g h$ and $\sigma^{r}: G \times G \rightarrow G,(g, h) \mapsto \sigma^{r}(g, h)=\rho_{g^{-1}}(h)=h g^{-1}$, we have $\mathbf{L}\left(\sigma^{l}\right) x=-x^{r}$ and $\mathbf{L}\left(\sigma^{r}\right)(x)=x_{l}$,\\
(iv) $x_{l} \mapsto-x_{r}$ yields a Lie algebra isomorphism between the Lie algebra of left invariant vector fields and the Lie algebra of right invariant vector fields on $G$.

Exercise 10.1.6. Let $\sigma: G \times M \rightarrow M$ be a smooth action of a Lie group $G$ on a smooth manifold $M$. Show that $\left(\sigma_{g}\right)_{*}(\mathbf{L}(\sigma)(x))=\mathbf{L}(\sigma)(\operatorname{Ad}(g) x)$ for all $g \in G$ and $x \in \mathfrak{g}$.

\subsubsection*{Transitive Actions and Homogeneous Spaces}
The main result of this section is that for any smooth action of a Lie group $G$ on a smooth manifold $M$, all orbits carry a natural manifold structure. First, we take a closer look at transitive actions, i.e., actions with a single orbit.

Definition 10.1.8. (a) Let $G$ be a group and $H$ a subgroup. We write

$$
G / H:=\{g H: g \in G\}
$$

for the set of left cosets of $H$ in $G$ and $q_{G / H}: G \rightarrow G / H, g \mapsto g H$ for the quotient map. Then

$$
\sigma: G \times G / H \rightarrow G / H, \quad(g, x H) \mapsto g x H
$$

defines a transitive action of $G$ on the set $G / H$ (easy exercise).\\
(b) Let $G$ be a group and $\sigma_{1}: G \times M_{1} \rightarrow M_{1}$ and $\sigma_{2}: G \times M_{2} \rightarrow M_{2}$ two actions of the group $G$ on sets. A map $f: M_{1} \rightarrow M_{2}$ is called $G$-equivariant if

$$
f(g . m)=g . f(m) \text { holds for all } g \in G, m \in M_{1} \text {. }
$$

Remark 10.1.9. Let $\sigma: G \times M \rightarrow M$ be an action of the group $G$ on the set $M$. Fix $m \in M$. Then the orbit map (cf. Definition 9.1.11)

$$
\sigma^{m}: G \rightarrow \mathcal{O}_{m} \subseteq M, \quad g \mapsto g . m
$$

factors through a bijective map

$$
\bar{\sigma}^{m}: G / G_{m} \rightarrow \mathcal{O}_{m}, \quad g G_{m} \mapsto g . m
$$

which is equivariant with respect to the $G$-actions on $G / G_{m}$ and $M$ (Exercise).

The preceding remark shows that if we want to obtain a manifold structure on orbits of smooth actions, it is natural to try to define a manifold structure on the coset spaces $G / H$ for closed subgroups $H$ of a Lie group $G$.

Theorem 10.1.10. Let $G$ be a Lie group and $H \leq G$ a closed subgroup. Then the coset space $G / H$, endowed with the quotient topology, carries a natural manifold structure for which the quotient map $q: G \rightarrow G / H, g \mapsto g H$ is a submersion.

Moreover, $\sigma: G \times G / H \rightarrow G / H,(g, x H) \mapsto g x H$ defines a smooth action of $G$ on $G / H$.

Proof. Let $E \subseteq \mathbf{L}(G)$ be a vector space complement of the subspace $\mathbf{L}(H)$ and $V_{E}$ be as in the Closed Subgroup Theorem 9.3.7(iii).

Step 1 (The topology on $G / H$ ): We endow $M:=G / H$ with the quotient topology. Since for each open subset $O \subseteq G$ the product $O H$ is open (Exercise 9.3.10), the openness of $O H=q^{-1}(q(O))$ shows that $q$ is an open map, i.e., maps open subsets to open subsets.

To see that $G / H$ is a Hausdorff space, let $g_{1}, g_{2} \in G$ with $g_{1} H \neq g_{2} H$, i.e., $g_{1} \notin g_{2} H$. Since $H$ is closed, there exists a $\mathbf{1}$-neighborhood $U_{1}$ in $G$ with $U_{1} g_{1} \cap g_{2} H=\emptyset$, and further a symmetric 1-neighborhood $U_{2}$ with $U_{2}^{-1} U_{2} \subseteq U_{1}$. Then $U_{2} g_{1} H$ and $U_{2} g_{2} H$ are disjoint $q$-saturated open subset of $G$, so that

$$
q\left(U_{2} g_{1} H\right)=q\left(U_{2} g_{1}\right) \quad \text { and } \quad q\left(U_{2} g_{2} H\right)=q\left(U_{2} g_{2}\right)
$$

are disjoint open subsets of $G / H$, separating $g_{1} H$ and $g_{2} H$. This shows that $G / H$ is a Hausdorff space.

We also observe that the action map $\sigma$ is continuous because $\operatorname{id}_{G} \times q: G \times G \rightarrow G \times G / H$ is a quotient map since $q$ is open (cf. Exercise 8.2.13) and

$$
\sigma \circ\left(\operatorname{id}_{G} \times q\right)=q \circ m_{G}: G \times G \rightarrow G / H, \quad(g, x) \mapsto g x H
$$

is continuous.\\
Step 2 (The atlas of $G / H$ ): Let $W:=q\left(\exp _{G}\left(V_{E}\right)\right)$ with $V_{E}$ as above and define a smooth map

$$
p_{E}: q^{-1}(W)=\exp _{G}\left(V_{E}\right) H \rightarrow V_{E} \quad \text { by } \quad \exp _{G}(x) h \mapsto x .
$$

Since $q^{-1}(W)$ is open in $G, W$ is open in $G / H$. Moreover, a subset $O \subseteq W$ is open if and only if $q^{-1}(O) \subseteq q^{-1}(W)$ is open. Since $q^{-1}(O)= \exp _{G}\left(p_{E}\left(q^{-1}(O)\right)\right) H$ and $q^{-1}(W) \cong V_{E} \times H$, this is equivalent to $p_{E}\left(q^{-1}(O)\right)$ being open in $V_{E}$. Therefore, the map $\psi: W \rightarrow V_{E}, q(g) \mapsto p_{E}(g)$ is a homeomorphism and ( $\psi, W$ ) is a chart of $G / H$.

For $g \in G$, we put $W_{g}:=g . W$ and $\psi_{g}(x)=\psi\left(g^{-1} . x\right)$. Since all maps $\sigma_{g}: G / H \rightarrow G / H$ are homeomorphisms (by Step 1), we thus get charts $\left(\psi_{g}, W_{g}\right)_{g \in G}$, and it is clear that $\bigcup_{g \in G} W_{g}=G / H$.

We claim that this collection of homeomorphisms is a smooth atlas of $G / H$. Let $g_{1}, g_{2} \in G$ and assume that $W_{g_{1}} \cap W_{g_{2}} \neq \emptyset$. We then have for $x \in V_{E}$ :

$$
\begin{aligned}
\psi_{g_{1}} \circ \psi_{g_{2}}^{-1}(x) & =\psi\left(g_{1}^{-1} g_{2} \cdot \psi^{-1}(x)\right)=\psi\left(g_{1}^{-1} g_{2} \cdot q\left(\exp _{G}(x)\right)\right) \\
& =\psi\left(q\left(g_{1}^{-1} g_{2} \exp _{G}(x)\right)\right)=p_{E}\left(g_{1}^{-1} g_{2} \exp _{G}(x)\right)
\end{aligned}
$$

Since $p_{E}$ is smooth, this map is smooth on its open domain, which shows that $\left(\psi_{g}, W_{g}\right)_{g \in G}$ is a smooth atlas of $G / H$.

Step 3 (Smoothness of the maps $\sigma_{g}$ ): For $g_{1}, g_{2} \in G$, we have $\sigma_{g_{1}}\left(W_{g_{2}}\right)=W_{g_{1} g_{2}}$ and $\psi_{g_{1} g_{2}} \circ \sigma_{g_{1}}=\psi_{g_{2}}$, which immediately implies that $\left.\sigma_{g_{1}}\right|_{W_{g_{2}}}: W_{g_{2}} \rightarrow W_{g_{1} g_{2}}$ is smooth. Since $g_{2}$ was arbitrary, all maps $\sigma_{g}, g \in G$, are smooth. From $\sigma_{g} \circ \sigma_{g^{-1}}=\operatorname{id}_{M}$, we further derive that they are diffeomorphisms.

Step 4 ( $q$ is a submersion): The smoothness of $q$ on $q^{-1}(W)$ follows from $\psi(q(g))=p_{E}(g)$ and the smoothness of $p_{E}$ on $q^{-1}(W)$. Moreover, $T_{\mathbf{1} H}(\psi) T_{\mathbf{1}}(q)=T_{\mathbf{1}}\left(p_{E}\right): \mathbf{L}(G) \rightarrow E$ is the linear projection onto $E$ with kernel $\mathbf{L}(H)$, hence surjective. This proves that $T_{\mathbf{1}}(q)$ is surjective.

For each $g \in G$, we have $q \circ \lambda_{g}=\sigma_{g} \circ q$, so that Step 3 implies that $q$ is smooth on all of $G$. Taking derivatives, we obtain

$$
T_{g}(q) \circ T_{\mathbf{1}}\left(\lambda_{g}\right)=T_{\mathbf{1} H}\left(\sigma_{g}\right) \circ T_{\mathbf{1}}(q)
$$

and since all $\sigma_{g}$ are diffeomorphisms, this implies that all differentials $T_{g}(q)$ are surjective, hence that $q$ is a submersion.

Step 5 (Smoothness of the action of $G / H$ ): Since $q$ is a submersion, the product map $\operatorname{id}_{G} \times q: G \times G \rightarrow G \times G / H$ also is a submersion. In view of Proposition 8.3.16, it therefore suffices to show that

$$
\sigma \circ\left(\operatorname{id}_{G} \times q\right): G \times G \rightarrow G / H
$$

is a smooth map, which follows from $\sigma \circ\left(\operatorname{id}_{G} \times q\right)=q \circ m_{G}$.

For later use we collect some facts from Step 1 of the proof of Theorem 10.1.10 into a separate corollary.

Corollary 10.1.11. Let $G$ be a Lie group and $H \leq G$ a closed subgroup. Then for any $x \in G / H$ there exists an open neighborhood $U \subseteq G / H$ and a smooth section $\sigma: U \rightarrow G$ for the quotient map $q: G \rightarrow G / H$ such that

$$
m_{\sigma}: U \times H \rightarrow \sigma(U) H, \quad(u, h) \mapsto \sigma(u) h
$$

is a diffeomorphism onto an open subset of $G$.\\
The following corollary shows that for each smooth group action, all orbits carry natural manifold structures. Not all these manifold structures turn these orbits into submanifolds, as the dense wind (see Example 9.3.12) shows.

Corollary 10.1.12. Let $\sigma: G \times M \rightarrow M$ be a smooth action of the Lie group $G$ on $M$. Then for each $m \in M$ the orbit map $\sigma^{m}: G \rightarrow M, g \mapsto g . m$ factors through a smooth injective equivariant map

$$
\bar{\sigma}^{m}: G / G_{m} \rightarrow M, \quad g G_{m} \mapsto g \cdot m,
$$

whose image is the set $\mathcal{O}_{m}$.\\
Proof. The existence of the map $\bar{\sigma}^{m}$ is clear (Remark 10.1.9). Since the quotient map $q: G \rightarrow G / G_{m}$ is a submersion, the smoothness of $\bar{\sigma}^{m}$ follows from the smoothness of the map $\bar{\sigma}^{m} \circ q=\sigma^{m}$ (Proposition 8.3.16).

The preceding corollary provides on each orbit $\mathcal{O}_{m}$ of a smooth Lie group action the structure of a smooth manifold. Its dimension is given by

$$
\begin{aligned}
\operatorname{dim}\left(G / G_{m}\right) & =\operatorname{dim} G-\operatorname{dim} G_{m}=\operatorname{dim} \mathbf{L}(G)-\operatorname{dim} \mathbf{L}\left(G_{m}\right) \\
& =\operatorname{dim} \mathbf{L}(\sigma)(\mathbf{L}(G))(m)
\end{aligned}
$$

because $\mathbf{L}\left(G_{m}\right)$ is the kernel of the linear map

$$
\mathbf{L}(G) \rightarrow T_{m}(M), \quad x \mapsto \mathbf{L}(\sigma)(x)(m)
$$

(Corollary 10.1.7). In this sense, we may identify the subspace $\mathbf{L}(\sigma)(\mathbf{L}(G))(m) \subseteq T_{m}(M)$ with the tangent space of the orbit $\mathcal{O}_{m}$.

We want to show that if $G$ has at most countably many connected components, then $\mathcal{O}_{m}$ is always an initial submanifold. For this we have a lemma.

Lemma 10.1.13. Let $\sigma: G \times M \rightarrow M$ be a smooth action and $m \in M$. Suppose that $N$ is a closed submanifold of an open neighborhood of $m$ in $M$ such that $m \in N$ and $T_{m}(M)=T_{m}(N) \oplus T_{\mathbf{1}}\left(\sigma^{m}\right)(\mathfrak{g})$. Further, suppose that $C$ is a closed submanifold of an open neighborhood of $\mathbf{1}$ in $G$ such that $\mathfrak{g}=T_{\mathbf{1}}(C) \oplus \mathfrak{g}_{m}$. Then there are open neighborhoods $C_{o}, N_{o}$, and $M_{o}$ of $\mathbf{1}$ and $m$ in $C, N$ and $M$, respectively, such that

$$
\Phi:=\left.\sigma\right|_{C_{o} \times N_{o}}: C_{o} \times N_{o} \rightarrow M_{o}:=\sigma\left(C_{o} \times N_{o}\right) \subseteq M
$$

is a diffeomorphism.

Proof. We have $T_{(\mathbf{1}, m)}(\Phi)(x, v)=v+T_{\mathbf{1}}\left(\sigma^{m}\right) x$. If this expression vanishes, then $v \in T_{m}(N) \cap T_{\mathbf{1}}\left(\sigma^{m}\right)(\mathfrak{g})=\{0\}$ also implies $T_{\mathbf{1}}\left(\sigma^{m}\right) x=0$. Therefore, $T_{(\mathbf{1}, m)}(\Phi)$ is injective. From

$$
T_{m}(M)=T_{m}(N)+T_{\mathbf{1}}\left(\sigma^{m}\right)(\mathfrak{g})=T_{m}(N)+T_{\mathbf{1}}\left(\sigma^{m}\right)\left(T_{\mathbf{1}}(C)\right)
$$

we further derive that it is bijective. Now the claim follows with the Inverse Function Theorem 8.1.5.

Proposition 10.1.14. Let $\sigma: G \times M \rightarrow M$ be a smooth action and $m \in M$. If $G$ has at most countably many connected components, then the immersion $\bar{\sigma}^{m}: G / G_{m} \rightarrow \mathcal{O}_{m} \subseteq M$ defines on the orbit $\mathcal{O}_{m}$ the structure of an initial submanifold of $M$.

Proof. Using Lemma 10.1.13 and Corollary 10.1.11, we find submanifolds $N$ and $C$ of $M$ and $G$, and a neighborhood $U$ of $m$ in $M$ such that the action gives a diffeomorphism $\Phi: C \times N \rightarrow U$. Each of the sets $\sigma(C \times\{n\})$ belongs to a single $G$-orbit. We claim that, if that orbit happens to be $\mathcal{O}_{m}, \sigma(C \times\{n\})$ contains an open subset of $\mathcal{O}_{m}$. In fact, the manifold structure on $\mathcal{O}_{m}$ is the one inherited from $G / G_{m}$, so that $\sigma(C \times\{n\})$ corresponds to a set of the form $C g_{o} G_{m} \subseteq G / G_{m}$, where $\sigma\left(g_{o}, m\right)=n$ and $T_{\mathbf{1}}(C) \cap \mathbf{L}\left(G_{n}\right)=\{0\}$. This implies that the map $C \rightarrow G / G_{m}, c \mapsto c g G_{m}$ has surjective differential at $\mathbf{1}$, so that its image $C g G_{m}$ contains an open subset of $G / G_{m}$. Next we recall from Proposition 9.1.15(iv)(d) that every pairwise disjoint family of open subsets of $G$ is countable, and this property is inherited by $G / G_{m}$ which carries the quotient topology. Since the sets $\sigma(C \times\{n\}), n \in N$, are pairwise disjoint, the preceding argument implies that $N \cap \mathcal{O}_{m}$ is countable.

Now let $f: M^{\prime} \rightarrow \mathcal{O}_{m}$ be a map for which the composition with the inclusion map $\iota: \mathcal{O}_{m} \cong G / G_{m} \rightarrow M$ is smooth. Let $m^{\prime} \in M$. We have to show that $\iota^{-1} \circ f$ is smooth in a neighborhood of $m^{\prime}$. Since $G$ acts on $\mathcal{O}_{m}$ and $M$ by smooth maps, we may w.l.o.g. assume that $f\left(m^{\prime}\right)=m, f\left(M^{\prime}\right) \subseteq \sigma(C \times N)$, and that $M^{\prime}$ is connected. Now $\Phi^{-1} \circ f: M^{\prime} \rightarrow C \times N$ is a smooth map, and for the $N$-projection $p_{N}: C \times N \rightarrow N$ the image $\left(p_{N} \circ \Phi^{-1} \circ f\right)\left(M^{\prime}\right)$ is connected by smooth arcs. On the other hand, it is contained in the countable subset $N \cap \mathcal{O}_{m}$, hence trivial because every non-constant smooth arc is uncountable. This proves that $f\left(M^{\prime}\right) \subseteq \sigma(C \times\{m\})$, and the map $h:=p_{C} \circ \Phi^{-1} \circ f: M^{\prime} \rightarrow C$ is smooth. Now $f(x)=\sigma(h(x), m)$, and the smoothness of $G$ on $\mathcal{O}_{m}$ implies that $\iota^{-1} \circ f$ is smooth. This proves that $\iota$ defines on the subset $\mathcal{O}_{m}$ of $M$ the structure of a smooth initial submanifold.

Remark 10.1.15. The assumption that $G$ has at most countably many connected components is crucial for $\mathcal{O}_{m}$ to be an initial submanifold. To see this, consider the group $\mathbb{R}_{d}$ which is the additive group $\mathbb{R}$ endowed with the discrete topology. This is a 0-dimensional Lie group, and $\sigma(x, y)=x+y$ defines a transitive action of $G=\mathbb{R}_{d}$ on $M=\mathbb{R}$. In this case, the spaces $G / G_{m} \cong G$ are discrete, but $\mathcal{O}_{m}=\mathbb{R}$ is not.

In some cases, the orbit $\mathcal{O}_{m}$ may already have another manifold structure, for instance, if it is a submanifold of $M$. In this case, the following proposition says that this manifold structure coincides with the one induced by identifying it with $G / G_{m}$.

Corollary 10.1.16. Let $\sigma: G \times M \rightarrow M$ be a smooth action and $m \in M$. We assume that $G$ has at most countably many connected components. If $\mathcal{O}_{m}$ is a submanifold of $M$, then the map $\bar{\sigma}^{m}: G / G_{m} \rightarrow \mathcal{O}_{m}$ is a diffeomorphism.

Proof. We recall from Lemma 8.6.5 that the submanifold $\mathcal{O}_{m}$ of $M$ is initial. On the other hand, we have seen in Proposition 10.1.14 that the map $\bar{\sigma}^{m}: G / G_{m} \rightarrow M$ also defines an initial submanifold structure on $\mathcal{O}_{m}$. Therefore, the assertion follows from the uniqueness of initial submanifold structures (Remark 8.6.2).

Corollary 10.1.17. If $\sigma: G \times M \rightarrow M$ is a transitive smooth action of the Lie group $G$ on the manifold $M$ and $m \in M$, then the orbit map $\bar{\sigma}^{m}: G / G_{m} \rightarrow M$ is a $G$-equivariant diffeomorphism.

Definition 10.1.18. The manifolds of the form $M=G / H$, where $H$ is a closed subgroup of a Lie group $G$, are called homogeneous spaces. We know already that the canonical action of $G$ on $G / H$ is smooth and transitive, and Corollary 10.1.17 shows the converse, i.e., that each transitive action is equivalent to the action on some $G / H$ because there exists an equivariant diffeomorphism.

\subsubsection*{Examples}
Example 10.1.19 (Graßmannians). Let $M:=\operatorname{Gr}_{k}\left(\mathbb{R}^{n}\right)$ denote the set of all $k$-dimensional subspaces of $\mathbb{R}^{n}$, the Graßmann manifold of degree $k$. We know from linear algebra that the natural action

$$
\sigma: \operatorname{GL}_{n}(\mathbb{R}) \times \operatorname{Gr}_{k}\left(\mathbb{R}^{n}\right) \rightarrow \operatorname{Gr}_{k}\left(\mathbb{R}^{n}\right), \quad(g, F) \mapsto g(F)
$$

is transitive (Exercise). Let $F:=\operatorname{span}\left\{e_{1}, \ldots, e_{k}\right\}$. Writing elements of $M_{n}(\mathbb{R})$ as $2 \times 2$-block matrices, according to

$$
M_{n}(\mathbb{R})=\left(\begin{array}{cc}
M_{k}(\mathbb{R}) & M_{k, n-k}(\mathbb{R}) \\
M_{n-k, k}(\mathbb{R}) & M_{n-k}(\mathbb{R})
\end{array}\right)
$$

the stabilizer of $F$ in $\mathrm{GL}_{n}(\mathbb{R})$ is

$$
\mathrm{GL}_{n}(\mathbb{R})_{F}:=\left\{\left(\begin{array}{cc}
a & b \\
0 & d
\end{array}\right): a \in \mathrm{GL}_{k}(\mathbb{R}), b \in M_{k, n-k}(\mathbb{R}), d \in \mathrm{GL}_{n-k}(\mathbb{R})\right\}
$$

which is a closed subgroup. Then the homogeneous space $\mathrm{GL}_{n}(\mathbb{R}) / \mathrm{GL}_{n}(\mathbb{R})_{F}$ carries a natural manifold structure, and since the orbit map of $F$ induces a bijection

$$
\bar{\sigma}^{F}: \mathrm{GL}_{n}(\mathbb{R}) / \mathrm{GL}_{n}(\mathbb{R})_{F} \rightarrow \mathrm{Gr}_{k}\left(\mathbb{R}^{n}\right), \quad g \mathrm{GL}_{n}(\mathbb{R})_{F} \mapsto g(F),
$$

we obtain a manifold structure on $\operatorname{Gr}_{k}\left(\mathbb{R}^{n}\right)$ for which the natural action of $\mathrm{GL}_{n}(\mathbb{R})$ is smooth.

The dimension of $\operatorname{Gr}_{k}\left(\mathbb{R}^{n}\right)$ is given by\\
$\operatorname{dim} \mathrm{GL}_{n}(\mathbb{R})-\operatorname{dim} \mathrm{GL}_{n}(\mathbb{R})_{F}=n^{2}-\left(k^{2}+(n-k)^{2}+k(n-k)\right)=k(n-k)$.\\
Note that for $k=1$ we obtain the manifold structure on the projective space $\mathbb{P}\left(\mathbb{R}^{n}\right)$.

Example 10.1.20 (Flag manifolds). A flag in $\mathbb{R}^{n}$ is a tuple

$$
\mathcal{F}=\left(F_{1}, \ldots, F_{m}\right)
$$

of subspaces of $\mathbb{R}^{n}$ with

$$
F_{1} \subseteq F_{2} \subseteq \cdots \subseteq F_{m}
$$

Let $k_{i}:=\operatorname{dim} F_{i}$ and call $\left(k_{1}, \ldots, k_{m}\right)$ the signature of the flag. We write $\operatorname{Fl}\left(k_{1}, \ldots, k_{m}\right)$ for the set of all flags of signature ( $k_{1}, \ldots, k_{m}$ ) in $\mathbb{R}^{n}$. Clearly,

$$
\operatorname{Fl}\left(k_{1}, \ldots, k_{m}\right) \subseteq \operatorname{Gr}_{k_{1}}\left(\mathbb{R}^{n}\right) \times \cdots \times \operatorname{Gr}_{k_{m}}\left(\mathbb{R}^{n}\right)
$$

We also have a natural action of $\mathrm{GL}_{n}(\mathbb{R})$ on the product of the Graßmann manifolds by

$$
g \cdot\left(F_{1}, \ldots, F_{m}\right):=\left(g\left(F_{1}\right), \ldots, g\left(F_{m}\right)\right)
$$

To describe a base point, let

$$
F_{i}^{0}:=\operatorname{span}\left\{e_{1}, \ldots, e_{k_{i}}\right\}
$$

and note that

$$
\mathcal{F}^{0}:=\left(F_{1}^{0}, \ldots, F_{m}^{0}\right) \in \operatorname{Fl}\left(k_{1}, k_{2}, \ldots, k_{m}\right)
$$

From basic linear algebra, it follows that the action of $\mathrm{GL}_{n}(\mathbb{R})$ on the subset $\operatorname{Fl}\left(k_{1}, \ldots, k_{m}\right)$ is transitive, which is shown by choosing for each flag $\mathcal{F}$ of the given signature a basis $\left(b_{i}\right)_{1 \leq i \leq n}$ for $\mathbb{R}^{n}$ such that

$$
F_{i}:=\operatorname{span}\left\{b_{1}, \ldots, b_{k_{i}}\right\} \quad \text { for } \quad i=1, \ldots, m
$$

Writing elements of $M_{n}(\mathbb{R})$ as ( $m \times m$ )-block matrices according to the partition

$$
n=k_{1}+\left(k_{2}-k_{1}\right)+\left(k_{3}-k_{2}\right)+\cdots+\left(k_{m}-k_{m-1}\right)+\left(n-k_{m}\right)
$$

the stabilizer of $\mathcal{F}^{0}$ is given by

$$
\mathrm{GL}_{n}(\mathbb{R})_{\mathcal{F}^{0}}=\left\{\left(g_{i j}\right)_{i, j=1, \ldots, m}:\left(i>j \Rightarrow g_{i j}=0\right) ; g_{i i} \in \mathrm{GL}_{k_{i}-k_{i-1}}(\mathbb{R})\right\}
$$

which is a closed subgroup of $\mathrm{GL}_{n}(\mathbb{R})$. We now proceed as above to get a manifold structure on the set $\operatorname{Fl}\left(k_{1}, \ldots, k_{m}\right)$, turning it into a homogeneous space, called a flag manifold. (Exercise: Calculate the dimension of $\left.\operatorname{Fl}(1,2,3,4)\left(\mathbb{R}^{6}\right).\right)$

Example 10.1.21. The orthogonal group $\mathrm{O}_{n}(\mathbb{R})$ acts smoothly on $\mathbb{R}^{n}$, and its orbits are the spheres

$$
S(r):=\left\{x \in \mathbb{R}^{n}:\|x\|=r\right\}, \quad r \geq 0 .
$$

We know already that all these spheres carry natural manifold structures. Therefore, Corollary 10.1.17 implies that for each $r>0$ we have

$$
S(r) \cong \mathbb{S}^{n-1} \cong \mathrm{O}_{n}(\mathbb{R}) / \mathrm{O}_{n}(\mathbb{R})_{e_{1}},
$$

where

$$
\mathrm{O}_{n}(\mathbb{R})_{e_{1}}=\left\{\left(\begin{array}{ll}
1 & 0 \\
0 & d
\end{array}\right): d \in \mathrm{O}_{n-1}(\mathbb{R})\right\} \cong \mathrm{O}_{n-1}(\mathbb{R}) .
$$

\subsection*{Frame Bundles}
Tensor bundles are obtained naturally from the tangent bundle using constructions from linear algebra. The sections of these bundles are tensor fields. Tensor fields of various kinds play an important role in differential geometry and its applications. In particular, we give a unified construction of tensor bundles as associated bundles for the frame bundle.

We start by introducing the general concept of a fiber bundle which we then specialize to vector and principal bundles This construction requires bundles whose fibers are not necessarily vector spaces.

\subsubsection*{Fiber Bundles}
Definition 10.2.1. (a) A quadruple ( $\mathcal{F}, M, F, \pi$ ), where $\mathcal{F}, M$, and $F$ are smooth manifolds and $\pi: \mathcal{F} \rightarrow M$ is a smooth map, is called a fiber bundle with a typical fiber $F$ over $M$ if there exists an open covering $\left(U_{i}\right)_{i \in I}$ of $M$ and diffeomorphisms $\Phi_{i}: \pi^{-1}\left(U_{i}\right) \rightarrow U_{i} \times F$, called local trivializations, such that the following diagram commutes\\
\includegraphics[max width=\textwidth, center]{2025_10_18_3eaf04e77f3cb364fce5g-381}

Then $\mathcal{F}$ is called the total space of ( $\mathcal{F}, \pi$ ), and $M$ is called the base space of $(\mathcal{F}, \pi)$. The preimage $\mathcal{F}_{m}:=\pi^{-1}(m)$ is called the fiber over $m$. We also refer to the family $\left(\Phi_{i}\right)_{i \in I}$ as a local trivialization of the bundle.\\
(b) A morphism of fiber bundles $\varphi:\left(\mathcal{F}_{1}, M_{1}, F_{1}, \pi_{1}\right) \rightarrow\left(\mathcal{F}_{2}, M_{2}, F_{2}, \pi_{2}\right)$ is a smooth map $\varphi: \mathcal{F}_{1} \rightarrow \mathcal{F}_{2}$ for which there exists a smooth map $\varphi_{M}: M_{1} \rightarrow M_{2}$ with $\pi_{2} \circ \varphi=\varphi_{M} \circ \pi_{1}$, i.e., the following diagram commutes:\\
\includegraphics[max width=\textwidth, center]{2025_10_18_3eaf04e77f3cb364fce5g-382}

A morphism $\varphi: \mathcal{F}_{1} \rightarrow \mathcal{F}_{2}$ of fiber bundles is called an isomorphism if there exists a second morphism of fiber bundles $\psi: \mathcal{F}_{2} \rightarrow \mathcal{F}_{1}$ with $\psi \circ \varphi=\operatorname{id}_{\mathcal{F}_{1}}$ and $\varphi \circ \psi=\operatorname{id}_{\mathcal{F}_{2}}$. It is easy to see that a morphism of fiber bundles is an isomorphism of fiber bundles if and only if it is a diffeomorphism (Exercise).

An isomorphism $\varphi: \mathcal{F}_{1} \rightarrow \mathcal{F}_{2}$ of fiber bundles over $M$ is called an equivalence if $\varphi_{M}=\operatorname{id}_{M}$.\\
(c) The pair ( $M \times F, \mathrm{pr}_{M}$ ) is called the trivial fiber bundle with fiber $F$.\\
(d) If ( $\mathcal{F}, M, F, \pi$ ) is a fiber bundle, then for each open subset $U \subseteq M$ and $\mathcal{F}_{U}:=\pi^{-1}(U)$, we obtain the fiber bundle $\left(\mathcal{F}_{U}, U, F,\left.\pi\right|_{\mathcal{F}_{U}}\right)$. A bundle chart ( $\psi, \mathcal{F}_{U}$ ) is an equivalence $\varphi: \mathcal{F}_{U} \rightarrow U \times F$ followed by a chart $\rho=\rho_{U} \times \rho_{F}$ with $\rho_{U}: U \rightarrow \mathbb{R}^{n}$ and $\rho_{F}: F \rightarrow \mathbb{R}^{m}$. For $i, j \in I$ we put $U_{i j}:=U_{i} \cap U_{j}$. Then the map $\Phi_{j} \circ \Phi_{i}^{-1} \in \operatorname{Diff}\left(U_{i j} \times F\right)$ is a self-equivalence of the trivial bundle. This implies that it has the form

$$
(x, f) \mapsto\left(x, \Phi_{j i}(x)(f)\right)
$$

for a function $\Phi_{j i}: U_{i j} \rightarrow \operatorname{Diff}(F)$, called the transition function for $\mathcal{F}$ from $\Phi_{i}$ to $\Phi_{j}$.

The construction of the tangent bundle and the tangent map yields examples of special bundles with vectors spaces as fibers. These will be call vector bundles.

Definition 10.2.2. (a) Let $M$ be a differentiable manifold and $V$ a finitedimensional vector space. A smooth vector bundle with a typical fiber $V$ on $M$ is a differentiable manifold $\mathbb{V}$ together with a smooth map $\pi: \mathbb{V} \rightarrow M$ such that\\
(i) All fibers $\mathbb{V}_{p}:=\pi^{-1}(p)$ carry the structure of a vector space.\\
(ii) For each $p \in M$, there are an open neighborhood $U$ of $p$ in $M$ and a diffeomorphism $\varphi_{U}: \pi^{-1}(U) \rightarrow U \times V$, called a local trivialization, with $\operatorname{pr}_{U} \circ \varphi_{U}=\pi$, and for all $q \in U$, the map

$$
\left.\operatorname{pr}_{V} \circ \varphi_{U}\right|_{\mathbb{V}_{q}}: \mathbb{V}_{q} \rightarrow V
$$

is a linear isomorphism, where $\operatorname{pr}_{U}(u, v)=u$ and $\operatorname{pr}_{V}(u, v)=v$.\\
(b) A smooth map $\sigma: M \rightarrow \mathbb{V}$ with $\pi \circ \sigma=\operatorname{id}_{M}$ is called a section of the bundle. The set $\Gamma(\mathbb{V})$ of sections of $\mathbb{V}$ is a vector space with respect to the pointwise addition and scalar multiplication.\\
(c) The vector bundle $\mathbb{V}=M \times V$ with the projection $\operatorname{pr}_{M}(m, v):=m$ and the obvious vector space structure on the fibers $\{p\} \times V$ is called the trivial vector bundle with a typical fiber $V$.

Definition 10.2.3. Let $\pi_{\mathbb{V}}: \mathbb{V} \rightarrow M$ and $\pi_{\mathbb{W}}: \mathbb{W} \rightarrow N$ be smooth vector bundles. A smooth map $\varphi: \mathbb{V} \rightarrow \mathbb{W}$ is called a morphism of vector bundles if there exists a smooth map $\varphi_{M}: M \rightarrow N$ with $\varphi_{M} \circ \pi_{\mathbb{V}}=\pi_{\mathbb{W}} \circ \varphi$ and the restrictions $\varphi_{p}: \mathbb{V}_{p} \rightarrow \mathbb{W}_{\varphi(p)}$ to the fibers are linear. We then have a commutative diagram\\
\includegraphics[max width=\textwidth, center]{2025_10_18_3eaf04e77f3cb364fce5g-383}

Example 10.2.4. The tangent bundle $T(M)$ of a smooth $n$-dimensional manifold $M$ is a smooth vector bundle with a typical fiber $V=\mathbb{R}^{n}$ and the tangent map $T(f): T(M) \rightarrow T(N)$ of a smooth map $f: M \rightarrow N$ is a morphism of vector bundles.

Remark 10.2.5. The fiberwise linearity of the local trivializations of a vector bundle (cf. Definition 10.2.2) shows that if a fiber bundle ( $\mathcal{F}, M, F, \pi$ ) is a vector bundle, then\\
(a) The typical fiber $F$ is a vector space.\\
(b) One can choose a local trivialization $\Phi_{i}$ in such a way that for $i, j \in I$ we always have $\Phi_{j i}\left(U_{i j}\right) \subseteq \mathrm{GL}(F)$.

If, conversely, (a) and (b) are satisfied, then

$$
\Phi_{i}^{-1}(x, v)+\Phi_{i}^{-1}(x, w):=\Phi_{i}^{-1}(x, v+w), \quad \lambda \Phi_{i}^{-1}(x, v):=\Phi_{i}^{-1}(x, \lambda v)
$$

defines on each fiber $\mathcal{F}_{x}, x \in U_{i}$, a vector space structure which does not depend on the choice of $i$. We thus obtain on $\mathcal{F}$ the structure of a vector bundle.

Example 10.2.6. Let $M$ be a manifold of dimension $n$ and $T M$ its tangent bundle. The construction of the manifold structure of $T M$ in Definition 8.3.6 shows that $T M$ is a vector bundle with a typical fiber $\mathbb{R}^{n}$ (cf. Example 10.2.4). A local trivialization is obtained from a smooth atlas $\left(\varphi_{i}, U_{i}\right)_{i \in I}$ of $M$ by

$$
\Phi_{i}^{-1}: U_{i} \times \mathbb{R}^{n} \rightarrow T\left(U_{i}\right)=T(M)_{U_{i}}, \quad(p, v) \mapsto T_{p}\left(\varphi_{i}\right)^{-1} v
$$

Thus the transition function is given by

$$
\Phi_{j i}(p)=T_{p}\left(\varphi_{j}\right) T_{p}\left(\varphi_{i}\right)^{-1}=\mathrm{d}\left(\varphi_{j} \circ \varphi_{i}^{-1}\right)\left(\varphi_{i}(p)\right) .
$$

Using the atlas $\left(\varphi_{i}, U_{i}\right)_{i \in I}$ and the identity as a chart for $\mathbb{R}^{n}$, the local trivialization yields the atlas for $T M$ described in Definition 8.3.6. If the $\varphi_{i^{-}}$ coordinates are denoted by $x_{1}, \ldots, x_{n}$ and the $\varphi_{j}$-coordinates are denoted by $y_{1}, \ldots, y_{n}$, the transition functions can be written in the form

$$
\Phi_{j i}(p)=\left(\frac{\partial y_{r}\left(\varphi_{i}(p)\right)}{\partial x_{s}}\right)_{1 \leq r, s \leq n} \in \mathrm{GL}_{n}(\mathbb{R}) \cong \mathrm{GL}\left(\mathbb{R}^{n}\right) .
$$

Definition 10.2.7 (Principal bundles). Let ( $\mathcal{F}, M, G, \pi$ ) be a fiber bundle over a Lie group $G$. If $\sigma: \mathcal{F} \times G \rightarrow \mathcal{F}$ is a smooth right action of $G$ on $\mathcal{F}$ (cf. Remark 9.1.12), then ( $\mathcal{F}, M, G, \pi, \sigma$ ) is called a principal bundle with structure group $G$ if there exists a local trivialization $\left(\Phi_{i}\right)_{i \in I}$ consisting of maps

$$
\Phi_{i}: \pi^{-1}\left(U_{i}\right) \rightarrow U_{i} \times G
$$

which are equivariant with respect to the natural right action $(x, h) g:=$ ( $x, h g$ ) of $G$ on $U_{i} \times G$, i.e.,

$$
\Phi_{i}^{-1}(x, h g)=\Phi_{i}^{-1}(x, h) g, \quad x \in U_{i}, h, g \in G
$$

Remark 10.2.8. For each principal bundle ( $\mathcal{F}, M, G, \pi, \sigma$ ), the action $\sigma$ of $G$ on $\mathcal{F}$ is free, i.e., $x g=x$ implies $g=\mathbf{1}$, and transitive on the fibers of $\pi$ (Exercise).

Example 10.2.9 (Frame bundles). Let ( $\mathbb{V}, M, V, \pi$ ) be a vector bundle. We set

$$
\operatorname{GL}(\mathbb{V}):=\coprod_{p \in M} \operatorname{Iso}\left(V, \mathbb{V}_{p}\right)
$$

where $\operatorname{Iso}(V, W)$ denotes the set of linear isomorphisms $V \rightarrow W$ of two vector spaces. Then the map

$$
\widetilde{\pi}: \mathrm{GL}(\mathbb{V}) \longrightarrow M, \quad \varphi \longmapsto p \text { if } \varphi \in \operatorname{Iso}\left(V, \mathbb{V}_{p}\right)
$$

is surjective, and

$$
\sigma: \mathrm{GL}(\mathbb{V}) \times \mathrm{GL}(V) \rightarrow \mathrm{GL}(\mathbb{V}), \quad(\varphi, g) \mapsto \varphi \circ g
$$

defines a right action of the Lie group $\mathrm{GL}(V)$ on $\mathrm{GL}(\mathbb{V})$ which is easily seen to be free.

We want to turn $\mathrm{GL}(\mathbb{V})$ into a $\mathrm{GL}(V)$-principal bundle. First, we have to construct a smooth manifold structure on $\mathrm{GL}(\mathbb{V})$. We start with a local trivialization $\left(\Phi_{i}\right)_{i \in I}$ of $\mathbb{V}$ so that $\Phi_{i}: \pi^{-1}\left(U_{i}\right) \rightarrow U_{i} \times V$. For each $i \in I$, we have a natural map

$$
\widetilde{\Psi}_{i}: U_{i} \times \mathrm{GL}(V) \rightarrow \mathrm{GL}(\mathbb{V}), \quad(x, g) \mapsto \Psi_{i, x} \circ g
$$

where $\Psi_{i, x}: V \rightarrow \mathbb{V}_{x}, v \mapsto \Phi_{i}^{-1}(x, v)$. Since the action of $\operatorname{GL}(V)$ on $\operatorname{GL}(\mathbb{V})$ is free and transitive on the fibers $\operatorname{Iso}\left(V, \mathbb{V}_{p}\right)$, the map $\widetilde{\Psi}_{i}$ is injective with image $\widetilde{\pi}^{-1}\left(U_{i}\right)$. We denote its inverse by $\widetilde{\Phi}_{i}$. For $i, j \in I$, the maps

$$
\begin{aligned}
\widetilde{\Psi}_{j}^{-1} \circ \widetilde{\Psi}_{i}: U_{i j} \times \mathrm{GL}(V) & \rightarrow U_{i j} \times \mathrm{GL}(V) \\
(x, g) & \mapsto\left(x,\left(\Psi_{j, x}^{-1} \circ \Psi_{i, x}\right) g\right)
\end{aligned}
$$

are smooth. In fact, $\Psi_{j, x}^{-1} \circ \Psi_{i, x}=\Phi_{j i}(x)$ is the transition function of $\mathbb{V}$ with respect to the local trivialization $\left(\Phi_{i}\right)_{i \in I}$.

We now endow $\mathrm{GL}(\mathbb{V})$ with the topology for which a subset $O \subseteq \mathrm{GL}(\mathbb{V})$ is open if and only if $\widetilde{\Psi}_{j}^{-1}(O)$ is open for each $j$. Since the maps $\widetilde{\Psi}_{j}^{-1} \circ \widetilde{\Psi}_{i}$ are homeomorphisms, each $\widetilde{\Psi}_{i}$ is an open embedding. Moreover, the projection $\widetilde{\pi}$ is continuous because all compositions

$$
\operatorname{pr}_{U_{i}}=\widetilde{\pi} \circ \widetilde{\Psi}_{i}: U_{i} \times \mathrm{GL}(V) \rightarrow U_{i}
$$

are continuous. From the continuity of $\widetilde{\pi}$ and the fact that the $\widetilde{\Psi}_{j}$ 's are open embeddings, it follows easily that $\mathrm{GL}(\mathbb{V})$ is a Hausdorff space.

From the smoothness of the maps $\widetilde{\Psi}_{j}^{-1} \circ \widetilde{\Psi}_{i}$ it further follows that $\operatorname{GL}(\mathbb{V})$ carries a unique smooth manifold structure for which the maps $\widetilde{\Psi}_{i}$ are diffeomorphisms onto open subsets. Thus the $\widetilde{\Phi}_{i}: \widetilde{\pi}^{-1}\left(U_{i}\right) \rightarrow U_{i} \times \operatorname{GL}(V)$ are $\mathrm{GL}(V)$-equivariant local trivializations, so that ( $\mathrm{GL}(\mathbb{V}), M, \mathrm{GL}(V), \widetilde{\pi}, \sigma$ ) is a principal bundle. It is called the frame bundle of $\mathbb{V}$. Note that the transition functions of the frame bundle with respect to the local trivialization $\left(\widetilde{\Phi}_{i}\right)_{i \in I}$ are given by

$$
\widetilde{\Phi}_{j i}(x)=\lambda_{\Phi_{j i}(x)}: \mathrm{GL}(V) \rightarrow \mathrm{GL}(V)
$$

where the $\Phi_{j i}$ are the transition functions of $\mathbb{V}$ with respect to the local trivialization $\left(\Phi_{i}\right)_{i \in I}$ and $\lambda_{g}$ denotes left multiplication by $g$.

If $\mathbb{V}=T M$ is the tangent bundle of a smooth manifold $M$ (Example 10.2 .4 ) with a typical fiber $V=\mathbb{R}^{n}$, we simply write $\operatorname{GL}(M):=\operatorname{GL}(T M)$ and call it the frame bundle of $M$. It is a principal bundle with the structure group $\mathrm{GL}\left(\mathbb{R}^{n}\right) \cong \mathrm{GL}_{n}(\mathbb{R})$.

Example 10.2.10 (Associated bundles). Let ( $P, M, G, \pi, \sigma$ ) be a principal bundle and $F$ a smooth manifold with a smooth left action $\tau$ : $G \times F \rightarrow F$. Then

$$
(p, f) \cdot g:=\left(p \cdot g, g^{-1} \cdot f\right)
$$

defines a smooth right $G$-action on $P \times F$. For each $x \in M$, each $G$-orbit in the invariant subset $P_{x} \times F \subseteq P \times F$ meets each of the sets $\left\{p_{0}\right\} \times F$ exactly once because $G$ acts freely on $P$.

Let

$$
P \times_{G} F:=P \times_{\tau} F:=(P \times F) / G:=\{G \cdot(p, f): p \in P, f \in F\}
$$

be the set of $G$-orbits in $P \times F$. We write $[p, f]:=G .(p, f)$ for the orbit of $(p, f)$ and $\widetilde{\pi}([p, f]):=\pi(p)$ for the projection to $M$.

To turn $P \times_{G} F$ into a smooth manifold, we start with a $G$-equivariant local trivialization $\left(\Phi_{i}\right)_{i \in I}$ of the principal bundle $P$ and consider the maps

$$
\widetilde{\Psi}_{i}: U_{i} \times F \rightarrow P \times_{G} F, \quad(x, f) \mapsto\left[\Phi_{i}^{-1}(x, \mathbf{1}), f\right]
$$

As we have already seen above, all these maps are injective. Moreover, if $\Phi_{j} \circ \Phi_{i}^{-1}(x, g)=\left(x, \Phi_{j i}(x) g\right)$, then

$$
\widetilde{\Psi}_{j}^{-1} \circ \widetilde{\Psi}_{i}(x, f)=\widetilde{\Psi}_{j}^{-1}\left[\Phi_{i}^{-1}(x, \mathbf{1}), f\right]=\left(x, \Phi_{j i}(x) \cdot f\right)
$$

follows from

$$
\begin{aligned}
\widetilde{\Psi}_{j}\left(x, \Phi_{j i}(x) \cdot f\right) & =\left[\Phi_{j}^{-1}(x, \mathbf{1}), \Phi_{j i}(x) \cdot f\right]=\left[\Phi_{j}^{-1}\left(x, \Phi_{j i}(x)\right), f\right] \\
& =\left[\Phi_{i}^{-1}(x, \mathbf{1}), f\right]
\end{aligned}
$$

Now one argues as in Example 10.2.9 to see that $P \times_{G} F$ endowed with the quotient topology is a Hausdorff space and that it carries a unique smooth structure for which all the sets $\widetilde{\pi}^{-1}\left(U_{i}\right) \subseteq P \times_{G} F$ are open and the maps $\widetilde{\Psi}_{i}: U_{i} \times F \rightarrow \widetilde{\pi}^{-1}\left(U_{i}\right)$ are diffeomorphisms. We thus obtain a fiber bundle ( $P \times_{G} F, M, F, \widetilde{\pi}$ ) for which the maps $\widetilde{\Phi}_{i}:=\widetilde{\Psi}_{i}^{-1}: \widetilde{\pi}^{-1}\left(U_{i}\right) \rightarrow U_{i} \times F$ form local trivializations. It is called the bundle associated with $P$ and the $G$-space $F$. Note that the transition functions of $P \times_{G} F$ with respect to the local trivialization $\left(\widetilde{\Phi}_{i}\right)_{i \in I}$ are given by

$$
\widetilde{\Phi}_{j i}(x)(f)=\Phi_{j i}(x) \cdot f
$$

where the $\Phi_{j i}$ are the transition functions of $P$ with respect to the local trivialization $\left(\Phi_{i}\right)_{i \in I}$.

Example 10.2.11. An interesting special case arises for the one point space $F=\{*\}$ (endowed with the trivial $G$-action). Then $P \times_{G} F \cong P / G$ is the set of $G$-orbits in $P$. Since the map $\pi: P \rightarrow M$ is a submersion whose fibers are the $G$-orbits in $P$, it follows that the map $\widetilde{\pi}: P \times_{G}\{*\} \rightarrow M$ is a diffeomorphism.

Example 10.2.12. Let $(\mathbb{V}, M, V, \pi)$ be a vector bundle as in Example 10.2.9. Then its frame bundle $\mathrm{GL}(\mathbb{V})$ is a $\mathrm{GL}(V)$-principal bundle and we have the canonical smooth action $\tau$ : $\mathrm{GL}(V) \times V \rightarrow V,(g, v) \mapsto g v$. From the definition of the frame bundle $\operatorname{GL}(\mathbb{V})$, we immediately obtain a smooth map

$$
\mathrm{ev}: \mathrm{GL}(\mathbb{V}) \times V \rightarrow \mathbb{V}, \quad(p, v) \mapsto p(v)
$$

which is a morphism of bundles over $M$. This map is constant on the $\mathrm{GL}(V)$ orbits in $\mathrm{GL}(\mathbb{V}) \times V$ under the action $g \cdot(p, v):=\left(p \circ g^{-1}, g v\right)$, hence factors through a smooth bundle morphism

$$
\overline{\mathrm{eV}}: \mathrm{GL}(\mathbb{V}) \times_{\mathrm{GL}(V)} V \rightarrow \mathbb{V}, \quad[p, v] \mapsto p(v)
$$

(cf. Proposition 8.3.16). On each fiber, the map

$$
\left(\operatorname{Iso}\left(V, \mathbb{V}_{x}\right) \times V\right) / \operatorname{GL}(V) \rightarrow \mathbb{V}_{x}, \quad[p, v] \mapsto p(v)
$$

is bijective because for any fixed isomorphism $p_{0}: V \rightarrow \mathbb{V}_{x}$, each $\operatorname{GL}(V)$-orbit meets the set $\left\{p_{0}\right\} \times V$ exactly once (cf. Exercise 10.2.8). This implies that $\overline{\mathrm{ev}}$ is a bijective submersion, and this implies that it is an equivalence of vector bundles (cf. Definition 10.2.3).

Corollary 10.2.13. Let $G$ be a Lie group, $H$ a closed subgroup and $q: G \rightarrow G / H$ the quotient map. We write $\mathfrak{g}:=\mathbf{L}(G)$ and $\mathfrak{h}:=\mathbf{L}(H)$. Let $p_{o}:=\mathbf{1} H=q(\mathbf{1})$ be the canonical base point for $G / H$ and $\sigma_{g}$ be the left translation by $g$ on $G / H$ (see Definition 10.1.8). Then $T_{p_{o}}(G / H) \cong \mathfrak{g} / \mathfrak{h}$ and with the $H$-action on $\mathfrak{g} / \mathfrak{h}$ given by $\operatorname{Ad}_{\mathfrak{g} / \mathfrak{h}}(h)(x+\mathfrak{h})=\operatorname{Ad}(h) x+\mathfrak{h}$, the map

$$
G \times_{H} \mathfrak{g} / \mathfrak{h} \rightarrow T(G / H), \quad[g, x+\mathfrak{h}] \mapsto T_{p_{o}}\left(\sigma_{g}\right) T_{\mathbf{1}}(q) x
$$

is a $G$-equivariant bundle isomorphism.\\
Proof. The $G$-equivariance is clear once we have shown that the map is welldefined. Thus, in view of Theorem 10.1.10, it only remains to show that

$$
T_{p_{o}}\left(\sigma_{g}\right) \circ T_{\mathbf{1}}(q)=T_{p_{o}}\left(\sigma_{g h}\right) \circ T_{\mathbf{1}}(q) \circ \operatorname{Ad}(h)^{-1}
$$

In view of $\sigma_{g} \circ q=q \circ \lambda_{g}$ and $q \circ \rho_{h}=q$, this follows from

$$
\begin{aligned}
T_{p_{o}}\left(\sigma_{g h}\right) \circ T_{\mathbf{1}}(q) \circ \operatorname{Ad}(h)^{-1} & =T_{\mathbf{1}}\left(q \circ \lambda_{g h} \circ c_{h}^{-1}\right)=T_{\mathbf{1}}\left(q \circ \lambda_{g} \circ \rho_{h}\right) \\
& =T_{\mathbf{1}}\left(q \circ \lambda_{g}\right)=T_{p_{o}}\left(\sigma_{g}\right) \circ T_{\mathbf{1}}(q) .
\end{aligned}
$$

\subsubsection*{Sections}
Sections of bundles can be viewed as generalizations of functions. In particular, in the case of vector bundles they unify a large number of concepts well-known in classical analysis: Functions, vector fields, differential forms, and tensor fields.

Definition 10.2.14. Let ( $\mathcal{F}, M, F, \pi$ ) be a fiber bundle. A map $s: M \rightarrow \mathcal{F}$ is called a section of $\mathcal{F}$ if $\pi \circ s=\operatorname{id}_{M}$. The set of all smooth sections of $\mathcal{F}$ is denoted by $\Gamma(\mathcal{F})$. If $\mathcal{F}$ is a vector bundle, then $\Gamma_{c}(\mathcal{F})$ denotes the vector space of all smooth sections $s$ with compact support, i.e., $s(m)=0$ for all $m$ outside of a compact set.

The following proposition explains how smooth sections are described in terms of local trivializations. It is an elementary consequence of the definitions.

Proposition 10.2.15. Let ( $\mathcal{F}, M, F, \pi$ ) be a fiber bundle with local trivializations $\Phi_{i}: \mathcal{F}_{U_{i}} \rightarrow U_{i} \times F$, transition functions $\Phi_{j i}$, and let $s: M \rightarrow \mathcal{F}$ be a smooth section. Then we obtain for each $i \in I$ a smooth function $s_{i}: U_{i} \rightarrow F$, determined by

$$
\Phi_{i}(s(x))=\left(x, s_{i}(x)\right) \quad \text { for } x \in U_{i}
$$

and these functions satisfy the relation

$$
s_{j}(x)=\Phi_{j i}(x)\left(s_{i}(x)\right) \quad \text { for } x \in U_{i} \cap U_{j}
$$

If, conversely, $\left(s_{i}\right)_{i \in I}$ is a family of smooth functions $s_{i}: U_{i} \rightarrow F$ satisfying (10.2), then (10.1) yields a well-defined smooth section of $\mathcal{F}$.

Example 10.2.16 (Density bundles). Let $M$ be a smooth manifold of dimension $n$ and $\widetilde{\pi}: \mathrm{GL}(M) \rightarrow M$ be the frame bundle of $M$. For $r>0$, consider the one-dimensional representation $\Delta_{r}: \mathrm{GL}_{n}(\mathbb{R}) \rightarrow \mathbb{R}_{+}^{\times}, \Delta_{r}(g):= |\operatorname{det} g|^{-r}$. Then the line bundle $|\Lambda|^{r}(M):=\mathrm{GL}(M) \times{ }_{\Delta_{r}} \mathbb{R}$ associated with this representation via Example 10.2.10 is called the $r$-density bundle of $M$. Using the notation from Example 10.2.6, the transition functions of $|\Lambda|^{r}(M)$ can be written

$$
\Phi_{j i}(p)=\left|\operatorname{det} \mathrm{d}\left(\varphi_{j} \circ \varphi_{i}^{-1}\right)\left(\varphi_{i}(p)\right)\right|^{-r} .
$$

The sections of $|\Lambda|^{r}(M)$ are called $r$-densities. For $r=1$, we will simply call them densities and write $|\Lambda|(M)$. Their transformation properties under coordinate changes is reminiscent of the change of variables formula in multivariable calculus, and indeed, they play an important role in the integration theory on manifolds.

We give an alternative description of the density bundles: Let $p \in M$ and $L_{p}^{(r)}$ be the space of $r$-densities on $T_{p}(M)$, i.e., all continuous functions $f: T_{p}(M)^{n} \rightarrow \mathbb{R}$ satisfying

$$
f\left(A v_{1}, \ldots, A v_{n}\right)=|\operatorname{det} A|^{r} f\left(v_{1}, \ldots, v_{n}\right)
$$

for all $v_{1}, \ldots, v_{n} \in T_{p}(M)$ and $A \in \operatorname{End}\left(T_{p}(M)\right)$. Note that $f \in L_{p}^{(r)}$ is uniquely determined by its values on a single basis for $T_{p}(M)$, so that $L_{p}^{(r)}$ is one-dimensional. Set $L^{(r)}(M)=\coprod_{p \in M} L_{p}^{(r)}$ and use a smooth atlas $\left(\varphi_{i}, U_{i}\right)_{i \in I}$ of $M$ to define maps $\Psi_{i}: U_{i} \times \mathbb{R} \rightarrow L^{(r)}(M)$ via

$$
\Psi_{i}(p, t)\left(\left.\frac{\partial}{\partial x_{1}^{(i)}}\right|_{p}, \ldots,\left.\frac{\partial}{\partial x_{n}^{(i)}}\right|_{p}\right)=t
$$

where the $x_{k}^{(i)}$ are the coordinate functions for the chart $\varphi_{i}$. Then $\Psi_{i}$ is injective with image $\pi^{-1}\left(U_{i}\right)$, where $\pi: L^{(r)}(M) \rightarrow M$ is the obvious base point projection. If now $p \in U_{i} \cap U_{j}$, we have $\Psi_{i}(p, t)=\Psi_{j}(p, s)$ if and only if

$$
\begin{aligned}
t & =\Psi_{j}(p, s)\left(\left.\frac{\partial}{\partial x_{1}^{(i)}}\right|_{p}, \ldots,\left.\frac{\partial}{\partial x_{n}^{(i)}}\right|_{p}\right) \\
& =\left|\operatorname{det}\left(\frac{\partial x_{l}^{(j)}}{\partial x_{k}^{(i)}}\left(\varphi_{i}(p)\right)\right)\right|^{r} \Psi_{j}(p, s)\left(\left.\frac{\partial}{\partial x_{1}^{(j)}}\right|_{p}, \ldots,\left.\frac{\partial}{\partial x_{n}^{(j)}}\right|_{p}\right) \\
& =\left|\operatorname{det}\left(\frac{\partial x_{l}^{(j)}}{\partial x_{k}^{(i)}}\left(\varphi_{i}(p)\right)\right)\right|^{r} s \\
& =\left|\operatorname{det} \mathrm{d}\left(\varphi_{j} \circ \varphi_{i}^{-1}\right)\left(\varphi_{i}(p)\right)\right|^{r} s .
\end{aligned}
$$

Thus we can equip $L^{(r)}(M)$ with a line bundle structure such that the $\Phi_{i}:= \Psi_{i}^{-1}$ form a local trivialization. The corresponding transition functions are

$$
\Phi_{j i}(p)=\left|\operatorname{det} \mathrm{d}\left(\varphi_{j} \circ \varphi_{i}^{-1}\right)\left(\varphi_{i}(p)\right)\right|^{-r}
$$

which shows that $L^{(r)}(M)$ is indeed the bundle of $r$-densities.

It is also possible to give an explicit isomorphism $|\Lambda|^{r}(M) \rightarrow L^{(r)}(M)$ : For $\varphi \in \mathrm{GL}(M)$ and $t \in \mathbb{R}$, let $\mu_{\varphi, t}$ be the $r$-density on $T_{\widetilde{\pi}(\varphi)}(M)$ satisfying

$$
\mu_{\varphi, t}\left(\varphi\left(e_{1}\right), \ldots, \varphi\left(e_{n}\right)\right)=t
$$

where $e_{1}, \ldots, e_{n}$ is the standard basis for $\mathbb{R}^{n}$. Then $\mu_{\varphi g, t}=\mu_{\varphi, \Delta_{r}(g)^{-1} t}$ and

$$
\operatorname{GL}(M) \times_{\Delta_{r}} \mathbb{R} \rightarrow L^{(r)}(M), \quad[\varphi, t] \mapsto \mu_{\varphi, t}
$$

is the desired bundle isomorphism.\\
Proposition 10.2.17. Let ( $P, M, G, q, \sigma$ ) be a principal bundle and $\tau: G \times F \rightarrow F$ be a smooth action. We consider the set

$$
C^{\infty}(P, F)^{G}:=\left\{\alpha \in C^{\infty}(P, F):(\forall p \in P)(\forall g \in G) \alpha(p g)=g^{-1} \cdot \alpha(p)\right\}
$$

of smooth $G$-equivariant maps. Then each $\alpha \in C^{\infty}(P, F)^{G}$ defines a smooth section

$$
s_{\alpha}: M \rightarrow P \times_{G} F, \quad q(p) \mapsto[p, \alpha(p)],
$$

and we thus obtain a bijection

$$
\Phi: C^{\infty}(P, F)^{G} \rightarrow \Gamma\left(P \times_{G} F\right), \quad \alpha \mapsto s_{\alpha}
$$

If, in addition, $F$ is a vector space and $\tau$ a linear action, i.e., a representation, then $\Phi$ is linear.

Proof. First, we note that $s_{\alpha}$ is well-defined because

$$
[p g, \alpha(p g)]=[p, g \cdot \alpha(p g)]=[p, \alpha(p)]
$$

for each $p \in P$. Since the map $P \rightarrow P \times_{G} F, p \mapsto[p, \alpha(p)]$ is smooth, the smoothness of $s_{\alpha}$ follows from the fact that $q: P \rightarrow M$ is a submersion (Proposition 8.3.16).\\
$\Phi$ is injective: Suppose that $\Phi(\alpha)=\Phi(\beta)$. Then, for each $p \in P$, we have $[p, \alpha(p)]=[p, \beta(p)]$, and this implies that $\alpha(p)=\beta(p)$ because the action of $G$ on $P$ is free.\\
$\Phi$ is surjective: Let $s: M \rightarrow P \times_{G} F$ be a smooth section. For each $p \in P$, there exists a unique element $\alpha(p)$ with $s(q(p))=[p, \alpha(p)]$. For each $g \in G$, we then have

$$
s(q(p))=s(q(p g))=[p g, \alpha(g p)]=[p, g \cdot \alpha(p g)]
$$

so that $\alpha: P \rightarrow F$ is equivariant. It remains to show that $\alpha$ is smooth. It suffices to show that for each local trivialization $\varphi: P_{U} \rightarrow U \times G$, the map $\alpha \circ \varphi^{-1}$ is smooth. For $x=q(p) \in U$, we have

$$
s(x)=\left[\varphi^{-1}(x, \mathbf{1}), \alpha\left(\varphi^{-1}(x, \mathbf{1})\right]=\widetilde{\varphi}^{-1}\left(x, \alpha\left(\varphi^{-1}(x, \mathbf{1})\right)\right)\right.
$$

Since $\widetilde{\varphi}: P \times_{G} F \rightarrow U \times F$ is a local trivialization, the map $U \rightarrow F, x \mapsto \alpha\left(\varphi^{-1}(x, \mathbf{1})\right)$ is smooth, so that $(x, g) \mapsto \alpha\left(\varphi^{-1}(x, g)\right)=g^{-1} \alpha\left(\varphi^{-1}(x, \mathbf{1})\right)$ is also smooth.

\subsubsection*{Tensor Bundles and Tensor Fields}
Recall the concept of a $V$-valued Pfaffian form $\omega: T M \rightarrow V$ from Definition 8.3.9. The map $\omega$ can also be interpreted as a map

$$
M \rightarrow \coprod_{p \in M} \operatorname{Hom}\left(T_{p}(M), V\right), \quad p \mapsto \omega_{p}
$$

such that $\omega_{p} \in \operatorname{Hom}\left(T_{p}(M), V\right)$. This suggests defining a suitable vector bundle for which $\omega$ is a section. We will do this in greater generality.

Definition 10.2.18. Let $M$ be a smooth manifold of dimension $n$ and view the tangent bundle $T M$ as the vector bundle $\mathrm{GL}(M) \times{ }_{\mathrm{GL}_{n}(\mathbb{R})} \mathbb{R}^{n}$ associated with the frame bundle of $M$ (cf. Example 10.2.9). For any representation of $\mathrm{GL}_{n}(\mathbb{R})$ on a vector space $W$, we obtain the vector bundle $\mathbb{W}:=\operatorname{GL}(M) \times_{\mathrm{GL}_{n}(\mathbb{R})} W \rightarrow M$. To make our notation less clumsy we write $E:=\mathbb{R}^{n}$. If $W$ is a $\mathrm{GL}_{n}(\mathbb{R})$-invariant subspace of the tensor algebra $\mathcal{T}\left(E \oplus E^{*}\right)$, we call $\pi_{\mathbb{W}}: \mathbb{W} \rightarrow M$ a tensor bundle over $M$.

On $W:=T^{r, s}(E):=\left(E^{*}\right)^{\otimes r} \otimes E^{\otimes s}$, we have a representation of $\mathrm{GL}(E)$ by\\
$g \cdot\left(\alpha_{1} \otimes \cdots \otimes \alpha_{r}\right) \otimes\left(\alpha_{1} \otimes \cdots \otimes v_{s}\right)=\left(\left(g^{-1}\right)^{*} \alpha_{1} \otimes \cdots \otimes\left(g^{-1}\right)^{*} \alpha_{r}\right) \otimes\left(g v_{1} \otimes \cdots \otimes g v_{s}\right)$.\\
We write $\mathcal{T}^{r, s}(M)$ for $\mathbb{W}:=\mathrm{GL}(M) \times{ }_{\mathrm{GL}(E)} W$ and call it the bundle of tensors of type ( $r, s$ ). The smooth sections of $\mathcal{T}^{r, s}(M)$ are called the tensor fields of type ( $r, s$ ). The bundle $T^{*}(M):=\mathcal{T}^{1,0}(M)$ is called the cotangent bundle of $M$.

Example 10.2.19. If, in the situation of Definition 10.2.18, $W:= \operatorname{Hom}(E, V)$, then the resulting vector bundle ( $\mathbb{W}, M, W, \pi_{\mathbb{W}}$ ) satisfies

$$
\mathbb{W}_{p} \cong \operatorname{Hom}\left(T_{p}(M), V\right)
$$

because the evaluation map

$$
\operatorname{Iso}\left(E, T_{p}(M)\right) \times \operatorname{Hom}(E, V) \rightarrow \operatorname{Hom}\left(T_{p}(M), V\right), \quad(\varphi, \alpha) \mapsto \alpha \circ \varphi^{-1}
$$

factors through a bijection

$$
\begin{aligned}
\left(\operatorname{Iso}\left(E, T_{p}(M)\right) \times \operatorname{Hom}(E, V)\right) / \operatorname{GL}(E) & \rightarrow \operatorname{Hom}\left(T_{p}(M), V\right),[(\varphi, \alpha)] \\
& \mapsto \alpha \circ \varphi^{-1}
\end{aligned}
$$

If $T M$ is trivial over an open subset $U$, then so is $\mathrm{GL}(M)$, and thus $\mathbb{W}_{U} \cong \mathrm{GL}(M)_{U} \times_{\mathrm{GL}(E)} W$ is also trivial. Therefore, $\Gamma\left(\mathbb{W}_{U}\right) \cong C^{\infty}(U, W)= C^{\infty}(U, \operatorname{Hom}(E, V))$.

Similarly, the elements of $\Omega^{1}(U, V)$ are given by smooth functions $\omega \in C^{\infty}(U \times E, V)$ satisfying $\omega_{p}:=\omega(p, \cdot) \in \operatorname{Hom}(E, V)$. This leads to a linear isomorphism

$$
\Omega^{1}(U, V) \rightarrow \Gamma\left(\mathbb{W}_{U}\right), \quad \omega \mapsto\left(p \mapsto \omega_{p}\right),
$$

which we use to identify these two spaces. Thus the $V$-valued Pfaffian forms are simply the sections of $\mathbb{W}$. In particular, the (scalar valued) Pfaffian forms are the tensor fields of type $(0,1)$.

Example 10.2.20. Let $M$ be a smooth manifold of dimension $n$ and ( $\varphi, U$ ) be a chart on $M$. In $T_{p}(M)$ and $T_{p}(M)^{*}$, we consider the $\varphi$-bases $\left(\left.\frac{\partial}{\partial x_{j}}\right|_{p}\right)_{j=1, \ldots, n}$ and $\left(\mathrm{d} x_{j}(p)\right)_{j=1, \ldots, n}$, respectively, and we define $n^{r+s}$ elements of $\mathcal{T}^{r, s}(M)_{p}$ by

$$
\left(\mathrm{d} x \otimes \frac{\partial}{\partial x}\right)_{(\varrho, \sigma)}(p):=\left.\left.\mathrm{d} x_{\varrho(1)}(p) \otimes \cdots \otimes \mathrm{d} x_{\varrho(r)}(p) \otimes \frac{\partial}{\partial x_{\sigma(1)}}\right|_{p} \otimes \cdots \otimes \frac{\partial}{\partial x_{\sigma(s)}}\right|_{p}
$$

where $\varrho:\{1, \ldots, r\} \rightarrow\{1, \ldots, n\}$ and $\sigma:\{1, \ldots, s\} \rightarrow\{1, \ldots, n\}$ are arbitrary functions. The resulting tensor fields $\left(\mathrm{d} x \otimes \frac{\partial}{\partial x}\right)_{(\varrho, \sigma)}$ are called the $\varphi$-basic fields on $U$. As for vector fields and Pfaffian forms, this terminology is doubly justified. On the one hand, the $\left(d x \otimes \frac{\partial}{\partial x}\right)_{(\varrho, \sigma)}(p)$ form a basis for $\mathcal{T}^{r, s}(M)_{p}$ for every $p \in U$. On the other hand, combining Example 8.4.12 and Example 8.4.10, we see that every local tensor field $T \in \Gamma\left(\mathcal{T}^{r, s}(U)\right)$ is of the form

$$
\left.T\right|_{U}=\sum_{(\varrho, \sigma) \in\{1, \ldots, n\}^{r+s}} c_{(\varrho, \sigma)} \cdot\left(\mathrm{d} x \otimes \frac{\partial}{\partial x}\right)_{(\varrho, \sigma)}
$$

where $c_{(\varrho, \sigma)} \in C^{\infty}(U)$ are uniquely determined by $T$.\\
Fix a smooth manifold $M$. Let $\omega \in \Omega^{1}(M)$ and $X \in \mathcal{V}(M)$. For every $p \in M$, we can apply $\omega(p)$ to $X(p)$ and obtain a number $\omega(p)(X(p))$. As usual, if we consider all $p \in M$ simultaneously, we get a function $\widetilde{\omega}(X): M \rightarrow \mathbb{R}$ which is defined by $\widetilde{\omega}(X)(p)=\omega(p)(X(p))$. This function is smooth, as we easily see in local coordinates (see Exercise 10.2.1).

The form $\omega$ assigns a smooth function to each vector field $X \in \mathcal{V}(M)$. Clearly, this assignment is an $\mathbb{R}$-linear map $\widetilde{\omega}: \mathcal{V}(M) \rightarrow C^{\infty}(M)$. Recall from Remark 8.4.14 that one can multiply vector fields by smooth functions and in this way turn $\mathcal{V}(M)$ into a $C^{\infty}(M)$-module. For $f \in C^{\infty}(M), \omega \in \Omega^{1}(M)$ and $X \in \mathcal{V}(M)$, we calculate:

$$
\widetilde{\omega}(f X)(p)=\omega(p)(f(p) X(p))=f(p)(\omega(p)(X(p)))=(f \cdot \widetilde{\omega}(X))(p) .
$$

This just means that $\widetilde{\omega}$ defines a $C^{\infty}(M)$-linear function $\mathcal{V}(M) \rightarrow C^{\infty}(M)$. We denote the set of all these maps by

$$
\mathcal{V}(M)^{*}:=\operatorname{Hom}_{C^{\infty}(M)}\left(\mathcal{V}(M), C^{\infty}(M)\right)
$$

In algebraic terminology, $\mathcal{V}(M)^{*}$ is the dual $C^{\infty}(M)$-module to $\mathcal{V}(M)$. Similarly as for the case of vector fields (see Theorem 8.4.18), we want to show that the map $\Omega^{1}(M) \rightarrow \mathcal{V}(M)^{*}$ is a bijection.

As before, injectivity is an easy verification:

Lemma 10.2.21. The map $\Omega^{1}(M) \rightarrow \mathcal{V}(M)^{*}, \omega \mapsto \widetilde{\omega}$ is injective.\\
Proof. Let $v \in T_{p}(M)$ and use Remark 8.4.19(iii) to find a smooth vector field $X \in \mathcal{V}(M)$ with $X(p)=v$. If $\widetilde{\omega}=0$, then $\omega_{p}(v)=\widetilde{\omega}(X)(p)=0$, which implies that $\omega_{p}=0$.

As in the case of vector fields, for a given $F \in \mathcal{V}(M)^{*}$, we want to construct a form $\omega$ with $\widetilde{\omega}=F$. For this, we need to have:

$$
\omega(p)(v)=\widetilde{\omega}(X)(p)=F(X)(p)
$$

for all $v \in T_{p} M$ and $X \in \mathcal{V}(M)$ with $X(p)=v$. We can define a map $\omega: M \rightarrow T^{*} M$ with $\pi \circ \omega=\operatorname{id}_{M}$ by (10.4) if we are able to show that $F(X)(p)$ depends only on the value $X(p)$. But, this follows from the following lemma.

Lemma 10.2.22. Let $F \in \mathcal{V}(M)^{*}$ and $X \in \mathcal{V}(M)$ with $X(p)=0$, then $F(X)(p)=0$.

Proof. Step 1: Suppose that $V$ is an open neighborhood of $p$ with $\left.X\right|_{V} \equiv 0$. Remark 8.4.19(i) yields a function $\chi \in C^{\infty}(M)$ such that $\chi(p)=1$ and $\left.\chi\right|_{M \backslash V} \equiv 0$. Then $X=(1-\chi) X$, and hence

$$
F(X)(p)=F((1-\chi) X)(p)=(1-\chi)(p) F(X)(p)=0
$$

Step 2: Let $(\varphi, U)$ be a chart on $M$ and $p \in U$. Then there exists a vector field $\widetilde{X} \in \mathcal{V}(M)$ with compact support contained in $U$ such that $X-\widetilde{X}$ vanishes on a neighborhood of $p$ (Remark 8.4.19(i)). Then Step 1 implies that $F(X)(p)=F(\widetilde{X})(p)$, so that we may w.lo.g. assume that $\operatorname{supp}(X)$ is a compact subset of $U$.\\
Step 3: Using Remark 8.4.19(ii), we write $X(q)=\sum a_{j}(q) X_{j}(q)$ for $q \in U$ with $a_{j} \in C^{\infty}(U)$ and $X_{j} \in \mathcal{V}(U)$, where $X_{j}(q)=\left.\frac{\partial}{\partial x_{j}}\right|_{q}$ for all $q \in U$. Then the functions $a_{j}$ are compactly supported in $U$, hence extend to smooth functions on $M$. Moreover, there exists a smooth function $f: M \rightarrow \mathbb{R}$ with $f(q)=1$ if $X(q) \neq 0$ and $\operatorname{supp}(f)$ is a compact subset of $U$. Then $X= f X=\sum_{j=1}^{n} a_{j}\left(f X_{j}\right)$, and each vector field $f X_{j}$ extends by 0 to a smooth vector field on $M$. Then $X(p)=0$ and $X_{j}(p)=\left.\frac{\partial}{\partial x_{j}}\right|_{p}$ imply $a_{j}(p)=0$ for all $j=1, \ldots, n$. We thus obtain

$$
F(X)(p)=F\left(\sum_{j=1}^{n} a_{j} f X_{j}\right)(p)=\sum_{j=1}^{n} a_{j}(p)\left(F\left(f X_{j}\right)(p)\right)=0
$$

which concludes the proof.\\
Note that in the proof of Lemma 10.2.22 it was decisive that $F$ is not only $\mathbb{R}$-linear, but even $C^{\infty}(M)$-linear.

To show the bijectivity of the map $\Omega^{1}(M) \rightarrow \operatorname{der}\left(C^{\infty}(M)\right)^{*}, \omega \mapsto \widetilde{\omega}$, it only remains to show that the map $\omega: M \rightarrow T^{*} M$ defined by (10.4) is smooth\\
for every $F \in \mathcal{V}(M)^{*}$. For this, we represent the map in local coordinates. From Remark 8.4.19, we get a chart ( $\varphi, U$ ) on $M$, as well as vector fields $X_{j} \in \operatorname{der}\left(C^{\infty}(M)\right)$ with $X_{j}(p)=\left.\frac{\partial}{\partial x_{j}}\right|_{p}$ for all $p \in U$. We obtain

$$
\omega(p)\left(\left.\frac{\partial}{\partial x_{j}}\right|_{p}\right)=F\left(X_{j}\right)(p)=\left(\sum_{k=1}^{n} F\left(X_{k}\right) \mathrm{d} x_{k}\right)(p)\left(\left.\frac{\partial}{\partial x_{j}}\right|_{p}\right)
$$

for all $j=1, \ldots, n$ and $p \in U$. Therefore, we have $\left.\omega\right|_{U}=\left.\sum_{k=1}^{n} F\left(X_{k}\right) \mathrm{d} x_{k}\right|_{U}$, so that $\omega$ is smooth. Thus, we have proved the following theorem:

Theorem 10.2.23. Let $M$ be a smooth manifold. Then the map

$$
\Omega^{1}(M) \rightarrow \mathcal{V}(M)^{*}, \quad \omega \mapsto \widetilde{\omega}
$$

defined by $\widetilde{\omega}(X)(p)=\omega(p)(X(p))$, is a bijection.\\
Just as in Remark 8.4.19, which deals with vector fields, the following remark expresses the corresponding facts for sections of any vector bundle. It applies in particular to Pfaffian forms which are sections of the cotangent bundle $T^{*} M$.

Remark 10.2.24. Let $M$ be a smooth manifold of dimension and ( $\mathbb{V}, M, V, \pi$ ) a vector bundle.\\
(i) Let $C$ be a compact subset of $M$, and let $U \supseteq C$ be an open subset of $M$. Then, for every section $s \in \Gamma\left(\mathbb{V}_{U}\right)$, there exists a global section $\widetilde{s} \in \Gamma \mathbb{V}$ which coincides with $s$ on $C$.\\
(ii) Suppose that $U \subseteq M$ is an open subset for which $\mathbb{V}_{U}$ is trivial. Then there exist sections $s_{1}, \ldots, s_{k} \in \Gamma\left(\mathbb{V}_{U}\right)$ such that $s_{1}(x), \ldots, s_{k}(x)$ form a basis for $\mathbb{V}_{x}$ for each $x \in U$. Then any section $s \in \Gamma\left(\mathbb{V}_{U}\right)$ is of the form

$$
s=\sum_{i=1}^{k} a_{i} \cdot s_{i}
$$

with uniquely determined $a_{i} \in C^{\infty}(U)$.\\
(iii) For every $v \in \mathbb{V}_{x}$, there exists a section $s \in \Gamma \mathbb{V}$ with $s(x)=v$.

\subsubsection*{Lie Derivatives}
Lie derivatives of tensor fields generalize directional derivatives of functions and Lie derivatives of vector fields. They describe the infinitesimal changes of tensor fields under flows of vector fields.

Remark 10.2.25. Let $M$ be a smooth manifold, $\operatorname{GL}(M)$ be the corresponding frame bundle, and $\tau: G \times F \rightarrow F$ a smooth action of the structure group $G:=\mathrm{GL}_{n}(\mathbb{R})$ of $\mathrm{GL}(M)$. Further, let $U \subseteq M$ be open and $\varphi: U \rightarrow M$ be a\\
diffeomorphism onto an open subset $\varphi(U)$ of $M$. Then $\varphi$ lifts to a diffeomor$\operatorname{phism} \operatorname{GL}(\varphi): \operatorname{GL}(U) \rightarrow \operatorname{GL}(\varphi(U))$ via

$$
\operatorname{GL}(\varphi)(\beta)=T(\varphi) \circ \beta \in \operatorname{Iso}\left(\mathbb{R}^{n}, T_{\varphi(x)}(M)\right)
$$

for $\beta \in \operatorname{Iso}\left(\mathbb{R}^{n}, T_{x}(M)\right) \subseteq \operatorname{GL}(U)$.\\
Applying this construction to a local flow $\Phi_{t}^{X}$, we see that a vector field $X \in \mathcal{V}(M)$ automatically lifts to a vector field $\widetilde{X} \in \mathcal{V}(\operatorname{GL}(M))$ defined by

$$
\widetilde{X}(\beta)=\left.\frac{d}{d t}\right|_{t=0} \operatorname{GL}\left(\Phi_{t}^{X}\right)(\beta)
$$

$\mathrm{GL}(\varphi)$ also enables us to let $\varphi$ act on sections. More precisely, suppose that $\alpha \in C^{\infty}(\mathrm{GL}(\varphi(U)), F)^{G}$ represents a section $s_{\alpha} \in \Gamma\left(\mathrm{GL}(\varphi(U)) \times_{G} F\right)$ as explained in Proposition 10.2.17. Then we set

$$
\varphi^{*} \alpha:=\alpha \circ \operatorname{GL}(\varphi) \in C^{\infty}(\operatorname{GL}(U), F)
$$

and note that

$$
\begin{aligned}
\left(\varphi^{*} \alpha\right)(\beta g) & =\alpha(T(\varphi) \circ(\beta g))=\alpha((T(\varphi) \circ \beta) g) \\
& =g^{-1} \cdot \alpha(T(\varphi) \circ \beta)=g^{-1} \cdot\left(\varphi^{*} \alpha\right)(\beta) .
\end{aligned}
$$

Therefore, $\varphi^{*} \alpha \in C^{\infty}(\operatorname{GL}(U), F)^{G}$ and $\varphi^{*} \alpha$ represents a section $\varphi^{*} s_{\alpha}:= s_{\varphi^{*} \alpha}$.

Definition 10.2.26. Let $M$ be a manifold, $G=\mathrm{GL}_{n}(\mathbb{R})$, and $U \subseteq M$ open. Given a linear representation $\tau: G \times F \rightarrow F$ and a vector field $X \in \mathcal{V}(M)$, the Lie derivative $\mathcal{L}_{X} \alpha$ of a section $\alpha \in C^{\infty}(\mathrm{GL}(U), F)^{G}$ is defined via the pull-back $\left(\Phi_{t}^{X}\right)^{*} \alpha$ of $\alpha$ by the local flow of $X$ (Remark 10.2.25). We set

$$
\mathcal{L}_{X}(\alpha)_{p}:=\lim _{t \rightarrow 0} \frac{1}{t}\left(\left(\Phi_{t}^{X}\right)^{*} \alpha-\alpha\right)_{p}
$$

whence $\mathcal{L}_{X}(\alpha)$ is again a section in $C^{\infty}(\operatorname{GL}(U), F)^{G}$.\\
Remark 10.2.27. It is possible to reformulate the Lie derivative $\mathcal{L}_{X}(Y)$ of a vector field $Y \in \mathcal{V}(M)$ with respect to $X \in \mathcal{V}(M)$ in an analogous way:

$$
\mathcal{L}_{X}(Y)(p)=\lim _{t \rightarrow 0} \frac{1}{t}\left(\left(\Phi_{t}^{X}\right)^{*}(Y)-Y\right)(p)
$$

Here we define the pull-back $\varphi^{*}(X) \in \mathcal{V}(M)$ of a vector-field $X \in \mathcal{V}(N)$ by a diffeomorphism $\varphi: M \rightarrow N$ via the commutative diagram\\
\includegraphics[max width=\textwidth, center]{2025_10_18_3eaf04e77f3cb364fce5g-394}\\
so that

$$
\varphi^{*}(X)=T\left(\varphi^{-1}\right) \circ X \circ \varphi .
$$

In order to verify that this is compatible with Remark 10.2.25, observe that, according to Example 10.2.12, the tangent map

$$
T(\varphi): \mathrm{GL}(M) \times_{\mathrm{GL}_{m}(\mathbb{R})} \mathbb{R}^{m} \rightarrow \mathrm{GL}(N) \times_{\mathrm{GL}_{n}(\mathbb{R})} \mathbb{R}^{n}
$$

of a smooth map $\varphi: M \rightarrow N$ in the realization of the tangent bundles as associated bundles is given by

$$
T(\varphi)[p, v]=\left[q_{0}, w\right]
$$

where $p \in \operatorname{Iso}\left(\mathbb{R}^{m}, T_{x}(M)\right), v \in \mathbb{R}^{m}$, and $q_{0} \in \operatorname{Iso}\left(\mathbb{R}^{n}, T_{\varphi(x)}(N)\right), w \in \mathbb{R}^{n}$ are such that $q_{0}(w)=(T(\varphi) \circ p)(v)=\mathrm{GL}(\varphi)(v)$. If now $\varphi: M \rightarrow N$ is a diffeomorphism and $\alpha \in C^{\infty}\left(\mathrm{GL}(N), \mathbb{R}^{n}\right)^{\mathrm{GL}_{n}(\mathbb{R})}$, then Proposition 10.2.17 yields

$$
\begin{aligned}
T(\varphi) \circ s_{\varphi^{*} \alpha}(x) & =T(\varphi)\left(\varphi^{*} \alpha(p)\right)=T(\varphi)[p, \alpha(T(\varphi) \circ p)] \\
& =[T(\varphi) \circ p, \alpha(T(\varphi) \circ p)]=s_{\alpha}(\varphi(x))
\end{aligned}
$$

since $n=m$ and $T(\varphi) \circ p \in \operatorname{Iso}\left(\mathbb{R}^{n}, T_{\varphi(x)}(N)\right)$.\\
Proposition 10.2.28. Let $M$ be a differentiable manifold of dimension $m$ and $\mu: E \times F \rightarrow W$ a bilinear equivariant map of $\mathrm{GL}_{m}(\mathbb{R})$-spaces. If $\alpha, \alpha_{1}, \alpha_{2} \in C^{\infty}(\mathrm{GL}(M), E)^{\mathrm{GL}_{m}(\mathbb{R})}, \beta \in C^{\infty}(\mathrm{GL}(M), F)^{\mathrm{GL}_{m}(\mathbb{R})}$, and $X \in \mathcal{V}(M)$, then\\
(i) $\mathcal{L}_{X}\left(\alpha_{1}+\alpha_{2}\right)=\mathcal{L}_{X}\left(\alpha_{1}\right)+\mathcal{L}_{X}\left(\alpha_{2}\right)$;\\
(ii) $\mathcal{L}_{X}(\mu(\alpha, \beta))=\mu\left(\mathcal{L}_{X}(\alpha), \beta\right)+\mu\left(\alpha, \mathcal{L}_{X}(\beta)\right)$.

In particular, setting $E=\mathbb{R}$ with the trivial action and $F=\mathbb{R}^{m}$ with the natural action, we find $\mathcal{L}_{X}(f Y)=(X f) Y+f \mathcal{L}_{X}(Y)$ for $f \in C^{\infty}(M)$ and $Y \in \mathcal{V}(M)$.

Proof. We leave the verification of (i) as an easy exercise to the reader. For (ii), we calculate

$$
\begin{aligned}
\mathcal{L}_{X}(\mu(\alpha, \beta))(p)= & \lim _{t \rightarrow 0} \frac{1}{t}\left(\left(\Phi_{t}^{X}\right)^{*} \mu(\alpha, \beta)(p)-\mu(\alpha, \beta)(p)\right) \\
= & \lim _{t \rightarrow 0} \frac{1}{t}\left(\mu(\alpha, \beta)\left(T\left(\Phi_{t}^{X}\right) \circ p\right)-\mu(\alpha, \beta)(p)\right) \\
= & \lim _{t \rightarrow 0} \frac{1}{t}\left(\mu\left(\alpha\left(T\left(\Phi_{t}^{X}\right) \circ p\right), \beta\left(T\left(\Phi_{t}^{X}\right) \circ p\right)\right)\right. \\
& \quad-\mu(\alpha(p), \beta(p))) \\
= & \lim _{t \rightarrow 0} \frac{1}{t}\left(\mu\left(\alpha\left(T\left(\Phi_{t}^{X}\right) \circ p\right)-\alpha(p), \beta\left(T\left(\Phi_{t}^{X}\right) \circ p\right)\right)\right. \\
& \left.\quad+\mu\left(\alpha(p), \beta\left(T\left(\Phi_{t}^{X}\right) \circ p\right)-\beta(p)\right)\right) \\
= & \mu\left(\mathcal{L}_{X}(\alpha), \beta\right)(p)+\mu\left(\alpha, \mathcal{L}_{X}(\beta)\right)(p)
\end{aligned}
$$

Remark 10.2.29. One can use Proposition 10.2.28, applied to the natural pairing of $\mathbb{R}^{m}$ and $\left(\mathbb{R}^{m}\right)^{*}$, to calculate the Lie derivatives of a Pfaffian form in local coordinates: For a vector field of the form $X=\sum_{k} a_{k} \frac{\partial}{\partial x_{k}}$, we find

$$
\begin{aligned}
\left(\mathcal{L}_{X}\left(\mathrm{~d} x_{j}\right)\right)\left(\frac{\partial}{\partial x_{i}}\right) & =-\mathrm{d} x_{j}\left(\mathcal{L}_{X} \frac{\partial}{\partial x_{i}}\right)=-\mathrm{d} x_{j}\left(\left[X, \frac{\partial}{\partial x_{i}}\right]\right) \\
& =\mathrm{d} x_{j}\left(\sum_{k} \frac{\partial a_{k}}{\partial x_{i}} \frac{\partial}{\partial x_{k}}\right)=\frac{\partial a_{j}}{\partial x_{i}}
\end{aligned}
$$

so that

$$
\mathcal{L}_{X}\left(\mathrm{~d} x_{j}\right)=\sum_{i} \frac{\partial a_{j}}{\partial x_{i}} \mathrm{~d} x_{i}
$$

Example 10.2.30. Let $M$ and $N$ be smooth manifolds and $\varphi: M \rightarrow N$ a diffeomorphism. Then the linear isomorphisms $T_{p}(\varphi): T_{p}(M) \rightarrow T_{\varphi(p)}(N)$ induce linear isomorphisms

$$
\mathcal{T}_{p}^{r, s}(\varphi): \mathcal{T}^{r, s}(M)_{p} \rightarrow \mathcal{T}^{r, s}(N)_{f(p)}, \quad \alpha \mapsto \varphi_{*} \alpha,
$$

where

$$
\begin{aligned}
& \left(\varphi_{*} \alpha\right)_{x}\left(\omega_{1} \otimes \cdots \otimes \omega_{r} \otimes v_{1} \otimes \cdots \otimes v_{s}\right) \\
& =\alpha_{\varphi(x)}\left(\omega_{1} \circ T_{x}(\varphi)^{-1}, \ldots, \omega_{r} \circ T_{x}(\varphi)^{-1}, T_{x}(\varphi) v_{1}, \ldots, T_{x}(\varphi) v_{s}\right)
\end{aligned}
$$

for $x \in M, v_{j} \in T_{x}(M)$ and $\omega_{j} \in T_{x}^{*}(M)$. Together they form a bundle isomorphism $\mathcal{T}^{r, s}(\varphi): \mathcal{T}^{r, s}(M) \rightarrow \mathcal{T}^{r, s}(N)$. This is easily read off from the isomorphisms\\
$\mathcal{T}^{r, s}(M) \cong \mathrm{GL}(M) \times_{\mathrm{GL}_{n}(\mathbb{R})} \mathcal{T}^{r, s}\left(\mathbb{R}^{n}\right), \quad \mathcal{T}^{r, s}(N) \cong \mathrm{GL}(N) \times_{\mathrm{GL}_{n}(\mathbb{R})} \mathcal{T}^{r, s}\left(\mathbb{R}^{n}\right)$\\
because in this picture

$$
\mathcal{T}^{r, s}(\varphi)([\alpha, v])=\left[T_{p}(\varphi) \circ \alpha, v\right], \quad \alpha \in \mathrm{GL}(M)_{p}=\operatorname{Iso}\left(\mathbb{R}^{n}, T_{p}(M)\right)
$$

In this picture, the pull-back $\varphi^{*} T \in \Gamma\left(\mathcal{T}^{r, s}(M)\right)$ of $T \in \Gamma\left(\mathcal{T}^{r, s}(N)\right)$ is obtained via the following commutative diagram:\\
\includegraphics[max width=\textwidth, center]{2025_10_18_3eaf04e77f3cb364fce5g-396}

The Lie derivative $\mathcal{L}_{X}(T)$ of $T$ with respect to a vector field $X \in \mathcal{V}(M)$ is then given by

$$
\mathcal{L}_{X}(T)(p):=\lim _{t \rightarrow 0} \frac{1}{t}\left(\left(\Phi_{t}^{X}\right)^{*} T-T\right)(p)
$$

for all $p \in M$, where $\Phi_{t}^{X}$ is the local flow of $X$.

Corollary 10.2.31. Let $M$ be a manifold and $X \in \mathcal{V}(M)$. Then for two tensor fields $T$ and $T^{\prime}$ on $M$ the following equality holds

$$
\mathcal{L}_{X}\left(T \otimes T^{\prime}\right)=\mathcal{L}_{X} T \otimes T^{\prime}+T \otimes \mathcal{L}_{X} T^{\prime}
$$

Proof. Apply Proposition 10.2.28(ii) to the bilinear map $\left(T, T^{\prime}\right) \mapsto T \otimes T^{\prime}$.\\
Remark 10.2.32. While the Lie derivative is a good notion of differentiation of a tensor field with respect to a vector field, it does not lead to a notion of differentiation of a tensor field with respect to a tangent vector at a point. The reason for this is that the value of the Lie derivative of a tensor field at a point depends on the value of the vector field not only at that point, but in a neighborhood. It is easy to see that from the algebraic point of view this is due to the fact that the map $X \mapsto \mathcal{L}_{X}$ is itself a differential operator and not a $C^{\infty}(M)$-linear map. For instance, if $f \in C^{\infty}(M)$ and $\alpha \in \Omega^{1}(M)$, we have (cf. Proposition 10.2.28)

$$
\begin{aligned}
\left(\mathcal{L}_{f X} \alpha\right)(Y) & =f X(\alpha(Y))-\alpha\left(\mathcal{L}_{f X} Y\right)=f X(\alpha(Y))-\alpha([f X, Y]) \\
& =f X(\alpha(Y))-\alpha(f[X, Y]-(Y f) X) \\
& =f\left(X(\alpha(Y))-\alpha\left(\mathcal{L}_{X} Y\right)\right)+(Y f) \cdot \alpha(X) \\
& =f \mathcal{L}_{X} \alpha(Y)+\alpha(X) \mathrm{d} f(Y) .
\end{aligned}
$$

Thus we have

$$
\mathcal{L}_{f X} \alpha=f \mathcal{L}_{X} \alpha+\alpha(X) \mathrm{d} f
$$

\subsubsection*{Exercises for Section 10.2}
Exercise 10.2.1. Show that the function $\widetilde{\omega}(X)$ is smooth for $\omega \in \Omega^{1}(M)$ and $X \in \mathcal{V}(M)$.

Exercise 10.2.2. Show that the bidual $C^{\infty}(M)$-module $\left(\mathcal{V}(M)^{*}\right)^{*}$ of $\mathcal{V}(M)$ is isomorphic to $\mathcal{V}(M)$ as a $C^{\infty}(M)$-module.

Exercise 10.2.3. Compute how the coordinates of $\omega \in T_{p}^{*}(M)$ change when one changes charts for $M$ and uses the associated charts for $T^{*}(M)$ constructed from $T^{*}(M) \cong \mathrm{GL}(M) \rtimes_{\mathrm{GL}_{n}(\mathbb{R})}\left(\mathbb{R}^{n}\right)^{*}$ and local trivializations of $\mathrm{GL}(M)$.

Exercise 10.2.4. We denote by $\left(\mathrm{d} x \otimes \frac{\partial}{\partial x}\right)_{(\varrho, \sigma)}(p)$ the element of $\mathcal{T}^{r, s}\left(T_{p} M\right)$ which is determined by $\varrho$ and $\sigma$. Show that the

$$
\left.\left.\mathrm{d} x_{\varrho(1)}(p) \otimes \cdots \otimes \mathrm{d} x_{\varrho(r)}(p) \otimes \frac{\partial}{\partial x_{\sigma(1)}}\right|_{p} \otimes \cdots \otimes \frac{\partial}{\partial x_{\sigma(s)}}\right|_{p}
$$

form a basis for $\mathcal{T}^{r, s}\left(T_{p} M\right)$.

Exercise 10.2.5. Compute how the coordinates of $t \in \mathcal{T}^{r, s}\left(T_{p} M\right)$ change when one changes charts for $M$ and uses the associated charts for $\mathcal{T}^{r, s}(M)$ constructed from $\mathcal{T}^{r, s}(M) \cong \mathrm{GL}(M) \rtimes_{\mathrm{GL}_{n}(\mathbb{R})} T^{r, s}\left(\mathbb{R}^{n}\right)$ and local trivializations of GL $(M)$.

Exercise 10.2.6. Generalize Theorem 10.2.23 to general tensor bundles $\mathcal{T}^{r, s}(M)$. More precisely, show that the space of sections $\Gamma\left(\mathcal{T}^{r, s}(M)\right)$ is naturally isomorphic to $\left(\mathcal{V}(M)^{*}\right)^{\otimes r} \otimes \mathcal{V}(M)^{\otimes s}$ as a $C^{\infty}(M)$-module.

Exercise 10.2.7. In Exercise 10.2.5, it was calculated how the coordinates vary under a change of coordinates. Conversely, show that a family $\left(a_{(\varrho, \sigma)}^{i}\right)_{i \in I}$ of coordinate functions transforming as in Exercise 10.2.5 define a tensor field which has these functions as coordinate functions ("a tensor is something that transforms like a tensor").

Exercise 10.2.8. Let $G$ be a group and $\tau: G \times F \rightarrow F,(g, f) \mapsto g \cdot f$ be an action of $G$ on $F$. Show that:\\
(a) The map $\tau: G \times F \rightarrow F$ factors through a bijection $(G \times F) / G$, where the $G$-action on $G \times F$ is defined by $g(h, f):=\left(h g^{-1}, g \cdot f\right)$.\\
(b) If $V$ is a vector space, then the evaluation map ev: $\mathrm{GL}(V) \times V \rightarrow V$ factors through a bijection ( $\operatorname{GL}(V) \times V) / \operatorname{GL}(V)$, where the GL $(V)$-action on $\operatorname{GL}(V) \times V$ is defined by $g(h, f):=\left(h g^{-1}, g \cdot f\right)$.\\
(c) If $V$ is an $n$-dimensional vector space, then the evaluation map

$$
\text { ev: } \operatorname{Iso}\left(\mathbb{R}^{n}, V\right) \times \mathbb{R}^{n} \rightarrow V, \quad(\varphi, v) \mapsto \varphi(v)
$$

factors through a bijection $\left(\mathrm{GL}_{n}(\mathbb{R}) \times \mathbb{R}^{n}\right) / \mathrm{GL}_{n}(\mathbb{R})$, where the $\mathrm{GL}_{n}(\mathbb{R})$ action on $\operatorname{Iso}\left(\mathbb{R}^{n}, V\right) \times \mathbb{R}^{n}$ is defined by $g(h, f):=\left(h g^{-1}, g \cdot f\right)$. Give an interpretation of the fibers of ev in terms of bases of $V$ and coordinates of a vector with respect to a basis.

Exercise 10.2.9. Let $M:=\mathbb{R}^{n}$. For a matrix $A \in M_{n}(\mathbb{R})$ and $b \in \mathbb{R}^{n}$, we consider the affine vector field

$$
X_{A, b}(x):=A x+b
$$

(1) Calculate the maximal flow $\Phi: \mathbb{R} \times \mathbb{R}^{n} \rightarrow \mathbb{R}^{n}$ of this vector field.\\
(2) Let

$$
\mathfrak{a f f}_{n}(\mathbb{R}):=\left(\begin{array}{cc}
\mathfrak{g} \mathfrak{l}_{n}(\mathbb{R}) & \mathbb{R}^{n} \\
0 & 0
\end{array}\right) \subseteq \mathfrak{g} \mathfrak{l}_{n+1}(\mathbb{R})
$$

be the affine Lie algebra on $\mathbb{R}^{n}$, realized as a Lie subalgebra of $\mathfrak{g l}_{n+1}(\mathbb{R})$, endowed with the commutator bracket. Show that the map

$$
\varphi: \mathfrak{a f f}_{n}(\mathbb{R}) \rightarrow \mathcal{V}\left(\mathbb{R}^{n}\right), \quad\left(\begin{array}{cc}
A & b \\
0 & 0
\end{array}\right) \mapsto-X_{A, b}
$$

is a homomorphism of Lie algebras.

\subsection*{Integration on Manifolds}
Integration on manifolds is built up from integration on coordinate patches. Since the integrals on coordinate patches are not invariant under coordinate changes, the naive approach to integrate functions does not work. What can be integrated are sections of bundles which transform in the same way as integrals, i.e., via the absolute value of the Jacobi determinants of the coordinate changes. There are two choices: One either uses densities which have precisely this transformation behavior or one uses differential forms of degree $\operatorname{dim} M$. In the latter case, one has to make sure that all coordinate changes have positive Jacobi determinants. This leads to the concept of an orientable manifold.

\subsubsection*{Differential Forms}
Differential forms are tensor fields with additional symmetry conditions. They are of central importance in the description of the topology of a manifold. Moreover, they can be used to build an integration theory for which one has generalizations of the classical Stokes Theorem.

Let $M$ be a smooth manifold of dimension $n$ and $\operatorname{GL}(M)$ its frame bundle. Realize the tangent bundle $T M$ as the associated bundle $\mathrm{GL}(M) \times{ }_{\mathrm{GL}_{n}(\mathbb{R})} \mathbb{R}^{n}$ (cf. Definition 10.2.18)

Definition 10.3.1. For $k \in \mathbb{N}_{0}$, we consider the $\mathrm{GL}_{n}(\mathbb{R})$-module Alt ${ }^{k}\left(\mathbb{R}^{n}, \mathbb{R}\right)$ of alternating $k$-linear forms on $E$, on which $\mathrm{GL}_{n}(\mathbb{R})$ acts by

$$
(g \cdot \alpha)\left(v_{1}, \ldots, v_{k}\right):=\alpha\left(g^{-1} v_{1}, \ldots, g^{-1} v_{k}\right) .
$$

The resulting associated bundle

$$
\operatorname{Alt}^{k}(T M):=\operatorname{GL}(M) \times_{\operatorname{GL}_{n}(\mathbb{R})} \operatorname{Alt}^{k}\left(\mathbb{R}^{n}, \mathbb{R}\right)
$$

is called the $k$-form bundle over $M$. Its sections are called alternating $k$-forms, or simply $k$-forms on $M$. One also speaks of differential forms of degree $k$. The space of differential forms of degree $k$ on $M$ will be denoted by $\Omega^{k}(M)$. Note that there are no non-zero differential forms of degree greater than the dimension of $M$. We set $\Omega(M)=\bigoplus_{k=0}^{\operatorname{dim} M} \Omega^{k}(M)=\bigoplus_{k=0}^{\infty} \Omega^{k}(M)$.

Example 10.3.2. If $n=\operatorname{dim} M$, then $\operatorname{Alt}^{n}\left(\mathbb{R}^{n}, \mathbb{R}\right) \cong \mathbb{R}$ and the $\mathrm{GL}_{n}(\mathbb{R})$ representation on $\operatorname{Alt}^{n}\left(\mathbb{R}^{n}, \mathbb{R}\right)$ is the inverse of the determinant. Thus the transition functions of the line bundle $\mathrm{Alt}^{n}(T M)$ are the inverses of the Jacobi determinants of the coordinate changes.

Remark 10.3.3. (cf. Exercise 10.2.6) Theorem 10.2.23 generalizes from $T^{*}(M)=\operatorname{Alt}^{1}(T M)$ to $\operatorname{Alt}^{k}(T M)$. More precisely, the space of differential forms $\Omega^{k}(M)$ is naturally isomorphic to the space of alternating $C^{\infty}(M)$ multilinear maps $\mathcal{V}(M)^{k} \rightarrow C^{\infty}(M)$.

Definition 10.3.4. For each vector space $E$ we have a natural bilinear product

$$
\wedge: \operatorname{Alt}^{p}(E, \mathbb{R}) \times \operatorname{Alt}^{q}(E, \mathbb{R}) \rightarrow \operatorname{Alt}^{p+q}(E, \mathbb{R})
$$

called the exterior or wedge product, defined by

$$
\begin{aligned}
& (\alpha \wedge \beta)\left(v_{1}, \ldots, v_{p+q}\right) \\
& =\frac{1}{p!q!} \sum_{\sigma \in S_{p+q}} \operatorname{sgn}(\sigma) \alpha\left(v_{\sigma(1)}, \ldots, v_{\sigma(p)}\right) \beta\left(v_{\sigma(p+1)}, \ldots, X_{\sigma(p+q)}\right)
\end{aligned}
$$

Taking into account that $\alpha$ and $\beta$ are alternating, this product can also be written with $\binom{p+q}{p}$ summands instead of $(p+q)$ !:

$$
\begin{aligned}
& (\alpha \wedge \beta)\left(v_{1}, \ldots, v_{p+q}\right) \\
& =\sum_{\sigma \in \operatorname{Sh}(p, q)} \operatorname{sgn}(\sigma) \alpha\left(v_{\sigma(1)}, \ldots, v_{\sigma(p)}\right) \beta\left(v_{\sigma(p+1)}, \ldots, v_{\sigma(p+q)}\right)
\end{aligned}
$$

where $\operatorname{Sh}(p, q)$ denotes the set of all $(p, q)$-shuffles in $S_{p+q}$, i.e., all permutations with

$$
\sigma(1)<\cdots<\sigma(p) \quad \text { and } \quad \sigma(p+1)<\cdots<\sigma(p+q) .
$$

This product induces a bilinear bundle map

$$
\wedge: \operatorname{Alt}^{p}(T M) \times \operatorname{Alt}^{q}(T M) \rightarrow \operatorname{Alt}^{p+q}(T M), \quad([\psi, \alpha],[\psi, \beta]) \mapsto[\psi, \alpha \wedge \beta]
$$

and hence a $C^{\infty}(M)$-bilinear product on the spaces of smooth sections

$$
\Omega^{p}(M) \times \Omega^{q}(M) \rightarrow \Omega^{p+q}(M), \quad(\alpha, \beta) \mapsto \alpha \wedge \beta,
$$

where $(\alpha \wedge \beta)_{p}=\alpha_{p} \wedge \beta_{p}$ for each $p \in M$.\\
Proposition 10.3.5. If $\alpha, \beta$ are differential forms on a manifold $M$, then

$$
\mathcal{L}_{X}(\alpha \wedge \beta)=\mathcal{L}_{X} \alpha \wedge \beta+\alpha \wedge \mathcal{L}_{X} \beta
$$

for each vector field $X \in \mathcal{V}(M)$.\\
Proof. Take $\mu$ in Proposition 10.2.28(ii) to be $(\alpha, \beta) \mapsto \alpha \wedge \beta$.\\
Remark 10.3.6. Let $X$ be a smooth vector field on a manifold $M$. From Definition 10.2.26, it is easy to deduce that for any section $\alpha$ of a vector bundle $\mathbb{V}$ associated with the frame bundle $\mathrm{GL}(M)$ of $M$, the equality $\mathcal{L}_{X} A=0$ implies that $A$ is invariant under the flow $\Phi_{t}^{X}$ of $X$, i.e., $\left(\Phi_{t}^{X}\right)^{*} \alpha=\alpha$. Hence we say such a section is invariant under $X$ if $\mathcal{L}_{X} A=0$.

Remark 10.3.7. Proposition 10.3.5 can be used to compute $\mathcal{L}_{X}$ on differential forms. We have

$$
\left(\mathcal{L}_{X} \omega\right)(Y)=X \omega(Y)-\omega([X, Y])
$$

for a 1 -form $\omega$, and

$$
\left(\mathcal{L}_{X} \alpha\right)\left(X_{1}, \ldots, X_{k}\right)=X \alpha\left(X_{1}, \ldots, X_{k}\right)-\sum_{i=1}^{k} \alpha\left(X_{1}, \ldots,\left[X, X_{i}\right], \ldots, X_{k}\right)
$$

for forms $\alpha$ of degree $k$ (cf. Lemma 7.5.8). This formula follows directly from Proposition 10.2.28 and $\mathcal{L}_{X} Y=[X, Y]$ for $X, Y \in \mathcal{V}(M)$.

Definition 10.3.8. Let $\alpha$ be a differential form of degree $k$. Then we define the exterior derivative $\mathrm{d} \alpha$ of $\alpha$ to be the differential form of degree $k+1$ given by the formula

$$
\begin{aligned}
(\mathrm{d} \alpha)\left(X_{0}, \ldots, X_{k}\right):= & \sum_{i=0}^{k}(-1)^{i} X_{i} \alpha\left(X_{0}, \ldots, \widehat{X}_{i}, \ldots, X_{k}\right)+ \\
& \sum_{0 \leq i<j \leq k}(-1)^{i+j} \alpha\left(\left[X_{i}, X_{j}\right], X_{0}, \ldots, \widehat{X}_{i}, \ldots, \widehat{X}_{j}, \ldots, X_{k}\right),
\end{aligned}
$$

where a hat over a symbol means that the symbol does not occur (cf. Definition 7.5.2).

Clearly, $\mathrm{d} \alpha$ defines a $(k+1)$-linear map $\mathcal{V}(M)^{k+1} \rightarrow C^{\infty}(M)$, and it is easily seen to be alternating (cf. Definition 7.5.2). We claim that it is also $C^{\infty}(M)$-linear. Since it is alternating, it suffices to verify the $C^{\infty}(M)$ linearity in $X_{0}$. The terms for which the $C^{\infty}(M)$-linearity is not obvious are

$$
\begin{aligned}
& \sum_{i=1}^{k}(-1)^{i} X_{i} \alpha\left(X_{0}, \ldots, \widehat{X}_{i}, \ldots, X_{k}\right) \\
& \quad+\sum_{0<i \leq k}(-1)^{i} \alpha\left(\left[X_{0}, X_{i}\right], X_{1}, \ldots, \widehat{X}_{i}, \ldots, X_{k}\right)
\end{aligned}
$$

But the relations

$$
X(f g)=X f \cdot g+f \cdot X g \quad \text { and } \quad[X, f Y]=(X f) \cdot Y+f[X, Y]
$$

imply that $\mathrm{d} \alpha$ is, in fact, $C^{\infty}(M)$-linear in each argument.\\
Example 10.3.9. If $\alpha$ is a form of degree 0 , i.e., a function $f$, then the definition of the exterior derivative gives $\mathrm{d} f(X)=X f$. If $\alpha$ is of degree 1 , then

$$
\mathrm{d} \alpha(X, Y)=X \alpha(Y)-Y \alpha(X)-\alpha([X, Y]) .
$$

Definition 10.3.10. If $X$ is a vector field and $\alpha$ is a differential form of degree $k$ on a manifold, we define the interior product or contraction of $\alpha$ with respect to $X$ to be a form of degree $k-1$ given by

$$
\left(i_{X} \alpha\right)\left(X_{1}, \ldots, X_{k-1}\right):=\alpha\left(X, X_{1}, \ldots, X_{k-1}\right)
$$

for $k>0$. For $k=0$, i.e., for functions, we set $i_{X} f=0$.\\
Proposition 10.3.11. For forms $\alpha, \beta$ of degree $k, m$ on a manifold

$$
i_{X}(\alpha \wedge \beta)=i_{X} \alpha \wedge \beta+(-1)^{k} \alpha \wedge i_{X} \beta
$$

Proof. This is a simple consequence of the definitions.\\
Proposition 10.3.12. The exterior derivative d on a manifold has the following properties.\\
(i) $\mathrm{d}: \Omega(M) \rightarrow \Omega(M)$ is linear.\\
(ii) Any vector field $X$ satisfies the Cartan formula

$$
\mathrm{d} \circ i_{X}+i_{X} \circ \mathrm{~d}=\mathcal{L}_{X} \quad \text { on } \quad \Omega(M)
$$

(iii) If $\alpha$ and $\beta$ are differential forms of degree $p$ and $q$, respectively, then

$$
\mathrm{d}(\alpha \wedge \beta)=\mathrm{d} \alpha \wedge \beta+(-1)^{p} \alpha \wedge \mathrm{~d} \beta
$$

(iv) $\mathrm{d} \circ \mathrm{d}=0$.

Proof. (i) This is obvious from the definition of d .\\
(ii) follows from Lemma 7.5.9.\\
(iii) This is proved by induction on $p+q$. If one of $p, q$ is 0 , then the assertion is easily seen by direct calculation. So let us assume that $p+q>0$. In order to prove the equality of the differential forms on both sides, it is enough to show that they are the same after applying $i_{X}$ for any $X$. Now we may use (ii) and the induction hypothesis to complete the proof. More precisely,

$$
\begin{aligned}
& \quad i_{X} \mathrm{~d}(\alpha \wedge \beta) \\
& =\mathcal{L}_{X}(\alpha \wedge \beta)-\mathrm{d}\left(i_{X}(\alpha \wedge \beta)\right) \\
& \stackrel{(10.6)}{=} \mathcal{L}_{X} \alpha \wedge \beta+\alpha \wedge \mathcal{L}_{X} \beta-\mathrm{d}\left(i_{X} \alpha \wedge \beta+(-1)^{p} \alpha \wedge i_{X} \beta\right) \\
& \stackrel{\text { ind. }}{=} i_{X}(\mathrm{~d} \alpha) \wedge \beta+(-1)^{p}\left(i_{X} \alpha\right) \wedge \mathrm{d} \beta+(-1)^{p+1} \mathrm{~d} \alpha \wedge\left(i_{X} \beta\right)+\alpha \wedge\left(i_{X} \mathrm{~d} \beta\right) \\
& =i_{X}(\mathrm{~d} \alpha \wedge \beta)+(-1)^{p} i_{X}(\alpha \wedge \mathrm{~d} \beta)=i_{X}\left(\mathrm{~d} \alpha \wedge \beta+(-1)^{p} \alpha \wedge \mathrm{~d} \beta\right)
\end{aligned}
$$

(iv) follows from Proposition 7.5.12.

Remark 10.3.13. Using the Cartan Formula, we find

$$
\begin{aligned}
\mathcal{L}_{f X} \alpha & =i_{f X} \mathrm{~d} \alpha+\mathrm{d}\left(i_{f X} \alpha\right)=f i_{X}(\mathrm{~d} \alpha)+\mathrm{d}\left(f i_{X} \alpha\right) \\
& =f i_{X}(\mathrm{~d} \alpha)+f \mathrm{~d}\left(i_{X} \alpha\right)+\mathrm{d} f \wedge i_{X} \alpha=f \mathcal{L}_{X} \alpha+\mathrm{d} f \wedge i_{X} \alpha
\end{aligned}
$$

i.e.,

$$
\mathcal{L}_{f X} \alpha=f \mathcal{L}_{X} \alpha+\mathrm{d} f \wedge i_{X} \alpha
$$

If $\alpha$ is of degree $n=\operatorname{dim}_{\mathbb{R}}(M)$, we have

$$
0=i_{X}(\mathrm{~d} f \wedge \alpha)=i_{X}(\mathrm{~d} f) \wedge \alpha-\mathrm{d} f \wedge i_{X} \alpha=(X f) \alpha-\mathrm{d} f \wedge i_{X} \alpha
$$

so that $\mathcal{L}_{f X}(\alpha)=f \mathcal{L}_{X} \alpha+(X f) \cdot \alpha$.\\
Definition 10.3.14. A differential form $\omega \in \Omega^{k}(M)$ satisfying $\mathrm{d} \omega=0$ is called closed. If $\omega=\mathrm{d} \nu$ for some $\nu \in \Omega^{k-1}(M)$, then $\omega$ is called an exact form. Then $\mathrm{d}^{2}=0$ implies that each exact form is closed. Putting $\Omega^{-1}(M):=\{0\}$, the quotient space

$$
H_{\mathrm{dR}}^{k}(M):=\frac{\operatorname{ker}\left(\mathrm{d}: \Omega^{k}(M) \rightarrow \Omega^{k+1}(M)\right)}{\operatorname{im}\left(\mathrm{d}: \Omega^{k-1}(M) \rightarrow \Omega^{k}(M)\right)}
$$

is called the $k$ th de Rham cohomology space of $M$.\\
Definition 10.3.15. Let $M$ and $N$ be smooth manifolds and $f: M \rightarrow N$ be a smooth map. Given a differential form $\omega \in \Omega^{k}(N)$, we can define the pull-back $f^{*} \omega$ by $\left(f^{*} \omega\right)_{p}:=T_{p}(f)^{*} \omega_{f(p)}$, i.e.,

$$
\left(f^{*} \omega\right)_{p}\left(v_{1}, \ldots, v_{p}\right):=\omega_{f(p)}\left(T_{p}(f) v_{1}, \ldots, T_{p}(f) v_{p}\right)
$$

The pull-back of $\omega$ is a differential form $f^{*} \omega \in \Omega^{k}(M)$.\\
Proposition 10.3.16. The pull-back of differential forms is compatible with the exterior derivative, i.e.,

$$
\mathrm{d}\left(f^{*} \omega\right)=f^{*}(\mathrm{~d} \omega)
$$

Proof. In fact, if $\omega=\varphi \in C^{\infty}(N)=\Omega^{0}(N)$, then

$$
f^{*}(\mathrm{~d} \varphi)=\mathrm{d} \varphi \circ T(f)=\mathrm{d}(\varphi \circ f)=\mathrm{d}\left(f^{*} \varphi\right)
$$

Now suppose that we can write $\omega$ as an exterior product of 1 -forms:

$$
\omega=\varphi \mathrm{d} y_{i_{1}} \wedge \cdots \wedge \mathrm{~d} y_{i_{k}}
$$

Then $f^{*}(\alpha \wedge \beta)=f^{*} \alpha \wedge f^{*} \beta$ (Exercise!) implies

$$
\begin{aligned}
& \mathrm{d}\left(f^{*} \omega\right) \\
& =\mathrm{d}\left((\varphi \circ f) f^{*} \mathrm{~d} y_{i_{1}} \wedge \cdots \wedge f^{*} \mathrm{~d} y_{i_{k}}\right)=\mathrm{d}\left((\varphi \circ f) \mathrm{d}\left(f^{*} y_{i_{1}}\right) \wedge \cdots \wedge \mathrm{d}\left(f^{*} y_{i_{k}}\right)\right) \\
& =\mathrm{d}(\varphi \circ f) \wedge \mathrm{d}\left(f^{*} y_{i_{1}}\right) \wedge \cdots \wedge \mathrm{d}\left(f^{*} y_{i_{k}}\right)=f^{*} \mathrm{~d} \varphi \wedge \mathrm{~d}\left(f^{*} y_{i_{1}}\right) \wedge \cdots \wedge \mathrm{d}\left(f^{*} y_{i_{k}}\right) \\
& =f^{*} \mathrm{~d} \varphi \wedge f^{*} \mathrm{~d} y_{i_{1}} \wedge \cdots \wedge f^{*} \mathrm{~d} y_{i_{k}}=f^{*}\left(\mathrm{~d} \varphi \wedge \mathrm{~d} y_{i_{1}} \wedge \cdots \wedge \mathrm{~d} y_{i_{k}}\right) \\
& =f^{*}(\mathrm{~d} \omega)
\end{aligned}
$$

Locally, all differential forms can be written as sums of terms of the type (10.7). Moreover, if $\omega$ vanishes on an open subset $U$ of $N$, then $\mathrm{d} \omega$ vanishes on $U$, and $f^{*} \omega$ vanishes on $f^{-1}(U)$. But then also $\mathrm{d}\left(f^{*} \omega\right)$ vanishes on $U$. Thus the claim can indeed be checked locally.

As a consequence of the preceding argument, pull-backs of closed/exact forms are closed/exact. We conclude that each smooth map $f: M \rightarrow N$ leads to well-defined linear maps (functoriality of de Rham cohomology)

$$
f^{*}: H_{\mathrm{dR}}^{k}(N) \rightarrow H_{\mathrm{dR}}^{k}(M), \quad[\alpha] \mapsto\left[f^{*} \alpha\right]
$$

Lemma 10.3.17 (Homotopy Lemma). Let $M$ and $N$ be smooth manifolds and $F: M \times I \rightarrow N$ be a smooth map, where $I \subseteq \mathbb{R}$ is an open interval containing 0 and 1 . Then the induced smooth maps $F_{0}, F_{1}: M \rightarrow N$ satisfy

$$
F_{0}^{*}=F_{1}^{*}: H_{\mathrm{dR}}^{k}(N) \rightarrow H_{\mathrm{dR}}^{k}(M), \quad k \in \mathbb{N}_{0}
$$

Proof. We write $i_{t}: M \rightarrow M \times\{t\} \subseteq M \times I$ for the canonical embeddings.\\
For $\omega \in \Omega^{k}(M \times I, \mathbb{R}), k \geq 0$, we define the fiber integral $I(\omega) \in \Omega^{k-1}(M)$ by

$$
I(\omega)_{x}\left(v_{1}, \ldots, v_{k-1}\right):=\int_{0}^{1} \omega_{(t, x)}\left(\frac{\partial}{\partial t}, v_{1}, \ldots, v_{k-1}\right) d t
$$

i.e.,

$$
I(\omega)_{x}:=\int_{0}^{1} i_{\frac{\partial}{\partial t}} \omega_{(t, x)} d t
$$

Let $X_{0}, \ldots, X_{k-1} \in \mathcal{V}(M)$ and extend these vector fields in the canonical fashion to vector fields $\widetilde{X}_{i}$ on $M \times I$, constant in the second component. Then we have $\left[\widetilde{X}_{i}, \widetilde{X}_{j}\right]=\left[X_{i}, X_{j}\right]$, and from Cartan's formula (see Proposition 10.3.12) we further get

$$
\begin{aligned}
& \left(\mathrm{d}_{M} I(\omega)\right)\left(X_{0}, \ldots, X_{k-1}\right) \\
& =\mathrm{d}_{M}\left(\int_{0}^{1} i_{\frac{\partial}{\partial t}} \omega(t, x) d t\right)\left(X_{0}, \ldots, X_{k-1}\right) \\
& =\int_{0}^{1}\left(\mathrm{~d}_{M \times I} i_{\frac{\partial}{\partial t}} \omega\right)\left(\widetilde{X}_{0}, \ldots, \widetilde{X}_{k-1}\right) d t \\
& =\int_{0}^{1}\left(\mathcal{L}_{\frac{\partial}{\partial t}} \omega\right)\left(\widetilde{X}_{0}, \ldots, \widetilde{X}_{k}\right) d t-\int_{0}^{1} i_{\frac{\partial}{\partial t}}\left(\mathrm{~d}_{M \times I} \omega\right)\left(\widetilde{X}_{0}, \ldots, \widetilde{X}_{k-1}\right) d t \\
& =\left(i_{1}^{*} \omega\right)\left(X_{0}, \ldots, X_{k-1}\right)-\left(i_{0}^{*} \omega\right)\left(X_{0}, \ldots, X_{k-1}\right)-I\left(\mathrm{~d}_{M \times I} \omega\right)\left(X_{0}, \ldots, X_{k-1}\right)
\end{aligned}
$$

This means that we have the homotopy formula

$$
\mathrm{d}_{M} I(\omega)+I\left(\mathrm{~d}_{M \times I} \omega\right)=i_{1}^{*} \omega-i_{0}^{*} \omega
$$

We apply this formula to $\omega=F^{*} \alpha$ for a closed form $\alpha$ of degree $k \geq 1$ on $N$ and obtain

$$
\left[F_{1}^{*} \alpha-F_{0}^{*} \alpha\right]=\left[i_{1}^{*} \omega-i_{0}^{*} \omega\right]=\left[I\left(\mathrm{~d}_{M \times I} \omega\right)\right]=\left[I\left(F^{*} \mathrm{~d} \alpha\right)\right]=0
$$

For degree $k=0$, the space $H_{\mathrm{dR}}^{0}(N)$ consists of locally constant functions $f$, and since $F_{t}^{*} f=f \circ F_{t}$ does not depend on $t$, we also get $F_{0}^{*}=F_{1}^{*}$ for $k=0$.

Corollary 10.3.18. If $M$ is a smooth manifold and $n \in \mathbb{N}_{0}$, then the projection $p_{M}: M \times \mathbb{R}^{n} \rightarrow M$ defines an isomorphism

$$
p_{M}^{*}: H_{\mathrm{dR}}^{k}(M) \rightarrow H_{\mathrm{dR}}^{k}\left(M \times \mathbb{R}^{n}\right), \quad[\alpha] \mapsto\left[p_{M}^{*} \alpha\right]
$$

for each $k \in \mathbb{N}_{0}$.\\
Proof. For $I=\mathbb{R}$, we consider the smooth map

$$
F: M \times \mathbb{R}^{n} \times I \rightarrow M \times \mathbb{R}^{n}, \quad(m, x, t) \mapsto(m, t x)
$$

Then $F_{1}=\operatorname{id}_{M \times \mathbb{R}^{n}}$, and $F_{0}$ is the projection onto $M \times\{0\}$.\\
The inclusion $i_{0}: M \rightarrow M \times \mathbb{R}^{n}, m \mapsto(m, 0)$ satisfies $p_{M} \circ i_{0}=\operatorname{id}_{M}$ and $i_{0} \circ p_{M}=F_{0}$. In view of Lemma 10.3.17, we have $F_{0}^{*}=F_{1}^{*}=\mathrm{id}$, so that the pull-back maps $p_{M}^{*}$ and $i_{0}^{*}$ induce mutually inverse isomorphisms between the spaces $H_{\mathrm{dR}}^{k}(M)$ and $H_{\mathrm{dR}}^{k}\left(M \times \mathbb{R}^{n}\right)$.\\
Definition 10.3.19. A smooth manifold $M$ is called smoothly contractible to a point $p \in M$ if there exists a smooth map $F: M \times I \rightarrow M$ with

$$
F(x, 0)=x \quad \text { and } \quad F(x, 1)=p
$$

for all $x \in M$, where $I \subseteq \mathbb{R}$ is an open interval containing $[0,1]$.\\
Theorem 10.3.20 (Poincaré Lemma). Suppose that $M$ is a smoothly contractible manifold. Then any closed form of degree at least one is exact.

Proof. Let $F: M \times I \rightarrow M$ be a smooth contraction of $M$ to $p \in M$ and $\alpha \in \Omega^{k}(M), k>0$. Since $F_{0}=\operatorname{id}_{M}$ and $F_{1}=p$ is the constant map, the Homotopy Lemma implies that

$$
[\alpha]=\left[F_{0}^{*} \alpha\right]=\left[F_{1}^{*} \alpha\right]
$$

But since $F_{1}$ is constant and $k>0$, we have $F_{1}^{*} \alpha=0$.\\
Corollary 10.3.21. Suppose that $M \subseteq \mathbb{R}^{n}$ is open and star-shaped. Then any closed form of degree at least one is exact.

Proof. Contract linearly to any point with respect to which $M$ is starshaped.

\subsubsection*{Integration of Densities}
Recall the density bundle $|\Lambda|(M):=|\Lambda|^{1}(M)$ from Example 10.2.16.\\
Remark 10.3.22 (Transformation formula for densities). Let $\left(\varphi_{i}, U_{i}\right)_{i \in I}$ be a smooth atlas of the $n$-dimensional smooth manifold $M$ and $\left(\Phi_{i}\right)_{i \in I}$ be the corresponding trivialization of GL $(M)$ constructed in Example 10.2.9. Then, according to Example 10.2.16, the transition functions of $|\Lambda|(M)$ with respect to the local trivialization $\left(\widetilde{\Phi}_{i}\right)_{i \in I}$ constructed from $\left(\Phi_{i}\right)_{i \in I}$ in Example 10.2.16 are given by

$$
\widetilde{\Phi}_{j i}(p)=\left|\operatorname{det} \mathrm{d}\left(\varphi_{j} \circ \varphi_{i}^{-1}\right)\left(\varphi_{i}(p)\right)\right|^{-1}=\left|\operatorname{det} \mathrm{d}\left(\varphi_{i} \circ \varphi_{j}^{-1}\right)\left(\varphi_{j}(p)\right)\right| .
$$

Let $\mu$ be a density on $M$. Then we obtain for each $i$ a density $\mu_{i}:=\left(\varphi_{i}^{-1}\right)^{*} \mu$ on the open subset $\varphi_{i}\left(U_{i}\right) \subseteq \mathbb{R}^{n}$. On this set, we write $\left|\mathrm{d} x_{1} \cdots \mathrm{~d} x_{n}\right|$ for the density assigning at each point the value 1 to the canonical basis for $\mathbb{R}^{n}$. Then we can write

$$
\mu_{i}=\left(\varphi_{i}^{-1}\right)^{*} \mu=f_{i} \cdot\left|\mathrm{~d} x_{1} \cdots \mathrm{~d} x_{n}\right|
$$

for a real-valued function $f_{i}: \varphi_{i}\left(U_{i}\right) \rightarrow \mathbb{R}$. In terms of the local trivializations, this means that

$$
\widetilde{\Phi}_{i} \circ \mu \circ \varphi_{i}^{-1}=\left(\varphi_{i}^{-1}, f_{i}\right): \varphi_{i}\left(U_{i}\right) \rightarrow U_{i} \times \mathbb{R}
$$

We say that the $f_{i}$ represent $\mu$ in the given atlas, resp., the corresponding local trivialization. The $\left(f_{i}\right)_{i \in I}$ satisfy the following transformation property

$$
\left(f_{j} \circ \varphi_{j}\right)(p)=\widetilde{\Phi}_{j i}(p) \cdot\left(f_{i} \circ \varphi_{i}\right)(p)
$$

which we can rewrite as

$$
f_{j}=\left(f_{i} \circ \varphi_{i} \circ \varphi_{j}^{-1}\right)\left|\operatorname{det}\left(\mathrm{d}\left(\varphi_{i} \circ \varphi_{j}^{-1}\right)\right)\right|
$$

on $\varphi_{j}\left(U_{i} \cap U_{j}\right)$, which corresponds to

$$
\mu_{j}=\left(\varphi_{i} \circ \varphi_{j}^{-1}\right)^{*} \mu_{i}=\left(\varphi_{i} \circ \varphi_{j}^{-1}\right)^{*} f_{j} \cdot\left(\left(\varphi_{i} \circ \varphi_{j}^{-1}\right)^{*}\left|\mathrm{~d} x_{1} \cdots \mathrm{~d} x_{n}\right|\right)
$$

with

$$
\left(\varphi_{i} \circ \varphi_{j}^{-1}\right)^{*}\left|\mathrm{~d} x_{1} \cdots \mathrm{~d} x_{n}\right|=\left|\operatorname{det}\left(\mathrm{d}\left(\varphi_{i} \circ \varphi_{j}^{-1}\right)\right)\right| \cdot\left|\mathrm{d} x_{1} \cdots \mathrm{~d} x_{n}\right| .
$$

Definition 10.3.23. The transformation behavior allows defining positive densities by requiring that they be represented by positive functions. Similarly, we call a density measurable if the representing functions are (Borel) measurable. Finally, we say that a positive measurable density $\mu$ is a locally bounded Borel density if the representing functions with respect to the charts are bounded on compact subsets.

If $\mu$ is a positive and measurable density, we set for any open subset $O \subseteq U_{i}$ :

$$
\int_{O} \mu:=\int_{\varphi_{i}(O)}\left(\varphi_{i}^{-1}\right)^{*} \mu=\int_{\varphi_{i}(O)} f_{i}(x)\left|\mathrm{d} x_{1} \cdots \mathrm{~d} x_{n}\right|:=\int_{\varphi_{i}(O)} f_{i}(x) d x
$$

The change of variables formula from multi-variable calculus allows us to calculate

$$
\begin{aligned}
\int_{\varphi_{j}\left(U_{i} \cap U_{j}\right)} f_{j}(y) d y & =\int_{\varphi_{j}\left(U_{i} \cap U_{j}\right)}\left(f_{i} \circ \varphi_{i} \circ \varphi_{j}^{-1}\right)(y)\left|\operatorname{det} \mathrm{d}\left(\varphi_{i} \circ \varphi_{j}^{-1}\right)(y)\right| d y \\
& =\int_{\varphi_{i}\left(U_{i} \cap U_{j}\right)} f_{i}(x) d x
\end{aligned}
$$

which shows that the definition is independent of the choice of the chart.\\
Definition 10.3.24. Let $M$ be a topological space and $\left(U_{i}\right)_{i \in I}$ be an open cover of $M$. An open cover $\left(V_{j}\right)_{j \in J}$ of $M$ is called a refinement of $\left(U_{i}\right)_{i \in I}$ if for each $j \in J$ there exists an $i \in I$ such that $V_{j} \subseteq U_{i}$. An open cover $\left(U_{i}\right)_{i \in I}$ is called locally finite if for each $p \in M$ there exists a neighborhood $U$ in $M$ such that $U \cap U_{i} \neq \emptyset$ only for finitely many $i \in I$. The space $M$ is called paracompact if it is Hausdorff and every open cover has a locally finite refinement. Further, $M$ is called $\sigma$-compact, if it is a countable union of compact subsets.\\
Proposition 10.3.25. Let $M$ be a second countable smooth manifold of dimension $n$. Then $M$ is $\sigma$-compact and paracompact.

Proof. Let $\left(W_{n}\right)_{n \in \mathbb{N}}$ be a countable basis for the topology of $M$. Since $M$ is locally compact, we may w.lo.g. assume that all the closures $\overline{W_{i}}$ are compact because the relatively compact subsets also form a basis for the topology.

Then choose an increasing family $\left(A_{k}\right)_{k \in \mathbb{N}}$ of compact sets inductively as follows: $A_{1}$ is the closure of $W_{1}$. If $A_{k}$ is defined, set

$$
j_{k}:=\min \left\{j \geq k+1 \mid A_{k} \subseteq \bigcup_{m=1}^{j} W_{m}\right\}
$$

and let $A_{k+1}$ be the closure of $\bigcup_{m=1}^{j_{k}} W_{m}$. Then $\bigcup_{k \in \mathbb{N}} A_{k}=M$ and $A_{k}$ is contained in the interior $A_{k+1}^{\circ}$ of $A_{k+1}$. Since each $A_{k}$ is compact, we see that $M$ is $\sigma$-compact.

Now let $\left(V_{j}\right)_{j \in J}$ be an open cover of $M$. Then any $p \in A_{k+1} \backslash A_{k}^{\circ}$ is contained in some $V_{j}$ and there exists an open neighborhood

$$
U_{p, j} \subseteq V_{j} \cap\left(A_{k+2}^{\circ} \backslash A_{k-1}\right)
$$

of $p$. Then the sets $U_{p, j}$ cover the compact set $A_{k+1} \backslash A_{k}^{\circ}$. Pick a finite subcover, and write $U_{i}$ for the elements of this subcover. Collecting these sets for all $k$ 's gives the desired covering. In fact, the relation $U_{j} \subseteq A_{k+2}^{\circ} \backslash A_{k-1}$ proves local finiteness of the cover.

Definition 10.3.26. Let $M$ be a topological space. A partition of unity is a collection $\left(\rho_{i}\right)_{i \in I}$ of continuous functions $\rho_{i}: M \rightarrow[0,1]$ such that each $p \in M$ has a neighborhood $U$ on which only finitely many functions $\rho_{i}$ are non-zero and

$$
\sum_{i \in I} \rho_{i}(p)=1
$$

If $\left(U_{i}\right)_{i \in I}$ is an open cover of $M$ and $\operatorname{supp}\left(\rho_{i}\right) \subseteq U_{i}$, then the partition of unity is said to be subordinate to the cover $\left(U_{i}\right)_{i \in I}$.

Theorem 10.3.27. If $M$ is second countable and $\left(U_{j}\right)_{j \in J}$ is a locally finite open cover of $M$, then there exists a smooth partition of unity on $M$, which is subordinate to $\left(U_{j}\right)_{j \in J}$.

Proof. The proof of Proposition 10.3.25 shows that we can find a sequence, $\left(K_{i}\right)_{i \in \mathbb{N}}$ of compact subsets of $M$ with $\bigcup_{n \in \mathbb{N}} K_{n}=M$ and $K_{n} \subseteq K_{n+1}^{\circ}$. Put $K_{0}:=\emptyset$. For $p \in M$ let $i_{p}$ be the largest integer with $p \in M \backslash K_{i_{p}}$. We then have $p \in K_{i_{p}+1} \subseteq K_{i_{p}+2}^{\circ}$. Choose a $j_{p} \in J$ with $p \in U_{j_{p}}$ and pick $\psi_{p} \in C^{\infty}(M, \mathbb{R})$ with $\psi_{p}(p)>0$ and

$$
\operatorname{supp}\left(\psi_{p}\right) \subseteq U_{j_{p}} \cap\left(K_{i_{p}+2}^{\circ} \backslash K_{i_{p}}\right)
$$

(Lemma 8.4.15). Then $W_{p}:=\psi_{p}^{-1}(] 0, \infty[$ ) is an open neighborhood of $p$. For each $i \geq 1$, choose a finite set of points $p$ in $M$ whose corresponding neighborhoods $W_{p}$ cover the compact set $K_{i} \backslash K_{i-1}^{\circ}$. We order the corresponding functions $\psi_{p}$ in a sequence $\left(\psi_{i}\right)_{i \in \mathbb{N}}$. Their supports form a locally finite family of subsets of $M$ because for only finitely many of them, the supports intersect a given set $K_{i}$. Moreover, the sets $\psi_{i}^{-1}(] 0, \infty[)$ cover $M$. Therefore,

$$
\psi:=\sum_{j} \psi_{j}
$$

is a smooth function which is everywhere positive (Exercise 10.3.5). Therefore we obtain smooth functions $\varphi_{i}:=\frac{\psi_{i}}{\psi}, i \in \mathbb{N}$. Then the functions $\varphi_{i}$ form a smooth partition of unity on $M$.

We now define a modified partition of unity, which will be subordinate to the open cover $\left(U_{j}\right)_{j \in J}$ : For each $i \in \mathbb{N}$, we pick a $j_{i} \in J$ with $\operatorname{supp}\left(\varphi_{i}\right) \subseteq U_{j_{i}}$ and define

$$
\alpha_{j}:=\sum_{j_{i}=j} \varphi_{i}
$$

As the sum on the right hand side is locally finite, the functions $\alpha_{j}$ are smooth and

$$
\operatorname{supp}\left(\alpha_{j}\right) \subseteq \bigcup_{j_{i}=j} \operatorname{supp}\left(\varphi_{i}\right) \subseteq U_{j}
$$

(Exercise 10.3.3). We further observe that only countably many of the $\alpha_{j}$ are non-zero, that $0 \leq \alpha_{j}, \sum_{j} \alpha_{j}=1$, and that the supports form a locally finite family because the cover $\left(U_{j}\right)_{j \in J}$ is locally finite.

The following corollary is a refinement of Lemma 8.4.15(i).\\
Corollary 10.3.28. Let $M$ be a paracompact smooth manifold, $K \subseteq M$ a closed subset, and $U \subseteq M$ an open neighborhood of $K$. Then there exists a smooth function $f: M \rightarrow \mathbb{R}$ with

$$
0 \leq f \leq 1,\left.\quad f\right|_{K}=1 \quad \text { and } \quad \operatorname{supp}(f) \subseteq U
$$

Proof. In view of Theorem 10.3.27, there exists a smooth partition of unity subordinate to the open cover $\{U, M \backslash K\}$. This is a pair of smooth functions $(f, g)$ with $\operatorname{supp}(f) \subseteq U, \operatorname{supp}(g) \subseteq M \backslash K, 0 \leq f, g$, and $f+g=1$. These properties immediately imply the claim.

Corollary 10.3.29. A second countable manifold $M$ always admits a nowhere vanishing density.

Proof. Recall from Example 10.2.16 how to construct a local trivialization $\left(\Phi_{i}\right)_{i \in I}$ of $|\Lambda|(M)$ from a smooth atlas $\left(\varphi_{i}, U_{i}\right)_{i \in I}$ of $M$. We assume that $\left(U_{i}\right)_{i \in I}$ is locally finite (Proposition 10.3.25). Pick a partition of unity $\left(\rho_{i}\right)_{i \in I}$ subordinate to $\left(U_{i}\right)_{i \in I}$ (Theorem 10.3.27). From the construction of $\left(\Phi_{i}\right)_{i \in I}$ it is clear that $\mu_{i}(p)\left(\left.\frac{\partial}{\partial x_{1}^{(i)}}\right|_{p}, \ldots,\left.\frac{\partial}{\partial x_{n}^{(i)}}\right|_{p}\right)=1$ defines a smooth density over $U_{i}$. Therefore, $\rho_{i} \mu_{i}$ is a density on $M$. We claim that $\mu:=\sum_{i \in I} \rho_{i} \mu_{i}$ is a nowhere vanishing smooth density on $M$. First, the local finiteness of the atlas guarantees that $\mu$ is a well defined density. To see that $\mu$ is nowhere vanishing it suffices to show that $\left.\mu_{j}\right|_{U_{j} \cap U_{i}}$ with respect to the chart $\varphi_{i}$ is represented by a non-negative function. But since $\mu_{j}$ is represented by the constant function 1 with respect to the chart $\varphi_{j}$, this is an immediate consequence of

$$
f_{j}=\left(f_{i} \circ \varphi_{i} \circ \varphi_{j}^{-1}\right)\left|\operatorname{det}\left(\mathrm{d}\left(\varphi_{i} \circ \varphi_{j}^{-1}\right)\right)\right|
$$

on $\varphi_{j}\left(U_{i} \cap U_{j}\right)$.\\
Definition 10.3.30. Let $\mathfrak{B}_{M}$ be the $\sigma$-algebra generated by the open subsets of the locally compact Hausdorff space $M$. It is called the Borel $\sigma$-algebra, and its elements are called Borel subsets or (Borel) measurable subsets. A measure $\mu$ on $\left(M, \mathfrak{B}_{M}\right)$, for which we have $\mu(K)<\infty$ for all compact subsets $K \subseteq M$, is called a Borel measure on $M$.

Definition 10.3.31. Let $M$ be a locally compact Hausdorff space and $C_{c}(M)$ be the space of continuous functions $f: M \rightarrow \mathbb{C}$ with compact support. A linear functional $I: C_{c}(M) \rightarrow \mathbb{C}$ is called positive if $I(f) \geq 0$ whenever $f \geq 0$.

Definition 10.3.32. Let $M$ be a second countable smooth manifold and $\left(\varphi_{\alpha}, U_{\alpha}\right)_{\alpha \in A}$ be an atlas of $M$ corresponding to a locally finite open cover. Further, assume that $\left(\rho_{\alpha}\right)_{\alpha \in A}$ is a subordinate partition of unity. Let $\mu \in \left|\Lambda^{n}\right|(M)$ be a positive measurable density on $M$. We define the integral of $\mu$ over $M$ by

$$
\int_{M} \mu:=\sum_{\alpha \in A} \int_{U_{\alpha}} \rho_{\alpha} \mu
$$

We have to show that this definition is independent of the choice of partition of unity and the underlying atlas. We know already that the definition of the integral over $U_{\alpha}$ is independent of the choice of the coordinates, once the atlas is chosen. So suppose $\left(V_{\lambda}\right)_{\lambda \in \Lambda}$ and $\left(\eta_{\lambda}\right)_{\lambda \in \Lambda}$ is another pair of an atlas and a subordinate partition of unity. Then the calculation

$$
\begin{aligned}
\sum_{\alpha \in A} \int_{U_{\alpha}} \rho_{\alpha} \mu & =\sum_{\alpha \in A} \int_{U_{\alpha}}\left(\sum_{\lambda \in \Lambda} \eta_{\lambda}\right) \rho_{\alpha} \mu=\sum_{\alpha \in A \lambda \in \Lambda} \int_{U_{\alpha}} \eta_{\lambda} \rho_{\alpha} \mu \\
& =\sum_{\lambda \in \Lambda} \sum_{\alpha \in A} \int_{U_{\alpha} \cap V_{\lambda}} \rho_{\alpha} \eta_{\lambda} \mu=\sum_{\lambda \in \Lambda} \int_{V_{\lambda}}\left(\sum_{\alpha \in A} \rho_{\alpha}\right) \eta_{\lambda} \mu=\sum_{\lambda \in \Lambda} \int_{V_{\lambda}} \eta_{\lambda} \mu
\end{aligned}
$$

proves that the integral of $\mu$ over $M$ is well-defined.\\
Proposition 10.3.33. Fix a positive measurable density $\mu$ and for a Borel measurable subset $E \subseteq M$ set

$$
\widetilde{\mu}(E):=\int_{E} \mu:=\int_{M} \chi_{E} \mu \in[0, \infty]
$$

where $\chi_{E}$ is the characteristic function of $E$. Then $\widetilde{\mu}$ is a positive measure on M. If $\mu$ is a locally bounded Borel density, then $\tilde{\mu}$ is a Borel measure.

Proof. Exercise.\\
Remark 10.3.34. Let $\mu$ a strictly positive continuous locally bounded Borel density on $M$, then $\int_{M} f \cdot \mu>0$ for each nonzero continuous function $f \geq 0$. In fact, the integral over any Borel subset of $M$ is nonnegative. Moreover, we can find a coordinate patch $U_{\alpha}$ on which $f$ and the function representing $\mu$ are strictly positive, so that the integral $\int_{U_{\alpha}} f \cdot \mu$ is also strictly positive.

Remark 10.3.35. (i) Given a locally bounded Borel density $\mu$, we can integrate functions $f \in C_{c}(M)$ with respect to $\widetilde{\mu}$ via

$$
I_{\mu}(f):=\int_{M} f(x) \mathrm{d} \widetilde{\mu}(x)=\int_{M} f \cdot \mu
$$

Note that $I_{\mu}: C_{c}(M) \rightarrow \mathbb{R}$ defines a positive linear functional.\\
(ii) Once one has the measure $\widetilde{\mu}$ one can also integrate functions $f: M \rightarrow V$ with values in finite-dimensional vector spaces using either a basis for $V$ and coordinates or linear functionals.\\
(iii) If $A: V \rightarrow W$ is a linear map, then $A$ commutes with integration:

$$
\int_{M}(A \circ f) \mu=A\left(\int_{M} f \mu\right) \quad \text { for } \quad f \in C_{c}(M, V) .
$$

Definition 10.3.36. Let $\varphi: M \rightarrow N$ be a smooth map between $n$-dimensional manifolds and $\mu$ an $r$-density on $N$. Then we define an $r$-density $f^{*} \mu$ on $M$ by

$$
\left(f^{*} \mu\right)(p)\left(v_{1}, \ldots, v_{n}\right)=\mu\left(T_{p} f\left(v_{1}\right), \ldots, T_{p} f\left(v_{n}\right)\right)
$$

for $v_{1}, \ldots, v_{n} \in T_{p}(M)$. We call $f^{*} \mu$ the pull-back of $\mu$ by $f$.\\
Remark 10.3.37 (Global transformation formula for densities). The well-definedness of the integral associated with a locally bounded Borel density shows that, given such a density $\mu$ on $N$, for a diffeomorphism $\varphi: M \rightarrow N$, we have

$$
\int_{\varphi^{-1}(U)}(f \circ \varphi) \cdot \varphi^{*} \mu=\int_{U} f \cdot \mu
$$

where $f \in C_{c}(N)$ and $U \subseteq N$ is a Borel set.

\subsubsection*{Some Technical Results on Integration}
In this section, we compile a few technical results on integration, which will be used in later chapters.

Lemma 10.3.38. Let $I: C_{c}(M) \rightarrow \mathbb{C}$ be a positive linear functional, and let $K \subseteq M$ be compact. Then there exists a constant $C_{K}$ such that

$$
|I(f)| \leq C_{K}\|f\|_{\infty} \quad \text { for } \quad f \in C_{c}(M), \operatorname{supp} f \subseteq K .
$$

Proof. First, we assume that $f$ is a real valued function. By Corollary 10.3.28, there is a function $\varphi \in C_{c}(M)$ with values in $[0,1]$ which is identical to the constant map 1 on $K$. Now, if $f \in C_{c}(M)$ with $\operatorname{supp} f \subseteq K$ then $|f| \leq\|f\|_{\infty} \varphi$. This gives $\|f\|_{\infty} \varphi \pm f \geq 0$, hence, $\|f\|_{\infty} I(\varphi) \pm I(f) \geq 0$, by assumption, and therefore, $|I(f)| \leq I(\varphi)\|f\|_{\infty}$.

If $f=f_{1}+i f_{2}$ is complex-valued, we now obtain with $C:=I(\varphi)$ :\\
$|I(f)|=\left|I\left(f_{1}\right)+i I\left(f_{2}\right)\right| \leq\left|I\left(f_{1}\right)\right|+\left|I\left(f_{2}\right)\right| \leq C\left(\left\|f_{1}\right\|_{\infty}+\left\|f_{2}\right\|_{\infty}\right) \leq 2 C\|f\|_{\infty}$.

Remark 10.3.39. If $\mu$ is a locally bounded Borel density on $M$, then by Proposition 10.3.33 it defines a Borel measure $\widetilde{\mu}$ on $M$ and Lemma 10.3.38 yields the following estimate for $K \subseteq M$ compact and $f \in C_{c}(M)$ :

$$
\left|\int_{M} f \mu\right| \leq c \sup _{x \in K}|f(x)|
$$

Here the constant $c$ depends only on $K$ and $\mu$.\\
Theorem 10.3.40. Let $M$ be a manifold and $h \in C_{c}(M \times M)$ with $\operatorname{supp}(h) \subseteq V \times V$ for a compact subset $V \subseteq M$. Then the following statements are true:\\
(i) The maps

$$
M \rightarrow C(V), \quad g \mapsto(x \mapsto h(g, x))
$$

and

$$
M \rightarrow C(V), \quad g \mapsto(x \mapsto h(x, g))
$$

are continuous, where $C(V)$ is equipped with the supremum norm $\|\cdot\|_{\infty}$.\\
(ii) (Fubini) Let $\mu$ and $\mu^{\prime}$ be positive Borel measures on $M$. Then the functions $y \mapsto \int_{M} h(x, y) d \mu(x)$ and $y \mapsto \int_{M} h(y, x) d \mu^{\prime}(x)$ are continuous with

$$
\int_{M} \int_{M} h(x, y) d \mu(x) d \mu^{\prime}(y)=\int_{M} \int_{M} h(x, y) d \mu^{\prime}(y) d \mu(x)
$$

Proof. (i) It suffices to prove the first of the two assertions. Write $h_{x}(y):= h(x, y)$. Let $x_{0} \in M$ and $\varepsilon>0$. Then

$$
U:=\left\{(x, y) \in M \times V:\left|h(x, y)-h\left(x_{0}, y\right)\right|<\varepsilon\right\}
$$

is an open subset of $M \times V$ containing the compact subset $\left\{x_{0}\right\} \times V$, hence also a set of the form $W \times V$, where $W$ is a neighborhood of $x_{0}$ (Exercise 10.3.6), and this means that for $x \in W$, we $\left\|h_{x}-h_{x_{0}}\right\|_{\infty} \leq \varepsilon$.\\
(ii) By (i) and Lemma 10.3.38, we get the continuity of both functions. Therefore, since the supports of both functions lie in $V$, both double integrals are defined because the integrands are compactly supported continuous functions. By Lemma 10.3.38, the linear maps

$$
\alpha: C(V \times V) \rightarrow \mathbb{R}, \quad f \mapsto \int_{M} \int_{M} f(x, y) d \mu(x) d \mu^{\prime}(y)
$$

and

$$
\beta: C(V \times V) \rightarrow \mathbb{R}, \quad f \mapsto \int_{M} \int_{M} f(x, y) d \mu^{\prime}(y) d \mu(x)
$$

are both continuous with respect to $\|f\|:=\max \{|f(x, y)|: x, y \in V\}$. As we immediately see, both linear maps coincide on functions of the form $f(x, y)= f_{1}(x) f_{2}(y)$. By the Stone-Weierstraß Theorem, the subspace which is spanned by these functions is dense in $C(V \times V)$. This implies $\alpha=\beta$.

Lemma 10.3.41 (Integrals with Parameters). Let ( $M, \mathfrak{M}, \widetilde{\mu}$ ) be a measure space, $I$ an interval of positive length, and let $f: M \times I \rightarrow \mathbb{C}$ be a function such that $f(\cdot, t): M \rightarrow \mathbb{C}$ is integrable for every $t \in I$. We set $F(t):=\int_{M} f(x, t) \mathrm{d} \widetilde{\mu}(x)$.\\
(i) For each integrable, nonnegative function $g$ on $M$ such that $|f(x, t)| \leq g(x)$ for all $x \in M$ and $t \in I$, and $\lim _{t \rightarrow t_{0}} f(x, t)=f\left(x, t_{0}\right)$ for all $x \in M$, we have

$$
\lim _{t \rightarrow t_{0}} F(t)=F\left(t_{0}\right)
$$

(ii) Assume that the partial derivative $\frac{\partial f}{\partial t}$ exists for $t \in I^{\circ}$, and assume furthermore that there exists an integrable, nonnegative function $g$ on $M$ such that $\left|\frac{\partial f}{\partial t}(x, t)\right| \leq g(x)$ for all $x \in M$ and $t \in I^{\circ}$. Then $F$ is differentiable in $] a, b[$, and

$$
F^{\prime}(t)=\int_{M} \frac{\partial f}{\partial t}(x, t) \mathrm{d} \widetilde{\mu}(x)
$$

Proof. (i) Set $f_{n}(x)=f\left(x, t_{n}\right)$ where $t_{n} \rightarrow t_{0} \in I$ for $n \rightarrow \infty$. By Lebesgue's Dominated Convergence Theorem, we get

$$
\begin{aligned}
F\left(t_{0}\right) & =\int_{M} f\left(x, t_{0}\right) \mathrm{d} \widetilde{\mu}(x)=\int_{M} \lim _{n \rightarrow \infty} f_{n}(x) \mathrm{d} \widetilde{\mu}(x) \\
& =\lim _{n \rightarrow \infty} \int_{M} f_{n}(x) \mathrm{d} \widetilde{\mu}(x)=\lim _{n \rightarrow \infty} \int_{M} f\left(x, t_{n}\right) \mathrm{d} \widetilde{\mu}(x)=\lim _{n \rightarrow \infty} F\left(t_{n}\right)
\end{aligned}
$$

(ii) For $h_{n}(x):=\frac{f\left(x, t_{n}\right)-f\left(x, t_{0}\right)}{t_{n}-t_{0}}$ we obtain $\lim _{n \rightarrow \infty} h_{n}(x)=\frac{\partial f}{\partial t}\left(x, t_{0}\right)$, hence, $\frac{\partial f}{\partial t}\left(\cdot, t_{0}\right)$ is measurable. The Mean Value Theorem implies

$$
\left|h_{n}(x)\right| \leq \sup _{t \in I^{\circ}}\left|\frac{\partial f}{\partial t}(x, t)\right| \leq g(x)
$$

for every $x \in M$. Using the Dominated Convergence Theorem again, we get

$$
\begin{aligned}
& \int_{M} \frac{\partial f}{\partial t}\left(x, t_{0}\right) \mathrm{d} \widetilde{\mu}(x)=\lim _{n \rightarrow \infty} \int_{M} h_{n}(x) \mathrm{d} \widetilde{\mu}(x) \\
= & \lim _{n \rightarrow \infty} \frac{1}{t_{n}-t_{0}} \int_{M}\left(f\left(x, t_{n}\right)-f\left(x, t_{0}\right)\right) \mathrm{d} \widetilde{\mu}(x)=\lim _{n \rightarrow \infty} \frac{F\left(t_{n}\right)-F\left(t_{0}\right)}{t_{n}-t_{0}},
\end{aligned}
$$

which implies the claim.\\
Proposition 10.3.42. Let $M$ and $N$ be smooth manifolds, $\mu$ a locally bounded Borel density on $M$, and $f: M \times N \rightarrow \mathbb{R}$ a continuous function which is smooth in the second argument. We make the following assumptions:\\
(a) Locally, all partial derivatives of $f$ with respect to the $N$-variables are also continuous as functions on $M \times N$.\\
(b) There exists a subset $C \subseteq N$ such that $\left.f\right|_{M \times C}$ is compactly supported.

Then

$$
y \mapsto \int_{M} f^{y} \mu, \quad \text { where } \quad f^{y}(x)=f(x, y)
$$

defines a smooth function on the interior $C^{\circ}$ of $C$ such that integration over $M$ and differentiation with respect to the $N$-variables commute.

Proof. From the hypotheses, we see that there exists a compact subset $K \subseteq M$ such that $\operatorname{supp} f \subseteq K \times N$. Pick finitely many open coordinate neighborhoods $U_{1}, \ldots, U_{k}$ in $M$ covering $K$ and a partition of unity $\left(\rho_{i}\right)_{i=0, \ldots, k}$\\
subordinate to $\left(U_{i}\right)_{i=0, \ldots, k}$ with $U_{0}:=M \backslash K$. Then, for each $y \in C$, the function $f^{y}$ is bounded and we can write

$$
\int_{M} f^{y} \mu=\sum_{j=1}^{k} \int_{U_{j}} \rho_{j} f^{y} \mu
$$

To prove the claim, we may now assume that $M \subseteq \mathbb{R}^{m}$ and $C^{\circ} \subseteq \mathbb{R}^{n}$ are open boxes. But then, by hypothesis, the functions $\frac{\partial f}{\partial y_{j}}: M \times C \rightarrow \mathbb{R}$ are compactly supported and smooth. Therefore, one can find a dominating function $g$ as required in Lemma 10.3.41(ii), whence

$$
\frac{\partial}{\partial y_{j}} \int_{M} f(x, y) \mathrm{d} \widetilde{\mu}(x)=\int_{M} \frac{\partial}{\partial y_{j}} f(x, y) \mathrm{d} \widetilde{\mu}(x)
$$

exists and is continuous for $j=1, \ldots, n$. In view of (a), Lemma 10.3.41 shows that $y \mapsto \int_{M} f^{y} \mu$ has continuous partial derivatives, hence is $C^{1}$. But (b) implies that we can apply the same argument to the functions $(x, y) \mapsto \frac{\partial}{\partial y_{j}} f(x, y)$, so we can use induction to complete the proof.

\subsubsection*{Orientations}
We have seen in Subsection 10.3.2 that, under suitable conditions like positivity of compact support, one can integrate densities. In this subsection, we show how to make the connection to the integration of differential forms which is used, for instance, in Stokes' Theorem. The objects that make such a connection possible are the orientations.

Definition 10.3.43. Let $V$ be an $n$-dimensional $\mathbb{R}$-vector space and $e_{1}$, $\ldots, e_{n}, \widetilde{e}_{1}, \ldots, \widetilde{e}_{n}$ be two bases for $V$. Define the base change matrix $A:= \left(a_{i j}\right)_{1 \leq i, j \leq n}$ in $\mathrm{GL}_{n}(\mathbb{R})$ via

$$
e_{i}=\sum_{j=1}^{n} a_{j i} \widetilde{e}_{j}
$$

If $\operatorname{det}(A)>0$, we say that the two bases have the same orientation, and if $\operatorname{det}(A)<0$, we say that the two bases have opposite orientation.

Definition 10.3.44. Let $M$ be an $n$-dimensional smooth manifold. Two charts ( $\varphi, U$ ) and ( $\psi, V$ ) of $M$ are said to be equally oriented if

$$
\operatorname{det}\left(\mathrm{d}\left(\psi \circ \varphi^{-1}\right)(x)\right)>0 \quad \text { for } \quad x \in \varphi(U \cap V) .
$$

An orientation of $M$ is an atlas $\left(\varphi_{\alpha}, U_{\alpha}\right)_{\alpha \in J}$ which is maximal with respect to the property that any two charts in this atlas are equally oriented. The manifold $M$ is called orientable, if it has an orientation (for connected $M$ in that case there are precisely two orientations).

Remark 10.3.45. Suppose that $x^{\alpha}$ and $x^{\beta}$ are the coordinate functions on $U_{\alpha}$ and $U_{\beta}$ for two equally oriented charts of $M$. If $x^{\alpha}=f_{\alpha \beta}\left(x^{\beta}\right)$ are corresponding coordinate changes, then

$$
\operatorname{det}\left(\frac{\partial x_{i}^{\beta}}{\partial x_{j}^{\alpha}}\right)>0
$$

This means that for $p \in U_{\alpha} \cap U_{\beta}$ the bases

$$
\left.\frac{\partial}{\partial x_{1}^{\alpha}}\right|_{p}, \ldots,\left.\frac{\partial}{\partial x_{n}^{\alpha}}\right|_{p} \text { and }\left.\frac{\partial}{\partial x_{1}^{\beta}}\right|_{p}, \ldots,\left.\frac{\partial}{\partial x_{n}^{\beta}}\right|_{p}
$$

of $T_{p}(M)$ have the same orientation.\\
Definition 10.3.46 (Orientation bundle). Let $M$ be a smooth manifold of dimension $n$ and $\operatorname{GL}(M) \rightarrow M$ be the frame bundle of $M$. Consider the one-dimensional representation $g \mapsto \operatorname{sign}(\operatorname{det} g)$ of $G=\mathrm{GL}_{n}(\mathbb{R})$. Then the line bundle $\mathrm{OR}(M):=\mathrm{GL}(M) \times_{G} \mathbb{R}$ associated with this representation via Definition 10.2.10 is called the orientation bundle of $M$. The transition functions of OR( $M$ ) are the signs of the Jacobi determinants of the coordinate changes.

The orientation bundle $\mathrm{OR}(M)$ can be used to give an elegant characterization of the orientability of $M$.

Proposition 10.3.47. A smooth manifold $M$ is orientable if and only if OR(M) admits a nowhere vanishing smooth section.

Proof. Given an orientation, we can cover $M$ by charts such that the Jacobi determinants of all coordinate changes are positive. Thus taking the constant function 1 on all trivializations defines a smooth nowhere vanishing section of the orientation bundle (Proposition 10.2.15).

Conversely, suppose we have a nowhere vanishing smooth section $\sigma$ of OR( $M$ ). Let $\left(U_{\alpha}\right)_{\alpha \in J}$ be a cover of $M$ by connected coordinate neighborhoods with coordinate functions $x_{i}^{\alpha}$ on $U_{\alpha}$ such that $\mathrm{OR}(M)$ is trivial over each $U_{\alpha}$. If $s_{\alpha}: \varphi_{\alpha}(U) \rightarrow \mathbb{R}$ is the function representing $\sigma$ on $U_{\alpha}$, then changing the first coordinate $x_{1}^{\alpha}$ to $-x_{1}^{\alpha}$ if necessary, we may assume that all $s_{\alpha}$ are strictly positive. But then

$$
s_{\alpha}\left(x^{\alpha}\right)=s_{\beta}\left(x^{\beta}\right) \operatorname{sign}\left(\operatorname{det}\left(\frac{\partial x_{i}^{\alpha}}{\partial x_{j}^{\beta}}\right)\right)
$$

now implies that $\operatorname{det}\left(\frac{\partial x_{i}^{\alpha}}{\partial x_{j}^{\beta}}\right)>0$, i.e., the charts $\left(\varphi_{\alpha}, U_{\alpha}\right)$ and $\left(\varphi_{\beta}, U_{\beta}\right)$ are equally oriented for all $\alpha, \beta$. Hence the cover defines an orientation on $M$.

Remark 10.3.48. Let $M$ be a smooth manifold of dimension $n$. Then the tensor product of the line bundles $\operatorname{OR}(M)$ and $\left|\Lambda^{n}\right|(M)$ is again a line bundle and it corresponds to the representation det: $\mathrm{GL}_{n}(\mathbb{R}) \rightarrow \mathbb{R}$. But then the transition functions are the Jacobi determinants, i.e., this is the bundle of $n$-forms. If $M$ is orientable and $\sigma$ is a nowhere vanishing smooth section of OR $(M)$, then multiplication by $\sigma$ defines isomorphisms between the spaces of densities (smooth, continuous, compactly supported, etc.) and $n$-forms (with the corresponding properties). Thus, given $\sigma$ we can integrate $n$-forms.

Definition 10.3.49. Let $M$ be a smooth manifold of dimension $n$. An element $\mu \in \Omega^{n}(M)$ is called a volume form if it does not vanish anywhere.

Corollary 10.3.50. A second countable manifold $M$ of dimension $n$ is orientable if and only if there exists a volume form $\mu \in \Omega^{n}(M)$.

Example 10.3.51. (i) Consider the 2 -torus $T=\mathbb{R}^{2} /\left(\mathbb{Z} e_{1}+\mathbb{Z} e_{2}\right)$ with the 1 -forms $\mathrm{d} x_{1}, \mathrm{~d} x_{2}$, where $x_{1}, x_{2}$ are the coordinates on $\mathbb{R}^{2}$. Then $\mathrm{d} x_{1} \wedge \mathrm{~d} x_{2} \in \Omega^{2}(T)$ is never zero.\\
(ii) Consider the sphere $\mathbb{S}^{n} \subset \mathbb{R}^{n+1}$ and define $\mu \in \Omega^{n}\left(\mathbb{S}^{n}\right)$ via $\mu(x):= i_{x}\left(\mathrm{~d} x_{1} \wedge \cdots \wedge \mathrm{~d} x_{n+1}\right)$, where we identify the tangent spaces $T_{x}\left(\mathbb{S}^{n}\right)$ with the corresponding subspaces $x^{\perp} \subseteq \mathbb{R}^{n+1}$. Then we find

$$
\begin{aligned}
\mu(x) & =i_{\sum x_{i} \frac{\partial}{\partial x_{j}}}\left(\mathrm{~d} x_{1} \wedge \cdots \wedge \mathrm{~d} x_{n+1}\right) \\
& =x_{1} \mathrm{~d} x_{2} \wedge \cdots \wedge \mathrm{~d} x_{n+1}+\cdots+(-1)^{n-1} x_{n} \mathrm{~d} x_{1} \wedge \cdots \wedge \mathrm{~d} x_{n} \\
& =\sum_{j=1}^{n+1}(-1)^{j-1} x_{j} \mathrm{~d} x_{1} \wedge \cdots \wedge \widehat{\mathrm{~d} x_{j}} \wedge \cdots \wedge \mathrm{~d} x_{n+1}
\end{aligned}
$$

(iii) The real projective spaces $\mathbb{P}_{n}(\mathbb{R})$ defined in Example 8.2.15 are not orientable for $n$ even. To prove this, assume to the contrary that $n$ is even and that $\eta \in \Omega^{n}\left(\mathbb{P}_{n}(\mathbb{R})\right)$ is nowhere zero. Let $\pi: \mathbb{S}^{n} \rightarrow \mathbb{P}_{n}(\mathbb{R})$ denote the quotient map. Then also $\pi^{*} \eta$ is nowhere zero. Recall that $\mu(x):=i_{x} \mathrm{~d} x_{1} \wedge \cdots \wedge \mathrm{~d} x_{n+1}$ defines an element $\mu \in \Omega^{n}\left(\mathbb{S}^{n}\right)$, which is nowhere zero. Therefore, $\pi^{*} \eta=f \mu$ for some $f \in C^{\infty}\left(\mathbb{S}^{n}\right)$ satisfying $f(p) \neq 0$ for all $p \in \mathbb{S}^{n}$. Since $\mathbb{S}^{n}$ is connected, we may assume that $f(p)>0$ for all $p$. Consider the map $\sigma: \mathbb{S}^{n} \rightarrow \mathbb{S}^{n}$ defined by $\sigma(x):=-x$. Then we have $\sigma^{*} \mu=(-1)^{n+1} \mu=-\mu$ (because $n$ is even) and since

$$
-\sigma^{*} f \cdot \mu=\sigma^{*} f \cdot \sigma^{*} \mu=\sigma^{*}(f \mu)=\sigma^{*} \pi^{*} \eta=(\pi \circ \sigma)^{*} \eta=\pi^{*} \eta=f \cdot \mu
$$

implies $-f \circ \sigma=f$, we arrive at a contradiction to the positivity of $f$.\\
Remark 10.3.52. If $\mu \in \Omega^{n}(M)$ is a volume form and $\varphi \in \operatorname{Diff}(M)$ is a diffeomorphism of $M$, then $\varphi$ is called orientation preserving if $\varphi^{*} \mu=f \mu$ for a strictly positive function $f \in C^{\infty}(M)$. If $M$ is connected, then this definition does not depend on the choice of the volume form. In fact, if $\widetilde{\mu}$ is\\
another volume form, then $\widetilde{\mu}=h \mu$ for a smooth function $h: M \rightarrow \mathbb{R}^{\times}$. Since $M$ is assumed to be connected, we either have $h>0$ or $h<0$ which implies $\frac{\varphi^{*} h}{h}>0$ and further

$$
\varphi^{*} \widetilde{\mu}=\left(\varphi^{*} h\right)\left(\varphi^{*} \mu\right)=\frac{\varphi^{*} h}{h} f h \mu=\frac{\varphi^{*} h}{h} f \widetilde{\mu}
$$

with $\frac{\varphi^{*} h}{h} f>0$.\\
Proposition 10.3.53 (Orientation of homogeneous spaces). Let $G$ be a Lie group and $H$ a closed subgroup of $G$. If there exists an $\operatorname{Ad}_{\mathfrak{g} / \mathfrak{h}} H$-invariant volume form on $\mathfrak{g} / \mathfrak{h}$, then there exists a $G$-invariant volume form on $G / H$. In particular, $G / H$ is orientable.

Proof. Let $\omega_{0} \in \operatorname{Alt}^{k}(\mathfrak{g} / \mathfrak{h}, \mathbb{R})$ be an Ad $H$-invariant volume form. The description of $\mathrm{Alt}^{k} T(G / H)$ as the associated bundle

$$
G \times_{H} \operatorname{Alt}^{k}(\mathfrak{g} / \mathfrak{h}, \mathbb{R})
$$

(see Definition 10.3.1 and Corollary 10.2.13) immediately shows that

$$
\omega(g H)=\left[g, \omega_{0}\right] \in G \times_{H} \operatorname{Alt}^{k}(\mathfrak{g} / \mathfrak{h}, \mathbb{R})
$$

for $g H \in G / H$ defines a $G$-invariant volume form on $G / H$.\\
Remark 10.3.54. Proposition 10.3.53 in particular shows that each Lie group is orientable and that for each non-zero $\omega_{0} \in \operatorname{Alt}^{k}(\mathfrak{g}, \mathbb{R})$ there is a unique choice of orientation in such a way that the left-invariant volume form $\omega$ associated with $\omega_{0}$ is positive.

\subsubsection*{Exercises for Section 10.3}
Exercise 10.3.1. Let $M$ be a manifold and ( $\mathbb{E}, M, E, \pi$ ) a vector bundle. A subset $\mathbb{F} \subseteq \mathbb{E}$ is called a vector subbundle of $\mathbb{E}$ if the sets $\mathbb{F}_{p}:=\mathbb{F} \cap \mathbb{E}_{p}$ are vector subspaces of the fibers $\mathbb{E}_{p}$ such that the restriction $\left.\pi\right|_{\mathbb{F}}: \mathbb{F} \rightarrow M$ defines a vector bundle in its own right.

Exercise 10.3.2. Let $X$ be a topological space and $\left(K_{n}\right)_{n \in \mathbb{N}}$ a sequence of compact subsets of $X$ with $\bigcup_{n \in \mathbb{N}} K_{n}=X$ and $K_{n} \subseteq K_{n+1}^{0}$ for each $n \in \mathbb{N}$. Show that for each compact subset $C \subseteq X$ there exists an $n \in \mathbb{N}$ with $C \subseteq K_{n}$.

Exercise 10.3.3. A family $\left(S_{i}\right)_{i \in I}$ of subsets of a topological space $X$ is said to be locally finite if each point $p \in X$ has a neighborhood intersecting only finitely many $S_{i}$. Show that if $\left(S_{i}\right)_{i \in I}$ is a locally finite family of closed subsets of $X$, then $\bigcup_{i \in I} S_{i}$ is closed.

Exercise 10.3.4. Let $\left(S_{i}\right)_{i \in I}$ be a locally finite family of subsets of the topological space $X$. Show that each compact subset $K \subseteq X$ intersects only finitely many of the sets $S_{i}$.

Exercise 10.3.5. Let $E$ be a finite-dimensional vector space. A family $\left(f_{j}\right)_{j \in J}$ of smooth $E$-valued functions on $M$ is called locally finite if each point $p \in M$ has a neighborhood $U$ for which the set $\left\{j \in J:\left.f_{j}\right|_{U} \neq 0\right\}$ is finite. Show that this implies that $f:=\sum_{j \in J} f_{j}$ defines a smooth $E$-valued function on $M$.

Exercise 10.3.6. Let $X$ and $Y$ be Hausdorff spaces and $K \subseteq X$, resp., $Q \subseteq Y$, be a compact subset. Then for each open subset $U \subseteq X \times Y$ containing $K \times Q$, there exist open subsets $U_{K} \subseteq X$ containing $K$ and $U_{Q} \subseteq Y$ containing $Y$ with

$$
K \times Q \subseteq U_{K} \times U_{Q} \subseteq U .
$$

\subsection*{Invariant Integration}
In this section, we show how to construct measures on Lie groups and certain of their homogeneous spaces which are invariant under the group action. Following the approach via densities laid out in Section 10.3, we start with a description of invariant densities which will then give rise to the desired measures.

\subsubsection*{Invariant Densities}
Definition 10.4.1. Let $M$ be a smooth manifold, $G$ a Lie group, and $\sigma: G \times M \rightarrow M$ a smooth action. An $r$-density $\mu$ on $M$ is called $G$-invariant, if $\sigma_{g}^{*} \mu=\mu$, where $\sigma_{g}^{*} \mu$ is the pull-back of $\mu$ by $\sigma_{g}$ (see Definition 10.3.36).

If $M=G$, then using left- and right translations as actions we obtain the notions of left- and right invariant densities on $G$.

Remark 10.4.2. (i) If a group $G$ acts linearly on a finite-dimensional vector space $V$, then this action induces a $G$-action on the space of $r$-densities on $V$ via

$$
(g \cdot \mu)\left(v_{1}, \ldots, v_{n}\right):=\mu\left(g^{-1} v_{1}, \ldots, g^{-1} v_{n}\right)=|\operatorname{det} g|^{-r}(\mu)\left(v_{1}, \ldots, v_{n}\right)
$$

for $v_{1}, \ldots, v_{n} \in V$. Thus it makes sense to talk about $G$-invariant densities on $V$.\\
(ii) Let $M$ be a smooth manifold, $G$ a Lie group, and $\sigma: G \times M \rightarrow M$ a smooth action. Recall the two realizations $|\Lambda|^{r}(M)$ and $L^{(r)}(M)$ of the $r$-density bundle from Definition 10.2.16. The action of $G$ lifts to actions on $|\Lambda|^{r}(M)$ and $L^{(r)}(M)$ in such a way that the isomorphism $[\varphi, t] \mapsto \mu_{\varphi, t}$ between the two spaces described in (10.3) is $G$-equivariant: The lift of $\sigma$ to\\
$\mathrm{GL}(M) \times{ }_{\Delta_{r}} \mathbb{R}$ is obtained by the lift of $\sigma$ to $\mathrm{GL}(M)$ which in turn is provided by the tangent maps of the $\sigma_{g}$. On $L^{(r)}(M)$, we set

$$
(g \cdot \mu)\left(v_{1}, \ldots, v_{n}\right)=\mu\left(T_{p}\left(\sigma_{g^{-1}}\right) v_{1}, \ldots, T_{p}\left(\sigma_{g^{-1}}\right) v_{n}\right)
$$

for $v_{1}, \ldots, v_{n} \in T_{\sigma_{g}(p)}(M)$ and $\mu \in L_{p}^{(r)}(M)$. To see the $G$-equivariance, we have to show that $\mu_{\sigma_{g}(\varphi), t}=g \cdot \mu_{\sigma_{g}(\varphi), t}$ for $g \in G$. Denote the canonical projection $\operatorname{GL}(M) \rightarrow M$ by $\widetilde{\pi}$. For $\widetilde{\pi}(\varphi)=p$ and $v_{j}=T_{p}\left(\sigma_{g}\right) \varphi\left(e_{j}\right)=\sigma_{g}(\varphi)\left(e_{j}\right)$, we have $\mu_{\sigma_{g}(\varphi), t}\left(v_{1}, \ldots, v_{n}\right)=t$. On the other hand,

$$
\left(g \cdot \mu_{\varphi, t}\right)\left(v_{1}, \ldots, v_{n}\right)=\mu_{\varphi, t}\left(\varphi\left(e_{1}\right), \ldots, \varphi\left(e_{n}\right)\right)=t
$$

so that the claim is proved. Note that the corresponding $G$-action on the sections of $L^{(r)}(M)$ is the pull-back of densities.

Lemma 10.4.3. Let $G$ be a Lie group and $H$ a closed subgroup. Then the following conditions are equivalent:\\
(1) There exists a $G$-invariant smooth density $\mu$ on $G / H$.\\
(2) There exists an $\operatorname{Ad}_{\mathfrak{g} / \mathfrak{h}}(H)$-invariant density $\mu_{0}$ on $\mathfrak{g} / \mathfrak{h}$.\\
(3) $\left|\operatorname{det} \circ \operatorname{Ad}_{G}\right|_{H}\left|=\left|\operatorname{det} \circ \operatorname{Ad}_{H}\right|\right.$.

Proof. Recall from Corollary 10.2.13 that we have a $G$-equivariant bundle isomorphism $G \times_{H} \mathfrak{g} / \mathfrak{h} \rightarrow T(G / H)$. Therefore, the frame bundle $\operatorname{GL}(G / H)$ can be obtained as $G \times_{H} \mathrm{GL}(\mathfrak{g} / \mathfrak{h})$, where $H$ acts on $\mathrm{GL}(\mathfrak{g} / \mathfrak{h})$ by $(h, A) \mapsto \operatorname{Ad}_{\mathfrak{g} / \mathfrak{h}}(h) A$. But then the density bundle is

$$
\left(G \times_{H} \operatorname{GL}(\mathfrak{g} / \mathfrak{h})\right) \times_{\mathrm{GL}(\mathfrak{g} / \mathfrak{h})} \mathbb{R}=G \times_{H} \mathbb{R}
$$

where the action of $H$ on $\mathbb{R}$ is given by

$$
(h, t) \mapsto\left|\operatorname{det}\left(\operatorname{Ad}_{\mathfrak{g} / \mathfrak{h}}(h)\right)\right|^{-1} t
$$

According to Proposition 10.2.17, the smooth sections of $G \times_{H} \mathbb{R}$ are of the form $g H \mapsto[g, F(g)]$, where $F: G \rightarrow \mathbb{R}$ is a smooth function such that $F(g h)=h^{-1} \cdot F(g)$ for all $g \in G$ and $h \in H$. Such a section is $G$-invariant if and only if $F$ is constant since $G$-invariance implies

$$
[g x, F(x)]=g \cdot[x, F(x)]=[g x, F(g x)]
$$

Thus invariant densities on $G / H$ are in one-to-one correspondence with $H$ fixed points in the fiber over the base point $\mathbf{1} H$, i.e., $\operatorname{Ad}_{\mathfrak{g} / \mathfrak{h}}(H)$-invariant densities on $\mathfrak{g} / \mathfrak{h}$. But then (10.9) implies the equivalence of (1), (2), and (3) since

$$
\operatorname{det}\left(\operatorname{Ad}_{\mathfrak{g} / \mathfrak{h}}(h)\right) \operatorname{det}\left(\operatorname{Ad}_{H}(h)\right)=\operatorname{det}\left(\operatorname{Ad}_{G}(h)\right) .
$$

In view of Proposition 10.3.33 and Remark 10.3.39, we see that any $\operatorname{Ad}_{\mathfrak{g} / \mathfrak{h}} H$-invariant density $\mu$ defines a positive Borel measure $\widetilde{\mu}$ on $G / H$.

Proposition 10.4.4. Let $G$ be a Lie group, $H$ a closed subgroup, and $\mu_{0}$ an $\operatorname{Ad}_{\mathfrak{g} / \mathfrak{h}}(H)$-invariant density on $\mathfrak{g} / \mathfrak{h}$ with positive values. Then the associated smooth $G$-invariant density $\mu$ is positive. The associated measure $\widetilde{\mu}$ (see Remark 10.3.32) is left invariant and satisfies

$$
\int_{G / H} \lambda_{g}^{*} f \mathrm{~d} \widetilde{\mu}=\int_{G / H} f \mathrm{~d} \widetilde{\mu} \quad \text { for } \quad f \in C_{c}(G / H) .
$$

Proof. The first claim follows immediately from the proof of Lemma 10.4.3. Thus it suffices to show that $\widetilde{\mu}$ is $G$-invariant if and only if $\mu$ is $G$-invariant. But this follows from the global transformation formula given in Remark 10.3.37, which allows us to write

$$
\int_{M} \lambda_{g}^{*} f \mathrm{~d}\left(\lambda_{g}^{*} \mu\right)^{\sim}=\int_{M} f \mathrm{~d} \widetilde{\mu} .
$$ \(\square\)

Definition 10.4.5 (Haar measure). Let $G$ be a Lie group. A positive Borel measure $\mu$ on $G$ is called a left Haar measure, if it is invariant under left translations. Similarly, it is called a right Haar measure, if it is invariant under right translations. We will denote Haar measures on $G$ by $\mu_{G}$ and write the corresponding integrals as

$$
\int_{G} f \mathrm{~d} \mu_{G}=\int_{G} f(g) \mathrm{d} \mu_{G}(g) .
$$

Remark 10.4.6. (i) Note that Proposition 10.4.4 guarantees that left Haar measures exist. The existence of right Haar measures is shown analogously using right invariant densities.\\
(ii) The question of uniqueness is more subtle. It can be shown that left/right Haar integrals $C_{c}(G) \rightarrow \mathbb{R}$ are unique up to positive scalars. On the level of measures, one first has to observe that on Lie groups with at most countably many connected components Borel measures automatically satisfy additional regularity conditions ([Ru86, Thm. 2.17]) ensuring that they are uniquely determined by the corresponding integrals of compactly supported continuous functions (which are in particular supported by finitely many connected components). The uniqueness of Haar measure up to positive scalar multiples then follows from a general theorem on locally compact groups. Together with Proposition 10.4.4 this shows that the left/right Haar integrals on a Lie group are in one-to-one correspondence with the left/right invariant densities. In this book, we shall avoid arguments using the uniqueness of Haar measures and replace them by arguments using concrete left/right invariant densities.\\
(iii) In view of Proposition 10.3.53 (see also Corollary 10.3.50), Remark 10.3.39 also shows that any $\operatorname{Ad}_{\mathfrak{g} / \mathfrak{h}} H$-invariant volume form $\omega_{0}$ on $\mathfrak{g} / \mathfrak{h}$ defines a left invariant nonvanishing differential form of degree $\operatorname{dim} G / H$ and then a positive Borel measure on $G / H$ which we denote by $\omega$. As in Proposition 10.4.4, one shows that $\omega$ is a left Haar measure. In many texts on Lie groups, this is the path taken for the construction of Haar measures.

\subsubsection*{Integration Formulas}
Proposition 10.4.7. If $\mu_{l}$ is a left invariant density on a Lie group $G$, then

$$
\mu_{r}(x):=|\operatorname{det} \operatorname{Ad}(x)| \mu_{l}(x)
$$

is a right invariant density on $G$ and

$$
\rho_{g}^{*} \mu_{l}=|\operatorname{det} \operatorname{Ad}(g)|^{-1} \mu_{l}
$$

Proof. In view of $\left(\rho_{g}^{*} \mu_{r}\right)(x)=|\operatorname{det} \operatorname{Ad}(x g)|\left(\rho_{g}^{*} \mu_{l}\right)(x)$, it remains to verify (10.10). As both sides of this equation are left invariant densities, it suffices to show that their values at $\mathbf{1}$ coincide. With $T_{\mathbf{1}}\left(c_{g}\right)=\operatorname{Ad}(g)$, this follows from

$$
\left(\rho_{g}^{*} \mu_{l}\right)(\mathbf{1})=\left(c_{g^{-1}}^{*} \mu_{l}\right)(\mathbf{1})=\left|\operatorname{det} \operatorname{Ad}\left(g^{-1}\right)\right| \mu_{l}(\mathbf{1})
$$

Corollary 10.4.8. If $\mu$ is a left invariant density on $G$, then

$$
\int_{G} f\left(g^{-1}\right) \mathrm{d} \widetilde{\mu}(g)=\int_{G} f(g)|\operatorname{det} \operatorname{Ad}(g)| \mathrm{d} \widetilde{\mu}(g)
$$

Proof. Let $\eta_{G}(g)=g^{-1}$ be the inversion on $G$ and note that $\eta_{G}^{*} \mu$ is a right invariant density with the same value as $\mu$ at $\mathbf{1}$. Now Proposition 10.4.7 implies that

$$
\eta_{G}^{*} \mu=|\operatorname{det} \circ \mathrm{Ad}| \cdot \mu
$$

so that the assertion follows from the transformation formula for integrals with respect to densities:

$$
\int_{G} f\left(g^{-1}\right) \mathrm{d} \widetilde{\mu}(g)=\int_{G} f(g) \mathrm{d}\left(\eta_{G}^{*} \mu\right)(g)=\int_{G} f(g)|\operatorname{det} \operatorname{Ad}(g)| \mathrm{d} \widetilde{\mu}(g)
$$

Definition 10.4.9. A Lie group $G$ is called unimodular if $|\operatorname{det} \operatorname{Ad}(x)|=1$ for all $x \in G$. The function $\Delta_{G}(g):=|\operatorname{det} \circ \operatorname{Ad}(g)|^{-1}$ on $G$ is called the modular function. It is a continuous group homomorphism $\Delta_{G}: G \rightarrow\left(\mathbb{R}_{+}^{\times}, \cdot\right)$.

If $G$ is unimodular, any left invariant density is also a right invariant density measure. Thus the corresponding left Haar measure is also a right Haar measure.

Remark 10.4.10. For a connected Lie group $G$, the following properties are equivalent:\\
(1) $G$ is unimodular.\\
(2) $\operatorname{det} \circ \mathrm{Ad} \equiv 1$.\\
(3) $\mathrm{tr} \circ \mathrm{ad} \equiv 0$.

Proposition 10.4.11. The following properties of a Lie group $G$ imply that it is unimodular:\\
(a) $G$ is compact.\\
(b) $G$ is abelian.\\
(c) The commutator group $G^{\prime}$ of $G$ is dense in $G$.\\
(d) $G$ is connected with nilpotent Lie algebra.

Proof. (a) $\Delta_{G}(G)$ is a compact subgroup of $\left(\mathbb{R}_{+}^{\times}, \cdot\right)$, hence trivial.\\
(b) If $G$ is abelian, then $\operatorname{Ad}(g)=\mathbf{L}\left(c_{g}\right)=\operatorname{id}_{\mathfrak{g}}$ for each $g \in G$.\\
(c) Since $\mathbb{R}^{\times}$is abelian, $G^{\prime} \subseteq \operatorname{ker} \Delta_{G}$, so, by continuity, the hypothesis implies that $\Delta_{G}=1$.\\
(d) Since $\mathfrak{g}$ is nilpotent, the operators $\operatorname{ad}(x)$ for $x \in \mathfrak{g}$ are nilpotent, so they have zero trace. Hence $\operatorname{Ad}(\exp x)=e^{\operatorname{ad}(x)}$ leads to $\operatorname{det} \operatorname{Ad}(\exp x)= e^{\operatorname{tr} \operatorname{ad}(x)}=1$, which implies (d) because $\exp \mathfrak{g}$ generates $G$.

Proposition 10.4.12. Let $G$ be a Lie group and $H$ a closed subgroup. Suppose that $\mu$ is the $G$-invariant density on $G / H$ associated with an $\operatorname{Ad}_{\mathfrak{g} / \mathfrak{h}}(H)$ invariant density $\mu_{0}$ on $\mathfrak{g} / \mathfrak{h}$. Then we can find Haar measures $\mu_{G}$ and $\mu_{H}$ such that

$$
\int_{G} f(g) \mathrm{d} \mu_{G}(g)=\int_{G / H}\left(\int_{H} f(g h) \mathrm{d} \mu_{H}(h)\right) \mathrm{d} \widetilde{\mu}(g H)
$$

for all $f \in C_{c}(G)$.\\
Proof. Using partitions of unity, it suffices to show formula (10.11) for $G$ replaced by open subsets of the form $q^{-1}\left(U_{\alpha}\right)$, where $\left(U_{\alpha}\right)_{\alpha \in A}$ is a locally finite open covering of $G / H$. Using Corollary 10.1.11, we can cover $G / H$ by open sets $U$ for which we have a smooth section $\sigma: U \rightarrow G$ for the quotient map $q: G \rightarrow G / H$ such that $m_{\sigma}: U \times H \rightarrow \sigma(U) H,(u, h) \mapsto \sigma(U) H$ is a diffeomorphism onto an open subset of $G$. Note that $\sigma(U) H=q^{-1}(U)$.

Pick a basis $v_{k+1}, \ldots, v_{n}$ for $\mathfrak{h}$ and complete it by vectors $v_{1}, \ldots, v_{k}$ to a basis for $\mathfrak{g}$. Then the $\bar{v}_{j}:=v_{j}+\mathfrak{h}, j=1, \ldots, k$, form a basis for $\mathfrak{g} / \mathfrak{h}$ which we identify with the tangent space $T_{p_{o}}(G / H)$ of $G / H$ at the base point $p_{o}=q(\mathbf{1})$. We may assume that $\mu\left(p_{o}\right)\left(\bar{v}_{1}, \ldots, \bar{v}_{k}\right)=1$. Let $\nu_{H}$, resp., $\nu_{G}$, be a left invariant density on $H$, resp., $G$, with $\mu_{G}=\widetilde{\nu}_{G}$, resp., $\mu_{H}=\widetilde{\nu}_{H}$. We normalize them in such a way that

$$
\nu_{H}(\mathbf{1})\left(v_{k+1}, \ldots, v_{n}\right)=\nu_{G}(\mathbf{1})\left(v_{1}, \ldots, v_{n}\right)=1
$$

Consider the product density $\mu \otimes \nu_{H}$ on $U \times H$.\\
Claim: $m_{\sigma}^{*} \nu_{G}=\mu \otimes \nu_{H}$.\\
Before we prove the claim, we show how it implies the proposition: Since $m_{\sigma}^{-1}(U H)=U \times H$ and $U H=q^{-1}(U)$, the transformation formula for densities shows that

$$
\int_{U \times H} f(\sigma(u) h) \mathrm{d} \widetilde{\mu}(u) \mathrm{d} \mu_{H}(h)=\int_{q^{-1}(U)} f(g) \mathrm{d} \mu_{G}(g),
$$

where $\mu_{H}$ and $\mu_{G}$ are the Haar measures corresponding to the densities $\nu_{H}$ and $\nu_{G}$. But this is precisely what we had to show since $\int_{H} f(\sigma(u) h) \mathrm{d} \mu_{H}(h)$ is independent of the choice of $\sigma(u)$.

To prove the claim, pick $u \in U$ and write $g=\sigma(u)$. Then $m_{\sigma}(u, h)=g h$ and we may choose as a basis for $T_{(u, h)}(U \times H)=T_{p_{o}}\left(\bar{\lambda}_{g}\right)(\mathfrak{g} / \mathfrak{h}) \times T_{\mathbf{1}}\left(\lambda_{h}\right)(\mathfrak{h})$

$$
T_{p_{o}}\left(\bar{\lambda}_{g h}\right) \bar{v}_{1}, \ldots, T_{p_{o}}\left(\bar{\lambda}_{g h}\right) \bar{v}_{k}, T_{\mathbf{1}}\left(\lambda_{h}\right) v_{k+1}, \ldots, T_{\mathbf{1}}\left(\lambda_{h}\right) v_{n},
$$

where $\bar{\lambda}_{g}$ is the (left) translation by $g$ on $G / H$ and $\bar{v}_{j}=v_{j}+\mathfrak{h}$. To calculate $m_{\sigma}^{*} \nu_{G}(u, h)$, we note first that

$$
T_{g h}\left(\lambda_{(g h)^{-1}}\right) T_{(u, h)}\left(m_{\sigma}\right) T_{p_{o}, \mathbf{1}}\left(\bar{\lambda}_{g h} \times \rho_{h}\right)=T_{p_{o}, \mathbf{1}}\left(\lambda_{(g h)^{-1}} \circ m_{\sigma} \circ\left(\bar{\lambda}_{g h} \times \lambda_{h}\right)\right)
$$

and

$$
\lambda_{(g h)^{-1}} \circ m_{\sigma} \circ\left(\bar{\lambda}_{g h} \times \lambda_{h}\right)\left(u^{\prime}, h^{\prime}\right)=(g h)^{-1} \sigma\left(g h . u^{\prime}\right) h h^{\prime} \in G .
$$

Since $g h . p_{o}=g . p_{o}=u$, this implies

$$
\left(p_{o}, \mathbf{1}\right) \mapsto c_{h^{-1}}\left(g^{-1} \sigma\left(g h . p_{o}\right)\right)=c_{h^{-1}}\left(g^{-1} \sigma(u)\right)=c_{h^{-1}}(\mathbf{1})=\mathbf{1} .
$$

We set

$$
\begin{array}{ll}
w_{j}:=T_{g h}\left(\lambda_{(g h)^{-1}}\right) T_{(u, h)}\left(m_{\sigma}\right) T_{p_{o}, \mathbf{1}}\left(\bar{\lambda}_{g h} \times \lambda_{h}\right)\left(\bar{v}_{j}, 0\right) \in \mathfrak{g}, & j=1, \ldots, k, \\
w_{j}:=T_{g h}\left(\lambda_{(g h)^{-1}}\right) T_{(u, h)}\left(m_{\sigma}\right) T_{p_{o}, \mathbf{1}}\left(\bar{\lambda}_{g h} \times \lambda_{h}\right)\left(0, v_{j}\right) \in \mathfrak{g}, & j=k+1, \ldots, n,
\end{array}
$$

and use

$$
q \circ \lambda_{(g h)^{-1}} \circ m_{\sigma} \circ\left(\bar{\lambda}_{g h} \times \lambda_{h}\right)\left(u^{\prime}, h^{\prime}\right)=q\left((g h)^{-1} \sigma\left(g h . u^{\prime}\right) h h^{\prime}\right)=u^{\prime} \in G / H
$$

to see that

$$
T_{\mathbf{1}}(q)\left(w_{j}\right)=v_{j}+\mathfrak{h}, \quad j=1, \ldots, k
$$

which implies the existence of $y_{j} \in \mathfrak{h}, j=1, \ldots, k$ such that

$$
w_{j}=v_{j}+y_{j}, \quad j=1, \ldots, k
$$

Moreover,\\
$w_{j}=\left.\frac{d}{d t}\right|_{t=0}(g h)^{-1} \sigma\left(g . p_{o}\right) h \exp t v_{j}=\left.\frac{d}{d t}\right|_{t=0} \exp t v_{j}=v_{j}, \quad j=k+1, \ldots, n$.\\
Thus the invariance of $\nu_{G}$ implies

$$
\begin{aligned}
& m_{\sigma}^{*} \nu_{G}(u, h)\left(T_{p_{o}}\left(\bar{\lambda}_{g h}\right) \bar{v}_{1}, \ldots, T_{p_{o}}\left(\bar{\lambda}_{g h}\right) \bar{v}_{k}, T_{\mathbf{1}}\left(\lambda_{h}\right) v_{k+1}, \ldots, T_{\mathbf{1}}\left(\lambda_{h}\right) v_{n}\right) \\
& =\nu_{G}(\mathbf{1})\left(w_{1}, \ldots, w_{n}\right)=\nu_{G}(\mathbf{1})\left(v_{1}+y_{1}, \ldots, v_{k}+y_{k}, v_{k+1}, \ldots, v_{n}\right) \\
& =\nu_{G}(\mathbf{1})\left(v_{1}, \ldots, v_{k}, v_{k+1}, \ldots, v_{n}\right)=1
\end{aligned}
$$

On the other hand, the invariance properties of $\mu \otimes \nu_{H}$ imply

$$
\begin{aligned}
& \left(\mu \otimes \nu_{H}\right)(u, h)\left(T_{p_{o}}\left(\bar{\lambda}_{g h}\right) \bar{v}_{1}, \ldots, T_{p_{o}}\left(\bar{\lambda}_{g h}\right) \bar{v}_{k}, T_{\mathbf{1}}\left(\lambda_{h}\right) v_{k+1}, \ldots, T_{\mathbf{1}}\left(\lambda_{h}\right) v_{n}\right) \\
& =\mu\left(p_{o}\right)\left(T_{u}\left(\bar{\lambda}_{g^{-1}}\right) T_{p_{o}}\left(\bar{\lambda}_{g h}\right) \bar{v}_{1}, \ldots, T_{u}\left(\bar{\lambda}_{g^{-1}}\right) T_{p_{o}}\left(\bar{\lambda}_{g h}\right) \bar{v}_{k}\right) \\
& \quad \cdot \nu_{H}(\mathbf{1})\left(T_{h}\left(\lambda_{h^{-1}}\right) T_{\mathbf{1}}\left(\lambda_{h}\right) v_{k+1}, \ldots, T_{h}\left(\lambda_{h^{-1}}\right) T_{\mathbf{1}}\left(\lambda_{h}\right) v_{n}\right) \\
& =\left(\mu \otimes \nu_{H}\right)\left(p_{o}, \mathbf{1}\right)\left(T_{p_{o}}\left(\bar{\lambda}_{h}\right) \bar{v}_{1}, \ldots, T_{p_{o}}\left(\bar{\lambda}_{h}\right) \bar{v}_{k}, v_{k+1}, \ldots, v_{n}\right) \\
& =\mu\left(p_{o}\right)\left(\overline{\operatorname{Ad}(h) v_{1}}, \ldots, \overline{\operatorname{Ad}(h) v_{k}}\right) \cdot \nu_{H}(\mathbf{1})\left(v_{k+1}, \ldots, v_{n}\right) \\
& =\mu\left(p_{o}\right)\left(\bar{v}_{1}, \ldots, \bar{v}_{k}\right) \cdot \nu_{H}(\mathbf{1})\left(v_{k+1}, \ldots, v_{n}\right)=1,
\end{aligned}
$$

and this proves the claim.\\
Proposition 10.4.13. Let $G$ be a Lie group and $A, B$ two integral subgroups with compact intersection such that $\mathfrak{g}=\mathfrak{a}+\mathfrak{b}$ and the multiplication induces an open map $\mu: A \times B \rightarrow G,(a, b) \mapsto a b$. If $G \backslash A B$ has Haar measure zero, then, for any two left Haar measures $\mu_{A}$ and $\mu_{B}$ on $A$ and $B$ we obtain a left Haar measure on $G$ by

$$
\int_{G} f(g) \mathrm{d} \mu_{G}(g)=\int_{B}\left|\frac{\operatorname{det} \operatorname{Ad}_{G}(b)}{\operatorname{det} \operatorname{Ad}_{B}(b)}\right| \int_{A} f(a b) \mathrm{d} \mu_{A}(a) \mathrm{d} \mu_{B}(b)
$$

for $f \in C_{c}(G)$.\\
Proof. The group $A \times B$ acts transitively on $A B$ via ( $a, b$ ) $\cdot x=a x b^{-1}$ and the stabilizer of $\mathbf{1}$ is the compact subgroup $K:=\{(a, a) \in A \times B: a \in A \cap B\}$. Therefore, the map

$$
\varphi:(A \times B) / K \rightarrow A B, \quad(a, b) K \mapsto a b^{-1}
$$

is a diffeomorphism, and for any left invariant density $\mu$ on $G$, we obtain the formula

$$
\int_{G} f \mathrm{~d} \widetilde{\mu}=\int_{(A \times B) / K}(f \circ \varphi) \mathrm{d}\left(\varphi^{*} \mu\right)^{\sim}
$$

Since $\varphi$ is a diffeomorphism onto an open subset of $G$ satisfying $\varphi \circ \lambda_{a}= \lambda_{a} \circ \varphi$ and $\varphi \circ \lambda_{b}=\rho_{b}^{-1} \circ \varphi$ for $a \in A$ and $b \in B$, the density $\varphi^{*} \mu$ on $(A \times B) / K$ is left $A$-invariant and satisfies

$$
\lambda_{b}^{*} \varphi^{*} \mu=\left(\varphi \circ \lambda_{b}\right)^{*} \mu=\left(\rho_{b}^{-1} \circ \varphi\right)^{*} \mu=\varphi^{*}\left(\rho_{b}^{-1}\right)^{*} \mu=\left|\operatorname{det}\left(\operatorname{Ad}_{G}(b)\right)\right| \cdot \varphi^{*} \mu
$$

(Proposition 10.4.7). Since the subgroup $A \cap B$ of $G$ is compact, $\Delta_{G}(A \cap B)$ is a compact subgroup of $\mathbb{R}_{+}^{\times}$, hence trivial, and therefore the function

$$
s:(A \times B) / K \rightarrow \mathbb{R}_{+}^{\times}, \quad s((a, b) K):=\left|\operatorname{det}\left(\operatorname{Ad}_{G}(b)\right)\right|
$$

is well defined and the preceding calculations imply that

$$
\omega:=s^{-1} \cdot \varphi^{*} \mu
$$

is an ( $A \times B$ )-invariant density on ( $A \times B$ ) / $K$. Since $K \cong A \cap B$ is compact, there exists a normalized Haar measure $\mu_{A \cap B}$, and, for a suitable normalization of Haar measure on $A \times B$, we then have (cf. Corollary 10.4.8 and Proposition 10.4.12)

$$
\begin{aligned}
& \int_{A \times B} f(a b) \mathrm{d} \mu_{A}(a) \mathrm{d} \mu_{B}(b) \\
& =\int_{A \times B} f\left(a b^{-1}\right)\left|\operatorname{det} \circ \operatorname{Ad}_{B}(b)\right| \mathrm{d} \mu_{A}(a) \mathrm{d} \mu_{B}(b) \\
& =\int_{(A \times B) / K} \int_{A \cap B} f\left((a h)(b h)^{-1}\right)\left|\operatorname{det} \circ \operatorname{Ad}_{B}(b)\right| \mathrm{d} \mu_{A \cap B}(h) \mathrm{d} \widetilde{\omega}((a, b) K) \\
& =\int_{(A \times B) / K} f\left(a b^{-1}\right)\left|\operatorname{det} \circ \operatorname{Ad}_{B}(b)\right| \mathrm{d} \widetilde{\omega}((a, b) K) \\
& =\int_{(A \times B) / K} f\left(a b^{-1}\right) \frac{\left|\operatorname{det}\left(\operatorname{Ad}_{B}(b)\right)\right|}{\left|\operatorname{det}\left(\operatorname{Ad}_{G}(b)\right)\right|} \mathrm{d}\left(\varphi^{*} \mu\right) \Upsilon((a, b) K)
\end{aligned}
$$

Comparing with (10.12), the assertion follows.

\subsubsection*{Averaging}
In this subsection, we give a few applications of invariant integration. We start with the construction of invariant inner products for representation spaces of compact Lie groups, which is a key step in the proof of Weyl's Trick which provides an equivalence between unitary representations of compact Lie groups and holomorphic representations of their complexifications.

Lemma 10.4.14 (Unitarity Lemma for Compact Groups). Let $\pi: G \rightarrow \mathrm{GL}(V)$ be a representation of the compact Lie group on the finitedimensional $\mathbb{K}$-vector space $V$. Then there exists a positive definite hermitian form $\beta$ on $V$ with $\pi(G) \subseteq \mathrm{U}(V, \beta)$.

Proof. Let $\mu_{G}$ be a normalized Haar measure on $G$. Using a basis, we obtain a linear isomorphism $\varphi: V \rightarrow \mathbb{K}^{n}$, so that we can use the standard hermitian form on $\mathbb{K}^{n}$ to obtain on $V$ a positive definite hermitian form $\widetilde{\beta}$. We now put

$$
\beta(v, w):=\int_{G} \widetilde{\beta}(\pi(g) v, \pi(g) w) d \mu_{G}(g) \quad \text { for } \quad x, y \in V
$$

This integral is defined because the integrand is a continuous function. It is clear that $\beta$ is sesquilinear and hermitian because

$$
\begin{aligned}
\overline{\beta(v, w)} & =\int_{G} \widetilde{\widetilde{\beta}}(\pi(g) v, \pi(g) w) \\
& =\int_{G} \widetilde{\beta}(\pi(g) w, \pi(g) v) d \mu_{G}(g)=\beta(w, v)
\end{aligned}
$$

For $v=w \neq 0$, we obtain $\beta(v, v)>0$ because $\widetilde{\beta}(\pi(g) v, \pi(g) v)>0$ holds for each $g \in G$. Finally, the right invariance of $\mu_{G}$ implies that

$$
\begin{aligned}
\beta(\pi(g) v, \pi(g) w) & =\int_{G} \widetilde{\beta}(\pi(h g) v, \pi(h g) w) d \mu_{G}(h) \\
& =\int_{G} \widetilde{\beta}(\pi(h) v, \pi(h) w) d \mu_{G}(h)=\beta(v, w)
\end{aligned}
$$

This proves that $\pi(G) \subseteq \mathrm{U}(V, \beta)$. \(\square\)

Next we turn to differential forms and de Rham cohomology classes.\\
Definition 10.4.15. Let $M$ be a smooth manifold and $G$ a Lie group acting smoothly on $M$. A differential form $\omega \in \Omega^{k}(M)$ is called G-invariant, if $\sigma_{g}^{*} \omega=\omega$ where $\sigma_{g}: M \rightarrow M$ is the diffeomorphism induced by $g$ on $M$ via the action, which we denote by $\sigma$. We write $\Omega^{k}(M)^{G}$ for the space of $G$-invariant $k$-forms on $M$.

Remark 10.4.16. If $\sigma: G \times M \rightarrow M$ is a smooth action, then Remark 10.3.16 implies that exterior differentiation commutes with the pullback maps $\sigma_{g}^{*}$, so that we obtain a map $\mathrm{d}: \Omega^{k}(M)^{G} \rightarrow \Omega^{k+1}(M)^{G}$. Thus we can define an invariant version of de Rham cohomology (cf. Definition 10.3.14) via

$$
H_{\mathrm{dR}, G}^{k}(M):=\frac{\operatorname{ker}\left(\mathrm{d}: \Omega^{k}(M)^{G} \rightarrow \Omega^{k+1}(M)^{G}\right)}{\operatorname{im}\left(\mathrm{d}: \Omega^{k-1}(M)^{G} \rightarrow \Omega^{k}(M)^{G}\right)}
$$

is called the $k$ th invariant de Rham cohomology space of $M$. It is clear that the inclusions $\Omega^{k}(M)^{G} \rightarrow \Omega^{k}(G)$ induce natural maps

$$
i_{*}: H_{\mathrm{dR}, G}^{k}(M) \rightarrow H_{\mathrm{dR}}^{k}(M) .
$$

Proposition 10.4.17. Let $\sigma: G \times M \rightarrow M$ be a smooth action of a compact Lie group $G$ on a smooth manifold $M$. Then

$$
\left(P_{k} \omega\right)_{x}\left(v_{1}, \ldots, v_{k}\right):=\int_{G}\left(\sigma_{g}^{*} \omega\right)_{x}\left(v_{1}, \ldots, v_{k}\right) \mathrm{d} \mu_{G}(g)
$$

for $x \in M, k \in \mathbb{N}_{0}$, and $v_{1}, \ldots, v_{k} \in T_{x}(M)$ defines linear projections $P_{k}: \Omega^{k}(M) \rightarrow \Omega^{k}(M)^{G}$ commuting with exterior differentiation.

Proof. The definition of the pull-back of differential form and the chain rule, together with the right invariance of the Haar measure, allow us to calculate

$$
\begin{aligned}
\left(\sigma_{h}^{*}\left(P_{k} \omega\right)\right)_{x}\left(v_{1}, \ldots, v_{k}\right) & =\left(P_{k} \omega\right)_{\sigma_{h}(x)}\left(T_{x}\left(\sigma_{h}\right) v_{1}, \ldots, T_{x}\left(\sigma_{h}\right) v_{k}\right) \\
& =\int_{G}\left(\sigma_{g}^{*} \omega\right)_{\sigma_{h}(x)}\left(T_{x}\left(\sigma_{h}\right) v_{1}, \ldots, T_{x}\left(\sigma_{h}\right) v_{k}\right) \mathrm{d} \mu_{G}(g) \\
& =\int_{G} \omega_{\sigma_{g h(x)}}\left(T_{x}\left(\sigma_{g h}\right) v_{1}, \ldots, T_{x}\left(\sigma_{g h}\right) v_{k}\right) \mathrm{d} \mu_{G}(g) \\
& =\int_{G} \omega_{\sigma_{g(x)}}\left(T_{x}\left(\sigma_{g}\right) v_{1}, \ldots, T_{x}\left(\sigma_{g}\right) v_{k}\right) \mathrm{d} \mu_{G}(g) \\
& =\left(P_{k} \omega\right)_{x}\left(v_{1}, \ldots, v_{k}\right) .
\end{aligned}
$$

This implies the first claim. The second claim follows by integration from the fact that the exterior derivative commutes with pullbacks (Remark 10.3.16).

The averaging operators $P_{k}$ can be used to show that the maps

$$
i_{*}: H_{\mathrm{dR}, G}^{k}(M) \rightarrow H_{\mathrm{dR}}^{k}(M)
$$

are injective:\\
Theorem 10.4.18. Let $M$ be a smooth manifold and $G$ a compact Lie group acting smoothly on $M$. Then the natural maps $i_{*}: H_{\mathrm{dR}, G}^{k}(M) \rightarrow H_{\mathrm{dR}}^{k}(M)$ are linear and injective.

Proof. Linearity is clear, so we assume that $\omega=\mathrm{d} \nu \in \Omega^{k}(M)^{G}$ with $\nu \in \Omega^{k-1}(M)$. Then $P_{k-1} \nu \in \Omega^{k-1}(M)^{G}$ and

$$
\mathrm{d}\left(P_{k-1} \nu\right)=P_{k}(\mathrm{~d} \nu)=P_{k} \omega=\omega
$$

Thus $\omega \in \operatorname{im}\left(\mathrm{d}: \Omega^{k-1}(M)^{G} \rightarrow \Omega^{k}(M)^{G}\right)$ induces the zero cohomology class in $H_{\mathrm{dR}}^{k}(M)^{G}$.

Remark 10.4.19. If, in the situation of Theorem 10.4.18, $G$ is connected, one can show that $i_{*}$ is indeed an isomorphism (see [GHV73, p. 151]).

Definition 10.4.20. Let $G$ be a Lie group with Lie algebra $\mathfrak{g}, V$ a vector space and $\rho: G \rightarrow \mathrm{GL}(V)$ be a smooth representation of $G$, so that the derived representation yields on $V$ the structure of a $\mathfrak{g}$-module.

We call a $p$-form $\omega \in \Omega^{p}(G, V)$ equivariant if we have for all $g \in G$ the relation

$$
\lambda_{g}^{*} \omega=\rho(g) \circ \omega
$$

If $V$ is a trivial module, then an equivariant $p$-form is a left invariant $V$-valued $p$-form on $G$. An equivariant $p$-form is uniquely determined by the corresponding element $\omega_{\mathbf{1}} \in C^{p}(\mathfrak{g}, V)$ via

$$
\omega_{g}\left(g x_{1}, \ldots, g x_{p}\right)=\rho(g) \omega_{1}\left(x_{1}, \ldots, x_{p}\right)
$$

for $g \in G, x_{i} \in \mathfrak{g} \cong T_{\mathbf{1}}(G)$, where $G \times T(G) \rightarrow T(G),(g, x) \mapsto g x$ denotes the natural action of $G$ on its tangent bundle $T(G)$ obtained by restricting the tangent map of the group multiplication.

Conversely, (10.13) provides for each $\omega \in C^{p}(\mathfrak{g}, V)$ a unique equivariant $p$-form $\omega^{\mathrm{eq}}$ on $G$ with $\omega_{1}^{\mathrm{eq}}=\omega$.

Remark 10.4.21. If $\omega \in C^{p}(\mathfrak{g}, V)$ and $\omega^{\mathrm{eq}} \in \Omega^{p}(G, V)$ is the corresponding left equivariant $p$-form on $G$, then each right multiplication $\rho_{g}: G \rightarrow G$ satisfies

$$
\begin{aligned}
\left(\rho^{*} \omega^{\mathrm{eq}}\right)_{\mathbf{1}}\left(x_{1}, \ldots, x_{p}\right) & =\omega^{\mathrm{eq}}\left(x_{1} g, \ldots, x_{p} g\right)=\rho(g) \omega\left(g^{-1} x_{1} g, \ldots, g^{-1} x_{p} g\right) \\
& =\rho(g)\left(\operatorname{Ad}\left(g^{-1}\right)^{*} \omega\right)\left(x_{1}, \ldots, x_{p}\right) .
\end{aligned}
$$

Since $\rho_{g}^{*} \omega^{\mathrm{eq}}$ is also left equivariant, we see that pullback with right multiplications correspond to the action of $G$ on $C^{p}(\mathfrak{g}, V)$ by $(g, \omega) \mapsto \operatorname{Ad}\left(g^{-1}\right)^{*} \omega$.

Proposition 10.4.22. For each $\omega \in C^{p}(\mathfrak{g}, V)$, the equation $\mathrm{d}\left(\omega^{\mathrm{eq}}\right)=\left(\mathrm{d}_{\mathfrak{g}} \omega\right)^{\mathrm{eq}}$ holds. Accordingly, the evaluation map

$$
\mathrm{ev}_{\mathbf{1}}: \Omega^{p}(G, V)^{\mathrm{eq}} \rightarrow C^{p}(\mathfrak{g}, V), \quad \omega \mapsto \omega_{\mathbf{1}} \quad \text { satisfies } \quad \mathrm{ev}_{\mathbf{1}} \circ \mathrm{d}=\mathrm{d}_{\mathfrak{g}} \circ \mathrm{ev}_{\mathbf{1}} .
$$

Proof. For $g \in G$, we have

$$
\lambda_{g}^{*} \mathrm{~d} \omega^{\mathrm{eq}}=\mathrm{d} \lambda_{g}^{*} \omega^{\mathrm{eq}}=\mathrm{d}\left(\rho(g) \circ \omega^{\mathrm{eq}}\right)=\rho(g) \circ\left(\mathrm{d} \omega^{\mathrm{eq}}\right),
$$

showing that $\mathrm{d} \omega^{\mathrm{eq}}$ is also equivariant.\\
For $x \in \mathfrak{g}$, we write $x_{l}$ for the corresponding left invariant vector field on $G$, i.e., $x_{l}(g)=g x$. It suffices to calculate the value of $\mathrm{d} \omega^{\mathrm{eq}}$ on ( $p+1$ )-tuples of left invariant vector fields in the identity element. In view of

$$
\omega^{\mathrm{eq}}\left(x_{1, l}, \ldots, x_{p, l}\right)(g)=\rho(g) \omega\left(x_{1}, \ldots, x_{p}\right),
$$

we obtain

$$
\left(x_{0, l} \omega^{\mathrm{eq}}\left(x_{1, l}, \ldots, x_{p, l}\right)\right)(\mathbf{1})=\mathbf{L}(\rho)\left(x_{0}\right) \omega\left(x_{1}, \ldots, x_{p}\right),
$$

and therefore,

$$
\begin{aligned}
& \left(\mathrm{d} \omega^{\mathrm{eq}}\left(x_{0, l}, \ldots, x_{p, l}\right)\right)(\mathbf{1}) \\
= & \sum_{i=0}^{p}(-1)^{i} \mathbf{L}(\rho)\left(x_{i, l}\right) \omega^{\mathrm{eq}}\left(x_{0, l}, \ldots, \widehat{x_{i, l}}, \ldots, x_{p, l}\right)(\mathbf{1}) \\
& +\sum_{i<j}(-1)^{i+j} \omega^{\mathrm{eq}}\left(\left[x_{i, l}, x_{j, l}\right], x_{0, l}, \ldots, \widehat{x_{i, l}}, \ldots, \widehat{x_{j, l}}, \ldots, x_{p, l}\right)(\mathbf{1}) \\
= & \sum_{i=0}^{p}(-1)^{i} \mathbf{L}(\rho)\left(x_{i}\right) \omega\left(x_{0}, \ldots, \widehat{x_{i}}, \ldots, x_{p}\right) \\
& +\sum_{i<j}(-1)^{i+j} \omega\left(\left[x_{i}, x_{j}\right], x_{0}, \ldots, \widehat{x_{i}}, \ldots, \widehat{x_{j}}, \ldots, x_{p}\right) \\
= & \left(\mathrm{d}_{\mathfrak{g}} \omega\right)\left(x_{0}, \ldots, x_{p}\right) .
\end{aligned}
$$

This proves our assertion. \(\square\)

Corollary 10.4.23. Every differential form $\omega \in \Omega^{p}(G, \mathbb{R})$ on a Lie group $G$ which is invariant under left and right translations is closed.

Proof. The right invariance of $\omega$ implies that $\omega_{\mathbf{1}} \in C^{p}(\mathfrak{g}, \mathbb{R})$ is invariant under the adjoint action, i.e.,

$$
\operatorname{Ad}(g)^{*} \omega_{\mathbf{1}}=\omega_{\mathbf{1}} \quad \text { for } \quad g \in G
$$

(Remark 10.4.21). Taking derivatives, we see that $\omega \in C^{p}(\mathfrak{g}, \mathbb{R})^{\mathfrak{g}}$, i.e., $\omega$ is $\mathfrak{g}$-invariant. Now Lemma 7.5.28 implies that $\mathrm{d}_{\mathfrak{g}} \omega_{\mathbf{1}}=0$, and Proposition 10.4.22 shows that $\omega^{\mathrm{eq}}$ is closed. \(\square\)

Corollary 10.4.24. Every non-zero differential form $\omega \in \Omega^{p}(G, \mathbb{R})$ on a compact Lie group $G$ which is invariant under left and right translations is closed, but not exact, so it defines a nonzero element in $H_{\mathrm{dR}}^{p}(G, \mathbb{R})$.

Proof. In view of Corollary 10.4.23, it only remains to show that an exact form $\omega \in \Omega^{p}(G, \mathbb{R})$ which is left and right invariant vanishes. To this end, pick such a form and let $G \times G$ act on $G$ by $\left(g_{1}, g_{2}\right) . x=g_{1} x g_{2}^{-1}$. This defines a smooth action for which $\omega$ is invariant, i.e., $\omega \in \Omega^{p}(G)^{G \times G}$. But then exactness by Theorem 10.4.18 implies that there exists a ( $G \times G$ )-invariant form $\nu \in \Omega^{p-1}(G)$ with $\omega=\mathrm{d} \nu$. On the other hand, Corollary 10.4.23 implies that $\nu$ is closed, so that $\omega=\mathrm{d} \nu=0$. \(\square\)

Corollary 10.4.25. If $G$ is a compact connected Lie group of dimension $n$, then $H_{\mathrm{dR}}^{n}(G, \mathbb{R}) \neq\{0\}$.

Proof. The compactness of the Lie algebra $\mathfrak{g}$ implies in particular that $G$ is unimodular, i.e., $|\operatorname{det}(\operatorname{Ad}(g))|=1$ for each $g \in G$. Taking derivatives, it follows that $\operatorname{tr}(\operatorname{ad} x)=0$ for each $x \in \mathfrak{g}$. Using Proposition 7.5.24, we find a $\mathfrak{g}$-invariant volume form $\omega \in C^{n}(\mathfrak{g}, \mathbb{R})$. As $G$ is connected, the $\mathfrak{g}$-invariance of $\omega$ implies that $\operatorname{Ad}(g)^{*} \omega=\omega$ holds for each $g \in G$. Then the corresponding left invariant differential form $\omega^{\mathrm{eq}} \in \Omega^{n}(G, \mathbb{R})$ is also invariant under right translations, and Corollary 10.4.24 shows that it defines a non-zero class in $H_{\mathrm{dR}}^{n}(G, \mathbb{R})$. \(\square\)

\subsubsection*{Exercises for Section 10.4}
Exercise 10.4.1. Let $G$ be a Lie group, $\mu_{G}$ a left Haar measure. For $f, h \in C_{c}(G)$, define the convolution product

$$
(f * h)(g):=\int_{G} f\left(g x^{-1}\right) h(x) d \mu_{G}(x) \quad \text { and } \quad\|f\|_{1}:=\int_{G}|f(x)| d \mu_{G}(x) .
$$

Show that:\\
(i) The convolution product turns $C_{c}(G)$ into an associative $\mathbb{R}$-algebra.\\
(ii) For $f, h \in C_{c}(G)$ we have $\|f * h\|_{1} \leq\|f\|_{1}\|h\|_{1}$, i.e., $\left(C_{c}(G), *,\|\cdot\|_{1}\right)$ is a normed algebra.\\
(iii) The completion of $C_{c}(G)$ with respect to the norm $\|\cdot\|_{1}$ is denoted by $L^{1}(G)$. It is possible to extend the convolution to an associative multiplication on $L^{1}(G)$, which turns $L^{1}(G)$ into a Banach algebra.\\
(iv) Write down the convolution for a finite group $G$.\\
(v) What is $L^{1}(G)$ for a discrete group, e.g., for $\mathbb{Z}$ ?

Exercise 10.4.2. Let $G$ be a locally compact group, and let $f \in C_{c}(G)$. Then $f$ is uniformly continuous in the sense that for every $\varepsilon>0$ there is a neighborhood $V \subseteq G$ of unity such that

$$
|f(x)-f(y)| \leq \varepsilon \quad \text { for } \quad y \in x V
$$

Exercise 10.4.3. Let $\mu$ be a Haar measure on the Lie group $G$ and $h \in C(G)$ with $\mu(f h)=0$ for all $f \in C_{c}(G)$. Show that $h=0$.

\subsection*{Integrating Lie Algebras of Vector Fields}
In Proposition 9.1.14, we obtained for each smooth action of a Lie group $G$ on a manifold $M$ a representation of the Lie algebra $\mathbf{L}(G)$ in the Lie algebra $\mathcal{V}(M)$ of vector fields on $M$. In this section, we study the converse problem of integrating such a Lie algebra representation to a smooth action.

\subsubsection*{Palais' Theorem}
We start with the case where all representation vector fields are known to be complete, a condition which is hard to verify.

Theorem 10.5.1 (Palais). Let $G$ be a 1-connected Lie group with Lie algebra $\mathfrak{g}, M$ a smooth manifold, and $\alpha: \mathfrak{g} \rightarrow \mathcal{V}(M)$ a morphism of Lie algebras whose image consists of complete vector fields. Then there exists a homomorphism $\beta: G \rightarrow \operatorname{Diff}(M)$ defining a smooth $G$-action $\sigma$ on $M$ with $\dot{\sigma}=\alpha$.

If $X \in \mathcal{V}(M)$ is a complete vector field, then we write

$$
\exp (X):=\Phi_{1}^{X} \in \operatorname{Diff}(M)
$$

for the time-1-flow of $X$.\\
Proposition 10.5.2. If $M$ is finite-dimensional and $E \subseteq \mathcal{V}(M)$ is a finitedimensional subspace consisting of complete vector fields, then the map

$$
E \times M \rightarrow M, \quad(X, m) \mapsto \exp (X)(m)
$$

is smooth.

Proof. Let $f: \mathbb{R}^{n} \rightarrow E$ be a linear isomorphism. Clearly, the map

$$
\mathbb{R}^{n} \times M \rightarrow T(M), \quad(x, m) \mapsto f(x)(m)
$$

is smooth. We thus obtain on the manifold $\widehat{M}:=\mathbb{R}^{n} \times M$ a smooth vector field $\widehat{X}(x, m):=(0, f(x)(m))$ with

$$
\exp (\widehat{X})(x, m)=(x, \exp (f(x))(m))
$$

so that the assertion follows from the smoothness of $\exp (\widehat{X})$ on $\widehat{M}$ (Theorem 8.5.12).

We have a natural action of $\operatorname{Diff}(M)$ on $\mathcal{V}(M)$ by

$$
\operatorname{Diff}(M) \times \mathcal{V}(M) \rightarrow \mathcal{V}(M), \quad(\varphi, X) \mapsto \varphi_{*} X:=T(\varphi) \circ X \circ \varphi^{-1}
$$

called the adjoint action.\\
Definition 10.5.3. Let $M$ and $N$ be smooth manifolds.\\
(a) Although $\operatorname{Diff}(M)$ is certainly not a finite-dimensional Lie group, we call a map $\varphi: N \rightarrow \operatorname{Diff}(M)$ smooth if the map

$$
\widehat{\varphi}: N \times M \rightarrow M \times M, \quad(n, x) \mapsto\left(\varphi(n)(x), \varphi(n)^{-1}(x)\right)
$$

is smooth. If $N$ is an interval in $\mathbb{R}$, we obtain in particular a notion of a smooth curve.\\
(b) Smooth curves $\varphi: J \subseteq \mathbb{R} \rightarrow \operatorname{Diff}(M)$ have (left) logarithmic derivatives

$$
\delta(\varphi): J \rightarrow \mathcal{V}(M), \quad \delta(\varphi)_{t}:=T\left(\varphi(t)^{-1}\right) \circ \varphi^{\prime}(t)
$$

which are curves in the Lie algebra $\mathcal{V}(M)$ of smooth vector fields on $M$, i.e., time-dependent vector fields. Here we consider $\varphi^{\prime}(t): M \rightarrow T M$ as the smooth map defined by

$$
\varphi^{\prime}(t)(m)=\frac{d}{d t} \varphi(t)(m), \quad m \in M, t \in J
$$

As we shall see below, for general $N$, the logarithmic derivatives are $\mathcal{V}(M)$ valued 1-forms on $N$ which are smooth in the sense that composition with any evaluation map $\mathrm{ev}_{p}: \mathcal{V}(M) \rightarrow T_{p}(M)$ yields a $T_{p}(M)$-valued differential form.

If $\varphi: N \rightarrow \operatorname{Diff}(M)$ is smooth and $\widehat{\varphi}_{1}: N \times M \rightarrow M,(n, x) \mapsto \varphi(n)(x)$, then we have a smooth tangent map

$$
T\left(\widehat{\varphi}_{1}\right): T(N \times M) \cong T(N) \times T(M) \rightarrow T(M)
$$

and obtain for each $v \in T_{p}(N)$ the partial map

$$
T_{p}(\varphi) v: M \rightarrow T(M), \quad m \mapsto T_{(p, m)}\left(\widehat{\varphi}_{1}\right)(v, 0)
$$

We define the (left) logarithmic derivative of $\varphi$ in $p$ by

$$
\delta(\varphi)_{p}: T_{p}(N) \rightarrow \mathcal{V}(M), \quad v \mapsto T\left(\varphi(p)^{-1}\right) \circ T_{p}(\varphi)(v) .
$$

Then $\delta(\varphi)$ is a $\mathcal{V}(M)$-valued 1-form on $N$.\\
Lemma 10.5.4. If $\mathfrak{g}$ is a finite-dimensional Lie algebra and $\alpha: \mathfrak{g} \rightarrow \mathcal{V}(M)$ is a homomorphism, then each $x \in \mathfrak{g}$, for which $\alpha(x)$ is complete, satisfies the relation

$$
\exp (-\alpha(x))_{*} \circ \alpha=\alpha \circ e^{\operatorname{ad} x} .
$$

Proof. We consider the smooth curve

$$
\psi(t):=\exp (t \alpha(x))_{*} \alpha\left(e^{t \operatorname{ad} x} y\right) .
$$

Then the definition of the Lie derivative implies that

$$
\begin{aligned}
\psi^{\prime}(t) & =\exp (t \alpha(x))_{*}\left(-\mathcal{L}_{\alpha(x)}\left(\alpha\left(e^{t \operatorname{ad} x} y\right)\right)\right)+\exp (t \alpha(x))_{*} \alpha\left(\left[x, e^{t \operatorname{ad} x} y\right]\right) \\
& =\exp (t \alpha(x))_{*}\left(-\left[\alpha(x), \alpha\left(e^{t \operatorname{ad} x} y\right)\right]\right)+\exp (t \alpha(x))_{*} \alpha\left(\left[x, e^{t \operatorname{ad} x} y\right]\right) \\
& =\exp (t \alpha(x))_{*} \alpha\left(\left[-x, e^{t \operatorname{ad} x} y\right]+\left[x, e^{t \operatorname{ad} x} y\right]\right)=0
\end{aligned}
$$

For $t=1$, we thus obtain $\exp (\alpha(x))_{*} \alpha\left(e^{\mathrm{ad} x} y\right)=\psi(1)=\psi(0)=\alpha(y)$, and the assertion follows. \(\square\)

For calculations, it is convenient to observe the Product- and Quotient Rules (cf. Lemma 9.2.27), both easy consequences of the Chain Rule:

Lemma 10.5.5. For two smooth maps $f, g: N \rightarrow \operatorname{Diff}(M)$, define

$$
(f g)(n):=f(n) \circ g(n) \quad \text { and } \quad\left(g^{-1}\right)(n):=g(n)^{-1} .
$$

Then the following assertions hold:\\
(i) Product Rule: $\delta(f g)=\delta(g)+\left(g^{-1}\right)_{*} \delta(f)$, and the\\
(ii) Quotient Rule: $\delta\left(f g^{-1}\right)=g_{*}(\delta(f)-\delta(g))$, where we write $\left(f_{*} \alpha\right)_{n}:=f(n)_{*} \circ \alpha_{n}$ for a $\mathcal{V}(M)$-valued 1 -form $\alpha$ on $N$.\\
(iii) If $\delta(f)=\delta(g)$, then there exists a $\varphi \in \operatorname{Diff}(M)$ with $f(n)=\varphi \circ g(n)$ for all $n \in N$.

Proof. (i), (ii) The smoothness of $f^{-1}$ and $g^{-1}$ follows directly from the definitions. For the smoothness of the product map $f g$, we observe that the maps

$$
N^{2} \times M \rightarrow M, \quad(t, s, x) \mapsto f(t)(g(s)(x))=(f(t) \circ g(s))(x)
$$

and

$$
N^{2} \times M \rightarrow M, \quad(t, s, x) \mapsto(f(t) \circ g(s))^{-1}(x)=g(s)^{-1}\left(f(t)^{-1}(x)\right)
$$

are smooth because they are compositions of smooth maps.\\
The Chain Rule now implies

$$
T_{p}(f g) v=T(f(p)) \circ\left(T_{p}(g) v\right)+\left(T_{p}(f) v\right) \circ g(p)
$$

which leads to

$$
\begin{aligned}
\delta(f g)_{p} & =T\left(g(p)^{-1}\right) T_{p}(g)+T\left(g(p)^{-1}\right) T\left(f(p)^{-1}\right) T_{p}(f) \circ g(p) \\
& =\delta(g)_{p}+\left(g(p)^{-1}\right)_{*} \delta(f)_{p}
\end{aligned}
$$

For $g=f^{-1}$, we obtain in particular

$$
0=\delta\left(f f^{-1}\right)=\delta\left(f^{-1}\right)+f_{*} \delta(f)
$$

Combining this with (i), we obtain the Quotient Rule.\\
(iii) From (ii), we immediately derive that the logarithmic derivative of $f g^{-1}$ vanishes, hence that it is constant, and this implies (iii).

Lemma 10.5.6. Let $X, Y \in \mathcal{V}(M)$ be two smooth vector fields with the property that $E:=\mathbb{R} X+\mathbb{R} Y$ consists of complete vector fields. Then $\operatorname{Exp}:=\left.\exp \right|_{E}$ is smooth and satisfies (in the pointwise sense)

$$
\delta(\operatorname{Exp})_{X}(Y)=\int_{0}^{1} \operatorname{Exp}(-t X)_{*} Y d t \quad \text { for } \quad X, Y \in E
$$

Proof. We claim that the smooth function

$$
\psi: I \rightarrow \mathcal{V}(M), \quad \psi(t):=\delta(\operatorname{Exp})_{t X}(t Y),
$$

satisfies the functional equation

$$
\psi(t+s)=\operatorname{Exp}(-s X)_{*} \psi(t)+\psi(s)
$$

We consider the three smooth functions $F, G, H: E \rightarrow \operatorname{Diff}(M)$, given by $F(X):=\operatorname{Exp}((t+s) X), G(X):=\operatorname{Exp}(t X)$, and $H(X):=\operatorname{Exp}(s X)$, satisfying $F=G \cdot H$ pointwise. The Product Rule (Lemma 10.5.5) implies that

$$
\delta(F)=\delta(H)+\left(H^{-1}\right)_{*} \delta(G)
$$

Now (10.15) follows from

$$
\delta(G)_{X}(Y)=\psi(t), \quad \delta(H)_{X}(Y)=\psi(s), \quad \text { and } \quad \delta(F)_{X}(Y)=\psi(s+t)
$$

Clearly, $\psi(0)=0$ and

$$
\psi^{\prime}(0)=\lim _{t \rightarrow 0} \delta(\operatorname{Exp})_{t X}(Y)=\delta(\operatorname{Exp})_{0}(Y)=Y
$$

Taking derivatives with respect to $t$ at 0 , (10.15) thus leads to

$$
\psi^{\prime}(s)=\operatorname{Exp}(-s X)_{*} Y
$$

which implies the assertion by integration:

$$
\delta(\operatorname{Exp})_{X}(Y)=\psi(1)=\int_{0}^{1} \psi^{\prime}(s) d s=\int_{0}^{1} \operatorname{Exp}(-s X)_{*} Y d s
$$

Proposition 10.5.7. Let $G$ be a Lie group with Lie algebra $\mathfrak{g}$ and $U \subseteq \mathfrak{g}$ a convex symmetric 0 -neighborhood such that $\left.\exp _{G}\right|_{U}$ is a diffeomorphism onto an open subset of $G$ and $V \subseteq U$ an open convex 0 -neighborhood with $\exp _{G}(V) \exp _{G}(V) \subseteq \exp _{G}(U)$. For the multiplication

$$
x * y:=\left(\left.\exp _{G}\right|_{U}\right)^{-1}\left(\exp _{G}(x) \exp _{G}(y)\right), \quad x, y \in V
$$

we then have in $\operatorname{Diff}(M)$

$$
\exp (-\alpha(x * y))=\exp (-\alpha(x)) \exp (-\alpha(y))
$$

for $x, y \in V$.\\
Proof. On $I:=[0,1]$, we define the two smooth functions

$$
\gamma_{1}: I \rightarrow \operatorname{Diff}(M), \quad \gamma_{1}(t):=\exp (-\alpha(x)) \exp (-t \alpha(y))
$$

and

$$
\gamma_{2}: I \rightarrow \operatorname{Diff}(M), \quad \gamma_{2}(t):=\exp (-\alpha(x * t y))
$$

Then $\gamma_{1}(0)=\gamma_{2}(0)=\exp (-\alpha(x))$ and $\delta\left(\gamma_{1}\right)(t)=-\alpha(y)$. For the smooth map $\exp \circ \alpha: \mathfrak{g} \rightarrow \operatorname{Diff}(M)$, we obtain with Lemmas 10.5.4 and 10.5.6

$$
\begin{aligned}
\delta(\exp \circ \alpha)_{x}(y) & =\int_{0}^{1} \exp (-t \alpha(x))_{*} \alpha(y) d t=\int_{0}^{1} \alpha\left(e^{t \operatorname{ad} x} y\right) d t \\
& =\alpha\left(\int_{0}^{1} e^{t \operatorname{ad} x} y d t\right)
\end{aligned}
$$

i.e.,

$$
\delta(\exp \circ \alpha)_{x}(y)=\alpha\left(\kappa_{\mathfrak{g}}(-x) y\right) \quad \text { for } \quad \kappa_{\mathfrak{g}}(x):=\int_{0}^{1} e^{-t \operatorname{ad} x} d t
$$

On the other hand, the relation $\exp _{G}(x * t y)=\exp _{G}(x) \exp _{G}(t y)$ leads with Proposition 9.2.29 to

$$
y=\delta\left(\exp _{G}\right)_{x * t y} \frac{d}{d t} x * t y=\kappa_{\mathfrak{g}}(x * t y) \frac{d}{d t} x * t y
$$

so that

$$
\frac{d}{d t} x * t y=\kappa_{\mathfrak{g}}(x * t y)^{-1} y
$$

We thus obtain

$$
\begin{aligned}
\delta\left(\gamma_{2}\right)(t) & =-\delta(\exp \circ \alpha)_{(-x * t y)} \alpha\left(\frac{d}{d t} x * t y\right)=-\alpha\left(\kappa_{\mathfrak{g}}(x * t y) \kappa_{\mathfrak{g}}(x * t y)^{-1} y\right) \\
& =-\alpha(y)
\end{aligned}
$$

This shows that $\gamma_{1}$ and $\gamma_{2}$ have the same logarithmic derivative, and since both curves start at the same point, they coincide (Lemma 10.5.5(iii)).

Proof. (of Palais' Theorem 10.5.1) Let $V$ be as in Proposition 10.5.7 and define

$$
f: \exp _{G}(V) \rightarrow \operatorname{Diff}(M), \quad f\left(\exp _{G}(x)\right):=\exp (-\alpha(x)) .
$$

Then

$$
f(g h)=f(g) \circ f(h) \quad \text { for } \quad g, h \in \exp _{G}(V)
$$

so that the Monodromy Principle (Proposition 9.5.8) implies the existence of an extension of $f$ to a homomorphism $G \rightarrow \operatorname{Diff}(M)$. Since $f$ is smooth on the identity neighborhood $\exp _{G}(V)$, it follows easily that $f$ defines a smooth action $\sigma$ of $G$ on $M$, and the definition of $f$ implies that $\mathbf{L}(\sigma)=\alpha$.

\subsubsection*{Lie Algebras Generated by Complete Vector Fields}
It is part of the assumptions of Theorem 10.5.1 that the Lie algebra $\alpha(\mathfrak{g})$ consists of complete vector fields. In this subsection, we show that this condition can be relaxed significantly to the requirement that $\alpha(\mathfrak{g})$ is merely generated by complete vector fields.

To this end, we consider a finite-dimensional Lie subalgebra $\mathfrak{g} \subseteq \mathcal{V}(M)$ and assume that $\mathfrak{g}$ is generated, as a Lie algebra, by complete vector fields. Let $C(\mathfrak{g})$ denote the subset of complete vector fields in $\mathfrak{g}$.

Lemma 10.5.8. If $X, Y \in C(\mathfrak{g})$, then $e^{\text {ad } X} Y \in C(\mathfrak{g})$.\\
Proof. We consider the smooth flow on $M$ defined by

$$
\gamma(t):=(\exp -X) \circ \exp (t Y) \circ(\exp X) \in \operatorname{Diff}(M)
$$

and note that

$$
\gamma^{\prime}(0)=(\exp -X)_{*} Y=e^{\operatorname{ad} X} Y
$$

by Lemma 10.5.4. Therefore, $e^{\text {ad } X} Y$ is a complete vector field.\\
Lemma 10.5.9. $\mathfrak{g}=\operatorname{span} C(\mathfrak{g})$.\\
Proof. Let $V:=\operatorname{span} C(\mathfrak{g}) \subseteq \mathfrak{g}$. By Lemma 10.5.8, we have $e^{\text {ad } X} V=V$ for each $X \in C(\mathfrak{g})$, hence $[C(\mathfrak{g}), V] \subseteq V$, so that $C(\mathfrak{g}) \subseteq \mathfrak{n}_{\mathfrak{g}}(V)$. Since $\mathfrak{n}_{\mathfrak{g}}(V)$ is a Lie subalgebra of $\mathfrak{g}$ and $C(\mathfrak{g})$ generates $\mathfrak{g}$, it follows that $[\mathfrak{g}, V] \subseteq V$, i.e., that $V$ is an ideal, hence in particular a subalgebra. Now the fact that $C(\mathfrak{g})$ generates $\mathfrak{g}$ yields $V=\mathfrak{g}$.

Proposition 10.5.10. If $\mathfrak{g} \subseteq \mathcal{V}(M)$ is a finite-dimensional Lie subalgebra generated by complete vector fields, then $\mathfrak{g}$ consists of complete vector fields.

Proof. In view of Lemma 10.5.9, $\mathfrak{g}$ is spanned by complete vector fields, so that there exists a basis $X_{1}, \ldots, X_{n}$, consisting of complete vector fields. Let $G$ be a connected Lie group with Lie algebra $\mathfrak{g}$ and $Y \in \mathfrak{g}$.

Using canonical coordinates of the second kind on $G$ with respect to the basis $X_{1}, \ldots, X_{n}$ (Lemma 9.2.6), we obtain smooth functions $\alpha_{i}$, defined on some 0 -neighborhood $[-\varepsilon, \varepsilon] \subseteq \mathbb{R}$, satisfying

$$
\exp _{G}(t Y)=\exp _{G}\left(\alpha_{1}(t) X_{1}\right) \cdots \exp _{G}\left(\alpha_{n}(t) X_{n}\right)
$$

for each $t \in[-\varepsilon, \varepsilon]$. We then have $\alpha_{i}(0)=0$ for each $i$ and $Y=\sum_{i=1}^{n} \alpha_{i}^{\prime}(0) X_{i}$.\\
Since the curve $\gamma(t):=\exp _{G}(t Y)$ in $G$ satisfies $\delta(\gamma)_{t}=Y$ for each $t \in \mathbb{R}$, the Product Rule for logarithmic derivatives (Lemma 9.2.27) yields

$$
\begin{aligned}
& Y=\delta\left(\exp _{G}\right)_{\alpha_{n}(t) X_{n}} \\
& +\sum_{i=1}^{n-1} \operatorname{Ad}\left(\exp _{G}\left(-\alpha_{n}(t) X_{n}\right)\right) \cdots \operatorname{Ad}\left(\exp _{G}\left(-\alpha_{i+1}(t) X_{i+1}\right)\right) \\
& \cdot \delta\left(\exp _{G}\right)_{\alpha_{i}(t) X_{i}} \alpha_{i}^{\prime}(t) X_{i}
\end{aligned}
$$

Since

$$
\delta\left(\exp _{G}\right)\left(\alpha_{i}(t) X_{i}\right) \alpha_{i}^{\prime}(t) X_{i}=\alpha_{i}^{\prime}(t) X_{i}
$$

follows from the Chain Rule, this simplifies to the relation

$$
Y=\sum_{i=1}^{n} \operatorname{Ad}\left(\exp _{G}\left(-\alpha_{n}(t) X_{n}\right)\right) \cdots \operatorname{Ad}\left(\exp _{G}\left(-\alpha_{i+1}(t) X_{i+1}\right)\right) \alpha_{i}^{\prime}(t) X_{i}
$$

We now consider the smooth curve

$$
\beta(t):=\exp \left(-\alpha_{1}(t) X_{1}\right) \cdots \exp \left(-\alpha_{n}(t) X_{n}\right) \in \operatorname{Diff}(M)
$$

and use the Product Formula from Lemma 10.5.5 to obtain

$$
\delta(\beta)_{t}=-\alpha_{n}^{\prime}(t) X_{n}-\sum_{i=1}^{n-1} \exp \left(\alpha_{n}(t) X_{n}\right)_{*} \cdots \exp \left(\alpha_{i+1}(t) X_{i+1}\right)_{*} \alpha_{i}^{\prime}(t) X_{i}
$$

Further, Lemma 10.5.4 implies that the inclusion homomorphism $\alpha: \mathfrak{g} \rightarrow \mathcal{V}(M)$ satisfies

$$
\begin{aligned}
Y & =\sum_{i=1}^{n} \alpha \circ \operatorname{Ad}\left(\exp _{G}\left(-\alpha_{n}(t) X_{n}\right)\right) \cdots \operatorname{Ad}\left(\exp _{G}\left(-\alpha_{i+1}(t) X_{i+1}\right)\right) \alpha_{i}^{\prime}(t) X_{i} \\
& =\sum_{i=1}^{n} \exp _{G}\left(\alpha_{n}(t) X_{n}\right)_{*} \cdots \exp \left(-\alpha_{i+1}(t) X_{i+1}\right)_{*} \alpha_{i}^{\prime}(t) X_{i}=-\delta(\beta)_{t}
\end{aligned}
$$

This shows that $\delta(\beta)_{t}=-Y$ for each $t \in[-\varepsilon, \varepsilon]$, and applying this to the two curves

$$
t \mapsto \beta(s) \circ \beta(t), \beta(s+t) \quad \text { for } \quad|t|,|s|,|t+s| \leq \varepsilon,
$$

we see that both have the same logarithmic derivative with respect to $t$ and the same values for $t=0$. Hence Lemma 10.5.5(iii) leads to

$$
\beta(s) \circ \beta(t)=\beta(s+t) \quad \text { for } \quad|t|,|s|,|t+s| \leq \varepsilon .
$$

We conclude that $\beta$ defines a local flow with generator $Y$ which is defined on $[-\varepsilon, \varepsilon] \times M$. This easily implies that $Y$ is complete (see Exercise 8.5.2).

\subsubsection*{Notes on Chapter 10}
Historically, the origin of smooth actions of Lie groups lies in Felix Klein's "Erlanger Programm" from $1872^{1}$, in which a geometry on a space $S$ was considered as the same structure as a group acting on this space. Then the geometric properties ${ }^{2}$ are those invariant under the action of the group.

Nowadays smooth actions of Lie groups on manifolds form a separate field of mathematics under the name of Lie transformation groups. In this chapter, we only collected what was needed in the further course of the book. This explains why there are quite different topics, none of which is explored systematically. One problem which occurs in this context is that smooth actions are needed to describe quotient structures and geometric properties of Lie groups, but, conversely, detailed knowledge of Lie groups is required for a systematic treatment of transformation groups. We refer to [Ko95] and [DK00] for more detailed information.

\footnotetext{${ }^{1}$ Christian Felix Klein (1849-1925) held the chair of geometry in Erlangen for a few years and the "Erlanger Programm" was his "Programmschrift", where he formulated his research plans when he came to Erlangen. Later he was a professor for mathematics in Munich, Leipzig, and eventually in Göttingen.\\
${ }^{2}$ A typical example is the group $\operatorname{Mot}\left(E_{2}\right)$ of motions (orientation preserving isometries) of the euclidian plane. The length of an interval or the area of a triangle are properties preserved by this group, hence geometric quantities. It was an important conceptual step to observe that changing the group means changing the geometry, resp., the notion of a geometric quantity. For example, the automorphism group Aff ( $A_{2}$ ) of the two-dimensional affine plane $A_{2}$ does not preserve the area of a triangle (it is larger than the euclidian group), hence the area of a triangle cannot be considered as an affine geometric quantity.
}


\pagebreak
\printbibliography
\end{document}